% Generated mechanically from frozen input p04/ko/morphisms.tex.
% Do not edit this generated file; edit the bound locale source instead.

% OK, start here.
%

\setcounter{chapter}{28}
\chapter{스킴의 사상}



\phantomsection
\label{morphisms-section-phantom}


\section{서론}
\label{morphisms-section-introduction}

\noindent
이 장에서는 스킴의 몇 가지 사상 유형을 도입한다.
기본 참고문헌은 \cite{EGA}이다.



\section{닫힌 몰입}
\label{morphisms-section-closed-immersions}

\noindent
이 절에서는 앞서 스킴의 닫힌 몰입에 관하여 얻은 결과 몇 가지를
설명한다. 스킴의 사상 $i : Z \to X$가 닫힌 몰입으로 정의된다는
것은 (a) $i$가 $X$의 닫힌 부분집합으로의 위상동형을 유도하고,
(b) $i^\sharp : \mathcal{O}_X \to i_*\mathcal{O}_Z$가 전사이며,
(c) $i^\sharp$의 핵이 절단들에 의해 국소적으로 생성된다는 뜻이다.
Schemes의 정의 \ref{schemes-definition-immersion} 및
\ref{schemes-definition-closed-immersion-locally-ringed-spaces}를 보라.
$Z$와 $X$가 스킴이라는 가정 아래에서는 닫힌 몰입을 특징짓는
서로 다른 방법이 많이 존재한다.

\begin{lemma}
\label{morphisms-lemma-closed-immersion}
스킴의 사상 $i : Z \to X$를 잡자.
다음 조건들은 서로 동치이다.
\begin{enumerate}
\item 사상 $i$는 닫힌 몰입이다.
\item 모든 아핀 열린집합 $\Spec(R) = U \subset X$에 대하여,
아이디얼 $I \subset R$가 존재하여
$i^{-1}(U) = \Spec(R/I)$가
$U = \Spec(R)$ 위의 스킴이 된다.
\item $X = \bigcup_{j \in J} U_j$인 아핀 열린 덮개가 존재하고,
$U_j = \Spec(R_j)$이며 각 $j \in J$에 대하여
아이디얼 $I_j \subset R_j$가 존재하여
$i^{-1}(U_j) = \Spec(R_j/I_j)$가
$U_j = \Spec(R_j)$ 위의 스킴이 된다.
\item 사상 $i$는 $Z$와 $X$의 닫힌 부분집합 사이의 위상동형을
유도하고 $i^\sharp : \mathcal{O}_X \to i_*\mathcal{O}_Z$는 전사이다.
\item 사상 $i$는 $Z$와 $X$의 닫힌 부분집합 사이의 위상동형을
유도하고, 사상 $i^\sharp : \mathcal{O}_X \to i_*\mathcal{O}_Z$는
전사이며, 핵 $\Ker(i^\sharp)\subset \mathcal{O}_X$는 준연접
아이디얼층이다.
\item 사상 $i$는 $Z$와 $X$의 닫힌 부분집합 사이의 위상동형을
유도하고, 사상 $i^\sharp : \mathcal{O}_X \to i_*\mathcal{O}_Z$는
전사이며, 핵 $\Ker(i^\sharp)\subset \mathcal{O}_X$는 절단들에 의해
국소적으로 생성되는 아이디얼층이다.
\end{enumerate}
\end{lemma}

\begin{proof}
조건 (6)은 닫힌 몰입의 정의이다. Schemes의 정의
\ref{schemes-definition-closed-immersion-locally-ringed-spaces} 및
\ref{schemes-definition-immersion}을 보라.
따라서 (6) $\Leftrightarrow$ (1)이다. Schemes의 보조정리
\ref{schemes-lemma-closed-subspace-scheme}에 의해
(1) $\Rightarrow$ (2)이다. (2) $\Rightarrow$ (3)은 자명하다.

\medskip\noindent
(3)을 가정하자. 각 사상
$\Spec(R_j/I_j) \to \Spec(R_j)$는 닫힌 몰입이다. 이는
Schemes의 예 \ref{schemes-example-closed-immersion-affines}를 보라.
따라서 $i^{-1}(U_j) \to U_j$는 그 상으로의 위상동형이고
$i^\sharp|_{U_j}$는 전사이다. 그러므로 $i$는 그 상으로의
위상동형이며 $i^\sharp$는 전사이다. 이는 국소적으로 확인할 수
있기 때문이다. 따라서 (3) $\Rightarrow$ (4)이다.

\medskip\noindent
(4) $\Rightarrow$ (1)은 Schemes의 보조정리
\ref{schemes-lemma-characterize-closed-immersions}이다.
(5) $\Rightarrow$ (6)은 자명하다.
또한 (6) $\Rightarrow$ (5)는 Schemes의 보조정리
\ref{schemes-lemma-closed-subspace-scheme}에서 따른다.
\end{proof}

\begin{lemma}
\label{morphisms-lemma-closed-immersion-ideals}
$X$를 스킴이라 하자. $i : Z \to X$와 $i' : Z' \to X$를
닫힌 몰입이라 하고, $\mathcal{I} = \Ker(i^\sharp)$와
$\mathcal{I}' = \Ker((i')^\sharp)$인 $\mathcal{O}_X$의 아이디얼층을
생각하자.
\begin{enumerate}
\item 사상 $i : Z \to X$가 $Z \to Z' \to X$로 분해되며
그 첫 번째 사상이 어떤 $a : Z \to Z'$인 것과
$\mathcal{I}' \subset \mathcal{I}$인 것은 동치이다.
이 경우 $a$는 닫힌 몰입이다.
\item $Z \cong Z'$가 $X$ 위에서 성립할 필요충분조건은
$\mathcal{I} = \mathcal{I}'$이다.
\end{enumerate}
\end{lemma}

\begin{proof}
이는 닫힌 부분공간에 관한 논의에서 따른다. 특히 Schemes의 절
\ref{schemes-section-closed-immersion} 및 Schemes의 보조정리
\ref{schemes-lemma-closed-immersion}와
\ref{schemes-lemma-characterize-closed-subspace}를 보라.
위의 보조정리 \ref{morphisms-lemma-closed-immersion}의 특징짓기 (3)에서도
곧바로 따른다.
\end{proof}

\begin{lemma}
\label{morphisms-lemma-closed-immersion-bijection-ideals}
$X$를 스킴이라 하자. $\mathcal{I} \subset \mathcal{O}_X$를
아이디얼층이라 하자. 다음 조건들은 서로 동치이다.
\begin{enumerate}
\item $\mathcal{I}$는 $\mathcal{O}_X$-가군층으로서 절단들에 의해
국소적으로 생성된다.
\item $\mathcal{I}$는 $\mathcal{O}_X$-가군층으로서 준연접이며,
\item 스킴의 닫힌 몰입 $i : Z \to X$가 존재하여 대응하는
아이디얼층 $\Ker(i^\sharp)$가 $\mathcal{I}$와 같다.
\end{enumerate}
\end{lemma}

\begin{proof}
(1)과 (2)의 동치는 Schemes의 보조정리
\ref{schemes-lemma-closed-subspace-scheme}에서 즉시 따른다.
(1)이 성립하면 Schemes의 정의
\ref{schemes-definition-closed-subspace}와 예
\ref{schemes-example-closed-subspace}에 의해
닫힌 부분공간 $i : Z \to X$가 존재하고
$\mathcal{I} = \Ker(i^\sharp)$이다. Schemes의 보조정리
\ref{schemes-lemma-closed-subspace-scheme}에 의해 이는 스킴의
닫힌 몰입이고 (3)이 성립한다.
반대로 (3)이 성립하면 Schemes의 보조정리
\ref{schemes-lemma-closed-subspace-scheme}에 의해 (2)가 성립한다.
(이 보조정리는 스킴의 닫힌 몰입이 국소 환 달린 공간의 닫힌
몰입이라는 사실 때문에 적용된다.)
\end{proof}

\begin{lemma}
\label{morphisms-lemma-base-change-closed-immersion}
닫힌 몰입의 밑변환은 닫힌 몰입이다.
\end{lemma}

\begin{proof}
Schemes의 보조정리
\ref{schemes-lemma-base-change-immersion}을 보라.
\end{proof}

\begin{lemma}
\label{morphisms-lemma-composition-closed-immersion}
닫힌 몰입들의 합성은 닫힌 몰입이다.
\end{lemma}

\begin{proof}
이는 Schemes의 보조정리
\ref{schemes-lemma-composition-immersion}에서 이미 보았다.
그러나 다른 증명을 제시하자. 즉, 보조정리
\ref{morphisms-lemma-closed-immersion}에서 닫힌 몰입의 특징짓기 (3)을
사용하면 따른다. 실제로 $I \subset R$가 아이디얼이고
$\overline{J} \subset R/I$가 아이디얼이면,
$\overline{J} = J/I$가 성립하는 아이디얼 $J \subset R$가
존재하고 이 아이디얼은 $I$를 포함하며
$(R/I)/\overline{J} = R/J$이다.
\end{proof}

\begin{lemma}
\label{morphisms-lemma-closed-immersion-quasi-compact}
닫힌 몰입은 준콤팩트이다.
\end{lemma}

\begin{proof}
이 보조정리는 Schemes의 보조정리
\ref{schemes-lemma-closed-immersion-quasi-compact}와 중복된다.
\end{proof}

\begin{lemma}
\label{morphisms-lemma-closed-immersion-separated}
닫힌 몰입은 분리이다.
\end{lemma}

\begin{proof}
이는 Schemes의 보조정리
\ref{schemes-lemma-immersions-monomorphisms}의 특수한 경우이다.
\end{proof}

\section{몰입}
\label{morphisms-section-immersions}

\noindent
이 절에서는 몰입에 관한 사실 몇 가지를 모은다.

\begin{lemma}
\label{morphisms-lemma-immersion-permanence}
스킴의 사상 $Z \to Y \to X$를 잡자.
\begin{enumerate}
\item $Z \to X$가 몰입이면 $Z \to Y$도 몰입이다.
\item $Z \to X$가 준콤팩트 몰입이고 $Y \to X$가 준분리이면,
$Z \to Y$는 준콤팩트 몰입이다.
\item $Z \to X$가 닫힌 몰입이고 $Y \to X$가 분리이면,
$Z \to Y$는 닫힌 몰입이다.
\end{enumerate}
\end{lemma}

\begin{proof}
각 경우의 증명은 다음 가환 다이어그램을 살펴보는 것이다.
$$
\xymatrix{
Z \ar[r] \ar[rd] & Y \times_X Z \ar[r] \ar[d] & Z \ar[d] \\
& Y \ar[r] & X
}
$$
위쪽 수평 화살표들의 합성은 항등사상이다.
(1)을 증명하자. 첫 번째 수평 화살표는
$Y \times_X Z \to Z$의 단면이므로 몰입이다. 이는 Schemes의
보조정리 \ref{schemes-lemma-section-immersion}에 따른다.
화살표 $Y \times_X Z \to Y$는 $Z \to X$의 밑변환이므로
몰입이다(Schemes의 보조정리
\ref{schemes-lemma-base-change-immersion}).
마지막으로 몰입들의 합성은 몰입이다
(Schemes의 보조정리
\ref{schemes-lemma-composition-immersion}). 이로써 (1)이 증명된다.
나머지 두 결과도 정확히 같은 방식으로 증명된다.
\end{proof}

\begin{lemma}
\label{morphisms-lemma-factor-quasi-compact-immersion}
몰입 $h : Z \to X$를 잡자.
$h$가 준콤팩트이면 $h = i \circ j$로 분해할 수 있으며, 여기서
$j : Z \to \overline{Z}$는 열린 몰입이고
$i : \overline{Z} \to X$는 닫힌 몰입이다.
\end{lemma}

\begin{proof}
$h$는 준콤팩트이고 준분리임을 유의하자(Schemes의 보조정리
\ref{schemes-lemma-immersions-monomorphisms}).
따라서 $h_*\mathcal{O}_Z$는 $\mathcal{O}_X$-가군의 준연접층이다.
이는 Schemes의 보조정리
\ref{schemes-lemma-push-forward-quasi-coherent}에 따른다.
그러므로
$\mathcal{I} = \Ker(\mathcal{O}_X \to h_*\mathcal{O}_Z)$는
준연접 아이디얼층이다. Schemes의 절
\ref{schemes-section-quasi-coherent}를 보라.
닫힌 부분스킴 $\overline{Z} \subset X$를 $\mathcal{I}$에 대응하는
것으로 잡자. 보조정리
\ref{morphisms-lemma-closed-immersion-bijection-ideals}를 보라.
Schemes의 보조정리 \ref{schemes-lemma-characterize-closed-subspace}에 의해
사상 $h$는
$h = i \circ j$로 분해되며, 여기서 $i : \overline{Z} \to X$는
포함 사상이다.
$j$가 열린 몰입임을 보이기 위해 $U \subset X$인 열린 부분스킴을
택하여 $h$가 $Z$를 $U$ 안으로 닫힌 몰입으로 유도하게 하자.
그러면 $\mathcal{I}|_U$는 닫힌 몰입 $Z \to U$에 대응하는
아이디얼층임이 분명하다.
따라서 $Z = \overline{Z} \cap U$임을 얻는다.
\end{proof}

\begin{lemma}
\label{morphisms-lemma-factor-reduced-immersion}
몰입 $h : Z \to X$를 잡자.
$Z$가 축약이면 $h = i \circ j$로 분해할 수 있으며, 여기서
$j : Z \to \overline{Z}$는 열린 몰입이고
$i : \overline{Z} \to X$는 닫힌 몰입이다.
\end{lemma}

\begin{proof}
$\overline{Z} \subset X$를 $h(Z)$의 폐포로 잡고, 축약 유도
닫힌 부분스킴 구조를 부여하자. Schemes의 정의
\ref{schemes-definition-reduced-induced-scheme}를 보라.
Schemes의 보조정리 \ref{schemes-lemma-map-into-reduction}에 의해
사상 $h$는 $h = i \circ j$로 분해되며,
$i : \overline{Z} \to X$는 포함 사상이고
$j : Z \to \overline{Z}$이다.
몰입의 정의로부터 $U \subset X$인 열린 부분스킴이 존재하여
$h$가 그 안의 닫힌 몰입을 통과하여 분해됨을 안다.
따라서 $U$에서 $\overline{Z} \cap U$와 $h(Z)$는 같은 기저
닫힌집합을 갖는 $U$의 축약 닫힌 부분스킴이다.
그러므로 Schemes의 보조정리
\ref{schemes-lemma-reduced-closed-subscheme}의 유일성에 의해
$h(Z) \cong \overline{Z} \cap U$임을 얻는다.
따라서 $j$는 $Z$와 $\overline{Z} \cap U$ 사이의 동형을 유도한다.
달리 말하면 $j$는 열린 몰입이다.
\end{proof}
\begin{example}
\label{morphisms-example-thibaut}
다음은 열린 몰입 뒤에 닫힌 몰입을 합성한 것이 아닌
몰입의 예이다.
$k$를 체라 하자.
$X = \Spec(k[x_1, x_2, x_3, \ldots])$라 하자.
$U = \bigcup_{n = 1}^{\infty} D(x_n)$라 하자.
그러면 $U \to X$는 열린 몰입이다.
아이디얼들을 생각하자.
$$
I_n =
(x_1^n, x_2^n, \ldots, x_{n - 1}^n, x_n - 1, x_{n + 1}, x_{n + 2}, \ldots)
\subset
k[x_1, x_2, x_3, \ldots][1/x_n].
$$
$I_n k[x_1, x_2, x_3, \ldots][1/x_nx_m] = (1)$임을 유의하자.
임의의 $m \not = n$에 대하여 성립한다.
따라서 준연접 아이디얼들 $\widetilde I_n$은
$D(x_n)$ 위에서 $D(x_nx_m)$ 위로 제한할 때 일치한다. 즉,
$\widetilde I_n|_{D(x_nx_m)} = \mathcal{O}_{D(x_n x_m)}$는
$n \not = m$일 때 성립한다.
따라서 이 아이디얼들은 준연접 아이디얼층
$\mathcal{I} \subset \mathcal{O}_U$로 붙는다.
닫힌 부분스킴 $Z \subset U$를 $\mathcal{I}$에 대응하는 것으로 잡자.
그러면 $Z \to X$는 몰입이다.

\medskip\noindent
$Z \to X$를
$Z \to \overline{Z} \to X$로 분해할 수 없다고 주장하자.
여기서 $\overline{Z} \to X$는 닫힌 몰입이고
$Z \to \overline{Z}$는 열린 몰입이다.
즉, $\overline{Z}$는 어떤 아이디얼
$I \subset k[x_1, x_2, x_3, \ldots]$에 의해 정의되어야 하고
$I_n = I k[x_1, x_2, x_3, \ldots][1/x_n]$이어야 한다.
하지만 $f \in k[x_1, x_2, x_3, \ldots]$ 중 모든 $I_n$에 들어가는
것은 $0$뿐이다!
따라서 $I$는 존재하지 않는다.
\end{example}
\begin{lemma}
\label{morphisms-lemma-check-immersion}
스킴의 사상 $f : Y \to X$를 잡자. 모든 $y \in Y$에 대하여
명제에 나타난 $f(y) \in U \subset X$인 열린 부분스킴이 존재하여
$f|_{f^{-1}(U)} : f^{-1}(U) \to U$가 몰입이면,
$f$는 몰입이다.
\end{lemma}

\begin{proof}
이 명제는 Schemes의 절
\ref{schemes-section-immersions} 끝에 있는 닫힌 부분스킴에 관한
논의에서 곧바로 따른다. 그러나 자세한 증명도 제시하자.
폐포 $Z \subset X$를 $f(Y)$의 것으로 잡자. 폐포를 취하는 연산은
열린집합으로 제한하는 것과 가환하므로, 가정에 의해
$f(Y) \subset Z$는 열린집합이다. 따라서
$Z' = Z \setminus f(Y)$는 닫힌집합이다.
그러므로 $X' = X \setminus Z'$는 $X$의 열린 부분스킴이고,
$f$는 포함사상에 이어지는 $f : Y \to X'$로 분해된다.
$y \in Y$이고 보조정리의 명제에 나타난 $U \subset X$가 주어졌다고 하자.
그러면 $U' = X' \cap U$는 $f'(y)$의 열린 근방이고
$(f')^{-1}(U') \to U'$는 몰입이다
(보조정리 \ref{morphisms-lemma-immersion-permanence}). 그 상은 닫혀 있다.
따라서 이는 닫힌 몰입이다. Schemes의 보조정리
\ref{schemes-lemma-immersion-when-closed}를 보라.
닫힌 몰입이라는 성질은 표적에 대해 국소적이므로
(예를 들어 보조정리 \ref{morphisms-lemma-closed-immersion}에 의해)
원하는 대로 $f'$가 닫힌 몰입임을 얻는다.
\end{proof}


\section{닫힌 몰입과 준연접층}
\label{morphisms-section-closed-immersions-quasi-coherent}

\noindent
다음 보조정리는 스킴 위의 준연접층에 대하여
Modules의 보조정리 \ref{modules-lemma-i-star-exact}가 아벨층에
대해 하는 일을 마무리한다. Modules의 절
\ref{modules-section-closed-immersion}에 있는 논의도 보라.

\begin{lemma}
\label{morphisms-lemma-i-star-equivalence}
스킴의 닫힌 몰입 $i : Z \to X$를 잡자.
준연접 아이디얼층 $\mathcal{I} \subset \mathcal{O}_X$를 잡자.
이는 $Z$를 정의한다.
함자
$$
i_* :
\QCoh(\mathcal{O}_Z)
\longrightarrow
\QCoh(\mathcal{O}_X)
$$
는 완전하고 충실하며, 본질적 상은 준연접
$\mathcal{O}_X$-가군 $\mathcal{G}$들 중
$\mathcal{I}\mathcal{G} = 0$인 것들이다.
\end{lemma}

\begin{proof}
닫힌 몰입은 준콤팩트이고 분리이다. 이는 보조정리
\ref{morphisms-lemma-closed-immersion-quasi-compact}와
\ref{morphisms-lemma-closed-immersion-separated}를 보라.
따라서 Schemes의 보조정리
\ref{schemes-lemma-push-forward-quasi-coherent}를 적용할 수 있고,
$Z$ 위의 준연접층의 직접상은 실제로 $X$ 위의 준연접층이다.

\medskip\noindent
Modules의 보조정리 \ref{modules-lemma-i-star-equivalence}에 의해
함자 $i_*$는 충실하다.

\medskip\noindent
이제 함자 $i_*$의 본질적 상을 설명하자.
예를 들어 Modules의 보조정리
\ref{modules-lemma-i-star-equivalence}에 의해
$\mathcal{I}(i_*\mathcal{F}) = 0$이다. 이는 임의의 준연접
$\mathcal{O}_Z$-가군에 대하여 성립한다.
다음으로 $\mathcal{G}$를 생각하자. 이는 임의의 준연접
$\mathcal{O}_X$-가군이고 $\mathcal{I}\mathcal{G} = 0$을 만족한다.
표준 사상
$$
\mathcal{G} \longrightarrow i_* i^*\mathcal{G}
$$
이 동형임을 보이면 충분하다\footnote{이는 더 일반적인 상황에서
Modules의 보조정리 \ref{modules-lemma-i-star-equivalence}의 증명으로
보였다.}.
스킴과 준연접 가군의 경우, $X$ 위에서 아핀 국소적으로 작업하고
보조정리 \ref{morphisms-lemma-closed-immersion} 및 Schemes의 보조정리
\ref{schemes-lemma-widetilde-pullback}을 사용하면 다음 대수적
명제를 증명하는 것으로 충분하다. $R$를 환, $I$를 아이디얼,
$R$-가군 $N$을 잡고 $IN = 0$이라고 하자. 표준 사상
$$
N \longrightarrow N \otimes_R R/I,\quad
n \longmapsto n \otimes 1
$$
은 $R$-가군의 동형이다. 이 간단한 대수 사실의 증명은 생략한다.
\end{proof}
\noindent
$i : Z \to X$를 닫힌 몰입이라 하자. 위 보조정리 때문에 표기법을
약간 남용하여 $\mathcal{F}$를 층 $i_*\mathcal{F}$라고 $X$ 위에서
쓰는 경우가 많다.

\begin{lemma}
\label{morphisms-lemma-largest-quasi-coherent-subsheaf}
$X$를 스킴이라 하자. $\mathcal{F}$를 준연접
$\mathcal{O}_X$-가군이라 하자. $\mathcal{G} \subset \mathcal{F}$를
$\mathcal{O}_X$-부분가군이라 하자. 다음 성질을 만족하는 유일한
준연접 $\mathcal{O}_X$-부분가군 $\mathcal{G}' \subset \mathcal{G}$가
존재한다. 모든 준연접 $\mathcal{O}_X$-가군 $\mathcal{H}$에 대하여
사상
$$
\Hom_{\mathcal{O}_X}(\mathcal{H}, \mathcal{G}')
\longrightarrow
\Hom_{\mathcal{O}_X}(\mathcal{H}, \mathcal{G})
$$
은 전단사이다. 특히 $\mathcal{G}'$는 $\mathcal{O}_X$-준연접
부분가군으로서 $\mathcal{F}$에 포함되고 $\mathcal{G}$에 포함된 것들 중
가장 큰 것이다.
\end{lemma}

\begin{proof}
$\mathcal{G}_a$, $a \in A$를 준연접 $\mathcal{O}_X$-부분가군 중
$\mathcal{G}$에 포함된 것들의 모임이라 하자.
그러면
$$
\bigoplus\nolimits_{a \in A} \mathcal{G}_a \longrightarrow \mathcal{F}
$$
의 상 $\mathcal{G}'$는 준연접이다. $X$ 위에서 준연접층 사상의 상은
준연접이고, 준연접층들의 직합도 준연접이기 때문이다. Schemes의 절
\ref{schemes-section-quasi-coherent}를 보라.
가군 $\mathcal{G}'$는 $\mathcal{G}$에 포함된다. 따라서 이는
$\mathcal{O}_X$-준연접 가군 중 $\mathcal{G}$에 포함된 가장 큰 것이다.

\medskip\noindent
공식을 증명하기 위해 준연접 $\mathcal{H}$인
$\mathcal{O}_X$-가군과 $\alpha : \mathcal{H} \to \mathcal{G}$인
$\mathcal{O}_X$-가군 사상을 잡자. 합성
$\mathcal{H} \to \mathcal{G} \to \mathcal{F}$의 상은 준연접층 사상의
상인 준연접층이다. 따라서 이는 $\mathcal{G}'$에 포함된다.
그러므로 $\alpha$는 $\mathcal{G}'$를 통하여 인수분해된다.
\end{proof}

\begin{lemma}
\label{morphisms-lemma-i-upper-shriek}
스킴의 닫힌 몰입 $i : Z \to X$를 잡자.
함자\footnote{이는 아마도 표준적이지 않은 표기이다.}
$i^! : \QCoh(\mathcal{O}_X) \to \QCoh(\mathcal{O}_Z)$가 존재하여
$i_*$의 오른쪽 수반이 된다. (Modules의 보조정리
\ref{modules-lemma-i-star-right-adjoint}와 비교하라.)
\end{lemma}

\begin{proof}
준연접 $\mathcal{O}_X$-가군 $\mathcal{G}$가 주어졌다고 하자.
$\mathcal{H}_Z(\mathcal{G})$를 $\mathcal{G}$의 부분층으로 생각하자.
이는 $\mathcal{I}$에 의해 소멸하는 국소 절단들로 이루어진다. 보조정리
\ref{morphisms-lemma-largest-quasi-coherent-subsheaf}에 의해 표준적인 최대
준연접 $\mathcal{O}_X$-부분가군 $\mathcal{H}_Z(\mathcal{G})'$가
존재한다. 구성에 의해
$$
\Hom_{\mathcal{O}_X}(i_*\mathcal{F}, \mathcal{H}_Z(\mathcal{G})')
=
\Hom_{\mathcal{O}_X}(i_*\mathcal{F}, \mathcal{G})
$$
는 임의의 준연접 $\mathcal{O}_Z$-가군 $\mathcal{F}$에 대하여 성립한다.
따라서
$i^!\mathcal{G} = i^*(\mathcal{H}_Z(\mathcal{G})')$로 둘 수 있다.
세부사항은 생략한다.
\end{proof}
\noindent
준연접 아이디얼층과 닫힌 부분스킴 사이의 $1$-대-$1$ 대응(보조정리
\ref{morphisms-lemma-closed-immersion-bijection-ideals}를 보라)을 사용하면
닫힌 부분스킴들의 스킴론적 교집합과 합집합을 정의할 수 있다.

\begin{definition}
\label{morphisms-definition-scheme-theoretic-intersection-union}
$X$를 스킴이라 하자. $Z, Y \subset X$를 준연접 아이디얼층
$\mathcal{I}, \mathcal{J} \subset \mathcal{O}_X$에 대응하는 닫힌
부분스킴이라 하자. $Z$와 $Y$의 {\it 스킴론적 교집합}은
$X$에서 $\mathcal{I} + \mathcal{J}$로 잘라낸 닫힌 부분스킴이다.
$Z$와 $Y$의 {\it 스킴론적 합집합}은
$X$에서 $\mathcal{I} \cap \mathcal{J}$로 잘라낸 닫힌 부분스킴이다.
\end{definition}

\begin{lemma}
\label{morphisms-lemma-scheme-theoretic-intersection}
$X$를 스킴이라 하자. $Z, Y \subset X$를 닫힌 부분스킴이라 하자.
$Z \cap Y$를 $Z$와 $Y$의 스킴론적 교집합이라 하자.
그러면 $Z \cap Y \to Z$와 $Z \cap Y \to Y$는 닫힌 몰입이고
$$
\xymatrix{
Z \cap Y \ar[r] \ar[d] & Z \ar[d] \\
Y \ar[r] & X
}
$$
는 스킴의 카르테지안 도식이다. 즉, $Z \cap Y = Z \times_X Y$이다.
\end{lemma}

\begin{proof}
$Z \cap Y \to Z$와 $Z \cap Y \to Y$는 보조정리
\ref{morphisms-lemma-closed-immersion-ideals}에 의해 닫힌 몰입이다.
$U = \Spec(A)$를 $X$의 아핀 열린집합이라 하고, $Z \cap U$와
$Y \cap U$가 아이디얼 $I \subset A$와 $J \subset A$에 대응한다고 하자.
그러면 $Z \cap Y \cap U$는 $I + J \subset A$에 대응한다.
$A/I \otimes_A A/J = A/(I + J)$이므로, Schemes의 절
\ref{schemes-section-fibre-products}에서 스킴의 파이버 곱을 기술한 것에
의해 이 도식은 카르테지안임을 알 수 있다.
\end{proof}

\begin{lemma}
\label{morphisms-lemma-scheme-theoretic-union}
$S$를 스킴이라 하자. $X, Y \subset S$를 닫힌 부분스킴이라 하자.
$X \cup Y$를 $X$와 $Y$의 스킴론적 합집합이라 하자.
$X \cap Y$를 $X$와 $Y$의 스킴론적 교집합이라 하자.
그러면 $X \to X \cup Y$와 $Y \to X \cup Y$는 닫힌 몰입이고,
다음 짧은 완전열이 존재한다.
$$
0 \to \mathcal{O}_{X \cup Y} \to \mathcal{O}_X \times \mathcal{O}_Y
\to \mathcal{O}_{X \cap Y} \to 0
$$
이는 $\mathcal{O}_S$-가군의 짧은 완전열이며, 다음 도식
$$
\xymatrix{
X \cap Y \ar[r] \ar[d] & X \ar[d] \\
Y \ar[r] & X \cup Y
}
$$
은 스킴의 범주에서 코카르테지안이다. 즉,
$X \cup Y = X \amalg_{X \cap Y} Y$이다.
\end{lemma}

\begin{proof}
$X \to X \cup Y$와 $Y \to X \cup Y$는 보조정리
\ref{morphisms-lemma-closed-immersion-ideals}에 의해 닫힌 몰입이다. 짧은 완전열에서는
보조정리 \ref{morphisms-lemma-i-star-equivalence}의 동치를 사용하여
$S$의 닫힌 부분스킴 위의 준연접 가군을 $S$ 위의 준연접 가군으로
생각한다. 열의 첫 사상에는 표준 사상
$\mathcal{O}_{X \cup Y} \to \mathcal{O}_X$와
$\mathcal{O}_{X \cup Y} \to \mathcal{O}_Y$를 사용하고,
둘째 사상에는 표준 사상 $\mathcal{O}_X \to \mathcal{O}_{X \cap Y}$와
표준 사상의 음인 $\mathcal{O}_Y \to \mathcal{O}_{X \cap Y}$를 사용한다.
그런 다음 완전성을 확인하기 위해 아핀 국소적으로 작업하면 된다.
$U = \Spec(A)$를 $S$의 아핀 열린집합이라 하고, $X \cap U$와
$Y \cap U$가 아이디얼 $I \subset A$와 $J \subset A$에 대응한다고 하자.
그러면 $(X \cup Y) \cap U$는 $I \cap J \subset A$에 대응하고
$X \cap Y \cap U$는 $I + J \subset A$에 대응한다.
따라서 완전성은 다음 열의 완전성에서 따른다.
$$
0 \to A/I \cap J \to A/I \times A/J \to A/(I + J) \to 0
$$
도식이 코카르테지안임을 보이기 위해 스킴 $T$와
스킴의 사상 $f : X \to T$, $g : Y \to T$가
$X \cap Y \to T$로서 일치한다고 하자. 유일한 사상
$h : X \cup Y \to T$가 존재하여 $f$와 $g$에 일치함을 보이는 것이
목표이다.
$h$를 구성하기 위해 $X \cup Y$에서 아핀 국소적으로 작업하면 된다.
Schemes의 절 \ref{schemes-section-glueing-schemes}를 보라.
$s \in X$, $s \not \in Y$이면 $X \to X \cup Y$는 $s$의 근방에서
동형이고 $h$를 구성하는 방법은 자명하다. $s \in Y$, $s \not \in X$인
경우도 마찬가지다. $s \in X \cap Y$에 대해서는
$V = \Spec(B) \subset T$를 $f(s) = g(s)$를 포함하는 아핀 열린집합으로
택할 수 있다. 그런 다음 $U = \Spec(A) \subset S$를 $s$를 포함하는
아핀 열린집합으로 택하여 $f(X \cap U)$와 $g(Y \cap U)$가 $V$에
포함되게 할 수 있다. $f|_{X \cap U}$와 $g|_{Y \cap V}$가 $V$로 가는
사상들은 다음 환 사상들에 대응한다.
$$
B \to A/I
\quad\text{and}\quad
B \to A/J
$$
이 사상들은 $A/(I + J)$로 가는 사상으로서 일치한다. 위에서 보인 짧은
완전열에 의해 이 환 준동형들의 유일한 리프트인
$B \to A/I \cap J$인 환 사상이 존재한다.
\end{proof}









\section{가군의 지지집합}
\label{morphisms-section-support}

\noindent
이 절에서는 스킴 위의 준연접 가군의 지지집합에 관한 몇 가지 기본 결과를
모은다.
가군층의 지지집합은 Modules의 절
\ref{modules-section-support}에서 정의했음을 상기하자.
한편 가군의 지지집합은 Algebra의 절
\ref{algebra-section-support}에서 정의했다.
이 둘은 일치한다.

\begin{lemma}
\label{morphisms-lemma-support-affine-open}
$X$를 스킴이라 하자. $\mathcal{F}$를 $X$ 위의 준연접층이라 하자.
$\Spec(A) = U \subset X$를 아핀 열린집합이라 하고
$M = \Gamma(U, \mathcal{F})$로 두자.
$x \in U$라 하고 $\mathfrak p \subset A$를 대응하는 소 아이디얼이라 하자.
다음 조건들은 동치이다.
\begin{enumerate}
\item $\mathfrak p$는 $M$의 지지집합에 속하고,
\item $x$는 $\mathcal{F}$의 지지집합에 속한다.
\end{enumerate}
\end{lemma}

\begin{proof}
이는 등식 $\mathcal{F}_x = M_{\mathfrak p}$와 정의들에서 따른다.
Schemes의 보조정리 \ref{schemes-lemma-spec-sheaves}를 보라.
\end{proof}
\begin{lemma}
\label{morphisms-lemma-support-closed-specialization}
$X$를 스킴이라 하자.
$\mathcal{F}$를 $X$ 위의 준연접층이라 하자.
$\mathcal{F}$의 지지집합은 특수화에 대해 닫혀 있다.
\end{lemma}

\begin{proof}
$x' \leadsto x$가 특수화이고 $\mathcal{F}_x = 0$이면
$\mathcal{F}_{x'}$는 0이다. $\mathcal{F}_{x'}$는
가군 $\mathcal{F}_x$의 국소화이기 때문이다. 따라서
$\text{Supp}(\mathcal{F})$의 여집합은 일반화에 대해 닫혀 있다.
\end{proof}

\noindent
유한형 준연접 가군의 지지집합은 닫혀 있고, 파이버에서 확인할 수 있으며,
기저변환과 가환한다.

\begin{lemma}
\label{morphisms-lemma-support-finite-type}
유한형 준연접 가군 $\mathcal{F}$를 스킴 $X$ 위에 잡자. 그러면
\begin{enumerate}
\item $\mathcal{F}$의 지지집합은 닫혀 있다.
\item $x \in X$에 대하여 다음이 성립한다.
$$
x \in \text{Supp}(\mathcal{F})
\Leftrightarrow
\mathcal{F}_x \not = 0
\Leftrightarrow
\mathcal{F}_x \otimes_{\mathcal{O}_{X, x}} \kappa(x) \not = 0.
$$
\item 스킴의 임의의 사상 $f : Y \to X$에 대하여 당김
$f^*\mathcal{F}$도 유한형이고
$\text{Supp}(f^*\mathcal{F}) = f^{-1}(\text{Supp}(\mathcal{F}))$이다.
\end{enumerate}
\end{lemma}

\begin{proof}
(1)은 보조정리
Modules, \ref{modules-lemma-support-finite-type-closed}의 재서술이다.
또한 보조정리 \ref{morphisms-lemma-support-affine-open},
Properties의 보조정리 \ref{properties-lemma-finite-type-module}, 및
Algebra의 보조정리 \ref{algebra-lemma-support-closed}를 결합해도
이를 볼 수 있다. (2)의 첫 번째 동치는 지지집합의 정의이고,
두 번째 동치는 나카야마 보조정리에서 따른다. Algebra의 보조정리
\ref{algebra-lemma-NAK}를 보라.
$f : Y \to X$를 스킴의 사상이라 하자. 보조정리
Modules, \ref{modules-lemma-pullback-finite-type}에 의해
$f^*\mathcal{F}$는 유한형임에 유의하자.
마지막 주장을 위해 $y \in Y$를 잡고 그 상을 $x \in X$라 하자.
다음을 상기하자.
$$
(f^*\mathcal{F})_y =
\mathcal{F}_x \otimes_{\mathcal{O}_{X, x}} \mathcal{O}_{Y, y},
$$
Sheaves의 보조정리 \ref{sheaves-lemma-stalk-pullback-modules}를 보라.
따라서 $(f^*\mathcal{F})_y \otimes \kappa(y)$가 0이 아닌 것은
$\mathcal{F}_x \otimes \kappa(x)$가 0이 아닌 것과 동치이다.
(2)에 의해 이는 $x \in \text{Supp}(\mathcal{F})$인 것과
$y \in \text{Supp}(f^*\mathcal{F})$인 것이 동치임을 뜻하며,
이는 (3)의 내용이다.
\end{proof}

\begin{lemma}
\label{morphisms-lemma-scheme-theoretic-support}
유한형 준연접 가군 $\mathcal{F}$를 스킴 $X$ 위에 잡자. $i : Z \to X$가
가장 작은 닫힌 부분스킴으로서, 어떤 준연접
$\mathcal{O}_Z$-가군 $\mathcal{G}$에 대해
$i_*\mathcal{G} \cong \mathcal{F}$가 되도록 하는 것이 존재한다. 또한:
\begin{enumerate}
\item $\Spec(A) \subset X$가 임의의 아핀 열린집합이고
$\mathcal{F}|_{\Spec(A)} = \widetilde{M}$이면
$Z \cap \Spec(A) = \Spec(A/I)$이며 $I = \text{Ann}_A(M)$이다.
\item 준연접층 $\mathcal{G}$는 유일한 동형사상까지 유일하다.
\item 준연접층 $\mathcal{G}$는 유한형이다.
\item $\mathcal{G}$와 $\mathcal{F}$의 지지집합은 $Z$이다.
\end{enumerate}
\end{lemma}

\begin{proof}
$i' : Z' \to X$가 이 보조정리의 열린 아핀에 대한 기술을 만족하는
닫힌 부분스킴이라고 하자. 그러면 보조정리
\ref{morphisms-lemma-i-star-equivalence}에 의해 어떤 유일한 준연접층이 존재하여
$\mathcal{F} \cong i'_*\mathcal{G}'$이고 그 층 $\mathcal{G}'$는
$Z'$ 위에 있음을 알 수 있다.
더욱이 (같은 보조정리에 의해) $Z'$가 이 성질을 갖는 가장 작은 닫힌
부분스킴임은 자명하다. 마지막으로 Properties의 보조정리
\ref{properties-lemma-finite-type-module}와 Algebra의 보조정리
\ref{algebra-lemma-finite-over-subring}를 사용하면
$\mathcal{G}'$가 유한형임을 얻는다. 또한
$\text{Supp}(\mathcal{G}') = Z$이다. Algebra의 보조정리
\ref{algebra-lemma-support-closed}를 보라.
따라서 보조정리를 증명하려면 (1)의 특성화가 실제로 닫힌 부분스킴을
정의한다는 것을 보이면 충분하다.
이를 위해서는 주어진 규칙이 준연접 아이디얼층을 만든다는 것을
증명하면 충분하다. 보조정리
\ref{morphisms-lemma-closed-immersion-bijection-ideals}를 보라.
이는 다음 대수적 사실로 귀결된다. $A$가 환이고 $f \in A$이며
$M$이 유한 $A$-가군이면
$\text{Ann}_A(M)_f = \text{Ann}_{A_f}(M_f)$이다.
증명은 생략한다.
\end{proof}

\begin{definition}
\label{morphisms-definition-scheme-theoretic-support}
$X$를 스킴이라 하자. 유한형 준연접 가군 $\mathcal{F}$를
$\mathcal{O}_X$-가군으로 잡자. $\mathcal{F}$의
{\it 스킴론적 지지집합}은 보조정리
\ref{morphisms-lemma-scheme-theoretic-support}에서 구성한 닫힌 부분스킴
$Z \subset X$이다.
\end{definition}

\noindent
이 상황에서 우리는 보조정리
\ref{morphisms-lemma-i-star-equivalence}의 범주 동치를 통해
$\mathcal{F}$를 $Z$ 위의 유한형 준연접층으로 생각하는 경우가 많다.

\section{스킴론적 상}
\label{morphisms-section-scheme-theoretic-image}

\noindent
주의: 이 절의 일부 내용은 매우 일반적이며 여러분이 예상하는 것과
다르게 작동한다.

\begin{lemma}
\label{morphisms-lemma-scheme-theoretic-image}
$f : X \to Y$를 스킴의 사상이라 하자. $Z \subset Y$인 닫힌
부분스킴이 존재하여 $f$가 $Z$를 통하여 인수분해되며, 임의의 다른 닫힌 부분스킴
$Z' \subset Y$에 대하여 $f$가 $Z'$를 통하여 인수분해되면
$Z \subset Z'$이다.
\end{lemma}

\begin{proof}
$\mathcal{I} = \Ker(\mathcal{O}_Y \to f_*\mathcal{O}_X)$로 두자.
$\mathcal{I}$가 준연접이면 $Z$를 $\mathcal{I}$가 정하는 닫힌
부분스킴으로 택하면 된다. 보조정리
\ref{morphisms-lemma-closed-immersion-bijection-ideals}를 보라. 이는 보조정리
\ref{morphisms-lemma-closed-immersion-ideals}에 의해 성립한다.
일반적으로 같은 보조정리를 적용하려면 $\mathcal{I}'$가
$\mathcal{I}$에 포함된 가장 큰 준연접 아이디얼층으로 존재함을
보여야 한다.
이는 보조정리 \ref{morphisms-lemma-largest-quasi-coherent-subsheaf}에서 따른다.
\end{proof}

\begin{definition}
\label{morphisms-definition-scheme-theoretic-image}
$f : X \to Y$의 {\it 스킴론적 상}, 즉 사상 $f$의 상은 가장 작은
닫힌 부분스킴 $Z \subset Y$로서 $f$가 이를 통하여 인수분해되는 것이다.
위 보조정리
\ref{morphisms-lemma-scheme-theoretic-image}를 보라.
\end{definition}
\noindent
스킴의 사상 $f : X \to Y$의 스킴론적 상이 $Z$일 때,
우리는 흔히 $f : X \to Z$를 $f$의 스킴론적 상을 통한
인수분해라고 쓴다. 사상 $f$가 준콤팩트가 아니면 (일반적으로)
\begin{enumerate}
\item 점들의 집합으로서의 포함 $\overline{f(X)} \subset Z$는
등식이 아니며, 즉 $f(X) \subset Z$는 조밀한 부분집합이 아니고,
\item 스킴론적 상의 구성은 $Y$의 열린 부분스킴으로의 제한과
가환하지 않는다.
\end{enumerate}
Examples의 절 \ref{examples-section-scheme-theoretic-image}에서
독자는 두 현상 모두의 예를 볼 수 있다.
이 현상들은 몰입에 대해서도 일어날 수 있다. Examples의 절
\ref{examples-section-strange-immersion}을 보라.
그러나 다음 보조정리는 사상이 준콤팩트일 때 이 두 재난이 피할 수
있음을 보여준다.

\begin{lemma}
\label{morphisms-lemma-quasi-compact-scheme-theoretic-image}
$f : X \to Y$를 스킴의 사상이라 하자.
$Z \subset Y$를 $f$의 스킴론적 상이라 하자.
$f$가 준콤팩트이면
\begin{enumerate}
\item 아이디얼층
$\mathcal{I} = \Ker(\mathcal{O}_Y \to f_*\mathcal{O}_X)$는
준연접이다.
\item 스킴론적 상 $Z$는 $\mathcal{I}$가 정하는 닫힌 부분스킴이다.
\item 임의의 열린집합 $U \subset Y$에 대하여
$f|_{f^{-1}(U)} : f^{-1}(U) \to U$의 스킴론적 상은
$Z \cap U$와 같다.
\item 상 $f(X) \subset Z$는 $Z$의 조밀한 부분집합이다. 즉,
사상 $X \to Z$는 우세하다
(정의 \ref{morphisms-definition-dominant}를 보라).
\end{enumerate}
\end{lemma}

\begin{proof}
(4)는 (3)에서 따른다. (3)을 보이려면 $\mathcal{I}$의 구성은
$Y$의 열린 부분스킴으로의 제한과 가환하므로 (1)을 증명하면 충분하다.
그리고 (1)이 성립하면 보조정리
\ref{morphisms-lemma-scheme-theoretic-image}의 증명에서 (2)를 보였다.
따라서 $\mathcal{I}$가 준연접임을 증명하면 충분하다.
준연접이라는 성질은 국소적이므로 $Y$가 아핀이라고 가정할 수 있다.
$f$가 준콤팩트이므로 유한한 아핀 열린 덮개
$X = \bigcup_{i = 1, \ldots, n} U_i$를 잡을 수 있다.
$f'$를 다음 합성으로 쓰자.
$$
X' = \coprod U_i \longrightarrow X \longrightarrow Y.
$$
그러면 $f_*\mathcal{O}_X$는 $f'_*\mathcal{O}_{X'}$의 부분층이고,
따라서 $\mathcal{I} = \Ker(\mathcal{O}_Y \to f'_*\mathcal{O}_{X'})$이다.
Schemes의 보조정리
\ref{schemes-lemma-push-forward-quasi-coherent}에 의해
층 $f'_*\mathcal{O}_{X'}$는 $Y$ 위에서 준연접이다. 따라서 끝난다.
\end{proof}
\begin{example}
\label{morphisms-example-scheme-theoretic-image}
$A \to B$가 핵이 $I$인 환 사상이면, $\Spec(B) \to \Spec(A)$의
스킴론적 상은 닫힌 부분스킴 $\Spec(A/I)$이며 이는 $\Spec(A)$ 안에
놓인다. 이는
보조정리 \ref{morphisms-lemma-quasi-compact-scheme-theoretic-image}에서 따른다.
\end{example}

\noindent
사상이 준콤팩트이면 스킴론적 상은 상에 있는 점들의 특수화인 점들만
추가한다.

\begin{lemma}
\label{morphisms-lemma-reach-points-scheme-theoretic-image}
$f : X \to Y$를 준콤팩트 사상이라 하자.
$Z$를 $f$의 스킴론적 상이라 하자.
$z \in Z$\footnote{보조정리
\ref{morphisms-lemma-quasi-compact-scheme-theoretic-image}에 의해 집합으로서
$Z$는 $f(X)$의 $Y$ 안에서의 폐포와 일치한다.}라 하자.
값매김환 $A$와 그 분수체 $K$, 그리고 다음 가환 도식이 존재한다.
$$
\xymatrix{
\Spec(K) \ar[rr] \ar[d] & & X \ar[d] \ar[ld] \\
\Spec(A) \ar[r] & Z \ar[r] & Y
}
$$
그러한 $\Spec(A)$의 닫힌점은 $z$로 사상된다. 특히 $Z$의 모든 점은
$f(X)$의 한 점의 특수화이다.
\end{lemma}

\begin{proof}
$z \in \Spec(R) = V \subset Y$를 $z$의 아핀 열린 근방이라 하자. 보조정리
\ref{morphisms-lemma-quasi-compact-scheme-theoretic-image}에 의해
$Z \cap V$는 $f^{-1}(V) \to V$의 스킴론적 상이다. 따라서 $Y$를
$V$로 바꾸고 $Y = \Spec(R)$가 아핀이라고 가정할 수 있다.
이 경우 $X$는 준콤팩트이다. $f$가 준콤팩트이기 때문이다.
$X = U_1 \cup \ldots \cup U_n$이 유한 아핀 열린 덮개라고 하자.
$U_i = \Spec(A_i)$로 쓰자.
$I = \Ker(R \to A_1 \times \ldots \times A_n)$로 두자.
보조정리 \ref{morphisms-lemma-quasi-compact-scheme-theoretic-image}를 다시
적용하면 $Z$는 닫힌 부분스킴 $\Spec(R/I)$에 대응하며 이는 $Y$ 안에
놓인다.
$\mathfrak p \subset R$를 $z$에 대응하는 소 아이디얼이라 하자.
그러면 $I_{\mathfrak p} \subset R_{\mathfrak p}$는 등식이 아니다.
국소화가 완전하므로 (Algebra의 명제
\ref{algebra-proposition-localization-exact}를 보라), 따라서
$R_{\mathfrak p} \to
(A_1)_{\mathfrak p} \times \ldots \times (A_n)_{\mathfrak p}$
는 영사상이 아니다. 따라서 어떤 환 $(A_i)_{\mathfrak p}$는 영환이 아니다.
그러므로 어떤 $i$와 소 아이디얼 $\mathfrak q_i \subset A_i$가 존재하여
$\mathfrak p_i \subset \mathfrak p$ 위에 놓인다.
Algebra의 보조정리 \ref{algebra-lemma-dominate}에 의해
$A \subset K = \kappa(\mathfrak q_i)$인 값매김환을 택할 수 있고, 이는
국소환
$R_{\mathfrak p}/\mathfrak p_iR_{\mathfrak p} \subset \kappa(\mathfrak q_i)$를
지배한다.
이것이 원하는 도식을 준다. 세부사항은 생략한다.
\end{proof}

\begin{lemma}
\label{morphisms-lemma-factor-factor}
다음과 같은 스킴의 가환 도식이 있다고 하자.
$$
\xymatrix{
X_1 \ar[d] \ar[r]_{f_1} & Y_1 \ar[d] \\
X_2 \ar[r]^{f_2} & Y_2
}
$$
$Z_i \subset Y_i$, $i = 1, 2$를 $f_i$의 스킴론적 상이라 하자.
그러면 $Y_1 \to Y_2$는 $Z_1 \to Z_2$를 유도하고 다음 가환 도식을 준다.
$$
\xymatrix{
X_1 \ar[r] \ar[d] & Z_1 \ar[d] \ar[r] & Y_1 \ar[d] \\
X_2 \ar[r] & Z_2 \ar[r] & Y_2
}
$$
\end{lemma}

\begin{proof}
$Z_2$의 $Y_1$에서의 스킴론적 역상이 $Y_1$의 닫힌 부분스킴이며,
$f_1$은 이를 통하여 인수분해된다. 따라서 $Z_1$은 여기에 포함된다.
이것으로 보조정리가 증명된다.
\end{proof}

\begin{lemma}
\label{morphisms-lemma-scheme-theoretic-image-reduced}
$f : X \to Y$를 스킴의 사상이라 하자.
$X$가 축약이면 $f$의 스킴론적 상은
$\overline{f(X)}$에 유도된 축약 스킴 구조이다.
\end{lemma}

\begin{proof}
$\overline{f(X)}$에 유도된 축약 스킴 구조가
$Y$의 닫힌 부분스킴 중 $f$가 인수분해되는 가장 작은 것이기 때문이다.
Schemes의 보조정리 \ref{schemes-lemma-map-into-reduction}를 보라.
\end{proof}

\begin{lemma}
\label{morphisms-lemma-scheme-theoretic-image-of-partial-section}
$f : X \to Y$를 스킴의 분리 사상이라 하자.
$V \subset Y$를 소콤팩트 열린집합이라 하자. $s : V \to X$를
$f \circ s = \text{id}_V$를 만족하는 사상이라 하자.
$Y'$를 $s$의 스킴론적 상이라 하자.
그러면 $Y' \to Y$는 $V$ 위에서 동형이다.
\end{lemma}

\begin{proof}
$V$가 $Y$에서 소콤팩트라는 가정은
(Topology의 정의 \ref{topology-definition-quasi-compact}에 의해)
$V \to Y$가 준콤팩트 사상이라는 뜻이다.
Schemes의 보조정리 \ref{schemes-lemma-quasi-compact-permanence}에 의해
사상 $s : V \to X$는 준콤팩트이다.
따라서 스킴론적 상 $Y'$의 구성은 $s$에 대해 열린집합으로의 제한과
가환한다. 이는 보조정리
\ref{morphisms-lemma-quasi-compact-scheme-theoretic-image}에 따른다.
특히 $Y' \cap f^{-1}(V)$는 분리 사상
$f^{-1}(V) \to V$의 단면의 스킴론적 상이다.
분리 사상의 단면은 닫힌 몰입이므로 (Schemes의 보조정리
\ref{schemes-lemma-section-immersion}를 보라), 원하는 대로
$Y' \cap f^{-1}(V) \to V$가 동형임을 얻는다.
\end{proof}
\section{스킴론적 폐포와 조밀성}
\label{morphisms-section-scheme-theoretic-closure}

\noindent
다음 정의는 \cite[IV, Definition 11.10.2]{EGA}에서 가져온다.

\begin{definition}
\label{morphisms-definition-scheme-theoretically-dense}
$X$를 스킴이라 하자. $U \subset X$를 열린 부분스킴이라 하자.
\begin{enumerate}
\item 사상 $U \to X$의 스킴론적 상을 {\it $U$의 $X$에서의 스킴론적 폐포}라 한다.
\item $U$가 {\it $X$에서 스킴론적으로 조밀하다}고 하는 것은 모든 열린집합
$V \subset X$에 대하여 $U \cap V$의 스킴론적 폐포가
$V$에서 $V$와 같다는 뜻이다.
\end{enumerate}
\end{definition}

\noindent
이 정의에서는 $U$가 $X$에서 스킴론적으로 조밀하다는 것이
$U$의 스킴론적 폐포가 $X$라는 것과 동치가 {\bf 아니다}. 예를 들어
\ref{morphisms-example-scheme-theretically-dense-not-dense}를 보라.
이는 다소 우아하지 않지만, 아래의 보조정리
\ref{morphisms-lemma-scheme-theoretically-dense-quasi-compact}와
\ref{morphisms-lemma-reduced-scheme-theoretically-dense}를 보라.
한편 이 정의에 따르면 $U$가 $X$에서 스킴론적으로 조밀한 것은
모든 열린 $V \subset X$에 대하여 환 사상
$\mathcal{O}_X(V) \to \mathcal{O}_X(U \cap V)$가 단사인 것과 동치이다.
아래 보조정리 \ref{morphisms-lemma-characterize-scheme-theoretically-dense}를 보라.
특히 스킴론적 조밀성은 조밀성을 함의하며, 이는 만족스럽다.

\begin{example}
\label{morphisms-example-scheme-theretically-dense-not-dense}
스킴론적 폐포가 $X$라는 것이 기저 위상공간에서의 조밀성을
함의하지 않는 예를 보자.
$k$를 체라 하자.
$A = k[x, z_1, z_2, \ldots]/(x^n z_n)$로 두자.
$I = (z_1, z_2, \ldots) \subset A$로 두자.
아핀 스킴 $X = \Spec(A)$와 열린 부분스킴
$U = X \setminus V(I)$를 생각하자.
$A \to \prod_n A_{z_n}$가 단사이므로 $U$의 스킴론적 폐포는
$X$임을 알 수 있다. 사상
$X \to \Spec(k[x])$를 생각하자. 이 사상은 전사이다
(모든 $z_n = 0$으로 두면 이를 알 수 있다). 그러나 이 사상을
$U$로 제한하면 전사가 아니다. 왜냐하면 점 $x = 0$으로 사상되기
때문이다. 따라서 $U$는 $X$에서 위상적으로 조밀할 수 없다.
\end{example}

\begin{lemma}
\label{morphisms-lemma-scheme-theoretically-dense-quasi-compact}
$X$를 스킴이라 하자.
$U \subset X$를 열린 부분스킴이라 하자.
포함 사상 $U \to X$가 준콤팩트이면 $U$가 $X$에서 스킴론적으로
조밀한 것은 $U$의 $X$에서의 스킴론적 폐포가 $X$라는 것과 동치이다.
\end{lemma}

\begin{proof}
이는 보조정리 \ref{morphisms-lemma-quasi-compact-scheme-theoretic-image}의
(3)에서 따른다.
\end{proof}

\begin{example}
\label{morphisms-example-scheme-theoretic-closure}
$A$를 환이라 하고 $X = \Spec(A)$라 하자.
$f_1, \ldots, f_n \in A$를 잡고
$U = D(f_1) \cup \ldots \cup D(f_n)$로 두자.
$I = \Ker(A \to \prod A_{f_i})$로 두자.
그러면 $U$의 $X$에서의 스킴론적 폐포는 닫힌 부분스킴
$\Spec(A/I)$이고, 이는 $X$의 부분스킴이다.
$U \to X$가 준콤팩트임에 유의하자. 따라서 보조정리
\ref{morphisms-lemma-scheme-theoretically-dense-quasi-compact}에 의해
$U$가 $X$에서 스킴론적으로 조밀한 것은 $I = 0$인 것과 동치이다.
\end{example}
\begin{lemma}
\label{morphisms-lemma-characterize-scheme-theoretically-dense}
$j : U \to X$를 스킴의 열린 몰입이라 하자.
그러면 $U$가 $X$에서 스킴론적으로 조밀한 것은
$\mathcal{O}_X \to j_*\mathcal{O}_U$가 단사인 것과 동치이다.
\end{lemma}

\begin{proof}
$\mathcal{O}_X \to j_*\mathcal{O}_U$가 단사이면
임의의 열린집합 $V$에 대해 $X$에서 제한해도 단사이다.
따라서 $U \cap V$의 $V$에서의 스킴론적 폐포는 $V$와 같다.
스킴론적 상에 관한 보조정리
\ref{morphisms-lemma-scheme-theoretic-image}의 증명을 보라.
역으로 $U \cap V$의 스킴론적 폐포가 $V$와 같다고 하자.
이는 모든 열린집합 $V$에 대하여 성립한다.
$\mathcal{O}_X \to j_*\mathcal{O}_U$가 단사가 아니라고 가정하자.
그러면 아핀 열린집합, 즉
$\Spec(A) = V \subset X$와, $f \in A$인 0이 아닌 원소를 택할 수 있고
$f$는 $\Gamma(V \cap U, \mathcal{O}_X)$에서 0으로 사상된다.
이 경우 $V \cap U$의 $V$에서의 스킴론적 폐포는 분명히
$\Spec(A/(f))$에 포함되므로 모순이다.
\end{proof}

\begin{lemma}
\label{morphisms-lemma-intersection-scheme-theoretically-dense}
$X$를 스킴이라 하자. $U$, $V$가 $X$의 스킴론적으로 조밀한
열린 부분스킴이면 $U \cap V$도 그러하다.
\end{lemma}

\begin{proof}
$W \subset X$를 임의의 열린집합이라 하자.
다음 사상을 생각하자.
$\mathcal{O}_X(W) \to \mathcal{O}_X(W \cap V)
\to \mathcal{O}_X(W \cap V \cap U)$.
보조정리 \ref{morphisms-lemma-characterize-scheme-theoretically-dense}에 의해
두 사상은 모두 단사이다. 따라서 합성도 단사이다.
다시 보조정리 \ref{morphisms-lemma-characterize-scheme-theoretically-dense}에 의해
$U \cap V$는 $X$에서 스킴론적으로 조밀하다.
\end{proof}
\begin{lemma}
\label{morphisms-lemma-quasi-compact-immersion}
몰입 $h : Z \to X$가 주어졌다고 하자. $h$가 준콤팩트이거나
$Z$가 축약이라고 가정하자. $\overline{Z} \subset X$를 $h$의
스킴론적 상이라 하자. 그러면 사상 $Z \to \overline{Z}$는 열린 몰입이며
$Z$를 $\overline{Z}$의 스킴론적으로 조밀한 열린 부분스킴으로
식별한다. 더 나아가 $Z$는 $\overline{Z}$에서 위상적으로 조밀하다.
\end{lemma}

\begin{proof}
보조정리 \ref{morphisms-lemma-factor-quasi-compact-immersion} 또는
보조정리 \ref{morphisms-lemma-factor-reduced-immersion}에 의해
$Z \to X$를 $Z \to \overline{Z}_1 \to X$로 인수분해할 수 있다.
여기서 $Z \to \overline{Z}_1$은 열린 몰입이고
$\overline{Z}_1 \to X$는 닫힌 몰입이다. 한편
$Z \to \overline{Z} \subset X$를 $Z \to X$의 스킴론적 상이라 하자.
그러면 $\overline{Z} \subset \overline{Z}_1$임을 얻는다.
$Z$는 $\overline{Z}_1$의 열린 부분스킴이므로
$Z$는 $\overline{Z}$의 열린 부분스킴이기도 하다.
$Z$가 축약인 경우, 보조정리
\ref{morphisms-lemma-factor-reduced-immersion}의 증명에서
$Z \subset \overline{Z}_1$은
$\overline{Z}_1$을 구성한 방식에 의해 위상적으로 조밀하다.
따라서 $\overline{Z}_1$과 $\overline{Z}$는 같은 기저 위상공간을 갖는다.
그러므로 $\overline{Z} \subset \overline{Z}_1$은 축약 스킴으로의
닫힌 몰입이며, 기저 위상공간에 전단사 사상을 유도하므로 동형이다.
$Z \to X$가 준콤팩트인 경우에는 다음과 같이 논증한다.
$Z$가 $\overline{Z}$에서 스킴론적으로 조밀하다는 주장은
보조정리 \ref{morphisms-lemma-quasi-compact-scheme-theoretic-image}의
(3)에서 따른다. 마지막 주장은
보조정리 \ref{morphisms-lemma-quasi-compact-scheme-theoretic-image}의
(4)에서 따른다.
\end{proof}

\begin{lemma}
\label{morphisms-lemma-reduced-scheme-theoretically-dense}
$X$를 축약 스킴이라 하고 $U \subset X$를 열린 부분스킴이라 하자.
그러면 다음 조건들은 서로 동치이다.
\begin{enumerate}
\item $U$는 $X$에서 위상적으로 조밀하다.
\item $U$의 $X$에서의 스킴론적 폐포는 $X$이고,
\item $U$는 $X$에서 스킴론적으로 조밀하다.
\end{enumerate}
\end{lemma}

\begin{proof}
이는 보조정리 \ref{morphisms-lemma-quasi-compact-immersion}과
$Z$를 $X$의 닫힌 부분스킴이라 하자. 그 기저 위상공간이
$X$와 같으면 $X$와 스킴으로서 같아야 한다는 사실에서 따른다.
\end{proof}

\begin{lemma}
\label{morphisms-lemma-reduced-subscheme-closure}
$X$를 스킴이라 하고 $U \subset X$를 축약 열린 부분스킴이라 하자.
그러면 다음 조건들은 서로 동치이다.
\begin{enumerate}
\item $U$의 $X$에서의 스킴론적 폐포는 $X$이고,
\item $U$는 $X$에서 스킴론적으로 조밀하다.
\end{enumerate}
이 조건들이 성립하면 $X$는 축약 스킴이다.
\end{lemma}

\begin{proof}
이는 보조정리 \ref{morphisms-lemma-quasi-compact-immersion}과
$U$의 $X$에서의 스킴론적 폐포가 보조정리
\ref{morphisms-lemma-scheme-theoretic-image-reduced}에 의해 축약이라는
사실에서 따른다.
\end{proof}

\begin{lemma}
\label{morphisms-lemma-equality-of-morphisms}
$S$를 스킴이라 하자. $X$, $Y$를 $S$ 위의 스킴이라 하자.
$f, g : X \to Y$를 $S$ 위의 스킴 사상이라 하자.
$U \subset X$를 열린 부분스킴이라 하고 $f|_U = g|_U$라고 하자.
$U$의 $X$에서의 스킴론적 폐포가 $X$이고 $Y \to S$가 분리라면
$f = g$이다.
\end{lemma}

\begin{proof}
정의와
Schemes의 보조정리 \ref{schemes-lemma-where-are-they-equal}에서 따른다.
\end{proof}

\section{우세 사상}
\label{morphisms-section-dominant}

\noindent
스킴 사상이 우세하다는 정의는 다양체의 우세 사상 개념에 익숙한
경우에 예상하는 것과 조금 다르다.

\begin{definition}
\label{morphisms-definition-dominant}
스킴의 사상 $f : X \to S$가 {\it 우세하다}고 함은
$f$의 상이 $S$의 조밀한 부분집합인 경우이다.
\end{definition}

\noindent
예를 들어 $k$가 무한체이고 $\lambda_1, \lambda_2, \ldots$가
$k$의 서로 다른 원소들의 가산 모임이라고 하자. 그러면 다음 사상
$$
\coprod\nolimits_{i = 1, 2, \ldots } \Spec(k)
\longrightarrow
\Spec(k[x])
$$
에서 $i$번째 인자는 점 $x = \lambda_i$로 사상되며 우세하다.

\begin{lemma}
\label{morphisms-lemma-generic-points-in-image-dominant}
$f : X \to S$를 스킴의 사상이라 하자.
$S$의 모든 기약 성분의 모든 일반점이 $f$의 상에 속하면
$f$는 우세하다.
\end{lemma}

\begin{proof}
이는 다음 두 사실에서 직접 따르는 위상적 사실이다. 스킴의
기저 위상공간은 소버이다(스킴의 보조정리
\ref{schemes-lemma-scheme-sober}를 보라). 또한 $S$의 모든 점은
$S$의 어떤 기약 성분에 포함된다(위상의 보조정리
\ref{topology-lemma-irreducible}를 보라).
\end{proof}

\noindent
사상이 준콤팩트이면, 사상의 일반점이 상에 속할 때에만 우세하다는
예상이 성립한다.

\begin{lemma}
\label{morphisms-lemma-quasi-compact-dominant}
\begin{slogan}
상이 일반점을 포함하는 사상은 우세하다
\end{slogan}
$f : X \to S$를 스킴의 준콤팩트 사상이라 하자.
그러면 $f$가 우세한 것은 모든 기약 성분
$Z \subset S$에 대해 $Z$의 일반점이 $f$의 상에 속하는 것과 동치이다.
\end{lemma}

\begin{proof}
$V \subset S$를 아핀 열린집합이라 하자.
$f$가 준콤팩트이므로 유한 개의 아핀 열린집합
$U_i \subset f^{-1}(V)$, $i = 1, \ldots, n$를 택하여
$f^{-1}(V)$를 덮게 할 수 있다.
아핀 스킴의 사상을 생각하자.
$$
f' :
\coprod\nolimits_{i = 1, \ldots, n} U_i
\longrightarrow
V.
$$
아핀들의 서로소 합집합은 아핀이다(스킴의 보조정리
\ref{schemes-lemma-disjoint-union-affines}를 보라).
$V$의 기약 성분의 일반점은 정확히 $S$의 기약 성분 중
$V$와 만나는 것들의 일반점이다. 또한 $f$가 우세한 것은
$f'$가 우세한 것과 동치이며, 이는 위에서 $V, n, U_i$를
어떻게 선택하더라도 성립한다.
따라서 이 보조정리는 아핀 스킴 사상의 경우로 환원된다.
아핀의 경우는 Algebra의 보조정리
\ref{algebra-lemma-image-dense-generic-points}이다.
\end{proof}

\begin{lemma}
\label{morphisms-lemma-base-change-dominant}
$f : X \to S$를 스킴의 준콤팩트 우세 사상이라 하자.
$g : S' \to S$를 스킴의 사상이라 하고, $f' : X' \to S'$를
$f$의 $g$에 의한 기저 변환이라 하자.
일반화가 $g$를 따라 올라가면 $f'$는 우세하다.
\end{lemma}

\noindent
뒤에서 $g$가 열린 사상(보조정리 \ref{morphisms-lemma-open-generizing})이거나
평탄 사상(보조정리 \ref{morphisms-lemma-generalizations-lift-flat})이면
일반화가 $g$를 따라 올라감을 보일 것이다.

\begin{proof}
스킴의 보조정리
\ref{schemes-lemma-quasi-compact-preserved-base-change}에 의해
$f'$가 준콤팩트임에 유의하자.
$\eta' \in S'$를 기약 성분의 일반점이라 하자.
$S'$의 기약 성분에 대해 일반화가 $g$를 따라 올라가면
$\eta = g(\eta')$는 $S$의
기약 성분의 일반점이다. 보조정리
\ref{morphisms-lemma-quasi-compact-dominant}에 의해 $\eta$는 $f$의 상에 속한다.
따라서 스킴의 보조정리
\ref{schemes-lemma-points-fibre-product}에 의해 $\eta'$는
$f'$의 상에 속한다.
그러므로 보조정리 \ref{morphisms-lemma-quasi-compact-dominant}에 의해
$f'$는 우세하다.
\end{proof}

\noindent
이 결과의 일부는 임의의 (반드시 준콤팩트일 필요는 없는) 우세 사상에도
성립한다.

\begin{lemma}
\label{morphisms-lemma-open-base-change-dominant}
$f : X \to S$를 스킴의 우세 사상이라 하자.
$g : S' \to S$를 스킴의 열린 사상이라 하고, $f' : X' \to S'$를
$f$의 $g$에 의한 기저 변환이라 하자. 그러면 $f'$는 우세하다.
\end{lemma}

\begin{proof}
$f'$가 우세하지 않다면 $f'(X')$가 $B$라는 어떤 닫힌
부분집합에 포함되며 이는 $S'$와 같지 않다. 그러면 $S'-B$는
$S'$에서 공집합이 아닌
열린집합이다. $g$가 열림에 의해 $g(S'-B)$는 $S$에서 공집합이 아닌
열린집합이다. $f(X)$가 $S$에서 조밀하므로
$x\in X$인 점으로서 $f(x) \in g(S'-B)$를 만족하는 것을 택할 수 있다.
$W = \Spec(\kappa (f(x)))$라 하자.
$X_W$, $S_W$, $S'_W$를 이 스킴들을 $S$에서 $W$로 당겨 얻은
당김들로 쓰자. 그러면 $X_W$와 $(S'-B)_W$는 공집합이 아니며,
$X_W \times_W (S'-B)_W$는 공집합이 아니다. 이는 $W$가 체의
스펙트럼이기 때문이다.
이는 $f'(X \times_S S')$가 $B \subset S'$에 포함된다는 것과 모순이다.
따라서 실제로 $f'$는 우세하다.
\end{proof}

\begin{lemma}
\label{morphisms-lemma-quasi-compact-generic-point-not-in-image}
$f : X \to S$를 스킴의 준콤팩트 사상이라 하자.
$\eta \in S$를 $S$의 기약 성분의 일반점이라 하자.
$\eta \not \in f(X)$이면 열린 근방 $V \subset S$가 존재하여
$\eta$를 포함하고 $f^{-1}(V) = \emptyset$이다.
\end{lemma}

\begin{proof}
$Z \subset S$를 $f$의 스킴론적 상이라 하자.
$\eta \not \in Z$임을 보여야 한다. 이는
보조정리 \ref{morphisms-lemma-reach-points-scheme-theoretic-image}에서 따르지만,
다음과 같이 볼 수도 있다.
보조정리 \ref{morphisms-lemma-quasi-compact-scheme-theoretic-image}에 의해
$X \to Z$는 우세하고, 보조정리
\ref{morphisms-lemma-quasi-compact-dominant}에 의해 이는 $Z$의 모든 기약
성분의 모든 일반점이 $X \to Z$의 상에 속한다는 뜻이다.
가정에 의해 $\eta \not \in Z$이다. $\eta$가 $Z$에 속한다면
$Z$의 어떤 기약 성분의 일반점일 것이기 때문이다.
\end{proof}

\noindent
우세가 기약 성분의 모든 일반점이 상에 속한다는 것과 동치가 되는
또 다른 경우가 있다.

\begin{lemma}
\label{morphisms-lemma-dominant-finite-number-irreducible-components}
$f : X \to S$를 스킴의 사상이라 하자.
$X$가 유한 개의 기약 성분을 가진다고 가정하자.
그러면 $f$는 우세하다 (그리고 우세할 필요충분조건은 다음과 같다):
모든 기약 성분 $Z \subset S$에 대해 $Z$의 일반점이 $f$의 상에 속한다.
그렇다면 $S$도 유한 개의 기약 성분을 가진다.
\end{lemma}

\begin{proof}
$f$가 우세하다고 가정하자.
$X = Z_1 \cup Z_2 \cup \ldots \cup Z_n$을 $X$를 기약 성분들로
분해한 것이라 하자. $\xi_i \in Z_i$를 그 일반점이라 하자.
그러면 $Z_i = \overline{\{\xi_i\}}$이다.
$f(Z_i)$는 $S$의 기약 부분집합임에 유의하자.
따라서
$$
S = \overline{f(X)} = \bigcup \overline{f(Z_i)} =
\bigcup \overline{\{f(\xi_i)\}}
$$
는 기약 부분집합들의 유한 합집합이며, 그 일반점들은 $f$의 상에 속한다.
이로부터 보조정리가 따른다.
\end{proof}

\begin{lemma}
\label{morphisms-lemma-dominant-between-integral}
$f : X \to Y$를 정수적 스킴 사이의 사상이라 하자. 다음 조건들은
서로 동치이다.
\begin{enumerate}
\item $f$는 우세하다.
\item $f$는 $X$의 일반점을 $Y$의 일반점으로 사상한다.
\item 어떤 공집합이 아닌 아핀 열린집합 $U \subset X$와
$V \subset Y$에 대해 $f(U) \subset V$이면 환 사상
$\mathcal{O}_Y(V) \to \mathcal{O}_X(U)$는 단사이다.
\item 모든 공집합이 아닌 아핀 열린집합 $U \subset X$와
$V \subset Y$에 대해 $f(U) \subset V$이면 환 사상
$\mathcal{O}_Y(V) \to \mathcal{O}_X(U)$는 단사이다.
\item 어떤 $x \in X$에 대해 그 상 $y = f(x) \in Y$에서 국소환 사상
$\mathcal{O}_{Y, y} \to \mathcal{O}_{X, x}$가 단사이다.
\item 모든 $x \in X$에 대해 그 상 $y = f(x) \in Y$에서 국소환 사상
$\mathcal{O}_{Y, y} \to \mathcal{O}_{X, x}$가 단사이다.
\end{enumerate}
\end{lemma}

\begin{proof}
(1)과 (2)의 동치성은 보조정리
\ref{morphisms-lemma-dominant-finite-number-irreducible-components}에서 따른다.
$U \subset X$와 $V \subset Y$를 공집합이 아닌 아핀 열린집합으로 잡자.
$f(U) \subset V$라고 하자. 환 $A = \mathcal{O}_X(U)$와
$B = \mathcal{O}_Y(V)$가 정역임을 상기하자.
사상 $f|_U : U \to V$는 환 사상 $\varphi : B \to A$에 대응한다.
$X$와 $Y$의 일반점은 각각 소 아이디얼
$(0) \subset A$와 $(0) \subset B$에 대응한다. 따라서 (2)는
$(0) = \varphi^{-1}((0))$이라는 조건, 즉 $\varphi$가 단사라는
조건과 동치이다.
이로써 (2), (3), (4)가 동치임을 알 수 있다.
마찬가지로 (5)의 $x$와 $y$가 주어지면 국소환
$\mathcal{O}_{X, x}$와 $\mathcal{O}_{Y, y}$는 정역이고,
소 아이디얼 $(0) \subset \mathcal{O}_{X, x}$와
$(0) \subset \mathcal{O}_{Y, y}$는 $X$와 $Y$의 일반점에 대응한다
(이는 $x$에서의 국소환의 스펙트럼을 $x$로 특수화되는 점들의 집합과
동일시하는 것을 통해서이다. 스킴의 보조정리
\ref{schemes-lemma-specialize-points}를 보라).
따라서 위와 정확히 같은 방식으로 논증하여 (2), (5), (6)이
동치임을 얻는다.
\end{proof}

\section{전사 사상}
\label{morphisms-section-surjective}

\begin{definition}
\label{morphisms-definition-surjective}
스킴의 사상이 {\it 전사}라고 함은 그 기저 위상공간들 사이의
사상이 전사인 경우이다.
\end{definition}

\begin{lemma}
\label{morphisms-lemma-composition-surjective}
전사 사상들의 합성은 전사이다.
\end{lemma}

\begin{proof}
생략한다.
\end{proof}

\begin{lemma}
\label{morphisms-lemma-when-point-maps-to-pair}
$X$와 $Y$를 기저 스킴 $S$ 위의 스킴이라 하자. 점
$x \in X$와 $y \in Y$가 주어졌다고 하자. $X \times_S Y$에
사영에 의해 $x$와 $y$로 가는 점이 존재하는 것은
$x$와 $y$가 $S$의 같은 점 위에 있을 때 그리고 그때뿐이다.
\end{lemma}

\begin{proof}
이 조건이 필요하다는 것은 명백하고, 역은 스킴의 보조정리
\ref{schemes-lemma-points-fibre-product}의 증명에서 따른다.
\end{proof}

\begin{lemma}
\label{morphisms-lemma-base-change-surjective}
전사 사상의 기저 변환은 전사이다.
\end{lemma}

\begin{proof}
$f: X \to Y$를 기저 스킴 $S$ 위의 스킴 사상이라 하자.
$S' \to S$가 스킴의 사상이라면 정준 사영
$p: X_{S'} \to X$와 $q: Y_{S'} \to Y$를 쓰자. 다음 가환
도식은
$$
\xymatrix{
X_{S'} \ar[d]_{f_{S'}} \ar[r]_p & X \ar[d]^{f} \\
Y_{S'} \ar[r]^{q} & Y.
}
$$
$X_{S'}$를 $X \to Y$와 $Y_{S'} \to Y$의 올곱으로 식별한다.
$Z$를 $X$의 기저 위상공간의 부분집합이라 하자. 그러면
$q^{-1}(f(Z)) = f_{S'}(p^{-1}(Z))$이다. 왜냐하면
$y' \in q^{-1}(f(Z))$인 것은 $q(y') = f(x)$인
$x \in Z$가 존재하는 것과 동치이고, 보조정리
\ref{morphisms-lemma-when-point-maps-to-pair}에 의해 이는
$x' \in X_{S'}$가 존재하여
$f_{S'}(x') = y'$이고 $p(x') = x$인 것과 동치이기 때문이다.
특히 $Z = X$로 두면 $f$가 전사일 때 기저 변환
$f_{S'}: X_{S'} \to Y_{S'}$도 전사임을 알 수 있다.
\end{proof}

\begin{example}
\label{morphisms-example-injective-not-preserved-base-change}
전단사성은 기저 변환에서 보존되지 않으므로 단사성도 보존되지 않는다.
예를 들어 전단사
$\Spec(\mathbf{C}) \to \Spec(\mathbf{R})$
를 생각하자. 그 기저 변환
$\Spec(\mathbf{C} \otimes_{\mathbf{R}} \mathbf{C}) \to
\Spec(\mathbf{C})$
는 단사가 아니다. 실제로 동형
$\mathbf{C} \otimes_{\mathbf{R}} \mathbf{C} \cong \mathbf{C} \times \mathbf{C}$
가 존재한다(이 분해는 멱등원
$\frac{1 \otimes 1 + i \otimes i}{2}$에서 온다).
따라서 $\Spec(\mathbf{C} \otimes_{\mathbf{R}} \mathbf{C})$는
두 점을 가진다.
\end{example}

\begin{lemma}
\label{morphisms-lemma-surjection-from-quasi-compact}
다음과 같은 스킴 사상들의 가환 도식이 있다고 하자.
$$
\xymatrix{
X \ar[rr]_f \ar[rd]_p & &
Y \ar[dl]^q \\
& Z
}
$$
$f$가 전사이고 $p$가 준콤팩트이면 $q$는 준콤팩트이다.
\end{lemma}

\begin{proof}
$W \subset Z$를 준콤팩트 열린집합이라 하자. 가정에 의해
$p^{-1}(W)$는 준콤팩트이다. 따라서 위상의 보조정리
\ref{topology-lemma-image-quasi-compact}에 의해
$q^{-1}(W) = f(p^{-1}(W))$도 준콤팩트이다.
이로써 보조정리가 증명된다.
\end{proof}

\section{순수 비분기 사상과 보편 단사 사상}
\label{morphisms-section-radicial}

\noindent
이 절에서는 스킴의 사상이 {\it 순수 비분기}라는 것과
스킴의 사상이 {\it 보편적으로 단사}라는 것이 무엇을 뜻하는지 정의한다.
그런 다음 이 두 개념이 일치함을 보인다.
두 개를 모두 도입하는 이유는 대수적 공간의 경우 이에 대응하는
개념들이 항상 일치하지 않을 수 있기 때문이다.

\begin{definition}
\label{morphisms-definition-universally-injective}
$f : X \to S$를 사상이라 하자.
\begin{enumerate}
\item $f$가 {\it 보편적으로 단사}라고 함은 임의의 스킴 사상
$S' \to S$에 대해 기저 변환 $f' : X_{S'} \to S'$가
(기저 위상공간에서) 단사인 경우, 그리고 그 경우뿐이다.
\item $f$가 {\it 순수 비분기}라고 함은 $f$가 위상공간의 사상으로서
단사이고, 모든 $x \in X$에 대해 체의 확대
$\kappa(x)/\kappa(f(x))$가 순수 비분리인 경우이다.
\end{enumerate}
\end{definition}

\begin{lemma}
\label{morphisms-lemma-universally-injective}
\begin{reference}
\cite[Proposition 3.7.1]{EGA1-second}
\end{reference}
$f : X \to S$를 스킴의 사상이라 하자.
다음 조건들은 서로 동치이다.
\begin{enumerate}
\item 모든 체 $K$에 대해 유도된 사상
$\Mor(\Spec(K), X) \to \Mor(\Spec(K), S)$가 단사이다.
\item 사상 $f$는 보편적으로 단사이다.
\item 사상 $f$는 순수 비분기이다.
\item 대각 사상 $\Delta_{X/S} : X \longrightarrow X \times_S X$가
전사이다.
\end{enumerate}
\end{lemma}

\begin{proof}
$K$를 체라 하고 $s : \Spec(K) \to S$를 사상이라 하자.
$x : \Spec(K) \to X$인 사상으로서 $f \circ x = s$를 만족하는 것을
주는 것은
사영 $X_K = \Spec(K) \times_S X \to \Spec(K)$의 단면을 주는 것과
동치이다. 이는 다시 점 $x \in X_K$로서 잉여체가 $K$인 것을 주는 것과
동치이다. 따라서 (2)가 (1)을 함의함을 알 수 있다.

\medskip\noindent
반대로 (1)이 성립한다고 하자. $x, x' \in X_{S'}$가 같은 점
$s' \in S'$로 사상된다고 하자. 다음 가환 도식을 택하자.
$$
\xymatrix{
K & \kappa(x) \ar[l] \\
\kappa(x') \ar[u] & \kappa(s') \ar[l] \ar[u]
}
$$
이는 체들의 도식이다. 스킴의 보조정리
\ref{schemes-lemma-characterize-points}에 의해
$a, a' : \Spec(K) \to X_{S'}$라는 두 사상을 얻는다.
하나는 점 $x$와 포함 $\kappa(x) \subset K$에 대응하고,
다른 하나는 점 $x'$와 포함 $\kappa(x') \subset K$에 대응한다.
또한 $f' \circ a = f' \circ a'$이다.
이제 (1)에서 $a$와 $a'$를
$X_{S'} \to X$와 합성한 것이 같음을 얻는다.
$X_{S'}$는 $S'$와 $X$의 $S$ 위 올곱이므로
$a = a'$이다. 따라서 $x = x'$이다.
그러므로 (1)은 (2)를 함의한다.

\medskip\noindent
서로 다른 점 $x, x' \in X$가 $s$의 같은 점으로 사상되면
(2)가 깨진다.
어떤 $s = f(x)$에 대해 $x \in X$에서 체의 확대
$\kappa(x)/\kappa(s)$가 순수 비분리라면, 다음 조건을 만족하는 체의
확대 $K/\kappa(s)$를 찾을 수 있다. 즉 $\kappa(x)$가
$\kappa(s)$-준동형을 $K$ 안으로 두 개 갖는다. 스킴의 보조정리
\ref{schemes-lemma-characterize-points}에 의해 이는 사상
$\Mor(\Spec(K), X) \to \Mor(\Spec(K), S)$가 단사가 아님을
뜻하므로 (1)이 깨진다.
따라서 동치인 조건 (1)과 (2)가 $f$가 순수 비분기임을 함의한다.
즉 이 조건들은 (3)을 함의한다.

\medskip\noindent
(3)을 가정하자. 스킴의 보조정리
\ref{schemes-lemma-characterize-points}에 의하면
$\Spec(K) \to X$인 사상은 순서쌍
$(x, \kappa(x) \to K)$로 주어진다.
(3)의 성질은 순서쌍
$(x, \kappa(x) \to K)$를 순서쌍
$(s, \kappa(s) \to \kappa(x) \to K)$에 대응시키는 것이
단사라는 것과 정확히 같다. 다르게 말하면 (1)이 성립한다.
이제 (1), (2), (3)이 모두 동치임을 알았다.

\medskip\noindent
마지막으로 (4)가 (1), (2), (3)과 동치임을 증명하자.
$X \times_S X$의 점은 네 쌍
$(x_1, x_2, s, \mathfrak p)$로 주어진다. 여기서
$x_1, x_2 \in X$이고,
$f(x_1) = f(x_2) = s$이며
$\mathfrak p \subset \kappa(x_1) \otimes_{\kappa(s)} \kappa(x_2)$가
소 아이디얼이다. 스킴의 보조정리
\ref{schemes-lemma-points-fibre-product}를 보라.
$f$가 보편적으로 단사이면, 보편적 단사의 정의에서 $S'=X$로
두어 얻는 바에 의해 $\Delta_{X/S}$는 전사여야 한다. 이는
단사 사상
$X \times_S X  \longrightarrow X$의 단면이기 때문이다.
반대로 $\Delta_{X/S}$가 전사라면 항상
$x_1 = x_2 = x$이고 그러한 소 아이디얼 $\mathfrak p$는 정확히
하나뿐이다. 이는 $\kappa(s) \subset \kappa(x)$가 순수 비분리라는
뜻이다. 따라서 $f$는 순수 비분기이다.
다르게는, $\Delta_{X/S}$가 전사라면 임의의 $S' \to S$에 대해
기저 변환 $\Delta_{X_{S'}/S'}$가 전사이고, 이는 $f$가
보편적으로 단사임을 함의한다. 이로써 보조정리의 증명이 끝난다.
\end{proof}

\begin{lemma}
\label{morphisms-lemma-universally-injective-separated}
보편적으로 단사인 사상은 분리 사상이다.
\end{lemma}

\begin{proof}
보조정리 \ref{morphisms-lemma-universally-injective}와 다음 주의를 결합하자.
$X \to S$가 분리라는 것은 대각 사상
$\Delta_{X/S}$의 상이 $X \times_S X$에서 닫힌 것과 동치이다.
스킴의 정의 \ref{schemes-definition-separated}와 그 뒤의 논의를 보라.
\end{proof}

\begin{lemma}
\label{morphisms-lemma-base-change-universally-injective}
보편적으로 단사인 사상의 기저 변환은 보편적으로 단사이다.
\end{lemma}

\begin{proof}
형식적인 결과이다.
\end{proof}

\begin{lemma}
\label{morphisms-lemma-composition-universally-injective}
순수 비분기 사상들의 합성은 순수 비분기이며, 따라서 동치인 조건인
보편적으로 단사라는 성질에 대해서도 같은 결과가 성립한다.
\end{lemma}

\begin{proof}
생략한다.
\end{proof}

\section{아핀 사상}
\label{morphisms-section-affine}

\begin{definition}
\label{morphisms-definition-affine}
스킴의 사상 $f : X \to S$가 {\it 아핀}이라고 함은
$S$의 모든 아핀 열린집합의 역상이 $X$의 아핀 열린집합인 경우이다.
\end{definition}

\begin{lemma}
\label{morphisms-lemma-affine-separated}
아핀 사상은 분리 사상이고 준콤팩트이다.
\end{lemma}

\begin{proof}
$f : X \to S$를 아핀이라 하자. 준콤팩트성은 스킴의 보조정리
\ref{schemes-lemma-quasi-compact-affine}에서 즉시 따른다.
스킴의 보조정리 \ref{schemes-lemma-characterize-separated}를 사용하여
$f$가 분리임을 보이자. $x_1, x_2 \in X$를 $X$의 점들 중
같은 점 $s \in S$로 사상되는 것이라 하자. $W \subset S$를
$s$를 포함하는 임의의 아핀 열린집합으로 택하자. 가정에 의해
$f^{-1}(W)$는 아핀이다. 인용한 보조정리에서
$U = V = f^{-1}(W)$로 두어 적용하면 된다.
\end{proof}

\begin{lemma}
\label{morphisms-lemma-characterize-affine}
\begin{reference}
\cite[II, Corollary 1.3.2]{EGA}
\end{reference}
$f : X \to S$를 스킴의 사상이라 하자.
다음 조건들은 서로 동치이다.
\begin{enumerate}
\item 사상 $f$는 아핀이다.
\item 아핀 열린 덮개 $S = \bigcup W_j$가 존재하여 각
$f^{-1}(W_j)$가 아핀이다.
\item $\mathcal{O}_S$-대수의 준연접 층 $\mathcal{A}$와
$X \cong \underline{\Spec}_S(\mathcal{A})$라는 동형이 존재한다.
이는 $S$ 위 스킴의 동형이다. 표기는 구성의 절
\ref{constructions-section-spec}을 보라.
\end{enumerate}
더 나아가 이 경우 $X = \underline{\Spec}_S(f_*\mathcal{O}_X)$이다.
\end{lemma}

\begin{proof}
(1)은 (2)를 함의한다는 것이 명백하다.

\medskip\noindent
$S = \bigcup_{j \in J} W_j$가 아핀 열린 덮개이고 각
$f^{-1}(W_j)$가 아핀이라고 하자. 스킴의 보조정리
\ref{schemes-lemma-quasi-compact-affine}에 의해 $f$는 준콤팩트이다.
스킴의 보조정리 \ref{schemes-lemma-characterize-quasi-separated}에 의해
사상 $f$는 준분리이다. 따라서 스킴의 보조정리
\ref{schemes-lemma-push-forward-quasi-coherent}에 의해
층 $\mathcal{A} = f_*\mathcal{O}_X$는 $\mathcal{O}_S$-대수의
준연접 층이다. 그러므로 스킴
$g : Y = \underline{\Spec}_S(\mathcal{A}) \to S$가 $S$ 위에 있다.
항등 사상
$\text{id} : \mathcal{A} = f_*\mathcal{O}_X \to f_*\mathcal{O}_X$는
$can : X \to Y$라는 $S$ 위의 사상을 상대 스펙트럼의 정의를 통해
준다. 구성의 보조정리
\ref{constructions-lemma-canonical-morphism}를 보라.
가정과 방금 인용한 보조정리에 의해 제한
$can|_{f^{-1}(W_j)} : f^{-1}(W_j) \to g^{-1}(W_j)$는
동형이다. 따라서 $can$은 동형이다. 이로써 (2)가 (3)을
함의함을 보였다.

\medskip\noindent
(3)을 가정하자. 구성의 보조정리
\ref{constructions-lemma-spec-properties}에 의해 모든 아핀 열린집합의
역상은 아핀이다. 따라서 정의에 의해 이 사상은 아핀이다.
\end{proof}

\begin{remark}
\label{morphisms-remark-direct-argument}
위 보조정리 \ref{morphisms-lemma-characterize-affine}에서 (2)가 (1)을
함의한다는 것을 다음과 같이 직접 논증할 수도 있다.
$S = \bigcup W_j$가 아핀 열린 덮개이고 각 $f^{-1}(W_j)$가
아핀이라고 하자. 먼저 위 증명에서와 같이
$\mathcal{A} = f_*\mathcal{O}_X$가 준연접임을 보이자.
$\Spec(R) = V \subset S$를 아핀 열린집합이라 하자.
$f^{-1}(V)$가 아핀임을 보여야 한다. 다음과 같이 두자.
$A = \mathcal{A}(V) = f_*\mathcal{O}_X(V) = \mathcal{O}_X(f^{-1}(V))$.
스킴의 보조정리 \ref{schemes-lemma-morphism-into-affine}에 의해
정준 사상 $\psi : f^{-1}(V) \to \Spec(A)$가
$\Spec(R) = V$ 위에 존재한다.
스킴의 보조정리 \ref{schemes-lemma-good-subcover}에 의해 정수
$n \geq 0$, 표준 열린 덮개
$V = \bigcup_{i = 1, \ldots, n} D(h_i)$, $h_i \in R$,
그리고 사상 $a : \{1, \ldots, n\} \to J$가 존재하여 각
$D(h_i)$가 아핀 스킴 $W_{a(i)}$의 표준 열린집합이기도 하다.
아핀 스킴의 사상에 의한 표준 열린집합의 역상은 표준 열린집합이다.
Algebra의 보조정리 \ref{algebra-lemma-spec-functorial}을 보라.
따라서 $f^{-1}(D(h_i))$는 $f^{-1}(W_{a(i)})$의 표준 열린집합이고,
특히 $f^{-1}(D(h_i))$는 아핀이다. $\mathcal{A}$가 준연접이므로
$A_{h_i} = \mathcal{A}(D(h_i)) = \mathcal{O}_X(f^{-1}(D(h_i)))$이고,
따라서 $f^{-1}(D(h_i))$는 $A_{h_i}$의 스펙트럼이다.
사상 $\psi$는 열린집합 $f^{-1}(D(h_i))$와
$\Spec(A_{h_i})$의 열린집합 사이에 동형을 유도한다.
여기서 후자는 $\Spec(A)$의 열린집합이다. $f^{-1}(V) = \bigcup f^{-1}(D(h_i))$이고
$\Spec(A) = \bigcup \Spec(A_{h_i})$이므로
결론을 얻는다.
\end{remark}

\begin{lemma}
\label{morphisms-lemma-affine-equivalence-algebras}
$S$를 스킴이라 하자. 다음 두 범주 사이에 반등가가 존재한다.
$$
\begin{matrix}
\text{Schemes affine} \\
\text{over }S
\end{matrix}
\longleftrightarrow
\begin{matrix}
\text{quasi-coherent sheaves} \\
\text{of }\mathcal{O}_S\text{-algebras}
\end{matrix}
$$
이 반등가는 $f : X \to S$를 층 $f_*\mathcal{O}_X$에 대응시킨다.
더 나아가 이 등가는 임의의 기저 변환과 양립한다.
\end{lemma}

\begin{proof}
오른쪽에서 왼쪽으로 가는 함자는 $\underline{\Spec}_S$로 주어진다.
두 함자는 보조정리 \ref{morphisms-lemma-characterize-affine}와 구성의 보조정리
\ref{constructions-lemma-spec-properties} (3)에 의해 서로 역이다.
마지막 주장은 구성의 보조정리
\ref{constructions-lemma-spec-properties} (2)이다.
\end{proof}

\begin{lemma}
\label{morphisms-lemma-quasi-coherence-scalar-restriction}
\begin{reference}
\cite[I, Proposition 9.6.1]{EGA}
\end{reference}
$S$를 스킴이라 하고 $\mathcal{A}$를 준연접
$\mathcal{O}_S$-대수라 하자. $\mathcal{A}$-가군이
$\mathcal{O}_S$-가군으로서 준연접인 것은
$\mathcal{A}$-가군으로서 준연접인 것과 동치이다.
\end{lemma}

\begin{proof}
$\mathcal{F}$를 $\mathcal{A}$-가군이라 하자.
$\mathcal{F}$가 $\mathcal{A}$-가군으로서 준연접이면
모든 $s \in S$에 대해 열린 근방 $U$가 존재하여 $s$를 포함하고 정합열
$$
\bigoplus\nolimits_J\mathcal{A}|_U\to
\bigoplus\nolimits_I\mathcal{A}|_U\to
\mathcal{F}|_U
$$
인 $\mathcal{A}|_U$-가군의 정합열이 존재한다.
그러면 이는 $\mathcal{O}_U$-가군의 정합열이기도 하다.
따라서 $\mathcal{F}|_U$는 스킴 위의 준연접
$\mathcal{O}_U$-가군들의 사상의 여핵으로서 준연접이다.
그러므로 $\mathcal{F}$는 $\mathcal{O}_X$-가군으로서 준연접이다.

\medskip\noindent
반대로 $\mathcal{F}$가 $\mathcal{O}_X$-가군으로서 준연접이라고 하자.
아핀 열린집합 $\Spec(R) = U \subset S$를 택하자.
$\mathcal{O}_U$-가군의 동형
$\mathcal{A}|_U \cong \widetilde{A}$와
$\mathcal{F}|_U \cong \widetilde{M}$이 존재하며, 어떤 $R$-대수 $A$와
어떤 $R$-가군 $M$에 대응한다. $\mathcal{A}$-가군 구조가
$\mathcal{F}$ 위에 주어져 있으며, 이는 $A$-가군 구조로 $M$ 위에 옮겨가며
주어진 $R$-가군 구조와 양립한다(세부사항은 생략한다).
$A$-가군의 정합열
$$
\bigoplus\nolimits_J A \to
\bigoplus\nolimits_I A \to
M \to 0
$$
을 택하자. 함자 $\widetilde{\ }$가 완전하므로 이는
$\widetilde{A}$-가군의 정합열
$$
\bigoplus\nolimits_J \widetilde{A}\to
\bigoplus\nolimits_I \widetilde{A}\to
\widetilde{M} \to 0
$$
을 준다. 이는 $\mathcal{F}$가 $\mathcal{A}$-가군으로서
준연접임을 뜻한다.
\end{proof}

\begin{lemma}
\label{morphisms-lemma-affine-equivalence-modules}
$f : X \to S$를 스킴의 아핀 사상이라 하자.
$\mathcal{A} = f_*\mathcal{O}_X$라 하자.
함자 $\mathcal{F} \mapsto f_*\mathcal{F}$는 다음 범주들의 등가를 유도한다.
$$
\left\{
\begin{matrix}
\text{category of quasi-coherent}\\
\mathcal{O}_X\text{-modules}
\end{matrix}
\right\}
\longrightarrow
\left\{
\begin{matrix}
\text{category of quasi-coherent}\\
\mathcal{A}\text{-modules}
\end{matrix}
\right\}
$$
더 나아가 $\mathcal{A}$-가군이 $\mathcal{O}_S$-가군으로서
준연접인 것은 $\mathcal{A}$-가군으로서 준연접인 것과 동치이다.
\end{lemma}

\begin{proof}
마지막 주장은 보조정리
\ref{morphisms-lemma-quasi-coherence-scalar-restriction}이다.
보조정리 \ref{morphisms-lemma-affine-separated}와 스킴의 보조정리
\ref{schemes-lemma-push-forward-quasi-coherent}에 의해
준연접 가군의 직접상 $f_*\mathcal{F}$는 준연접
$\mathcal{O}_X$-가군 $\mathcal{F}$에 대응하며,
준연접 $\mathcal{O}_S$-가군이다.
따라서 보조정리의 서술에 나타난 함자가 존재한다.

\medskip\noindent
이 함자의 준역함자 $h$를 구성하자. $\mathcal{G}$를
준연접 $\mathcal{A}$-가군이라 하자. 다음과 같이 놓는다.
$$
h(\mathcal{G}) = f^*\mathcal{G} \otimes_{f^*\mathcal{A}} \mathcal{O}_X =
f^*\mathcal{G} \otimes_{f^*f_*\mathcal{O}_X} \mathcal{O}_X
$$
설명하자면, 당김 $f^*\mathcal{A} = f^*f_*\mathcal{O}_X$는
$\mathcal{O}_X$-대수이고, 인접 사상
$f^*f_*\mathcal{O}_X \to \mathcal{O}_X$는 대수 준동형이며,
당김 $f^*\mathcal{G}$는 $f^*\mathcal{A}$-가군이다. 준연접성은
당김과 환 변경에 의해 보존되므로 $h(\mathcal{G})$는 준연접이다.
또한 다음 함자적 사상이 있다.
$$
h(f_*\mathcal{F}) =
f^*f_*\mathcal{F} \otimes_{f^*f_*\mathcal{O}_X} \mathcal{O}_X \to \mathcal{F}
$$
이는 인접 사상 $f^*f_*\mathcal{F} \to \mathcal{F}$에서 온다.
그리고 다음 함자적 사상이 있다.
$$
\mathcal{G} \to f_*h(\mathcal{G}) \quad\text{adjoint to the map}\quad
f^*\mathcal{G} \to f^*\mathcal{G} \otimes_{f^*f_*\mathcal{O}_X} \mathcal{O}_X
$$
이는 국소 단면 $s$를 $f^*\mathcal{F}$에서 $s \otimes 1$로
보낸다.
증명을 끝내려면 위의 $\mathcal{F}$와 $\mathcal{G}$에 대해
이 사상들이 동형임을 보이면 충분하다. 이는 아핀 덮개의 원소들에서,
즉 $X$와 $S$가 아핀인 경우에 확인할 수 있다.

\medskip\noindent
이 논증을 가능하게 하는 핵심 대수 관찰은 다음과 같다.
$R \to A$를 환 사상이라 하고 $N$을 $A$-가군이라 하자. 그러면
$$
(N \otimes_R A) \otimes_{(A \otimes_R A)} A = N
$$
이다. 실제로 이는 함자 $h$를 준연접 $\widetilde{A}$-가군
$\widetilde{N}$에 적용하여 $\Spec(R)$ 위에서 얻은 결과이다.
세부사항은 생략한다.
\end{proof}

\begin{lemma}
\label{morphisms-lemma-composition-affine}
아핀 사상들의 합성은 아핀이다.
\end{lemma}

\begin{proof}
$f : X \to Y$와 $g : Y \to Z$를 아핀 사상이라 하자.
$U \subset Z$를 아핀 열린집합이라 하자. 가정에 의해
$g^{-1}(U)$는 $g$에 대한 가정에 의해 아핀이다. 따라서
$f^{-1}(g^{-1}(U))$도 $f$에 대한 가정에 의해 아핀이다.
그러므로 $(g \circ f)^{-1}(U)$는
아핀이다.
\end{proof}

\begin{lemma}
\label{morphisms-lemma-base-change-affine}
아핀 사상의 기저 변환은 아핀이다.
\end{lemma}

\begin{proof}
$f : X \to S$를 아핀 사상이라 하자. $S' \to S$를 임의의
사상이라 하자. $f' : X_{S'} = S' \times_S X \to S'$를
$f$의 기저 변환이라 하자.
모든 $s' \in S'$에 대해 어떤 아핀 열린 근방
$s' \in V \subset S'$가 존재하며, 이는 어떤 아핀 열린집합
$U \subset S$로 사상된다.
가정에 의해 $f^{-1}(U)$는 아핀이다. 스킴, 절
\ref{schemes-section-fibre-products}의 내용에 의해
$f^{-1}(U)_V = V \times_U f^{-1}(U)$가 아핀이고
$(f')^{-1}(V)$와 같음을 알 수 있다. 따라서 $S'$는
$f'$에 의한 역상이 아핀인 아핀 열린집합들로 열린 덮개를 갖는다.
위의 보조정리 \ref{morphisms-lemma-characterize-affine}에 의해 결론을 얻는다.
\end{proof}

\begin{lemma}
\label{morphisms-lemma-closed-immersion-affine}
닫힌 몰입은 아핀이다.
\end{lemma}

\begin{proof}
이 사실의 첫 단서는 스킴, 보조정리
\ref{schemes-lemma-closed-immersion-affine-case}이다.
완전한 서술은 스킴, 보조정리
\ref{schemes-lemma-closed-subspace-scheme}를 보라.
\end{proof}
\begin{lemma}
\label{morphisms-lemma-affine-s-open}
$X$를 스킴이라 하자.
$\mathcal{L}$을 가역 $\mathcal{O}_X$-가군이라 하자.
$s \in \Gamma(X, \mathcal{L})$라 하자.
포함 사상 $j : X_s \to X$는 아핀이다.
\end{lemma}

\begin{proof}
이는 성질, 보조정리
\ref{properties-lemma-affine-cap-s-open}과 정의에 따른다.
\end{proof}

\begin{lemma}
\label{morphisms-lemma-affine-permanence}
$g : X \to Y$가 $S$ 위의 스킴 사상이라고 하자.
\begin{enumerate}
\item $X$가 $S$ 위에서 아핀이고
$\Delta : Y \to Y \times_S Y$가 아핀이면 $g$는 아핀이다.
\item $X$가 $S$ 위에서 아핀이고 $Y$가 $S$ 위에서 분리되어 있으면
$g$는 아핀이다.
\item 아핀 스킴에서 아핀 대각선을 갖는 스킴으로 가는 사상은 아핀이다.
\item 아핀 스킴에서 분리된 스킴으로 가는 사상은 아핀이다.
\end{enumerate}
\end{lemma}

\begin{proof}
(1)의 증명. 기저 변환 $X \times_S Y \to Y$는
보조정리 \ref{morphisms-lemma-base-change-affine}에 의해 아핀이다.
사상 $(1, g) : X \to X \times_S Y$는
$Y \to Y \times_S Y$의 기저 변환이며,
이는 사상 $X \times_S Y \to Y \times_S Y$에 의해 취해진다.
따라서 이는 보조정리
\ref{morphisms-lemma-base-change-affine}에 의해 아핀이다.
아핀 사상들의 합성은 아핀이므로
(보조정리 \ref{morphisms-lemma-composition-affine} 참조) (1)이 따른다.
(2)는 (1)에서 따른다. 닫힌 몰입은 아핀이며
(보조정리 \ref{morphisms-lemma-closed-immersion-affine} 참조),
$Y/S$가 분리되었다는 것은 $\Delta$가 닫힌 몰입이라는 뜻이다.
(3)과 (4)는 (1)과 (2)의 특수한 경우이다.
\end{proof}

\begin{lemma}
\label{morphisms-lemma-morphism-affines-affine}
아핀 스킴들 사이의 사상은 아핀이다.
\end{lemma}

\begin{proof}
$S = \Spec(\mathbf{Z})$로 두고 보조정리
\ref{morphisms-lemma-affine-permanence}를 적용하면 즉시 따른다.
이는 보조정리 \ref{morphisms-lemma-characterize-affine}의
(1)과 (2)의 동치에서도 직접 따른다.
\end{proof}

\begin{lemma}
\label{morphisms-lemma-Artinian-affine}
$S$를 스킴이라 하자.
$A$를 아르틴 환이라 하자.
임의의 사상 $\Spec(A) \to S$는 아핀이다.
\end{lemma}

\begin{proof}
생략한다.
\end{proof}

\begin{lemma}
\label{morphisms-lemma-get-affine}
$j : Y \to X$를 스킴의 몰입이라 하자.
열린집합 $U \subset X$가 존재하며 그 여집합은
$Z = X \setminus U$이고,
다음을 만족한다고 하자.
\begin{enumerate}
\item $U \to X$는 아핀이다.
\item $j^{-1}(U) \to U$는 아핀이다.
\item $j(Y) \cap Z$는 닫혀 있다.
\end{enumerate}
그러면 $j$는 아핀이다. 특히 $X$가 아핀이면 $Y$도 아핀이다.
\end{lemma}

\begin{proof}
스킴, 정의 \ref{schemes-definition-immersion}에 의해
열린 부분스킴 $W \subset X$를 통해 $j$가 인수분해되어
닫힌 몰입 $i : Y \to W$ 다음에 포함 사상 $W \to X$가
오는 형태가 존재한다.
닫힌 몰입은 아핀이므로
(보조정리 \ref{morphisms-lemma-closed-immersion-affine}),
모든 $x \in W$에 대해 $x$의 $X$ 안의 아핀 열린 근방이 존재하고
그 근방의 $j$에 의한 역상은 아핀이다.
$x \in U$이면 가정 (2)에 의해 같은 사실이 성립한다.
이제 $x \in Z$이고 $x \not \in W$라고 하자.
그러면 $x \not \in j(Y) \cap Z$이다.
가정 (3)에 의해 $V \subset X$인 아핀 열린집합을 찾을 수 있으며,
이는 $x$를 포함하고 $j(Y) \cap Z$와 만나지 않는다.
그러면 $j^{-1}(V) = j^{-1}(V \cap U)$이고,
가정 (1)과 (2)에 의해 이는 아핀이다.
따라서 보조정리 \ref{morphisms-lemma-characterize-affine}에 의해
$j$는 아핀이다.
\end{proof}
\section{풍부한 가역 가군의 족}
\label{morphisms-section-families-ample-invertible-modules}

\noindent
가역 가군의 족이라는 개념에 관한 짧은 절이다.

\begin{definition}
\label{morphisms-definition-family-ample-invertible-modules}
\begin{reference}
\cite[II Definition 2.2.4]{SGA6}
\end{reference}
$X$를 스킴이라 하자. $\{\mathcal{L}_i\}_{i \in I}$를
가역 $\mathcal{O}_X$-가군들의 족이라 하자. 다음 조건을 만족하면
$\{\mathcal{L}_i\}_{i \in I}$를
{\it $X$ 위의 풍부한 가역 가군의 족}이라고 부른다.
\begin{enumerate}
\item $X$는 준콤팩트이다.
\item 모든 $x \in X$에 대해 어떤 $i \in I$, $n \geq 1$ 및
$s \in \Gamma(X, \mathcal{L}_i^{\otimes n})$이 존재하여
$x \in X_s$이고 $X_s$가 아핀이다.
\end{enumerate}
\end{definition}

\noindent
$\{\mathcal{L}_i\}_{i \in I}$가 스킴 $X$ 위의 풍부한 가역 가군의
족이면, 유한 부분집합 $I' \subset I$가 존재하여
$\{\mathcal{L}_i\}_{i \in I'}$도 $X$ 위의 풍부한 가역 가군의
족이 된다 (이는 준콤팩트성에서 즉시 따른다).
풍부한 가역 가군의 족을 갖는 스킴은 다음 보조정리에 의해 아핀
대각선을 가지며, 따라서 특히 준분리이다.

\begin{lemma}
\label{morphisms-lemma-affine-diagonal}
스킴 $X$에 대해 모든 점 $x \in X$에 대해 어떤 가역 $\mathcal{O}_X$-가군
$\mathcal{L}$과 전역 단면 $s \in \Gamma(X, \mathcal{L})$가 존재하여
$x \in X_s$이고 $X_s$가 아핀이라고 하자.
그러면 $X$의 대각선은 아핀 사상이다.
\end{lemma}

\begin{proof}
가역 $\mathcal{O}_X$-가군 $\mathcal{L}$, $\mathcal{M}$과 전역 단면
$s \in \Gamma(X, \mathcal{L})$, $t \in \Gamma(X, \mathcal{M})$이
주어져 $X_s$와 $X_t$가 아핀이라고 하자. $X_s \cap X_t$가
아핀임을 보여야 한다.
실제로 대각선 $\Delta : X \to X \times X$에 보조정리
\ref{morphisms-lemma-characterize-affine}를 적용하고
$\Delta^{-1}(X_s \times X_t) = X_s \cap X_t$임을 사용하면
$\Delta$가 아핀임을 얻는다. $X_s \cap X_t$가 아핀이라는 사실은
성질, 보조정리 \ref{properties-lemma-affine-cap-s-open}에 따른다.
\end{proof}

\begin{remark}
\label{morphisms-remark-affine-s-opens-cover-family}
성질, 보조정리
\ref{properties-lemma-affine-s-opens-cover-quasi-separated}에서
풍부한 가역 가군을 갖는 스킴은 분리된다는 것을 알 수 있다.
그러나 풍부한 가역 가군의 족을 갖는 스킴에 대해서는 이는 틀리다.
구체적으로, 스킴, 예
\ref{schemes-example-affine-space-zero-doubled}에 나타난 $X$에서
$n = 1$인 경우, 즉 영점이 두 개로 겹친 아핀 직선을 택하자.
그 예의 표기를 사용하되 $x$를 $x_1$ 대신 쓰고, $y$를 $y_1$ 대신
쓰기로 하자. 모든 정수 $n$에 대해 가역 층 $\mathcal{L}_n$이
$X$ 위에 존재하며, 이는 $X_1$과 $X_2$에서 자명하고
그 전이함수 $U_{12} \to U_{21}$는
$f(x) \mapsto y^n f(y)$이다.
$\mathcal{L}_n$의 전역 단면은
$(f(x), g(y)) \in k[x] \oplus
k[y]$인 쌍으로서 $y^n f(y) = g(y)$를 만족하는 것들이다.
단면 $s = (1,
y)$인 $\mathcal{L}_1$과 단면 $t = (x, 1)$인
$\mathcal{L}_{-1}$는
$X_s = X_1$이고 $X_t = X_2$이므로 열린 아핀 덮개를 결정한다.
따라서 $X$는 풍부한 가역 가군의 족을 갖지만 분리되어 있지 않다.
\end{remark}
\section{준아핀 사상}
\label{morphisms-section-quasi-affine}

\noindent
스킴 $X$가 준콤팩트이고 아핀 스킴의 열린 부분스킴과 동형이면
{\it 준아핀}이라고 부른다. 성질, 정의
\ref{properties-definition-quasi-affine}를 보라.

\begin{definition}
\label{morphisms-definition-quasi-affine}
스킴의 사상 $f : X \to S$가 {\it 준아핀}이라는 것은
$S$의 모든 아핀 열린집합의 역상이 준아핀 스킴인 경우를 말한다.
\end{definition}

\begin{lemma}
\label{morphisms-lemma-quasi-affine-separated}
준아핀 사상은 분리되고 준콤팩트이다.
\end{lemma}

\begin{proof}
$f : X \to S$를 준아핀이라 하자.
준콤팩트성은 스킴, 보조정리
\ref{schemes-lemma-quasi-compact-affine}에서 즉시 따른다.
$U \subset S$를 아핀 열린집합이라 하자.
$f^{-1}(U)$가 분리된 스킴임을 보일 수 있다면 $f$가 분리됨을
얻는다 (스킴, 보조정리
\ref{schemes-lemma-characterize-separated}에 의해 분리성은
기저 위에서 국소적이다).
가정에 의해 $f^{-1}(U)$는 아핀 스킴의 열린 부분스킴과 동형이다.
아핀 스킴은 분리되어 있으므로 아핀 스킴의 모든 열린 부분스킴도
분리되어 있고, 원하는 결론을 얻는다.
\end{proof}
\begin{lemma}
\label{morphisms-lemma-characterize-quasi-affine}
$f : X \to S$를 스킴의 사상이라 하자.
다음 조건들은 서로 동치이다.
\begin{enumerate}
\item 사상 $f$는 준아핀이다.
\item 아핀 열린 덮개 $S = \bigcup W_j$가 존재하여 각
$f^{-1}(W_j)$가 준아핀이 된다.
\item 준연접 $\mathcal{O}_S$-대수 층 $\mathcal{A}$와
준콤팩트 열린 몰입
$$
\xymatrix{
X \ar[rr] \ar[rd] & & \underline{\Spec}_S(\mathcal{A}) \ar[dl] \\
& S &
}
$$
이 $S$ 위에 존재한다.
\item (3)과 같지만 $\mathcal{A} = f_*\mathcal{O}_X$이고
수평 화살표가 구성, 보조정리
\ref{constructions-lemma-canonical-morphism}의 정준 사상인 경우.
\end{enumerate}
\end{lemma}

\begin{proof}
(1)이 (2)를 함의하고 (4)가 (3)을 함의한다는 것은 자명하다.

\medskip\noindent
$S = \bigcup_{j \in J} W_j$가 아핀 열린 덮개이고 각
$f^{-1}(W_j)$가 준아핀이라고 하자.
스킴, 보조정리 \ref{schemes-lemma-quasi-compact-affine}에 의해
$f$는 준콤팩트임을 알 수 있다. 스킴, 보조정리
\ref{schemes-lemma-characterize-quasi-separated}에 의해
사상 $f$는 준분리임을 알 수 있다. 따라서 스킴, 보조정리
\ref{schemes-lemma-push-forward-quasi-coherent}에 의해
층 $\mathcal{A} = f_*\mathcal{O}_X$는
$\mathcal{O}_X$-대수의 준연접 층이다. 그러므로
$g : Y = \underline{\Spec}_S(\mathcal{A}) \to S$인
$S$ 위의 스킴을 갖는다.
항등 사상
$\text{id} : \mathcal{A} = f_*\mathcal{O}_X \to f_*\mathcal{O}_X$는
상대 스펙트럼의 정의를 통해
$can : X \to Y$인 $S$ 위의 사상을 주며,
구성, 보조정리
\ref{constructions-lemma-canonical-morphism}을 보라.
가정과 방금 인용한 보조정리 및 성질, 보조정리
\ref{properties-lemma-quasi-affine}에 의해
제한 $can|_{f^{-1}(W_j)} : f^{-1}(W_j) \to g^{-1}(W_j)$는
준콤팩트 열린 몰입이다. 따라서 $can$은 준콤팩트 열린 몰입이다.
이로써 (2)가 (4)를 함의함을 보였다.

\medskip\noindent
(3)을 가정하자. 임의의 아핀 열린집합 $U \subset S$를 택하자.
구성, 보조정리 \ref{constructions-lemma-spec-properties}에 의해
상대 스펙트럼에서 $U$의 역상은 아핀이다.
따라서 $f^{-1}(U)$는 준아핀이라고 결론내린다
((3)에는 준콤팩트성도 포함되어 있음에 유의하라).
그러므로 (3)이 (1)을 함의한다.
\end{proof}
\begin{lemma}
\label{morphisms-lemma-composition-quasi-affine}
준아핀 사상들의 합성은 준아핀이다.
\end{lemma}

\begin{proof}
$f : X \to Y$와 $g : Y \to Z$를 준아핀 사상이라 하자.
$U \subset Z$를 아핀 열린집합이라 하자. 그러면
$g^{-1}(U)$는 $g$에 대한 가정에 의해 준아핀이다.
$j : g^{-1}(U) \to V$를 아핀 스킴 $V$로의
준콤팩트 열린 몰입이라 하자.
위의 보조정리 \ref{morphisms-lemma-characterize-quasi-affine}에 의해
$f^{-1}(g^{-1}(U))$는 상대 스펙트럼
$\underline{\Spec}_{g^{-1}(U)}(\mathcal{A})$의 준콤팩트 열린
부분스킴이며, 이는 어떤 준연접
$\mathcal{O}_{g^{-1}(U)}$-대수 층 $\mathcal{A}$에 대한 것이다.
스킴, 보조정리 \ref{schemes-lemma-push-forward-quasi-coherent}에 의해
층 $\mathcal{A}' = j_*\mathcal{A}$는
$\mathcal{O}_V$-대수의 준연접 층이고
$j^*\mathcal{A}' = \mathcal{A}$를 만족한다.
따라서 다음 가환 도표를 얻는다.
$$
\xymatrix{
f^{-1}(g^{-1}(U)) \ar[r] &
\underline{\Spec}_{g^{-1}(U)}(\mathcal{A})
\ar[r] \ar[d] &
\underline{\Spec}_V(\mathcal{A}') \ar[d] \\
& g^{-1}(U) \ar[r]^j & V
}
$$
이때 정사각형은 스킴의 올곱이다.
구성, 보조정리 \ref{constructions-lemma-spec-properties}를 보라.
오른쪽 위 꼭짓점은 아핀 스킴임에 유의하라.
그러므로 $(g \circ f)^{-1}(U)$는 준아핀이다.
\end{proof}
\begin{lemma}
\label{morphisms-lemma-base-change-quasi-affine}
준아핀 사상의 기저 변환은 준아핀이다.
\end{lemma}

\begin{proof}
$f : X \to S$를 준아핀 사상이라 하자.
위의 보조정리 \ref{morphisms-lemma-characterize-quasi-affine}에 의해
$\mathcal{O}_S$-대수의 준연접 층 $\mathcal{A}$와
$X \to \underline{\Spec}_S(\mathcal{A})$인 준콤팩트 열린 몰입을
$S$ 위에서 찾을 수 있다.
$g : S' \to S$를 임의의 사상이라 하자.
$f' : X_{S'} = S' \times_S X \to S'$를 $f$의
기저 변환이라 하자.
준콤팩트 열린 몰입의 기저 변환은 준콤팩트 열린 몰입이므로,
$X_{S'} \to \underline{\Spec}_{S'}(g^*\mathcal{A})$가
준콤팩트 열린 몰입임을 알 수 있다
(스킴, 보조정리
\ref{schemes-lemma-quasi-compact-preserved-base-change}와
\ref{schemes-lemma-base-change-immersion}, 그리고 구성,
보조정리 \ref{constructions-lemma-spec-properties}를 사용했다).
다시 보조정리 \ref{morphisms-lemma-characterize-quasi-affine}에 의해
$X_{S'} \to S'$가 준아핀이라고 결론내린다.
\end{proof}

\begin{lemma}
\label{morphisms-lemma-quasi-compact-immersion-quasi-affine}
준콤팩트 몰입은 준아핀이다.
\end{lemma}

\begin{proof}
$X \to S$를 준콤팩트 몰입이라 하자.
모든 아핀 열린집합의 역상이 준아핀임을 보여야 한다.
따라서 $S$가 아핀 스킴이라고 가정하면 $X$가 준아핀임을 보여야 한다.
보조정리 \ref{morphisms-lemma-quasi-compact-immersion}에 의해
사상 $X \to S$는 $X \to Z \to S$로 인수분해되며,
여기서 $Z$는 $S$의 닫힌 부분스킴이고 $X \subset Z$는
준콤팩트 열린집합이다.
$S$가 아핀이므로 보조정리 \ref{morphisms-lemma-closed-immersion}에 의해
$Z$는 아핀이다. 그러므로 증명되었다.
\end{proof}

\begin{lemma}
\label{morphisms-lemma-affine-quasi-affine}
$S$를 스킴이라 하자. $X$를 아핀 스킴이라 하자.
사상 $f : X \to S$가 준아핀인 것은 그것이 준콤팩트인 것과 동치이다.
특히 아핀 스킴에서 준분리 스킴으로 가는 임의의 사상은 준아핀이다.
\end{lemma}

\begin{proof}
$V \subset S$를 아핀 열린집합이라 하자.
그러면 $f^{-1}(V)$는 아핀 스킴 $X$의 열린 부분스킴이므로
준아핀인 것은 준콤팩트인 것과 동치이다. 이것이 첫째 주장을 증명한다.
임의의 $f : X \to S$가 준콤팩트라는 사실은 $X$가 아핀이고
$S$가 준분리일 때 스킴, 보조정리
\ref{schemes-lemma-quasi-compact-permanence}를
$X \to S \to \Spec(\mathbf{Z})$에 적용하면 따른다.
\end{proof}

\begin{lemma}
\label{morphisms-lemma-quasi-affine-permanence}
$g : X \to Y$가 $S$ 위의 스킴 사상이라고 하자.
$X$가 $S$ 위에서 준아핀이고 $Y$가 $S$ 위에서 준분리이면
$g$는 준아핀이다. 특히 준아핀 스킴에서 준분리 스킴으로 가는
임의의 사상은 준아핀이다.
\end{lemma}

\begin{proof}
기저 변환 $X \times_S Y \to Y$는 보조정리
\ref{morphisms-lemma-base-change-quasi-affine}에 의해 준아핀이다.
사상 $X \to X \times_S Y$는 $Y \to S$가 준분리이므로
준콤팩트 몰입이다. 스킴, 보조정리
\ref{schemes-lemma-section-immersion}을 보라.
준콤팩트 몰입은 보조정리
\ref{morphisms-lemma-quasi-compact-immersion-quasi-affine}에 의해
준아핀이고, 준아핀 사상들의 합성은 준아핀이다
(보조정리 \ref{morphisms-lemma-composition-quasi-affine} 참조).
따라서 증명되었다.
\end{proof}
\section{환 사상의 성질로 정의되는 사상의 유형}
\label{morphisms-section-properties-ring-maps}

\noindent
이 절에서는 환 사상의 어떤 성질이 스킴 사상의 국소적 성질을
정의하게 하는지 연구한다.

\begin{definition}
\label{morphisms-definition-property-local}
$P$를 환 사상의 성질이라 하자.
\begin{enumerate}
\item 다음 조건들이 성립하면 $P$를 {\it 국소적}이라고 부른다.
\begin{enumerate}
\item 임의의 환 사상 $R \to A$와 임의의 $f \in R$에 대해
$P(R \to A) \Rightarrow P(R_f \to A_f)$이다.
\item 임의의 환 $R$, $A$, 임의의 $f \in R$, $a\in A$ 및 임의의
환 사상 $R_f \to A$에 대해
$P(R_f \to A) \Rightarrow P(R \to A_a)$이다.
\item 환 사상 $R \to A$와 $a_i \in A$가 주어지고
$(a_1, \ldots, a_n) = A$를 만족하면
$\forall i, P(R \to A_{a_i}) \Rightarrow P(R \to A)$이다.
\end{enumerate}
\item $P$를 {\it 기저 변환에 대해 안정적}이라고 부르는 것은
임의의 환 사상 $R \to A$, $R \to R'$에 대해
$P(R \to A) \Rightarrow P(R' \to R' \otimes_R A)$가 성립하는
경우이다.
\item $P$를 {\it 합성에 대해 안정적}이라고 부르는 것은
임의의 환 사상 $A \to B$, $B \to C$에 대해
$P(A \to B) \wedge P(B \to C) \Rightarrow P(A \to C)$가
성립하는 경우이다.
\end{enumerate}
\end{definition}

\begin{definition}
\label{morphisms-definition-locally-P}
$P$를 환 사상의 성질이라 하자.
$f : X \to S$를 스킴의 사상이라 하자.
사상 $f$가 {\it $P$ 유형 국소적}이라고 부르는 것은 모든
$x \in X$에 대해 아핀 열린 근방 $U$가 $x$를 포함하고,
$X$에서 이 열린집합의 상이 아핀 열린집합 $V \subset S$에 포함되며,
유도된 환 사상
$\mathcal{O}_S(V) \to \mathcal{O}_X(U)$가 성질 $P$를 갖는
경우이다.
\end{definition}
\noindent
이는 성질 $P$가 국소적이지 않으면 “좋은” 정의가 아니다.
$P$가 국소적인 성질이라 하더라도, 다른 곳에서 정의를 명시하지
않는 한 사상이 “$P$ 유형 국소적”이라고 말하기 위해 이 정의를
자동으로 사용하지는 않는다.

\begin{lemma}
\label{morphisms-lemma-locally-P}
$f : X \to S$를 스킴의 사상이라 하자.
$P$를 환 사상의 성질이라 하자.
$U$를 $X$의 아핀 열린집합, $V$를 $S$의 아핀 열린집합으로서
$f(U) \subset V$가 되게 하자.
$f$가 $P$ 유형 국소적이고 $P$가 국소적이면
$P(\mathcal{O}_S(V) \to \mathcal{O}_X(U))$가 성립한다.
\end{lemma}

\begin{proof}
$f$가 $P$ 유형 국소적이므로 모든 $u \in U$에 대해
$U_u \subset X$인 아핀 열린집합이 존재하고, 이는
$V_u \subset S$인 아핀 열린집합으로 사상되며
$P(\mathcal{O}_S(V_u) \to \mathcal{O}_X(U_u))$가 성립한다.
$U'_u \subset U \cap U_u$인 $u$의 근방을 택하되,
이는 $U$와 $U_u$ 모두에서 표준 아핀 열린집합이 되게 하자.
스킴, 보조정리
\ref{schemes-lemma-standard-open-two-affines}를 보라.
정의 \ref{morphisms-definition-property-local} (1)(b)에 의해
$P(\mathcal{O}_S(V_u) \to \mathcal{O}_X(U'_u))$가 성립한다.
따라서 $U_u \subset U$가 표준 아핀 열린집합이라고 가정할 수 있다.
$V'_u \subset V \cap V_u$인 $f(u)$의 근방을 택하되,
이는 $V$와 $V_u$ 모두에서 표준 아핀 열린집합이 되게 하자.
스킴, 보조정리
\ref{schemes-lemma-standard-open-two-affines}를 보라.
그러면 $U'_u = f^{-1}(V'_u) \cap U_u$는 $U_u$ (따라서 $U$)의
표준 아핀 열린집합이고,
정의 \ref{morphisms-definition-property-local} (1)(a)에 의해
$P(\mathcal{O}_S(V'_u) \to \mathcal{O}_X(U'_u))$가 성립한다.
따라서 $U_u \subset U$와 $V_u \subset V$가 모두 표준 아핀
열린집합이라고 가정할 수 있다. 정의
\ref{morphisms-definition-property-local} (1)(b)를 한 번 더 적용하면
$P(\mathcal{O}_S(V) \to \mathcal{O}_X(U_u))$가 성립한다.
$U$가 준콤팩트이므로
$$
U = U_{u_1} \cup \ldots \cup U_{u_n}.
$$
가 되도록 유한 개의 점 $u_1, \ldots, u_n \in U$를 택할 수 있다.
정의 \ref{morphisms-definition-property-local} (1)(c)에 의해
$P(\mathcal{O}_S(V) \to \mathcal{O}_X(U))$가 성립한다.
\end{proof}
\begin{lemma}
\label{morphisms-lemma-locally-P-characterize}
$P$를 환 사상의 국소적 성질이라 하자.
$f : X \to S$를 스킴의 사상이라 하자.
다음 조건들은 서로 동치이다.
\begin{enumerate}
\item 사상 $f$는 $P$ 유형 국소적이다.
\item 모든 아핀 열린집합 $U \subset X$, $V \subset S$에 대해
$f(U) \subset V$이면
$P(\mathcal{O}_S(V) \to \mathcal{O}_X(U))$가 성립한다.
\item 열린 덮개 $S = \bigcup_{j \in J} V_j$와
열린 덮개 $f^{-1}(V_j) = \bigcup_{i \in I_j} U_i$가 존재하여
각 사상 $U_i \to V_j$, $j\in J, i\in I_j$가
$P$ 유형 국소적이다.
\item 아핀 열린 덮개 $S = \bigcup_{j \in J} V_j$와
아핀 열린 덮개 $f^{-1}(V_j) = \bigcup_{i \in I_j} U_i$가 존재하여
성질 $P(\mathcal{O}_S(V_j) \to \mathcal{O}_X(U_i))$가
모든 $j\in J, i\in I_j$에 대해 성립한다.
\end{enumerate}
더 나아가 $f$가 $P$ 유형 국소적이면 임의의 열린 부분스킴
$U \subset X$, $V \subset S$에 대해 $f(U) \subset V$일 때 제한
$f|_U : U \to V$도 $P$ 유형 국소적이다.
\end{lemma}

\begin{proof}
이는 위의 보조정리 \ref{morphisms-lemma-locally-P}에 따른다.
\end{proof}
\begin{lemma}
\label{morphisms-lemma-composition-type-P}
$P$를 환 사상의 성질이라 하자.
$P$가 국소적이고 합성에 대해 안정적이라고 가정하자.
$P$ 유형 국소적 사상들의 합성은
$P$ 유형 국소적이다.
\end{lemma}

\begin{proof}
$f : X \to Y$와 $g : Y \to Z$를 $P$ 유형 국소적 사상이라 하자.
$x \in X$를 택하자. $W \subset Z$를
$g(f(x))$의 아핀 열린 근방으로 택하자.
$V \subset g^{-1}(W)$를 $f(x)$의 아핀 열린 근방으로 택하자.
$U \subset f^{-1}(V)$를 $x$의 아핀 열린 근방으로 택하자.
보조정리 \ref{morphisms-lemma-locally-P-characterize}에 의해 환 사상
$\mathcal{O}_Z(W) \to \mathcal{O}_Y(V)$와
$\mathcal{O}_Y(V) \to \mathcal{O}_X(U)$는 모두 $P$를 만족한다.
따라서 합성에 대해 안정적이라는 가정에 의해
$\mathcal{O}_Z(W) \to \mathcal{O}_X(U)$도 $P$를 만족한다.
이는 $P$가 합성에 대해 안정적이라는 가정에 따른다.
\end{proof}

\begin{lemma}
\label{morphisms-lemma-base-change-type-P}
$P$를 환 사상의 성질이라 하자.
$P$가 국소적이고 기저 변환에 대해 안정적이라고 가정하자.
$P$ 유형 국소적 사상의 기저 변환은
$P$ 유형 국소적이다.
\end{lemma}

\begin{proof}
$f : X \to S$를 $P$ 유형 국소적 사상이라 하자.
$S' \to S$를 임의의 사상이라 하자. 다음과 같이 놓는다.
$f' : X_{S'} = S' \times_S X \to S'$를 $f$의
기저 변환이라 하자.
모든 $s' \in S'$에 대해 $s' \in V' \subset S'$인 아핀 열린 근방이
존재하며, 이는 어떤 아핀 열린집합 $V \subset S$로 사상된다.
보조정리 \ref{morphisms-lemma-locally-P-characterize}에 의해 열린집합
$f^{-1}(V)$는 아핀 열린집합 $U_i$들의 합집합이며,
환 사상 $\mathcal{O}_S(V) \to \mathcal{O}_X(U_i)$가 모두
$P$를 만족한다.
스킴, 절 \ref{schemes-section-fibre-products}의 내용에 의해
$f^{-1}(U)_{V'} = V' \times_V f^{-1}(V)$는
아핀 열린집합 $V' \times_V U_i$들의 합집합임을 알 수 있다.
그리고
$\mathcal{O}_{X_{S'}}(V' \times_V U_i) =
\mathcal{O}_{S'}(V') \otimes_{\mathcal{O}_S(V)} \mathcal{O}_X(U_i)$
이므로 환 사상
$\mathcal{O}_{S'}(V') \to \mathcal{O}_{X_{S'}}(V' \times_V U_i)$는
기저 변환에 대해 안정적이라는 $P$의 가정에 의해 $P$를 만족한다.
\end{proof}
\begin{lemma}
\label{morphisms-lemma-properties-local}
환 사상 $R \to A$의 다음 성질들은 국소적이다.
\begin{enumerate}
\item (국소환에서의 동형.)
$\mathfrak q$가 $A$의 소 아이디얼이고
$\mathfrak p \subset R$ 위에 놓일 때, 환 사상 $R \to A$는
동형
$R_{\mathfrak p} \to A_{\mathfrak q}$를 유도한다.
\item (열린 몰입.)
$\mathfrak q$가 $A$의 모든 소 아이디얼일 때 어떤 $f \in R$가
존재하여 $\varphi(f) \not \in \mathfrak q$이고,
환 사상 $\varphi : R \to A$가 동형
$R_f \to A_f$를 유도한다.
\item (기약인 올.)
$\mathfrak p$가 $R$의 모든 소 아이디얼일 때 올 환
$A \otimes_R \kappa(\mathfrak p)$는 기약이다.
\item (차원이 $n$ 이하인 올.)
$\mathfrak p$가 $R$의 모든 소 아이디얼일 때 올 환
$A \otimes_R \kappa(\mathfrak p)$의 크룰 차원은 $n$ 이하이다.
\item (표적 위에서 국소적으로 뇌터.)
환 사상 $R \to A$가 $A$가 뇌터 환이라는 성질을 갖는다.
\item 필요에 따라 더 추가한다\footnote{단, 다른 곳에서 별도로
이미 다루어진 성질은 제외한다.}.
\end{enumerate}
\end{lemma}

\begin{proof}
생략한다.
\end{proof}

\begin{lemma}
\label{morphisms-lemma-properties-base-change}
환 사상의 다음 성질들은 기저 변환에 대해 안정적이다.
\begin{enumerate}
\item (국소환에서의 동형.)
$\mathfrak q$가 $A$의 소 아이디얼이고
$\mathfrak p \subset R$ 위에 놓일 때, 환 사상 $R \to A$는
동형 $R_{\mathfrak p} \to A_{\mathfrak q}$를 유도한다.
\item (열린 몰입.)
$\mathfrak q$가 $A$의 모든 소 아이디얼일 때 어떤 $f \in R$가
존재하여 $\varphi(f) \not \in \mathfrak q$이고,
환 사상 $\varphi : R \to A$가 동형 $R_f \to A_f$를 유도한다.
\item 필요에 따라 더 추가한다\footnote{단, 다른 곳에서 별도로
이미 다루어진 성질은 제외한다.}.
\end{enumerate}
\end{lemma}

\begin{proof}
생략한다.
\end{proof}

\begin{lemma}
\label{morphisms-lemma-properties-composition}
환 사상의 다음 성질들은 합성에 대해 안정적이다.
\begin{enumerate}
\item (국소환에서의 동형.)
$\mathfrak q$가 $A$의 소 아이디얼이고
$\mathfrak p \subset R$ 위에 놓일 때, 환 사상 $R \to A$는
동형 $R_{\mathfrak p} \to A_{\mathfrak q}$를 유도한다.
\item (열린 몰입.)
$\mathfrak q$가 $A$의 모든 소 아이디얼일 때 어떤 $f \in R$가
존재하여 $\varphi(f) \not \in \mathfrak q$이고,
환 사상 $\varphi : R \to A$가 동형 $R_f \to A_f$를 유도한다.
\item (표적 위에서 국소적으로 뇌터.)
환 사상 $R \to A$가 $A$가 뇌터 환이라는 성질을 갖는다.
\item 필요에 따라 더 추가한다\footnote{단, 다른 곳에서 별도로
이미 다루어진 성질은 제외한다.}.
\end{enumerate}
\end{lemma}

\begin{proof}
생략한다.
\end{proof}
\section{유한형 사상}
\label{morphisms-section-finite-type}

\noindent
환 사상 $R \to A$가 유한형이라고 하는 것은 $A$가
$R[x_1, \ldots, x_n]$의 몫과 동형인 $R$-대수라는 뜻이다. 대수,
정의 \ref{algebra-definition-finite-type}을 보라.

\begin{definition}
\label{morphisms-definition-finite-type}
$f : X \to S$를 스킴의 사상이라 하자.
\begin{enumerate}
\item $f$가 {\it $x \in X$에서 유한형}이라고 하는 것은
아핀 열린집합 $\Spec(A) = U \subset X$가 $x$를 포함하고,
아핀 열린집합 $\Spec(R) = V \subset S$가 존재하며 $f(U) \subset V$이고
존재하여 유도된 환 사상 $R \to A$가 유한형인 경우이다.
\item $f$가 {\it 국소적으로 유한형}이라고 하는 것은
$X$의 모든 점에서 유한형인 경우이다.
\item $f$가 {\it 유한형}이라고 하는 것은 국소적으로 유한형이고
준콤팩트인 경우이다.
\end{enumerate}
\end{definition}

\begin{lemma}
\label{morphisms-lemma-locally-finite-type-characterize}
$f : X \to S$를 스킴의 사상이라 하자.
다음 조건들은 서로 동치이다.
\begin{enumerate}
\item 사상 $f$는 국소적으로 유한형이다.
\item 모든 아핀 열린집합 $U \subset X$, $V \subset S$에 대해
$f(U) \subset V$이면 환 사상
$\mathcal{O}_S(V) \to \mathcal{O}_X(U)$가 유한형이다.
\item 열린 덮개 $S = \bigcup_{j \in J} V_j$와
열린 덮개 $f^{-1}(V_j) = \bigcup_{i \in I_j} U_i$가 존재하여
각 사상 $U_i \to V_j$, $j\in J, i\in I_j$가
국소적으로 유한형이다.
\item 아핀 열린 덮개 $S = \bigcup_{j \in J} V_j$와
아핀 열린 덮개 $f^{-1}(V_j) = \bigcup_{i \in I_j} U_i$가 존재하여
환 사상 $\mathcal{O}_S(V_j) \to \mathcal{O}_X(U_i)$가
모든 $j\in J, i\in I_j$에 대해 유한형이다.
\end{enumerate}
더 나아가 $f$가 국소적으로 유한형이면 임의의 열린 부분스킴
$U \subset X$, $V \subset S$에 대해 $f(U) \subset V$일 때 제한
$f|_U : U \to V$도 국소적으로 유한형이다.
\end{lemma}

\begin{proof}
이는 다음 사실을 보이면 위의 보조정리
\ref{morphisms-lemma-locally-P}에서 따른다. 즉 성질
“$R \to A$가 유한형이다”가 국소적임을 보이자.
정의 \ref{morphisms-definition-property-local}의 조건 (a), (b), (c)를 확인한다.
대수, 보조정리 \ref{algebra-lemma-base-change-finiteness}에 의해
유한형이라는 성질은 기저 변환에 대해 안정적이므로 (a)가 성립한다.
대수, 보조정리 \ref{algebra-lemma-compose-finite-type}에 의해
유한형이라는 성질은 합성에 대해 안정적이고, 임의의 환
$R$에 대해 환 사상 $R \to R_f$가 유한형인 것은 자명하다.
따라서 (b)가 성립한다. 마지막으로 (c)는 대수, 보조정리
\ref{algebra-lemma-cover-upstairs}에 따른다.
\end{proof}
\begin{lemma}
\label{morphisms-lemma-composition-finite-type}
국소적으로 유한형인 두 사상의 합성은 국소적으로 유한형이다.
유한형 사상들에 대해서도 같은 결론이 성립한다.
\end{lemma}

\begin{proof}
보조정리 \ref{morphisms-lemma-locally-finite-type-characterize}의 증명에서
유한형이라는 것이 환 사상의 국소적 성질임을 보였다.
따라서 첫째 주장은 보조정리
\ref{morphisms-lemma-composition-type-P}와 유한형이 합성에 대해 안정적인
환 사상의 성질이라는 사실을 결합하면 따른다. 대수, 보조정리
\ref{algebra-lemma-compose-finite-type}을 보라.
위의 결과와 준콤팩트 사상들의 합성이 준콤팩트라는 사실
(스킴, 보조정리 \ref{schemes-lemma-composition-quasi-compact} 참조)에
의해 유한형 사상들의 합성은 유한형이다.
\end{proof}

\begin{lemma}
\label{morphisms-lemma-base-change-finite-type}
국소적으로 유한형인 사상의 기저 변환은 국소적으로 유한형이다.
유한형 사상들에 대해서도 같은 결론이 성립한다.
\end{lemma}

\begin{proof}
보조정리 \ref{morphisms-lemma-locally-finite-type-characterize}의 증명에서
유한형이라는 것이 환 사상의 국소적 성질임을 보였다.
따라서 첫째 주장은 보조정리
\ref{morphisms-lemma-base-change-type-P}와 유한형이 기저 변환에 대해
안정적인 환 사상의 성질이라는 사실을 결합하면 따른다. 대수,
보조정리 \ref{algebra-lemma-base-change-finiteness}를 보라.
위의 결과와 준콤팩트 사상의 기저 변환이 준콤팩트라는 사실
(스킴, 보조정리
\ref{schemes-lemma-quasi-compact-preserved-base-change} 참조)에 의해
유한형 사상의 기저 변환은 유한형 사상이다.
\end{proof}

\begin{lemma}
\label{morphisms-lemma-immersion-locally-finite-type}
닫힌 몰입은 유한형이다.
몰입은 국소적으로 유한형이다.
\end{lemma}

\begin{proof}
이는 열린 몰입이 국소 동형이고 닫힌 몰입이 자명하게 유한형이기
때문이다.
\end{proof}

\begin{lemma}
\label{morphisms-lemma-finite-type-noetherian}
$f : X \to S$를 사상이라 하자.
$S$가 (국소적으로) 뇌터이고 $f$가 (국소적으로) 유한형이면
$X$는 (국소적으로) 뇌터이다.
\end{lemma}

\begin{proof}
이는 뇌터 환 위의 유한형 환이 뇌터라는 사실에서 즉시 따른다.
대수, 보조정리 \ref{algebra-lemma-Noetherian-permanence}를 보라.
(또한 준콤팩트 사상의 표적이 준콤팩트이면 그 원천이
준콤팩트라는 사실을 사용하라.)
\end{proof}
\begin{lemma}
\label{morphisms-lemma-finite-type-Noetherian-quasi-separated}
$f : X \to S$를 국소적으로 유한형 사상이고 $S$가 국소적으로
뇌터라고 하자. 그러면 $f$는 준분리이다.
\end{lemma}

\begin{proof}
사실 $X$가 준분리라는 것도 성립한다. 성질, 보조정리
\ref{properties-lemma-locally-Noetherian-quasi-separated}와
위의 보조정리 \ref{morphisms-lemma-finite-type-noetherian}을 보라.
그런 다음 스킴, 보조정리
\ref{schemes-lemma-compose-after-separated}를 적용하면
$f$가 준분리임을 결론내릴 수 있다.
\end{proof}

\begin{lemma}
\label{morphisms-lemma-permanence-finite-type}
$X \to Y$를 기저 스킴 $S$ 위의 스킴 사상이라 하자.
$X$가 $S$ 위에서 국소적으로 유한형이면,
$X \to Y$는 국소적으로 유한형이다.
\end{lemma}

\begin{proof}
보조정리 \ref{morphisms-lemma-locally-finite-type-characterize}에 의해
이는 다음 대수적 사실로 옮겨진다. 환 사상
$A \to B \to C$가 주어지고 $A \to C$가 유한형이면
$B \to C$도 유한형이다.
(대수, 보조정리 \ref{algebra-lemma-compose-finite-type}를 보라.)
\end{proof}
\section{유한형 점과 야콥슨 스킴}
\label{morphisms-section-points-finite-type}

\noindent
$S$를 스킴이라 하자. 점 $s$를 $S$의 유한형 점이라 하자.
사상 $\Spec(\kappa(s)) \to S$가 유한형인 점을 말한다.
이를 연구하는 이유는 어떤 의미와 어떤 상황에서 유한형 점이
닫힌 점을 대신할 수 있기 때문이다. 예를 들어 항상 충분히 많은
유한형 점이 존재한다. 더 나아가 모든 유한형 점이 닫힌 점일
필요충분조건은 스킴이 야콥슨인 것이다.

\begin{lemma}
\label{morphisms-lemma-point-finite-type}
$S$를 스킴이라 하자. $k$를 체라 하자.
$f : \Spec(k) \to S$를 사상이라 하자.
다음 조건들은 서로 동치이다.
\begin{enumerate}
\item 사상 $f$는 유한형이다.
\item 사상 $f$는 국소적으로 유한형이다.
\item 아핀 열린집합 $U = \Spec(R)$가 $S$ 안에 존재하여
$f$가 유한 환 사상 $R \to k$에 대응한다.
\item 아핀 열린집합 $U = \Spec(R)$가 $S$ 안에 존재하여
$f$의 상이 닫힌점 $u$ 하나로 이루어지고 그 점이 $U$ 안에 있으며,
체의 확장 $k/\kappa(u)$가 유한이다.
\end{enumerate}
\end{lemma}

\begin{proof}
(1)과 (2)의 동치는 자명하다. $\Spec(k)$가 한 점만 가지므로
그로부터 나오는 모든 사상은 준콤팩트이기 때문이다.

\medskip\noindent
$f$가 국소적으로 유한형이라고 하자. 그 사상의 상이 포함되는
아핀 열린집합 $\Spec(R) = U \subset S$를 하나 택하자. $f$의 상이
이 열린집합 $U$ 안에 포함되고, 환 사상 $R \to k$가 유한형이라고 하자.
핵을 $\mathfrak p \subset R$라 하자. 그러면
$R/\mathfrak p \subset k$는 유한형이다. 대수, 보조정리
\ref{algebra-lemma-field-finite-type-over-domain}에 의해
$\overline{f} \in R/\mathfrak p$를 택할 수 있어
$(R/\mathfrak p)_{\overline{f}}$가 체이고
$(R/\mathfrak p)_{\overline{f}} \to k$가 유한 체 확장이 된다.
$f \in R$가 $\overline{f}$의 올림이면
$k$가 유한 $R_f$-가군임을 알 수 있다. 따라서
(2) $\Rightarrow$ (3)이다.

\medskip\noindent
$\Spec(R) = U \subset S$가 아핀 열린집합이고
$f$가 유한 환 사상 $R \to k$에 대응한다고 하자.
그러면 보조정리 \ref{morphisms-lemma-locally-finite-type-characterize}에 의해
$f$는 국소적으로 유한형이다. 따라서 (3) $\Rightarrow$ (2)이다.

\medskip\noindent
$R \to k$가 유한이라고 하자. $R \to k$의 상은
그 위에서 $k$가 유한인 체이며, 이는 대수, 보조정리
\ref{algebra-lemma-integral-under-field}에 따른다.
따라서 $R \to k$의 핵은 극대 아이디얼이다.
그러므로 (3) $\Rightarrow$ (4)이다.

\medskip\noindent
함의 (4) $\Rightarrow$ (3)은 즉시 따른다.
\end{proof}
\begin{lemma}
\label{morphisms-lemma-artinian-finite-type}
$S$를 스킴이라 하자.
$A$를 잉여체가 $\kappa$인 아르틴 국소환이라 하자.
$f : \Spec(A) \to S$를 스킴의 사상이라 하자.
그러면 $f$가 유한형인 것은 합성
$\Spec(\kappa) \to \Spec(A) \to S$가 유한형인 것과 동치이다.
\end{lemma}

\begin{proof}
사상 $\Spec(\kappa) \to \Spec(A)$는 유한형이므로,
$f$가 유한형이면 합성
$\Spec(\kappa) \to S$도 유한형이라는 것은 자명하다
(보조정리 \ref{morphisms-lemma-composition-finite-type}를 보라).
역을 보이려면, $\Spec(A) \to S$의 상이 $U = \Spec(B)$인
$S$의 어떤 아핀 열린집합 안에 들어간다는 것을, $\Spec(A)$가
한 점만 가지므로 관찰하자.
마지막으로 대수, 보조정리
\ref{algebra-lemma-essentially-of-finite-type-into-artinian-local}를
$B \to A$에 적용하면 된다.
\end{proof}

\noindent
스킴의 점 $s$가 속한 $S$가 주어지면 표준 사상
$\Spec(\kappa(s)) \to S$가 존재함을 상기하자. 이는
스킴, 절 \ref{schemes-section-points}를 보라.

\begin{definition}
\label{morphisms-definition-finite-type-point}
$S$를 스킴이라 하자.
점 $s$가 $S$의 {\it 유한형 점}이라고 하자.
이는 표준 사상 $\Spec(\kappa(s)) \to S$가 유한형인 경우이다.
$S_{\text{ft-pts}}$를 $S$의 유한형 점들의 집합으로 표시한다.
\end{definition}

\noindent
유한형 점들의 집합은 다음과 같이 기술할 수 있다.

\begin{lemma}
\label{morphisms-lemma-identify-finite-type-points}
$S$를 스킴이라 하자. 다음이 성립한다.
$$
S_{\text{ft-pts}} = \bigcup\nolimits_{U \subset S\text{ open}} U_0
$$
여기서 $U_0$는 $U$의 닫힌점들의 집합이다.
여기서 $U$를 모든 열린집합에 걸쳐 취해도 되고
$S$의 모든 아핀 열린집합에 걸쳐 취해도 된다.
\end{lemma}

\begin{proof}
이는 보조정리 \ref{morphisms-lemma-point-finite-type}에서 즉시 따른다.
\end{proof}

\begin{lemma}
\label{morphisms-lemma-finite-type-points-morphism}
$f : T \to S$를 스킴의 사상이라 하자.
$f$가 국소적으로 유한형이면
$f(T_{\text{ft-pts}}) \subset S_{\text{ft-pts}}$이다.
\end{lemma}

\begin{proof}
$T$가 체의 스펙트럼이면 이는 보조정리
\ref{morphisms-lemma-point-finite-type}이다.
일반적인 경우에는 국소적으로 유한형 사상들의 합성이
국소적으로 유한형이라는 사실(보조정리
\ref{morphisms-lemma-composition-finite-type})으로부터 따른다.
\end{proof}

\begin{lemma}
\label{morphisms-lemma-finite-type-points-surjective-morphism}
$f : T \to S$를 스킴의 사상이라 하자.
$f$가 국소적으로 유한형이고 전사이면
$f(T_{\text{ft-pts}}) = S_{\text{ft-pts}}$이다.
\end{lemma}

\begin{proof}
보조정리 \ref{morphisms-lemma-finite-type-points-morphism}에 의해
$f(T_{\text{ft-pts}}) \subset S_{\text{ft-pts}}$이다.
$s \in S$를 유한형 점이라 하자. $f$가 전사이므로 스킴
$T_s = \Spec(\kappa(s)) \times_S T$는 공집합이 아니다.
따라서 보조정리 \ref{morphisms-lemma-identify-finite-type-points}에 의해
이 스킴에는 유한형 점 $t \in T_s$가 존재한다.
이제 $T_s \to T$는 $s \to S$의 기저변환이므로 유한형 사상이다
(보조정리 \ref{morphisms-lemma-base-change-finite-type}).
따라서 보조정리 \ref{morphisms-lemma-finite-type-points-morphism}에 의해
$t$의 $T$에서의 상은 유한형 점이고,
구성에 의해 그 상은 $s$로 사상된다.
\end{proof}

\begin{lemma}
\label{morphisms-lemma-enough-finite-type-points}
$S$를 스킴이라 하자.
임의의 국소 폐집합 $T \subset S$에 대하여
$$
T \not = \emptyset
\Rightarrow
T \cap S_{\text{ft-pts}} \not = \emptyset.
$$
특히 임의의 닫힌집합 $T \subset S$에 대하여
$T \cap S_{\text{ft-pts}}$가 $T$에서 조밀함을 알 수 있다.
\end{lemma}

\begin{proof}
$T$는 스킴 구조를 가진다는 것을 관찰하자(스킴, 보조정리
\ref{schemes-lemma-reduced-closed-subscheme} 참조).
그 구조에서 $T \to S$는 국소 폐쇄 몰입이다.
임의의 국소 폐쇄 몰입은 국소적으로 유한형이다.
이는 보조정리 \ref{morphisms-lemma-immersion-locally-finite-type}를 보라.
따라서 보조정리 \ref{morphisms-lemma-finite-type-points-morphism}에 의해
$T_{\text{ft-pts}} \subset S_{\text{ft-pts}}$임을 얻는다.
마지막으로 $T$의 공집합이 아닌 아핀 열린집합은 적어도 하나의
닫힌점을 가지며, 보조정리 \ref{morphisms-lemma-identify-finite-type-points}에
의해 그 점은 $T$의 유한형 점이다.
\end{proof}

\noindent
따라서 위상수학, 절 \ref{topology-section-space-jacobson}의
대부분의 내용은 닫힌점들의 집합을 유한형 점들의 집합으로
바꾸어도 성립한다.
실제로 $S$가 야콥슨이면 유한형 점들이 닫힌점으로 되돌아온다.

\begin{lemma}
\label{morphisms-lemma-jacobson-finite-type-points}
$S$를 스킴이라 하자. 다음 조건들은 서로 동치이다.
\begin{enumerate}
\item 스킴 $S$는 야콥슨이다.
\item $S_{\text{ft-pts}}$는 $S$의 닫힌점들의 집합이다.
\item 국소적으로 유한형인 모든 $T \to S$에 대하여
닫힌점은 닫힌점으로 사상된다.
\item 국소적으로 유한형인 모든 $T \to S$에 대하여
닫힌점 $t \in T$는 닫힌점 $s \in S$로 사상되고
$\kappa(s) \subset \kappa(t)$가 유한이다.
\end{enumerate}
\end{lemma}

\begin{proof}
(4) $\Rightarrow$ (3) $\Rightarrow$ (2)는 자명하게 성립한다.
보조정리 \ref{morphisms-lemma-enough-finite-type-points}에 의해
(2)는 (1)을 함의한다. 그러므로 (1)이 (4)를 함의함을
보이는 것으로 충분하다.
$T \to S$가 국소적으로 유한형이라고 하자.
$t \in T$를 닫힌점으로 택하고 그 상을 $s \in S$라 하자.
아핀 열린 근방 $\Spec(R) = U \subset S$를 $s$의 근방으로,
아핀 열린 근방 $\Spec(A) = V \subset T$를 $t$의 근방으로 택하되,
$V$가 $U$ 안으로 사상되도록 하자.
유도된 환 사상 $R \to A$는 유한형이다
(보조정리 \ref{morphisms-lemma-locally-finite-type-characterize}를 보라).
또한 성질, 보조정리
\ref{properties-lemma-locally-jacobson}에 의해 $R$은 야콥슨이다.
따라서 결과는 대수, 명제
\ref{algebra-proposition-Jacobson-permanence}로부터 따른다.
\end{proof}

\begin{lemma}
\label{morphisms-lemma-Jacobson-universally-Jacobson}
$S$를 야콥슨 스킴이라 하자.
$S$ 위에서 국소적으로 유한형인 임의의 스킴은 야콥슨이다.
\end{lemma}

\begin{proof}
이는
대수, 명제 \ref{algebra-proposition-Jacobson-permanence}
(및 성질, 보조정리 \ref{properties-lemma-locally-jacobson}와
보조정리 \ref{morphisms-lemma-locally-finite-type-characterize})로부터
자명하다.
\end{proof}

\begin{lemma}
\label{morphisms-lemma-ubiquity-Jacobson-schemes}
다음 유형의 스킴들은 야콥슨이다.
\begin{enumerate}
\item 체 위에서 국소적으로 유한형인 임의의 스킴.
\item $\mathbf{Z}$ 위에서 국소적으로 유한형인 임의의 스킴.
\item 무한히 많은 소수를 갖는 $1$차원 뇌터 환 위에서
국소적으로 유한형인 임의의 스킴.
\item $\Spec(R) \setminus \{\mathfrak m\}$ 꼴의 스킴이며,
여기서 $(R, \mathfrak m)$가 뇌터 국소환이다.
그 위에서 국소적으로 유한형인 임의의 스킴도 그렇다.
\end{enumerate}
\end{lemma}

\begin{proof}
언급하지 않고 보조정리
\ref{morphisms-lemma-Jacobson-universally-Jacobson}을 사용할 것이다.
체의 스펙트럼은 분명히 야콥슨이다.
$\mathbf{Z}$의 스펙트럼은 야콥슨이며, 대수, 보조정리
\ref{algebra-lemma-pid-jacobson}을 보라.
(3)에 대해서는 대수, 보조정리
\ref{algebra-lemma-noetherian-dim-1-Jacobson}을 보라.
(4)에 대해서는 성질, 보조정리
\ref{properties-lemma-complement-closed-point-Jacobson}을 보라.
\end{proof}
\section{보편적으로 카테나리인 스킴}
\label{morphisms-section-universally-catenary}

\noindent
위상공간 $X$를 모든 기약 닫힌 부분집합의 쌍 $T \subset T'$에 대하여
기약 닫힌 부분집합들의 최대 사슬
$$
T = T_0 \subset T_1 \subset \ldots \subset T_e = T'
$$
이 존재하고 이러한 사슬이 모두 같은 길이를 가질 때 {\it 카테나리}라고
부른다. 위상수학, 정의 \ref{topology-definition-catenary}를 보라.
스킴은 그 기저 위상공간이 카테나리일 때 카테나리라고 한다. 성질,
정의 \ref{properties-definition-catenary}를 보라.

\begin{definition}
\label{morphisms-definition-universally-catenary}
$S$를 국소적으로 뇌터인 스킴이라 하자.
$S$가 {\it 보편적으로 카테나리}라고 함은 국소적으로 유한형인
$S$를 보편적으로 카테나리라고 하며, 모든 사상에 대하여
$X \to S$인 스킴 $X$가 카테나리인 것을 말한다.
\end{definition}

\noindent
이는 카테나리보다 “더 나은” 개념이다. 카테나리이지만 보편적으로
카테나리가 아닌 뇌터 스킴들이 존재하기 때문이다. 예, 절
\ref{examples-section-non-catenary-Noetherian-local}을 보라.
많은 스킴이 보편적으로 카테나리이다. 아래 보조정리
\ref{morphisms-lemma-ubiquity-uc}를 보라.

\medskip\noindent
환 $A$를 임의의 소 아이디얼 쌍 $\mathfrak p \subset \mathfrak q$에 대하여
소 아이디얼들의 최대 사슬
$$
\mathfrak p =
\mathfrak p_0 \subset \ldots \subset \mathfrak p_e
= \mathfrak q
$$
이 존재하고 이 사슬들이 모두 같은 길이를 가질 때 {\it 카테나리}라고
부른다. 대수, 정의 \ref{algebra-definition-catenary}를 보라.
성질, 절 \ref{properties-section-catenary}에서 카테나리 스킴과
카테나리 환 사이의 관계를 살펴보았다.
환 $A$를 $A$가 뇌터이고 모든 유한형 환 사상 $A \to B$에 대하여
환 $B$가 카테나리일 때 {\it 보편적으로 카테나리}라고 부른다.
대수, 정의 \ref{algebra-definition-universally-catenary}를 보라.
대수기하에서 나타나는 흥미로운 환의 대부분은 이 성질을 만족한다.

\begin{lemma}
\label{morphisms-lemma-universally-catenary-local}
$S$를 국소적으로 뇌터인 스킴이라 하자. 다음 조건들은 서로 동치이다.
\begin{enumerate}
\item $S$는 보편적으로 카테나리이다.
\item $S$의 열린 덮개가 존재하여 그 모든 원소가 보편적으로
카테나리인 스킴이다.
\item 모든 아핀 열린집합 $\Spec(R) = U \subset S$에 대하여
환 $R$이 보편적으로 카테나리이고,
\item 아핀 열린 덮개 $S = \bigcup U_i$가 존재하여 각 $U_i$가
보편적으로 카테나리인 환의 스펙트럼이다.
\end{enumerate}
더 나아가 이 경우 $S$ 위에서 국소적으로 유한형인 임의의 스킴도
보편적으로 카테나리이다.
\end{lemma}

\begin{proof}
보조정리 \ref{morphisms-lemma-immersion-locally-finite-type}에 의해 열린 몰입은
국소적으로 유한형이다. 국소적으로 유한형인 사상들의 합성도 국소적으로
유한형이다(보조정리 \ref{morphisms-lemma-composition-finite-type}).
따라서 $S$가 보편적으로 카테나리이면 모든 열린집합과
$S$ 위에서 국소적으로 유한형인 모든 스킴도 보편적으로 카테나리라는
것은 자명하다. 이것이 보조정리의 마지막 명제와 (1)이 (2)를
함의한다는 것을 증명한다.

\medskip\noindent
$\Spec(R)$가 보편적으로 카테나리인 스킴이면,
모든 스킴 $\Spec(A)$가, $A$가 $R$ 위에서 유한형인 대수일 때,
카테나리이다.
따라서 이 모든 환 $A$는 대수, 보조정리
\ref{algebra-lemma-catenary}에 의해 카테나리이다.
그러므로 $R$은 보편적으로 카테나리이다. 위의 관찰과 결합하면
(1)이 (3)을, (2)가 (4)를 함의함을 얻는다. 물론 (3)은 자명하게
(4)를 함의한다.

\medskip\noindent
증명을 마치기 위해 (4)가 (1)을 함의함을 보인다.
(4)를 가정하고 $X \to S$를 국소적으로 유한형인 사상이라 하자.
아핀 열린 덮개 $X = \bigcup V_j$를 택할 수 있어 각 $V_j \to S$가
어떤 $U_i$ 안으로 사상된다. 보조정리
\ref{morphisms-lemma-locally-finite-type-characterize}에 의해 유도된 환 사상
$\mathcal{O}(U_i) \to \mathcal{O}(V_j)$는 유한형이다.
따라서 $\mathcal{O}(V_j)$는 카테나리이다. 그러므로 성질, 보조정리
\ref{properties-lemma-catenary-local}에 의해 $X$는 카테나리이다.
\end{proof}

\begin{lemma}
\label{morphisms-lemma-universally-catenary-local-rings-universally-catenary}
$S$를 국소적으로 뇌터인 스킴이라 하자.
다음 조건들은 서로 동치이다.
\begin{enumerate}
\item $S$는 보편적으로 카테나리이고,
\item 모든 국소환 $\mathcal{O}_{S, s}$가 $S$에서 보편적으로 카테나리이다.
\end{enumerate}
\end{lemma}

\begin{proof}
$S$의 모든 국소환이 보편적으로 카테나리라고 가정하자.
$f : X \to S$를 국소적으로 유한형인 사상이라 하자.
$X$가 카테나리인 것은 $\mathcal{O}_{X, x}$가 카테나리인 것이
모든 $x \in X$에 대하여 성립하는 것과 동치임을 안다.
$f(x) = s$이면 $\mathcal{O}_{X, x}$는
$\mathcal{O}_{S, s}$ 위에서 본질적으로 유한형이다.
따라서 $\mathcal{O}_{X, x}$는 카테나리이며,
$\mathcal{O}_{S, s}$가 보편적으로 카테나리라는 가정에 따른다.

\medskip\noindent
역으로 $S$가 보편적으로 카테나리라고 가정하자. $s \in S$라 하자.
보조정리 \ref{morphisms-lemma-universally-catenary-local}에 의해
$S$를 $s$의 아핀 열린 근방으로 바꾸어도 된다. 즉
$S = \Spec(R)$이고 $s$는 소 아이디얼 $\mathfrak p$에 대응한다고 하자.
유한형 $R_{\mathfrak p}$-대수 $A'$는 어떤 $A_{\mathfrak p}$ 꼴이며,
여기서 $R$ 위의 유한형 대수 $A$이다.
가정에 의해(원한다면 보조정리
\ref{morphisms-lemma-universally-catenary-local}도 사용하여) 환 $A$는 카테나리이고,
따라서 $A'$도 $A$의 국소화이므로 카테나리이다.
그러므로 $R_{\mathfrak p}$는 보편적으로 카테나리이다.
\end{proof}

\begin{lemma}
\label{morphisms-lemma-catenary-check-irreducible}
$S$를 국소적으로 뇌터인 스킴이라 하자. $S$가 보편적으로 카테나리인 것은
$S$의 기약 성분들이 보편적으로 카테나리인 것과 동치이다.
\end{lemma}

\begin{proof}
생략한다. 아핀인 경우에는 대수, 보조정리
\ref{algebra-lemma-catenary-check-irreducible}를 보라.
\end{proof}

\begin{lemma}
\label{morphisms-lemma-ubiquity-uc}
다음 유형의 스킴들은 보편적으로 카테나리이다.
\begin{enumerate}
\item 체 위에서 국소적으로 유한형인 임의의 스킴.
\item 코헨--매콜리 스킴 위에서 국소적으로 유한형인 임의의 스킴.
\item $\mathbf{Z}$ 위에서 국소적으로 유한형인 임의의 스킴.
\item $1$차원 뇌터 정역 위에서 국소적으로 유한형인 임의의 스킴.
\item 기타.
\end{enumerate}
\end{lemma}

\begin{proof}
이 모든 결과는 코헨--매콜리 환이 보편적으로 카테나리라는 사실에서
따른다. 대수, 보조정리 \ref{algebra-lemma-CM-ring-catenary}를 보라.
또한 보조정리 \ref{morphisms-lemma-universally-catenary-local}의 마지막
주장을 사용한다. 세부 사항은 생략한다.
\end{proof}
\section{나가타 스킴의 재고찰}
\label{morphisms-section-nagata}

\noindent
정의와 나가타 및 보편적으로 일본적인 스킴의 기본 성질은 성질,
절 \ref{properties-section-nagata}를 보라.

\begin{lemma}
\label{morphisms-lemma-finite-type-nagata}
$f : X \to S$를 사상이라 하자.
$S$가 나가타이고 $f$가 국소적으로 유한형이면 $X$는 나가타이다.
$S$가 보편적으로 일본적이고
$f$가 국소적으로 유한형이면 $X$는 보편적으로 일본적이다.
\end{lemma}

\begin{proof}
“보편적으로 일본적”인 경우에는 대수, 보조정리
\ref{algebra-lemma-universally-japanese}에서 따른다.
“나가타”인 경우에는 대수, 명제
\ref{algebra-proposition-nagata-universally-japanese}에서 따른다.
\end{proof}

\begin{lemma}
\label{morphisms-lemma-ubiquity-nagata}
다음 유형의 스킴들은 나가타이다.
\begin{enumerate}
\item 체 위에서 국소적으로 유한형인 임의의 스킴.
\item 뇌터 완비 국소환 위에서 국소적으로 유한형인 임의의 스킴.
\item $\mathbf{Z}$ 위에서 국소적으로 유한형인 임의의 스킴.
\item 표수가 0인 데데킨트 환 위에서 국소적으로 유한형인 임의의 스킴.
\item 기타.
\end{enumerate}
\end{lemma}

\begin{proof}
보조정리 \ref{morphisms-lemma-finite-type-nagata}에 의해 위에서 언급한
환들이 나가타 환임을 보이면 충분하다. 이는 대수, 명제
\ref{algebra-proposition-ubiquity-nagata}를 보라.
\end{proof}

\section{특이 궤적의 재고찰}
\label{morphisms-section-singular-locus}

\noindent
기저 위에서 국소적으로 유한형인 임의의 스킴에 대해 정칙 궤적이
열린집합임을 보장하는 기준을 찾는다. 정의는 다음과 같다.

\begin{definition}
\label{morphisms-definition-J}
$X$를 국소적으로 뇌터인 스킴이라 하자. $X$가 모든 국소적으로 유한형인
사상 $Y \to X$에 대하여 정칙 궤적 $\text{Reg}(Y)$가
$Y$에서 열린집합이면 {\it J-2}라고 한다.
\end{definition}

\noindent
이는 뇌터 환에 대한 대응 개념이다. 더 많은 대수, 정의
\ref{more-algebra-definition-J}를 보라.

\begin{lemma}
\label{morphisms-lemma-J}
$X$를 국소적으로 뇌터인 스킴이라 하자. 다음 조건들은 서로 동치이다.
\begin{enumerate}
\item $X$는 J-2이다.
\item $X$의 열린 덮개가 존재하여 그 모든 원소가 J-2 스킴이다.
\item 모든 아핀 열린집합 $\Spec(R) = U \subset X$에 대하여 환
$R$이 J-2이고,
\item 아핀 열린 덮개 $S = \bigcup U_i$가 존재하여 각
$\mathcal{O}(U_i)$가 모든 $i$에 대하여 J-2이다.
\end{enumerate}
더 나아가 이 경우 $X$ 위에서 국소적으로 유한형인 임의의 스킴도
J-2이다.
\end{lemma}

\begin{proof}
보조정리 \ref{morphisms-lemma-immersion-locally-finite-type}에 의해 열린 몰입은
국소적으로 유한형이다. 국소적으로 유한형인 사상들의 합성도
국소적으로 유한형이다
(보조정리 \ref{morphisms-lemma-composition-finite-type}).
따라서 $X$가 J-2이면 모든 열린집합과 $X$ 위에서 국소적으로
유한형인 임의의 스킴도 J-2라는 것은 자명하다. 이것이 보조정리의
마지막 명제를 증명한다.

\medskip\noindent
$\Spec(R)$가 J-2이면 모든 유한형 $R$-대수 $A$에 대하여
스킴 $\Spec(A)$의 정칙 궤적은 열린집합이다. 따라서 정의에 의해
$R$은 J-2이다(대수 추가, 정의
\ref{more-algebra-definition-J}를 보라).
위의 관찰과 결합하면 (1)이 (3)을, (2)가 (4)를 함의함을 얻는다.
물론 (1) $\Rightarrow$ (2)와 (3) $\Rightarrow$ (4)는 자명하다.

\medskip\noindent
증명을 마치기 위해 (4)가 (1)을 함의함을 보인다.
(4)를 가정하고 $Y \to X$를 국소적으로 유한형인 사상이라 하자.
아핀 열린 덮개 $Y = \bigcup V_j$를 택할 수 있어 각 $V_j \to X$가
어떤 $U_i$ 안으로 사상된다. 보조정리
\ref{morphisms-lemma-locally-finite-type-characterize}에 의해 유도된 환 사상
$\mathcal{O}(U_i) \to \mathcal{O}(V_j)$는 유한형이다.
따라서 $V_j = \Spec(\mathcal{O}(V_j))$의 정칙 궤적은 열린집합이다.
$\text{Reg}(Y) \cap V_j = \text{Reg}(V_j)$이므로
$\text{Reg}(Y)$가 원하는 대로 열린집합임을 얻는다.
\end{proof}

\begin{lemma}
\label{morphisms-lemma-ubiquity-J-2}
다음 유형의 스킴들은 J-2이다.
\begin{enumerate}
\item 체 위에서 국소적으로 유한형인 임의의 스킴.
\item 뇌터 완비 국소환 위에서 국소적으로 유한형인 임의의 스킴.
\item $\mathbf{Z}$ 위에서 국소적으로 유한형인 임의의 스킴.
\item 차원이 $1$인 뇌터 국소환 위에서 국소적으로 유한형인 임의의 스킴.
\item 차원이 $1$인 나가타 환 위에서 국소적으로 유한형인 임의의 스킴.
\item 표수가 0인 데데킨트 환 위에서 국소적으로 유한형인 임의의 스킴.
\item 기타.
\end{enumerate}
\end{lemma}

\begin{proof}
보조정리 \ref{morphisms-lemma-J}에 의해 위에서 언급한 환들이 J-2임을
보이면 충분하다. 이는 대수 추가, 명제
\ref{more-algebra-proposition-ubiquity-J-2}를 보라.
\end{proof}
\section{우수 스킴}
\label{morphisms-section-excellent}

\noindent
환이 G-환이고 J-2이면 준우수라고 한다.
환이 준우수이고 보편적으로 카테나리이면 우수라고 한다.

\begin{definition}
\label{morphisms-definition-excellent}
$X$를 스킴이라 하자.
\begin{enumerate}
\item $X$를 {\it 준우수}라고 함은 모든 $x \in X$에 대하여
$x \in U \subset X$인 아핀 열린 근방이 존재하고
그 환 $\mathcal{O}_X(U)$가 준우수인 경우이다
(대수 추가, 정의 \ref{more-algebra-definition-excellent} 참조).
\item $X$를 {\it 우수}라고 함은 모든 $x \in X$에 대하여
$x \in U \subset X$인 아핀 열린 근방이 존재하고
그 환 $\mathcal{O}_X(U)$가 우수인 경우이다
(대수 추가, 정의 \ref{more-algebra-definition-excellent} 참조).
\end{enumerate}
\end{definition}

\noindent
준우수 스킴은 국소적으로 뇌터이다(G-환은 뇌터이다).

\begin{lemma}
\label{morphisms-lemma-characterize-quasi-excellent}
$X$를 스킴이라 하자. 다음 조건들은 서로 동치이다.
\begin{enumerate}
\item $X$는 준우수이고,
\item $X$는 G-스킴이고 J-2이다.
\end{enumerate}
\end{lemma}

\begin{proof}
(1) $\Rightarrow$ (2)는 준우수 스킴의 정의
(정의 \ref{morphisms-definition-excellent}), 준우수 환의 정의
(대수 추가, 정의 \ref{more-algebra-definition-excellent}),
G-스킴의 정의(성질, 정의
\ref{properties-definition-G-scheme})와 보조정리 \ref{morphisms-lemma-J}를
결합하면 따른다.
역으로 $X$가 G-스킴이고 J-2이면 모든 $x \in X$에 대하여
$x \in U \subset X$인 아핀 열린 근방을 택할 수 있고
$\mathcal{O}_X(U)$는 G-환이다. 보조정리 \ref{morphisms-lemma-J}에 의해
환 $\mathcal{O}_X(U)$도 J-2이므로 준우수이다.
\end{proof}

\begin{lemma}
\label{morphisms-lemma-quasi-excellent-nagata}
준우수 스킴은 나가타이다.
\end{lemma}

\begin{proof}
대수 추가, 보조정리
\ref{more-algebra-lemma-quasi-excellent-nagata}를 보라.
\end{proof}

\begin{lemma}
\label{morphisms-lemma-characterize-excellent}
$X$를 스킴이라 하자. 다음 조건들은 서로 동치이다.
\begin{enumerate}
\item $X$는 우수이고,
\item $X$는 준우수이며 보편적으로 카테나리이다.
\end{enumerate}
\end{lemma}

\begin{proof}
(1)을 가정하자. 우수 환은 준우수이므로 $X$가 준우수인 것은 자명하다.
우수 환은 보편적으로 카테나리이므로 보조정리
\ref{morphisms-lemma-universally-catenary-local}에 의해 $X$도 보편적으로
카테나리이다.
(2)를 가정하자. $x \in X$를 택하자. $X$가 준우수이므로
$x \in U \subset X$인 아핀 열린 근방이 존재하여
$\mathcal{O}_X(U)$가 준우수이다.
$X$가 보편적으로 카테나리이므로 보조정리
\ref{morphisms-lemma-universally-catenary-local}에 의해
환 $\mathcal{O}_X(U)$도 보편적으로 카테나리이고, 따라서 우수이다.
\end{proof}

\begin{lemma}
\label{morphisms-lemma-locally-excellent}
$X$를 스킴이라 하자. 다음 조건들은 서로 동치이다.
\begin{enumerate}
\item 스킴 $X$는 (준)우수이다.
\item 모든 아핀 열린집합 $U \subset X$에 대하여 환
$\mathcal{O}_X(U)$가 (준)우수이다.
\item 아핀 열린 덮개 $X = \bigcup U_i$가 존재하여 각
$\mathcal{O}_X(U_i)$가 (준)우수이다.
\item 열린 덮개 $X = \bigcup X_j$가 존재하여 각 열린 부분스킴
$X_j$가 (준)우수이다.
\end{enumerate}
더 나아가 $X$가 (준)우수이면 모든 열린 부분스킴도 (준)우수이다.
\end{lemma}

\begin{proof}
보조정리 \ref{morphisms-lemma-characterize-quasi-excellent}에 의해 준우수인 것은
G-스킴이고 J-2인 것과 같다. 따라서 준우수의 경우에는 보조정리
\ref{morphisms-lemma-J}와 성질, 보조정리
\ref{properties-lemma-locally-G-scheme}에서 따른다.
우수의 경우에는 보조정리 \ref{morphisms-lemma-characterize-excellent},
준우수의 경우 및 보조정리
\ref{morphisms-lemma-universally-catenary-local}을 사용한다.
\end{proof}

\begin{lemma}
\label{morphisms-lemma-finite-type-excellent}
$f : X \to S$를 사상이라 하자.
$S$가 (준)우수이고 $f$가 국소적으로 유한형이면
$X$는 (준)우수이다.
\end{lemma}

\begin{proof}
대수 추가, 보조정리
\ref{more-algebra-lemma-finite-type-over-excellent}를 보라.
\end{proof}

\begin{lemma}
\label{morphisms-lemma-ubiquity-excellent}
다음 유형의 스킴들은 우수이다.
\begin{enumerate}
\item 체 위에서 국소적으로 유한형인 임의의 스킴.
\item 뇌터 완비 국소환 위에서 국소적으로 유한형인 임의의 스킴.
\item $\mathbf{Z}$ 위에서 국소적으로 유한형인 임의의 스킴.
\item 표수가 0인 데데킨트 환 위에서 국소적으로 유한형인 임의의 스킴.
\end{enumerate}
\end{lemma}

\begin{proof}
보조정리 \ref{morphisms-lemma-locally-excellent}와
\ref{morphisms-lemma-finite-type-excellent}에 의해 위에서 언급한 환들이
우수임을 보이면 충분하다.
이는 대수 추가, 명제
\ref{more-algebra-proposition-ubiquity-excellent}를 보라.
\end{proof}
\section{준유한 사상}
\label{morphisms-section-quasi-finite}

\noindent
준유한 사상을 철저히 다루는 것은 이후의 많은 전개를 위한 토대이다.
이는 자리스키의 주정리의 여러 형태, 올의 차원 거동, 에탈 사상의 하강
등으로 이어진다. 이 절을 읽기 전에 대수, 절
\ref{algebra-section-quasi-finite}의 대수 결과를 살펴보는 것이 좋다.

\medskip\noindent
유한형 환 사상 $R \to A$가 소 아이디얼 $\mathfrak q$에서 준유한하다는 것은
$\mathfrak q$가 그 올의 고립점을 정의한다는 뜻이다. 대수, 정의
\ref{algebra-definition-quasi-finite}를 보라.

\begin{definition}
\label{morphisms-definition-quasi-finite}
\begin{reference}
\cite[II Definition 6.2.3]{EGA}
\end{reference}
$f : X \to S$를 스킴의 사상이라 하자.
\begin{enumerate}
\item 사상 $f$가 점 $x \in X$에서 {\it 준유한}이라고 함은
아핀 근방 $\Spec(A) = U \subset X$가 $x$를 포함하고
아핀 열린집합 $\Spec(R) = V \subset S$가 존재하여
$f(U) \subset V$이고 환 사상 $R \to A$가 유한형이며,
$R \to A$가 $A$의 소 아이디얼 가운데 $x$에 대응하는 것에서 준유한인 경우이다
(위의 정의를 보라).
\item $f$가 {\it 국소적으로 준유한}이라고 함은 $f$가
각 점 $x$에서 준유한이고 이 점들이 $X$의 모든 점을 달리는 경우이다.
\item $f$가 {\it 준유한}이라고 함은 $f$가 유한형이고
모든 점 $x$가 그 올의 고립점인 경우이다.
\end{enumerate}
\end{definition}

\noindent
자명하게 국소적으로 준유한인 사상은 국소적으로 유한형이다.
아래에서 국소적으로 유한형인 사상 $f$가 점 $x$에서 준유한인 것은
$x$가 그 올에서 고립점인 것과 동치임을 보일 것이다.
더 나아가 사상이 준유한인 점들의 집합은 열린집합이다.
이는 자리스키의 주정리를 다루는 절 \ref{morphisms-section-Zariski}에서 보인다.

\begin{lemma}
\label{morphisms-lemma-algebraic-residue-field-extension-closed-point-fibre}
$f : X \to S$를 스킴의 사상이라 하자.
점 $x \in X$를 택하고 $s = f(x)$로 두자.
$\kappa(x)/\kappa(s)$가 대수적 체 확장이면
\begin{enumerate}
\item $x$는 그 올의 닫힌점이다.
\item 추가로 $s$가 $S$의 닫힌점이면
$x$는 $X$의 닫힌점이다.
\end{enumerate}
\end{lemma}

\begin{proof}
두 번째 명제는 초등 위상수학에 의해 첫 번째 명제에서 따른다.
첫 번째 명제를 증명하기 위해 스킴, 보조정리
\ref{schemes-lemma-fibre-topological}에 따라
$X$를 $X_s$로, $S$를 $\Spec(\kappa(s))$로 바꾸어도 된다.
따라서 $S = \Spec(k)$가 체의 스펙트럼이라고 가정할 수 있다.
이 경우 아핀 열린집합 $\Spec(A) = U \subset X$ 중 $x$를 포함하는 것을
하나 택하자. 점 $x$는
$\mathfrak q \subset A$인 소 아이디얼에 대응하고
$\kappa(\mathfrak q)/k$는 대수적 체 확장이다. 대수, 보조정리
\ref{algebra-lemma-finite-residue-extension-closed}에 의해
$\mathfrak q$는 극대 아이디얼이다. 즉 $x \in U$는 닫힌점이다.
아핀 열린집합들이 $X$의 위상의 기저를 이루므로
$\{x\}$가 닫힌집합임을 결론내린다.
\end{proof}

\noindent
다음 보조정리는 힐베르트 영점정리의 한 형태이다.

\begin{lemma}
\label{morphisms-lemma-closed-point-fibre-locally-finite-type}
$f : X \to S$를 스킴의 사상이라 하자.
점 $x \in X$를 택하고 $s = f(x)$로 두자.
$f$가 국소적으로 유한형이라고 가정하자.
그러면 $x$가 그 올의 닫힌점인 것은
$\kappa(x)/\kappa(s)$가 유한 체 확장인 것과 동치이다.
\end{lemma}

\begin{proof}
확장이 유한이면 위의 보조정리
\ref{morphisms-lemma-algebraic-residue-field-extension-closed-point-fibre}에 의해
$x$는 올의 닫힌점이다. 역으로 $x$가 그 올의 닫힌점이라고 가정하자.
아핀 열린집합 $\Spec(A) = U \subset X$와 $\Spec(R) = V \subset S$를
$f(U) \subset V$가 되도록 택하자.
보조정리 \ref{morphisms-lemma-locally-finite-type-characterize}에 의해
환 사상 $R \to A$는 유한형이다. 소 아이디얼 $\mathfrak q \subset A$와
$\mathfrak p \subset R$를 각각 $x$와 $s$에 대응하도록 택하자.
올 환을
$\overline{A} = A \otimes_R \kappa(\mathfrak p)$로 두자.
$\overline{\mathfrak q}$를 $\overline{A}$의 소 아이디얼 중
$\mathfrak q$에 대응하는 것으로 두자.
$x$가 그 올의 닫힌점이라는 가정은
$\overline{\mathfrak q}$가 $\overline{A}$의 극대 아이디얼임을 뜻한다.
$\overline{A}$는 체 $\kappa(\mathfrak p)$ 위의 유한형 대수이므로
힐베르트 영점정리에 의해, 대수, 정리
\ref{algebra-theorem-nullstellensatz}를 보라.
$\kappa(\overline{\mathfrak q})$는 $\kappa(\mathfrak p)$의
유한 확장이다.
$\kappa(s) = \kappa(\mathfrak p)$이고
$\kappa(x) = \kappa(\mathfrak q) = \kappa(\overline{\mathfrak q})$이므로
결과를 얻는다.
\end{proof}

\begin{lemma}
\label{morphisms-lemma-base-change-closed-point-fibre-locally-finite-type}
$f : X \to S$를 국소적으로 유한형인 스킴의 사상이라 하자.
$g : S' \to S$를 임의의 사상이라 하자. 기저변환을
$f' : X' \to S'$로 표시하자.
$x' \in X'$가 점 $x \in X$로 사상되고 $X_{f(x)}$에서 닫힌점이면
$x'$는 $X'_{f'(x')}$에서 닫힌점이다.
\end{lemma}

\begin{proof}
잉여체 $\kappa(x')$는
$\kappa(f'(x')) \otimes_{\kappa(f(x))} \kappa(x)$의 몫이다.
스킴, 보조정리 \ref{schemes-lemma-points-fibre-product}를 보라.
따라서 $\kappa(f'(x'))$의 유한 확장이다. 실제로
$\kappa(x)$는 $\kappa(f(x))$의 유한 확장이다. 이는 보조정리
\ref{morphisms-lemma-closed-point-fibre-locally-finite-type}에 따른다.
그러므로 그 보조정리를 한 번 더 적용하면
$x'$가 그 올에서 닫힌점임을 얻는다.
\end{proof}

\begin{lemma}
\label{morphisms-lemma-residue-field-quasi-finite}
$f : X \to S$를 스킴의 사상이라 하자.
점 $x \in X$를 택하고 $s = f(x)$로 두자.
$f$가 $x$에서 준유한이면 잉여체 확장
$\kappa(x)/\kappa(s)$는 유한이다.
\end{lemma}

\begin{proof}
이는 대수, 정의 \ref{algebra-definition-quasi-finite}에서 자명하다.
\end{proof}

\begin{lemma}
\label{morphisms-lemma-quasi-finite-at-point-characterize}
$f : X \to S$를 스킴의 사상이라 하자.
점 $x \in X$를 택하고 $s = f(x)$로 두자.
$X_s$를 $f$의 $s$에서의 올이라 하자.
$f$가 국소적으로 유한형이라고 가정하자.
다음 조건들은 서로 동치이다.
\begin{enumerate}
\item 사상 $f$는 $x$에서 준유한이다.
\item 점 $x$는 $X_s$의 고립점이다.
\item 점 $x$는 $X_s$에서 닫혀 있고,
$x' \in X_s$, $x' \not = x$이면서 $x$로 특수화되는 점은 없다.
\item $\Spec(A) = U \subset X$, $\Spec(R) = V \subset S$인 임의의
두 아핀 열린집합이 $f(U) \subset V$를 만족하고, $x \in U$가
$\mathfrak q \subset A$에 대응하면 환 사상 $R \to A$는
$\mathfrak q$에서 준유한이다.
\end{enumerate}
\end{lemma}

\begin{proof}
$f$가 $x$에서 준유한이라고 가정하자. 가정에 따라
$U \subset X$, $V \subset S$인 열린집합들이 존재하여
$f(U) \subset V$, $x \in U$이고 $x$는 $U_s$의 고립점이다.
따라서 $\{x\} \subset U_s$는 열린 부분집합이다.
$U_s = U \cap X_s \subset X_s$ 역시 열린집합이므로
$\{x\} \subset X_s$도 열린 부분집합이다. 따라서
$x$는 $X_s$의 고립점이다.

\medskip\noindent
$X_s$는 보조정리 \ref{morphisms-lemma-ubiquity-Jacobson-schemes}와
보조정리 \ref{morphisms-lemma-base-change-finite-type}에 의해 야콥슨 스킴임에 유의하자.
$x$가 $X_s$에서 고립되어 있다면, 즉 $\{x\} \subset X_s$가 열려 있다면,
$\{x\}$는 야콥슨 성질에 의해 닫힌점을 포함한다. 따라서
$x$는 $X_s$에서 닫혀 있다. 또한 $x' \in X_s$가 $x$와 다르면서
$x$로 특수화되는 일은 없음이 자명하다.

\medskip\noindent
$x$가 $X_s$에서 닫혀 있고 $x' \in X_s$가 $x$와 다르면서
$x$로 특수화되는 일은 없다고 가정하자. 아핀 열린집합
$\Spec(A) = U \subset X$, $\Spec(R) = V \subset S$가
$f(U) \subset V$와 $x \in U$를 만족한다고 하자.
$\mathfrak q \subset A$는 $x$에, $\mathfrak p \subset R$는 $s$에 대응한다고 하자.
보조정리 \ref{morphisms-lemma-locally-finite-type-characterize}에 의해 환 사상
$R \to A$는 유한형이다. 올 환
$\overline{A} = A \otimes_R \kappa(\mathfrak p)$를 생각하자.
$\overline{\mathfrak q}$를 $\overline{A}$의 소 아이디얼 중
$\mathfrak q$에 대응하는 것이라 하자. $\Spec(\overline{A})$는
올 $X_s$의 열린 부분스킴이므로 $\overline{q}$는
$\overline{A}$의 극대 아이디얼이고, $\Spec(\overline{A})$에는
$\overline{\mathfrak q}$로 특수화되는 점이 없다. 이는
$\dim(\overline{A}_{\overline{q}}) = 0$임을 뜻한다.
따라서 대수, 정의 \ref{algebra-definition-quasi-finite}에 의해
$R \to A$는 $\mathfrak q$에서 준유한이다. 즉 정의에 따라
$X \to S$는 $x$에서 준유한이다.

\medskip\noindent
이제 조건 (1)--(3)이 모두 동치임을 보였다.
(4)가 (1)을 함의함은 자명하다. 또한 (2)가 (4)를 함의함도 자명하다.
실제로 $x$가 $X_s$의 고립점이면 그 점은 $U_s$에서도 고립점이다.
여기서 $U$는 그 점을 포함하는 임의의 열린집합이다.
\end{proof}

\begin{lemma}
\label{morphisms-lemma-finite-fibre}
$f : X \to S$를 스킴의 사상이라 하자.
$s \in S$라 하자. 다음을 가정하자.
\begin{enumerate}
\item $f$는 국소적으로 유한형이다.
\item $f^{-1}(\{s\})$는 유한집합이다.
\end{enumerate}
그러면 $X_s$는 유한 이산 위상공간이고,
$f$는 $X$에서 $s$ 위에 놓이는 각 점에서 준유한이다.
\end{lemma}

\begin{proof}
$T$가 (a) 체 $k$ 위에서 국소적으로 유한형이고 (b) 점이 유한개인
스킴이라고 하자. 그러면 보조정리
\ref{morphisms-lemma-ubiquity-Jacobson-schemes}에 의해 $T$는 야콥슨 스킴이다.
유한 야콥슨 공간은 이산공간이다. 위상수학, 보조정리
\ref{topology-lemma-finite-jacobson}를 보라.
이 관찰을 올 $X_s$에 적용하자. 이 올은 $\Spec(\kappa(s))$ 위에서
국소적으로 유한형이므로 첫 번째 명제를 얻는다. 마지막으로 보조정리
\ref{morphisms-lemma-quasi-finite-at-point-characterize}를 적용하면 두 번째 명제를 얻는다.
\end{proof}

\begin{lemma}
\label{morphisms-lemma-locally-quasi-finite-fibres}
\begin{slogan}
이산 올을 갖는 유한형 사상은 준유한이다.
\end{slogan}
$f : X \to S$를 스킴의 사상이라 하자.
$f$가 국소적으로 유한형이라고 가정하자.
다음 조건들은 서로 동치이다.
\begin{enumerate}
\item $f$는 국소적으로 준유한이다.
\item 모든 $s \in S$에 대해 올 $X_s$는 이산 위상공간이다.
\item 모든 사상 $\Spec(k) \to S$에 대해, 여기서 $k$는 체이며,
기저변환 $X_k$의 바탕 위상공간은 이산공간이다.
\end{enumerate}
\end{lemma}

\begin{proof}
(3)이 (2)를 함의함은 즉시 알 수 있다.
보조정리 \ref{morphisms-lemma-quasi-finite-at-point-characterize}에 의해
(2)와 (1)은 동치이다. (2)를 가정하고 (3)에서와 같이
$\Spec(k) \to S$가 주어졌다고 하자.
그 상을 $s \in S$라 하자. 이 점은 $\Spec(k) \to S$의 상이다.
그러면 $X_k$는 $X_s$를 $\Spec(k) \to \Spec(\kappa(s))$로
기저변환하여 얻어진다. 따라서 보조정리
\ref{morphisms-lemma-base-change-closed-point-fibre-locally-finite-type}에 의해
$X_k$의 모든 점은 닫혀 있다. 또한 $X_k \to \Spec(k)$는
보조정리 \ref{morphisms-lemma-base-change-finite-type}에 의해 국소적으로 유한형이므로,
보조정리 \ref{morphisms-lemma-quasi-finite-at-point-characterize}를 적용하면
$X_k$의 모든 점이 고립되어 있음을, 즉 $X_k$의 바탕 위상공간이
이산공간임을 알 수 있다.
\end{proof}

\begin{lemma}
\label{morphisms-lemma-quasi-finite-locally-quasi-compact}
$f : X \to S$를 스킴의 사상이라 하자.
$f$가 준유한인 것은 $f$가 국소적으로 준유한이고 준콤팩트인 것과 동치이다.
\end{lemma}

\begin{proof}
$f$가 준유한이라고 가정하자. 정의
\ref{morphisms-definition-finite-type}에 의해 이는 준콤팩트이다. $x \in X$라 하자.
보조정리 \ref{morphisms-lemma-quasi-finite-at-point-characterize}에 의해
$f$는 $x$에서 준유한이다. 따라서 $f$는 준콤팩트이고 국소적으로 준유한이다.

\medskip\noindent
$f$가 준콤팩트이고 국소적으로 준유한이라고 가정하자.
그러면 $f$는 유한형이다. 점 $x \in X$를 택하자.
보조정리 \ref{morphisms-lemma-quasi-finite-at-point-characterize}에 의해
$x$는 그 올의 고립점이다. 이로써 보조정리가 증명되었다.
\end{proof}

\begin{lemma}
\label{morphisms-lemma-quasi-finite}
$f : X \to S$를 스킴의 사상이라 하자.
다음 조건들은 서로 동치이다.
\begin{enumerate}
\item $f$는 준유한이다.
\item $f$는 국소적으로 유한형이고 준콤팩트이며 올이 유한하다.
\end{enumerate}
\end{lemma}

\begin{proof}
$f$가 준유한이라고 가정하자. 특히 $f$는 국소적으로 유한형이고
준콤팩트이다(유한형이기 때문이다). $s \in S$라 하자.
모든 $x \in X_s$가 $X_s$의 고립점이므로
$X_s = \bigcup_{x \in X_s} \{x\}$는 열린 덮개이다. $f$가
준콤팩트이므로 올 $X_s$는 준콤팩트이다. 따라서
$X_s$는 유한하다.

\medskip\noindent
역으로 $f$가 국소적으로 유한형이고 준콤팩트이며 유한한 올을 갖는다고 하자.
그러면 보조정리 \ref{morphisms-lemma-finite-fibre}에 의해 국소적으로 준유한하다.
따라서 보조정리 \ref{morphisms-lemma-quasi-finite-locally-quasi-compact}에 의해
준유한하다.
\end{proof}

\noindent
환 사상 $R \to A$가 유한형이고 $A$의 {\it 모든} 소 아이디얼에서
준유한일 때 이를 준유한이라고 한다는 것을 상기하자.
대수, 정의 \ref{algebra-definition-quasi-finite}를 보라.

\begin{lemma}
\label{morphisms-lemma-locally-quasi-finite-characterize}
$f : X \to S$를 스킴의 사상이라 하자.
다음 조건들은 서로 동치이다.
\begin{enumerate}
\item 사상 $f$는 국소적으로 준유한이다.
\item $U \subset X$, $V \subset S$인 임의의 두 아핀 열린집합이
$f(U) \subset V$를 만족하면 환 사상
$\mathcal{O}_S(V) \to \mathcal{O}_X(U)$는 준유한이다.
\item 열린 덮개 $S = \bigcup_{j \in J} V_j$와 열린 덮개
$f^{-1}(V_j) = \bigcup_{i \in I_j} U_i$가 존재하여 각 사상
$U_i \to V_j$, $j\in J, i\in I_j$는 국소적으로 준유한이다.
\item 아핀 열린 덮개 $S = \bigcup_{j \in J} V_j$와 아핀 열린 덮개
$f^{-1}(V_j) = \bigcup_{i \in I_j} U_i$가 존재하여 모든
환 사상 $\mathcal{O}_S(V_j) \to \mathcal{O}_X(U_i)$가
$j\in J, i\in I_j$에 대해 준유한이다.
\end{enumerate}
더 나아가 $f$가 국소적으로 준유한이면 $U \subset X$, $V \subset S$인
임의의 열린 부분스킴이 $f(U) \subset V$를 만족할 때 제한
$f|_U : U \to V$도 국소적으로 준유한이다.
\end{lemma}

\begin{proof}
환 사상 $R \to A$에 대해 $P(R \to A)$를
``$R \to A$는 준유한이다''라는 뜻으로 정의하자
(보조정리 앞의 주의를 보라).
$P$가 환 사상의 국소적 성질임을 보이겠다.
정의 \ref{morphisms-definition-property-local}의 조건 (a), (b), (c)를 확인한다.
보조정리 \ref{morphisms-lemma-locally-finite-type-characterize}의 증명에서
``유한형''이라는 성질에 대해 (a), (b), (c)가 성립함을 보였다.
유한형 환 사상 $R \to A$에 대해 $R \to A$가
$\mathfrak q$에서 준유한이라는 성질은 국소환 $A_{\mathfrak q}$의
$R_{\mathfrak p}$-대수 구조에만 의존한다. 여기서
$\mathfrak p = R \cap \mathfrak q$이다(통상적인 기호의 남용이다).
이 관찰들로부터 정의 \ref{morphisms-definition-property-local}의
(a), (b), (c)가 곧바로 따른다. 예를 들어 환 사상 $R \to A$가 주어지고,
모든 환 사상 $R \to A_{a_i}$가 준유한이라고 하자. 여기서
$a_1, \ldots, a_n \in A$는 단위원 아이디얼을 생성한다.
그러면 $R \to A$는 유한형이다.
또한 임의의 소 아이디얼 $\mathfrak q \subset A$에 대해 국소환
$A_{\mathfrak q}$는 $R$-대수로서 국소환
$(A_{a_i})_{\mathfrak q_i}$와 동형이다. 여기서 적당한 $i$와
$\mathfrak q_i \subset A_{a_i}$를 택한다.
따라서 $R \to A$는 $\mathfrak q$에서 준유한이다.

\medskip\noindent
따라서 앞 단락의 $P$에 보조정리 \ref{morphisms-lemma-locally-P}를 적용할 수 있다.
그러므로 $f$가 국소적으로 준유한인 것과 $f$가 국소적으로 $P$형인 것이
동치임을 보이면 충분하다. $P(R \to A)$는
``$R \to A$는 준유한이다''라는 뜻이며, 이는 $R \to A$가
$A$의 모든 소 아이디얼에서 준유한이라는 뜻이다. 따라서 이 동치는
보조정리 \ref{morphisms-lemma-quasi-finite-at-point-characterize}에서 따른다.
\end{proof}

\begin{lemma}
\label{morphisms-lemma-composition-quasi-finite}
국소적으로 준유한인 두 사상의 합성은 국소적으로 준유한이다.
준유한 사상들에 대해서도 같은 명제가 성립한다.
\end{lemma}

\begin{proof}
보조정리 \ref{morphisms-lemma-locally-quasi-finite-characterize}의 증명에서
$P = $``준유한''이 환 사상의 국소적 성질이고, 스킴의 사상이
국소적으로 준유한인 것과 정의 \ref{morphisms-definition-locally-P}의 의미에서
국소적으로 $P$형인 것이 동치임을 보였다.
따라서 보조정리의 첫 번째 명제는 보조정리
\ref{morphisms-lemma-composition-type-P}와 준유한성이 합성에 대해 안정적인
환 사상의 성질이라는 사실을 함께 적용하면 따른다. 대수, 보조정리
\ref{algebra-lemma-quasi-finite-composition}를 보라.
위의 사실, 보조정리 \ref{morphisms-lemma-quasi-finite-locally-quasi-compact},
그리고 준콤팩트 사상의 합성이 준콤팩트라는 사실을 함께 적용하자.
스킴, 보조정리 \ref{schemes-lemma-composition-quasi-compact}를 보라.
그러면 준유한 사상들의 합성이 준유한임을 알 수 있다.
\end{proof}

\noindent
뒤에서(보조정리 \ref{morphisms-lemma-quasi-finite-points-open}) 다음 보조정리의
집합 $U$가 열린집합임을 보일 것이다.

\begin{lemma}
\label{morphisms-lemma-base-change-quasi-finite}
\begin{slogan}
(국소적으로) 준유한인 사상은 기저변환에 대해 안정적이다.
\end{slogan}
$f : X \to S$를 스킴의 사상이라 하자.
$g : S' \to S$를 스킴의 사상이라 하자.
$f' : X' \to S'$를 $f$의 $g$에 의한 기저변환이라 하고,
$g' : X' \to X$를 사영이라 하자.
$X$가 $S$ 위에서 국소적으로 유한형이라고 가정하자.
\begin{enumerate}
\item $U \subset X$(각각 $U' \subset X'$)를 $f$(각각 $f'$)가
준유한인 점들의 집합이라 하자. 그러면
$U' = U \times_S S' = (g')^{-1}(U)$이다.
\item 국소적으로 준유한인 사상의 기저변환은 국소적으로 준유한이다.
\item 준유한 사상의 기저변환은 준유한이다.
\end{enumerate}
\end{lemma}

\begin{proof}
첫 번째와 두 번째 명제는 대응하는 대수 결과에서 따른다. 대수, 보조정리
\ref{algebra-lemma-quasi-finite-base-change}를 보라. 여기에는 보조정리
\ref{morphisms-lemma-base-change-finite-type}에 의해 $f'$도 국소적으로 유한형이라는
사실을 함께 사용한다. 위의 사실과 보조정리
\ref{morphisms-lemma-quasi-finite-locally-quasi-compact}, 그리고 준콤팩트 사상의
기저변환이 준콤팩트라는 사실을 함께 적용하자. 스킴, 보조정리
\ref{schemes-lemma-quasi-compact-preserved-base-change}를 보라.
그러면 준유한 사상의 기저변환이 준유한임을 알 수 있다.
\end{proof}
\begin{lemma}
\label{morphisms-lemma-quasi-finite-at-a-finite-number-of-points}
$f : X \to S$를 유한형인 스킴의 사상이라 하자.
$s \in S$라 하자. $X$에서 $s$ 위에 놓이고 $f$가 준유한인 점은
유한개 이하이다.
\end{lemma}

\begin{proof}
올 $X_s$는 체 위의 유한형 스킴이므로 뇌터 스킴이다
(보조정리 \ref{morphisms-lemma-finite-type-noetherian}). 따라서 $X_s$의 위상은
뇌터 위상이고(성질, 보조정리 \ref{properties-lemma-Noetherian-topology}),
고립점은 유한개 이하일 수밖에 없다(초등 위상수학에 따른다).
그러므로 보조정리 \ref{morphisms-lemma-quasi-finite-at-point-characterize}에서
이 보조정리가 따른다.
\end{proof}

\begin{lemma}
\label{morphisms-lemma-monomorphism-loc-finite-type-loc-quasi-finite}
$f : X \to Y$를 스킴의 사상이라 하자.
$f$가 국소적으로 유한형인 모노사상이면 $f$는 분리이고
국소적으로 준유한이다.
\end{lemma}

\begin{proof}
모노사상은 분리이다. 스킴, 보조정리
\ref{schemes-lemma-monomorphism-separated}를 보라.
모노사상은 단사이므로, 예를 들어 보조정리
\ref{morphisms-lemma-quasi-finite-at-point-characterize}에 의해 $f$는
모든 $x \in X$에서 준유한이다.
\end{proof}

\begin{lemma}
\label{morphisms-lemma-immersion-locally-quasi-finite}
모든 몰입은 국소적으로 준유한이다.
\end{lemma}

\begin{proof}
열린몰입은 국소 동형이고 닫힌몰입은 자명하게 준유한이므로 성립한다.
\end{proof}

\begin{lemma}
\label{morphisms-lemma-permanence-quasi-finite}
$X \to Y$를 기저 스킴 $S$ 위의 스킴 사상이라 하자.
$x \in X$라 하자. $X \to S$가 $x$에서 준유한이면
$X \to Y$도 $x$에서 준유한이다.
$X$가 $S$ 위에서 국소적으로 준유한이면 $X \to Y$도
국소적으로 준유한이다.
\end{lemma}

\begin{proof}
보조정리 \ref{morphisms-lemma-locally-quasi-finite-characterize}를 통하면 이는 다음
대수적 사실로 옮겨진다. 환 사상 $A \to B \to C$가 주어지고
$A \to C$가 준유한이면 $B \to C$도 준유한이다.
이는 대수, 보조정리 \ref{algebra-lemma-four-rings}에서
$R = A$, $S = S' = C$, $R' = B$로 두면 따른다.
\end{proof}

\begin{lemma}
\label{morphisms-lemma-quasi-finite-local-source}
$f : X \to Y$와 $g : Y \to S$를 스킴의 사상이라 하자.
$f$가 전사이고 $g \circ f$가 국소적으로 준유한이며
$g$가 국소적으로 유한형이면 $g : Y \to S$는 국소적으로 준유한이다.
\end{lemma}

\begin{proof}
$x \in X$를 택하고 그 상들을 $y \in Y$, $s \in S$라 하자.
$g \circ f$가 국소적으로 준유한이므로 보조정리
\ref{morphisms-lemma-residue-field-quasi-finite}에 의해
$\kappa(x)/\kappa(s)$는 유한 확장이다.
따라서 $\kappa(y)/\kappa(s)$도 유한 확장이다. 그러므로 보조정리
\ref{morphisms-lemma-algebraic-residue-field-extension-closed-point-fibre}에 의해
$y$는 $Y_s$의 닫힌점이다. $f$가 전사이므로 $Y$의 모든 점은
$S$ 위의 자기 올에서 닫혀 있다. 따라서 보조정리
\ref{morphisms-lemma-quasi-finite-at-point-characterize}에 의해
$g$는 모든 점에서 준유한이다.
\end{proof}














\section{유한 표시의 사상}
\label{morphisms-section-finite-presentation}

\noindent
환 사상 $R \to A$에 대해 $A$가
$R[x_1, \ldots, x_n]/(f_1, \ldots, f_m)$과 $R$-대수로서
동형이고, 여기서 어떤 $n, m$과 어떤 다항식 $f_j$를 택할 수 있을 때
이를 유한 표시라고 한다는 것을 상기하자.
대수, 정의 \ref{algebra-definition-finite-type}를 보라.

\begin{definition}
\label{morphisms-definition-finite-presentation}
$f : X \to S$를 스킴의 사상이라 하자.
\begin{enumerate}
\item $f$가 $x \in X$에서 {\it 유한 표시}라고 함은,
아핀 열린 근방 $\Spec(A) = U \subset X$가 $x$를 포함하고,
아핀 열린집합 $\Spec(R) = V \subset S$가 존재하여
$f(U) \subset V$이고 유도된 환 사상 $R \to A$가 유한 표시인 경우이다.
\item $f$가 $X$의 모든 점에서 유한 표시이면 이를
{\it 국소적으로 유한 표시}라고 한다.
\item $f$가 국소적으로 유한 표시이고 준콤팩트이며 준분리이면
이를 {\it 유한 표시}라고 한다.
\end{enumerate}
\end{definition}

\noindent
유한 표시인 사상은 단지 국소적으로 유한 표시인 준콤팩트 사상이
{\bf 아님}에 유의하자. 뒤에서 국소적으로 유한 표시인 사상을,
아핀 스킴들의 임의의 유향계 $T_i$가 $S$ 위에 있을 때
$$
\colim \Mor_S(T_i, X) = \Mor_S(\lim T_i, X)
$$
를 만족하는 사상으로 특징지을 것이다. 극한,
명제 \ref{limits-proposition-characterize-locally-finite-presentation}를 보라.
극한, 절 \ref{limits-section-descending-relative}에서는
$S = \lim_i S_i$가 아핀 스킴들의 극한이고 $X$가 $S$ 위에서
유한 표시이면 $X_i$를 $S_i$ 위의 스킴이라 할 때 적당한 $i$에 대해
그 스킴으로 하강함을 보인다.

\begin{lemma}
\label{morphisms-lemma-locally-finite-presentation-characterize}
$f : X \to S$를 스킴의 사상이라 하자.
다음 조건들은 서로 동치이다.
\begin{enumerate}
\item 사상 $f$는 국소적으로 유한 표시이다.
\item 아핀 열린집합 $U \subset X$, $V \subset S$가
$f(U) \subset V$를 만족하면 환 사상
$\mathcal{O}_S(V) \to \mathcal{O}_X(U)$는 유한 표시이다.
\item 열린 덮개 $S = \bigcup_{j \in J} V_j$와 열린 덮개
$f^{-1}(V_j) = \bigcup_{i \in I_j} U_i$가 존재하여 각 사상
$U_i \to V_j$, $j\in J, i\in I_j$가 국소적으로 유한 표시이다.
\item 아핀 열린 덮개 $S = \bigcup_{j \in J} V_j$와 아핀 열린 덮개
$f^{-1}(V_j) = \bigcup_{i \in I_j} U_i$가 존재하여 환 사상
$\mathcal{O}_S(V_j) \to \mathcal{O}_X(U_i)$가 모든
$j\in J, i\in I_j$에 대해 유한 표시이다.
\end{enumerate}
더 나아가 $f$가 국소적으로 유한 표시이면 $U \subset X$, $V \subset S$인
임의의 열린 부분스킴이 $f(U) \subset V$를 만족할 때 제한
$f|_U : U \to V$도 국소적으로 유한 표시이다.
\end{lemma}

\begin{proof}
``$R \to A$는 유한 표시이다''라는 성질이 국소적임을 보이면
보조정리 \ref{morphisms-lemma-locally-P-characterize}에서 따른다.
정의 \ref{morphisms-definition-property-local}의 조건 (a), (b), (c)를 확인하자.
유한 표시라는 성질은 기저변환에 대해 안정적이므로 대수, 보조정리
\ref{algebra-lemma-base-change-finiteness}에 의해 (a)가 성립한다.
대수, 보조정리 \ref{algebra-lemma-compose-finite-type}에 의해
유한 표시는 합성에 대해 안정적이고, 임의의 환 $R$에 대해 환 사상
$R \to R_f$가 유한 표시임은 자명하다. 따라서 (b)가 성립한다.
마지막으로 대수, 보조정리 \ref{algebra-lemma-cover-upstairs}에 의해
성질 (c)가 성립한다.
\end{proof}

\begin{lemma}
\label{morphisms-lemma-composition-finite-presentation}
국소적으로 유한 표시인 두 사상의 합성은 국소적으로 유한 표시이다.
유한 표시인 사상들에 대해서도 같은 명제가 성립한다.
\end{lemma}

\begin{proof}
보조정리 \ref{morphisms-lemma-locally-finite-presentation-characterize}의 증명에서
유한 표시라는 성질이 환 사상의 국소적 성질임을 보였다.
따라서 보조정리의 첫 번째 명제는 보조정리
\ref{morphisms-lemma-composition-type-P}와 유한 표시라는 성질이 합성에 대해
안정적인 환 사상의 성질이라는 사실을 함께 적용하면 따른다.
대수, 보조정리 \ref{algebra-lemma-compose-finite-type}를 보라.
위의 사실과 준콤팩트이고 준분리인 사상들의 합성이 다시 준콤팩트이고
준분리라는 사실을 함께 사용하자. 스킴, 보조정리
\ref{schemes-lemma-composition-quasi-compact}와
\ref{schemes-lemma-separated-permanence}를 보라. 그러면 유한 표시인
사상들의 합성도 유한 표시임을 알 수 있다.
\end{proof}

\begin{lemma}
\label{morphisms-lemma-base-change-finite-presentation}
국소적으로 유한 표시인 사상의 기저변환은 국소적으로 유한 표시이다.
유한 표시인 사상에 대해서도 같은 명제가 성립한다.
\end{lemma}

\begin{proof}
보조정리 \ref{morphisms-lemma-locally-finite-presentation-characterize}의 증명에서
유한 표시라는 성질이 환 사상의 국소적 성질임을 보였다.
따라서 보조정리의 첫 번째 명제는 보조정리
\ref{morphisms-lemma-composition-type-P}와 유한 표시라는 성질이 기저변환에 대해
안정적인 환 사상의 성질이라는 사실을 함께 적용하면 따른다.
대수, 보조정리 \ref{algebra-lemma-base-change-finiteness}를 보라.
위의 사실과 준콤팩트이고 준분리인 사상의 기저변환이 다시 준콤팩트이고
준분리라는 사실을 함께 사용하자. 스킴, 보조정리
\ref{schemes-lemma-quasi-compact-preserved-base-change}와
\ref{schemes-lemma-separated-permanence}를 보라. 그러면 유한 표시인
사상의 기저변환도 유한 표시임을 알 수 있다.
\end{proof}

\begin{lemma}
\label{morphisms-lemma-open-immersion-locally-finite-presentation}
모든 열린몰입은 국소적으로 유한 표시이다.
\end{lemma}

\begin{proof}
열린몰입은 국소 동형이므로 성립한다.
\end{proof}

\begin{lemma}
\label{morphisms-lemma-quasi-compact-open-immersion-finite-presentation}
열린몰입이 유한 표시인 것은 준콤팩트인 것과 동치이다.
\end{lemma}

\begin{proof}
열린몰입이 국소적으로 유한 표시임은 이미 보였다
(보조정리 \ref{morphisms-lemma-open-immersion-locally-finite-presentation}).
몰입이 분리이고 따라서 준분리임도 이미 보았다
(스킴, 보조정리 \ref{schemes-lemma-immersions-monomorphisms}).
이 사실들과 정의 \ref{morphisms-definition-finite-presentation}에서 보조정리가 따른다.
\end{proof}

\begin{lemma}
\label{morphisms-lemma-closed-immersion-finite-presentation}
\begin{slogan}
유한 표시인 닫힌몰입은 유한형 준연접 아이디얼 층에 대응한다.
\end{slogan}
닫힌몰입 $i : Z \to X$가 유한 표시인 것은 대응하는 준연접 아이디얼 층
$\mathcal{I} = \Ker(\mathcal{O}_X \to i_*\mathcal{O}_Z)$가
$\mathcal{O}_X$-가군으로서 유한형인 것과 동치이다.
\end{lemma}

\begin{proof}
임의의 아핀 열린집합 $\Spec(R) \subset X$ 위에서
$i^{-1}(\Spec(R)) = \Spec(R/I)$이고
$\mathcal{I} = \widetilde{I}$이다. 더 나아가 $\mathcal{I}$가 유한형인 것은
모든 이러한 아핀 열린집합에 대해 $I$가 유한 $R$-가군인 것과 동치이다
(성질, 보조정리 \ref{properties-lemma-finite-type-module}를 보라).
또한 $R/I$가 $R$ 위에서 유한 표시인 것은 $I$가 유한 $R$-가군인 것과
동치이다. 따라서 결론을 얻는다.
\end{proof}

\begin{lemma}
\label{morphisms-lemma-finite-presentation-finite-type}
국소적으로 유한 표시인 사상은 국소적으로 유한형이다.
유한 표시인 사상은 유한형이다.
\end{lemma}

\begin{proof}
생략한다.
\end{proof}

\begin{lemma}
\label{morphisms-lemma-noetherian-finite-type-finite-presentation}
\begin{slogan}
국소 뇌터 기저 위에서는 유한형과 유한 표시가 같다.
\end{slogan}
$f : X \to S$를 사상이라 하자.
\begin{enumerate}
\item $S$가 국소 뇌터이고 $f$가 국소적으로 유한형이면
$f$는 국소적으로 유한 표시이다.
\item $S$가 국소 뇌터이고 $f$가 유한형이면 $f$는 유한 표시이다.
\end{enumerate}
\end{lemma}

\begin{proof}
첫 번째 명제는 뇌터 환 위의 유한형 환이 유한 표시라는 사실에서 따른다.
대수, 보조정리
\ref{algebra-lemma-Noetherian-finite-type-is-finite-presentation}를 보라.
$f$가 유한형이고 $S$가 국소 뇌터라고 하자. 그러면 (1)에 의해
$f$는 준콤팩트이고 국소적으로 유한 표시이다. 따라서 $f$가
준분리임을 보이면 충분하다. 이는 보조정리
\ref{morphisms-lemma-finite-type-Noetherian-quasi-separated}와 보조정리
\ref{morphisms-lemma-finite-presentation-finite-type}에서 따른다.
\end{proof}

\begin{lemma}
\label{morphisms-lemma-finite-presentation-quasi-compact-quasi-separated}
$S$를 준콤팩트이고 준분리인 스킴이라 하자.
$X$가 $S$ 위에서 유한 표시이면 $X$는 준콤팩트이고 준분리이다.
\end{lemma}

\begin{proof}
생략한다.
\end{proof}

\begin{lemma}
\label{morphisms-lemma-finite-presentation-permanence}
$f : X \to Y$를 $S$ 위의 스킴 사상이라 하자.
\begin{enumerate}
\item $X$가 $S$ 위에서 국소적으로 유한 표시이고 $Y$가 $S$ 위에서
국소적으로 유한형이면 $f$는 국소적으로 유한 표시이다.
\item $X$가 $S$ 위에서 유한 표시이고 $Y$가 준분리이며 $S$ 위에서
국소적으로 유한형이면 $f$는 유한 표시이다.
\end{enumerate}
\end{lemma}

\begin{proof}
(1)을 증명하자. 보조정리
\ref{morphisms-lemma-locally-finite-presentation-characterize}를 통하면 이는 다음
대수적 사실로 옮겨진다. 환 사상 $A \to B \to C$가 주어지고
$A \to C$가 유한 표시이며 $A \to B$가 유한형이면
$B \to C$는 유한 표시이다. 대수, 보조정리
\ref{algebra-lemma-compose-finite-type}를 보라.

\medskip\noindent
(2)는 (1)과 스킴, 보조정리
\ref{schemes-lemma-compose-after-separated} 및
\ref{schemes-lemma-quasi-compact-permanence}에서 따른다.
\end{proof}

\begin{lemma}
\label{morphisms-lemma-diagonal-morphism-finite-type}
$f : X \to Y$를 대각사상 $\Delta : X \to X \times_Y X$를 갖는
스킴의 사상이라 하자. $f$가 국소적으로 유한형이면 $\Delta$는
국소적으로 유한 표시이다. $f$가 준분리이고 국소적으로 유한형이면
$\Delta$는 유한 표시이다.
\end{lemma}

\begin{proof}
$\Delta$는 $X$ 위의 스킴 사상임에 유의하자(두 번째 사영
$X \times_Y X \to X$를 통한다). $f$가 국소적으로 유한형이라고 하자.
$X$는 $X$ 위에서 유한 표시이고, $X \times_Y X$는
$X$ 위에서 국소적으로 유한형이다(보조정리
\ref{morphisms-lemma-base-change-finite-type}). 따라서 첫 번째 명제는
보조정리 \ref{morphisms-lemma-finite-presentation-permanence}에서 따른다.
두 번째 명제는 첫 번째 명제와 정의들, 그리고 대각사상이 모노사상이므로
분리라는 사실에서 따른다(스킴, 보조정리
\ref{schemes-lemma-monomorphism-separated}).
\end{proof}








\section{구성가능 집합}
\label{morphisms-section-constructible}

\noindent
스킴의 구성가능 집합과 국소 구성가능 집합은 성질, 절
\ref{properties-section-constructible}에서 다루었다. 이 절에서는
(국소) 구성가능 집합의 상과 역상에 관한 결과들을 증명한다.
주요 결과는 유한 표시인 사상 아래 국소 구성가능 집합의 상이
국소 구성가능임을 말하는 슈발레 정리이다.

\begin{lemma}
\label{morphisms-lemma-inverse-image-constructible}
$f : X \to Y$를 스킴의 사상이라 하자.
$E \subset Y$를 부분집합이라 하자.
$E$가 $Y$에서 (국소) 구성가능이면 $f^{-1}(E)$는 $X$에서
(국소) 구성가능이다.
\end{lemma}

\begin{proof}
모든 구성가능 부분집합의 역상이 구성가능임을 보이려면 모든
역콤팩트 열린집합 $V$가 $Y$에 있을 때 그 역상이 $X$에서
역콤팩트임을 보이면 충분하다.
위상수학, 보조정리
\ref{topology-lemma-inverse-images-constructibles}를 보라.
$V$가 $Y$에서 역콤팩트라는 것은 열린몰입 $V \to Y$가
준콤팩트라는 뜻이다. 따라서 그 기저변환
$f^{-1}(V) = X \times_Y V \to X$도 준콤팩트이다. 스킴, 보조정리
\ref{schemes-lemma-quasi-compact-preserved-base-change}를 보라.
그러므로 $f^{-1}(V)$는 $X$에서 역콤팩트이다.
$E$가 $Y$에서 국소 구성가능이라고 하자. $x \in X$를 택하자.
아핀 근방 $V$를 $f(x)$에 대해 택하고, 아핀 근방 $U \subset X$를
$x$에 대해 택하되
$f(U) \subset V$가 되도록 택하자. 따라서
$f|_U : U \to V$를 $V$로 가는 사상으로 생각한다. 성질, 보조정리
\ref{properties-lemma-locally-constructible}에 의해 $E \cap V$는
$V$에서 구성가능하다. 구성가능한 경우에 의해
$(f|_U)^{-1}(E \cap V)$는 $U$에서 구성가능하다.
마지막으로 $(f|_U)^{-1}(E \cap V) = f^{-1}(E) \cap U$이므로 끝난다.
\end{proof}

\begin{lemma}
\label{morphisms-lemma-chevalley}
$f : X \to Y$를 스킴의 사상이라 하자. 다음을 가정하자.
\begin{enumerate}
\item $f$는 준콤팩트이고 국소적으로 유한 표시이다.
\item $Y$는 준콤팩트이고 준분리이다.
\end{enumerate}
그러면 $X$의 모든 구성가능 부분집합의 상은 $Y$에서 구성가능이다.
\end{lemma}

\begin{proof}
성질, 보조정리
\ref{properties-lemma-constructible-quasi-compact-quasi-separated}에 의해
$Y$가 아핀인 경우에만 이 보조정리를 증명하면 충분하다.
이 경우 $X$는 준콤팩트이다. 따라서
$X = U_1 \cup \ldots \cup U_n$로 쓸 수 있으며, 각 $U_i$는
$X$의 아핀 열린집합이다.
$E \subset X$가 구성가능이면 각 $E \cap U_i$도 구성가능이다.
위상수학, 보조정리
\ref{topology-lemma-open-immersion-constructible-inverse-image}를 보라.
$f(E) = \bigcup f(E \cap U_i)$이고 구성가능 집합들의 유한 합집합은
구성가능이므로, $X$가 아핀인 경우로 환원된다. 이 경우에는
대수, 정리 \ref{algebra-theorem-chevalley}가 바로 그 결과이다.
\end{proof}

\begin{theorem}[슈발레 정리]
\label{morphisms-theorem-chevalley}
\begin{reference}
\cite[IV, Theorem 1.8.4]{EGA}
\end{reference}
$f : X \to Y$를 스킴의 사상이라 하자.
$f$가 준콤팩트이고 국소적으로 유한 표시라고 가정하자.
그러면 모든 국소 구성가능 부분집합의 상은 국소 구성가능이다.
\end{theorem}

\begin{proof}
$E \subset X$를 국소 구성가능이라 하자.
$f(E)$도 국소 구성가능임을 보여야 한다. $f(E) \cap V$가
임의의 아핀 열린집합 $V \subset Y$에 대해 구성가능임을 보이겠다.
따라서 $Y$가 아핀인 경우로 환원된다. 이 경우 $X$는 준콤팩트이다.
그러므로 $X = U_1 \cup \ldots \cup U_n$로 쓸 수 있으며,
각 $U_i$는 $X$의 아핀 열린집합이다.
$E \subset X$가 국소 구성가능이면 각 $E \cap U_i$는 구성가능이다.
성질, 보조정리 \ref{properties-lemma-locally-constructible}를 보라.
$f(E) = \bigcup f(E \cap U_i)$이고 구성가능 집합들의 유한 합집합은
구성가능이므로, $X$가 아핀인 경우로 환원된다. 이 경우에는
대수, 정리 \ref{algebra-theorem-chevalley}가 바로 그 결과이다.
\end{proof}

\begin{lemma}
\label{morphisms-lemma-constructible-containing-open}
$X$를 스킴이라 하고 $x \in X$라 하자. $E \subset X$를
국소 구성가능 부분집합이라 하자. $\{x' \mid x' \leadsto x\} \subset E$이면
$E$는 $x$의 열린 근방을 포함한다.
\end{lemma}

\begin{proof}
$\{x' \mid x' \leadsto x\} \subset E$라고 가정하자.
$X$가 아핀이라고 가정해도 된다. 이 경우 $E$는 구성가능이다.
성질, 보조정리 \ref{properties-lemma-locally-constructible}를 보라.
특히 여집합 $E^c$도 구성가능이다. 대수, 보조정리
\ref{algebra-lemma-constructible-is-image}에 의해 아핀 스킴의 사상
$f : Y \to X$를 $E^c = f(Y)$가 되도록 찾을 수 있다.
$Z \subset X$를 $f$의 스킴론적 상이라 하자. 보조정리
\ref{morphisms-lemma-reach-points-scheme-theoretic-image}와 가정
$\{x' \mid x' \leadsto x\} \subset E$에 의해 $x \not \in Z$이다.
따라서 $X \setminus Z \subset E$는 $x$의 열린 근방이며 $E$에 포함된다.
\end{proof}




\section{열린 사상}
\label{morphisms-section-open}

\begin{definition}
\label{morphisms-definition-open}
$f : X \to S$를 사상이라 하자.
\begin{enumerate}
\item 바탕 위상공간 사이의 사상이 열린사상이면 $f$를 {\it 열린} 사상이라 한다.
\item $f$가 {\it 보편적으로 열린} 사상이라고 함은 스킴의 임의의 사상
$S' \to S$에 대해 기저변환 $f' : X_{S'} \to S'$가 열린사상인 경우이다.
\end{enumerate}
\end{definition}

\noindent
위상수학, 보조정리
\ref{topology-lemma-closed-open-map-specialization}에 따르면 위상공간의
특정한 열린사상을 따라 일반화를 들어 올릴 수 있다. 사실 스킴의 임의의
열린 사상을 따라서도 일반화를 들어 올릴 수 있다(보조정리
\ref{morphisms-lemma-open-generizing}를 보라). 또한 스킴의 평탄 사상을 따라서도
일반화를 들어 올릴 수 있음을 보일 것이다(보조정리
\ref{morphisms-lemma-generalizations-lift-flat}). 이 사실이 다시 그 사상이
열린사상임을 뜻하는 경우가 있다.

\begin{lemma}
\label{morphisms-lemma-locally-finite-presentation-universally-open}
$f : X \to S$를 사상이라 하자.
\begin{enumerate}
\item $f$가 국소적으로 유한 표시이고 $f$를 따라 일반화를 들어 올릴 수
있으면 $f$는 열린사상이다.
\item $f$가 국소적으로 유한 표시이고 $f$의 모든 기저변환을 따라
일반화를 들어 올릴 수 있으면 $f$는 보편적으로 열린 사상이다.
\end{enumerate}
\end{lemma}

\begin{proof}
첫 번째 명제만 증명하면 충분하다. $X$와 $S$가 모두 아핀인 경우로
환원된다. 이 경우에는 대수, 보조정리
\ref{algebra-lemma-going-up-down-specialization}와 명제
\ref{algebra-proposition-fppf-open}에서 결과가 따른다.
\end{proof}

\noindent
유한 표시인 평탄 사상의 경우에는 보조정리 \ref{morphisms-lemma-fppf-open}도 보라.

\begin{lemma}
\label{morphisms-lemma-composition-open}
(보편적으로) 열린 사상들의 합성은 (보편적으로) 열린 사상이다.
\end{lemma}

\begin{proof}
생략한다.
\end{proof}

\begin{lemma}
\label{morphisms-lemma-scheme-over-field-universally-open}
$k$를 체라 하고 $X$를 $k$ 위의 스킴이라 하자.
구조사상 $X \to \Spec(k)$는 보편적으로 열린 사상이다.
\end{lemma}

\begin{proof}
$S \to \Spec(k)$를 사상이라 하자.
기저변환 $X_S \to S$가 열린사상임을 보여야 한다.
이 문제는 $S$와 $X$에 대해 국소적이므로 $S$와 $X$가 아핀이라고
가정해도 된다. 이 경우에는 대수, 보조정리
\ref{algebra-lemma-map-into-tensor-algebra-open}에서 결과가 따른다.
\end{proof}

\begin{lemma}
\label{morphisms-lemma-open-generizing}
\begin{reference}
\cite[IV, Corollary 1.10.4]{EGA}의 (a) $\Rightarrow$ (b)에서 따른다.
\end{reference}
$\varphi : X \to Y$를 스킴의 사상이라 하자.
$\varphi$가 열린사상이면 $\varphi$는 일반화 보존 사상이다
(즉 $\varphi$를 따라 일반화를 들어 올릴 수 있다).
$\varphi$가 보편적으로 열린 사상이면 $\varphi$는 보편적으로
일반화 보존 사상이다.
\end{lemma}

\begin{proof}
$\varphi$가 열린사상이라고 가정하자.
점들의 특수화 $y' \leadsto y$가 $Y$에서 주어졌다고 하자.
$x \in X$가 $\varphi(x) = y$를 만족한다고 하자.
아핀 열린집합 $U \subset X$와 $V \subset Y$를
$\varphi(U) \subset V$와 $x \in U$가 성립하도록 택하자.
그러면 $y' \in V$도 성립한다. 따라서 $X$를 $U$로, $Y$를 $V$로
바꾸어 $X$, $Y$가 아핀이라고 가정해도 된다. 아핀인 경우는
대수, 보조정리 \ref{algebra-lemma-open-going-down}와
대수, 보조정리 \ref{algebra-lemma-going-up-down-specialization}에서 따른다.
\end{proof}

\begin{lemma}
\label{morphisms-lemma-descent-quasi-compact}
$f : X \to Y$를 스킴의 사상이라 하자.
$g : Y' \to Y$가 열린 전사이고 기저변환
$f' : X' \to Y'$가 준콤팩트라고 하자. 그러면 $f$는 준콤팩트이다.
\end{lemma}

\begin{proof}
$V \subset Y$를 준콤팩트 열린집합이라 하자. $g$가 열린 전사이므로
준콤팩트 열린집합 $W' \subset W$를 $g(W') = V$가 되도록 찾을 수 있다.
가정에 의해 $(f')^{-1}(W')$는 준콤팩트이다.
$(f')^{-1}(W')$의 $X$에서의 상은 $f^{-1}(V)$와 같다. 보조정리
\ref{morphisms-lemma-when-point-maps-to-pair}를 보라. 따라서 $f^{-1}(V)$는
준콤팩트 공간의 상이므로 준콤팩트이다. 위상수학, 보조정리
\ref{topology-lemma-image-quasi-compact}를 보라. 따라서 $f$는 준콤팩트이다.
\end{proof}





\section{몫사상}
\label{morphisms-section-submersive}

\begin{definition}
\label{morphisms-definition-submersive}
$f : X \to Y$를 스킴의 사상이라 하자.
\begin{enumerate}
\item 바탕 위상공간 사이의 연속사상이 몫사상이면 $f$를
{\it 몫사상}\footnote{이는 미분다양체의 침몰사상 개념과 매우 다르다.}이라 한다.
위상수학, 정의 \ref{topology-definition-submersive}를 보라.
\item $f$를 {\it 보편적 몫사상}이라고 함은 스킴의 모든 사상
$Y' \to Y$에 대해 기저변환 $Y' \times_Y X \to Y'$가 몫사상인 경우이다.
\end{enumerate}
\end{definition}

\noindent
몫사상은 특히 전사임에 유의하자.

\begin{lemma}
\label{morphisms-lemma-base-change-universally-submersive}
스킴의 보편적 몫사상을 스킴의 임의의 사상으로 기저변환하면
다시 보편적 몫사상이다.
\end{lemma}

\begin{proof}
정의에서 곧바로 따른다.
\end{proof}

\begin{lemma}
\label{morphisms-lemma-composition-universally-submersive}
스킴의 두 (보편적) 몫사상의 합성은 (보편적) 몫사상이다.
\end{lemma}

\begin{proof}
생략한다.
\end{proof}










\section{평탄 사상}
\label{morphisms-section-flat}

\noindent
평탄성은 대수기하학에서 가장 중요한 기술적 도구 가운데 하나이다.
이 절에서는 이 개념을 도입한다. 보조정리 \ref{morphisms-lemma-fppf-open}을 제외하고는
의도적으로 직접적인 관찰들만 다룬다. 매우 중요한 부류의 결과인 평탄성의
판정법은 대수, 절 \ref{algebra-section-criteria-flatness},
\ref{algebra-section-flatness-artinian},
\ref{algebra-section-more-flatness-criteria}와 사상 추가, 절
\ref{more-morphisms-section-criterion-flat-fibres}에서 다룬다.
스킴의 평탄 사상에 관한 고급 내용을 전담하는 장도 있다. 평탄성 추가, 절
\ref{flat-section-introduction}을 보라.

\medskip\noindent
가군 $M$이 환 $R$ 위에 있고 함자
$-\otimes_R M : \text{Mod}_R \to \text{Mod}_R$가 완전할 때
{\it 평탄}이라고 한다는 것을 상기하자. 환 사상 $R \to A$는
$A$가 $R$-가군으로서 평탄일 때 {\it 평탄}이라고 한다.
대수, 정의 \ref{algebra-definition-flat}을 보라.

\begin{definition}
\label{morphisms-definition-flat}
$f : X \to S$를 스킴의 사상이라 하자.
$\mathcal{F}$를 $\mathcal{O}_X$-가군의 준연접층이라 하자.
\begin{enumerate}
\item $f$가 점 $x \in X$에서 {\it 평탄}이라고 함은 국소환
$\mathcal{O}_{X, x}$가 국소환 $\mathcal{O}_{S, f(x)}$ 위에서
평탄인 경우이다.
\item $\mathcal{F}$가 $S$ 위에서 점 $x \in X$에서 {\it 평탄}이라고 함은
줄기 $\mathcal{F}_x$가 평탄 $\mathcal{O}_{S, f(x)}$-가군인 경우이다.
\item $f$가 {\it 평탄}이라고 함은 $f$가 $X$의 모든 점에서 평탄인 경우이다.
\item $\mathcal{F}$가 $S$ 위에 {\it 평탄}이라고 함은 $\mathcal{F}$가
$S$ 위에서 평탄이고, 이를 모든 점 $x$가 $X$를 달릴 때 요구하는 경우이다.
\end{enumerate}
\end{definition}

\noindent
따라서 $f$가 평탄인 것은 구조층 $\mathcal{O}_X$가 $S$ 위에서
평탄인 것과 동치이다.

\begin{lemma}
\label{morphisms-lemma-flat-module-characterize}
$f : X \to S$를 스킴의 사상이라 하자.
$\mathcal{F}$를 $\mathcal{O}_X$-가군의 준연접층이라 하자.
다음 조건들은 서로 동치이다.
\begin{enumerate}
\item 층 $\mathcal{F}$는 $S$ 위에서 평탄이다.
\item 아핀 열린집합 $U \subset X$, $V \subset S$가
$f(U) \subset V$를 만족하면 $\mathcal{O}_S(V)$-가군
$\mathcal{F}(U)$는 평탄이다.
\item 열린 덮개 $S = \bigcup_{j \in J} V_j$와 열린 덮개
$f^{-1}(V_j) = \bigcup_{i \in I_j} U_i$가 존재하여 각 가군
$\mathcal{F}|_{U_i}$가 $V_j$ 위에서 평탄이다. 이는 모든
$j\in J, i\in I_j$에 대해 성립한다.
\item 아핀 열린 덮개 $S = \bigcup_{j \in J} V_j$와 아핀 열린 덮개
$f^{-1}(V_j) = \bigcup_{i \in I_j} U_i$가 존재하여
$\mathcal{F}(U_i)$가 평탄 $\mathcal{O}_S(V_j)$-가군이다. 이는 모든
$j\in J, i\in I_j$에 대해 성립한다.
\end{enumerate}
더 나아가 $\mathcal{F}$가 $S$ 위에서 평탄이면 $U \subset X$,
$V \subset S$인 임의의 열린 부분스킴이 $f(U) \subset V$를 만족할 때
제한 $\mathcal{F}|_U$는 $V$ 위에서 평탄이다.
\end{lemma}

\begin{proof}
$R \to A$를 환 사상이라 하고 $M$을 $A$-가군이라 하자.
$M$이 $R$-평탄이면 모든 소 아이디얼 $\mathfrak q$에 대해
$M_{\mathfrak q}$는 $R_{\mathfrak p}$ 위에서 평탄이다. 여기서
$\mathfrak p$는 $R$의 소 아이디얼 중 $\mathfrak q$ 아래에 놓이는 것이다.
역으로 $M_{\mathfrak q}$가 $R_{\mathfrak p}$ 위에서 평탄이고 이것이
모든 소 아이디얼 $\mathfrak q$가 $A$를 달릴 때 성립하면
$M$은 $R$ 위에서 평탄이다. 대수, 보조정리
\ref{algebra-lemma-flat-localization}을 보라.
이 동치에서 보조정리의 명제들이 쉽게 따른다.
\end{proof}

\begin{lemma}
\label{morphisms-lemma-flat-characterize}
$f : X \to S$를 스킴의 사상이라 하자.
다음 조건들은 서로 동치이다.
\begin{enumerate}
\item 사상 $f$는 평탄이다.
\item 아핀 열린집합 $U \subset X$, $V \subset S$가
$f(U) \subset V$를 만족하면 환 사상
$\mathcal{O}_S(V) \to \mathcal{O}_X(U)$는 평탄이다.
\item 열린 덮개 $S = \bigcup_{j \in J} V_j$와 열린 덮개
$f^{-1}(V_j) = \bigcup_{i \in I_j} U_i$가 존재하여 각 사상
$U_i \to V_j$, $j\in J, i\in I_j$가 평탄이다.
\item 아핀 열린 덮개 $S = \bigcup_{j \in J} V_j$와 아핀 열린 덮개
$f^{-1}(V_j) = \bigcup_{i \in I_j} U_i$가 존재하여
$\mathcal{O}_S(V_j) \to \mathcal{O}_X(U_i)$가 모든
$j\in J, i\in I_j$에 대해 평탄이다.
\end{enumerate}
더 나아가 $f$가 평탄이면 $U \subset X$, $V \subset S$인
임의의 열린 부분스킴이 $f(U) \subset V$를 만족할 때 제한
$f|_U : U \to V$도 평탄이다.
\end{lemma}

\begin{proof}
이는 위의 보조정리 \ref{morphisms-lemma-flat-module-characterize}의 특수한 경우이다.
\end{proof}

\begin{lemma}
\label{morphisms-lemma-pushforward-flat-affine}
$f : X \to Y$를 기저 스킴 $S$ 위의 아핀 사상이라 하자.
$\mathcal{F}$를 준연접 $\mathcal{O}_X$-가군이라 하자.
$\mathcal{F}$가 $S$ 위에서 평탄인 것은 $f_*\mathcal{F}$가
$S$ 위에서 평탄인 것과 동치이다.
\end{lemma}

\begin{proof}
보조정리 \ref{morphisms-lemma-flat-module-characterize}와 $f$가 아핀 사상이라는
사실에 의해 아핀인 경우로 환원된다.
$X \to Y \to S$가 환 사상 $C \leftarrow B \leftarrow A$에
대응한다고 하자. $N$을 $C$-가군 중 $\mathcal{F}$에 대응하는 것이라 하자.
$f_*\mathcal{F}$는 $N$에 대응하며, 여기서 이를 $B$-가군으로 본다는 것을
상기하자.
스킴, 보조정리 \ref{schemes-lemma-widetilde-pullback}를 보라.
따라서 결과는 자명하다.
\end{proof}

\begin{lemma}
\label{morphisms-lemma-composition-module-flat}
$X \to Y \to Z$를 스킴의 사상이라 하자. $\mathcal{F}$를 준연접
$\mathcal{O}_X$-가군이라 하자. $x \in X$를 택하고 그 상 $y$가
$Y$에 놓인다고 하자. $\mathcal{F}$가 $Y$ 위에서 $x$에서 평탄이고
$Y$가 $Z$ 위에서 $y$에서 평탄이면 $\mathcal{F}$는 $Z$ 위에서
$x$에서 평탄이다.
\end{lemma}

\begin{proof}
대수, 보조정리 \ref{algebra-lemma-composition-flat}을 보라.
\end{proof}

\begin{lemma}
\label{morphisms-lemma-composition-flat}
평탄 사상들의 합성은 평탄하다.
\end{lemma}

\begin{proof}
이는 보조정리 \ref{morphisms-lemma-composition-module-flat}의 특수한 경우이다.
\end{proof}

\begin{lemma}
\label{morphisms-lemma-base-change-module-flat}
$f : X \to S$를 스킴의 사상이라 하자.
$\mathcal{F}$를 $\mathcal{O}_X$-가군의 준연접층이라 하자.
$g : S' \to S$를 스킴의 사상이라 하자.
$g' : X' = X_{S'} \to X$를 사영이라 하자.
$x' \in X'$를 택하고 그 상이 $x = g'(x') \in X$라고 하자.
$\mathcal{F}$가 $S$ 위에서 $x$에서 평탄이면
$(g')^*\mathcal{F}$는 $S'$ 위에서 $x'$에서 평탄이다.
특히 $\mathcal{F}$가 $S$ 위에서 평탄이면
$(g')^*\mathcal{F}$는 $S'$ 위에서 평탄이다.
\end{lemma}

\begin{proof}
대수, 보조정리 \ref{algebra-lemma-flat-base-change}를 보라.
\end{proof}

\begin{lemma}
\label{morphisms-lemma-base-change-flat}
평탄 사상의 기저변환은 평탄하다.
\end{lemma}

\begin{proof}
이는 보조정리 \ref{morphisms-lemma-base-change-module-flat}의 특수한 경우이다.
\end{proof}

\begin{lemma}
\label{morphisms-lemma-generalizations-lift-flat}
$f : X \to S$를 스킴의 평탄 사상이라 하자.
그러면 $f$를 따라 일반화를 들어 올릴 수 있다. 위상수학, 정의
\ref{topology-definition-lift-specializations}를 보라.
\end{lemma}

\begin{proof}
대수, 절 \ref{algebra-section-going-up}을 보라.
\end{proof}

\begin{lemma}
\label{morphisms-lemma-fppf-open}
국소적으로 유한 표시인 평탄 사상은 보편적으로 열린 사상이다.
\end{lemma}

\begin{proof}
이는 위의 보조정리 \ref{morphisms-lemma-generalizations-lift-flat}과
보조정리 \ref{morphisms-lemma-locally-finite-presentation-universally-open}에서 따른다.
다음과 같이 직접 논증할 수도 있다.

\medskip\noindent
$f : X \to S$가 평탄이고 국소적으로 유한 표시라고 하자.
보조정리 \ref{morphisms-lemma-base-change-flat}과
\ref{morphisms-lemma-base-change-finite-presentation}에 의해 $f$의 모든 기저변환은
평탄이고 국소적으로 유한 표시이다. 따라서 $f$가 열린사상임을 보이면 충분하다.
$f$가 열린사상임을 보이려면 $X$를 아핀 열린집합들
$X = \bigcup U_i$로 덮어서 $U_i \to S$가 열린사상이 되게 하면 충분하다.
$X$를 아핀 열린집합 $U_i \subset X$들로 덮되 각 $U_i$가 아핀 열린집합
$V_i \subset S$ 안으로 사상되고, 유도된 환 사상
$\mathcal{O}_S(V_i) \to \mathcal{O}_X(U_i)$가 평탄이고 유한 표시가 되게
할 수 있다(보조정리 \ref{morphisms-lemma-flat-characterize}와
\ref{morphisms-lemma-locally-finite-presentation-characterize}). 그러면
$U_i \to V_i$는 열린사상이다. 대수, 명제
\ref{algebra-proposition-fppf-open}을 보라. 이로써 증명이 완성된다.
\end{proof}

\begin{lemma}
\label{morphisms-lemma-pf-flat-module-open}
$f : X \to Y$를 스킴의 사상이라 하자.
$\mathcal{F}$를 준연접 $\mathcal{O}_X$-가군이라 하자.
$f$가 국소적으로 유한 표시이고, $\mathcal{F}$가 유한형이며,
$X = \text{Supp}(\mathcal{F})$이고, $\mathcal{F}$가 $Y$ 위에서
평탄이라고 가정하자. 그러면 $f$는 보편적으로 열린 사상이다.
\end{lemma}

\begin{proof}
보조정리 \ref{morphisms-lemma-base-change-module-flat},
\ref{morphisms-lemma-base-change-finite-presentation}, 그리고
\ref{morphisms-lemma-support-finite-type}에 의해 가정들은 기저변환 아래 보존된다.
보조정리 \ref{morphisms-lemma-locally-finite-presentation-universally-open}에 의해
$f$를 따라 일반화를 들어 올릴 수 있음을 보이면 충분하다.
이는 대수, 보조정리 \ref{algebra-lemma-going-down-flat-module}에서 따른다.
\end{proof}

\begin{lemma}
\label{morphisms-lemma-fpqc-quotient-topology}
\begin{reference}
\cite[Expose VIII, Corollaire 4.3]{SGA1} 및
\cite[IV, Corollaire 2.3.12]{EGA}
\end{reference}
$f : X \to Y$를 준콤팩트인 평탄 전사라 하자.
부분집합 $T \subset Y$가 열린집합(각각 닫힌집합)인 것은
$f^{-1}(T)$가 열린집합(각각 닫힌집합)인 것과 동치이다.
달리 말하면 $f$는 몫사상이다.
\end{lemma}

\begin{proof}
이 문제는 $Y$에 대해 국소적이므로 $Y$가 아핀이라고 가정해도 된다.
이 경우 $X$는 $f$가 준콤팩트이므로 준콤팩트이다.
$X = X_1 \cup \ldots \cup X_n$를 아핀 열린집합들의 유한 합집합으로 쓰자.
그러면 $f' : X' = X_1 \amalg \ldots \amalg X_n \to Y$는 아핀 스킴의
평탄 전사이다. $T \subset Y$에 대해
$(f')^{-1}(T) = f^{-1}(T) \cap X_1 \amalg \ldots \amalg f^{-1}(T) \cap X_n$
임에 유의하자. 따라서 $f^{-1}(T)$가 열린집합인 것은
$(f')^{-1}(T)$가 열린집합인 것과 동치이다. 그러므로 $X$와 $Y$가
모두 아핀이라고 가정해도 된다.

\medskip\noindent
$f : \Spec(B) \to \Spec(A)$를 평탄 환 사상 $A \to B$에 대응하는
아핀 스킴의 전사라 하자. $f^{-1}(T)$가 닫힌집합이라고 하자.
이를 $f^{-1}(T) = V(J)$로 쓰자. 여기서 $J \subset B$는 아이디얼이다.
그러면 $T = f(f^{-1}(T)) = f(V(J))$는
$\Spec(B/J) \to \Spec(A)$의 상이다(여기서는 $f$가 전사임을 사용한다).
한편 $f$를 따라 일반화를 들어 올릴 수 있다(보조정리
\ref{morphisms-lemma-generalizations-lift-flat}). 따라서 위상수학, 보조정리
\ref{topology-lemma-lift-specializations-images}에 의해
$Y \setminus T = f(X \setminus f^{-1}(T))$는 일반화 아래 안정적이다.
따라서 $T$는 특수화 아래 안정적이다(위상수학, 보조정리
\ref{topology-lemma-open-closed-specialization}). 그러므로 대수, 보조정리
\ref{algebra-lemma-image-stable-specialization-closed}에 의해 $T$는 닫혀 있다.
\end{proof}

\begin{lemma}
\label{morphisms-lemma-flat-permanence}
$h : X \to Y$를 $S$ 위의 스킴 사상이라 하자.
$\mathcal{G}$를 $Y$ 위의 준연접층이라 하자.
$x \in X$를 택하고 $y = h(x) \in Y$라 하자. $h$가 $x$에서 평탄이면
$$
\mathcal{G}\text{ flat over }S\text{ at }y
\Leftrightarrow
h^*\mathcal{G}\text{ flat over }S\text{ at }x.
$$
특히 $h$가 평탄 전사이면 $\mathcal{G}$가 $S$ 위에서 평탄인 것은
$h^*\mathcal{G}$가 $S$ 위에서 평탄인 것과 동치이다. $h$가 평탄 전사이고
$X$가 $S$ 위에서 평탄이면 $Y$도 $S$ 위에서 평탄이다.
\end{lemma}

\begin{proof}
대수, 보조정리
\ref{algebra-lemma-flatness-descends-more-general}을 적용하여 증명할 수 있다.
직접적인 증명을 제시하자. $s \in S$를 $y$의 상이라 하자.
국소환 사상들
$\mathcal{O}_{S, s} \to \mathcal{O}_{Y, y} \to \mathcal{O}_{X, x}$를 생각하자.
가정에 의해 환 사상 $\mathcal{O}_{Y, y} \to \mathcal{O}_{X, x}$는
충실 평탄이다. 대수, 보조정리 \ref{algebra-lemma-local-flat-ff}를 보라.
$N = \mathcal{G}_y$로 두자. 다음에 유의하자.
$h^*\mathcal{G}_x = N \otimes_{\mathcal{O}_{Y, y}} \mathcal{O}_{X, x}$.
층, 보조정리 \ref{sheaves-lemma-stalk-pullback-modules}를 보라.
$M' \to M$을 $\mathcal{O}_{S, s}$-가군의 단사사상이라 하자.
위에서 말한 충실 평탄성에 의해
\begin{align*}
\Ker(
M' \otimes_{\mathcal{O}_{S, s}} N \to M \otimes_{\mathcal{O}_{S, s}} N)
\otimes_{\mathcal{O}_{Y, y}} \mathcal{O}_{X, x} \\
=
\Ker(
M' \otimes_{\mathcal{O}_{S, s}} N
\otimes_{\mathcal{O}_{Y, y}} \mathcal{O}_{X, x}
\to
M \otimes_{\mathcal{O}_{S, s}} N
\otimes_{\mathcal{O}_{Y, y}} \mathcal{O}_{X, x})
\end{align*}
이다. 따라서 보조정리의 동치는 대수, 보조정리
\ref{algebra-lemma-flat}의 두 번째 평탄성 특징에서 따른다.
\end{proof}

\begin{lemma}
\label{morphisms-lemma-flat-pullback-support}
$f : Y \to X$를 스킴의 사상이라 하자. $\mathcal{F}$를 유한형 준연접
$\mathcal{O}_X$-가군이라 하고 그 스킴론적 지지집합을 $Z \subset X$라 하자.
$f$가 평탄이면 $f^{-1}(Z)$는
$f^*\mathcal{F}$의 스킴론적 지지집합이다.
\end{lemma}

\begin{proof}
보조정리 \ref{morphisms-lemma-scheme-theoretic-support}에 주어진 아핀 위의
스킴론적 지지집합 특징을 사용하면 대수, 보조정리
\ref{algebra-lemma-annihilator-flat-base-change}로 환원된다.
\end{proof}

\begin{lemma}
\label{morphisms-lemma-flat-morphism-scheme-theoretically-dense-open}
$f : X \to Y$를 스킴의 평탄 사상이라 하자. $V \subset Y$를
스킴론적으로 조밀한 역콤팩트 열린집합이라 하자. 그러면 $f^{-1}V$는
$X$에서 스킴론적으로 조밀하다.
\end{lemma}

\begin{proof}
보조정리 \ref{morphisms-lemma-characterize-scheme-theoretically-dense}의 특징을 사용한다.
임의의 열린집합 $U \subset X$에 대해 사상
$\mathcal{O}_X(U) \to \mathcal{O}_X(U \cap f^{-1}V)$가 단사임을 보여야 한다.
$U$가 아핀 열린집합이고 어떤 아핀 열린집합 $W \subset Y$ 안으로
사상되는 경우에만 증명하면 충분하다. $W = \Spec(A)$ 및
$U = \Spec(B)$라고 쓰자.
$V \cap W = D(f_1) \cup \ldots \cup D(f_n)$이고 여기서
$f_i \in A$이다. 대수, 보조정리
\ref{algebra-lemma-qc-open}을 보라. 따라서
$B \to B_{f_1} \times \ldots \times B_{f_n}$이 단사임을 보여야 한다.
$A \to A_{f_1} \times \ldots \times A_{f_n}$이 단사이고
$A \to B$가 평탄임이 주어졌다. $B_{f_i} = A_{f_i} \otimes_A B$이므로 끝난다.
\end{proof}

\begin{lemma}
\label{morphisms-lemma-flat-base-change-scheme-theoretic-image}
\begin{slogan}
준콤팩트한 경우 스킴론적 상을 취하는 것은 평탄 기저변환과 교환한다.
\end{slogan}
$f : X \to Y$를 스킴의 평탄 사상이라 하자. $g : V \to Y$를 스킴의
준콤팩트 사상이라 하자. $Z \subset Y$를 $g$의 스킴론적 상이라 하고,
$Z' \subset X$를 기저변환 $V \times_Y X \to X$의 스킴론적 상이라 하자.
그러면 $Z' = f^{-1}Z$이다.
\end{lemma}

\begin{proof}
$Z$는
$\mathcal{I} = \Ker(\mathcal{O}_Y \to g_*\mathcal{O}_V)$로 잘려 나오고,
$Z'$는
$\mathcal{I}' = \Ker(\mathcal{O}_X \to
(V \times_Y X \to X)_*\mathcal{O}_{V \times_Y X})$로 잘려 나옴을 상기하자.
보조정리 \ref{morphisms-lemma-quasi-compact-scheme-theoretic-image}를 보라.
따라서 문제는 $X$와 $Y$에 대해 국소적이고, $X$와 $Y$가 아핀이라고
가정해도 된다. $V$를 $\coprod V_i$로 바꾸어도 됨에 유의하자. 여기서
$V = V_1 \cup \ldots \cup V_n$는 유한 아핀 열린 덮개이다.
따라서 $g$가 아핀이라고 가정해도 된다. 이 경우
$(V \times_Y X \to X)_*\mathcal{O}_{V \times_Y X}$는
$g_*\mathcal{O}_V$의 $f$에 의한 당김이다. $f$가 평탄이므로
$f^*\mathcal{I} = \mathcal{I}'$이고 보조정리가 성립한다.
\end{proof}





\section{평탄 닫힌몰입}
\label{morphisms-section-flat-closed-immersions}

\noindent
스킴의 연결 성분은 항상 열린집합인 것은 아니다. 그러나 항상 표준적인
스킴 구조를 갖는다. 이 절에서 이를 설명한다.

\begin{lemma}
\label{morphisms-lemma-characterize-flat-closed-immersions}
$X$를 스킴이라 하자. $X$의 닫힌 부분스킴에 그 바탕 닫힌 부분집합을
대응시키는 규칙은 다음 전단사를 정의한다.
$$
\left\{
\begin{matrix}
\text{closed subschemes }Z \subset X \\
\text{such that }Z \to X\text{ is flat}
\end{matrix}
\right\}
\leftrightarrow
\left\{
\begin{matrix}
\text{closed subsets }Z \subset X \\
\text{closed under generalizations}
\end{matrix}
\right\}
$$
$Z \subset X$가 그러한 닫힌 부분스킴이면, 스킴의 모든 사상
$g : Y \to X$가 집합론적으로 $g(Y) \subset Z$를 만족할 때
스킴론적으로 $Z$를 통해 분해된다.
\end{lemma}

\begin{proof}
전단사의 아핀 경우는 대수, 보조정리
\ref{algebra-lemma-pure-open-closed-specializations}이다. 일반 스킴 $X$의
경우에는 $X$를 아핀들로 덮고 붙이기를 하여 전단사를 얻는다.
세부사항은 생략한다. 마지막 명제를 위해 사영
$Z \times_{X, g} Y \to Y$를 보자. 이는 평탄인 닫힌몰입이고
(보조정리 \ref{morphisms-lemma-base-change-flat}), 바탕 위상공간 사이에서
전단사이므로 증명의 첫 부분에서 확립한 전단사에 의해 동형이어야 한다.
\end{proof}

\begin{lemma}
\label{morphisms-lemma-flat-closed-immersions-finite-presentation}
유한 표시인 평탄 닫힌몰입은 열린닫힌 부분스킴의 열린몰입이다.
\end{lemma}

\begin{proof}
아핀인 경우는 대수, 보조정리
\ref{algebra-lemma-finitely-generated-pure-ideal}이다.
일반적인 경우에는 $X$를 아핀들로 덮으면 보조정리가 따른다.
세부사항은 생략한다.
\end{proof}

\noindent
연결 성분 $T$가 스킴 $X$에 속하면 이는 일반화 아래 안정적인 닫힌 부분집합임에
유의하자. 따라서 다음 정의는 잘 정의된다.

\begin{definition}
\label{morphisms-definition-scheme-structure-connected-component}
$X$를 스킴이라 하고 $T \subset X$를 연결 성분이라 하자.
$T$ 위의 {\it 표준 스킴 구조}란 $T$ 위의 유일한 스킴 구조 가운데
닫힌몰입 $T \to X$가 평탄이 되는 것을 말한다. 보조정리
\ref{morphisms-lemma-characterize-flat-closed-immersions}를 보라.
\end{definition}

\noindent
앞의 보조정리를 사용하면 모든 유한 평탄 $\mathcal{O}_X$-가군이 언제
유한 국소 자유 가군인지 판정할 수 있다.

\begin{lemma}
\label{morphisms-lemma-finite-flat-is-finite-locally-free}
$X$를 스킴이라 하자. 다음 조건들은 서로 동치이다.
\begin{enumerate}
\item 모든 유한 평탄 준연접 $\mathcal{O}_X$-가군은 유한 국소 자유이다.
\item 일반화 아래 닫혀 있는 모든 닫힌 부분집합 $Z \subset X$는 열려 있다.
\end{enumerate}
\end{lemma}

\begin{proof}
아핀인 경우는 대수, 보조정리
\ref{algebra-lemma-finite-flat-module-finitely-presented}이다.
스킴인 경우가 아핀인 경우에서 직접 따르지는 않으므로 논증을 되풀이한다.

\medskip\noindent
(1)을 가정하자. 닫힌몰입 $i : Z \to X$를 생각하고 $i$가 평탄이라고 하자.
그러면 $i_*\mathcal{O}_Z$는 준연접이고 평탄이므로 (1)에 의해
유한 국소 자유이다. 따라서 $Z = \text{Supp}(i_*\mathcal{O}_Z)$도 열려 있고
(2)가 성립한다. 그러므로 (1) $\Rightarrow$ (2)는 보조정리
\ref{morphisms-lemma-characterize-flat-closed-immersions}의 평탄 닫힌몰입 특징에서 따른다.

\medskip\noindent
역으로 $X$가 (2)를 만족한다고 가정하자.
$\mathcal{F}$를 유한 평탄 준연접 $\mathcal{O}_X$-가군이라 하자.
$Z = \text{Supp}(\mathcal{F})$는 닫혀 있다. 이는 $\mathcal{F}$의 지지집합이다.
가군, 보조정리 \ref{modules-lemma-support-finite-type-closed}를 보라.
한편 $x \leadsto x'$가 특수화이면 대수, 보조정리
\ref{algebra-lemma-finite-flat-local}에 의해 가군 $\mathcal{F}_{x'}$는
$\mathcal{O}_{X, x'}$ 위에서 자유이고,
$$
\mathcal{F}_x =
\mathcal{F}_{x'} \otimes_{\mathcal{O}_{X, x'}} \mathcal{O}_{X, x}.
$$
따라서
$x' \in \text{Supp}(\mathcal{F}) \Rightarrow x \in \text{Supp}(\mathcal{F})$이다.
달리 말하면 지지집합은 일반화 아래 닫혀 있다. $X$가 (2)를 만족하므로
$\mathcal{F}$의 지지집합은 열린닫힌집합이다.
가군 $\wedge^i(\mathcal{F})$, $i = 1, 2, 3, \ldots$도 유한 평탄 준연접
$\mathcal{O}_X$-가군이다. 가군, 절
\ref{modules-section-symmetric-exterior}를 보라. 다음에 유의하자.
$\text{Supp}(\wedge^{i + 1}(\mathcal{F})) \subset
\text{Supp}(\wedge^i(\mathcal{F}))$.
따라서 열린닫힌 부분집합들에 의한 분해
$$
X = U_0 \amalg U_1 \amalg U_2 \amalg \ldots
$$
가 존재하여 $\wedge^i(\mathcal{F})$의 지지집합은
$U_i \cup U_{i + 1} \cup \ldots$이다. 이는 모든 $i$에 대해 성립한다.
$x$를 $X$의 점이라 하고
$x \in U_r$이라 하자. 다음에 유의하자.
$\wedge^i(\mathcal{F})_x \otimes \kappa(x) =
\wedge^i(\mathcal{F}_x \otimes \kappa(x))$.
따라서 $x \in U_r$이면 $\mathcal{F}_x \otimes \kappa(x)$는
$r$차원 벡터공간이다. 나카야마 보조정리에 의해, 대수, 보조정리
\ref{algebra-lemma-NAK}를 보라, 아핀 열린 근방
$U \subset U_r \subset X$를 $x$의 근방으로 택하고 단면
$s_1, \ldots, s_r \in \mathcal{F}(U)$를
택하여 유도된 사상
$$
\mathcal{O}_U^{\oplus r} \longrightarrow \mathcal{F}|_U, \quad
(f_1, \ldots, f_r) \longmapsto \sum f_i s_i
$$
이 전사가 되게 할 수 있다. 이는 $\wedge^r(\mathcal{F}|_U)$가
유한 평탄 준연접 $\mathcal{O}_U$-가군이고 그 지지집합이 $U$ 전체라는 뜻이다.
위의 사실에 의해 이는 한 원소, 곧 $s_1 \wedge \ldots \wedge s_r$로 생성된다.
따라서 $\wedge^r(\mathcal{F}|_U) \cong \mathcal{O}_U/\mathcal{I}$이다.
여기서 $\mathcal{I}$는 어떤 준연접 아이디얼 층이다.
여기서 $\mathcal{O}_U/\mathcal{I}$는 $\mathcal{O}_U$ 위에서 평탄이고
$V(\mathcal{I}) = U$이다. 보조정리
\ref{morphisms-lemma-characterize-flat-closed-immersions}를 적용하면
$\mathcal{I} = 0$을 얻는다. 따라서 $s_1 \wedge \ldots \wedge s_r$는
$\wedge^r(\mathcal{F}|_U)$의 기저이고, 위에 표시한 사상은 전사일 뿐 아니라
단사이다. 이로써 원하는 대로 $\mathcal{F}$가 유한 국소 자유임을 보였다.
\end{proof}






\section{일반점 평탄성}
\label{morphisms-section-generic-flatness}

\noindent
정수적 기저 위의 유한형 스킴은 기저의 조밀한 열린집합 위에서 평탄하다.
대수, 절 \ref{algebra-section-generic-flatness}에서는 뇌터인 판본,
유한 표시 사상에 대한 판본, 그리고 일반적인 판본을 증명했다.
여기서는 일반적인 판본만 진술하고 증명한다. 하지만 기저가 축약이라는
가정만으로도 결과가 성립함을 보이는 명제
\ref{morphisms-proposition-generic-flatness-reduced}가 이 결과를 대체하게 된다.

\begin{proposition}[일반점 평탄성]
\label{morphisms-proposition-generic-flatness}
$f : X \to S$를 스킴의 사상이라 하자.
$\mathcal{F}$를 $\mathcal{O}_X$-가군의 준연접층이라 하자.
다음을 가정하자.
\begin{enumerate}
\item $S$는 정수적이다.
\item $f$는 유한형이다.
\item $\mathcal{F}$는 유한형 $\mathcal{O}_X$-가군이다.
\end{enumerate}
그러면 조밀한 열린 부분스킴 $U \subset S$가 존재하여
$X_U \to U$는 평탄이고 유한 표시이며, $\mathcal{F}|_{X_U}$는
$U$ 위에서 평탄이고 $\mathcal{O}_{X_U}$ 위에서 유한 표시이다.
\end{proposition}

\begin{proof}
$S$가 정수적이므로 기약이고(성질, 보조정리
\ref{properties-lemma-characterize-integral}을 보라), 공집합이 아닌 모든
열린집합은 조밀하다. 따라서 $S$를 $S$의 아핀 열린집합으로 바꾸어
$S = \Spec(A)$가 아핀이라고 가정해도 된다. $S$가 정수적이므로
$A$는 정역이다. $f$가 유한형이므로 준콤팩트이고, 따라서 $X$는
준콤팩트이다. 그러므로 유한 아핀 열린 덮개
$X = \bigcup_{i = 1, \ldots, n} X_i$를 찾을 수 있다.
$X_i = \Spec(B_i)$라고 쓰자. 그러면 $B_i$는 유한형 $A$-대수이다.
보조정리 \ref{morphisms-lemma-locally-finite-type-characterize}를 보라.
더 나아가 유한형 $B_i$-가군 $M_i$들이 존재하여
$\mathcal{F}|_{X_i}$는 $B_i$-가군 $M_i$에 대응하는 준연접층이다.
성질, 보조정리 \ref{properties-lemma-finite-type-module}을 보라.
다음으로 각 지표쌍 $i, j$에 대해 아이디얼 $I_{ij} \subset B_i$를
$X_i \setminus X_i \cap X_j = V(I_{ij})$가 $X_i = \Spec(B_i)$ 안에서
성립하도록 택하자.
$M_{ij} = B_i/I_{ij}$로 두고 이를 $B_i$-가군으로 생각하자.
그러면 $V(I_{ij}) = \text{Supp}(M_{ij})$이고 $M_{ij}$는 유한
$B_i$-가군이다.

\medskip\noindent
이제 대수, 보조정리 \ref{algebra-lemma-generic-flatness}를 쌍
$(A \to B_i, M_{ij})$들과 쌍 $(A \to B_i, M_i)$들에 적용한다.
그러면 영이 아닌 $f_{ij}, f_i \in A$를 다음 성질을 갖도록 얻는다.
(a) $A_{f_{ij}} \to B_{i, f_{ij}}$는 평탄이고 유한 표시이며,
$M_{ij, f_{ij}}$는 $A_{f_{ij}}$ 위에서 평탄이고
$B_{i, f_{ij}}$ 위에서 유한 표시이다. 또한
(b) $B_{i, f_i}$는 $A_f$ 위에서 평탄이고 유한 표시이며,
$M_{i, f_i}$는 $B_{i, f_i}$ 위에서 평탄이고 유한 표시이다.
$f = (\prod f_i) (\prod f_{ij})$로 두자.
$U = D(f)$로 두면 된다고 주장한다.

\medskip\noindent
주장을 증명하기 위해 $A$를 $A_f$로 바꾸어도 된다. 즉
$U = \Spec(A_f) \to S$에 의해 기저변환해도 된다.
이 기저변환 뒤 각 $A \to B_i$는 평탄이고 유한 표시이며,
$M_i$, $M_{ij}$는 $A$ 위에서 평탄이고 $B_i$ 위에서 유한 표시이다.
이로부터 이미 $X \to S$가 준콤팩트이고 국소적으로 유한 표시이며
평탄이고, $\mathcal{F}$가 $S$ 위에서 평탄이고 $\mathcal{O}_X$ 위에서
유한 표시임을 알 수 있다. 보조정리
\ref{morphisms-lemma-locally-finite-presentation-characterize}와 성질, 보조정리
\ref{properties-lemma-finite-presentation-module}을 보라.
$M_{ij}$가 $B_i$ 위에서 유한 표시이므로
$X_i \cap X_j = X_i \setminus \text{Supp}(M_{ij})$는 $X_i$의
준콤팩트 열린집합이다. 대수, 보조정리
\ref{algebra-lemma-support-finite-presentation-constructible}를 보라.
따라서 스킴, 보조정리
\ref{schemes-lemma-characterize-quasi-separated}에 의해 $X \to S$는
준분리이다. 이로써 명제가 증명되었다.
\end{proof}

\noindent
사실 임의의 축약 기저 위에서도 일반점 평탄성의 판본이 성립한다.
그 판본은 다음과 같다.

\begin{proposition}[일반점 평탄성, 축약인 경우]
\label{morphisms-proposition-generic-flatness-reduced}
$f : X \to S$를 스킴의 사상이라 하자.
$\mathcal{F}$를 $\mathcal{O}_X$-가군의 준연접층이라 하자.
다음을 가정하자.
\begin{enumerate}
\item $S$는 축약이다.
\item $f$는 유한형이다.
\item $\mathcal{F}$는 유한형 $\mathcal{O}_X$-가군이다.
\end{enumerate}
그러면 조밀한 열린 부분스킴 $U \subset S$가 존재하여
$X_U \to U$는 평탄이고 유한 표시이며, $\mathcal{F}|_{X_U}$는
$U$ 위에서 평탄이고 $\mathcal{O}_{X_U}$ 위에서 유한 표시이다.
\end{proposition}

\begin{proof}
성급한 독자를 위해 말하면, 이 증명은 명제
\ref{morphisms-proposition-generic-flatness}의 증명을 되풀이하되 대수, 보조정리
\ref{algebra-lemma-generic-flatness-reduced}를 대수, 보조정리
\ref{algebra-lemma-generic-flatness} 대신 사용하는 것이다.

\medskip\noindent
평탄성과 유한 표시는 기저에 대해 국소적이므로, 보조정리
\ref{morphisms-lemma-flat-module-characterize}와
\ref{morphisms-lemma-locally-finite-presentation-characterize}를 보라,
$S$ 위에서 아핀 국소적으로 작업해도 된다. 따라서
$S = \Spec(A)$라고 가정해도 되며, 여기서 $A$는 축약환이다(성질, 보조정리
\ref{properties-lemma-characterize-reduced}를 보라).
$f$가 유한형이므로 준콤팩트이고, 따라서 $X$는 준콤팩트이다.
그러므로 유한 아핀 열린 덮개
$X = \bigcup_{i = 1, \ldots, n} X_i$를 찾을 수 있다.
$X_i = \Spec(B_i)$라고 쓰자. 그러면 $B_i$는 유한형 $A$-대수이다.
보조정리 \ref{morphisms-lemma-locally-finite-type-characterize}를 보라.
더 나아가 유한형 $B_i$-가군 $M_i$들이 존재하여
$\mathcal{F}|_{X_i}$는 $B_i$-가군 $M_i$에 대응하는 준연접층이다.
성질, 보조정리 \ref{properties-lemma-finite-type-module}을 보라.
다음으로 각 지표쌍 $i, j$에 대해 아이디얼 $I_{ij} \subset B_i$를
$X_i \setminus X_i \cap X_j = V(I_{ij})$가 $X_i = \Spec(B_i)$ 안에서
성립하도록 택하자. $M_{ij} = B_i/I_{ij}$로 두고 이를
$B_i$-가군으로 생각하자. 그러면
$V(I_{ij}) = \text{Supp}(M_{ij})$이고 $M_{ij}$는 유한 $B_i$-가군이다.

\medskip\noindent
이제 대수, 보조정리 \ref{algebra-lemma-generic-flatness-reduced}를 쌍
$(A \to B_i, M_{ij})$들과 쌍 $(A \to B_i, M_i)$들에 적용한다.
그러면 조밀한 열린집합 $U(A \to B_i, M_{ij}) \subset S$와
조밀한 열린집합 $U(A \to B_i, M_i) \subset S$를 얻는다. 기호는
대수, 식 (\ref{algebra-equation-good-locus})에서와 같다.
조밀한 열린집합들의 유한 교집합은 조밀한 열린집합이므로
$$
U =
\bigcap\nolimits_{i, j} U(A \to B_i, M_{ij})
\quad\cap\quad
\bigcap\nolimits_i U(A \to B_i, M_i)
$$
는 $S$에서 열리고 조밀하다. $U$가 원하는 열린집합이라고 주장한다.

\medskip\noindent
$u \in U$를 택하자. 자취 $U(A \to B_i, M_{ij})$와
$U(A \to B, M_i)$의 정의에 의해 $f_{ij}, f_i \in A$가 존재하여
(a) $u \in D(f_i)$ 및 $u \in D(f_{ij})$이고,
(b) $A_{f_{ij}} \to B_{i, f_{ij}}$는 평탄이고 유한 표시이며
$M_{ij, f_{ij}}$는 $A_{f_{ij}}$ 위에서 평탄이고
$B_{i, f_{ij}}$ 위에서 유한 표시이고,
(c) $B_{i, f_i}$는 $A_{f_i}$ 위에서 평탄이고 유한 표시이며
$M_{i, f_i}$는 $B_{i, f_i}$ 위에서 유한 표시이고 $A_{f_i}$ 위에서
평탄하다. 다음과 같이 두자.
$f = (\prod f_i) (\prod f_{ij})$.
이제 $X \to S$가 $D(f)$ 위에서 평탄이고 유한 표시이며,
$\mathcal{F}$를 $X_{D(f)}$에 제한한 것이 $D(f)$ 위에서 평탄이고
$X_{D(f)}$의 구조층 위에서 유한 표시임을 보이면 충분하다.

\medskip\noindent
따라서 $A$를 $A_f$로 바꾸어도 된다. 즉
$\Spec(A_f) \to S$에 의해 기저변환해도 된다.
이 기저변환 뒤 각 $A \to B_i$는 평탄이고 유한 표시이며,
$M_i$, $M_{ij}$는 $A$ 위에서 평탄이고 $B_i$ 위에서 유한 표시이다.
이로부터 이미 $X \to S$가 준콤팩트이고 국소적으로 유한 표시이며
평탄이고, $\mathcal{F}$가 $S$ 위에서 평탄이고 $\mathcal{O}_X$ 위에서
유한 표시임을 알 수 있다. 보조정리
\ref{morphisms-lemma-locally-finite-presentation-characterize}와 성질, 보조정리
\ref{properties-lemma-finite-presentation-module}을 보라.
$M_{ij}$가 $B_i$ 위에서 유한 표시이므로
$X_i \cap X_j = X_i \setminus \text{Supp}(M_{ij})$는 $X_i$의
준콤팩트 열린집합이다. 대수, 보조정리
\ref{algebra-lemma-support-finite-presentation-constructible}를 보라.
따라서 스킴, 보조정리
\ref{schemes-lemma-characterize-quasi-separated}에 의해 $X \to S$는
준분리이다. 이로써 명제가 증명되었다.
\end{proof}










\begin{remark}
\label{morphisms-remark-flattening}
위의 결과들은 스킴 사상에 대한 더 정교한 평탄화 기법을 향한 첫걸음이다.
레이노와 그뤼종의 논문 \cite{GruRay}에는 이 방향의 훌륭한 결과들이 많다.
\end{remark}







\section{사상과 올의 차원}
\label{morphisms-section-dimension-fibres}

\noindent
$X$를 위상공간이라 하고 $x \in X$라 하자.
$\dim_x(X)$를 $x$의 $X$ 안의 열린 근방들의 차원 중 최솟값으로
정의했음을 상기하자. 위상수학, 정의
\ref{topology-definition-Krull}을 보라.

\begin{lemma}
\label{morphisms-lemma-dimension-fibre-at-a-point}
$f : X \to S$를 스킴의 사상이라 하자.
$x \in X$라 하고 $s = f(x)$로 두자.
$f$가 국소적으로 유한형이라고 가정하자. 그러면
$$
\dim_x(X_s) =
\dim(\mathcal{O}_{X_s, x}) + \text{trdeg}_{\kappa(s)}(\kappa(x)).
$$
\end{lemma}

\begin{proof}
이는 곧바로 $S = s$이고 $X$가 아핀인 경우로 환원된다.
이 경우 결과는 대수, 보조정리
\ref{algebra-lemma-dimension-at-a-point-finite-type-field}에서 따른다.
\end{proof}

\begin{lemma}
\label{morphisms-lemma-dimension-fibre-at-a-point-additive}
$f : X \to Y$와 $g : Y \to S$를 스킴의 사상이라 하자.
$x \in X$라 하고 $y = f(x)$, $s = g(y)$로 두자.
$f$와 $g$가 국소적으로 유한형이라고 가정하자. 그러면
$$
\dim_x(X_s) \leq \dim_x(X_y) + \dim_y(Y_s).
$$
더 나아가 $\mathcal{O}_{X_s, x}$가 $\mathcal{O}_{Y_s, y}$ 위에서
평탄이면 등호가 성립한다. 예를 들어 $\mathcal{O}_{X, x}$가
$\mathcal{O}_{Y, y}$ 위에서 평탄이면 이 조건이 성립한다.
\end{lemma}

\begin{proof}
다음에 유의하자.
$\text{trdeg}_{\kappa(s)}(\kappa(x)) =
\text{trdeg}_{\kappa(y)}(\kappa(x)) + \text{trdeg}_{\kappa(s)}(\kappa(y))$.
따라서 보조정리 \ref{morphisms-lemma-dimension-fibre-at-a-point}에 의해 명제는
다음과 동치이다.
$$
\dim(\mathcal{O}_{X_s, x})
\leq
\dim(\mathcal{O}_{X_y, x}) + \dim(\mathcal{O}_{Y_s, y}).
$$
이에 대해서는 대수, 보조정리
\ref{algebra-lemma-dimension-base-fibre-total}을 보라. 평탄인 경우에는
대수, 보조정리 \ref{algebra-lemma-dimension-base-fibre-equals-total}을 보라.
\end{proof}

\begin{lemma}
\label{morphisms-lemma-dimension-fibre-after-base-change}
스킴의 올곱 도식
$$
\xymatrix{
X' \ar[r]_{g'} \ar[d]_{f'} & X \ar[d]^f \\
S' \ar[r]^g & S
}
$$
이 주어졌다고 하자. $f$가 국소적으로 유한형이라고 가정하자.
$x' \in X'$, $x = g'(x')$, $s' = f'(x')$, 그리고
$s = g(s') = f(x)$라고 하자. 그러면
\begin{enumerate}
\item $\dim_x(X_s) = \dim_{x'}(X'_{s'})$이다.
\item $F$를 사상 $X'_{s'} \to X_s$의 $x$ 위의 올이라 하면
$$
\dim(\mathcal{O}_{F, x'}) =
\dim(\mathcal{O}_{X'_{s'}, x'}) - \dim(\mathcal{O}_{X_s, x}) =
\text{trdeg}_{\kappa(s)}(\kappa(x)) -
\text{trdeg}_{\kappa(s')}(\kappa(x'))
$$
이다. 특히 $\dim(\mathcal{O}_{X'_{s'}, x'}) \geq \dim(\mathcal{O}_{X_s, x})$이고
$\text{trdeg}_{\kappa(s)}(\kappa(x)) \geq
\text{trdeg}_{\kappa(s')}(\kappa(x'))$이다.
\item $s', s, x$가 주어지면 점 $x'$를 택하여
$\dim(\mathcal{O}_{X'_{s'}, x'}) = \dim(\mathcal{O}_{X_s, x})$ 및
$\text{trdeg}_{\kappa(s)}(\kappa(x)) = \text{trdeg}_{\kappa(s')}(\kappa(x'))$
를 만족하게 할 수 있다.
\end{enumerate}
\end{lemma}

\begin{proof}
(1)은 대수, 보조정리
\ref{algebra-lemma-dimension-at-a-point-preserved-field-extension}에서
곧바로 따른다. (2)와 (3)은 대수, 보조정리
\ref{algebra-lemma-inequalities-under-field-extension}에서 따른다.
\end{proof}

\noindent
다음 보조정리는 자명하지 않은 대수적 결과, 즉 자리스키 주정리의
대수적 판본에서 따른다.

\begin{lemma}
\label{morphisms-lemma-openness-bounded-dimension-fibres}
\begin{reference}
\cite[IV Theorem 13.1.3]{EGA}
\end{reference}
$f : X \to S$를 스킴의 사상이라 하자.
$n \geq 0$이라 하자. $f$가 국소적으로 유한형이라고 가정하자.
집합
$$
U_n = \{x \in X \mid \dim_x X_{f(x)} \leq n\}
$$
은 $X$에서 열린집합이다.
\end{lemma}

\begin{proof}
이는 대수, 보조정리
\ref{algebra-lemma-dimension-fibres-bounded-open-upstairs}에서 곧바로 따른다.
\end{proof}

\begin{lemma}
\label{morphisms-lemma-morphism-finite-type-bounded-dimension}
$f : X \to Y$를 유한형 사상이라 하고 $Y$가 준콤팩트라고 하자.
그러면 $f$의 올들의 차원은 유계이다.
\end{lemma}

\begin{proof}
보조정리 \ref{morphisms-lemma-openness-bounded-dimension-fibres}에 의해 점들의 집합
$U_n \subset X$는 열려 있다. 이 집합에서는 올의 차원이 $\leq n$이다.
$f$가 유한형이므로 모든 점은 어떤 $U_n$에 포함된다
(체 위의 유한형 대수의 차원은 유한하기 때문이다).
$Y$가 준콤팩트이고 $f$가 유한형이므로 $X$가 준콤팩트임을 알 수 있다.
따라서 $X = U_n$인 어떤 $n$이 존재한다.
\end{proof}

\begin{lemma}
\label{morphisms-lemma-openness-bounded-dimension-fibres-finite-presentation}
$f : X \to S$를 스킴의 사상이라 하자.
$n \geq 0$이라 하자. $f$가 국소적으로 유한 표시라고 가정하자.
보조정리 \ref{morphisms-lemma-openness-bounded-dimension-fibres}의 열린집합
$$
U_n = \{x \in X \mid \dim_x X_{f(x)} \leq n\}
$$
은 $X$에서 역콤팩트이다(위상수학, 정의
\ref{topology-definition-quasi-compact}를 보라).
\end{lemma}

\begin{proof}
위상공간 $X$는 아핀 열린집합 $U \subset X$들로 이루어진 위상 기저를
갖고, 유도된 사상 $f|_U : U \to S$는 어떤 아핀 열린집합
$V \subset S$를 통해 분해된다. 따라서 $U \cap U_n$이
그러한 $U$에 대해 준콤팩트임을 보이면 충분하다. $U_n \cap U$는 열린집합
$\{x \in U \mid \dim_x U_{f(x)} \leq n\}$와 같다.
그러므로 $X$와 $S$가 아핀인 경우로 환원된다. 이 경우 보조정리는
대수, 보조정리
\ref{algebra-lemma-dimension-fibres-bounded-quasi-compact-open-upstairs}와
보조정리 \ref{morphisms-lemma-locally-finite-presentation-characterize}에서 따른다.
\end{proof}

\begin{lemma}
\label{morphisms-lemma-dimension-fibre-specialization}
$f : X \to S$를 스킴의 사상이라 하자.
$x \leadsto x'$를 $X$에서 같은 점 $s \in S$ 위에 놓이는 점들의
자명하지 않은 특수화라 하자. $f$가 국소적으로 유한형이라고 가정하자.
그러면
\begin{enumerate}
\item $\dim_x(X_s) \leq \dim_{x'}(X_s)$,
\item $\dim(\mathcal{O}_{X_s, x}) < \dim(\mathcal{O}_{X_s, x'})$, 그리고
\item $\text{trdeg}_{\kappa(s)}(\kappa(x)) >
\text{trdeg}_{\kappa(s)}(\kappa(x'))$이다.
\end{enumerate}
\end{lemma}

\begin{proof}
(1)은 $X_s$의 모든 열린집합이 $x'$를 포함하면 $x$도 포함한다는
사실에서 따른다. (2)는 $\mathcal{O}_{X_s, x}$가
$\mathcal{O}_{X_s, x'}$를 한 소 아이디얼에서 국소화한 것이므로,
$\mathcal{O}_{X_s, x}$의 모든 소 아이디얼 사슬이
$\mathcal{O}_{X_s, x'}$의 진정으로 더 긴 소 아이디얼 사슬의 일부라는
사실에서 따른다. 마지막 부등식은 대수, 보조정리
\ref{algebra-lemma-tr-deg-specialization}에서 따른다.
\end{proof}






\section{주어진 상대차원의 사상}
\label{morphisms-section-relative-dimension}

\noindent
주어진 상대차원의 사상을 편리하게 말하기 위해 다음 개념을 도입한다.

\begin{definition}
\label{morphisms-definition-relative-dimension-d}
$f : X \to S$를 스킴의 사상이라 하자.
$f$가 국소적으로 유한형이라고 가정하자.
\begin{enumerate}
\item $f$가 {\it 상대차원 $\leq d$인 점 $x$}를 갖는다고 함은
$\dim_x(X_{f(x)}) \leq d$인 경우이다.
\item $f$가 {\it 상대차원 $\leq d$}라고 함은
$\dim_x(X_{f(x)}) \leq d$가 모든 $x \in X$에 대해 성립하는 경우이다.
\item $f$가 {\it 상대차원 $d$}라고 함은 공집합이 아닌 모든 올
$X_s$가 $d$차원 순수차원인 경우이다.
\end{enumerate}
\end{definition}

\noindent
이 개념은 특별히 좋은 거동을 보이지는 않지만 여러 상황에서 잘 작동한다.

\begin{lemma}
\label{morphisms-lemma-base-change-relative-dimension-d}
$f : X \to S$를 국소적으로 유한형인 스킴의 사상이라 하자.
$f$가 상대차원 $d$이면 $f$의 모든 기저변환도 같은 상대차원을 갖는다.
상대차원 $\leq d$인 경우도 마찬가지이다.
\end{lemma}

\begin{proof}
보조정리 \ref{morphisms-lemma-dimension-fibre-after-base-change}에서 곧바로 따른다.
\end{proof}

\begin{lemma}
\label{morphisms-lemma-composition-relative-dimension-d}
$f : X \to Y$, $g : Y \to Z$가 국소적으로 유한형이라고 하자.
$f$가 상대차원 $\leq d$이고 $g$가 상대차원 $\leq e$이면
$g \circ f$는 상대차원 $\leq d + e$이다. 다음을 가정하자.
\begin{enumerate}
\item $f$는 상대차원 $d$이다.
\item $g$는 상대차원 $e$이다.
\item $f$는 평탄이다.
\end{enumerate}
그러면 $g \circ f$는 상대차원 $d + e$이다.
\end{lemma}

\begin{proof}
보조정리 \ref{morphisms-lemma-dimension-fibre-at-a-point-additive}에서 곧바로 따른다.
\end{proof}

\noindent
일반적으로 사상을 상대차원이 각각 주어진 값인 조각들로 분해할 수는 없다.
그러나 사상이 코언--매콜리 올을 갖고 유한 표시인 평탄 사상이면 가능하다.

\begin{lemma}
\label{morphisms-lemma-flat-finite-presentation-CM-fibres-relative-dimension}
\begin{slogan}
코언--매콜리 사상은 순수 상대차원을 갖는 열린닫힌 조각들로 분해된다.
\end{slogan}
$f : X \to S$를 스킴의 사상이라 하자. 다음을 가정하자.
\begin{enumerate}
\item $f$는 평탄이다.
\item $f$는 국소적으로 유한 표시이다.
\item 모든 $s \in S$에 대해 올 $X_s$는 코언--매콜리이다
(성질, 정의 \ref{properties-definition-Cohen-Macaulay}).
\end{enumerate}
그러면 열린닫힌 부분스킴 $X_d \subset X$들이 존재하여
$X = \coprod_{d \geq 0} X_d$이고 $f|_{X_d} : X_d \to S$는
상대차원 $d$이다.
\end{lemma}

\begin{proof}
대수, 보조정리 \ref{algebra-lemma-relative-dimension-CM}에서 곧바로 따른다.
\end{proof}

\begin{lemma}
\label{morphisms-lemma-locally-quasi-finite-rel-dimension-0}
$f : X \to S$를 스킴의 사상이라 하자.
$f$가 국소적으로 유한형이라고 가정하자.
$x \in X$를 택하고 $s = f(x)$로 두자.
$f$가 $x$에서 준유한인 것은 $\dim_x(X_s) = 0$인 것과 동치이다.
특히 $f$가 국소적으로 준유한인 것은 $f$가 상대차원 $0$인 것과 동치이다.
\end{lemma}

\begin{proof}
첫 번째 증명.
$f$가 $x$에서 준유한이면 $\kappa(x)$는 $\kappa(s)$의 유한 확장이고
(보조정리 \ref{morphisms-lemma-residue-field-quasi-finite}), $x$는 $X_s$의 고립점이다
(보조정리 \ref{morphisms-lemma-quasi-finite-at-point-characterize}). 따라서 보조정리
\ref{morphisms-lemma-dimension-fibre-at-a-point}에 의해 $\dim_x(X_s) = 0$이다.
역으로 $\dim_x(X_s) = 0$이면 보조정리
\ref{morphisms-lemma-dimension-fibre-at-a-point}에 의해
$\kappa(s) \subset \kappa(x)$는 대수적이고 $X_s$에서 $x$로 특수화되는
다른 점은 없다. 따라서 보조정리
\ref{morphisms-lemma-algebraic-residue-field-extension-closed-point-fibre}에 의해
$x$는 자기 올에서 닫혀 있고, 보조정리
\ref{morphisms-lemma-quasi-finite-at-point-characterize}의 (3)에 의해
$f$는 $x$에서 준유한이다.

\medskip\noindent
두 번째 증명. 올 $X_s$는 체 위에서 국소적으로 유한형인 스킴이므로
국소 뇌터이다(보조정리 \ref{morphisms-lemma-finite-type-noetherian}). 이제 결과는
보조정리 \ref{morphisms-lemma-quasi-finite-at-point-characterize}와 성질, 보조정리
\ref{properties-lemma-isolated-dim-0-locally-noetherian}에서 따른다.
\end{proof}

\begin{lemma}
\label{morphisms-lemma-rel-dimension-dimension}
$f : X \to Y$를 평탄이고 국소적으로 유한형이며 상대차원 $d$인
국소 뇌터 스킴의 사상이라 하자. 모든 점 $x$가 $X$에 있고 그 상
$y$가 $Y$에 있으면 $\dim_x(X) = \dim_y(Y) + d$이다.
\end{lemma}

\begin{proof}
$X$와 $Y$를 $x$와 $y$의 열린 근방으로 줄인 뒤, 점에서의 스킴 차원
정의에 의해 $\dim(X) = \dim_x(X)$ 및 $\dim(Y) = \dim_y(Y)$라고
가정할 수 있다(성질, 정의 \ref{properties-definition-dimension}).
보조정리 \ref{morphisms-lemma-noetherian-finite-type-finite-presentation}과
\ref{morphisms-lemma-fppf-open}에 의해 사상 $f$는 열려 있다.
따라서 $Y$를 줄여 $f$가 전사가 되게 할 수 있다.
$\dim(X) = \dim(Y) + d$임을 보이는 일만 남았다.

\medskip\noindent
$a$를 $X$의 점이라 하고 그 상 $b$가 $Y$에 있다고 하자.
대수, 보조정리 \ref{algebra-lemma-dimension-base-fibre-equals-total}에 의해
$$
\dim(\mathcal{O}_{X,a}) = \dim(\mathcal{O}_{Y,b}) + \dim(\mathcal{O}_{X_b, a}).
$$
$a$가 $X$의 모든 점을 달릴 때 상한을 취하면 원하는 대로
$\dim(X) = \dim(Y) + d$를 얻는다. 성질, 보조정리
\ref{properties-lemma-dimension}을 보라.
\end{proof}











\section{신토믹 사상}
\label{morphisms-section-syntomic}

\noindent
$A$를 체 $k$ 위의 대수라 하자. 이 대수가
{\it $k$ 위의 대역 완전교차}라고 함은
$A \cong k[x_1, \ldots, x_n]/(f_1, \ldots, f_c)$이고
$\dim(A) = n - c$인 경우이다. 대수 $A$가 체 $k$ 위에서
{\it 국소 완전교차}라고 함은 $\Spec(A)$가 각각 $k$ 위의 대역
완전교차인 표준 열린집합들로 덮이는 경우이다. 대수, 절
\ref{algebra-section-lci}을 보라. 환 준동형 $R \to A$가 유한 표시이고
평탄이며 그 올들이 국소 완전교차환이면 이를 {\it 신토믹}이라 한다는
것을 상기하자. 대수, 정의 \ref{algebra-definition-lci}을 보라.

\begin{definition}
\label{morphisms-definition-syntomic}
$f : X \to S$를 스킴의 사상이라 하자.
\begin{enumerate}
\item $f$가 {\it $x \in X$에서 신토믹}이라고 함은 아핀 열린 근방
$\Spec(A) = U \subset X$로서 $x$를 포함하는 것과 아핀 열린집합
$\Spec(R) = V \subset S$가 존재하여 $f(U) \subset V$이고 유도된 환 준동형
$R \to A$가 신토믹인 경우이다.
\item $f$가 $X$의 모든 점에서 신토믹이면 이를 {\it 신토믹}이라 한다.
\item $S = \Spec(k)$이고 $f$가 신토믹이면 $X$를
{\it $k$ 위의 국소 완전교차}라 한다.
\item 아핀 스킴의 사상 $f : X \to S$에 대하여 대역 상대 완전교차
$R \to R[x_1, \ldots, x_n]/(f_1, \ldots, f_c)$가 존재하여(대수, 정의
\ref{algebra-definition-relative-global-complete-intersection}을 보라)
$X \to S$가
$$
\Spec(R[x_1, \ldots, x_n]/(f_1, \ldots, f_c)) \to \Spec(R).
$$
와 동형이면 이를 {\it 표준 신토믹}이라 한다.
\end{enumerate}
\end{definition}

\noindent
문헌에서는 신토믹 사상을 때때로
{\it 평탄 국소 완전교차 사상}이라고 부른다.
이는 편리한 사상류임이 드러난다. 예를 들어 이들로 신토믹 위상을
정의할 수 있는데, 이는 매끄러운 위상과 \'etale 위상보다 세밀하지만
그들과 같은 형식적 성질을 많이 갖는다.

\medskip\noindent
대역 상대 완전교차(표준 신토믹 환 준동형을 정의하는 데 사용했다)는
특히 평탄이다. 사상의 추가 내용, 절 \ref{more-morphisms-section-lci}에서는
국소적으로 사상 $X \to S$가
$$
\Spec(R[x_1, \ldots, x_n]/(f_1, \ldots, f_c)) \to \Spec(R).
$$
꼴인 경우를 고찰할 것이다. 여기서
$f_1, \ldots, f_r$는 $R[x_1, \ldots, x_n]$의 코쥘-정칙열이다.
이러한 사상을 {\it 국소 완전교차 사상}이라 부를 것이다.
이 정의를 마련하고 나면, 사상이 신토믹인 것은 평탄한 국소
완전교차 사상인 것과 동치가 된다.

\medskip\noindent
신토믹 사상의 정의에는 분리성이나 준콤팩트성 가정이 없다는 데
유의하자. 따라서 신토믹 여부는 원천에서 국소적인 문제이다.
정확한 결과는 다음과 같다.

\begin{lemma}
\label{morphisms-lemma-syntomic-characterize}
$f : X \to S$를 스킴의 사상이라 하자. 다음은 동치이다.
\begin{enumerate}
\item 사상 $f$는 신토믹이다.
\item 모든 아핀 열린집합 $U \subset X$, $V \subset S$에 대하여
$f(U) \subset V$이면 환 준동형
$\mathcal{O}_S(V) \to \mathcal{O}_X(U)$가 신토믹이다.
\item 열린 덮개 $S = \bigcup_{j \in J} V_j$와 열린 덮개
$f^{-1}(V_j) = \bigcup_{i \in I_j} U_i$가 존재하여 각 사상
$U_i \to V_j$, $j\in J, i\in I_j$가 신토믹이다.
\item 아핀 열린 덮개 $S = \bigcup_{j \in J} V_j$와 아핀 열린 덮개
$f^{-1}(V_j) = \bigcup_{i \in I_j} U_i$가 존재하여 모든
환 준동형 $\mathcal{O}_S(V_j) \to \mathcal{O}_X(U_i)$가
$j\in J, i\in I_j$에 대해 신토믹이다.
\end{enumerate}
또한 $f$가 신토믹이면 임의의 열린 부분스킴 $U \subset X$, $V \subset S$에
대해 $f(U) \subset V$일 때 제한
$f|_U : U \to V$도 신토믹이다.
\end{lemma}

\begin{proof}
``$R \to A$가 신토믹이다''라는 성질이 국소적임을 보이면 이 결과는
보조정리 \ref{morphisms-lemma-locally-P}에서 따른다. 정의
\ref{morphisms-definition-property-local}의 조건 (a), (b), (c)를 확인하자.
대수, 보조정리 \ref{algebra-lemma-base-change-syntomic}에 의해 신토믹성은
기저변환으로 보존되므로 (a)가 성립한다. 대수, 보조정리
\ref{algebra-lemma-composition-syntomic}에 의해 신토믹성은 합성으로
보존되고, 임의의 환 $R$에 대해 환 준동형 $R \to R_f$가 신토믹임은
자명하다. 따라서 (b)가 성립한다. 마지막으로 대수, 보조정리
\ref{algebra-lemma-local-syntomic}에 의해 성질 (c)가 참이다.
\end{proof}

\begin{lemma}
\label{morphisms-lemma-composition-syntomic}
신토믹인 두 사상의 합성은 신토믹이다.
\end{lemma}

\begin{proof}
보조정리 \ref{morphisms-lemma-syntomic-characterize}의 증명에서 신토믹성이 환
준동형의 국소적 성질임을 보았다. 따라서 보조정리
\ref{morphisms-lemma-composition-type-P}와 신토믹성이 합성으로 보존되는 환
준동형의 성질이라는 사실을 결합하면 보조정리의 첫 번째 명제를 얻는다.
대수, 보조정리 \ref{algebra-lemma-composition-syntomic}을 보라.
\end{proof}

\begin{lemma}
\label{morphisms-lemma-base-change-syntomic}
신토믹 사상의 기저변환은 신토믹이다.
\end{lemma}

\begin{proof}
보조정리 \ref{morphisms-lemma-syntomic-characterize}의 증명에서 신토믹성이 환
준동형의 국소적 성질임을 보았다. 따라서 보조정리
\ref{morphisms-lemma-composition-type-P}와 신토믹성이 기저변환으로 보존되는 환
준동형의 성질이라는 사실을 결합하면 결과가 따른다. 대수, 보조정리
\ref{algebra-lemma-base-change-syntomic}을 보라.
\end{proof}

\begin{lemma}
\label{morphisms-lemma-open-immersion-syntomic}
모든 열린 몰입은 신토믹이다.
\end{lemma}

\begin{proof}
열린 몰입은 국소 동형사상이므로 참이다.
\end{proof}

\begin{lemma}
\label{morphisms-lemma-syntomic-locally-finite-presentation}
신토믹 사상은 국소적으로 유한 표시이다.
\end{lemma}

\begin{proof}
신토믹 환 준동형은 정의상 유한 표시이므로 참이다.
\end{proof}

\begin{lemma}
\label{morphisms-lemma-syntomic-flat}
신토믹 사상은 평탄이다.
\end{lemma}

\begin{proof}
신토믹 환 준동형은 정의상 평탄이므로 참이다.
\end{proof}

\begin{lemma}
\label{morphisms-lemma-syntomic-open}
신토믹 사상은 보편 열림이다.
\end{lemma}

\begin{proof}
보조정리 \ref{morphisms-lemma-syntomic-locally-finite-presentation},
\ref{morphisms-lemma-syntomic-flat}, \ref{morphisms-lemma-fppf-open}을 결합하라.
\end{proof}

\noindent
$k$를 체라 하자. $A$를 $k$ 위에서 본질적으로 유한형인 국소
$k$-대수라 하자. $A$가 {\it $k$ 위의 완전교차}라고 함은
$A \cong R/(f_1, \ldots, f_c)$로 쓸 수 있고, 여기서 $R$은 $k$ 위에서
본질적으로 유한형인 정칙 국소환이며 $f_1, \ldots, f_c$는 $R$의
정칙열인 경우임을 상기하자. 대수, 정의
\ref{algebra-definition-lci-local-ring}을 보라.

\begin{lemma}
\label{morphisms-lemma-local-complete-intersection}
$k$를 체라 하자. $X$를 $k$ 위에서 국소적으로 유한형인 스킴이라 하자.
다음은 동치이다.
\begin{enumerate}
\item $X$는 $k$ 위의 국소 완전교차이다.
\item 모든 $x \in X$에 대해 아핀 열린집합
$U = \Spec(R) \subset X$가 존재하여 $x$를 포함하고,
$R \cong k[x_1, \ldots, x_n]/(f_1, \ldots, f_c)$는 $k$ 위의 대역
완전교차이다.
\item 모든 $x \in X$에 대해 국소환 $\mathcal{O}_{X, x}$는
$k$ 위의 완전교차이다.
\end{enumerate}
\end{lemma}

\begin{proof}
이에 대응하는 대수적 결과는 대수, 보조정리
\ref{algebra-lemma-lci-at-prime}과 \ref{algebra-lemma-lci-global}에 있다.
\end{proof}

\noindent
다음 보조정리는 국소적으로 모든 신토믹 사상이 표준 신토믹이라고 말한다.
따라서 표준 신토믹 사상을 신토믹 사상의 {\it 국소 모형}으로 사용할 수
있다. 또한 유한 표시인 평탄 사상이 신토믹인 것은 그 올들이 국소
완전교차 스킴인 것과 동치임을 말한다.

\begin{lemma}
\label{morphisms-lemma-syntomic-locally-standard-syntomic}
$f : X  \to S$를 스킴의 사상이라 하자. $x \in X$를 상이
$s = f(x)$인 점이라 하자. $V \subset S$를 $s$의 아핀 열린 근방이라 하자.
다음은 동치이다.
\begin{enumerate}
\item 사상 $f$는 $x$에서 신토믹이다.
\item $U \subset X$인 아핀 열린집합으로서 $x \in U$이고
$f(U) \subset V$이며 $f|_U : U \to V$가 표준 신토믹인 것이 존재한다.
\item 사상 $f$는 $x$에서 유한 표시이고, 국소환 준동형
$\mathcal{O}_{S, s} \to \mathcal{O}_{X, x}$는 평탄이며
$\mathcal{O}_{X, x}/\mathfrak m_s \mathcal{O}_{X, x}$는
$\kappa(s)$ 위의 완전교차이다(대수, 정의
\ref{algebra-definition-lci-local-ring}을 보라).
\end{enumerate}
\end{lemma}

\begin{proof}
정의들과 대수, 보조정리 \ref{algebra-lemma-syntomic}에서 따른다.
\end{proof}

\begin{lemma}
\label{morphisms-lemma-syntomic-flat-fibres}
$f : X \to S$를 스킴의 사상이라 하자. $f$가 평탄이고 국소적으로
유한 표시이며 모든 올 $X_s$가 국소 완전교차이면 $f$는 신토믹이다.
\end{lemma}

\begin{proof}
보조정리 \ref{morphisms-lemma-local-complete-intersection},
\ref{morphisms-lemma-syntomic-locally-standard-syntomic}과 국소환의 동형사상
$
\mathcal{O}_{X, x}/\mathfrak m_s \mathcal{O}_{X, x}
\cong
\mathcal{O}_{X_s, x}
$에서 명백하다.
\end{proof}

\begin{lemma}
\label{morphisms-lemma-set-points-where-fibres-lci}
$f : X \to S$를 스킴의 사상이라 하자. $f$가 국소적으로 유한형이라고
가정하자. 집합
$$
T = \{x \in X \mid \mathcal{O}_{X_{f(x)}, x}
\text{ is a complete intersection over }\kappa(f(x))\}
$$
의 구성은 임의의 기저변환과 가환한다. 즉 임의의 사상
$g : S' \to S$에 대해 $f' : X' \to S'$를 $f$의 기저변환이라 하고 사영
$g' : X' \to X$를 고찰하자. 그러면 대응하는 집합 $T'$는 사상 $f'$에 대해
$T' = (g')^{-1}(T)$와 같다. 특히 $f$가 평탄이고 국소적으로 유한
표시라고 가정하면, $f$가 신토믹인 점들의 열린집합에도 같은 결론이
성립한다.
\end{lemma}

\begin{proof}
$s' \in S'$를 점이라 하고 $s = g(s')$로 두자. 그러면
$$
X'_{s'} =
\Spec(\kappa(s')) \times_{\Spec(\kappa(s))} X_s
$$
이다. 다시 말해 기저변환의 올은 원래 올의 기저변환이다. 따라서 첫째
부분은 대수, 보조정리 \ref{algebra-lemma-lci-field-change-local}과
동치이다. 이 경우 $T$가 보조정리
\ref{morphisms-lemma-syntomic-locally-standard-syntomic}에 따라 $f$가 신토믹인
점들의 집합이므로 둘째 부분은 첫째 부분에서 따른다.
\end{proof}

\begin{lemma}
\label{morphisms-lemma-standard-syntomic-relative-dimension}
$R$을 환이라 하자. $R \to A = R[x_1, \ldots, x_n]/(f_1, \ldots, f_c)$를
상대 대역 완전교차라 하자. $S = \Spec(R)$ 및 $X = \Spec(A)$로 두자.
사상 $f : X \to S$를 고찰하자. 이는 환 준동형 $R \to A$에 대응한다.
함수 $x \mapsto \dim_x(X_{f(x)})$는 상수이고 그 값은 $n - c$이다.
\end{lemma}

\begin{proof}
대수, 정의 \ref{algebra-definition-relative-global-complete-intersection}에
따르면 $R \to A$가 상대 대역 완전교차라는 것은 영이 아닌 모든 올환의
차원이 $n - c$라는 뜻이다. 따라서 소 아이디얼 $\mathfrak p$를 $R$에서 택하면
대한 올환
$\kappa(\mathfrak p)[x_1, \ldots, x_n]/(\overline{f}_1, \ldots, \overline{f}_c)$는
영이거나 차원 $n - c$인 대역 완전교차환이다. 대수, 정의
\ref{algebra-definition-lci-field} 뒤의 논의에 의해 이는 차원
$n - c$인 순수차원환이다. 이로부터 보조정리가 따른다.
\end{proof}

\begin{lemma}
\label{morphisms-lemma-syntomic-relative-dimension}
$f : X \to S$를 신토믹 사상이라 하자. 함수
$x \mapsto \dim_x(X_{f(x)})$는 $X$ 위에서 국소상수이다.
\end{lemma}

\begin{proof}
보조정리 \ref{morphisms-lemma-syntomic-locally-standard-syntomic}에 의해 사상 $f$는
국소적으로 아핀 스킴 사이의 표준 신토믹 사상처럼 보인다. 따라서 결과는
보조정리 \ref{morphisms-lemma-standard-syntomic-relative-dimension}에서 따른다.
\end{proof}

\noindent
보조정리 \ref{morphisms-lemma-syntomic-relative-dimension}에 의해 다음 정의가
의미를 갖는다.

\begin{definition}
\label{morphisms-definition-syntomic-relative-dimension}
$d \geq 0$을 정수라 하자. 스킴의 사상 $f : X \to S$가
{\it 상대차원 $d$인 신토믹}이라고 함은 $f$가 신토믹이고 함수
$\dim_x(X_{f(x)}) = d$가 모든 $x \in X$에 대해 성립하는 경우이다.
\end{definition}

\noindent
다시 말해 $f$는 신토믹이고 공집합 아닌 올들은 차원 $d$인 순수차원이다.

\begin{lemma}
\label{morphisms-lemma-syntomic-permanence}
스킴의 사상들의 가환 그림
$$
\xymatrix{
X \ar[rr]_f \ar[rd]_p & &
Y \ar[dl]^q \\
& S
}
$$
가 주어졌다고 하자. 다음을 가정하자.
\begin{enumerate}
\item $f$는 전사이고 신토믹이다.
\item $p$는 신토믹이다.
\item $q$는 국소적으로 유한 표시이다\footnote{사실 $f$가 전사이고
평탄이며 국소적으로 유한 표시이고 $p$가 신토믹이면 $q$와 $f$가 모두
신토믹이다. 하강, 보조정리 \ref{descent-lemma-syntomic-permanence}을 보라.}.
\end{enumerate}
그러면 $q$는 신토믹이다.
\end{lemma}

\begin{proof}
보조정리 \ref{morphisms-lemma-flat-permanence}에 의해 $q$는 평탄이다. 따라서
보조정리 \ref{morphisms-lemma-syntomic-flat-fibres}에 의해 $Y \to S$의 올들이
국소 완전교차임을 보이면 충분하다. $s \in S$라 하자. 사상
$X_s \to Y_s$를 고찰하자. 이는 사상 $X \to Y$의 기저변환이므로 전사이고
신토믹이다(보조정리 \ref{morphisms-lemma-base-change-syntomic}). 같은 이유로
$X_s$는 $\kappa(s)$ 위에서 신토믹이다. 또한 $Y_s$는 $\kappa(s)$ 위에서
국소적으로 유한형이다(보조정리 \ref{morphisms-lemma-base-change-finite-type}).
이렇게 하여 $S$가 체의 스펙트럼인 경우로 환원된다.

\medskip\noindent
$S = \Spec(k)$라고 가정하자. $y \in Y$라 하자. 아핀 열린집합
$\Spec(A) \subset Y$로서 $y$를 포함하는 것을 택하자. 아핀 열린집합
$\Spec(B) \subset X$로서 $f(\Spec(B)) \subset \Spec(A)$이고 점
$x \in X$를 포함하며 $f(x) = y$인 것을 택하자. 전사
$k[x_1, \ldots, x_n] \to A$를 택하고 그 핵을 $I$라 하자. 전사
$A[y_1, \ldots, y_m] \to B$를 택하면 다시 전사
$k[x_i, y_j] \to B$가 생기며, 그 핵을 $J$라 하자.
$\mathfrak q \subset k[x_i, y_j]$를
$y \in \Spec(B)$에 대응하는 소 아이디얼이라 하고,
$\mathfrak p \subset k[x_i]$를 $x \in \Spec(A)$에 대응하는 소 아이디얼이라
하자. $x$가 $y$로 가므로 $\mathfrak p = \mathfrak q \cap k[x_i]$이다.
다음 국소환들의 가환 그림을 고찰하자.
$$
\xymatrix{
\mathcal{O}_{X, x} \ar@{=}[r] &
B_{\mathfrak q} &
k[x_1, \ldots, x_n, y_1, \ldots, y_m]_{\mathfrak q} \ar[l] \\
\mathcal{O}_{Y, y} \ar@{=}[r] \ar[u] & A_{\mathfrak p} \ar[u] &
k[x_1, \ldots, x_n]_{\mathfrak p} \ar[l] \ar[u]
}
$$
대수, 보조정리 \ref{algebra-lemma-lci-permanence-initial}의 가정들이
성립한다고 주장한다. 조건 (1)과 (2)는 자명하다. 조건 (4)는
$X \to Y$가 평탄이므로 따른다. 조건 (3)은 환
$\mathcal{O}_{Y, y}$와
$\mathcal{O}_{X_y, x} = \mathcal{O}_{X, x}/\mathfrak m_y\mathcal{O}_{X, x}$가
$f$와 $p$가 신토믹이라는 가정에 의해 완전교차환이므로 따른다.
보조정리 \ref{morphisms-lemma-syntomic-locally-standard-syntomic}을 보라.
대수, 보조정리 \ref{algebra-lemma-lci-permanence-initial}의 결론은 정확히
$\mathcal{O}_{Y, y}$가 완전교차환이라는 것이다! 따라서 다시 보조정리
\ref{morphisms-lemma-syntomic-locally-standard-syntomic}에 의해 $Y$는 $k$ 위에서
$y$에서 신토믹이고, 이것이 원하는 결론이다.
\end{proof}
\section{몰입의 코노멀 층}
\label{morphisms-section-conormal-sheaf}

\noindent
$i : Z \to X$를 닫힌 몰입이라 하자. $\mathcal{I} \subset \mathcal{O}_X$를
이에 대응하는 준연접 아이디얼 층이라 하자. $X$ 위 준연접 층들의 짧은
완전열
$$
0 \to \mathcal{I}^2 \to \mathcal{I} \to \mathcal{I}/\mathcal{I}^2 \to 0
$$
을 고찰하자. 층 $\mathcal{I}/\mathcal{I}^2$는 $\mathcal{I}$에 의해
소멸되므로 보조정리 \ref{morphisms-lemma-i-star-equivalence}에 의해 $Z$ 위의 층에
대응한다. 이 준연접 $\mathcal{O}_Z$-가군을
{\it $Z$의 $X$ 안에서의 코노멀 층}이라 하며, 절
\ref{morphisms-section-closed-immersions-quasi-coherent}에서 언급한 표기 남용에
따라 흔히 단순히 $\mathcal{I}/\mathcal{I}^2$로 쓴다.

\medskip\noindent
$i : Z \to X$가 (국소 닫힌) 몰입인 경우, $i$의 코노멀 층을 닫힌 몰입
$i : Z \to X \setminus \partial Z$의 코노멀 층으로 정의한다. 여기서
$\partial Z = \overline{Z} \setminus Z$이다. 이는 흔히
$\mathcal{I}/\mathcal{I}^2$로 쓰며, 이때 $\mathcal{I}$는 닫힌 몰입
$i : Z \to X \setminus \partial Z$의 아이디얼 층이다.

\begin{definition}
\label{morphisms-definition-conormal-sheaf}
$i : Z \to X$를 몰입이라 하자. {\it 코노멀 층 $\mathcal{C}_{Z/X}$, 즉
$Z$의 $X$ 안에서의 코노멀 층} 또는 {\it $i$의 코노멀 층}은 위에서
설명한 준연접 $\mathcal{O}_Z$-가군 $\mathcal{I}/\mathcal{I}^2$이다.
\end{definition}

\noindent
\cite[IV Definition 16.1.2]{EGA}에서는 이 층을 $\mathcal{N}_{Z/X}$로
쓴다. 우리는 $\mathcal{N}_{Z/X}$라는 표기를
{\it 몰입의 노멀 층}에 남겨 두고자 하므로 이 관례를 따르지 않는다.
코노멀 층이 유한 표시인 경우 이를
$$
\mathcal{N}_{Z/X} =
\SheafHom_{\mathcal{O}_Z}(\mathcal{C}_{Z/X}, \mathcal{O}_Z) =
\SheafHom_{\mathcal{O}_Z}(\mathcal{I}/\mathcal{I}^2, \mathcal{O}_Z)
$$
로 정의한다(그렇지 않으면 노멀 층은 준연접조차 아닐 수 있다).
노멀 층은 뒤에서 다시 다룰 것이다(여기에 나중 참조를 넣을 것).

\begin{lemma}
\label{morphisms-lemma-affine-conormal}
$i : Z \to X$를 몰입이라 하자. $i$의 코노멀 층은 다음 성질을 갖는다.
\begin{enumerate}
\item $U \subset X$를 임의의 열린 부분스킴이라 하고, $i$가
$Z \xrightarrow{i'} U \to X$로 분해되며 $i'$가 닫힌 몰입이라고 하자.
$\mathcal{I} = \Ker((i')^\sharp) \subset \mathcal{O}_U$로 두자. 그러면
$$
\mathcal{C}_{Z/X} = (i')^*\mathcal{I}\quad\text{and}\quad
i'_*\mathcal{C}_{Z/X} = \mathcal{I}/\mathcal{I}^2
$$
이다.
\item 임의의 아핀 열린집합 $\Spec(R) = U \subset X$ 중
$Z \cap U = \Spec(R/I)$인 것에 대해 표준 동형사상
$\Gamma(Z \cap U, \mathcal{C}_{Z/X}) = I/I^2$가 존재한다.
\end{enumerate}
\end{lemma}

\begin{proof}
대부분 정의에서 명백하다. 환 $R$과 아이디얼 $I$가 주어지고 후자가
$R$의 아이디얼이면
$I/I^2 = I \otimes_R R/I$임에 유의하자. 세부사항은 생략한다.
\end{proof}

\begin{lemma}
\label{morphisms-lemma-conormal-functorial}
스킴 범주에서 가환 그림
$$
\xymatrix{
Z \ar[r]_i \ar[d]_f & X \ar[d]^g \\
Z' \ar[r]^{i'} & X'
}
$$
이 주어졌다고 하자. $i$, $i'$가 몰입이라고 가정하자. 다음 표준
$\mathcal{O}_Z$-가군 사상이 존재한다.
$$
f^*\mathcal{C}_{Z'/X'}
\longrightarrow
\mathcal{C}_{Z/X}
$$
이 사상은 다음 성질로 특징지어진다. 모든 아핀 열린집합의 쌍
$(\Spec(R) = U \subset X, \Spec(R') = U' \subset X')$ 중
$g(U) \subset U'$이고 $Z \cap U = \Spec(R/I)$ 및
$Z' \cap U' = \Spec(R'/I')$인 것에 대해 유도된 사상
$$
\Gamma(Z' \cap U', \mathcal{C}_{Z'/X'}) = I'/I'^2
\longrightarrow
I/I^2 = \Gamma(Z \cap U, \mathcal{C}_{Z/X})
$$
은 환 준동형 $g^\sharp : R' \to R$에서 유도된 사상이며, 이 환 준동형은
$g^\sharp(I') \subset I$를 만족한다.
\end{lemma}

\begin{proof}
$\partial Z' = \overline{Z'} \setminus Z'$ 및
$\partial Z = \overline{Z} \setminus Z$로 두자. 이들은 각각 $X'$와 $X$의
닫힌 부분집합이다. $X'$를 $X' \setminus \partial Z'$로, $X$를
$X \setminus \big(g^{-1}(\partial Z') \cup \partial Z\big)$로 바꾸면
$i$와 $i'$가 닫힌 몰입이라고 가정할 수 있다.

\medskip\noindent
$g \circ i$가 $i'$를 통해 분해된다는 사실은 $g^*\mathcal{I}'$가
$\mathcal{I}$ 안으로 사상됨을 뜻한다. 여기서 사상은 표준 사상
$g^*\mathcal{I}' \to \mathcal{O}_X$에서 얻는다. 스킴, 보조정리
\ref{schemes-lemma-characterize-closed-subspace}와
\ref{schemes-lemma-restrict-map-to-closed}를 보라. 따라서 준연접 층들의
유도된 사상
$g^*(\mathcal{I}'/(\mathcal{I}')^2) \to \mathcal{I}/\mathcal{I}^2$를 얻는다.
$i$로 당겨오면
$i^*g^*(\mathcal{I}'/(\mathcal{I}')^2) \to i^*(\mathcal{I}/\mathcal{I}^2)$를
얻는다. $i^*(\mathcal{I}/\mathcal{I}^2) = \mathcal{C}_{Z/X}$임에 유의하자.
한편
$i^*g^*(\mathcal{I}'/(\mathcal{I}')^2) =
f^*(i')^*(\mathcal{I}'/(\mathcal{I}')^2) = f^*\mathcal{C}_{Z'/X'}$이다.
이것이 원하는 사상을 준다.

\medskip\noindent
이 사상이 국소적으로 주어진 사상 $I'/(I')^2 \to I/I^2$로 기술됨은
정의를 풀어 확인하면 되므로 생략한다. 또 임의의 $x \in i(Z)$가 주어지면
아핀 열린 근방 $U$, $U'$가 실제로 존재하여 $f(U) \subset U'$이고
$Z \cap U$ 및 $U' \cap Z'$가 닫히며 $x \in U$이다. 증명은 생략한다. 따라서
보조정리의 조건은 정말로 이 사상을 특징짓는다(그리고 이 조건으로
사상을 정의할 수도 있었다).
\end{proof}

\begin{lemma}
\label{morphisms-lemma-conormal-functorial-flat}
스킴 범주에서 올곱 그림
$$
\xymatrix{
Z \ar[r]_i \ar[d]_f & X \ar[d]^g \\
Z' \ar[r]^{i'} & X'
}
$$
이 주어지고 $i$, $i'$가 몰입이라고 하자. 그러면 보조정리
\ref{morphisms-lemma-conormal-functorial}의 표준 사상
$f^*\mathcal{C}_{Z'/X'} \to \mathcal{C}_{Z/X}$는 전사이다.
$g$가 평탄이면 이 사상은 동형사상이다.
\end{lemma}

\begin{proof}
$R' \to R$을 환 준동형이라 하고 $I' \subset R'$를 아이디얼이라 하자.
$I = I'R$로 두자. 그러면 $I'/(I')^2 \otimes_{R'} R \to I/I^2$는
전사이다. $R' \to R$이 평탄이면 $I = I' \otimes_{R'} R$이고
$I^2 = (I')^2 \otimes_{R'} R$이므로 이 사상은 동형사상이다.
\end{proof}

\begin{lemma}
\label{morphisms-lemma-transitivity-conormal}
$Z \to Y \to X$를 스킴의 몰입들이라 하자. 그러면 표준 완전열
$$
i^*\mathcal{C}_{Y/X} \to
\mathcal{C}_{Z/X} \to
\mathcal{C}_{Z/Y} \to 0
$$
이 존재한다. 여기서 사상들은 보조정리
\ref{morphisms-lemma-conormal-functorial}에서 오고 $i : Z \to Y$는 첫째 사상이다.
\end{lemma}

\begin{proof}
보조정리 \ref{morphisms-lemma-conormal-functorial}을 통해 이는 다음 대수적 사실로
옮겨진다. $C \to B \to A$가 전사인 환 준동형들이라고 하자.
$I = \Ker(B \to A)$, $J = \Ker(C \to A)$ 및
$K = \Ker(C \to B)$로 두자. 그러면 완전열
$$
K/K^2 \otimes_B A \to J/J^2 \to I/I^2 \to 0.
$$
이 존재한다. 이는 $I = J/K$라는 관찰에서 곧바로 따른다.
\end{proof}
\section{사상의 미분층}
\label{morphisms-section-sheaf-differentials}

\noindent
독자에게 가환대수 장의 해당 절(대수, 절
\ref{algebra-section-differentials})과 가군층 장의 해당 절(가군, 절
\ref{modules-section-differentials})을 살펴볼 것을 권한다.

\begin{definition}
\label{morphisms-definition-sheaf-differentials}
$f : X \to S$를 스킴의 사상이라 하자.
{\it 미분층 $\Omega_{X/S}$, 즉 $X$의 $S$ 위에서의 미분층}은 $f$를
환 달린 공간의 사상으로 보았을 때의 미분층이며(가군, 정의
\ref{modules-definition-differentials}), 그 {\it 보편 $S$-미분}
$$
\text{d}_{X/S} : \mathcal{O}_X \longrightarrow \Omega_{X/S}.
$$
과 함께 주어진다.
\end{definition}

\noindent
$\Omega_{X/S}$는 준연접 $\mathcal{O}_X$-가군임이 드러난다. 예를 들어
이는 대각 사상 $\Delta : X \to X \times_S X$의 코노멀 층과 동형이다
(보조정리 \ref{morphisms-lemma-differentials-diagonal}). 우리는 $X$의 $S$ 위에서의
미분 가군을 보편 성질, 즉 보편 미분을 받아들이는 대상으로 정의했다.
상대 미분층의 다른 구성이 이 보편 성질을 만족한다면 요네다 보조정리에
의해 위에서 정의한 것과 표준적으로 동형이다. 편의를 위해 보편 성질을
다시 적는다.

\begin{lemma}
\label{morphisms-lemma-universal-derivation-universal}
$f : X \to S$를 스킴의 사상이라 하자. 사상
$$
\Hom_{\mathcal{O}_X}(\Omega_{X/S}, \mathcal{F})
\longrightarrow
\text{Der}_S(\mathcal{O}_X, \mathcal{F}),\quad
\alpha \longmapsto \alpha \circ \text{d}_{X/S}
$$
은 함자 $\textit{Mod}(\mathcal{O}_X) \to \textit{Sets}$의 동형사상이다.
\end{lemma}

\begin{proof}
이는 정의를 다시 말한 것뿐이다.
\end{proof}

\begin{lemma}
\label{morphisms-lemma-differentials-restrict-open}
$f : X \to S$를 스킴의 사상이라 하자. $U \subset X$, $V \subset S$를
$f(U) \subset V$인 열린 부분스킴이라 하자. 다음을 만족하는 동형사상
$\Omega_{X/S}|_U = \Omega_{U/V}$가 $\mathcal{O}_U$-가군 사이에 유일하게
존재한다.
$\text{d}_{X/S}|_U = \text{d}_{U/V}$.
\end{lemma}

\begin{proof}
표준 동일시
$f^{-1}\mathcal{O}_S|_U = (f|_U)^{-1}\mathcal{O}_V$를 사용하면 이는
가군, 보조정리 \ref{modules-lemma-localize-differentials}의 특수한 경우이다.
\end{proof}

\noindent
이제부터 이 표준 동일시들을 사용하고 $\Omega_{U/S}$ 또는
$\Omega_{U/V}$로 $\Omega_{X/S}$의 $U$에 대한 제한을 간단히 나타낸다.

\begin{lemma}
\label{morphisms-lemma-affine-case-derivation}
$R \to A$를 환 준동형이라 하자. $\mathcal{F}$를 $\mathcal{O}_X$-가군층이라
하고, 이 층은 $X = \Spec(A)$ 위에 놓인다고 하자. $S = \Spec(R)$로 두자.
$S$-미분으로서 $\mathcal{F}$ 위에 주어진 것을 그 대역 단면에 대한
작용으로 보내는 규칙은 $S$-미분 중 $\mathcal{F}$ 위에 주어진 것들의
집합과 $R$-미분 중 $M = \Gamma(X, \mathcal{F})$ 위에 주어진 것들의
집합 사이의 전단사를 정의한다.
\end{lemma}

\begin{proof}
$D : A \to M$을 $R$-미분이라 하자. $S$-미분으로서 $\mathcal{F}$ 위에
주어지고 대역 단면 위에서 $D$를 주는 것이 유일하게 존재함을 보여야 한다.
$U = D(f) \subset X$를 표준 아핀 열린집합이라 하자.
$\Gamma(U, \mathcal{O}_X)$의 모든 원소는 $a/f^n$ 꼴이며, 여기서
$a \in A$이고 $n \geq 0$이다. 라이프니츠 법칙에 의해
$$
D(a)|_U = a/f^n D(f^n)|_U + f^n D(a/f^n)
$$
가 $\Gamma(U, \mathcal{F})$에서 성립한다. $f$는
$\Gamma(U, \mathcal{F})$ 위에서 가역적으로 작용하므로 이는
$D(a/f^n) \in \Gamma(U, \mathcal{F})$의 값을 완전히 결정한다.
이로써 유일성이 증명된다. 존재성은 단순히
$$
D(a/f^n) := (1/f^n) D(a)|_U - a/f^{2n} D(f^n)|_U
$$
로 정의하고 이것이 원하는 모든 성질을 가짐을 ($X$의 표준 열린집합
기저 위에서) 보이면 따른다. 세부사항은 생략한다.
\end{proof}

\begin{lemma}
\label{morphisms-lemma-differentials-affine}
$f : X \to S$를 스킴의 사상이라 하자. 임의의 아핀 열린집합의 쌍
$\Spec(A) = U \subset X$, $\Spec(R) = V \subset S$ 중
$f(U) \subset V$인 것에 대해
유일한 동형사상
$$
\Gamma(U, \Omega_{X/S}) = \Omega_{A/R}.
$$
이 존재하고, 이는 $\text{d}_{X/S}$ 및 $\text{d} : A \to \Omega_{A/R}$와
호환된다.
\end{lemma}

\begin{proof}
보조정리 \ref{morphisms-lemma-differentials-restrict-open}에 의해 $X$와 $S$를 각각
$U$와 $V$로 바꿀 수 있다. 따라서 $X = \Spec(A)$와 $S = \Spec(R)$라고
가정할 수 있고, $U = X$ 및 $V = S$인 경우에 보조정리를 보이면 된다.
$A$-가군 $M = \Gamma(X, \Omega_{X/S})$과 $R$-미분
$\text{d}_{X/S} : A \to M$을 고찰하자. $N$을 또 다른 $A$-가군이라 하고
$\widetilde{N}$으로 준연접 $\mathcal{O}_X$-가군 중 $N$에 대응하는 것을
나타내자. 스킴, 절 \ref{schemes-section-quasi-coherent-affine}을 보라.
$\text{d}_{X/S} : A \to M$과 미리 합성하여 사상
$$
\alpha : \Hom_A(M, N) \longrightarrow \text{Der}_R(A, N)
$$
을 얻는다. 보조정리 \ref{morphisms-lemma-universal-derivation-universal}과
\ref{morphisms-lemma-affine-case-derivation}을 사용하면 동일시
$$
\Hom_{\mathcal{O}_X}(\Omega_{X/S}, \widetilde{N}) =
\text{Der}_S(\mathcal{O}_X, \widetilde{N}) =
\text{Der}_R(A, N)
$$
을 얻는다. 대역 단면을 취하면 사상
$\Hom_{\mathcal{O}_X}(\Omega_{X/S}, \widetilde{N}) \to \Hom_R(M, N)$이
정해진다. 이 사상과 위의 동일시들을 결합하면 사상
$$
\beta : \text{Der}_R(A, N) \longrightarrow \Hom_R(M, N)
$$
을 얻는다. 대역 단면에서 일어나는 일을 확인하면 $\alpha$와 $\beta$가
서로 역임을 알 수 있다. 따라서 $\text{d}_{X/S} : A \to M$은
$\text{d} : A \to \Omega_{A/R}$와 같은 보편 성질을 만족한다.
대수, 보조정리 \ref{algebra-lemma-universal-omega}를 보라. 그러므로 요네다
보조정리(범주, 보조정리 \ref{categories-lemma-yoneda})에 의해 미분과
호환되는 유일한 $A$-가군 동형사상 $M \cong \Omega_{A/R}$가 존재한다.
\end{proof}
\begin{remark}
\label{morphisms-remark-differentials-glue}
위 보조정리는 미분 가군을 구성하는 두 번째 방법을 준다. 즉
$f : X \to S$를 스킴의 사상이라 하자. 모든 아핀 열린집합 $U \subset X$ 중
$S$의 어떤 아핀 열린집합 안으로 사상되는 것들의 모임을 고찰하자. 이들은
$X$의 위상의 기저를 이룬다. 따라서 $\Gamma(U, \Omega_{X/S})$를 그러한
$U$에 대해 정의하면 충분하다. 단순히
$\Gamma(U, \Omega_{X/S}) = \Omega_{A/R}$로 두며, 여기서 $A$, $R$은 위
보조정리 \ref{morphisms-lemma-differentials-affine}에서와 같다. 이 방법은 작동하지만,
제한 사상을 구성하고 이 방안으로 층 조건을 확인하려면 대수적 준비가
조금 더 필요하다.
\end{remark}

\noindent
다음 보조정리는 미분층을 정의하는 또 다른 방법을 주며, 특히
$\Omega_{X/S}$가 $X$와 $S$가 스킴일 때 준연접임을 보여 준다.

\begin{lemma}
\label{morphisms-lemma-differentials-diagonal}
$f : X \to S$를 스킴의 사상이라 하자. $\Omega_{X/S}$와 대각 사상
$\Delta_{X/S} : X \longrightarrow X \times_S X$의 코노멀 층 사이에
표준 동형사상이 존재한다.
\end{lemma}

\begin{proof}
먼저 두 개의 ``대역'' 층과 층들의 대역 사상들이 존재함을 보이고,
뒤에서 몇몇 아핀 열린집합 위에서 이 구성들을 기술한다.

\medskip\noindent
$\Delta = \Delta_{X/S} : X \to X \times_S X$가 몰입임을 상기하자.
스킴, 보조정리 \ref{schemes-lemma-diagonal-immersion}을 보라.
$\mathcal{J}$를 이 몰입의 아이디얼 층이라 하자. 이는
어떤 열린 부분스킴 $W$, 즉 $X \times_S X$의 열린 부분스킴으로서
$\Delta(X) \subset W$가 닫힌 것 위에 놓인다. 스킴, 보조정리
\ref{schemes-lemma-diagonal-immersion}의 증명에서 찾은 것을 택하자.
환의 층 $\mathcal{O}_W/\mathcal{J}^2$는 $\Delta(X)$ 위에 지지됨에
유의하자. 또한 이는 층들의 짧은 완전열
$$
0 \to \mathcal{J}/\mathcal{J}^2
\to \mathcal{O}_W/\mathcal{J}^2
\to \Delta_*\mathcal{O}_X
\to 0.
$$
에 놓인다. $\Delta^{-1}$을 사용하면 이를 $f^{-1}\mathcal{O}_S$-대수층의
전사로 생각할 수 있으며, 그 핵은 $\Delta$의 코노멀 층이다(정의
\ref{morphisms-definition-conormal-sheaf}와 보조정리 \ref{morphisms-lemma-affine-conormal}을
보라).
$$
0 \to \mathcal{C}_{X/X \times_S X}
\to \Delta^{-1}(\mathcal{O}_W/\mathcal{J}^2)
\to \mathcal{O}_X
\to 0
$$
이로써 가군, 보조정리
\ref{modules-lemma-double-structure-gives-derivation}의 상황에 놓인다.
사영 사상들 $p_i : X \times_S X \to X$, $i = 1, 2$는 환의 층의 사상
$(p_i)^\sharp : (p_i)^{-1}\mathcal{O}_X \to \mathcal{O}_{X \times_S X}$를
유도한다. $W$로 제한하고 $\mathcal{J}^2$로 나누면
$(p_i)^{-1}\mathcal{O}_X \to \mathcal{O}_W/\mathcal{J}^2$를 얻는다.
$\Delta^{-1}p_i^{-1}\mathcal{O}_X = \mathcal{O}_X$이므로 사상
$$
s_i : \mathcal{O}_X \to \Delta^{-1}(\mathcal{O}_W/\mathcal{J}^2).
$$
을 얻는다. $s_1$과 $s_2$는 모두 가군, 보조정리
\ref{modules-lemma-double-structure-gives-derivation}에서와 같이 사상
$\Delta^{-1}(\mathcal{O}_W/\mathcal{J}^2) \to \mathcal{O}_X$의 단면이다.
따라서 $S$-미분
$\text{d} = s_2 - s_1 : \mathcal{O}_X \to \mathcal{C}_{X/X \times_S X}$를
얻는다. 미분 가군의 보편 성질에 의해 유일한
$\mathcal{O}_X$-선형 사상
$$
\Omega_{X/S} \longrightarrow \mathcal{C}_{X/X \times_S X},\quad
f\text{d}g \longmapsto fs_2(g) - fs_1(g)
$$
을 얻는다. 이 사상이 동형사상임을 보이기 위해 적절한 아핀 열린집합
위에서 구체적으로 계산하자. $X$를 아핀 열린집합
$\Spec(A) = U \subset X$들로 덮을 수 있으며, 각각의 상은 아핀 열린집합
$\Spec(R) = V \subset S$에 포함된다. 스킴, 보조정리
\ref{schemes-lemma-diagonal-immersion}의 증명에 따르면
$U \times_V U \subset X \times_S X$는 위에서 언급한 열린집합 $W$에
포함되는 아핀 열린집합이다. 또한
$U \times_V U = \Spec(A \otimes_R A)$이다. 층 $\mathcal{J}$는 아이디얼
$J = \Ker(A \otimes_R A \to A)$에 대응한다. 앞의 짧은 완전열은
$A \otimes_R A$-가군들의 짧은 완전열
$$
0 \to J/J^2 \to (A \otimes_R A)/J^2 \to A \to 0
$$
이 된다. 단면 $s_i$는 환 준동형들
$$
A \longrightarrow (A \otimes_R A)/J^2,\quad
s_1 : a \mapsto a \otimes 1,\quad
s_2 : a \mapsto 1 \otimes a.
$$
에 대응한다. 보조정리 \ref{morphisms-lemma-affine-conormal}에 의해
$\Gamma(U, \mathcal{C}_{X/X \times_S X}) = J/J^2$이고, 보조정리
\ref{morphisms-lemma-differentials-affine}에 의해
$\Gamma(U, \Omega_{X/S}) = \Omega_{A/R}$이다. 위의 사상은
$a \text{d}b \mapsto a \otimes b - ab \otimes 1$이며, 대수, 보조정리
\ref{algebra-lemma-differentials-diagonal}에서 이 사상이 동형사상임을
보인다.
\end{proof}
\begin{lemma}
\label{morphisms-lemma-functoriality-differentials}
스킴들의 가환 그림
$$
\xymatrix{
X' \ar[d] \ar[r]_f & X \ar[d] \\
S' \ar[r] & S
}
$$
이 주어졌다고 하자. 표준 사상 $\mathcal{O}_X \to f_*\mathcal{O}_{X'}$와
사상 $f_*\text{d}_{X'/S'} : f_*\mathcal{O}_{X'} \to f_*\Omega_{X'/S'}$의
합성은 $S$-미분이다. 따라서 표준 $\mathcal{O}_X$-가군 사상
$\Omega_{X/S} \to f_*\Omega_{X'/S'}$를 얻고, $f_*$와 $f^*$의 수반성에
의해 표준 $\mathcal{O}_{X'}$-가군 준동형
$$
c_f : f^*\Omega_{X/S} \longrightarrow \Omega_{X'/S'}.
$$
을 얻는다. 이는 $f^*\text{d}_{X/S}(h)$가
$\text{d}_{X'/S'}(f^* h)$로 사상된다는 성질로 유일하게 특징지어진다.
여기서 $h$는 $\mathcal{O}_X$의 임의의 국소 단면이다.
\end{lemma}

\begin{proof}
이는 가군, 보조정리
\ref{modules-lemma-functoriality-differentials-ringed-spaces}의 특수한
경우이다. 스킴의 경우에는 코노멀 층의 함자성(보조정리
\ref{morphisms-lemma-conormal-functorial}을 보라)과 보조정리
\ref{morphisms-lemma-differentials-diagonal}을 사용하여 $c_f$를 정의할 수도 있다.
또는 보조정리의 마지막 줄에 있는 특징화를 사용하여 아핀 조각 위에서
정의된 사상들을 붙일 수도 있다(대수, 식
(\ref{algebra-equation-functorial-omega})을 보라).
\end{proof}

\begin{lemma}
\label{morphisms-lemma-triangle-differentials}
$f : X \to Y$, $g : Y \to S$를 스킴의 사상이라 하자. 그러면 표준 완전열
$$
f^*\Omega_{Y/S} \to \Omega_{X/S} \to \Omega_{X/Y} \to 0
$$
이 존재하며, 사상들은 보조정리
\ref{morphisms-lemma-functoriality-differentials}를 적용하여 얻는다.
\end{lemma}

\begin{proof}
이는 대수, 보조정리 \ref{algebra-lemma-exact-sequence-differentials}를
층화한 판이다. 또는 환 달린 공간의 사상에 대한 일반판은 가군, 보조정리
\ref{modules-lemma-triangle-differentials}에 있다.
\end{proof}

\begin{lemma}
\label{morphisms-lemma-base-change-differentials}
$X \to S$를 스킴의 사상이라 하자. $g : S' \to S$를 스킴의 사상이라
하자. $X' = X_{S'}$를 $X$의 기저변환이라 하고, $g' : X' \to X$를
사영이라 하자. 그러면 보조정리
\ref{morphisms-lemma-functoriality-differentials}의 사상
$$
(g')^*\Omega_{X/S} \to \Omega_{X'/S'}
$$
은 동형사상이다.
\end{lemma}

\begin{proof}
이는 대수, 보조정리 \ref{algebra-lemma-differentials-base-change}를
층화한 판이다.
\end{proof}

\begin{lemma}
\label{morphisms-lemma-differential-product}
$f : X \to S$와 $g : Y \to S$를 같은 목표를 갖는 스킴의 사상이라 하자.
$p : X \times_S Y \to X$와 $q : X \times_S Y \to Y$를 사영 사상이라 하자.
보조정리 \ref{morphisms-lemma-functoriality-differentials}에서 오는 사상
$$
p^*\Omega_{X/S} \oplus q^*\Omega_{Y/S}
\longrightarrow
\Omega_{X \times_S Y/S}
$$
은 동형사상이다.
\end{lemma}

\begin{proof}
보조정리 \ref{morphisms-lemma-base-change-differentials}에 의해 합성
$p^*\Omega_{X/S} \to \Omega_{X \times_S Y/S} \to \Omega_{X \times_S Y/Y}$는
동형사상이고, $q$에 대해서도 마찬가지이다. 또한 사상
$p^*\Omega_{X/S} \to \Omega_{X \times_S Y/S}$의 여핵은 보조정리
\ref{morphisms-lemma-triangle-differentials}에 의해 $\Omega_{X \times_S Y/X}$이다.
결과가 따른다.
\end{proof}

\begin{lemma}
\label{morphisms-lemma-finite-type-differentials}
$f : X \to S$를 스킴의 사상이라 하자. $f$가 국소적으로 유한형이면
$\Omega_{X/S}$는 유한형 $\mathcal{O}_X$-가군이다.
\end{lemma}

\begin{proof}
대수, 보조정리 \ref{algebra-lemma-differentials-finitely-generated},
보조정리 \ref{morphisms-lemma-differentials-affine},
보조정리 \ref{morphisms-lemma-locally-finite-type-characterize}, 성질, 보조정리
\ref{properties-lemma-finite-type-module}에서 곧바로 따른다.
\end{proof}

\begin{lemma}
\label{morphisms-lemma-finite-presentation-differentials}
$f : X \to S$를 스킴의 사상이라 하자. $f$가 국소적으로 유한 표시이면
$\Omega_{X/S}$는 유한 표시인 $\mathcal{O}_X$-가군이다.
\end{lemma}

\begin{proof}
대수, 보조정리 \ref{algebra-lemma-differentials-finitely-presented},
보조정리 \ref{morphisms-lemma-differentials-affine},
보조정리 \ref{morphisms-lemma-locally-finite-presentation-characterize}, 성질,
보조정리 \ref{properties-lemma-finite-presentation-module}에서 곧바로 따른다.
\end{proof}
\begin{lemma}
\label{morphisms-lemma-immersion-differentials}
$X \to S$가 몰입이거나, 더 일반적으로 모노사상이면
$\Omega_{X/S}$는 영이다.
\end{lemma}

\begin{proof}
이 경우 $\Delta_{X/S}$는 동형사상이어서 자명한 코노멀 층을 가지므로
참이다. 따라서 보조정리 \ref{morphisms-lemma-differentials-diagonal}에 의해
$\Omega_{X/S} = 0$이다. 대수적 판은 대수, 보조정리
\ref{algebra-lemma-trivial-differential-surjective}이다.
\end{proof}

\begin{lemma}
\label{morphisms-lemma-differentials-relative-immersion}
$i : Z \to X$를 $S$ 위의 스킴들의 몰입이라 하자. 표준 완전열
$$
\mathcal{C}_{Z/X} \to i^*\Omega_{X/S} \to \Omega_{Z/S} \to 0
$$
이 존재한다. 첫째 화살표는 $\text{d}_{X/S}$에서 유도되고 둘째 화살표는
보조정리 \ref{morphisms-lemma-functoriality-differentials}에서 온다.
\end{lemma}

\begin{proof}
이는 대수, 보조정리 \ref{algebra-lemma-differential-seq}를 층화한 판이다.
다만 첫째 화살표를 대역적으로 정의할 수 있는지 확인해야 한다. 따라서
여기서 ``$\text{d}_{X/S}$에서 유도된다''는 말의 뜻을 설명한다.
$i$가 닫힌 몰입이라고 가정할 수 있도록 $X$를 줄일 수 있다.
$\mathcal{I} \subset \mathcal{O}_X$를 $Z \subset X$에 대응하는 아이디얼
층이라 하자. 그러면 $\text{d}_{X/S} : \mathcal{I} \to \Omega_{X/S}$는
부분층 $\mathcal{I}^2 \subset \mathcal{I}$를
$\mathcal{I}\Omega_{X/S}$로 보낸다. 따라서
$\mathcal{I}/\mathcal{I}^2 \to \Omega_{X/S}/\mathcal{I}\Omega_{X/S}$인
$\mathcal{O}_X/\mathcal{I}$-선형 사상을 유도한다. 보조정리
\ref{morphisms-lemma-i-star-equivalence}에 의해 이는 원하는 사상
$\mathcal{C}_{Z/X} \to i^*\Omega_{X/S}$에 대응한다.
\end{proof}

\begin{lemma}
\label{morphisms-lemma-differentials-relative-immersion-section}
$i : Z \to X$를 $S$ 위의 스킴들의 몰입이라 하고, $i$가 (국소적으로)
왼쪽 역을 갖는다고 가정하자. 보조정리
\ref{morphisms-lemma-differentials-relative-immersion}의 표준 열
$$
0 \to \mathcal{C}_{Z/X} \to i^*\Omega_{X/S} \to \Omega_{Z/S} \to 0
$$
은 (국소적으로) 분할 완전열이다. 특히 $s : S \to X$가 구조 사상
$X \to S$의 단면이면 사상
$\mathcal{C}_{S/X} \to s^*\Omega_{X/S}$는 $\text{d}_{X/S}$에서 유도되며
동형사상이다.
\end{lemma}

\begin{proof}
대수, 보조정리 \ref{algebra-lemma-differential-seq-split}에 의해
$X, Z, S$가 아핀인 경우에 결과가 성립한다. 일반적으로 이 열은
국소적으로 분할 완전열임을 얻는다. $g : X \to Z$가 $i$의 왼쪽 역이면
$i^*c_g$는 가군, 보조정리
\ref{modules-lemma-check-functoriality-differentials}에 의해 사상
$i^*\Omega_{X/S} \to \Omega_{Z/S}$의 오른쪽 역이다. 그러므로 호몰로지,
보조정리 \ref{homology-lemma-ses-split}에 의해 열이 분할된다.
마지막으로 $s$가 단면이면 이는 몰입 $s : Z = S \to X$, 즉 $S$ 위의
몰입이고
(스킴, 보조정리 \ref{schemes-lemma-section-immersion}을 보라), 이 경우
$\Omega_{Z/S} = 0$이다.
\end{proof}

\begin{remark}
\label{morphisms-remark-differentials-diagonal}
$X \to S$를 스킴의 사상이라 하자. 보조정리
\ref{morphisms-lemma-differential-product}에 의해
$$
\Omega_{X \times_S X/S} =
\text{pr}_1^*\Omega_{X/S} \oplus \text{pr}_2^*\Omega_{X/S}
$$
이다. 한편 대각 사상 $\Delta : X \to X \times_S X$는 국소적으로 왼쪽
역을 갖는 몰입이다. 따라서 보조정리
\ref{morphisms-lemma-differentials-relative-immersion-section}에 의해 표준 짧은
완전열
$$
0 \to \mathcal{C}_{X/X \times_S X} \to \Omega_{X/S} \oplus \Omega_{X/S}
\to \Omega_{X/S} \to 0
$$
을 얻는다. 오른쪽 화살표는 실제로 분할 전사인 $(1, 1)$임에 유의하자.
한편 보조정리 \ref{morphisms-lemma-differentials-diagonal}에 의해 동일시
$\Omega_{X/S} = \mathcal{C}_{X/X \times_S X}$가 있다. 이 동일시에서
$\text{d}_{X/S}(f) = s_2(f) - s_1(f)$를 택했으므로 왼쪽 화살표는
$(-1, 1)$이다\footnote{즉 국소 단면
$\text{d}_{X/S}(f) = 1 \otimes f - f \otimes 1$, 즉 $\Delta$의 아이디얼
층의 이 단면은
$\text{d}_{X \times_S X/X}$에 의해 국소 단면
$1 \otimes 1 \otimes 1 \otimes f - 1 \otimes f \otimes 1 \otimes 1
-1 \otimes 1 \otimes f \otimes 1 + f \otimes 1 \otimes 1 \otimes 1 =
\text{pr}_2^*\text{d}_{X/S}(f) - \text{pr}_1^*\text{d}_{X/S}(f)$로 사상된다.}.
\end{remark}

\begin{lemma}
\label{morphisms-lemma-two-immersions}
스킴들의 가환 그림
$$
\xymatrix{
Z \ar[r]_i \ar[rd]_j & X \ar[d] \\
& Y
}
$$
이 주어지고 $i$와 $j$가 몰입이라고 하자. 그러면 표준 완전열
$$
\mathcal{C}_{Z/Y} \to
\mathcal{C}_{Z/X} \to
i^*\Omega_{X/Y} \to 0
$$
이 존재한다. 첫째 화살표는 보조정리
\ref{morphisms-lemma-conormal-functorial}에서 오고, 둘째 화살표는 보조정리
\ref{morphisms-lemma-differentials-relative-immersion}에서 온다.
\end{lemma}

\begin{proof}
이 명제의 대수적 판은 대수, 보조정리
\ref{algebra-lemma-application-NL}이다.
\end{proof}
\section{유한 차수 미분 작용소}
\label{morphisms-section-differential-operators}

\noindent
독자에게 가환대수 장의 해당 절(대수, 절
\ref{algebra-section-differential-operators})과 가군층 장의 해당 절
(가군, 절 \ref{modules-section-differential-operators})을 살펴볼 것을
권한다.

\begin{lemma}
\label{morphisms-lemma-affine-case-differential-operators}
$R \to A$를 환 준동형이라 하고, $f : X \to S$로 이에 대응하는 아핀
스킴의 사상을 나타내자. $\mathcal{F}$와 $\mathcal{G}$를
$\mathcal{O}_X$-가군이라 하자. $\mathcal{F}$가 준연접이면 미분 작용소를
그 대역 단면 위의 작용으로 보내는 사상
$$
\text{Diff}^k_{X/S}(\mathcal{F}, \mathcal{G}) \to
\text{Diff}^k_{A/R}(\Gamma(X, \mathcal{F}), \Gamma(X, \mathcal{G}))
$$
은 전단사이다.
\end{lemma}

\begin{proof}
$\mathcal{F} = \widetilde{M}$로 쓰자. 여기서 $A$-가군 $M$을 적절히 택한다.
$N = \Gamma(X, \mathcal{G})$로 두자. $D : M \to N$을 차수 $k$의 미분
작용소라 하자. 유일한 미분 작용소 $\mathcal{F} \to \mathcal{G}$로서
차수 $k$이고 대역 단면 위에서 $D$를 주는 것이
존재함을 보여야 한다.
$U = D(f) \subset X$를 표준 아핀 열린집합이라 하자. 그러면 국소화는
$\mathcal{F}(U) = M_f$이다. 대수, 보조정리
\ref{algebra-lemma-invert-system-differential-operators}에 의해 미분 작용소
$D$는 유일한 미분 작용소
$$
D_f : \mathcal{F}(U) = \widetilde{M}(U) = M_f \to N_f = \widetilde{N}(U)
$$
로 확장된다. 유일성에 의해 이 사상들 $D_f$는 서로 붙어서 층의 사상
$\widetilde{M} \to \widetilde{N}$을 $X$의 모든 표준 열린집합으로 된
기저 위에서 준다. 따라서 층, 절
\ref{sheaves-section-bases}의 내용에 의해 이 사상들과 일치하는 유일한
층의 사상 $\widetilde{D} : \widetilde{M} \to \widetilde{N}$을 얻는다.
$\widetilde{D}$는 표준 열린집합 위에서 차수 $k$의 미분 작용소들로
주어지므로 $\widetilde{D}$도 차수 $k$의 미분 작용소이다(작은 세부사항은
생략한다). 마지막으로 표준 $\mathcal{O}_X$-가군 사상
$c : \widetilde{N} \to \mathcal{G}$와 뒤에서 합성하면(스킴, 보조정리
\ref{schemes-lemma-compare-constructions}을 보라)
$c \circ \widetilde{D} : \mathcal{F} \to \mathcal{G}$를 얻으며, 이는
가군, 보조정리 \ref{modules-lemma-composition-differential-operators}에
의해 차수 $k$의 미분 작용소이다. 이로써 존재성이 증명된다. 유일성의
증명은 생략한다.
\end{proof}

\begin{lemma}
\label{morphisms-lemma-base-change-differential-operators}
$a : X \to S$와 $b : Y \to S$를 스킴의 사상이라 하자.
$\mathcal{F}$와 $\mathcal{F}'$를 준연접 $\mathcal{O}_X$-가군이라 하자.
$D : \mathcal{F} \to \mathcal{F}'$를 차수 $k$이고 $X/S$ 위에 있는 미분
작용소라 하자. $\mathcal{G}$를 준연접 $\mathcal{O}_Y$-가군이라 하자.
그러면 유일한 미분 작용소
$$
D' :
\text{pr}_1^*\mathcal{F} \otimes_{\mathcal{O}_{X \times_S Y}}
\text{pr}_2^*\mathcal{G}
\longrightarrow
\text{pr}_1^*\mathcal{F}' \otimes_{\mathcal{O}_{X \times_S Y}}
\text{pr}_2^*\mathcal{G}
$$
가 차수 $k$로 $X \times_S Y / Y$ 위에 존재하여
$
D'(s \otimes t) = D(s) \otimes t
$
를 만족한다. 여기서 $s$는 $\mathcal{F}$의 국소 단면이고 $t$는
$\mathcal{G}$의 국소 단면이다.
\end{lemma}

\begin{proof}
$X$, $Y$, $S$가 아핀인 경우 이는 보조정리
\ref{morphisms-lemma-affine-case-differential-operators}를 통해 대응하는 대수적
결과인 대수, 보조정리
\ref{algebra-lemma-base-change-differential-operators}에서 따른다.
일반적인 경우에는 아핀 열린집합들로 된 덮개를 사용하여(예를 들어 스킴,
보조정리 \ref{schemes-lemma-affine-covering-fibre-product}에서처럼)
$D'$를 대역적으로 구성한다. 세부사항은 생략한다.
\end{proof}

\begin{remark}
\label{morphisms-remark-base-change-differential-operators}
$a : X \to S$와 $b : Y \to S$를 스킴의 사상이라 하자.
$p : X \times_S Y \to X$와 $q : X \times_S Y \to Y$로 사영들을 나타내자.
이 주석에서 $\mathcal{O}_X$-가군 $\mathcal{F}$와 $\mathcal{O}_Y$-가군
$\mathcal{G}$가 주어지면
$$
\mathcal{F} \boxtimes \mathcal{G} =
p^*\mathcal{F} \otimes_{\mathcal{O}_{X \times_S Y}} q^*\mathcal{G}
$$
로 두자. $\mathcal{A}_{X/S}$로 대상이 준연접 $\mathcal{O}_X$-가군이고
사상이 $X/S$ 위의 유한 차수 미분 작용소인 가법 범주를 나타내자.
$\mathcal{A}_{Y/S}$와 $\mathcal{A}_{X \times_S Y/S}$도 마찬가지이다.
보조정리 \ref{morphisms-lemma-base-change-differential-operators}의 구성은 함자
$$
\boxtimes : 
\mathcal{A}_{X/S} \times \mathcal{A}_{Y/S} \longrightarrow
\mathcal{A}_{X \times_S Y/S},
\quad
(\mathcal{F}, \mathcal{G}) \longmapsto \mathcal{F} \boxtimes \mathcal{G}
$$
를 결정하며, 이는 사상에 대해 쌍선형이다. $X = \Spec(A)$,
$Y = \Spec(B)$, $S = \Spec(R)$이면 준연접 층과 가군의 동일시를 통해
이 함자는 대상 위에서 $(M, N) \mapsto M \otimes_R N$으로 주어지고,
사상 $(D, D') : (M, N) \to (M', N')$을
$D \otimes D' : M \otimes_R N \to M' \otimes_R N'$으로 보낸다.
\end{remark}
\section{매끄러운 사상}
\label{morphisms-section-smooth}

\noindent
$f : X \to Y$를 위상공간의 연속 사상이라 하자. 다음 조건을 고찰하자.
모든 $x \in X$에 대해 열린 근방 $x \in U \subset X$,
$f(x) \in V \subset Y$와 정수 $d$가 존재하여 $f(U) \subset V$이고
다음 가환 그림을 얻는다.
$$
\xymatrix{
X \ar[d] &
U \ar[l] \ar[d] \ar[r]_-\pi &
V \times \mathbf{R}^d \ar[ld] \\
Y & V \ar[l]
}
$$
여기서 $\pi$는 어떤 열린 부분집합 위로의 위상동형사상이다.
스킴의 매끄러운 사상은 스킴 범주에서 이러한 사상에 대응하는 개념이다.
보조정리 \ref{morphisms-lemma-smooth-locally-standard-smooth}와 보조정리
\ref{morphisms-lemma-smooth-etale-over-affine-space}를 보라.

\medskip\noindent
(어쩌면) 예상과 달리 매끄러운 환 준동형의 개념은 미분 가군만으로
정의되지 않는다. 즉 $R \to A$가 {\it 매끄러운 환 준동형}이라고 함은
$A$가 $R$ 위에서 유한 표시이고 $A$의 $R$ 위에서의 순진한 공변접 복합체가
차수 $0$에 놓인 사영 가군과 준동형인 경우임을 상기하자. 대수, 정의
\ref{algebra-definition-smooth}를 보라.

\begin{definition}
\label{morphisms-definition-smooth}
$f : X \to S$를 스킴의 사상이라 하자.
\begin{enumerate}
\item $f$가 {\it $x \in X$에서 매끄럽다}고 함은 아핀 열린 근방
$\Spec(A) = U \subset X$로서 $x$를 포함하는 것과 아핀 열린집합
$\Spec(R) = V \subset S$가 존재하여 $f(U) \subset V$이고 유도된 환 준동형
$R \to A$가 매끄러운 경우이다.
\item $f$가 $X$의 모든 점에서 매끄러우면 이를 {\it 매끄럽다}고 한다.
\item 아핀 스킴의 사상 $f : X \to S$에 대해 표준 매끄러운 환 준동형
$R \to R[x_1, \ldots, x_n]/(f_1, \ldots, f_c)$가 존재하여(대수, 정의
\ref{algebra-definition-standard-smooth}를 보라) $X \to S$가
$$
\Spec(R[x_1, \ldots, x_n]/(f_1, \ldots, f_c)) \to \Spec(R).
$$
와 동형이면 이를 {\it 표준 매끄럽다}고 한다.
\end{enumerate}
\end{definition}

\noindent
이 정의의 좋은 점 하나는 사상이 매끄러운 점들의 집합이 자동으로
열린집합이라는 것이다.

\medskip\noindent
정의에는 분리성이나 준콤팩트성 가정이 없다는 데 유의하자. 따라서
매끄러움은 원천에서 국소적인 문제이다. 정확한 결과는 다음과 같다.

\begin{lemma}
\label{morphisms-lemma-smooth-characterize}
$f : X \to S$를 스킴의 사상이라 하자. 다음은 동치이다.
\begin{enumerate}
\item 사상 $f$는 매끄럽다.
\item 모든 아핀 열린집합 $U \subset X$, $V \subset S$에 대해
$f(U) \subset V$이면 환 준동형
$\mathcal{O}_S(V) \to \mathcal{O}_X(U)$가 매끄럽다.
\item 열린 덮개 $S = \bigcup_{j \in J} V_j$와 열린 덮개
$f^{-1}(V_j) = \bigcup_{i \in I_j} U_i$가 존재하여 각 사상
$U_i \to V_j$, $j\in J, i\in I_j$가 매끄럽다.
\item 아핀 열린 덮개 $S = \bigcup_{j \in J} V_j$와 아핀 열린 덮개
$f^{-1}(V_j) = \bigcup_{i \in I_j} U_i$가 존재하여 환 준동형
$\mathcal{O}_S(V_j) \to \mathcal{O}_X(U_i)$가 모든
$j\in J, i\in I_j$에 대해 매끄럽다.
\end{enumerate}
또한 $f$가 매끄러우면 임의의 열린 부분스킴 $U \subset X$, $V \subset S$에
대해 $f(U) \subset V$일 때 제한 $f|_U : U \to V$도 매끄럽다.
\end{lemma}

\begin{proof}
``$R \to A$가 매끄럽다''라는 성질이 국소적임을 보이면 보조정리
\ref{morphisms-lemma-locally-P}에서 따른다. 정의
\ref{morphisms-definition-property-local}의 조건 (a), (b), (c)를 확인하자.
대수, 보조정리 \ref{algebra-lemma-base-change-smooth}에 의해 매끄러움은
기저변환으로 보존되므로 (a)가 성립한다. 대수, 보조정리
\ref{algebra-lemma-compose-smooth}에 의해 매끄러움은 합성으로 보존되고,
임의의 환 $R$에 대해 환 준동형 $R \to R_f$는 (표준) 매끄럽다. 따라서
(b)가 성립한다. 마지막으로 대수, 보조정리
\ref{algebra-lemma-locally-smooth}에 의해 성질 (c)가 참이다.
\end{proof}

\noindent
다음 보조정리는 매끄러운 사상을 올이 매끄러운 평탄 유한 표시 사상으로
특징짓는다. 체 위에서 매끄러운 스킴은 다형체, 절
\ref{varieties-section-smooth}에서 더 자세히 다룬다는 데 유의하자.

\begin{lemma}
\label{morphisms-lemma-smooth-flat-smooth-fibres}
$f : X \to S$를 스킴의 사상이라 하자. $f$가 평탄이고 국소적으로 유한
표시이며 모든 올 $X_s$가 매끄러우면 $f$는 매끄럽다.
\end{lemma}

\begin{proof}
대수, 보조정리 \ref{algebra-lemma-flat-fibre-smooth}에서 따른다.
\end{proof}

\begin{lemma}
\label{morphisms-lemma-composition-smooth}
매끄러운 두 사상의 합성은 매끄럽다.
\end{lemma}

\begin{proof}
보조정리 \ref{morphisms-lemma-smooth-characterize}의 증명에서 매끄러움이 환
준동형의 국소적 성질임을 보았다. 따라서 보조정리
\ref{morphisms-lemma-composition-type-P}와 매끄러움이 합성으로 보존되는 환
준동형의 성질이라는 사실을 결합하면 보조정리의 첫째 명제를 얻는다.
대수, 보조정리 \ref{algebra-lemma-compose-smooth}를 보라.
\end{proof}

\begin{lemma}
\label{morphisms-lemma-base-change-smooth}
매끄러운 사상의 기저변환은 매끄럽다.
\end{lemma}

\begin{proof}
보조정리 \ref{morphisms-lemma-smooth-characterize}의 증명에서 매끄러움이 환
준동형의 국소적 성질임을 보았다. 따라서 보조정리
\ref{morphisms-lemma-composition-type-P}와 매끄러움이 기저변환으로 보존되는 환
준동형의 성질이라는 사실을 결합하면 결과가 따른다. 대수, 보조정리
\ref{algebra-lemma-base-change-smooth}를 보라.
\end{proof}

\begin{lemma}
\label{morphisms-lemma-open-immersion-smooth}
모든 열린 몰입은 매끄럽다.
\end{lemma}

\begin{proof}
열린 몰입은 국소 동형사상이므로 참이다.
\end{proof}

\begin{lemma}
\label{morphisms-lemma-smooth-syntomic}
매끄러운 사상은 신토믹이다.
\end{lemma}

\begin{proof}
대수, 보조정리 \ref{algebra-lemma-smooth-syntomic}을 보라.
\end{proof}

\begin{lemma}
\label{morphisms-lemma-smooth-locally-finite-presentation}
매끄러운 사상은 국소적으로 유한 표시이다.
\end{lemma}

\begin{proof}
매끄러운 환 준동형은 정의상 유한 표시이므로 참이다.
\end{proof}

\begin{lemma}
\label{morphisms-lemma-smooth-flat}
매끄러운 사상은 평탄이다.
\end{lemma}

\begin{proof}
보조정리 \ref{morphisms-lemma-syntomic-flat}과 \ref{morphisms-lemma-smooth-syntomic}을 결합하라.
\end{proof}

\begin{lemma}
\label{morphisms-lemma-smooth-open}
매끄러운 사상은 보편 열림이다.
\end{lemma}

\begin{proof}
보조정리 \ref{morphisms-lemma-smooth-flat},
\ref{morphisms-lemma-smooth-locally-finite-presentation}, \ref{morphisms-lemma-fppf-open}을
결합하라. 또는 보조정리 \ref{morphisms-lemma-smooth-syntomic},
\ref{morphisms-lemma-syntomic-open}을 결합하라.
\end{proof}

\noindent
다음 보조정리는 국소적으로 모든 매끄러운 사상이 표준 매끄럽다고 말한다.
따라서 표준 매끄러운 사상을 매끄러운 사상의 {\it 국소 모형}으로 사용할
수 있다.

\begin{lemma}
\label{morphisms-lemma-smooth-locally-standard-smooth}
\begin{slogan}
매끄러운 사상은 국소적으로 표준 매끄럽다.
\end{slogan}
$f : X  \to S$를 스킴의 사상이라 하자. $x \in X$를 점이라 하자.
$V \subset S$를 $f(x)$의 아핀 열린 근방이라 하자. 다음은 동치이다.
\begin{enumerate}
\item 사상 $f$는 $x$에서 매끄럽다.
\item 아핀 열린집합 $U \subset X$가 존재하여 $x \in U$이고
$f(U) \subset V$이며, 유도된 사상 $f|_U : U \to V$는 표준 매끄럽다.
\end{enumerate}
\end{lemma}

\begin{proof}
정의들과 대수, 보조정리 \ref{algebra-lemma-standard-smooth},
\ref{algebra-lemma-smooth-syntomic}에서 따른다.
\end{proof}
\begin{lemma}
\label{morphisms-lemma-smooth-omega-finite-locally-free}
$f : X \to S$를 스킴의 사상이라 하자. $f$가 매끄럽다고 가정하자.
그러면 미분 가군 $\Omega_{X/S}$, 즉 $X$의 $S$ 위에서의 미분 가군은
유한 국소 자유이고
$$
\text{rank}_x(\Omega_{X/S}) = \dim_x(X_{f(x)})
$$
가 모든 $x \in X$에 대해 성립한다.
\end{lemma}

\begin{proof}
명제는 $X$와 $S$ 위에서 국소적이다. 위 보조정리
\ref{morphisms-lemma-smooth-locally-standard-smooth}에 의해 $f$가 아핀 스킴 사이의
표준 매끄러운 사상이라고 가정할 수 있다. 이 경우 결과는 대수, 보조정리
\ref{algebra-lemma-standard-smooth}에서 따른다(상대 대역 완전교차의 정의,
즉 대수, 정의
\ref{algebra-definition-relative-global-complete-intersection}도 보라).
\end{proof}

\noindent
보조정리 \ref{morphisms-lemma-smooth-omega-finite-locally-free}에 의해 다음 정의가
의미를 갖는다.

\begin{definition}
\label{morphisms-definition-smooth-relative-dimension}
$d \geq 0$을 정수라 하자. 스킴의 사상 $f : X \to S$가
{\it 상대차원 $d$로 매끄럽다}고 함은 $f$가 매끄럽고
$\Omega_{X/S}$가 상수 계수 $d$인 유한 국소 자유인 경우이다.
\end{definition}

\noindent
다시 말해 $f$는 매끄럽고 공집합 아닌 올들은 차원 $d$인 순수차원이다.
아래 보조정리 \ref{morphisms-lemma-smooth-at-point}에 의해 이는 다음을 요구하는 것과
같다. (a) $f$는 국소적으로 유한 표시이고, (b) $f$는 평탄이며, (c) 공집합
아닌 모든 올은 차원 $d$인 순수차원이고, (d) $\Omega_{X/S}$는 계수 $d$인
유한 국소 자유이다. 단순히 $f$가 평탄이고 유한 표시이며
$\Omega_{X/S}$가 계수 $d$인 유한 국소 자유라고 가정하는 것으로는
충분하지 않다. 반례는
$\Spec(\mathbf{F}_p[t]) \to \Spec(\mathbf{F}_p[t^p])$이다.

\medskip\noindent
점에서 매끄러움에 대한 미분적 판정법은 다음과 같다. 이 결과에는 여러
변형이 있고 언젠가 모두 유용할 수 있다. 필요할 때 여기에 추가하겠다.

\begin{lemma}
\label{morphisms-lemma-smooth-at-point}
$f : X \to S$를 스킴의 사상이라 하자. $x \in X$라 하고
$s = f(x)$로 두자. $f$가 국소적으로 유한 표시라고 가정하자.
다음은 동치이다.
\begin{enumerate}
\item 사상 $f$는 $x$에서 매끄럽다.
\item 국소환 준동형 $\mathcal{O}_{S, s} \to \mathcal{O}_{X, x}$는 평탄이고
$X_s \to \Spec(\kappa(s))$는 $x$에서 매끄럽다.
\item 국소환 준동형 $\mathcal{O}_{S, s} \to \mathcal{O}_{X, x}$는 평탄이고
$\mathcal{O}_{X, x}$-가군 $\Omega_{X/S, x}$는 많아야
$\dim_x(X_{f(x)})$개의 원소로 생성된다.
\item 국소환 준동형 $\mathcal{O}_{S, s} \to \mathcal{O}_{X, x}$는 평탄이고
$\kappa(x)$-벡터공간
$$
\Omega_{X_s/s, x} \otimes_{\mathcal{O}_{X_s, x}} \kappa(x) =
\Omega_{X/S, x} \otimes_{\mathcal{O}_{X, x}} \kappa(x)
$$
은 많아야 $\dim_x(X_{f(x)})$개의 원소로 생성된다.
\item 아핀 열린집합 $U \subset X$와 $V \subset S$가 존재하여
$x \in U$, $f(U) \subset V$이고 유도된 사상 $f|_U : U \to V$는
표준 매끄럽다.
\item 아핀 열린집합 $\Spec(A) = U \subset X$와
$\Spec(R) = V \subset S$가 존재하여 $x \in U$는
$\mathfrak q \subset A$에 대응하고 $f(U) \subset V$이며, 다음 표시가
존재한다.
$$
A = R[x_1, \ldots, x_n]/(f_1, \ldots, f_c)
$$
여기서
$$
g =
\det
\left(
\begin{matrix}
\partial f_1/\partial x_1 &
\partial f_2/\partial x_1 &
\ldots &
\partial f_c/\partial x_1 \\
\partial f_1/\partial x_2 &
\partial f_2/\partial x_2 &
\ldots &
\partial f_c/\partial x_2 \\
\ldots & \ldots & \ldots & \ldots \\
\partial f_1/\partial x_c &
\partial f_2/\partial x_c &
\ldots &
\partial f_c/\partial x_c
\end{matrix}
\right)
$$
의 상은 $A$의 원소이고 $\mathfrak q$에 속하지 않는다.
\end{enumerate}
\end{lemma}

\begin{proof}
$f$가 $x$에서 매끄러우면 보조정리
\ref{morphisms-lemma-smooth-locally-standard-smooth}에서 (5)가 성립함을 알 수 있고,
(6)은 (5)를 약간 약화한 판이다. 또한 $f$가 매끄러우면 환 준동형
$\mathcal{O}_{S, s} \to \mathcal{O}_{X, x}$는 평탄하고(보조정리
\ref{morphisms-lemma-smooth-flat}을 보라), $\Omega_{X/S}$는 계수가
$\dim_x(X_s)$와 같은 유한 국소 자유이다(보조정리
\ref{morphisms-lemma-smooth-omega-finite-locally-free}를 보라). 따라서 (1)은
(3)과 (4)를 함의한다. 보조정리 \ref{morphisms-lemma-base-change-smooth}에 의해
(1)은 (2)도 함의한다.

\medskip\noindent
보조정리 \ref{morphisms-lemma-base-change-differentials}에 의해 미분 가군
$\Omega_{X_s/s}$, 즉 올 $X_s$의 $\kappa(s)$ 위에서의 미분 가군은
미분 가군 $\Omega_{X/S}$, 즉 $X$의 $S$ 위에서의 미분 가군의 당김이다.
따라서 보조정리의 (4)에 표시된
등식이 성립한다. 보조정리 \ref{morphisms-lemma-finite-type-differentials}에 의해
이 가군들은 유한형이다. 그러므로 가군 $\Omega_{X/S, x}$와
$\Omega_{X_s/s, x}$의 최소 생성원 수는 같고, 나카야마 보조정리(대수,
보조정리 \ref{algebra-lemma-NAK})에 의해 이 $\kappa(x)$-벡터공간의
차원과 같다. 특히 (3)과 (4)가 동치임을 보인다.

\medskip\noindent
대수, 보조정리 \ref{algebra-lemma-flat-fibre-smooth}는 (2)가 (1)을
함의함을 보인다. 대수, 보조정리
\ref{algebra-lemma-characterize-smooth-over-field}는 (3)과 (4)가 (2)를
함의함을 보인다. 마지막으로 예를 들어 대수, 예
\ref{algebra-example-make-standard-smooth}에 의해 (6)은 (5)를 함의하고,
대수, 보조정리 \ref{algebra-lemma-standard-smooth}에 의해 (5)는 (1)을
함의한다.
\end{proof}
\begin{lemma}
\label{morphisms-lemma-set-points-where-fibres-smooth}
스킴들의 카르테시안 그림
$$
\xymatrix{
X' \ar[r]_{g'} \ar[d]_{f'} & X \ar[d]^f \\
S' \ar[r]^g & S
}
$$
이 주어졌다고 하자. $W \subset X$, 각각 $W' \subset X'$를 $f$, 각각
$f'$가 매끄러운 점들의 열린 부분스킴이라 하자. 다음 중 하나가 성립하면
$W' = (g')^{-1}(W)$이다.
\begin{enumerate}
\item $f$는 평탄이고 국소적으로 유한 표시이다.
\item $f$는 국소적으로 유한 표시이고 $g$는 평탄이다.
\end{enumerate}
\end{lemma}

\begin{proof}
먼저 $f$가 국소적으로 유한형이라고 가정하자. 집합
$$
T = \{x \in X \mid X_{f(x)}\text{ is smooth over }\kappa(f(x))\text{ at }x\}
$$
과 대응하는 집합 $T' \subset X'$를 $f'$에 대해 고찰하자. 그러면
$T' = (g')^{-1}(T)$라고 주장한다. 실제로 $s' \in S'$를 점이라 하고
$s = g(s')$로 두자. 그러면
$$
X'_{s'} =
\Spec(\kappa(s')) \times_{\Spec(\kappa(s))} X_s
$$
이다. 다시 말해 기저변환의 올은 원래 올의 기저변환이다. 따라서 주장은
대수, 보조정리 \ref{algebra-lemma-smooth-field-change-local}과 동치이다.

\medskip\noindent
따라서 (1)의 경우 $T$는 보조정리 \ref{morphisms-lemma-smooth-at-point}에 의해
$f$가 매끄러운 점들의 (열린) 집합이므로 (1)이 따른다.

\medskip\noindent
(2)의 경우 $x' \in W'$라 하자. 그러면 $g'$는 $x'$에서 평탄하고
(보조정리 \ref{morphisms-lemma-base-change-module-flat}), $g \circ f'$도 $x'$에서
평탄하다(보조정리 \ref{morphisms-lemma-composition-module-flat}). 따라서 보조정리
\ref{morphisms-lemma-flat-permanence}에 의해 $f$는 $x = g'(x')$에서 평탄이다.
한편 $x' \in T'$이므로(보조정리 \ref{morphisms-lemma-base-change-smooth})
$x \in T$임을 알 수 있다. 따라서 보조정리 \ref{morphisms-lemma-smooth-at-point}에
의해 $f$는 $x$에서 매끄럽다.
\end{proof}

\noindent
다음은 매끄러운 환 준동형의 순진한 공변접 복합체에서 $H^{-1}$의
소멸을 실제로 사용하는 보조정리이다.

\begin{lemma}
\label{morphisms-lemma-triangle-differentials-smooth}
$f : X \to Y$, $g : Y \to S$를 스킴의 사상이라 하자. $f$가 매끄럽다고
가정하자. 그러면
$$
0 \to f^*\Omega_{Y/S} \to \Omega_{X/S} \to \Omega_{X/Y} \to 0
$$
(보조정리 \ref{morphisms-lemma-triangle-differentials}를 보라)은 짧은 완전열이다.
\end{lemma}

\begin{proof}
이 보조정리의 대수적 판은 다음과 같다. 환 준동형
$A \to B \to C$가 주어지고 $B \to C$가 매끄러우면 대수, 보조정리
\ref{algebra-lemma-exact-sequence-differentials}의 열
$$
0 \to C \otimes_B \Omega_{B/A} \to \Omega_{C/A} \to \Omega_{C/B} \to 0
$$
은 완전하다. 이것이 대수, 보조정리
\ref{algebra-lemma-triangle-differentials-smooth}이다.
\end{proof}

\begin{lemma}
\label{morphisms-lemma-differentials-relative-immersion-smooth}
$i : Z \to X$를 $S$ 위의 스킴들의 몰입이라 하자. $Z$가 $S$ 위에서
매끄럽다고 가정하자. 그러면 보조정리
\ref{morphisms-lemma-differentials-relative-immersion}의 표준 완전열
$$
0 \to \mathcal{C}_{Z/X} \to i^*\Omega_{X/S} \to \Omega_{Z/S} \to 0
$$
은 짧은 완전열이다.
\end{lemma}

\begin{proof}
이 보조정리의 대수적 판은 다음과 같다. 환 준동형
$A \to B \to C$가 주어지고 $A \to C$가 매끄러우며 $B \to C$가
핵 $J$를 갖는 전사이면 대수, 보조정리
\ref{algebra-lemma-differential-seq}의 열
$$
0 \to J/J^2 \to C \otimes_B \Omega_{B/A} \to \Omega_{C/A} \to 0
$$
은 완전하다. 이것이 대수, 보조정리
\ref{algebra-lemma-differential-seq-smooth}이다.
\end{proof}

\begin{lemma}
\label{morphisms-lemma-two-immersions-smooth}
스킴들의 가환 그림
$$
\xymatrix{
Z \ar[r]_i \ar[rd]_j & X \ar[d] \\
& Y
}
$$
이 주어지고 $i$와 $j$가 몰입이며 $X \to Y$가 매끄럽다고 하자.
그러면 보조정리 \ref{morphisms-lemma-two-immersions}의 표준 완전열
$$
0 \to  \mathcal{C}_{Z/Y} \to \mathcal{C}_{Z/X} \to i^*\Omega_{X/Y} \to 0
$$
은 완전하다.
\end{lemma}

\begin{proof}
이 보조정리의 대수적 판은 다음과 같다. 환 준동형
$A \to B \to C$가 주어지고 $A \to C$가 전사이며 $A \to B$가
매끄러우면 대수, 보조정리 \ref{algebra-lemma-application-NL}의 열
$$
0 \to I/I^2 \to J/J^2 \to C \otimes_B \Omega_{B/A} \to 0
$$
은 완전하다. 이것이 대수, 보조정리
\ref{algebra-lemma-application-NL-smooth}이다.
\end{proof}
\begin{lemma}
\label{morphisms-lemma-smooth-permanence}
스킴의 사상들의 가환 그림
$$
\xymatrix{
X \ar[rr]_f \ar[rd]_p & &
Y \ar[dl]^q \\
& S
}
$$
이 주어졌다고 하자. 다음을 가정하자.
\begin{enumerate}
\item $f$는 전사이고 매끄럽다.
\item $p$는 매끄럽다.
\item $q$는 국소적으로 유한 표시이다\footnote{사실 이는 (1)과 (2)에서
따른다. 하강, 보조정리
\ref{descent-lemma-flat-finitely-presented-permanence}를 보라. 더 나아가
$f$가 전사이고 평탄이며 국소적으로 유한 표시라고 가정하는 것으로
충분하다. 하강, 보조정리 \ref{descent-lemma-smooth-permanence}를 보라.}.
\end{enumerate}
그러면 $q$는 매끄럽다.
\end{lemma}

\begin{proof}
보조정리 \ref{morphisms-lemma-flat-permanence}에 의해 $q$는 평탄이다.
점 $y \in Y$를 택하고, 점 $x \in X$를 택하여 $y$로 사상되게 하자.
$f$가 상대차원 $a$인 점이 $x$이고 $p$가 상대차원 $b$인 점도 $x$라고 하자.
보조정리 \ref{morphisms-lemma-smooth-omega-finite-locally-free}에 의해 이는
$\Omega_{X/S, x}$가 계수 $b$인 자유 가군이고 $\Omega_{X/Y, x}$가 계수
$a$인 자유 가군이라는 뜻이다. 보조정리
\ref{morphisms-lemma-triangle-differentials-smooth}의 짧은 완전열에 의해
$(f^*\Omega_{Y/S})_x$는 계수 $b - a$인 자유 가군이다. 나카야마
보조정리에 의해 $\Omega_{Y/S, y}$는 $b - a$개의 원소로 생성된다.
또한 보조정리 \ref{morphisms-lemma-dimension-fibre-at-a-point-additive}에 의해
$\dim_y(Y_s) = b - a$이다. 따라서 보조정리 \ref{morphisms-lemma-smooth-at-point}의
(2)에 의해 $Y \to S$는 $y$에서 매끄럽다.
\end{proof}

\noindent
다음 보조정리의 상황에서 $\sigma$의 상은 $X$ 위에서 국소적으로 정칙열에
의해 잘린다. 인자, 보조정리
\ref{divisors-lemma-section-smooth-regular-immersion}을 보라.

\begin{lemma}
\label{morphisms-lemma-section-smooth-morphism}
$f : X \to S$를 스킴의 사상이라 하자. $\sigma : S \to X$를 $f$의
단면이라 하자. $s \in S$를 $f$가 $x = \sigma(s)$에서 매끄럽도록 하는
점이라 하자. 그러면 아핀 열린 근방 $\Spec(A) = U \subset S$로서 $s$를
포함하는 것과 아핀 열린 근방 $\Spec(B) = V \subset X$로서 $x$를
포함하는 것이 존재하여 다음을 만족한다.
\begin{enumerate}
\item $f(V) \subset U$이고 $\sigma(U) \subset V$이다.
\item $I = \Ker(\sigma^\# : B \to A)$로 두면 가군 $I/I^2$는 자유
$A$-가군이다.
\item $B^\wedge \cong A[[x_1, \ldots, x_d]]$가 $A$-대수의 동형사상이고,
$B^\wedge$는 $B$를 $I$에 관해 완비화한 것이다.
\end{enumerate}
\end{lemma}

\begin{proof}
아핀 열린집합 $U \subset S$를 택하여 $s$를 포함하게 하자.
아핀 열린집합 $V \subset f^{-1}(U)$를 택하여 $x$를 포함하게 하자.
아핀 열린집합 $U' \subset \sigma^{-1}(V)$를 택하여 $s$를 포함하게 하자.
$V' = f^{-1}(U') \cap V$는 올곱 $V' = U' \times_U V$와 같으므로 아핀이다.
그러면 $U'$와 $V'$는 (1)을 만족한다. $U' = \Spec(A')$와
$V' = \Spec(B')$로 쓰자. 대수, 보조정리
\ref{algebra-lemma-section-smooth}에 의해 가군 $I'/(I')^2$는 유한 국소
자유 $A'$-가군이다. 따라서 $U'$를 더 작은 아핀 열린집합
$U'' \subset U'$로, $V'$를 $V'' = V' \cap f^{-1}(U'')$로 바꾸면
$I''/(I'')^2$가 자유인 상황, 즉 (2)를 얻는다. 이 경우 (3)도 대수,
보조정리 \ref{algebra-lemma-section-smooth}에 의해 성립한다.
\end{proof}

\noindent
스킴 $X$의 점 $x$에서의 차원(성질, 정의
\ref{properties-definition-dimension})은 위상공간으로서 $X$의 점
$x$에서의 차원일 뿐이다. 위상, 정의 \ref{topology-definition-Krull}을 보라.
이는 일반적으로 국소환 $\mathcal{O}_{X,x}$의 차원이 아니다.

\begin{lemma}
\label{morphisms-lemma-smoothness-dimension}
$f : X \to Y$를 국소 뇌터 스킴의 매끄러운 사상이라 하자. 모든 점
$x$가 $X$에 있고 그 상 $y$가 $Y$에 있으면
$$
\dim_x(X) = \dim_y(Y) + \dim_x(X_y),
$$
이다. 여기서 $X_y$는 $y$ 위의 올을 나타낸다.
\end{lemma}

\begin{proof}
$X$를 $x$의 열린 근방으로 바꾸면, 자연수 $d$가 존재하여 $X \to Y$의
모든 올이 모든 점에서 차원 $d$이다. 보조정리
\ref{morphisms-lemma-smooth-omega-finite-locally-free}를 보라. 그러면 $f$는
평탄이고(보조정리 \ref{morphisms-lemma-smooth-flat}), 국소적으로 유한형이며
(보조정리 \ref{morphisms-lemma-smooth-locally-finite-presentation}), 상대차원
$d$이다. 따라서 결과는 보조정리 \ref{morphisms-lemma-rel-dimension-dimension}에서
따른다.
\end{proof}
\section{비분기 사상}
\label{morphisms-section-unramified}

\noindent
(어쩌면) 더 흥미로운 \'etale 사상류보다 앞서 비분기 사상을 간단히
논의한다. 환 준동형 $R \to A$가 유한형이고 $\Omega_{A/R} = 0$이면 이를
{\it 비분기}라고 함을 상기하자(이는 \cite{Henselian}의 정의이다).
환 준동형 $R \to A$가 유한 표시이고 $\Omega_{A/R} = 0$이면 이를
{\it G-비분기}라고 한다(이는 \cite{EGA}의 정의이다). 대수, 정의
\ref{algebra-definition-unramified}를 보라.

\begin{definition}
\label{morphisms-definition-unramified}
$f : X \to S$를 스킴의 사상이라 하자.
\begin{enumerate}
\item $f$가 {\it $x \in X$에서 비분기}라고 함은 아핀 열린 근방
$\Spec(A) = U \subset X$로서 $x$를 포함하는 것과 아핀 열린집합
$\Spec(R) = V \subset S$가 존재하여 $f(U) \subset V$이고 유도된 환 준동형
$R \to A$가 비분기인 경우이다.
\item $f$가 {\it $x \in X$에서 G-비분기}라고 함은 아핀 열린 근방
$\Spec(A) = U \subset X$로서 $x$를 포함하는 것과 아핀 열린집합
$\Spec(R) = V \subset S$가 존재하여 $f(U) \subset V$이고 유도된 환 준동형
$R \to A$가 G-비분기인 경우이다.
\item $f$가 $X$의 모든 점에서 비분기이면 이를 {\it 비분기}라고 한다.
\item $f$가 $X$의 모든 점에서 G-비분기이면 이를 {\it G-비분기}라고 한다.
\end{enumerate}
\end{definition}

\noindent
G-비분기 사상은 비분기임에 유의하자. 따라서 비분기 사상에 대한 모든
결과는 G-비분기 사상에 대한 대응하는 결과를 함의한다. 또한 $S$가
국소 뇌터이면 G-비분기와 비분기 사상 사이에 차이가 없다. 보조정리
\ref{morphisms-lemma-noetherian-unramified}를 보라. 이 정의의 좋은 점 하나는
사상이 비분기인(각각 G-비분기인) 점들의 집합이 자동으로 열린다는 것이다.

\begin{lemma}
\label{morphisms-lemma-unramified-omega-zero}
$f : X \to S$를 스킴의 사상이라 하자. 그러면
\begin{enumerate}
\item $f$가 비분기인 것은 $f$가 국소적으로 유한형이고
$\Omega_{X/S} = 0$인 것과 동치이다.
\item $f$가 G-비분기인 것은 $f$가 국소적으로 유한 표시이고
$\Omega_{X/S} = 0$인 것과 동치이다.
\end{enumerate}
\end{lemma}

\begin{proof}
정의상 환 준동형 $R \to A$가 비분기(각각 G-비분기)인 것은 유한형
(각각 유한 표시)이고 $\Omega_{A/R} = 0$인 것과 동치이다. 따라서
보조정리는 정의들과 보조정리 \ref{morphisms-lemma-differentials-affine}에서
곧바로 따른다.
\end{proof}

\noindent
정의에는 분리성이나 준콤팩트성 가정이 없다는 데 유의하자. 따라서
비분기 여부는 원천에서 국소적인 문제이다. 정확한 결과는 다음과 같다.

\begin{lemma}
\label{morphisms-lemma-unramified-characterize}
$f : X \to S$를 스킴의 사상이라 하자. 다음은 동치이다.
\begin{enumerate}
\item 사상 $f$는 비분기이다(각각 G-비분기이다).
\item 모든 아핀 열린집합 $U \subset X$, $V \subset S$에 대해
$f(U) \subset V$이면 환 준동형
$\mathcal{O}_S(V) \to \mathcal{O}_X(U)$가 비분기이다(각각 G-비분기이다).
\item 열린 덮개 $S = \bigcup_{j \in J} V_j$와 열린 덮개
$f^{-1}(V_j) = \bigcup_{i \in I_j} U_i$가 존재하여 각 사상
$U_i \to V_j$, $j\in J, i\in I_j$가 비분기이다(각각 G-비분기이다).
\item 아핀 열린 덮개 $S = \bigcup_{j \in J} V_j$와 아핀 열린 덮개
$f^{-1}(V_j) = \bigcup_{i \in I_j} U_i$가 존재하여 환 준동형
$\mathcal{O}_S(V_j) \to \mathcal{O}_X(U_i)$가 모든
$j\in J, i\in I_j$에 대해 비분기이다(각각 G-비분기이다).
\end{enumerate}
또한 $f$가 비분기이면(각각 G-비분기이면) 임의의 열린 부분스킴
$U \subset X$, $V \subset S$에 대해 $f(U) \subset V$일 때 제한
$f|_U : U \to V$도 비분기이다(각각 G-비분기이다).
\end{lemma}

\begin{proof}
``$R \to A$가 비분기이다''라는 성질이 국소적임을 보이면 보조정리
\ref{morphisms-lemma-locally-P}에서 따른다. 정의
\ref{morphisms-definition-property-local}의 조건 (a), (b), (c)를 확인하자.
이 성질들은 대수, 보조정리 \ref{algebra-lemma-unramified}에서 증명된다.
\end{proof}

\begin{lemma}
\label{morphisms-lemma-composition-unramified}
비분기인 두 사상의 합성은 비분기이다. G-비분기 사상에도 같은 결론이
성립한다.
\end{lemma}

\begin{proof}
보조정리 \ref{morphisms-lemma-unramified-characterize}의 증명은 비분기성(각각
G-비분기성)이 환 준동형의 국소적 성질임을 보인다. 따라서 보조정리
\ref{morphisms-lemma-composition-type-P}와 비분기성(각각 G-비분기성)이 합성으로
보존되는 환 준동형의 성질이라는 사실을 결합하면 보조정리의 첫째 명제가
따른다. 대수, 보조정리 \ref{algebra-lemma-unramified}를 보라.
\end{proof}

\begin{lemma}
\label{morphisms-lemma-base-change-unramified}
비분기 사상의 기저변환은 비분기이다. G-비분기 사상에도 같은 결론이
성립한다.
\end{lemma}

\begin{proof}
보조정리 \ref{morphisms-lemma-unramified-characterize}의 증명은 비분기성(각각
G-비분기성)이 환 준동형의 국소적 성질임을 보인다. 따라서 보조정리
\ref{morphisms-lemma-base-change-type-P}와 비분기성(각각 G-비분기성)이
기저변환으로 보존되는 환 준동형의 성질이라는 사실을 결합하면 결과가
따른다. 대수, 보조정리 \ref{algebra-lemma-unramified}를 보라.
\end{proof}

\begin{lemma}
\label{morphisms-lemma-noetherian-unramified}
$f : X \to S$를 스킴의 사상이라 하자. $S$가 국소 뇌터라고 가정하자.
그러면 $f$가 비분기인 것은 $f$가 G-비분기인 것과 동치이다.
\end{lemma}

\begin{proof}
정의들과 보조정리
\ref{morphisms-lemma-noetherian-finite-type-finite-presentation}에서 따른다.
\end{proof}


\begin{lemma}
\label{morphisms-lemma-open-immersion-unramified}
모든 열린 몰입은 G-비분기이다.
\end{lemma}

\begin{proof}
열린 몰입은 국소 동형사상이므로 참이다.
\end{proof}

\begin{lemma}
\label{morphisms-lemma-closed-immersion-unramified}
닫힌 몰입 $i : Z \to X$는 비분기이다. 이것이 G-비분기인 것은 대응하는
준연접 아이디얼 층
$\mathcal{I} = \Ker(\mathcal{O}_X \to i_*\mathcal{O}_Z)$가
($\mathcal{O}_X$-가군으로서) 유한형인 것과 동치이다.
\end{lemma}

\begin{proof}
보조정리 \ref{morphisms-lemma-closed-immersion-finite-presentation}과 대수, 보조정리
\ref{algebra-lemma-unramified}에서 따른다.
\end{proof}

\begin{lemma}
\label{morphisms-lemma-unramified-locally-finite-type}
비분기 사상은 국소적으로 유한형이다. G-비분기 사상은 국소적으로 유한
표시이다.
\end{lemma}

\begin{proof}
비분기 환 준동형은 정의상 유한형이다. G-비분기 환 준동형은 정의상 유한
표시이다.
\end{proof}

\begin{lemma}
\label{morphisms-lemma-unramified-quasi-finite}
$f : X \to S$를 스킴의 사상이라 하자. $f$가 $x$에서 비분기이면 $f$는
$x$에서 준유한이다. 특히 비분기 사상은 국소적으로 준유한이다.
\end{lemma}

\begin{proof}
대수, 보조정리 \ref{algebra-lemma-unramified-quasi-finite}를 보라.
\end{proof}

\begin{lemma}
\label{morphisms-lemma-unramified-over-field}
비분기 사상의 올에 대해 다음이 성립한다.
\begin{enumerate}
\item $X$를 체 $k$ 위의 스킴이라 하자. 구조 사상
$X \to \Spec(k)$가 비분기인 것은 $X$가 $k$의 유한 분리 체확장의
스펙트럼들의 서로소 합집합인 것과 동치이다.
\item $f : X \to S$가 비분기 사상이면 모든 $s \in S$에 대해 올 $X_s$는
$\kappa(s)$의 유한 분리 체확장의 스펙트럼들의 서로소 합집합이다.
\end{enumerate}
\end{lemma}

\begin{proof}
(2)는 (1)과 보조정리 \ref{morphisms-lemma-base-change-unramified}에서 따른다.
(1)을 증명하자. 먼저 대수, 보조정리
\ref{algebra-lemma-characterize-unramified}를 사용한다. 이 보조정리에
따르면 $X$가 $k$의 유한 분리 체확장의 스펙트럼들의 서로소 합집합이면
$X \to \Spec(k)$는 비분기이다. 역으로 $X \to \Spec(k)$가 비분기라고
가정하자. 대수, 보조정리 \ref{algebra-lemma-unramified-at-prime}에 의해
모든 $x \in X$에 대해 잉여체 확장 $\kappa(x)/k$는 유한 분리이다.
$X \to \Spec(k)$는 국소적으로 준유한이므로(보조정리
\ref{morphisms-lemma-unramified-quasi-finite}) $X$의 모든 점은 고립된 닫힌점이다.
보조정리 \ref{morphisms-lemma-quasi-finite-at-point-characterize}를 보라. 따라서
$X$는 이산공간이고, 특히 그 국소환들의 스펙트럼의 서로소 합집합이다.
대수, 보조정리 \ref{algebra-lemma-unramified-at-prime}에 의해 이 국소환들은
체이고, 결론을 얻는다.
\end{proof}
\noindent
다음 보조정리는 비분기 사상을 올들이 비분기인 국소 유한형 사상으로
특징짓는다.

\begin{lemma}
\label{morphisms-lemma-unramified-etale-fibres}
$f : X \to S$를 스킴의 사상이라 하자.
\begin{enumerate}
\item $f$가 비분기이면 임의의 $x \in X$에 대해 체확장
$\kappa(x)/\kappa(f(x))$는 유한 분리이다.
\item $f$가 국소적으로 유한형이고 모든 $s \in S$에 대해 올 $X_s$가
$\kappa(s)$의 유한 분리 체확장의 스펙트럼들의 서로소 합집합이면 $f$는
비분기이다.
\item $f$가 국소적으로 유한 표시이고 모든 $s \in S$에 대해 올 $X_s$가
$\kappa(s)$의 유한 분리 체확장의 스펙트럼들의 서로소 합집합이면 $f$는
G-비분기이다.
\end{enumerate}
\end{lemma}

\begin{proof}
대수, 보조정리 \ref{algebra-lemma-unramified-at-prime}과
\ref{algebra-lemma-characterize-unramified}에서 따른다.
\end{proof}

\noindent
대각 사상을 사용한 비분기 사상의 특징화는 다음과 같다.

\begin{lemma}
\label{morphisms-lemma-diagonal-unramified-morphism}
$f : X \to S$를 사상이라 하자.
\begin{enumerate}
\item $f$가 비분기이면 대각 사상 $\Delta : X \to X \times_S X$는
열린 몰입이다.
\item $f$가 국소적으로 유한형이고 $\Delta$가 열린 몰입이면 $f$는
비분기이다.
\item $f$가 국소적으로 유한 표시이고 $\Delta$가 열린 몰입이면 $f$는
G-비분기이다.
\end{enumerate}
\end{lemma}

\begin{proof}
첫째 명제는 대수, 보조정리
\ref{algebra-lemma-diagonal-unramified}에서 따른다. 둘째 명제는
$\Omega_{X/S}$가 대각 사상의 코노멀 층이라는 사실(보조정리
\ref{morphisms-lemma-differentials-diagonal})과 $\Delta$가 열린 몰입이면 이 층이
분명 영이라는 사실에서 따른다.
\end{proof}

\begin{lemma}
\label{morphisms-lemma-unramified-at-point}
$f : X \to S$를 스킴의 사상이라 하자. $x \in X$라 하고
$s = f(x)$로 두자. $f$가 국소적으로 유한형이라고(각각 국소적으로 유한
표시라고) 가정하자. 다음은 동치이다.
\begin{enumerate}
\item 사상 $f$는 $x$에서 비분기이다(각각 G-비분기이다).
\item 올 $X_s$는 $\kappa(s)$ 위에서 $x$에서 비분기이다.
\item $\mathcal{O}_{X, x}$-가군 $\Omega_{X/S, x}$는 영이다.
\item $\mathcal{O}_{X_s, x}$-가군 $\Omega_{X_s/s, x}$는 영이다.
\item $\kappa(x)$-벡터공간
$$
\Omega_{X_s/s, x} \otimes_{\mathcal{O}_{X_s, x}} \kappa(x) =
\Omega_{X/S, x} \otimes_{\mathcal{O}_{X, x}} \kappa(x)
$$
은 영이다.
\item $\mathfrak m_s\mathcal{O}_{X, x} = \mathfrak m_x$이고 체확장
$\kappa(x)/\kappa(s)$는 유한 분리이다.
\end{enumerate}
\end{lemma}

\begin{proof}
$f$가 $x$에서 비분기이면 정의들과 절
\ref{morphisms-section-sheaf-differentials}의 미분 가군에 관한 결과에 의해
$\Omega_{X/S} = 0$이 $x$의 한 근방에서 성립한다. 따라서 (1)은 (3)과
(5)의 오른쪽 벡터공간의 소멸을 함의한다. 또한 보조정리
\ref{morphisms-lemma-base-change-differentials}에 의해 미분 가군
$\Omega_{X_s/s}$, 즉 올 $X_s$의 $\kappa(s)$ 위에서의 미분 가군은
미분 가군 $\Omega_{X/S}$, 즉 $X$의 $S$ 위에서의 미분 가군의 당김이므로
(1)은 (2)를 함의한다. 이 미분 가군에 관한 사실은 (4)에 표시된
벡터공간의 등식도 함의한다. 보조정리
\ref{morphisms-lemma-finite-type-differentials}에 의해 가군
$\Omega_{X/S, x}$와 $\Omega_{X_s/s, x}$는 유한형이다. 따라서 나카야마
보조정리(대수, 보조정리 \ref{algebra-lemma-NAK})에 의해 가군
$\Omega_{X/S, x}$와 $\Omega_{X_s/s, x}$가 영인 것은 (4)의 대응하는
$\kappa(x)$-벡터공간이 영인 것과 동치이다. 특히 (3), (4), (5)가
동치이다. $\Omega_{X/S}$의 지지는 $X$에서 닫혀 있다. 가군, 보조정리
\ref{modules-lemma-support-finite-type-closed}를 보라. 가정 (3)은 $x$가
지지에 속하지 않음을 함의한다. 따라서 $\Omega_{X/S}$는 $x$의 한
근방에서 영이고, 이는 (1)을 함의한다. $X_s \to s$에 적용한 (1)과
(3)의 동치는 (2)와 (4)의 동치를 함의한다. 이로써 (1)--(5)가
동치임을 보았다.

\medskip\noindent
또는 대수, 보조정리 \ref{algebra-lemma-unramified}를 사용하면
(1)--(5)의 동치를 더 직접 볼 수 있다.

\medskip\noindent
(1)과 (6)의 동치는 보조정리 \ref{morphisms-lemma-unramified-etale-fibres}에서
따른다. 대수, 보조정리 \ref{algebra-lemma-unramified-at-prime}과
\ref{algebra-lemma-characterize-unramified}에서도 더 직접적으로 따른다.
\end{proof}

\begin{lemma}
\label{morphisms-lemma-set-points-where-fibres-unramified}
$f : X \to S$를 스킴의 사상이라 하자. $f$가 국소적으로 유한형이라고
가정하자. 열린집합
\begin{align*}
T
& =
\{x \in X \mid X_{f(x)}\text{ is unramified over }\kappa(f(x))\text{ at }x\} \\
& =
\{x \in X \mid X\text{ is unramified over }S\text{ at }x\}
\end{align*}
의 구성은 임의의 기저변환과 가환한다. 즉 임의의 사상
$g : S' \to S$에 대해 $f' : X' \to S'$를 $f$의 기저변환이라 하고 사영
$g' : X' \to X$를 고찰하자. 그러면 대응하는 집합 $T'$는 사상 $f'$에 대해
$T' = (g')^{-1}(T)$와 같다. $f$가 국소적으로 유한 표시라고 가정하면,
$f$가 G-비분기인 점들의 열린집합에도 같은 결론이 성립한다.
\end{lemma}

\begin{proof}
$s' \in S'$를 점이라 하고 $s = g(s')$로 두자. 그러면
$$
X'_{s'} =
\Spec(\kappa(s')) \times_{\Spec(\kappa(s))} X_s
$$
이다. 다시 말해 기저변환의 올은 원래 올의 기저변환이다. 특히
$$
\Omega_{X_s/s, x} \otimes_{\mathcal{O}_{X_s, x}} \kappa(x')
=
\Omega_{X'_{s'}/s', x'} \otimes_{\mathcal{O}_{X'_{s'}, x'}} \kappa(x')
$$
이다. 보조정리 \ref{morphisms-lemma-base-change-differentials}를 보라. 따라서
보조정리 \ref{morphisms-lemma-unramified-at-point}에 의해 $x' \in T'$인 것은
$x \in T$인 것과 동치이다. 이 경우 $T$는 보조정리
\ref{morphisms-lemma-unramified-at-point}에 따라 $f$가 G-비분기인 점들의 (열린)
집합이므로 둘째 부분은 첫째 부분에서 따른다.
\end{proof}

\begin{lemma}
\label{morphisms-lemma-unramified-permanence}
$f : X \to Y$를 $S$ 위의 스킴들의 사상이라 하자.
\begin{enumerate}
\item $X$가 $S$ 위에서 비분기이면 $f$는 비분기이다.
\item $X$가 $S$ 위에서 G-비분기이고 $Y$가 $S$ 위에서 국소적으로
유한형이면 $f$는 G-비분기이다.
\end{enumerate}
\end{lemma}

\begin{proof}
$X$가 $S$ 위에서 비분기라고 가정하자. 보조정리
\ref{morphisms-lemma-permanence-finite-type}에 의해 $f$는 국소적으로 유한형이다.
가정에 의해 $\Omega_{X/S} = 0$이다. 따라서 보조정리
\ref{morphisms-lemma-triangle-differentials}에 의해 $\Omega_{X/Y} = 0$이다.
그러므로 $f$는 비분기이다. $X$가 $S$ 위에서 G-비분기이고 $Y$가
$S$ 위에서 국소적으로 유한형이면 보조정리
\ref{morphisms-lemma-finite-presentation-permanence}에 의해 $f$가 국소적으로 유한
표시임을 알고, $f$가 G-비분기라고 결론 내린다.
\end{proof}

\begin{lemma}
\label{morphisms-lemma-value-at-one-point}
$S$를 스킴이라 하자. $X$, $Y$를 $S$ 위의 스킴이라 하자.
$f, g : X \to Y$를 $S$ 위의 사상이라 하고 $x \in X$라 하자.
다음을 가정하자.
\begin{enumerate}
\item 구조 사상 $Y \to S$는 비분기이다.
\item $f(x) = g(x)$가 $Y$에서 성립하고, 이를 $y = f(x) = g(x)$라 쓰자.
\item 유도된 사상 $f^\sharp, g^\sharp : \kappa(y) \to \kappa(x)$는 같다.
\end{enumerate}
그러면 $x$의 열린 근방이 $X$ 안에 존재하여 그 위에서 $f$와 $g$가
같다.
\end{lemma}

\begin{proof}
사상 $(f, g) : X \to Y \times_S Y$를 고찰하자. 가정 (1)과 보조정리
\ref{morphisms-lemma-diagonal-unramified-morphism}에 의해 $\Delta_{Y/S}(Y)$의
역상은 $X$에서 열린다. 그리고 가정 (2)와 (3)에 의해 $x$는 이 열린
부분집합에 속한다.
\end{proof}
\section{\'Etale 사상}
\label{morphisms-section-etale}

\noindent
스킴의 자리스키 위상은 매우 거친 위상이다. 이는 $\mathbf{C}$ 위의
다형체를 보면 특히 분명하다. \'etale 사상을 위상수학에서의 국소
동형사상에 대응하는 개념으로 선언하면 훨씬 더 세밀한 위상이 생긴다.
$\mathbf{C}$ 위의 다형체에서 이 위상은 유한 계수 코호몰로지를 계산할 때
``올바른'' 베티 수를 준다. 또 관찰할 수 있는 사실은
$f : X \to Y$가 $\mathbf{C}$ 위 다형체의 \'etale 사상이고 $x$가 $X$의
닫힌점이면 $f$가 완비 국소환의 동형사상
$\mathcal{O}^{\wedge}_{Y, f(x)} \to \mathcal{O}^{\wedge}_{X, x}$를
유도한다는 것이다.

\medskip\noindent
이 절에서 이러한 문제의 연구를 시작한다. 실제로 더 흥미로운 결과들은
다른 곳에서 논의할 것이므로 여기서는 의도적으로 논의를 최소한으로
제한한다. 환 준동형 $R \to A$가 매끄럽고 $\Omega_{A/R} = 0$이면 이를
{\it \'etale}이라 함을 상기하자. 대수, 정의
\ref{algebra-definition-etale}을 보라.

\begin{definition}
\label{morphisms-definition-etale}
$f : X \to S$를 스킴의 사상이라 하자.
\begin{enumerate}
\item $f$가 {\it $x \in X$에서 \'etale}이라고 함은 아핀 열린 근방
$\Spec(A) = U \subset X$로서 $x$를 포함하는 것과 아핀 열린집합
$\Spec(R) = V \subset S$가 존재하여 $f(U) \subset V$이고 유도된 환 준동형
$R \to A$가 \'etale인 경우이다.
\item $f$가 $X$의 모든 점에서 \'etale이면 이를 {\it \'etale}이라 한다.
\item 아핀 스킴의 사상 $f : X \to S$에 대해 $X \to S$가
$$
\Spec(R[x]_h/(g)) \to \Spec(R)
$$
와 동형이고, 여기서 $R \to R[x]_h/(g)$가 표준 \'etale 환 준동형이면
(대수, 정의 \ref{algebra-definition-standard-etale}을 보라) 이를
{\it 표준 \'etale}이라 한다. 즉 $g$는 monic이고 $g'$는
$R[x]_h/(g)$에서 가역이다.
\end{enumerate}
\end{definition}

\noindent
사상이 \'etale인 것은 상대차원 $0$으로 매끄러운 것과 동치이다(정의
\ref{morphisms-definition-smooth-relative-dimension}을 보라). 이 정의의 좋은 점
하나는 사상이 \'etale인 점들의 집합이 자동으로 열린다는 것이다.

\medskip\noindent
정의에는 분리성이나 준콤팩트성 가정이 없다는 데 유의하자. 따라서
\'etale 여부는 원천에서 국소적인 문제이다. 정확한 결과는 다음과 같다.

\begin{lemma}
\label{morphisms-lemma-etale-characterize}
$f : X \to S$를 스킴의 사상이라 하자. 다음은 동치이다.
\begin{enumerate}
\item 사상 $f$는 \'etale이다.
\item 모든 아핀 열린집합 $U \subset X$, $V \subset S$에 대해
$f(U) \subset V$이면 환 준동형
$\mathcal{O}_S(V) \to \mathcal{O}_X(U)$가 \'etale이다.
\item 열린 덮개 $S = \bigcup_{j \in J} V_j$와 열린 덮개
$f^{-1}(V_j) = \bigcup_{i \in I_j} U_i$가 존재하여 각 사상
$U_i \to V_j$, $j\in J, i\in I_j$가 \'etale이다.
\item 아핀 열린 덮개 $S = \bigcup_{j \in J} V_j$와 아핀 열린 덮개
$f^{-1}(V_j) = \bigcup_{i \in I_j} U_i$가 존재하여 환 준동형
$\mathcal{O}_S(V_j) \to \mathcal{O}_X(U_i)$가 모든
$j\in J, i\in I_j$에 대해 \'etale이다.
\end{enumerate}
또한 $f$가 \'etale이면 임의의 열린 부분스킴 $U \subset X$, $V \subset S$에
대해 $f(U) \subset V$일 때 제한 $f|_U : U \to V$도 \'etale이다.
\end{lemma}

\begin{proof}
``$R \to A$가 \'etale이다''라는 성질이 국소적임을 보이면 보조정리
\ref{morphisms-lemma-locally-P}에서 따른다. 정의
\ref{morphisms-definition-property-local}의 조건 (a), (b), (c)를 확인하자.
이들은 모두 대수, 보조정리 \ref{algebra-lemma-etale}에서 따른다.
\end{proof}

\begin{lemma}
\label{morphisms-lemma-composition-etale}
\'Etale인 두 사상의 합성은 \'etale이다.
\end{lemma}

\begin{proof}
보조정리 \ref{morphisms-lemma-etale-characterize}의 증명에서 \'etale성이 환 준동형의
국소적 성질임을 보았다. 따라서 보조정리 \ref{morphisms-lemma-composition-type-P}와
\'etale성이 합성으로 보존되는 환 준동형의 성질이라는 사실을 결합하면
보조정리의 첫째 명제가 따른다. 대수, 보조정리
\ref{algebra-lemma-etale}을 보라.
\end{proof}

\begin{lemma}
\label{morphisms-lemma-base-change-etale}
\'Etale 사상의 기저변환은 \'etale이다.
\end{lemma}

\begin{proof}
보조정리 \ref{morphisms-lemma-etale-characterize}의 증명에서 \'etale성이 환 준동형의
국소적 성질임을 보았다. 따라서 보조정리 \ref{morphisms-lemma-composition-type-P}와
\'etale성이 기저변환으로 보존되는 환 준동형의 성질이라는 사실을
결합하면 결과가 따른다. 대수, 보조정리 \ref{algebra-lemma-etale}을 보라.
\end{proof}

\begin{lemma}
\label{morphisms-lemma-etale-smooth-unramified}
$f : X \to S$를 스킴의 사상이라 하고 $x \in X$라 하자. 그러면 $f$가
$x$에서 \'etale인 것은 $f$가 $x$에서 매끄럽고 비분기인 것과 동치이다.
\end{lemma}

\begin{proof}
정의에서 곧바로 따른다.
\end{proof}

\begin{lemma}
\label{morphisms-lemma-etale-locally-quasi-finite}
\'Etale 사상은 국소적으로 준유한이다.
\end{lemma}

\begin{proof}
보조정리 \ref{morphisms-lemma-etale-smooth-unramified}에 의해 \'etale 사상은
비분기이다. 보조정리 \ref{morphisms-lemma-unramified-quasi-finite}에 의해 비분기
사상은 국소적으로 준유한이다.
\end{proof}

\begin{lemma}
\label{morphisms-lemma-etale-over-field}
\begin{slogan}
체 위의 \'etale 스킴과 \'etale 사상의 올의 기술.
\end{slogan}
\'Etale 사상의 올에 대해 다음이 성립한다.
\begin{enumerate}
\item $X$를 체 $k$ 위의 스킴이라 하자. 구조 사상
$X \to \Spec(k)$가 \'etale인 것은 $X$가 $k$의 유한 분리 체확장의
스펙트럼들의 서로소 합집합인 것과 동치이다.
\item $f : X \to S$가 \'etale 사상이면 모든 $s \in S$에 대해 올 $X_s$는
$\kappa(s)$의 유한 분리 체확장의 스펙트럼들의 서로소 합집합이다.
\end{enumerate}
\end{lemma}

\begin{proof}
보조정리 \ref{morphisms-lemma-unramified-over-field}에서, 위 보조정리
\ref{morphisms-lemma-etale-smooth-unramified}를 통해 이를 추론할 수 있다. 직접적인
증명은 다음과 같다.

\medskip\noindent
대수, 보조정리 \ref{algebra-lemma-etale-over-field}를 사용하겠다. 따라서
$X$가 $k$의 유한 분리 체확장의 스펙트럼들의 서로소 합집합이면
$X \to \Spec(k)$가 \'etale임은 명백하다. 역으로 $X \to \Spec(k)$가
\'etale이라고 가정하자. 그러면 임의의 아핀 열린집합 $U \subset X$에
대해 $U$는 $k$의 유한 분리 체확장의 스펙트럼들의 유한 서로소 합집합이다.
따라서 $X$의 모든 점은 닫힌점이다(예를 들어 보조정리
\ref{morphisms-lemma-algebraic-residue-field-extension-closed-point-fibre}를 보라).
그러므로 $X$는 이산공간이고 결론을 얻는다.
\end{proof}

\noindent
다음 보조정리는 \'etale 사상을 ``\'etale 올''을 갖는 평탄 유한 표시
사상으로 특징짓는다.

\begin{lemma}
\label{morphisms-lemma-etale-flat-etale-fibres}
$f : X \to S$를 스킴의 사상이라 하자. $f$가 평탄이고 국소적으로 유한
표시이며 모든 $s \in S$에 대해 올 $X_s$가 $\kappa(s)$의 유한 분리
체확장의 스펙트럼들의 서로소 합집합이면 $f$는 \'etale이다.
\end{lemma}

\begin{proof}
대수, 보조정리 \ref{algebra-lemma-characterize-etale}에서 이를 추론할 수
있다. 다른 증명은 다음과 같다.

\medskip\noindent
보조정리 \ref{morphisms-lemma-etale-over-field}에 의해 올 $X_s$는 \'etale이고 따라서
$s$ 위에서 매끄럽다. 보조정리 \ref{morphisms-lemma-smooth-flat-smooth-fibres}에 의해
$X \to S$가 매끄러움을 안다. 보조정리
\ref{morphisms-lemma-unramified-etale-fibres}에 의해 $f$가 비분기임을 안다.
보조정리 \ref{morphisms-lemma-etale-smooth-unramified}로 결론을 얻는다.
\end{proof}

\begin{lemma}
\label{morphisms-lemma-open-immersion-etale}
모든 열린 몰입은 \'etale이다.
\end{lemma}

\begin{proof}
열린 몰입은 국소 동형사상이므로 성립한다.
\end{proof}

\begin{lemma}
\label{morphisms-lemma-etale-syntomic}
\'Etale 사상은 신토믹이다.
\end{lemma}

\begin{proof}
대수, 보조정리 \ref{algebra-lemma-smooth-syntomic}을 보고, \'etale 사상은
상대차원 $0$인 매끄러운 사상과 같다는 사실을 사용하라.
\end{proof}

\begin{lemma}
\label{morphisms-lemma-etale-locally-finite-presentation}
\'Etale 사상은 국소적으로 유한 표시이다.
\end{lemma}

\begin{proof}
정의에 의해 \'etale 환 준동형은 유한 표시이므로 성립한다.
\end{proof}

\begin{lemma}
\label{morphisms-lemma-etale-flat}
\'Etale 사상은 평탄하다.
\end{lemma}

\begin{proof}
보조정리 \ref{morphisms-lemma-syntomic-flat}와
\ref{morphisms-lemma-etale-syntomic}을 결합하라.
\end{proof}

\begin{lemma}
\label{morphisms-lemma-etale-open}
\'Etale 사상은 열린 사상이다.
\end{lemma}

\begin{proof}
보조정리 \ref{morphisms-lemma-etale-flat},
\ref{morphisms-lemma-etale-locally-finite-presentation}, 그리고
\ref{morphisms-lemma-fppf-open}을 결합하라.
\end{proof}

\noindent
다음 보조정리는 모든 \'etale 사상이 국소적으로 표준 \'etale임을 말한다.
완전한 일반성 아래 이 결과를 증명하는 일은 실제로 다소 까다롭다.
까다로운 부분은 가환대수 장에 들어 있다. 따라서 표준 \'etale 사상은
일반적인 \'etale 사상의 {\it 국소 모형}이다.

\begin{lemma}
\label{morphisms-lemma-etale-locally-standard-etale}
$f : X  \to S$를 스킴의 사상이라 하자.
$x \in X$를 한 점이라 하자.
$V \subset S$를 $f(x)$의 아핀 열린 근방이라 하자.
다음은 동치이다.
\begin{enumerate}
\item 사상 $f$는 $x$에서 \'etale이다.
\item 아핀 열린집합 $U \subset X$가 존재하여 $x \in U$이고
$f(U) \subset V$이며, 유도된 사상 $f|_U : U \to V$가 표준 \'etale이다
(정의 \ref{morphisms-definition-etale}을 보라).
\end{enumerate}
\end{lemma}

\begin{proof}
정의와 대수, 명제
\ref{algebra-proposition-etale-locally-standard}에서 따른다.
\end{proof}

\noindent
한 점에서 \'etale일 미분 판정법을 제시한다. 이 결과에는 여러 변형이
있고, 어느 시점에는 모두 유용할 수 있다. 필요할 때마다 여기에
추가하겠다.

\begin{lemma}
\label{morphisms-lemma-etale-at-point}
$f : X \to S$를 스킴의 사상이라 하자.
$x \in X$라 하자.
$s = f(x)$로 놓자.
$f$가 국소적으로 유한 표시라고 가정하자.
다음은 동치이다.
\begin{enumerate}
\item 사상 $f$는 $x$에서 \'etale이다.
\item 국소환 준동형 $\mathcal{O}_{S, s} \to \mathcal{O}_{X, x}$는
평탄하고 $X_s \to \Spec(\kappa(s))$는 $x$에서 \'etale이다.
\item 국소환 준동형 $\mathcal{O}_{S, s} \to \mathcal{O}_{X, x}$는
평탄하고 $X_s \to \Spec(\kappa(s))$는 $x$에서 비분기이다.
\item 국소환 준동형 $\mathcal{O}_{S, s} \to \mathcal{O}_{X, x}$는
평탄하고 $\mathcal{O}_{X, x}$-가군 $\Omega_{X/S, x}$는 영이다.
\item 국소환 준동형 $\mathcal{O}_{S, s} \to \mathcal{O}_{X, x}$는
평탄하고 $\kappa(x)$-벡터공간
$$
\Omega_{X_s/s, x} \otimes_{\mathcal{O}_{X_s, x}} \kappa(x) =
\Omega_{X/S, x} \otimes_{\mathcal{O}_{X, x}} \kappa(x)
$$
은 영이다.
\item 국소환 준동형 $\mathcal{O}_{S, s} \to \mathcal{O}_{X, x}$는
평탄하고, $\mathfrak m_s\mathcal{O}_{X, x} = \mathfrak m_x$이며,
체확장 $\kappa(x)/\kappa(s)$는 유한 분리이다.
\item 아핀 열린집합 $U \subset X$와 $V \subset S$가 존재하여
$x \in U$, $f(U) \subset V$이고, 유도된 사상
$f|_U : U \to V$는 상대차원 $0$인 표준 매끄러운 사상이다.
\item 아핀 열린집합 $\Spec(A) = U \subset X$와
$\Spec(R) = V \subset S$가 존재하고, $x \in U$는
$\mathfrak q \subset A$에 대응하며 $f(U) \subset V$이고,
다음과 같은 표시가 존재한다.
$$
A = R[x_1, \ldots, x_n]/(f_1, \ldots, f_n)
$$
여기서
$$
g =
\det
\left(
\begin{matrix}
\partial f_1/\partial x_1 &
\partial f_2/\partial x_1 &
\ldots &
\partial f_n/\partial x_1 \\
\partial f_1/\partial x_2 &
\partial f_2/\partial x_2 &
\ldots &
\partial f_n/\partial x_2 \\
\ldots & \ldots & \ldots & \ldots \\
\partial f_1/\partial x_n &
\partial f_2/\partial x_n &
\ldots &
\partial f_n/\partial x_n
\end{matrix}
\right)
$$
는 $A$의 원소 가운데 $\mathfrak q$에 속하지 않는 원소로 사상된다.
\item 아핀 열린집합 $U \subset X$와 $V \subset S$가 존재하여
$x \in U$, $f(U) \subset V$이고, 유도된 사상
$f|_U : U \to V$는 표준 \'etale이다.
\item 아핀 열린집합 $\Spec(A) = U \subset X$와
$\Spec(R) = V \subset S$가 존재하고, $x \in U$는
$\mathfrak q \subset A$에 대응하며 $f(U) \subset V$이고,
다음과 같은 표시가 존재한다.
$$
A = R[x]_Q/(P) = R[x, 1/Q]/(P)
$$
여기서 $P, Q \in R[x]$이고, $P$는 모닉이며,
$P' = \text{d}P/\text{d}x$는 $A$의 원소 가운데 $\mathfrak q$에
속하지 않는 원소로 사상된다.
\end{enumerate}
\end{lemma}

\begin{proof}
보조정리 \ref{morphisms-lemma-etale-locally-standard-etale}와 정의를 사용하면
(1)이 나머지 모든 조건을 함의함을 알 수 있다. 조건 (2) -- (10) 각각에
대해 보조정리 \ref{morphisms-lemma-smooth-at-point}와
\ref{morphisms-lemma-unramified-at-point}를 결합하여, $f$가 $x$에서 매끄럽고
비분기임을 보이고 보조정리 \ref{morphisms-lemma-etale-smooth-unramified}을 적용하면
(1)이 성립함을 알 수 있다. 일부 세부사항은 생략한다.
\end{proof}

\begin{lemma}
\label{morphisms-lemma-flat-unramified-etale}
한 점에서 사상이 \'etale인 것은 그 점에서 평탄하고 G-비분기인 것과
동치이다. 사상이 \'etale인 것은 평탄하고 G-비분기인 것과 동치이다.
\end{lemma}

\begin{proof}
보조정리 \ref{morphisms-lemma-etale-at-point}와
\ref{morphisms-lemma-unramified-at-point}에서 명백하다.
\end{proof}

\begin{lemma}
\label{morphisms-lemma-set-points-where-fibres-etale}
다음을 스킴의 카르테시안 도표라 하자.
$$
\xymatrix{
X' \ar[r]_{g'} \ar[d]_{f'} & X \ar[d]^f \\
S' \ar[r]^g & S
}
$$
$W \subset X$, 각각 $W' \subset X'$를 $f$, 각각 $f'$가 \'etale인
점들로 이루어진 열린 부분스킴이라 하자. 다음 중 하나가 성립하면
$W' = (g')^{-1}(W)$이다.
\begin{enumerate}
\item $f$가 평탄이고 국소적으로 유한 표시이다.
\item $f$가 국소적으로 유한 표시이고 $g$가 평탄하다.
\end{enumerate}
\end{lemma}

\begin{proof}
먼저 $f$가 국소적으로 유한형이라고 가정하자. 집합
$$
T = \{x \in X \mid f\text{ is unramified at }x\}
$$
과 대응하는 집합 $T' \subset X'$를 생각하자. 이 집합은 $f'$에 대한
것이다. 그러면
보조정리 \ref{morphisms-lemma-set-points-where-fibres-unramified}에 의해
$T' = (g')^{-1}(T)$이다.

\medskip\noindent
따라서 경우 (1)에는 보조정리 \ref{morphisms-lemma-flat-unramified-etale}에 의해
$T$가 $f$가 \'etale인 점들의 (열린) 집합이므로 결론이 따른다.

\medskip\noindent
경우 (2)에는 $x' \in W'$라 하자. 그러면 $g'$는 $x'$에서 평탄하고
(보조정리 \ref{morphisms-lemma-base-change-module-flat}), $g \circ f'$는 $x'$에서
평탄하다(보조정리 \ref{morphisms-lemma-composition-module-flat}). 따라서 보조정리
\ref{morphisms-lemma-flat-permanence}에 의해 $f$는 $x = g'(x')$에서 평탄하다.
한편 $x' \in T'$이므로(보조정리 \ref{morphisms-lemma-base-change-smooth})
$x \in T$임을 안다. 따라서 보조정리 \ref{morphisms-lemma-etale-at-point}에 의해
$f$는 $x$에서 \'etale이다.
\end{proof}

\begin{lemma}
\label{morphisms-lemma-etale-permanence-better}
\begin{reference}
예를 들어 \cite[Remark 18.30 (2).]{GWII}를 보라.
\end{reference}
$f : X \to Y$를 $S$ 위의 스킴의 사상이라 하자.
$X$가 $S$ 위에서 \'etale이고 $Y$가 $S$ 위에서 비분기이면
$f$는 \'etale이다.
\end{lemma}

\begin{proof}
분해 $X \to X \times_S Y \to Y$를 생각하자. 여기서 첫 번째 화살표는
$\text{id}_X$와 $f$로 주어지고 두 번째 화살표는 사영이다. 두 화살표가
모두 \'etale임을 주장하며, 그러면 보조정리
\ref{morphisms-lemma-composition-etale}에 의해 $f$는 \'etale이다. 실제로 사영은
$X \to S$의 기저변환이므로 \'etale이다. 보조정리
\ref{morphisms-lemma-base-change-etale}을 보라. 첫 번째 화살표는 다음 사각형이
카르테시안이므로 대각 사상 $Y \to Y \times_S Y$의 기저변환이다.
$$
\xymatrix{
X \ar[d] \ar[r] & X \times_S Y \ar[d] \\
Y \ar[r] & Y \times_S Y
}
$$
보조정리 \ref{morphisms-lemma-diagonal-unramified-morphism}에 의해 대각 사상
$Y \to Y \times_S Y$는 열린 몰입이다. 열린 몰입의 기저변환은 열린
몰입이고(스킴, 보조정리 \ref{schemes-lemma-base-change-immersion}),
열린 몰입은 \'etale이다(보조정리 \ref{morphisms-lemma-open-immersion-etale}).
\end{proof}

\begin{lemma}
\label{morphisms-lemma-etale-permanence}
\begin{slogan}
\'Etale 사상의 소거법
\end{slogan}
$f : X \to Y$를 $S$ 위의 스킴의 사상이라 하자.
$X$와 $Y$가 $S$ 위에서 \'etale이면 $f$는 \'etale이다.
\end{lemma}

\begin{proof}
보조정리 \ref{morphisms-lemma-etale-smooth-unramified}에 의해 $Y \to S$가
비분기이므로 보조정리 \ref{morphisms-lemma-etale-permanence-better}를 적용할 수
있다. 대수, 보조정리 \ref{algebra-lemma-map-between-etale}도 보라.
\end{proof}

\begin{lemma}
\label{morphisms-lemma-etale-permanence-two}
다음을 스킴의 사상들로 이루어진 가환 도표라 하자.
$$
\xymatrix{
X \ar[rr]_f \ar[rd]_p & &
Y \ar[dl]^q \\
& S
}
$$
다음을 가정하자.
\begin{enumerate}
\item $f$는 전사이고 \'etale이다.
\item $p$는 \'etale이다.
\item $q$는 국소적으로 유한 표시이다\footnote{실제로 이는 (1)과
(2)에서 함의된다. 하강, 보조정리
\ref{descent-lemma-flat-finitely-presented-permanence}를 보라. 더 나아가
$f$가 전사이고 평탄하며 국소적으로 유한 표시라고 가정하는 것으로
충분하다. 하강, 보조정리 \ref{descent-lemma-smooth-permanence}를 보라.}.
\end{enumerate}
그러면 $q$는 \'etale이다.
\end{lemma}

\begin{proof}
보조정리 \ref{morphisms-lemma-smooth-permanence}에 의해 $q$가 매끄러움을 안다.
따라서 $q$의 상대차원이 $0$임만 보이면 된다. 이는 보조정리
\ref{morphisms-lemma-dimension-fibre-at-a-point-additive}와 $f$ 및 $p$의 상대차원이
$0$이라는 사실에서 따른다.
\end{proof}

\noindent
매끄러운 사상의 마지막 특징짓기는 다음과 같다. 매끄러운 사상
$f : X \to S$는 국소적으로 \'etale 사상과 사영
$\mathbf{A}_S^d \to S$의 합성이다.

\begin{lemma}
\label{morphisms-lemma-smooth-etale-over-affine-space}
\begin{slogan}
매끄러운 스킴은 \'etale-국소적으로 아핀공간과 같다.
\end{slogan}
$\varphi : X \to Y$를 스킴의 사상이라 하자. $x \in X$라 하자.
$V \subset Y$를 $\varphi(x)$의 아핀 열린 근방이라 하자.
$\varphi$가 $x$에서 매끄러우면, 정수 $d \geq 0$과 아핀 열린집합
$U \subset X$가 존재하여 $x \in U$, $\varphi(U) \subset V$이고,
다음 가환 도표가 존재한다.
$$
\xymatrix{
X \ar[d] & U \ar[l] \ar[d] \ar[r]_-\pi & \mathbf{A}^d_V \ar[ld] \\
Y & V \ar[l]
}
$$
여기서 $\pi$는 \'etale이다.
\end{lemma}

\begin{proof}
보조정리 \ref{morphisms-lemma-smooth-locally-standard-smooth}에 의해 보조정리에
나오는 아핀 열린집합 $U$를 찾아서
$\varphi|_U : U \to V$가 표준 매끄러운 사상이 되게 할 수 있다.
$U = \Spec(A)$와 $V = \Spec(R)$로 쓰면 다음과 같이 쓸 수 있다.
$$
A = R[x_1, \ldots, x_n]/(f_1, \ldots, f_c)
$$
여기서
$$
g =
\det
\left(
\begin{matrix}
\partial f_1/\partial x_1 &
\partial f_2/\partial x_1 &
\ldots &
\partial f_c/\partial x_1 \\
\partial f_1/\partial x_2 &
\partial f_2/\partial x_2 &
\ldots &
\partial f_c/\partial x_2 \\
\ldots & \ldots & \ldots & \ldots \\
\partial f_1/\partial x_c &
\partial f_2/\partial x_c &
\ldots &
\partial f_c/\partial x_c
\end{matrix}
\right)
$$
는 $A$의 가역원으로 사상된다. 그러면
$R[x_{c + 1}, \ldots, x_n] \to A$가 상대차원 $0$인 표준 매끄러운
사상임은 명백하다. 따라서 이는 상대차원 $0$인 매끄러운 사상이다.
달리 말해 환 준동형 $R[x_{c + 1}, \ldots, x_n] \to A$는 \'etale이다.
$\mathbf{A}^{n - c}_V = \Spec(R[x_{c + 1}, \ldots, x_n])$이므로
$d = n - c$인 보조정리.
\end{proof}
















\section{상대적으로 충분한 층}
\label{morphisms-section-relatively-ample}

\noindent
$X$를 스킴이라 하고 $\mathcal{L}$을 $X$ 위의 가역층이라 하자.
$\mathcal{L}$이 $X$ 위에서 충분하다는 것은 $X$가 준콤팩트이고
$X$의 모든 점이 $X_s$ 꼴의 아핀 열린집합에 포함된다는 뜻이다. 여기서
$s \in \Gamma(X, \mathcal{L}^{\otimes n})$이고 $n \geq 1$이다.
성질, 정의 \ref{properties-definition-ample}을 보라. 이를 다음과 같이
상대적 개념으로 바꾼다.

\begin{definition}
\label{morphisms-definition-relatively-ample}
\begin{reference}
\cite[II Definition 4.6.1]{EGA}
\end{reference}
$f : X \to S$를 스킴의 사상이라 하자.
$\mathcal{L}$을 가역 $\mathcal{O}_X$-가군이라 하자.
$\mathcal{L}$을 {\it 상대적으로 충분하다}, 또는 {\it $f$-상대적으로
충분하다}, 또는 {\it $X/S$ 위에서 충분하다}, 또는 {\it $f$-충분하다}고
하는 것은 다음을 뜻한다. $f : X \to S$가 준콤팩트이고 모든 아핀
열린집합 $V \subset S$에 대해 $\mathcal{L}$을 열린 부분스킴
$f^{-1}(V)$로 제한한 것이 충분하다. 이 부분스킴은 $X$ 안에 있다.
\end{definition}

\noindent
$X$ 위에 상대적으로 충분한 층이 존재한다고 해서 사상 $X \to S$가
유한형일 필요는 없음을 지적해 둔다.

\begin{lemma}
\label{morphisms-lemma-ample-power-ample}
$X \to S$를 스킴의 사상이라 하자.
$\mathcal{L}$을 가역 $\mathcal{O}_X$-가군이라 하자.
$n \geq 1$이라 하자. 그러면 $\mathcal{L}$이 $f$-충분할 필요충분조건은
$\mathcal{L}^{\otimes n}$이 $f$-충분한 것이다.
\end{lemma}

\begin{proof}
성질, 보조정리 \ref{properties-lemma-ample-power-ample}에서 따른다.
\end{proof}

\begin{lemma}
\label{morphisms-lemma-relatively-ample-separated}
$f : X \to S$를 스킴의 사상이라 하자.
$f$-충분한 가역층이 존재하면 $f$는 분리되어 있다.
\end{lemma}

\begin{proof}
분리성은 밑에서 국소적이다(예를 들어 스킴, 보조정리
\ref{schemes-lemma-characterize-separated}을 보라. 정의에서도 쉽게
따른다). 따라서 $S$가 아핀이고 $X$가 충분한 가역층을 갖는다고
가정해도 된다. 이 경우 결과는 성질, 보조정리
\ref{properties-lemma-ample-separated}에서 따른다.
\end{proof}

\noindent
성질, 명제 \ref{properties-proposition-characterize-ample}의 동치 조건들과
비슷하게, 상대적으로 충분한 가역층을 특징짓는 방법은 많다. 필요할
때마다 여기에 추가하겠다.

\begin{lemma}
\label{morphisms-lemma-characterize-relatively-ample}
\begin{reference}
\cite[II, Proposition 4.6.3]{EGA}
\end{reference}
$f : X \to S$를 스킴의 준콤팩트 사상이라 하자.
$\mathcal{L}$을 $X$ 위의 가역층이라 하자.
다음은 동치이다.
\begin{enumerate}
\item 가역층 $\mathcal{L}$은 $f$-충분하다.
\item 열린 덮개 $S = \bigcup V_i$가 존재하여 각각의
$\mathcal{L}|_{f^{-1}(V_i)}$는 $f^{-1}(V_i) \to V_i$에 대해 충분하다.
\item 아핀 열린 덮개 $S = \bigcup V_i$가 존재하여 각각의
$\mathcal{L}|_{f^{-1}(V_i)}$가 충분하다.
\item 준연접 등급 $\mathcal{O}_S$-대수 $\mathcal{A}$와 등급
$\mathcal{O}_X$-대수의 사상
$\psi : f^*\mathcal{A} \to \bigoplus_{d \geq 0} \mathcal{L}^{\otimes d}$가
존재하여 $U(\psi) = X$이고
$$
r_{\mathcal{L}, \psi} :
X
\longrightarrow
\underline{\text{Proj}}_S(\mathcal{A})
$$
가 열린 몰입이다(기호에 대해서는 구성, 보조정리
\ref{constructions-lemma-invertible-map-into-relative-proj}를 보라).
\item 사상 $f$는 준분리이고, 위의 (4)가
$\mathcal{A} = f_*(\bigoplus_{d \geq 0} \mathcal{L}^{\otimes d})$와
수반사상 $\psi$에 대해 성립한다.
\item (4)와 같되 $r_{\mathcal{L}, \psi}$가 몰입이기만 하면 된다.
\end{enumerate}
\end{lemma}

\begin{proof}
정의에서 (1)이 (2)를 함의하고 (2)가 (3)을 함의함은 즉시 알 수 있다.
(5)가 (4)를 함의함은 명백하다.

\medskip\noindent
(3)이 아핀 열린 덮개 $S = \bigcup V_i$에 대해 성립한다고 가정하자.
(5)가 성립함을 보이겠다. 각각의 $f^{-1}(V_i)$가 충분한 가역층을
가지므로 $f^{-1}(V_i)$는 분리되어 있다(성질, 보조정리
\ref{properties-lemma-ample-separated}). 따라서 $f$는 분리되어 있다.
스킴, 보조정리 \ref{schemes-lemma-push-forward-quasi-coherent}에 의해
$\mathcal{A} = f_*(\bigoplus_{d \geq 0} \mathcal{L}^{\otimes d})$는
준연접 등급 $\mathcal{O}_S$-대수이다. 수반사상을
$\psi : f^*\mathcal{A} \to \bigoplus_{d \geq 0} \mathcal{L}^{\otimes d}$로
나타내자. 구성, 절
\ref{constructions-section-invertible-relative-proj}에 나오는 열린집합
$U(\psi)$의 기술과 $\mathcal{L}|_{f^{-1}(V_i)}$의 충분성 정의에 의해
$U(\psi) = X$이다. 또한 구성, 보조정리
\ref{constructions-lemma-invertible-map-into-relative-proj}의 (3)에 의하면
$r_{\mathcal{L}, \psi}$를 $f^{-1}(V_i)$로 제한한 사상은 성질, 보조정리
\ref{properties-lemma-map-into-proj}의 사상과 같고, 이는 성질, 보조정리
\ref{properties-lemma-ample-immersion-into-proj}에 의해 열린 몰입이다.
따라서 (5)가 성립한다.

\medskip\noindent
(4)가 (1)을 함의함을 보이자. (4)를 가정하자. 구조 사상을
$\pi : \underline{\text{Proj}}_S(\mathcal{A}) \to S$로 나타내자.
아핀 열린집합 $V \subset S$를 택하자. 구성, 정의
\ref{constructions-definition-relative-proj}에 의해
$\pi^{-1}(V) \subset \underline{\text{Proj}}_S(\mathcal{A})$는
$\text{Proj}(A)$와 같으며, 여기서 등급환 $A = \mathcal{A}(V)$이다. 따라서
$r_{\mathcal{L}, \psi}$는 $f^{-1}(V)$를 $\text{Proj}(A)$의 준콤팩트
열린집합 위로 동형사상으로 보낸다. 더욱이 $\mathcal{L}^{\otimes d}$는
$\mathcal{O}_{\text{Proj}(A)}(d)$의 역상과 동형이며, 이는 어떤
$d \geq 1$에 대해 성립한다. (구성, 보조정리
\ref{constructions-lemma-invertible-map-into-relative-proj}의 (3)과 구성,
보조정리 \ref{constructions-lemma-invertible-map-into-proj}의 마지막 명제를
보라.) 이는 성질, 보조정리
\ref{properties-lemma-open-in-proj-ample}와
\ref{properties-lemma-ample-power-ample}에 의해
$\mathcal{L}|_{f^{-1}(V)}$가 충분함을 뜻한다.

\medskip\noindent
(6)을 가정하자. 위의 (1) - (5)의 동치에 의해 $X/S$ 위에서
상대적으로 충분하다는 성질은 $S$에서 국소적이다. 따라서 $S$가
아핀이라고 가정해도 되며, $\mathcal{L}$이 $X$ 위에서 충분함을 보여야
한다. 이 경우 사상 $r_{\mathcal{L}, \psi}$는 사상
$r_{\mathcal{L}, \psi} : X \to \text{Proj}(A)$와 동일시된다. 이 사상도
같은 기호로 쓰며, 사상
$\psi : A = \mathcal{A}(V) \to \Gamma_*(X, \mathcal{L})$에 딸린 것이다.
(위의 참고문헌을 보라.) 위와 마찬가지로 $\mathcal{L}^{\otimes d}$는
층 $\mathcal{O}_{\text{Proj}(A)}(d)$의 역상이며, 이는 어떤 $d \geq 1$에
대해 성립함을 안다. 더욱이 $X$가 준콤팩트이므로 $X$는 준콤팩트 열린
부분스킴 $Y \subset \text{Proj}(A)$의 닫힌 부분스킴과 동일시된다.
구성, 보조정리 \ref{constructions-lemma-ample-on-proj}(성질, 보조정리
\ref{properties-lemma-open-in-proj-ample}도 보라)에 의해
$\mathcal{O}_Y(d')$는 $Y$ 위의 충분한 가역층이며, 이는 어떤
$d' \geq 1$에 대해 성립한다.
충분한 층을 닫힌 부분스킴으로 제한하면 충분하므로, 성질, 보조정리
\ref{properties-lemma-ample-on-closed}를 보라, $\mathcal{O}_Y(d')$의
역상은 충분하다. 이 결과들과 성질, 보조정리
\ref{properties-lemma-ample-power-ample}를 결합하면 원하는 대로
$\mathcal{L}$이 충분하다는 결론을 얻는다.
\end{proof}

\begin{lemma}
\label{morphisms-lemma-ample-over-affine}
\begin{reference}
\cite[II Corollary 4.6.6]{EGA}
\end{reference}
$f : X \to S$를 스킴의 사상이라 하자.
$\mathcal{L}$을 가역 $\mathcal{O}_X$-가군이라 하자.
$S$가 아핀이라고 가정하자. 그러면 $\mathcal{L}$이 $f$-상대적으로
충분할 필요충분조건은 $\mathcal{L}$이 $X$ 위에서 충분한 것이다.
\end{lemma}

\begin{proof}
보조정리 \ref{morphisms-lemma-characterize-relatively-ample}와 정의에서 즉시 따른다.
\end{proof}

\begin{lemma}
\label{morphisms-lemma-quasi-affine-O-ample}
\begin{reference}
\cite[II Proposition 5.1.6]{EGA}
\end{reference}
$f : X \to S$를 스킴의 사상이라 하자. 그러면 $f$가 준아핀일
필요충분조건은 $\mathcal{O}_X$가 $f$-상대적으로 충분한 것이다.
\end{lemma}

\begin{proof}
성질, 보조정리 \ref{properties-lemma-quasi-affine-O-ample}와 정의에서
따른다.
\end{proof}

\begin{lemma}
\label{morphisms-lemma-pullback-ample-tensor-relatively-ample}
$f : X \to Y$를 스킴의 사상, $\mathcal{M}$을 가역
$\mathcal{O}_Y$-가군, $\mathcal{L}$을 가역 $\mathcal{O}_X$-가군이라 하자.
\begin{enumerate}
\item $\mathcal{L}$이 $f$-충분하고 $\mathcal{M}$이 충분하면
$\mathcal{L} \otimes f^*\mathcal{M}^{\otimes a}$는 $a \gg 0$에 대해
충분하다.
\item $\mathcal{M}$이 충분하고 $f$가 준아핀이면
$f^*\mathcal{M}$은 충분하다.
\end{enumerate}
\end{lemma}

\begin{proof}
$\mathcal{L}$이 $f$-충분하고 $\mathcal{M}$이 충분하다고 가정하자.
가정에 의해 $Y$와 $f$는 준콤팩트이다(정의
\ref{morphisms-definition-relatively-ample}과 성질, 정의
\ref{properties-definition-ample}을 보라). 따라서 $X$는 준콤팩트이다.
성질, 보조정리 \ref{properties-lemma-ample-separated}에 의해 스킴 $Y$는
분리되어 있고, 보조정리 \ref{morphisms-lemma-relatively-ample-separated}에 의해
사상 $f$는 분리되어 있다. 따라서 스킴, 보조정리
\ref{schemes-lemma-separated-permanence}에 의해 $X$는 분리되어 있다.
$x \in X$를 택하자. $m \geq 1$과
$t \in \Gamma(Y, \mathcal{M}^{\otimes m})$를 택하여 $Y_t$가 아핀이고
$f(x) \in Y_t$가 되게 할 수 있다. $\mathcal{L}$을
$f^{-1}(Y_t) = X_{f^*t}$로 제한하면 충분한 가역층이므로, $n \geq 1$과
$s \in \Gamma(X_{f^*t}, \mathcal{L}^{\otimes n})$를 택하여
$x \in (X_{f^*t})_s$이고 $(X_{f^*t})_s$가 아핀이 되게 할 수 있다.
위에서 가정들이 충족되었으므로 성질, 보조정리
\ref{properties-lemma-invert-s-sections}의 (2)에 의해 정수 $e \geq 1$과
절단
$s' \in \Gamma(X, \mathcal{L}^{\otimes n} \otimes f^*\mathcal{M}^{\otimes em})$
이 존재하여 $s(f^*t)^e$로 제한된다. 이 제한은 $X_{f^*t}$ 위에서의
것이다. 임의의 $b > 0$에 대해 절단 $s'' = s'(f^*t)^b$를 생각하자.
이는 $\mathcal{L}^{\otimes n} \otimes f^*\mathcal{M}^{\otimes (e + b)m}$의
절단이다. 그러면
$X_{s''} = (X_{f^*t})_s$는 $X$의 아핀 열린집합이고 $x$를 포함한다.
$b$를 택하되 $n$이 $e + b$를 나누게 하면
$\mathcal{L}^{\otimes n} \otimes f^*\mathcal{M}^{\otimes (e + b)m}$은
$n$제곱이다. 그 밑은 $\mathcal{L} \otimes f^*\mathcal{M}^{\otimes a}$이고,
이는 어떤 $a$에 대해 성립한다. 이 방법으로 충분히 큰 임의의 $a$를
얻을 수 있고, 그것은 $m$의 배수이다.
$X$가 준콤팩트이므로 이러한 아핀 열린집합 유한 개가 $X$를 덮는다.
따라서 충분히 나누어지고 충분히 큰 어떤 $a$에 대해 가역층
$\mathcal{L} \otimes f^*\mathcal{M}^{\otimes a}$는 $X$ 위에서 충분하다.
한편 $\mathcal{M}^{\otimes c}$와 그 $X$로의 역상은 모든 $c \gg 0$에
대해 대역 생성됨을 성질, 명제
\ref{properties-proposition-characterize-ample}에 의해 안다. 따라서
$\mathcal{L} \otimes f^*\mathcal{M}^{\otimes a + c}$는 $c \gg 0$에 대해
충분하다
(성질, 보조정리
\ref{properties-lemma-ample-tensor-globally-generated}). 이로써 (1)을
증명했다.

\medskip\noindent
(2)는 보조정리 \ref{morphisms-lemma-quasi-affine-O-ample}, 성질, 보조정리
\ref{properties-lemma-ample-power-ample}, 그리고 (1)에서 따른다.
\end{proof}

\begin{lemma}
\label{morphisms-lemma-ample-composition}
$g : Y \to S$와 $f : X \to Y$를 스킴의 사상이라 하자.
$\mathcal{M}$을 가역 $\mathcal{O}_Y$-가군이라 하자.
$\mathcal{L}$을 가역 $\mathcal{O}_X$-가군이라 하자.
$S$가 준콤팩트이고, $\mathcal{M}$이 $g$-충분하며,
$\mathcal{L}$이 $f$-충분하면
$\mathcal{L} \otimes f^*\mathcal{M}^{\otimes a}$는 $g \circ f$-충분하며,
이는 $a \gg 0$에 대해 성립한다.
\end{lemma}

\begin{proof}
$S = \bigcup_{i = 1, \ldots, n} V_i$를 유한 아핀 열린 덮개라 하자.
보조정리 \ref{morphisms-lemma-characterize-relatively-ample}에 의해
$\mathcal{L} \otimes f^*\mathcal{M}^{\otimes a}$가
$(g \circ f)^{-1}(V_i)$ 위에서 충분함을 각 $i = 1, \ldots, n$에 대해
증명하면 된다. 따라서 보조정리
\ref{morphisms-lemma-pullback-ample-tensor-relatively-ample}에서 결론이 따른다.
\end{proof}

\begin{lemma}
\label{morphisms-lemma-ample-base-change}
$f : X \to S$를 스킴의 사상이라 하자.
$\mathcal{L}$을 가역 $\mathcal{O}_X$-가군이라 하자.
$S' \to S$를 스킴의 사상이라 하자. $f' : X' \to S'$를 $f$의
기저변환이라 하고, $\mathcal{L}'$로 $\mathcal{L}$의 $X'$로의 역상을
나타내자. $\mathcal{L}$이 $f$-충분하면 $\mathcal{L}'$은 $f'$-충분하다.
\end{lemma}

\begin{proof}
보조정리 \ref{morphisms-lemma-characterize-relatively-ample}에 의해 아핀 열린 덮개
$S' = \bigcup U'_i$를 찾아서, $\mathcal{L}'$을
$(f')^{-1}(U_i')$로 제한한 것이 모든 $i$에 대해 충분하게 하면 된다.
$U'_i$가 아핀
열린집합 $U_i \subset S$ 안으로 사상되도록 택할 수 있다. 이 경우 사상
$(f')^{-1}(U'_i) \to f^{-1}(U_i)$는 아핀 사상 $U'_i \to U_i$의
기저변환이므로 아핀이다(보조정리 \ref{morphisms-lemma-base-change-affine}). 따라서
$\mathcal{L}'|_{(f')^{-1}(U'_i)}$은 보조정리
\ref{morphisms-lemma-pullback-ample-tensor-relatively-ample}에 의해 충분하다.
\end{proof}

\begin{lemma}
\label{morphisms-lemma-ample-permanence}
$g : Y \to S$와 $f : X \to Y$를 스킴의 사상이라 하자.
$\mathcal{L}$을 가역 $\mathcal{O}_X$-가군이라 하자.
$\mathcal{L}$이 $g \circ f$-충분하고 $f$가 준콤팩트이면\footnote{$g$가
준분리이면 이는 스킴, 보조정리
\ref{schemes-lemma-quasi-compact-permanence}에서 따른다.}
$\mathcal{L}$은 $f$-충분하다.
\end{lemma}

\begin{proof}
$f$가 준콤팩트이고 $\mathcal{L}$이 $g \circ f$-충분하다고 가정하자.
$U \subset S$를 아핀 열린집합이라 하고, $V \subset Y$를
$g(V) \subset U$인 아핀 열린집합이라 하자. 그러면 가정에 의해
$\mathcal{L}|_{(g \circ f)^{-1}(U)}$는 $(g \circ f)^{-1}(U)$ 위에서
충분하다. $f^{-1}(V) \subset (g \circ f)^{-1}(U)$이므로 성질, 보조정리
\ref{properties-lemma-ample-on-locally-closed}에 의해
$\mathcal{L}|_{f^{-1}(V)}$는 $f^{-1}(V)$ 위에서 충분하다. 실제로
$f^{-1}(V) \to (g \circ f)^{-1}(U)$는 스킴, 보조정리
\ref{schemes-lemma-quasi-compact-permanence}에 의해 준콤팩트 열린
몰입이다. 이는 $(g \circ f)^{-1}(U)$가 분리되어 있고(성질, 보조정리
\ref{properties-lemma-ample-separated}), $f^{-1}(V)$가 준콤팩트이기
때문이다($f$가 준콤팩트이다). 따라서
보조정리 \ref{morphisms-lemma-characterize-relatively-ample}에 의해
$\mathcal{L}$은 $f$-충분하다.
\end{proof}







\section{매우 충분한 층}
\label{morphisms-section-very-ample}

\noindent
준연접층 $\mathcal{E}$가 스킴 $S$ 위에 주어질 때, $\mathcal{E}$에 딸린
{\it 사영다발}은 사상 $\mathbf{P}(\mathcal{E}) \to S$임을 상기하자.
여기서
$\mathbf{P}(\mathcal{E}) = \underline{\text{Proj}}_S(\text{Sym}(\mathcal{E}))$
이다. 구성, 정의 \ref{constructions-definition-projective-bundle}을 보라.

\begin{definition}
\label{morphisms-definition-very-ample}
$f : X \to S$를 스킴의 사상이라 하자.
$\mathcal{L}$을 가역 $\mathcal{O}_X$-가군이라 하자.
$\mathcal{L}$을 {\it 상대적으로 매우 충분하다}, 더 정확히는
{\it $f$-상대적으로 매우 충분하다}, 또는 {\it $X/S$ 위에서 매우
충분하다}, 또는 {\it $f$-매우 충분하다}고 하는 것은 다음을 뜻한다.
준연접 $\mathcal{O}_S$-가군 $\mathcal{E}$와 몰입
$i : X \to \mathbf{P}(\mathcal{E})$가 $S$ 위에서 존재하여
$\mathcal{L} \cong i^*\mathcal{O}_{\mathbf{P}(\mathcal{E})}(1)$이다.
\end{definition}



\noindent
이 정의에는 준콤팩트성 가정이 없으므로, 일반적으로 상대적으로 매우
충분한 가역층이 상대적으로 충분한 가역층이라고 할 수는 없다.

\begin{lemma}
\label{morphisms-lemma-ample-very-ample}
\begin{reference}
\cite[II, Proposition 4.6.2]{EGA}
\end{reference}
$f : X \to S$를 스킴의 사상이라 하자.
$\mathcal{L}$을 가역 $\mathcal{O}_X$-가군이라 하자.
$f$가 준콤팩트이고 $\mathcal{L}$이 상대적으로 매우 충분한 가역층이면,
$\mathcal{L}$은 상대적으로 충분한 가역층이다.
\end{lemma}

\begin{proof}
정의에 의해 준연접 $\mathcal{O}_S$-가군 $\mathcal{E}$와 몰입
$i : X \to \mathbf{P}(\mathcal{E})$가 $S$ 위에서 존재하여
$\mathcal{L} \cong i^*\mathcal{O}_{\mathbf{P}(\mathcal{E})}(1)$이다.
$\mathcal{A} = \text{Sym}(\mathcal{E})$로 놓으면 정의에 의해
$\mathbf{P}(\mathcal{E}) = \underline{\text{Proj}}_S(\mathcal{A})$이다.
등급 $\mathcal{O}_S$-대수 $\mathcal{A}$에는 다음 사상이 주어진다.
$$
\psi :
\mathcal{A} \to
\bigoplus\nolimits_{n \geq 0}
\pi_*\mathcal{O}_{\mathbf{P}(\mathcal{E})}(n) \to
\bigoplus\nolimits_{n \geq 0}
f_*\mathcal{L}^{\otimes n}
$$
여기서 두 번째 화살표는 동일시
$\mathcal{L} \cong i^*\mathcal{O}_{\mathbf{P}(\mathcal{E})}(1)$을 사용한다.
$f_*$와 $f^*$의 수반성에 의해 사상
$\psi : f^*\mathcal{A} \to \bigoplus_{n \geq 0}\mathcal{L}^{\otimes n}$을
얻는다. 이 사상에 딸린 사상 $r_{\mathcal{L}, \psi}$가 정확히 몰입
$i$임을 확인하는 일은 생략한다. 따라서 보조정리
\ref{morphisms-lemma-characterize-relatively-ample}의 (6)에서 결과가 따른다.
\end{proof}

\noindent
이 보조정리의 올바른 역을 얻기 위해, 상대적으로 충분한 가역층
$\mathcal{L}$이 주어졌을 때 정수 $n \geq 1$이 존재하여
$\mathcal{L}^{\otimes n}$이 상대적으로 매우 충분하게 되는지를 묻자. 일반적으로 이는
거짓이다. 이를 막는 몇 가지 이유가 있다.
\begin{enumerate}
\item $S$가 아핀인 경우에도 작동하는 유한 정수 $n$이 없을 수 있다.
이는 $X \to S$가 유한형이 아니기 때문이다. 예 \ref{morphisms-example-not-finite-type-proj}를
보라.
\item 밑이 준콤팩트가 아니면 $f$가 유한형이어도 결과가 성립하지 않을
수 있다. 실제로 체 $k$가 주어지면 유한형 스킴 $X_d$와 충분한
가역층이 존재하는데, 이 스킴은 $k$ 위에 있고 그 가역층은
$\mathcal{O}_{X_d}(1)$이다. 이때 $\mathcal{O}_{X_d}(1)$의 텐서
거듭제곱 가운데 매우 충분한 최소 거듭제곱은 $d$제곱이다. 예
\ref{morphisms-example-not-bounded}를 보라. $f$를 스킴들 $X_d$의 서로소 합집합에서
$\Spec(k)$의 사본들의 서로소 합집합으로 가는 사상으로 택하면 예를
얻는다.
\end{enumerate}
우리의 역을 보려면 아래의 보조정리
\ref{morphisms-lemma-finite-type-ample-very-ample}을 보라. 이를 증명하기 전에 몇
가지 준비를 하겠다.

\begin{example}
\label{morphisms-example-very-ample}
$S$를 스킴이라 하자. $\mathcal{A}$를 준연접 등급
$\mathcal{O}_S$-대수라 하자. 이 대수는 $\mathcal{A}_1$로
$\mathcal{A}_0$ 위에서 생성된다고 가정한다.
$X = \underline{\text{Proj}}_S(\mathcal{A})$로 놓자. 이 경우
$\mathcal{O}_X(1)$은 가역이고 $X/S$ 위에서 매우 충분하다. 실제로 등급
$\mathcal{O}_S$-대수의 사상
$$
\text{Sym}_{\mathcal{O}_X}^*(\mathcal{A}_1)
\longrightarrow
\mathcal{A}
$$
에 딸린 사상은 닫힌 몰입 $X \to \mathbf{P}(\mathcal{A}_1)$이고,
$\mathcal{O}_{\mathbf{P}(\mathcal{A}_1)}(1)$을 $\mathcal{O}_X(1)$로
당긴다. 구성, 보조정리
\ref{constructions-lemma-surjective-generated-degree-1-map-relative-proj}를
보라.
\end{example}

\begin{example}
\label{morphisms-example-not-finite-type-proj}
$k$를 체라 하자. 다음 등급 $k$-대수를 생각하자.
$$
A = k[U, V, Z_1, Z_2, Z_3, \ldots]/I
\quad
\text{with}
\quad
I = (U^2 - Z_1^2, U^4 - Z_2^2, U^6 - Z_3^2, \ldots)
$$
등급은 $\deg(U) = \deg(V) = \deg(Z_1) = 1$ 및
$\deg(Z_d) = d$로 주어진다. $X = \text{Proj}(A)$는 $D_{+}(U)$와
$D_{+}(V)$로 덮임에 유의하자. 따라서 층 $\mathcal{O}_X(n)$은 모두
가역이고 $\mathcal{O}_X(1)^{\otimes n}$과 동형이다. 특히
$\mathcal{O}_X(1)$은 충분하고 $f$-충분하다. 여기서 사상은
$f : X \to \Spec(k)$이다. $\mathcal{O}_X(1)$의 어떠한 거듭제곱도 $f$-상대적으로
매우 충분하지 않다고 주장한다. 실제로
$\Gamma(X, \mathcal{O}_X(n))$은 차수 $n$ 성분이며, 이는 대수 $A$의
성분임을 쉽게 알
수 있다. 따라서 $\mathcal{O}_X(n)$이 매우 충분하다면 $X$는 $k$ 위의
사영공간의 닫힌 부분스킴이고, 따라서 $k$ 위의 유한형이다. 한편
$D_{+}(V)$는
$k[t, t_1, t_2, \ldots]/(t^2 - t_1^2, t^4 - t_2^2, t^6 - t_3^2, \ldots)$
의 스펙트럼이며, 이는 $k$ 위의 유한형이 아니다.
\end{example}

\begin{example}
\label{morphisms-example-not-bounded}
$k$를 무한체라 하자. $\lambda_1, \lambda_2, \lambda_3, \ldots$를
$k^*$의 서로 다른 원소들이라 하자. (이는 엄밀히 필요하지 않으며,
실제로 모든 $\lambda_i$가 $1$과 같아도 예는 완전히 잘 작동한다.)
다음 등급 $k$-대수를 생각하자.
$$
A_d = k[U, V, Z]/I_d
\quad
\text{with}
\quad
I_d = (Z^2 - \prod\nolimits_{i = 1}^{2d} (U - \lambda_i V)).
$$
등급은 $\deg(U) = \deg(V) = 1$ 및 $\deg(Z) = d$로 주어진다.
그러면 $X_d = \text{Proj}(A_d)$는 충분한 가역층
$\mathcal{O}_{X_d}(1)$을 갖는다. $\mathcal{O}_{X_d}(n)$이 매우
충분하면 $n \geq d$라고 주장한다. 그 이유는 $Z$의 차수가 $d$이고,
따라서 $\Gamma(X_d, \mathcal{O}_{X_d}(n)) =
k[U, V]_n$이며 이는 $n < d$일 때 성립하기 때문이다.
세부사항은 생략한다.
\end{example}

\begin{lemma}
\label{morphisms-lemma-relatively-very-ample-separated}
$f : X \to S$를 스킴의 사상이라 하자.
$\mathcal{L}$을 $X$ 위의 가역층이라 하자.
$\mathcal{L}$이 $X/S$ 위에서 상대적으로 매우 충분하면 $f$는 분리되어
있다.
\end{lemma}

\begin{proof}
분리성은 밑에서 국소적이다(스킴, 절
\ref{schemes-section-separation-axioms}을 보라). 몰입은 분리되어 있다
(스킴, 보조정리 \ref{schemes-lemma-immersions-monomorphisms}을 보라).
따라서 국소적으로 $X$가 등급환의 동차 스펙트럼 안으로 몰입하고 이
동차 스펙트럼이 분리되어 있으므로 보조정리가 따른다. 구성, 보조정리
\ref{constructions-lemma-proj-separated}를 보라.
\end{proof}

\begin{lemma}
\label{morphisms-lemma-relatively-very-ample}
$f : X \to S$를 스킴의 사상이라 하자.
$\mathcal{L}$을 $X$ 위의 가역층이라 하자.
$f$가 준콤팩트라고 가정하자. 다음은 동치이다.
\begin{enumerate}
\item $\mathcal{L}$은 $X/S$ 위에서 상대적으로 매우 충분하다.
\item 열린 덮개 $S = \bigcup V_j$가 존재하여
$\mathcal{L}|_{f^{-1}(V_j)}$는 $f^{-1}(V_j)/V_j$ 위에서 상대적으로 매우
충분하며, 이는 모든 $j$에 대해 성립한다.
\item 준연접 등급 $\mathcal{O}_S$-대수층 $\mathcal{A}$가 존재하여 차수
$1$에서 $\mathcal{O}_S$ 위로 생성되고, 등급 $\mathcal{O}_X$-대수의 사상
$\psi : f^*\mathcal{A} \to \bigoplus_{n \geq 0} \mathcal{L}^{\otimes n}$이
존재하여, $f^*\mathcal{A}_1 \to \mathcal{L}$은 전사이고 딸린 사상
$r_{\mathcal{L}, \psi} : X \to \underline{\text{Proj}}_S(\mathcal{A})$는
몰입이다.
\item $f$는 준분리이고, 표준 사상
$\psi : f^*f_*\mathcal{L} \to \mathcal{L}$은 전사이며, 딸린 사상
$r_{\mathcal{L}, \psi} : X \to \mathbf{P}(f_*\mathcal{L})$는 몰입이다.
\end{enumerate}
\end{lemma}

\begin{proof}
(1)이 (2)를 함의함은 명백하다. (4)가 (1)을 함의함도 명백하다. (4)의
준분리 가정은 스킴, 보조정리
\ref{schemes-lemma-push-forward-quasi-coherent}를 통해
$f_*\mathcal{L}$이 준연접임을 보장하는 데 쓰인다.

\medskip\noindent
(2)를 가정하자. (4)를 증명하겠다. $S = \bigcup V_j$를 (2)에 나오는
열린 덮개라 하자. $X_j = f^{-1}(V_j)$로 놓고 $f_j : X_j \to V_j$를
$f$의 제한이라 하자. 보조정리
\ref{morphisms-lemma-relatively-very-ample-separated}에 의해 $f$가 분리되어 있음을
안다(분리성은 밑에서 국소적이다). 가정에 의해 준연접
$\mathcal{O}_{V_j}$-가군 $\mathcal{E}_j$와 몰입
$i_j : X_j \to \mathbf{P}(\mathcal{E}_j)$가 존재하여
$\mathcal{L}|_{X_j} \cong i_j^*\mathcal{O}_{\mathbf{P}(\mathcal{E}_j)}(1)$이다.
사상 $i_j$는 전사사상 $f_j^*\mathcal{E}_j \to \mathcal{L}|_{X_j}$에
대응한다. 구성, 절 \ref{constructions-section-projective-bundle}을 보라.
이 사상은 사상 $\mathcal{E}_j \to f_*\mathcal{L}|_{V_j}$에 수반되며,
합성
$$
f_j^*\mathcal{E}_j \to (f^*f_*\mathcal{L})|_{X_j} \to \mathcal{L}|_{X_j}
$$
은 전사이다. 따라서
$\psi : f^*f_*\mathcal{L} \to \mathcal{L}$은 전사이다. 딸린 사상을
$r_{\mathcal{L}, \psi} : X \to \mathbf{P}(f_*\mathcal{L})$라 하자.
아직 $r_{\mathcal{L}, \psi}$가 몰입임을 보여야 한다. 독자에게 직접
증명할 것을 권한다. $\mathcal{O}_{V_j}$-가군 사상
$\mathcal{E}_j \to f_*\mathcal{L}|_{V_j}$는 대칭대수의 준동형을 결정하고,
이는 다시 사상
$$
\mathbf{P}(f_*\mathcal{L}|_{V_j}) \supset U_j \longrightarrow
\mathbf{P}(\mathcal{E}_j)
$$
을 정의한다. 여기서 $U_j$는 구성, 보조정리
\ref{constructions-lemma-morphism-relative-proj}의 열린 부분스킴이다.
$\psi$와 $\mathcal{E}_j \to f_*\mathcal{L}|_{V_j}$의 호환성에 의해
$r_{\mathcal{L}, \psi}(X_j) \subset U_j$이고 다음 분해가 존재한다.
$$
\xymatrix{
X_j \ar[r]^-{r_{\mathcal{L}, \psi}} & U_j \ar[r] &
\mathbf{P}(\mathcal{E}_j)
}
$$
확인은 생략한다. 이로부터 $r_{\mathcal{L}, \psi}$가 몰입임을 알 수 있다.

\medskip\noindent
이 시점에서 (1), (2), (4)가 동치임을 안다. (4)가 (3)을 함의함은
명백하다. (3)을 가정하자. (1)을 증명하겠다. $\mathcal{A}$를 준연접
등급 $\mathcal{O}_S$-대수층이라 하자. 이는 차수 $1$에서
$\mathcal{O}_S$ 위로 생성된다. 등급 $\mathcal{O}_S$-대수의 사상
$\text{Sym}(\mathcal{A}_1) \to \mathcal{A}$를 생각하자. 이는 가정에 의해
전사이므로 닫힌 몰입
$$
\underline{\text{Proj}}_S(\mathcal{A})
\longrightarrow
\mathbf{P}(\mathcal{A}_1)
$$
을 유도하며, 이는 $\mathcal{O}(1)$을 $\mathcal{O}(1)$로 당긴다.
구성, 보조정리
\ref{constructions-lemma-surjective-generated-degree-1-map-relative-proj}를
보라. 따라서 (3)이 (1)을 함의함은 명백하다.
\end{proof}

\begin{lemma}
\label{morphisms-lemma-very-ample-base-change}
$f : X \to S$를 스킴의 사상이라 하자.
$\mathcal{L}$을 가역 $\mathcal{O}_X$-가군이라 하자.
$S' \to S$를 스킴의 사상이라 하자. $f' : X' \to S'$를 $f$의
기저변환이라 하고, $\mathcal{L}'$로 $\mathcal{L}$의 $X'$로의 역상을
나타내자. $\mathcal{L}$이 $f$-매우 충분하면 $\mathcal{L}'$은
$f'$-매우 충분하다.
\end{lemma}

\begin{proof}
정의 \ref{morphisms-definition-very-ample}에 의해 준연접 $\mathcal{O}_S$-가군
$\mathcal{E}$와 몰입 $i : X \to \mathbf{P}(\mathcal{E})$가 $S$ 위에서
존재하여
$\mathcal{L} \cong i^*\mathcal{O}_{\mathbf{P}(\mathcal{E})}(1)$이다.
$\mathbf{P}(\mathcal{E})$를 $S'$로 기저변환한 것은 역상
$\mathcal{E}'$에 딸린 사영다발이다. 이 층은 $\mathcal{E}$의 역상이고,
$\mathcal{O}_{\mathbf{P}(\mathcal{E})}(1)$의 역상은
$\mathcal{O}_{\mathbf{P}(\mathcal{E}')}(1)$이다. 구성, 보조정리
\ref{constructions-lemma-relative-proj-base-change}를 보라. 마지막으로
몰입의 기저변환은 몰입이다(스킴, 보조정리
\ref{schemes-lemma-base-change-immersion}).
\end{proof}







\section{유한형 사상에 대한 충분층과 매우 충분층}
\label{morphisms-section-ample-finite-type}

\noindent
실제로 이 절의 내용 대부분은 아직 정의하지 않은 (준)사영 사상의
개념에 관한 것이다.

\begin{lemma}
\label{morphisms-lemma-very-ample-finite-type-over-affine}
$f : X \to S$를 스킴의 사상이라 하자.
$\mathcal{L}$을 $X$ 위의 가역층이라 하자.
다음을 가정하자.
\begin{enumerate}
\item 가역층 $\mathcal{L}$은 $X/S$ 위에서 매우 충분하다.
\item 사상 $X \to S$는 유한형이다.
\item $S$는 아핀이다.
\end{enumerate}
그러면 $n \geq 0$과 몰입 $i : X \to \mathbf{P}^n_S$가 $S$ 위에서
존재하여 $\mathcal{L} \cong i^*\mathcal{O}_{\mathbf{P}^n_S}(1)$이다.
\end{lemma}

\begin{proof}
(1), (2), (3)을 가정하자. 조건 (3)은 $S = \Spec(R)$임을 뜻하며,
여기서 $R$은 어떤 환이다. 조건 (1)은 정의에 의해 준연접
$\mathcal{O}_S$-가군 $\mathcal{E}$와 몰입
$\alpha : X \to \mathbf{P}(\mathcal{E})$가 존재하여
$\mathcal{L} = \alpha^*\mathcal{O}_{\mathbf{P}(\mathcal{E})}(1)$임을
뜻한다. $\mathcal{E} = \widetilde{M}$로 쓰자. 여기서 $R$ 위의 어떤
가군을 $M$으로 나타냈다.
그러면
$$
\mathbf{P}(\mathcal{E}) = \text{Proj}(\text{Sym}_R(M)).
$$
$\alpha$는 몰입이고 $\text{Proj}(\text{Sym}_R(M))$의 위상은 표준
열린집합 $D_{+}(f)$, $f \in \text{Sym}_R^d(M)$, $d \geq 1$로
생성되므로, 각 $x \in X$에 대해 $f \in \text{Sym}_R^d(M)$,
$d \geq 1$을 찾아서 $\alpha(x) \in D_{+}(f)$이고
$$
\alpha|_{\alpha^{-1}(D_{+}(f))} : \alpha^{-1}(D_{+}(f)) \to D_{+}(f)
$$
가 닫힌 몰입이 되게 할 수 있다. 조건 (2)에 의해 $X$는
준콤팩트이다. 따라서 원소들의 유한 모음
$f_j \in \text{Sym}_R^{d_j}(M)$, $d_j \geq 1$을 찾아서 각
$f = f_j$에 대해 위에 표시한 사상이 닫힌 몰입이고
$\alpha(X) \subset \bigcup D_{+}(f_j)$가 되게 할 수 있다.
$U_j = \alpha^{-1}(D_{+}(f_j))$로 쓰자. $U_j$는 아핀 스킴
$D_{+}(f_j)$의 닫힌 부분스킴이므로 아핀임에 유의하자.
$U_j = \Spec(A_j)$로 쓰자. 조건 (2)는 또한 $A_j$가 $R$ 위의
유한형임을 뜻한다. 보조정리
\ref{morphisms-lemma-locally-finite-type-characterize}를 보라. 유한 개의
$x_{j, k} \in A_j$를 택하여 $A_j$를 $R$-대수로 생성하게 하자.
$\alpha|_{U_j}$가 닫힌 몰입이므로 $x_{j, k}$는 다음 원소의 상이다.
$$
f_{j, k}/f_j^{e_{j, k}} \in \text{Sym}_R(M)_{(f_j)}
=
\Gamma(D_{+}(f_j), \mathcal{O}_{\text{Proj}(\text{Sym}_R(M))}).
$$
마지막으로 $n \geq 1$과 원소 $y_0, \ldots, y_n \in M$을 택하여
다항식 $f_j, f_{j, k} \in \text{Sym}_R(M)$ 각각이 $y_t$들의
다항식이 되게 하자. 그 계수는 $R$에 속한다. 등급환 준동형
$$
\psi : R[Y_0, \ldots, Y_n] \longrightarrow \text{Sym}_R(M),
\quad Y_i \longmapsto y_i.
$$
을 생각하자. $F_j$, $F_{j, k}$로 $R[Y_0, \ldots, Y_n]$의 원소 가운데
$\psi(F_j) = f_j$ 및 $\psi(F_{j, k}) = f_{j, k}$를 만족하는 것들을
나타내자. 구성, 보조정리 \ref{constructions-lemma-morphism-proj}에 의해
열린 부분스킴
$$
U(\psi) \subset \text{Proj}(\text{Sym}_R(M))
$$
과 사상 $r_\psi : U(\psi) \to \mathbf{P}^n_R$을 얻는다. 이 사상은
$r_\psi^{-1}(D_{+}(F_j)) = D_{+}(f_j)$를 만족하므로
$\alpha(X) \subset U(\psi)$임을 안다. 더욱이
$$
i = r_\psi \circ \alpha : X \longrightarrow \mathbf{P}^n_R
$$
는 여전히 몰입이다. 실제로 구성에 의해
$i^\sharp(F_{j, k}/F_j^{e_{j, k}}) = x_{j, k} \in
A_j = \Gamma(U_j, \mathcal{O}_X)$이기 때문이다. 또한 사상 $r_\psi$에는
사상
$\theta : r_\psi^*\mathcal{O}_{\mathbf{P}^n_R}(1)
\to \mathcal{O}_{\text{Proj}(\text{Sym}_R(M))}(1)|_{U(\psi)}$가 주어지고,
이 경우 이는 동형사상이다($\theta$의 구성은 위에서 인용한 보조정리를
보라. 일부 세부사항은 생략한다). 원래 사상 $\alpha$는
$\mathcal{L} = \alpha^*\mathcal{O}_{\text{Proj}(\text{Sym}_R(M))}(1)$을
만족한다고 가정했으므로 결론을 얻는다.
\end{proof}

\begin{lemma}
\label{morphisms-lemma-quasi-affine-finite-type-over-S}
$\pi : X \to S$를 스킴의 사상이라 하자.
$X$가 준아핀이고 $\pi$가 국소적으로 유한형이라고 가정하자. 그러면
$n \geq 0$과 몰입 $i : X \to \mathbf{A}^n_S$가 $S$ 위에서 존재한다.
\end{lemma}

\begin{proof}
$A = \Gamma(X, \mathcal{O}_X)$로 놓자. 가정에 의해 $X$는 준콤팩트이고
$\Spec(A)$의 열린 부분스킴과 동일시된다. 성질, 보조정리
\ref{properties-lemma-quasi-affine}을 보라. 더욱이 $X_f$ 가운데
$f \in A$이고 $X_f$가 아핀인 것들의 모음은 $X$의 위상의 기저를 이룬다.
성질, 보조정리 \ref{properties-lemma-quasi-affine}의 증명을 보라. 따라서
유한 개의 $f_j \in A$, $j = 1, \ldots, m$을 찾아서
$X = \bigcup X_{f_j}$이고, $\pi(X_{f_j}) \subset V_j$가 되게 할 수 있다.
여기서 $V_j \subset S$는 어떤 아핀 열린집합이다. 보조정리
\ref{morphisms-lemma-locally-finite-type-characterize}에 의해 환 준동형
$\mathcal{O}(V_j) \to \mathcal{O}(X_{f_j}) = A_{f_j}$는 유한형이다.
따라서 $a_1, \ldots, a_N \in A$를 택하여 원소
$a_1, \ldots, a_N, 1/f_j$가 $A_{f_j}$를 $\mathcal{O}(V_j)$ 위에서
각 $j$에 대해 생성하게 할 수 있다. $n = m + N$으로 놓고
$$
i : X \longrightarrow \mathbf{A}^n_S
$$
를 전역 절단 $f_1, \ldots, f_m, a_1, \ldots, a_N$으로 주어진
사상이라 하자. 이들은 $X$의 구조층의 절단이다.
$D(x_j) \subset \mathbf{A}^n_S$를 $j$번째
좌표함수가 영이 아닌 열린 부분스킴이라 하자. 그러면
$1 \leq j \leq m$에 대해 $i^{-1}(D(x_j)) = X_{f_j}$이고, 유도된 사상
$X_{f_j} \to D(x_j)$는 아핀 열린집합
$\Spec(\mathcal{O}(V_j)[x_1, \ldots, x_n, 1/x_j])$를 통하여 분해된다.
이 열린집합은 $D(x_j)$ 안에 있다.
환 준동형
$\mathcal{O}(V_j)[x_1, \ldots, x_n, 1/x_j] \to A_{f_j}$가 구성에 의해
전사이므로 $i^{-1}(D(x_j)) \to D(x_j)$는 원하는 대로 몰입이다.
\end{proof}

\begin{lemma}
\label{morphisms-lemma-quasi-projective-finite-type-over-S}
$f : X \to S$를 스킴의 사상이라 하자.
$\mathcal{L}$을 $X$ 위의 가역층이라 하자.
다음을 가정하자.
\begin{enumerate}
\item 가역층 $\mathcal{L}$은 $X$ 위에서 충분하다.
\item 사상 $X \to S$는 국소적으로 유한형이다.
\end{enumerate}
그러면 $d_0 \geq 1$이 존재하여 모든 $d \geq d_0$에 대해
$n \geq 0$과 몰입 $i : X \to \mathbf{P}^n_S$가 $S$ 위에서 존재하고
$\mathcal{L}^{\otimes d} \cong i^*\mathcal{O}_{\mathbf{P}^n_S}(1)$이다.
\end{lemma}

\begin{proof}
다음과 같이 놓자.
$A = \Gamma_*(X, \mathcal{L}) =
\bigoplus_{d \geq 0} \Gamma(X, \mathcal{L}^{\otimes d})$.
성질, 명제 \ref{properties-proposition-characterize-ample}에 의해
아핀 열린집합 $X_a$ 가운데 $a \in A_{+}$가 동차인 것들의 모음은 $X$의
위상의 기저를 이룬다. 따라서 유한 개의 원소
$a_0, \ldots, a_n \in A_{+}$를 찾아서 다음을 만족하게 할 수 있다.
\begin{enumerate}
\item $X = \bigcup_{i = 0, \ldots, n} X_{a_i}$이다.
\item 각각의 $X_{a_i}$는 아핀이다.
\item 각각의 $X_{a_i}$는 아핀 열린집합 $V_i \subset S$ 안으로 사상된다.
\end{enumerate}
보조정리 \ref{morphisms-lemma-locally-finite-type-characterize}에 의해 환 준동형
$\mathcal{O}_S(V_i) \to \mathcal{O}_X(X_{a_i})$가 유한형임을 안다.
따라서 유한 개의 원소
$f_{ij} \in \mathcal{O}_X(X_{a_i})$, $j = 1, \ldots, n_i$를 찾아서
$\mathcal{O}_X(X_{a_i})$를 $\mathcal{O}_S(V_i)$-대수로 생성하게 할 수
있다. 성질, 보조정리 \ref{properties-lemma-invert-s-sections}에 의해
각 $f_{ij}$를 $a_{ij}/a_i^{e_{ij}}$로 쓸 수 있으며, 여기서
$a_{ij} \in A_{+}$는 어떤 동차 원소이다. $N$을 모든 원소 $a_i$, $a_{ij}$의
차수들의 양의 공배수라 하자. 원소
$$
a_i^{N/\deg(a_i)}, \ a_{ij}a_i^{(N/\deg(a_i)) - e_{ij}} \in A_N.
$$
를 생각하자. 구성에 의해 이들은 가역층 $\mathcal{L}^{\otimes N}$을
$X$ 위에서 생성한다. 따라서 이들은 사상
$$
j : X \longrightarrow
\mathbf{P}_S^{m}
\quad
\text{with } m = n + \sum n_i
$$
을 준다. 이 사상은 $S$ 위에 있다. 구성, 보조정리
\ref{constructions-lemma-projective-space}와 정의
\ref{constructions-definition-projective-space}를 보라. 더욱이
$j^*\mathcal{O}_{\mathbf{P}_S}(1) = \mathcal{L}^{\otimes N}$이다.
동차좌표를 $T_0, \ldots, T_n, T_{ij}$로 명명하자. 이는
$T_0, \ldots, T_m$ 대신 쓰는 기호이다. $i = 0, \ldots, n$에 대해
$j^{-1}(D_{+}(T_i)) = X_{a_i}$이다. 또한 원소 $T_{ij}/T_i$를
$j^\sharp$로 당기면 원소 $f_{ij} \in \mathcal{O}_X(X_{a_i})$를 얻는다.
따라서 사상 $j$를 $X_{a_i}$로 제한하면 $X_{a_i}$에서 아핀 열린집합
$D_{+}(T_i) \cap \mathbf{P}^m_{V_i}$로 가는 닫힌 몰입을 얻는다.
이 열린집합은 $\mathbf{P}^N_S$ 안에 있다.
그러므로 사상 $j$는 몰입이다. 이는 보조정리가 어떤 $d$와 $n$에 대해
성립함을 뜻하며, 사실상 모든 응용에는 이것으로 충분하다.

\medskip\noindent
이는 하나의 $d_2 \geq 1$(즉, 위의 $d_2 = N$)과 어떤 $m \geq 0$에 대해
어떤 몰입 $j : X \to \mathbf{P}^m_S$가 존재함을 증명한다. 이 몰입은
전역 절단 $s'_0, \ldots, s'_m \in \Gamma(X, \mathcal{L}^{\otimes d_2})$로
주어진다.
성질, 명제 \ref{properties-proposition-characterize-ample}에 의해 정수
$d_1$이 존재하여 $\mathcal{L}^{\otimes d}$가 모든 $d \geq d_1$에 대해
대역 생성됨을 안다. $d_0 = d_1 + d_2$로 놓자. 이 $d_0$의 값에 대해
보조정리가 성립한다고 주장한다. 실제로 정수 $d \geq d_0$가 주어지면
절단 $s''_1, \ldots, s''_t
\in \Gamma(X, \mathcal{L}^{\otimes d - d_2})$를 택하여
$\mathcal{L}^{\otimes d - d_2}$를 $X$ 위에서 생성하게 할 수 있다.
$k = (m + 1)t$로 놓고 $s_0, \ldots, s_k$로 절단들의 모음
$s'_\alpha s''_\beta$, $\alpha = 0, \ldots, m$,
$\beta = 1, \ldots, t$를 나타내자. 이들은 $\mathcal{L}^{\otimes d}$를
$X$ 위에서 생성하므로 사상
$$
i : X \longrightarrow \mathbf{P}^{k - 1}_S
$$
을 정의하며, 이는
$i^*\mathcal{O}_{\mathbf{P}^n_S}(1) \cong \mathcal{L}^{\otimes d}$를
만족한다. $i$가 몰입임을 보이려면 $i$가 다음 합성임을 관찰하라.
$$
X \longrightarrow \mathbf{P}^m_S \times_S \mathbf{P}^{t - 1}_S
\longrightarrow \mathbf{P}^{k - 1}_S
$$
여기서 첫 번째 사상은 $(j, j')$이고, $j'$는
$s''_1, \ldots, s''_t$로 주어지며, 두 번째 사상은 세그레 매장이다
(구성, 보조정리 \ref{constructions-lemma-segre-embedding}).
$j$가 몰입이므로 $(j, j')$도 몰입이다(보조정리
\ref{morphisms-lemma-immersion-permanence}를
$X \to \mathbf{P}^m_S \times_S \mathbf{P}^{t - 1}_S
\to \mathbf{P}^m_S$에 적용하라). 따라서 $i$는 몰입들의 합성이므로
몰입이다(스킴, 보조정리 \ref{schemes-lemma-composition-immersion}).
\end{proof}

\begin{lemma}
\label{morphisms-lemma-finite-type-over-affine-ample-very-ample}
$f : X \to S$를 스킴의 사상이라 하자.
$\mathcal{L}$을 가역 $\mathcal{O}_X$-가군이라 하자.
$S$가 아핀이고 $f$가 유한형이라고 가정하자. 다음은 동치이다.
\begin{enumerate}
\item $\mathcal{L}$은 $X$ 위에서 충분하다.
\item $\mathcal{L}$은 $f$-충분하다.
\item $\mathcal{L}^{\otimes d}$는 $f$-매우 충분하며, 이는 어떤
$d \geq 1$에 대해 성립한다.
\item $\mathcal{L}^{\otimes d}$는 $f$-매우 충분하며, 이는 모든
$d \gg 1$에 대해 성립한다.
\item 어떤 $d \geq 1$에 대해 $n \geq 1$과 몰입
$i : X \to \mathbf{P}^n_S$가 존재하여
$\mathcal{L}^{\otimes d} \cong i^*\mathcal{O}_{\mathbf{P}^n_S}(1)$이다.
\item 모든 $d \gg 1$에 대해 $n \geq 1$과 몰입
$i : X \to \mathbf{P}^n_S$가 존재하여
$\mathcal{L}^{\otimes d} \cong i^*\mathcal{O}_{\mathbf{P}^n_S}(1)$이다.
\end{enumerate}
\end{lemma}

\begin{proof}
(1)과 (2)의 동치는 보조정리 \ref{morphisms-lemma-ample-over-affine}이다.
(2) $\Rightarrow$ (6)은 보조정리
\ref{morphisms-lemma-quasi-projective-finite-type-over-S}이다. (6)이 (5)를 함의함은
자명하다. $\mathbf{P}^n_S$는 $S$ 위의 사영다발이므로(구성, 보조정리
\ref{constructions-lemma-projective-space-bundle}을 보라), 상대적으로 매우
충분한 층의 정의에서 (5)가 (3)을 함의하고 (6)이 (4)를 함의함을 안다.
(4)가 (3)을 함의함은 자명하다. 끝으로 (3)이 (2)를 함의함을 보여야
하는데, 이는 보조정리 \ref{morphisms-lemma-ample-very-ample}와
\ref{morphisms-lemma-ample-power-ample}에서 따른다.
\end{proof}

\begin{lemma}
\label{morphisms-lemma-finite-type-ample-very-ample}
$f : X \to S$를 스킴의 사상이라 하자.
$\mathcal{L}$을 가역 $\mathcal{O}_X$-가군이라 하자.
$S$가 준콤팩트이고 $f$가 유한형이라고 가정하자. 다음은 동치이다.
\begin{enumerate}
\item $\mathcal{L}$은 $f$-충분하다.
\item $\mathcal{L}^{\otimes d}$는 $f$-매우 충분하며, 이는 어떤
$d \geq 1$에 대해 성립한다.
\item $\mathcal{L}^{\otimes d}$는 $f$-매우 충분하며, 이는 모든
$d \gg 1$에 대해 성립한다.
\end{enumerate}
\end{lemma}

\begin{proof}
(3)이 (2)를 함의함은 자명하다. 유한형 사상은 정의에 의해
준콤팩트이므로, 보조정리 \ref{morphisms-lemma-ample-very-ample}는 (2)가 (1)을
함의함을 보장한다. $\mathcal{L}$이 $f$-충분하다고 가정하자. 유한 아핀
열린 덮개 $S = V_1 \cup \ldots \cup V_m$를 택하자.
$X_i = f^{-1}(V_i)$로 쓰자. 위의 보조정리
\ref{morphisms-lemma-finite-type-over-affine-ample-very-ample}에 의해 $d_0$가
존재하여 $\mathcal{L}^{\otimes d}$가 $X_i/V_i$ 위에서 상대적으로 매우
충분하며, 이는 모든 $d \geq d_0$에 대해 성립함을 안다. 따라서 보조정리
\ref{morphisms-lemma-relatively-very-ample}에 의해 (1)이 (3)을 함의한다.
\end{proof}

\noindent
다음 두 보조정리는 사상이 유한형일 때 상대적으로 매우 충분한 가역층과
상대적으로 충분한 가역층에 대해 가장 많이 쓰이고 가장 유용한
특징짓기를 준다.

\begin{lemma}
\label{morphisms-lemma-characterize-very-ample-on-finite-type}
$f : X \to S$를 스킴의 사상이라 하자.
$\mathcal{L}$을 $X$ 위의 가역층이라 하자.
$f$가 유한형이라고 가정하자. 다음은 동치이다.
\begin{enumerate}
\item $\mathcal{L}$은 $f$-상대적으로 매우 충분하다.
\item 열린 덮개 $S = \bigcup V_j$와, 각 $j$에 대해 정수 $n_j$ 및 몰입
$$
i_j :
X_j = f^{-1}(V_j) = V_j \times_S X
\longrightarrow
\mathbf{P}^{n_j}_{V_j}
$$
이 존재하여, 이 몰입은 $V_j$ 위에 있고
$\mathcal{L}|_{X_j} \cong i_j^*\mathcal{O}_{\mathbf{P}^{n_j}_{V_j}}(1)$이다.
\end{enumerate}
\end{lemma}

\begin{proof}
$S$의 아핀 열린 덮개를 택하고 $f$ 및 $\mathcal{L}$의 각 제한에 보조정리
\ref{morphisms-lemma-very-ample-finite-type-over-affine}를 적용하면 (1)이 (2)를
함의함을 안다. 보조정리 \ref{morphisms-lemma-relatively-very-ample}에 의해 (2)가
(1)을 함의함을 안다.
\end{proof}

\begin{lemma}
\label{morphisms-lemma-characterize-ample-on-finite-type}
$f : X \to S$를 스킴의 사상이라 하자.
$\mathcal{L}$을 $X$ 위의 가역층이라 하자.
$f$가 유한형이라고 가정하자. 다음은 동치이다.
\begin{enumerate}
\item $\mathcal{L}$은 $f$-상대적으로 충분하다.
\item 열린 덮개 $S = \bigcup V_j$와, 각 $j$에 대해 정수
$d_j \geq 1$, $n_j \geq 0$ 및 몰입
$$
i_j :
X_j = f^{-1}(V_j) = V_j \times_S X
\longrightarrow
\mathbf{P}^{n_j}_{V_j}
$$
이 존재하여, 이 몰입은 $V_j$ 위에 있고
$\mathcal{L}^{\otimes d_j}|_{X_j} \cong
i_j^*\mathcal{O}_{\mathbf{P}^{n_j}_{V_j}}(1)$이다.
\end{enumerate}
\end{lemma}

\begin{proof}
$S$의 아핀 열린 덮개를 택하고 $f$ 및 $\mathcal{L}$의 각 제한에 보조정리
\ref{morphisms-lemma-finite-type-over-affine-ample-very-ample}를 적용하면 (1)이
(2)를 함의함을 안다. 보조정리
\ref{morphisms-lemma-characterize-relatively-ample}에 의해 (2)가 (1)을 함의함을
안다.
\end{proof}

\begin{lemma}
\label{morphisms-lemma-invertible-add-enough-ample-very-ample}
$f : X \to S$를 스킴의 사상이라 하자.
$\mathcal{N}$, $\mathcal{L}$을 가역 $\mathcal{O}_X$-가군이라 하자.
$S$가 준콤팩트이고, $f$가 유한형이며, $\mathcal{L}$이 $f$-충분하다고
가정하자. 그러면
$\mathcal{N} \otimes_{\mathcal{O}_X} \mathcal{L}^{\otimes d}$는
$f$-매우 충분하며, 이는 모든 $d \gg 1$에 대해 성립한다.
\end{lemma}

\begin{proof}
보조정리 \ref{morphisms-lemma-characterize-very-ample-on-finite-type}에 의해 $S$가
아핀인 경우로 환원된다. 보조정리
\ref{morphisms-lemma-finite-type-over-affine-ample-very-ample}와 성질, 명제
\ref{properties-proposition-characterize-ample}을 결합하면 정수 $d_0$를
찾아서 $\mathcal{N} \otimes \mathcal{L}^{\otimes d_0}$이 대역 생성되게
할 수 있다. 이를 생성하는 전역 절단
$s_0, \ldots, s_n$을 $\mathcal{N} \otimes \mathcal{L}^{\otimes d_0}$에서
택하자. 이는 사상 $j : X \to \mathbf{P}^n_S$를 결정하며, 이 사상은
$S$ 위에 있다.
보조정리 \ref{morphisms-lemma-finite-type-over-affine-ample-very-ample}에 의해 정수
$d_1$도 택할 수 있어서 모든 $d \geq d_1$에 대해 절단
$t_{d, 0}, \ldots, t_{d, n(d)}$가 존재하고, 이들은
$\mathcal{L}^{\otimes d}$을 생성하며 다음 몰입을 정의한다.
$$
j_d = \varphi_{\mathcal{L}^{\otimes d}, t_{d, 0}, \ldots, t_{d, n(d)}} :
X
\longrightarrow
\mathbf{P}^{n(d)}_S
$$
이 몰입은 $S$ 위에 있다. 그러면 $d \geq d_0 + d_1$에 대해 사상
$$
\varphi_{\mathcal{N} \otimes \mathcal{L}^{\otimes d},
s_j \otimes t_{d - d_0, i}} :
X
\longrightarrow
\mathbf{P}^{(n + 1)(n(d - d_0) + 1) - 1}_S
$$
을 생각할 수 있다. 이 사상은 다음 합성이므로 몰입이다.
$$
X \to \mathbf{P}^n_S \times_S \mathbf{P}^{n(d - d_0)}_S
\to \mathbf{P}^{(n + 1)(n(d - d_0) + 1) - 1}_S
$$
여기서 첫 번째 사상은 $(j, j_{d - d_0})$이고 두 번째 사상은 세그레
매장이다(구성, 보조정리 \ref{constructions-lemma-segre-embedding}).
$j$가 몰입이므로 $(j, j_{d - d_0})$도 몰입이다(보조정리
\ref{morphisms-lemma-immersion-permanence}를 적용하라). 몰입들의 합성이므로
몰입이다(스킴, 보조정리 \ref{schemes-lemma-composition-immersion}).
\end{proof}

\section{준사영 사상}
\label{morphisms-section-quasi-projective}

\noindent
앞 절의 논의는 다음 정의들을 암시한다. 준사영의 정의는 \cite{EGA}에서
가져온다. 문자 ``H''가 붙은 판본은 \cite{H}의 정의이다.

\begin{definition}
\label{morphisms-definition-quasi-projective}
\begin{reference}
\cite[II, Definition 5.3.1]{EGA} 및 \cite[page 103]{H}
\end{reference}
$f : X \to S$를 스킴의 사상이라 하자.
\begin{enumerate}
\item $f$를 {\it 준사영}이라고 하는 것은 $f$가 유한형이고
$f$-상대적으로 충분한 가역 $\mathcal{O}_X$-가군이 존재할 때이다.
\item $f$를 {\it H-준사영}이라고 하는 것은 준콤팩트 몰입
$X \to \mathbf{P}^n_S$가 $S$ 위에서 어떤 $n$에 대해 존재할 때이다.\footnote{이는 Hartshorne의 정의와 정확히 같지는 않다. 즉,
Hartshorne의 정의(1997년 제8차 교정 인쇄)는 $f$가 열린 몰입에 이어
H-사영 사상(정의 \ref{morphisms-definition-projective}을 보라)을 합성한 것이어야
한다는 것인데, 이는 $f$가 준콤팩트임을 함의하지 않는다. 역방향 함의에
대해서는 보조정리
\ref{morphisms-lemma-H-quasi-projective-open-H-projective}를 보라.}
\item $f$를 {\it 국소 준사영}이라고 하는 것은 열린 덮개
$S = \bigcup V_j$가 존재하여 각각의 $f^{-1}(V_j) \to V_j$가
준사영일 때이다.
\end{enumerate}
\end{definition}

\noindent
이 정의가 암시하듯 준사영이라는 성질은 $S$에서 국소적이지 않다.
나중에는 준사영 스킴의 범주에 대해 더 말할 수 있게 된다. 사상론 심화,
절 \ref{more-morphisms-section-quasi-projective}를 보라.

\begin{lemma}
\label{morphisms-lemma-base-change-quasi-projective}
준사영 사상의 기저변환은 준사영이다.
\end{lemma}

\begin{proof}
보조정리 \ref{morphisms-lemma-base-change-finite-type}와
\ref{morphisms-lemma-ample-base-change}에서 따른다.
\end{proof}

\begin{lemma}
\label{morphisms-lemma-composition-quasi-projective}
$f : X \to Y$와 $g : Y \to S$를 스킴의 사상이라 하자.
$S$가 준콤팩트이고 $f$와 $g$가 준사영이면 $g \circ f$는 준사영이다.
\end{lemma}

\begin{proof}
보조정리 \ref{morphisms-lemma-composition-finite-type}와
\ref{morphisms-lemma-ample-composition}에서 따른다.
\end{proof}

\begin{lemma}
\label{morphisms-lemma-quasi-projective-properties}
$f : X \to S$를 스킴의 사상이라 하자. $f$가 준사영이거나,
H-준사영이거나, 국소 준사영이면 $f$는 분리 유한형이다.
\end{lemma}

\begin{proof}
생략한다.
\end{proof}

\begin{lemma}
\label{morphisms-lemma-H-quasi-projective-quasi-projective}
H-준사영 사상은 준사영이다.
\end{lemma}

\begin{proof}
생략한다.
\end{proof}

\begin{lemma}
\label{morphisms-lemma-characterize-locally-quasi-projective}
$f : X \to S$를 스킴의 사상이라 하자.
다음은 동치이다.
\begin{enumerate}
\item 사상 $f$는 국소 준사영이다.
\item 열린 덮개 $S = \bigcup V_j$가 존재하여 각각의
$f^{-1}(V_j) \to V_j$는 H-준사영이다.
\end{enumerate}
\end{lemma}

\begin{proof}
보조정리 \ref{morphisms-lemma-H-quasi-projective-quasi-projective}에 의해 (2)가
(1)을 함의함을 안다. (1)을 가정하자. 문제는 $S$에서 국소적이므로
$S$가 아핀이고, $X$가 $S$ 위에서 유한형이며, $\mathcal{L}$이 $X/S$
위의 상대적으로 충분한 가역층이라고 가정해도 된다. 보조정리
\ref{morphisms-lemma-finite-type-over-affine-ample-very-ample}에 의해
$\mathcal{L}$이 $X$ 위에서 충분하다고 가정해도 된다. 보조정리
\ref{morphisms-lemma-quasi-projective-finite-type-over-S}에 의해 $X$에서 $S$ 위의
사영공간으로 가는 몰입이 존재함을 안다. 즉, 원하는 대로 $X$는 $S$
위에서 H-준사영이다.
\end{proof}

\begin{lemma}
\label{morphisms-lemma-quasi-affine-finite-type-quasi-projective}
\begin{reference}
\cite[II, Proposition 5.3.4 (i)]{EGA}
\end{reference}
유한형 준아핀 사상은 준사영이다.
\end{lemma}

\begin{proof}
보조정리 \ref{morphisms-lemma-quasi-affine-O-ample}에서 따른다.
\end{proof}

\begin{lemma}
\label{morphisms-lemma-quasi-projective-permanence}
$g : Y \to S$와 $f : X \to Y$를 스킴의 사상이라 하자.
$g \circ f$가 준사영이고 $f$가 준콤팩트이면\footnote{$g$가 준분리이면
이는 스킴, 보조정리 \ref{schemes-lemma-quasi-compact-permanence}에서
따른다.} $f$는 준사영이다.
\end{lemma}

\begin{proof}
보조정리 \ref{morphisms-lemma-permanence-finite-type}에 의해 $f$가 유한형임을
관찰하라. 따라서 보조정리 \ref{morphisms-lemma-ample-permanence}와 정의에서
결론이 따른다.
\end{proof}

\section{고유 사상}
\label{morphisms-section-proper}

\noindent
고유 사상의 개념은 대수기하학에서 중요한 역할을 한다. 고유 사상의 중요한
예로 사상 $\mathbf{P}^n_S \to S$가 있는데, 이는 사영 $n$-공간의 구조
사상이며 실제로 정의에 이르게 한 동기적 예이다.

\begin{definition}
\label{morphisms-definition-proper}
$f : X \to S$를 스킴의 사상이라 하자.
$f$를 {\it 고유}라고 하는 것은 $f$가 분리이고 유한형이며 보편적으로
닫혀 있을 때이다.
\end{definition}

\noindent
원점이 이중화된 아핀 직선에서 아핀 직선으로 가는 사상은 유한형이고
보편적으로 닫혀 있으므로, 위 정의에서 분리 조건은 필요하다. 이 절의
나머지에서는 고유 사상과 보편적으로 닫힌 사상의 몇 가지 기본 성질을
증명한다.

\begin{lemma}
\label{morphisms-lemma-universally-closed-local-on-the-base}
$f : X \to S$를 스킴의 사상이라 하자.
다음은 동치이다.
\begin{enumerate}
\item 사상 $f$는 보편적으로 닫혀 있다.
\item 열린 덮개 $S = \bigcup V_j$가 존재하여
$f^{-1}(V_j) \to V_j$가 모든 지표 $j$에 대해 보편적으로 닫혀 있다.
\end{enumerate}
\end{lemma}

\begin{proof}
정의에서 명백하다.
\end{proof}

\begin{lemma}
\label{morphisms-lemma-proper-local-on-the-base}
$f : X \to S$를 스킴의 사상이라 하자.
다음은 동치이다.
\begin{enumerate}
\item 사상 $f$는 고유이다.
\item 열린 덮개 $S = \bigcup V_j$가 존재하여
$f^{-1}(V_j) \to V_j$가 모든 지표 $j$에 대해 고유이다.
\end{enumerate}
\end{lemma}

\begin{proof}
생략한다.
\end{proof}

\begin{lemma}
\label{morphisms-lemma-composition-proper}
고유 사상들의 합성은 고유이다. 보편적으로 닫힌 사상들에 대해서도 같은
명제가 성립한다.
\end{lemma}

\begin{proof}
닫힌 사상들의 합성은 닫힌 사상이다. $X \to Y \to Z$가 보편적으로 닫힌
사상들이고 $Z' \to Z$가 임의의 사상이면,
$Z' \times_Z X = (Z' \times_Z Y) \times_Y X  \to Z' \times_Z Y$가
닫혀 있고 $Z' \times_Z Y \to Z'$도 닫혀 있음을 안다. 따라서
보편적으로 닫힌 사상에 대한 결론이 따른다. ``분리''와 ``유한형''이
합성에 의해 보존됨은 이미 보았다(스킴, 보조정리
\ref{schemes-lemma-separated-permanence} 및 보조정리
\ref{morphisms-lemma-composition-finite-type}). 따라서 고유 사상에 대한 결론도
따른다.
\end{proof}

\begin{lemma}
\label{morphisms-lemma-base-change-proper}
고유 사상의 기저변환은 고유이다. 보편적으로 닫힌 사상에 대해서도 같은
명제가 성립한다.
\end{lemma}

\begin{proof}
보편적으로 닫힌 사상에 대해서는 정의에 의해 성립한다. 분리 사상에
대해서도 성립한다(스킴, 보조정리
\ref{schemes-lemma-separated-permanence}). 유한형 사상에 대해서도
성립한다(보조정리 \ref{morphisms-lemma-base-change-finite-type}). 따라서 고유
사상에 대해서도 성립한다.
\end{proof}

\begin{lemma}
\label{morphisms-lemma-closed-immersion-proper}
닫힌 몰입은 고유이고, 따라서 더욱이 보편적으로 닫혀 있다.
\end{lemma}

\begin{proof}
닫힌 몰입의 기저변환은 닫힌 몰입이다(스킴, 보조정리
\ref{schemes-lemma-base-change-immersion}). 따라서 이는 보편적으로
닫혀 있다. 닫힌 몰입은 분리이다(스킴, 보조정리
\ref{schemes-lemma-immersions-monomorphisms}). 닫힌 몰입은 유한형이다
(보조정리 \ref{morphisms-lemma-immersion-locally-finite-type}). 따라서 닫힌 몰입은
고유이다.
\end{proof}

\begin{lemma}
\label{morphisms-lemma-image-proper-scheme-closed}
다음과 같은 스킴의 가환 그림이 주어졌다고 하자.
$$
\xymatrix{
X \ar[rr] \ar[rd] & &
Y \ar[ld] \\
& S &
}
$$
$Y$가 $S$ 위에서 분리라고 하자.
\begin{enumerate}
\item $X \to S$가 보편적으로 닫혀 있으면 사상 $X \to Y$도
보편적으로 닫혀 있다.
\item $X$가 $S$ 위에서 고유이면 사상 $X \to Y$도 고유이다.
\end{enumerate}
특히 두 경우 모두 $X$의 $Y$에서의 상은 닫혀 있다.
\end{lemma}

\begin{proof}
$X \to S$가 보편적으로 닫혀 있다고(각각 고유이라고) 가정하자. 사상을
$X \to X \times_S Y \to Y$로 분해한다. 첫 번째 사상은 닫힌 몰입이다.
스킴, 보조정리 \ref{schemes-lemma-semi-diagonal}을 보라. 따라서 첫 번째
사상은 고유이다(보조정리 \ref{morphisms-lemma-closed-immersion-proper}). 사영
$X \times_S Y \to Y$는 보편적으로 닫힌(각각 고유인) 사상의 기저변환이며,
따라서 보편적으로 닫혀 있다(각각 고유이다). 보조정리
\ref{morphisms-lemma-base-change-proper}을 보라. 그러므로 $X \to Y$는 보편적으로
닫힌(각각 고유인) 사상들의 합성이므로 보편적으로 닫혀 있다(각각 고유이다)
(보조정리 \ref{morphisms-lemma-composition-proper}).
\end{proof}

\noindent
다음 보조정리의 증명은 Bjorn Poonen에 의한다.
\href{https://mathoverflow.net/questions/23337/is-a-universally-closed-morphism-of-schemes-quasi-compact/23528#23528}{이 위치}를
보라.

\begin{lemma}
\label{morphisms-lemma-universally-closed-quasi-compact}
\begin{reference}
Bjorn Poonen에 의한다.
\end{reference}
보편적으로 닫힌 스킴 사상은 준콤팩트이다.
\end{lemma}

\begin{proof}
$f : X \to S$를 사상이라 하자. $f$가 준콤팩트가 아니라고 가정한다.
목표는 $f$가 보편적으로 닫혀 있지 않음을 보이는 것이다. 스킴, 보조정리
\ref{schemes-lemma-quasi-compact-affine}에 의해 아핀 열린집합
$V \subset S$가 존재하여 $f^{-1}(V)$는 준콤팩트가 아니다. 목표를
달성하려면 $f^{-1}(V) \to V$가 보편적으로 닫혀 있지 않음을 보이면
충분하므로, $S = \Spec(A)$라고 가정해도 된다. 여기서 $A$는 어떤
환이다.

\medskip\noindent
$X = \bigcup_{i \in I} X_i$라고 쓰되 $X_i$들은 $X$의 아핀 열린
부분스킴이라고 하자. $T = \Spec(A[y_i ; i \in I])$라 하자.
$T_i = D(y_i) \subset T$라 하자. $Z$를 닫힌집합
$(X \times_S T) - \bigcup_{i \in I} (X_i \times_S T_i)$라 하자.
$f_T(Z)$가 닫혀 있지 않음을 증명하면 충분하다. 여기서 이는 $Z$의
$f_T : X \times_S T \to T$ 아래의 상이다.

\medskip\noindent
점 $s \in S$가 존재하여, $U$가 $s$의 $S$에서의 근방일 때 $X_U$가
준콤팩트인 경우는 하나도 없다. 그렇지 않으면 $S$를 이러한 $U$들 중
유한 개로 덮을 수 있고, 스킴, 보조정리
\ref{schemes-lemma-quasi-compact-affine}에 의해 $f$가 준콤팩트가 된다.
그러한 $s \in S$를 고정하자.

\medskip\noindent
먼저 $f_T(Z_s) \ne T_s$임을 확인한다. $t \in T$를 $s$ 위에 놓이고
$\kappa(t) = \kappa(s)$이며 $y_i = 1$이 $\kappa(t)$에서 모든 $i$에
대해 성립하는 점으로 택하자. 그러면 $t \in T_i$가 모든 $i$에 대해
성립하고, $Z_s \to T_s$의 $t$ 위의 올은
$(X - \bigcup_{i \in I} X_i)_s$와 동형인데 이는 공집합이다. 따라서
$t \in T_s - f_T(Z_s)$이다.

\medskip\noindent
$f_T(Z)$가 $T$에서 닫혀 있다고 가정하자. 그러면 원소
$g \in A[y_i; i \in I]$가 존재하여 $f_T(Z) \subset V(g)$이지만
$t \not \in V(g)$이다. 따라서 $g$의 $\kappa(t)$에서의 상은 영이
아니다. 특히 $g$의 어떤 계수는 $\kappa(s)$에서 영이 아닌 상을 가진다.
그러므로 $U$가 $s$의 어떤 근방일 때 이 계수는 그 근방에서 가역이다.
$J$를, $j \in I$이고 $y_j$가 $g$에 나타나는 모든 지표의 유한집합이라
하자. $X_U$가 준콤팩트가 아니므로, 어떤 점
$x \in X - \bigcup_{j \in J} X_j$를 택하여 어떤 $u \in U$ 위에 놓이게
할 수 있다. $g$의 어떤 계수가 $U$에서 가역이므로, 점 $t' \in T$를
찾아서 $u$ 위에 놓이고 $t' \not \in V(g)$이며
$t' \in V(y_i)$가 모든 $i \notin J$에 대해 성립하게 할 수 있다. 이는
$V(y_i; i \in I, i \not\in J) = \Spec(A[y_j; j\in J])$이고 이 스킴의
$u$ 위에 놓이는 점들의 집합이
$\Spec(\kappa(u)[y_j; j \in J])$와 전단사이기 때문이다. 다시 말해
$t' \notin T_i$가 각 $i \notin J$에 대해 성립한다. 스킴, 보조정리
\ref{schemes-lemma-points-fibre-product}에 의해 점 $z$를
$X \times_S T$에서 찾아서 $x \in X$와 $t' \in T$로 사상되게 할 수
있다. $x \not \in X_j$가 $j \in J$에 대해 성립하고
$t' \not \in T_i$가 $i \in I \setminus J$에 대해 성립하므로
$z \in Z$임을 안다. 한편 $f_T(z) = t' \not \in V(g)$인데, 이는
$f_T(Z) \subset V(g)$에 모순이다. 따라서 ``$f_T(Z)$가 닫혀 있다''라는
가정은 잘못이며, 원하는
대로 $f_T$가 닫힌 사상이 아니라고 결론짓는다.
\end{proof}

\noindent
다음 보조정리는 고유 스킴의 상이 (기저 위에서 유한형인 분리 스킴 안에서)
고유임을 말한다.

\begin{lemma}
\label{morphisms-lemma-image-proper-is-proper}
$S$를 스킴이라 하자. $f : X \to Y$를 $S$ 위의 스킴 사상이라 하자.
$X$가 $S$ 위에서 보편적으로 닫혀 있고 $f$가 전사이면 $Y$는 $S$
위에서 보편적으로 닫혀 있다. 특히 $Y$가 또한 $S$ 위에서 분리이고
국소 유한형이면 $Y$는 $S$ 위에서 고유이다.
\end{lemma}

\begin{proof}
$X$가 보편적으로 닫혀 있고 $f$가 전사라고 가정하자. 구조 사상들을
$p : X \to S$, $q : Y \to S$로 나타내자. $S' \to S$를 스킴의
사상이라 하자. 기저변환 $f' : X_{S'} \to Y_{S'}$은 전사이고
(보조정리 \ref{morphisms-lemma-base-change-surjective}), 기저변환
$p' : X_{S'} \to S'$은 닫힌 사상이다. $T \subset Y_{S'}$가
닫혀 있으면 $(f')^{-1}(T) \subset X_{S'}$도 닫혀 있으므로
$p'((f')^{-1}(T)) = q'(T)$는 닫혀 있다. 따라서 $q'$는 닫힌 사상이다.
이는 첫 번째 명제를 증명한다. 그러므로 보조정리
\ref{morphisms-lemma-universally-closed-quasi-compact}에 의해 $Y \to S$는
준콤팩트이고, 따라서 $Y \to S$가 추가로 국소 유한형이고 분리이면
정의에 의해 $Y \to S$는 고유이다.
\end{proof}

\begin{lemma}
\label{morphisms-lemma-scheme-theoretic-image-is-proper}
다음과 같은 스킴의 가환 그림이 주어졌다고 하자.
$$
\xymatrix{
X \ar[rr]_h \ar[rd]_f & & Y \ar[ld]^g \\
& S
}
$$
다음을 가정하자.
\begin{enumerate}
\item $X \to S$는 보편적으로 닫힌(예를 들어 고유인) 사상이다.
\item $Y \to S$는 분리이고 국소 유한형이다.
\end{enumerate}
그러면 $Z \subset Y$로 나타낸 $h$의 스킴론적 상은 $S$ 위에서 고유이고
$X \to Z$는 전사이다.
\end{lemma}

\begin{proof}
$h$의 스킴론적 상은 절 \ref{morphisms-section-scheme-theoretic-image}에서
구성하였다. $f$는 준콤팩트이므로(보조정리
\ref{morphisms-lemma-universally-closed-quasi-compact}) $h$도 준콤팩트임을 안다
(스킴, 보조정리 \ref{schemes-lemma-quasi-compact-permanence}). 따라서
$h(X) \subset Z$는 조밀하다(보조정리
\ref{morphisms-lemma-quasi-compact-scheme-theoretic-image}). 한편 $h(X)$는
$Y$에서 닫혀 있으므로(보조정리
\ref{morphisms-lemma-image-proper-scheme-closed}) $X \to Z$는 전사이다. 따라서
$Z \to S$는 고유이다(보조정리 \ref{morphisms-lemma-image-proper-is-proper}).
\end{proof}

\noindent
전사인 보편적으로 닫힌 사상 아래에서 분리 스킴의 표적은 분리이다.

\begin{lemma}
\label{morphisms-lemma-image-universally-closed-separated}
$S$를 스킴이라 하자. $f : X \to Y$를 $S$ 위의 전사인 보편적으로
닫힌 스킴 사상이라 하자.
\begin{enumerate}
\item $X$가 준분리이면 $Y$도 준분리이다.
\item $X$가 분리이면 $Y$도 분리이다.
\item $X$가 $S$ 위에서 준분리이면 $Y$도 $S$ 위에서 준분리이다.
\item $X$가 $S$ 위에서 분리이면 $Y$도 $S$ 위에서 분리이다.
\end{enumerate}
\end{lemma}

\begin{proof}
(1)과 (2)는 $S = \Spec(\mathbf{Z})$에 대해 (3)과 (4)를 적용한 결과이다
(스킴, 정의 \ref{schemes-definition-separated}을 보라). 다음 가환 그림을
생각하자.
$$
\xymatrix{
X \ar[d] \ar[rr]_{\Delta_{X/S}} & & X \times_S X \ar[d] \\
Y \ar[rr]^{\Delta_{Y/S}} & & Y \times_S Y
}
$$
왼쪽 수직 화살표는 전사이다(즉 보편적으로 전사이다). 오른쪽 수직 화살표는
보편적으로 닫힌 사상들
$X \times_S X \to X \times_S Y \to Y \times_S Y$의 합성이므로
보편적으로 닫혀 있다. 따라서 이 사상은 준콤팩트이기도 하다. 보조정리
\ref{morphisms-lemma-universally-closed-quasi-compact}를 보라.

\medskip\noindent
$X$가 $S$ 위에서 준분리라고, 즉 $\Delta_{X/S}$가 준콤팩트라고 가정하자.
$V \subset Y \times_S Y$가 준콤팩트 열린집합이면
$V \times_{Y \times_S Y} X \to \Delta_{Y/S}^{-1}(V)$는 전사이고,
위의 설명에 의해 $V \times_{Y \times_S Y} X$는 준콤팩트이다. 따라서
$\Delta_{Y/S}$가 준콤팩트임을 결론짓는다. 즉 $Y$는 $S$ 위에서
준분리이다.

\medskip\noindent
$X$가 $S$ 위에서 분리라고, 즉 $\Delta_{X/S}$가 닫힌 몰입이라고
가정하자. 그러면 $X \to Y \times_S Y$는 닫힌 사상들의 합성이므로
닫혀 있다. $X \to Y$가 전사이므로 $\Delta_{Y/S}(Y)$가
$Y \times_S Y$에서 닫혀 있음을 얻는다. 따라서 스킴, 정의
\ref{schemes-definition-separated} 다음의 논의에 의해 $Y$는 $S$ 위에서
분리이다.
\end{proof}

\section{값매김 판정법}
\label{morphisms-section-valuative-criteria}

\noindent
보편 닫힘과 분리성의 값매김 판정법은 이미 스킴, 절
\ref{schemes-section-valuative-criterion-universal-closedness} 및
\ref{schemes-section-valuative-separatedness}에서 논의하였다. 이 절에서는
몇 가지 귀결과 변형을 논의한다. 극한, 절
\ref{limits-section-Noetherian-valuative-criterion}에서는 국소 뇌터
스킴과 유한형 사상을 다룰 때 이산 값매김환만 생각하면 충분함을 보일
것이다.

\begin{lemma}[고유성의 값매김 판정법]
\label{morphisms-lemma-characterize-proper}
\begin{reference}
\cite[II Theorem 7.3.8]{EGA}
\end{reference}
$S$를 스킴이라 하자. $f : X \to Y$를 $S$ 위의 스킴 사상이라 하자.
$f$가 유한형이고 준분리라고 가정하자. 다음은 동치이다.
\begin{enumerate}
\item $f$는 고유이다.
\item $f$는 값매김 판정법을 만족한다
(스킴, 정의 \ref{schemes-definition-valuative-criterion}).
\item 임의의 가환 실선 그림
$$
\xymatrix{
\Spec(K) \ar[r] \ar[d] & X \ar[d] \\
\Spec(A) \ar[r] \ar@{-->}[ru] & Y
}
$$
이 주어졌을 때, 여기서 $A$는 분수체가 $K$인 값매김환이며, 그림을
가환으로 만드는 점선 화살표가 유일하게 존재한다.
\end{enumerate}
\end{lemma}

\begin{proof}
(3)은 (2)를 다시 서술한 것이다. 따라서 이 보조정리는 스킴, 명제
\ref{schemes-proposition-characterize-universally-closed}, 보조정리
\ref{schemes-lemma-separated-implies-valuative}, 보조정리
\ref{schemes-lemma-valuative-criterion-separatedness} 및 정의들의 형식적
귀결이다.
\end{proof}

\noindent
값매김 판정법을 시험할 때 보통 가능한 모든 그림을 생각할 필요는 없다.
다음 보조정리와 같은 값매김 판정법을 ``정제된 값매김 판정법''이라고
부르겠다.

\begin{lemma}
\label{morphisms-lemma-refined-valuative-criterion-universally-closed}
$f : X \to S$와 $h : U \to X$를 스킴의 사상이라 하자. $f$와 $h$가
준콤팩트이고 $h(U)$가 $X$에서 조밀하다고 가정하자. 임의의 가환 실선
그림
$$
\xymatrix{
\Spec(K) \ar[r] \ar[d] & U \ar[r]^h & X \ar[d]^f \\
\Spec(A) \ar[rr] \ar@{-->}[rru] & & S
}
$$
이 주어졌을 때, 여기서 $A$는 분수체가 $K$인 값매김환이며, 그림을
가환으로 만드는 점선 화살표가 유일하게 존재한다고 하자. 그러면 $f$는
보편적으로 닫혀 있다. 또한 $f$가 준분리이면 $f$는 분리이다.
\end{lemma}

\begin{proof}
$f$가 보편적으로 닫혀 있음을 증명하기 위해 $f$에 대한 값매김 판정법의
존재 부분을 확인하겠다. 이는 스킴, 명제
\ref{schemes-proposition-characterize-universally-closed}에 의해 충분하다.
이를 위해 다음 가환 그림을 생각하자.
$$
\xymatrix{
\Spec(K) \ar[r] \ar[d] & X \ar[d] \\
\Spec(A) \ar[r] & S
}
$$
여기서 $A$는 값매김환이고 $K$는 $A$의 분수체이다. 값매김환과 체는
기약이므로, 스킴, 보조정리 \ref{schemes-lemma-map-into-reduction}에 의해
$U$, $X$, $S$를 각각의 축소로 대체해도 됨에 유의하라. 이 경우 $h(U)$가
조밀하다는 가정은 $h : U \to X$의 스킴론적 상이 $X$라는 뜻이다.
보조정리 \ref{morphisms-lemma-scheme-theoretic-image-reduced}를 보라. 또한 $S$를
사상 $\Spec(A) \to S$가 통과하는 아핀 열린집합으로 대체해도 된다.
따라서 $S = \Spec(R)$이라고 가정해도 된다.

\medskip\noindent
$\Spec(B) \subset X$를 사상 $\Spec(K) \to X$가 통과하는 아핀
열린집합이라 하자. $P$를 $B$ 위의 다항식 대수라 하고 $B$-대수 전사
$P \to K$를 택하자. 그러면 $\Spec(P) \to X$는 평탄하다. 따라서
보조정리 \ref{morphisms-lemma-flat-base-change-scheme-theoretic-image}에 의해 사상
$U \times_X \Spec(P) \to \Spec(P)$의 스킴론적 상은 $\Spec(P)$이다.
보조정리 \ref{morphisms-lemma-reach-points-scheme-theoretic-image}에 의해 다음 가환
그림을 찾을 수 있다.
$$
\xymatrix{
\Spec(K') \ar[r] \ar[d] & U \times_X \Spec(P) \ar[d] \\
\Spec(A') \ar[r] & \Spec(P)
}
$$
여기서 $A'$은 값매김환이고 $K'$은 $A'$의 분수체이며,
$\Spec(A')$의 닫힌점은 $\Spec(K) \subset \Spec(P)$로 사상된다. 다시
말해 $B$-대수 사상
$\varphi : K \to A'/\mathfrak m_{A'}$가 존재한다. 값매김환
$A'' \subset A'/\mathfrak m_{A'}$을 택하여 $\varphi(A)$를 지배하고
그 분수체가 $K'' = A'/\mathfrak m_{A'}$이 되게 하자(대수학, 보조정리
\ref{algebra-lemma-dominate}). 다음과 같이 놓는다.
$$
C = \{\lambda \in A' \mid \lambda \bmod \mathfrak m_{A'} \in A''\}.
$$
이는 대수학, 보조정리 \ref{algebra-lemma-stack-valuation-rings}에 의해
값매김환이다. $C$는 $R$-대수이고 분수체는 $K'$이므로, 보조정리의 명제에
나오는 다음 가환 그림을 얻는다.
$$
\xymatrix{
\Spec(K') \ar[r] \ar[d] & U \ar[r] & X \ar[d] \\
\Spec(C) \ar[rr] \ar@{-->}[rru] & & S
}
$$
따라서 표시된 대로 그림에 들어가는 점선 화살표가 존재한다. 보조정리의
유일성 가정에 의해 합성 $\Spec(A') \to \Spec(C) \to X$는 주어진 사상
$\Spec(A') \to \Spec(P) \to \Spec(B) \subset X$와 일치한다. 그러므로
이 사상을 $C/\mathfrak m_{A'} = A''$의 스펙트럼에 제한하면 주어진 사상
$\Spec(K'') = \Spec(A'/\mathfrak m_{A'}) \to \Spec(K) \to X$를 유도한다.
$x \in X$를 $\Spec(A'') \to X$ 아래에서 닫힌점의 상이라 하자. 유도된
환 사상 $\mathcal{O}_{X, x} \to A''$의 상은 $K \subset K''$에 포함되는
국소 부분환이다. $A$는 $K$ 안에서 지배 관계에 관하여 극대이고
$A \subset A''$이므로 $A = K \cap A''$이다. 따라서
$\mathcal{O}_{X, x} \to A''$는 $A \subset A''$를 통해 분해된다. 이로써
원하는 화살표 $\Spec(A) \to X$를 얻는다.

\medskip\noindent
마지막으로 $f$가 준분리라고 가정하자. 그러면
$\Delta : X \to X \times_S X$는 준콤팩트이다. 실선 그림
$$
\xymatrix{
\Spec(K) \ar[r] \ar[d] & U \ar[r]^h & X \ar[d]^\Delta \\
\Spec(A) \ar[rr] \ar@{-->}[rru] & & X \times_S X
}
$$
이 주어졌을 때, 여기서 $A$는 분수체가 $K$인 값매김환이며, 그림을
가환으로 만드는 점선 화살표가 유일하게 존재한다. 실제로 아래쪽 수평
화살표는 보조정리의 그림에서 점선 화살표 역할을 할 수 있는 사상들의 쌍
$\Spec(A) \to X$와 같은 것이다. 따라서 요구되는 유일성은 아래쪽 수평
화살표가 $\Delta$를 통해 분해됨을 보인다. 그러므로 방금 증명한 결과를
$\Delta : X \to X \times_S X$와 $h : U \to X$에 적용하여 $\Delta$가
보편적으로 닫혀 있음을 결론짓는다. 이것은 분명히 $f$가 분리라는
뜻이다.
\end{proof}










\begin{remark}
\label{morphisms-remark-check-val-on-open}
보조정리 \ref{morphisms-lemma-refined-valuative-criterion-universally-closed}에서 점선
화살표의 유일성에 대한 가정은 필요하다(세부사항은 생략한다). 물론 $f$가
분리이면 유일성은 보장된다(스킴, 보조정리
\ref{schemes-lemma-separated-implies-valuative}).
\end{remark}

\begin{lemma}
\label{morphisms-lemma-morphism-defined-local-ring}
$S$를 스킴이라 하자. $X$, $Y$를 $S$ 위의 스킴이라 하자. $s \in S$이고
$x \in X$, $y \in Y$가 $s$ 위의 점들이라 하자.
\begin{enumerate}
\item $f, g : X \to Y$를 $S$ 위의 사상이라 하고
$f(x) = g(x) = y$이며
$f^\sharp_x = g^\sharp_x : \mathcal{O}_{Y, y} \to \mathcal{O}_{X, x}$라고
하자. 다음 각 경우에는 열린 근방 $U \subset X$가 존재하여
$f|_U = g|_U$이다.
\begin{enumerate}
\item $Y$가 $S$ 위에서 국소 유한형이다.
\item $X$가 정역 스킴이다.
\item $X$가 국소 뇌터이다.
\item $X$가 기약 성분을 유한 개 갖는 축소 스킴이다.
\end{enumerate}
\item $\varphi : \mathcal{O}_{Y, y} \to \mathcal{O}_{X, x}$를 국소
$\mathcal{O}_{S, s}$-대수 사상이라 하자. 다음 각 경우에는 열린집합
$U \subset X$가 $x$를 포함하고 사상 $f : U \to Y$가 존재하여 $x$를 $y$로
보내고 $f^\sharp_x = \varphi$가 된다.
\begin{enumerate}
\item $Y$가 $S$ 위에서 국소 유한 표시이다.
\item $Y$가 국소 유한형이고 $X$가 정역 스킴이다.
\item $Y$가 국소 유한형이고 $X$가 국소 뇌터이다.
\item $Y$가 국소 유한형이고 $X$가 기약 성분을 유한 개 갖는 축소
스킴이다.
\end{enumerate}
\end{enumerate}
\end{lemma}

\begin{proof}
(1)의 증명. $X$, $Y$, $S$를 각각 $x$, $y$, $s$의 적당한 아핀 열린
근방으로 대체하면 다음 대수 문제로 환원된다. 환 $R$과 두 $R$-대수 사상
$\varphi, \psi : B \to A$가 주어졌고 다음을 만족한다고 하자.
\begin{enumerate}
\item $R \to B$가 유한형이거나, $A$가 정역이거나, $A$가 뇌터이거나,
$A$가 축소이고 극소 소아이디얼을 유한 개 가진다.
\item 두 사상 $B \to A_\mathfrak p$가 어떤 소아이디얼
$\mathfrak p \subset A$에 대해 같다.
\end{enumerate}
$\varphi, \psi$가 같은 사상 $B \to A_g$를 정의함을 보이면 된다. 여기서
$g \in A$, $g \not \in \mathfrak p$는 적당히 택한다. $R \to B$가 유한형이면
$t_1, \ldots, t_m \in B$를 $B$의 $R$-대수 생성원으로 택하자.
각 $j$에 대해 $g_j \in A$, $g_j \not \in \mathfrak p$를 찾아서
$\varphi(t_j)$와 $\psi(t_j)$의 $A_{g_j}$에서의 상이 같게 할 수 있다.
그런 다음 $g = \prod g_j$로 놓는다. 다른 경우들($A$가 정역이거나,
뇌터이거나, 축소이고 극소 소아이디얼을 유한 개 갖는 경우)에는
$g \in A$, $g \not \in \mathfrak p$를 찾아서
$A_g \subset A_\mathfrak p$가 되게 할 수 있다. 대수학, 보조정리
\ref{algebra-lemma-subring-of-local-ring}을 보라. 따라서 원하는 대로
사상들 $B \to A_g$가 같다.

\medskip\noindent
(2)의 증명. 이를 위해 $X$, $Y$, $S$를 적당한 아핀 열린집합으로
대체해도 된다. $X = \Spec(A)$, $Y = \Spec(B)$,
$S = \Spec(R)$이라 하자. $\mathfrak p \subset A$를 $x$에 대응하는
소아이디얼이라 하고, $\mathfrak q \subset B$를 $y$에 대응하는
소아이디얼이라 하자. 그러면 $\varphi$는 국소 $R$-대수 사상
$\varphi : B_\mathfrak q \to A_\mathfrak p$이다. $R \to B$가 유한
표시인 환 사상이면, 어떤 $g \in A \setminus \mathfrak p$와 어떤
$R$-대수 사상 $B \to A_g$가 존재하여
$$
\xymatrix{
B_\mathfrak q \ar[r]_\varphi & A_\mathfrak p \\
B \ar[u] \ar[r] & A_g \ar[u]
}
$$
가환이다. 대수학, 보조정리
\ref{algebra-lemma-characterize-finite-presentation} 및
\ref{algebra-lemma-localization-colimit}를 보라. 유도된 사상
$\Spec(A_g) \to \Spec(B)$가 원하는 사상이다. $B$가 $R$ 위에서
유한형이면 $t_1, \ldots, t_m \in B$를 $B$의 $R$-대수
생성원으로 택하자. 그러면 $g_j \in A$, $g_j \not \in \mathfrak p$를
택하여 $\varphi(t_j) \in \Im(A_{g_j} \to A_\mathfrak p)$가 되게 할 수
있다. 따라서 $A$를 $A[1/\prod g_j]$로 대체한 뒤에는 $B$가
$A \to A_\mathfrak p$의 상 안으로 사상된다고 가정해도 된다. 어떤
$g \in A$, $g \not \in \mathfrak p$를 찾아서
$A_g \to A_\mathfrak p$가 단사이게 할 수 있다면, 원하는 $R$-대수 사상
$B \to A_g$를 얻는다. 따라서 대수학, 보조정리
\ref{algebra-lemma-subring-of-local-ring}을 한 번 더 적용하면 증명이
끝난다.
\end{proof}

\begin{lemma}
\label{morphisms-lemma-extend-across}
$S$를 스킴이라 하자. $X$, $Y$를 $S$ 위의 스킴이라 하자. $x \in X$라
하자. $U \subset X$를 열린집합이라 하고 $f : U \to Y$를 $S$ 위의
사상이라 하자. 다음을 가정하자.
\begin{enumerate}
\item $x$가 $U$의 폐포에 속한다.
\item $X$가 기약 성분을 유한 개 갖는 축소 스킴이거나 $X$가 뇌터이다.
\item $\mathcal{O}_{X, x}$가 값매김환이다.
\item $Y \to S$가 고유이다.
\end{enumerate}
그러면 열린집합 $U \subset U' \subset X$가 $x$를 포함하고, $S$-사상
$f' : U' \to Y$가 존재하여 $f$를 확장한다.
\end{lemma}

\begin{proof}
$X$를 $x$의 $X$에서의 열린 근방으로 대체해도 무방하다(작은 세부사항은
생략한다). 성질, 보조정리
\ref{properties-lemma-ring-affine-open-injective-local-ring}에 의해 $X$가
아핀이고 $\Gamma(X, \mathcal{O}_X) \subset \mathcal{O}_{X, x}$라고
가정해도 된다. 특히 $X$는 유일한 일반점 $\xi$를 갖는 정역 스킴이고,
그 잉여체 $K$는 값매김환 $\mathcal{O}_{X, x}$의 분수체이다. $x$가
$U$의 폐포에 속하므로 $U$는 공집합이 아니며, 따라서 $U$는 $\xi$를
포함한다. 그러므로 고유성의 값매김 판정법(보조정리
\ref{morphisms-lemma-characterize-proper})에 의해 사상
$t : \Spec(\mathcal{O}_{X, x}) \to Y$가 존재하여 다음 그림에 들어간다.
$$
\xymatrix{
\Spec(K) \ar[d]_\xi \ar[r] & \Spec(\mathcal{O}_{X, x}) \ar[d]_t \\
U \ar[r]^f & Y
}
$$
이는 $S$ 위의 스킴 사상들의 가환 그림이다. 보조정리
\ref{morphisms-lemma-morphism-defined-local-ring}을 $y = t(x)$ 및
$\varphi = t^\sharp_x$와 함께 적용하면, 열린집합 $V \subset X$인
$x$의 근방과 사상 $g : V \to Y$를 $S$ 위에서 얻는다. 이 사상은
$x$를 $y$로 보내며 $g^\sharp_x = t^\sharp_x$이다. $Y \to S$가
분리이므로 등화자 $E$, 즉 $f|_{U \cap V}$와 $g|_{U \cap V}$의 등화자는
$U \cap V$의 닫힌 부분스킴이다. 스킴, 보조정리
\ref{schemes-lemma-where-are-they-equal}을 보라. 구성에 의해 $f$와 $g$가
같은 사상 $\Spec(K) \to Y$를 결정하므로, $E$는 정역 스킴
$U \cap V$의 일반점을 포함한다. 따라서 $E = U \cap V$이고, $f$와
$g$를 붙여 원하는 사상 $U' = U \cup V \to Y$를 얻는다.
\end{proof}

\section{사영 사상}
\label{morphisms-section-projective}

\noindent
사영 사상의 정의로 \cite{EGA}의 정의를 사용하겠다. ``H''가 붙은 정의는
\cite{H}의 정의이다. 이렇게 얻는 두 정의는 서로 다르며 둘 다 유용하다.

\begin{definition}
\label{morphisms-definition-projective}
$f : X \to S$를 스킴의 사상이라 하자.
\begin{enumerate}
\item $f$를 {\it 사영}이라고 하는 것은 $X$가 $S$-스킴으로서 사영다발
$\mathbf{P}(\mathcal{E})$의 닫힌 부분스킴과 동형일 때이다. 여기서 이
다발은 어떤 준연접 유한형 $\mathcal{O}_S$-가군 $\mathcal{E}$에 대한 것이다.
\item $f$를 {\it H-사영}이라고 하는 것은 어떤 정수 $n$과 닫힌 몰입
$X \to \mathbf{P}^n_S$가 $S$ 위에서 존재할 때이다.
\item $f$를 {\it 국소 사영}이라고 하는 것은 열린 덮개
$S = \bigcup U_i$가 존재하여 각각의 $f^{-1}(U_i) \to U_i$가 사영일
때이다.
\end{enumerate}
\end{definition}

\noindent
예상대로 사영 사상은 준사영이다. 보조정리
\ref{morphisms-lemma-projective-quasi-projective}를 보라. 역으로 준사영 사상은
흔히 열린 몰입과 사영 사상의 합성이다. 보조정리
\ref{morphisms-lemma-quasi-projective-open-projective}를 보라. 준사영 기저 위의
사영 사상의 성질을 개관하려면 사상론 심화, 절
\ref{more-morphisms-section-projective}를 보라.

\begin{example}
\label{morphisms-example-projective}
$S$를 스킴이라 하자. $\mathcal{A}$를 준연접 등급 $\mathcal{O}_S$-대수라
하고, 이 대수는 $\mathcal{A}_1$이 $\mathcal{A}_0$ 위에서 생성한다고 하자.
또한 $\mathcal{A}_1$이 $\mathcal{O}_S$ 위에서 유한형이라고 가정하자.
$X = \underline{\text{Proj}}_S(\mathcal{A})$로 놓는다. 이 경우
$X \to S$는 사영이다. 실제로 등급 $\mathcal{O}_S$-대수 사상
$$
\text{Sym}_{\mathcal{O}_X}^*(\mathcal{A}_1)
\longrightarrow
\mathcal{A}
$$
에 대응하는 사상은 닫힌 몰입이다. 구성, 보조정리
\ref{constructions-lemma-surjective-generated-degree-1-map-relative-proj}를
보라.
\end{example}
 
\begin{lemma}
\label{morphisms-lemma-H-projective}
H-사영 사상은 H-준사영이다. H-사영 사상은 사영이다.
\end{lemma}

\begin{proof}
첫 번째 명제는 정의에서 즉시 따른다. 두 번째는
$\mathbf{P}^n_S$가 $S$ 위의 사영다발이므로 성립한다. 구성, 보조정리
\ref{constructions-lemma-projective-space-bundle}을 보라.
\end{proof}

\begin{lemma}
\label{morphisms-lemma-characterize-locally-projective}
$f : X \to S$를 스킴의 사상이라 하자. 다음은 동치이다.
\begin{enumerate}
\item 사상 $f$는 국소 사영이다.
\item 열린 덮개 $S = \bigcup U_i$가 존재하여 각각의
$f^{-1}(U_i) \to U_i$가 H-사영이다.
\end{enumerate}
\end{lemma}

\begin{proof}
보조정리 \ref{morphisms-lemma-H-projective}에 의해 (2)는 (1)을 함의한다. (1)을
가정하자. 모든 점 $s \in S$에 대해 아핀 열린집합
$\Spec(R) = U \subset S$를 $s$의 근방으로 찾아서 $X_U$가
$\mathbf{P}(\mathcal{E})$의 닫힌 부분스킴과 동형이게 할 수 있다. 여기서
$\mathcal{O}_U$-가군 $\mathcal{E}$는 유한형이고 준연접이다.
$\mathcal{E} = \widetilde{M}$이라고 쓰되, 어떤 유한형 $R$-가군 $M$을 택한다
(성질, 보조정리
\ref{properties-lemma-finite-type-module}를 보라).
$x_0, \ldots, x_n \in M$을 $M$의 $R$-가군 생성원으로 택하자.
다음 전사인 등급 $R$-대수 사상을 생각하자.
$$
R[X_0, \ldots, X_n] \longrightarrow \text{Sym}_R(M).
$$
구성, 보조정리
\ref{constructions-lemma-surjective-graded-rings-map-proj}에 따르면 대응하는
사상
$$
\mathbf{P}(\mathcal{E}) \to \mathbf{P}^n_R
$$
은 닫힌 몰입이다. 따라서 $f^{-1}(U)$가 $\mathbf{P}^n_U$의 닫힌
부분스킴과 ($U$ 위의 스킴으로서) 동형임을 결론짓는다. 다시 말해 (2)가
성립한다.
\end{proof}

\begin{lemma}
\label{morphisms-lemma-locally-projective-proper}
국소 사영 사상은 고유이다.
\end{lemma}

\begin{proof}
$f : X \to S$를 국소 사영이라 하자. $f$가 고유임을 보이기 위해 기저에서
국소적으로 작업해도 된다. 보조정리
\ref{morphisms-lemma-proper-local-on-the-base}를 보라. 그러므로 위의 보조정리
\ref{morphisms-lemma-characterize-locally-projective}에 의해 닫힌 몰입
$X \to \mathbf{P}^n_S$가 존재한다고 가정해도 된다. 보조정리
\ref{morphisms-lemma-composition-proper} 및 \ref{morphisms-lemma-closed-immersion-proper}에
의해 $\mathbf{P}^n_S \to S$가 고유임을 증명하면 충분하다.
$\mathbf{P}^n_S \to S$는
$\mathbf{P}^n_{\mathbf{Z}} \to \Spec(\mathbf{Z})$의 기저변환이므로,
$\mathbf{P}^n_{\mathbf{Z}} \to \Spec(\mathbf{Z})$가 고유임을 보이면
충분하다. 보조정리 \ref{morphisms-lemma-base-change-proper}을 보라. 구성, 보조정리
\ref{constructions-lemma-proj-separated}에 의해 스킴
$\mathbf{P}^n_{\mathbf{Z}}$는 분리이다. 구성, 보조정리
\ref{constructions-lemma-proj-quasi-compact}에 의해 스킴
$\mathbf{P}^n_{\mathbf{Z}}$는 준콤팩트이다.
$\mathbf{P}^n_{\mathbf{Z}} \to \Spec(\mathbf{Z})$가 국소 유한형임은
명백하다. 실제로 $\mathbf{P}^n_{\mathbf{Z}}$는 아핀 열린집합
$D_{+}(X_i)$들로 덮이며, 각각은 유한형 $\mathbf{Z}$-대수
$$
\mathbf{Z}[X_0/X_i, \ldots, X_n/X_i].
$$
의 스펙트럼이다. 마지막으로
$\mathbf{P}^n_{\mathbf{Z}} \to \Spec(\mathbf{Z})$가 보편적으로 닫혀
있음을 보여야 한다. 이는 구성, 보조정리
\ref{constructions-lemma-proj-valuative-criterion} 및 값매김 판정법
(스킴, 명제 \ref{schemes-proposition-characterize-universally-closed})에서
따른다.
\end{proof}

\begin{lemma}
\label{morphisms-lemma-proper-ample-locally-projective}
$f : X \to S$를 고유인 스킴 사상이라 하자. 어떤 $f$-충분한 가역층이
$X$ 위에
존재하면 $f$는 국소 사영이다.
\end{lemma}

\begin{proof}
$f$-충분한 가역층이 존재하면 $S$에서 국소적으로 몰입
$i : X \to \mathbf{P}^n_S$를 찾을 수 있다. 보조정리
\ref{morphisms-lemma-finite-type-over-affine-ample-very-ample}를 보라. $X \to S$가
고유이므로 사상 $i$는 닫힌 몰입이다. 보조정리
\ref{morphisms-lemma-image-proper-scheme-closed}를 보라.
\end{proof}

\begin{lemma}
\label{morphisms-lemma-H-projective-composition}
H-사영 사상들의 합성은 H-사영이다.
\end{lemma}

\begin{proof}
$X \to Y$와 $Y \to Z$가 H-사영이라고 하자. 그러면 닫힌 몰입
$X \to \mathbf{P}^n_Y$가 $Y$ 위에서, 그리고 닫힌 몰입
$Y \to \mathbf{P}^m_Z$가 $Z$ 위에서 존재한다. 다음 그림을 생각하자.
$$
\xymatrix{
X \ar[r] \ar[d] &
\mathbf{P}^n_Y \ar[r] \ar[dl] &
\mathbf{P}^n_{\mathbf{P}^m_Z} \ar[dl] \ar@{=}[r] &
\mathbf{P}^n_Z \times_Z \mathbf{P}^m_Z \ar[r] &
\mathbf{P}^{nm + n + m}_Z \ar[ddllll] \\
Y \ar[r] \ar[d] & \mathbf{P}^m_Z \ar[dl] & \\
Z & &
}
$$
여기서 오른쪽 끝의 위쪽 수평 화살표는 세그레 매장이다. 구성, 보조정리
\ref{constructions-lemma-segre-embedding}을 보라. 이 그림은 원하는 대로
$X$를 $\mathbf{P}^{nm + n + m}_Z$의 닫힌 부분스킴과 동일시한다.
\end{proof}

\begin{lemma}
\label{morphisms-lemma-H-projective-base-change}
H-사영 사상의 기저변환은 H-사영이다.
\end{lemma}

\begin{proof}
스킴 위의 사영공간의 기저변환은 사영공간이고 닫힌 몰입의 기저변환은
닫힌 몰입이므로 성립한다. 스킴, 보조정리
\ref{schemes-lemma-base-change-immersion}을 보라.
\end{proof}

\begin{lemma}
\label{morphisms-lemma-base-change-projective}
(국소) 사영 사상의 기저변환은 (국소) 사영이다.
\end{lemma}

\begin{proof}
스킴 위의 사영다발의 기저변환은 사영다발이고, 유한형
$\mathcal{O}$-가군의 당김은 유한형이며(가군, 보조정리
\ref{modules-lemma-pullback-finite-type}), 닫힌 몰입의 기저변환은 닫힌
몰입이므로 성립한다. 스킴, 보조정리
\ref{schemes-lemma-base-change-immersion}을 보라. 일부 세부사항은
생략한다.
\end{proof}

\begin{lemma}
\label{morphisms-lemma-projective-quasi-projective}
사영 사상은 준사영이다.
\end{lemma}

\begin{proof}
$f : X \to S$를 사영 사상이라 하자. 닫힌 몰입
$i : X \to \mathbf{P}(\mathcal{E})$를 택하자. 여기서 $\mathcal{E}$는
준연접 유한형 $\mathcal{O}_S$-가군이다. 그러면
$\mathcal{L} = i^*\mathcal{O}_{\mathbf{P}(\mathcal{E})}(1)$은 $f$-매우
충분하다. $f$는 고유이므로(보조정리
\ref{morphisms-lemma-locally-projective-proper}) 준콤팩트이다. 따라서 보조정리
\ref{morphisms-lemma-ample-very-ample}에 의해 $\mathcal{L}$은 $f$-충분하다.
$f$가 고유이므로 유한형이다. 이로써 준사영의 모든 정의 성질이 성립함을
확인했으므로 결론을 얻는다.
\end{proof}

\begin{lemma}
\label{morphisms-lemma-H-quasi-projective-open-H-projective}
$f : X \to S$를 H-준사영 사상이라 하자. 그러면 $f$는
$X \to X' \to S$로 분해되며, 여기서 $X \to X'$은 열린 몰입이고
$X' \to S$는 H-사영이다.
\end{lemma}

\begin{proof}
정의에 의해 $f$를 준콤팩트 몰입 $i : X \to \mathbf{P}^n_S$와 사영
$\mathbf{P}^n_S \to S$의 합성으로 분해할 수 있다. 보조정리
\ref{morphisms-lemma-quasi-compact-immersion}에 의해 닫힌 부분스킴
$X' \subset \mathbf{P}^n_S$가 존재하여 $i$는 열린 몰입 $X \to X'$을
통해 분해된다. 이로써 보조정리가 따른다.
\end{proof}

\begin{lemma}
\label{morphisms-lemma-quasi-projective-open-projective}
$f : X \to S$를 준사영 사상이라 하고 $S$가 준콤팩트이고 준분리라고
하자. 그러면 $f$는 $X \to X' \to S$로 분해되며, 여기서
$X \to X'$은 열린 몰입이고 $X' \to S$는 사영이다.
\end{lemma}

\begin{proof}
$\mathcal{L}$을 $f$-충분하다고 하자. $f$가 유한형이고 $S$가
준콤팩트이므로, $\mathcal{L}^{\otimes n}$이 $f$-매우 충분하게 되는
$n > 0$이 존재한다. 보조정리
\ref{morphisms-lemma-finite-type-ample-very-ample}를 보라. $\mathcal{L}$을
$\mathcal{L}^{\otimes n}$으로 대체하자. $\mathcal{F} = f_*\mathcal{L}$로
쓰자. 이는 준연접 $\mathcal{O}_S$-가군이다. 스킴, 보조정리
\ref{schemes-lemma-push-forward-quasi-coherent}를 보라(준사영 사상은
준콤팩트이고 분리이다. 보조정리
\ref{morphisms-lemma-quasi-projective-properties}를 보라). 성질, 보조정리
\ref{properties-lemma-directed-colimit-finite-presentation}에 의해 유향집합
$I$와 유한형 준연접 $\mathcal{O}_S$-가군들의 계 $\mathcal{E}_i$를 찾아서,
이를 $I$로 첨자화했을 때 $\mathcal{F} = \colim \mathcal{E}_i$가 되게 할
수 있다. 합성
$\psi_i : f^*\mathcal{E}_i \to f^*\mathcal{F} \to \mathcal{L}$을
생각하자. 유한 아핀 열린 덮개
$S = \bigcup_{j = 1, \ldots, m} V_j$를 택하자. 각 $j$에 대해 절단
$$
s_{j, 0}, \ldots, s_{j, n_j} \in
\Gamma(f^{-1}(V_j), \mathcal{L}) = f_*\mathcal{L}(V_j) = \mathcal{F}(V_j)
$$
들을 택하여 $\mathcal{L}$을 $f^{-1}V_j$ 위에서 생성하고 몰입
$$
f^{-1}V_j \longrightarrow \mathbf{P}^{n_j}_{V_j},
$$
을 정의하게 할 수 있다. 보조정리
\ref{morphisms-lemma-very-ample-finite-type-over-affine}를 보라. $i$를 택하여 절단
$e_{j, t} \in \mathcal{E}_i(V_j)$들이 존재하고, 그 상이 $s_{j, t}$와
$\mathcal{F}$에서 같으며, 이는 모든 $j = 1, \ldots, m$ 및
$t = 1, \ldots, n_j$에 대해 성립하게 하자. 그러면 사상 $\psi_i$는
전사이다. 실제로 절단 $f^*e_{j, t}$의 상은 절단 $s_{j, t}$의 상과 같고
후자들이 $\mathcal{L}|_{f^{-1}V_j}$을 생성한다.
따라서 사상
$$
r_{\mathcal{L}, \psi_i} : X \longrightarrow \mathbf{P}(\mathcal{E}_i)
$$
을 얻으며, 이는 $S$ 위의 사상이고 $V_j$ 위에서는 위에서 얻은 몰입이
다음과 같이 분해된다.
$$
f^{-1}V_j \to \mathbf{P}(\mathcal{E}_i)|_{V_j} \to \mathbf{P}^{n_j}_{V_j}
$$
따라서 $r_{\mathcal{L}, \psi_i}|_{V_j}$는 몰입이다. 보조정리
\ref{morphisms-lemma-immersion-permanence}를 보라. $S = \bigcup V_j$이므로
$r_{\mathcal{L}, \psi_i}$는 몰입이라고 결론짓는다.
$r_{\mathcal{L}, \psi_i}$는 준콤팩트임에 유의하라. 실제로
$X \to S$는 준콤팩트이고 $\mathbf{P}(\mathcal{E}_i) \to S$는 분리이다
(스킴, 보조정리 \ref{schemes-lemma-quasi-compact-permanence}를 보라).
보조정리 \ref{morphisms-lemma-quasi-compact-immersion}에 의해 닫힌 부분스킴
$X' \subset \mathbf{P}(\mathcal{E}_i)$가 존재하여 $i$는 열린 몰입
$X \to X'$을 통해 분해된다. 그러면 $X' \to S$는 정의에 의해 사영이고
결론을 얻는다.
\end{proof}

\begin{lemma}
\label{morphisms-lemma-projective-is-quasi-projective-proper}
$S$를 준콤팩트이고 준분리인 스킴이라 하자. $f : X \to S$를 스킴의
사상이라 하자. 그러면
\begin{enumerate}
\item $f$가 사영일 필요충분조건은 $f$가 준사영이고 고유인 것이다.
\item $f$가 H-사영일 필요충분조건은 $f$가 H-준사영이고 고유인 것이다.
\end{enumerate}
\end{lemma}

\begin{proof}
$f$가 사영이면 보조정리 \ref{morphisms-lemma-projective-quasi-projective}에 의해
$f$는 준사영이고, 보조정리 \ref{morphisms-lemma-locally-projective-proper}에 의해
고유이다. 역으로 $X \to S$가 준사영이고 고유이면, 보조정리
\ref{morphisms-lemma-quasi-projective-open-projective}에 의해 열린 몰입
$X \to X'$을 택하여 $X' \to S$가 사영이게 할 수 있다. $X \to S$가
고유이므로 $X$는 $X'$에서 닫혀 있다(보조정리
\ref{morphisms-lemma-image-proper-scheme-closed}). 즉 $X \to X'$은 (열리고) 닫힌
몰입이다. $X'$은 $S$ 위의 사영다발의 닫힌 부분스킴과 동형이므로(정의
\ref{morphisms-definition-projective}), $X$에 대해서도 같은 명제가 성립한다. 즉
$X \to S$는 사영 사상이다. 이로써 (1)이 증명되었다. (2)의 증명도
같으며, 다만 보조정리 \ref{morphisms-lemma-H-projective}와
\ref{morphisms-lemma-H-quasi-projective-open-H-projective}를 사용한다.
\end{proof}

\begin{lemma}
\label{morphisms-lemma-composition-projective}
$f : X \to Y$와 $g : Y \to S$를 스킴의 사상이라 하자. $S$가
준콤팩트이고 준분리이며 $f$와 $g$가 사영이면 $g \circ f$는 사영이다.
\end{lemma}

\begin{proof}
보조정리 \ref{morphisms-lemma-projective-quasi-projective} 및
\ref{morphisms-lemma-locally-projective-proper}에 의해 $f$와 $g$는 준사영이고
고유이다. 보조정리 \ref{morphisms-lemma-composition-proper} 및
\ref{morphisms-lemma-composition-quasi-projective}에 의해 $g \circ f$는 고유이고
준사영이다. 따라서 보조정리
\ref{morphisms-lemma-projective-is-quasi-projective-proper}에 의해 $g \circ f$는
사영이다.
\end{proof}

\begin{lemma}
\label{morphisms-lemma-projective-permanence}
$g : Y \to S$와 $f : X \to Y$를 스킴의 사상이라 하자. $g \circ f$가
사영이고 $g$가 분리이면 $f$는 사영이다.
\end{lemma}

\begin{proof}
닫힌 몰입 $X \to \mathbf{P}(\mathcal{E})$를 택하자. 여기서
$\mathcal{E}$는 준연접 유한형 $\mathcal{O}_S$-가군이다. 그러면 사상
$X \to \mathbf{P}(\mathcal{E}) \times_S Y$를 얻는다. 이 사상은 다음
합성이므로 닫힌 몰입이다.
$$
X \to X \times_S Y \to \mathbf{P}(\mathcal{E}) \times_S Y
$$
여기서 첫 번째 사상은 스킴, 보조정리
\ref{schemes-lemma-semi-diagonal}에 의해 닫힌 몰입이고($g$가 분리라는
사실도 사용한다), 두 번째 사상은 닫힌 몰입의 기저변환이다. 마지막으로
올곱 $\mathbf{P}(\mathcal{E}) \times_S Y$는
$\mathbf{P}(g^*\mathcal{E})$와 동형이고, 당김은 준연접 유한형 가군을
보존한다.
\end{proof}

\begin{lemma}
\label{morphisms-lemma-projective-over-quasi-projective-is-H-projective}
$S$를 충분한 가역층을 허용하는 스킴이라 하자. 그러면
\begin{enumerate}
\item 임의의 사영 사상 $X \to S$는 H-사영이다.
\item 임의의 준사영 사상 $X \to S$는 H-준사영이다.
\end{enumerate}
\end{lemma}

\begin{proof}
$S$에 대한 가정은 $S$가 준콤팩트이고 분리임을 함의한다. 성질, 정의
\ref{properties-definition-ample}, 보조정리
\ref{properties-lemma-ample-immersion-into-proj} 및 구성, 보조정리
\ref{constructions-lemma-proj-separated}을 보라. 따라서 보조정리
\ref{morphisms-lemma-quasi-projective-open-projective}를 적용할 수 있고 (1)이
(2)를 함의함을 안다. $\mathcal{E}$를 유한형 준연접
$\mathcal{O}_S$-가군이라 하자. 사영 사상의 정의에 의해
$\mathbf{P}(\mathcal{E}) \to S$가 H-사영임을 보이면 충분하다.
$\mathcal{E}$가 유한 개의 전역 절단으로 생성되면, 대응하는 전사
$\mathcal{O}_S^{\oplus n} \to \mathcal{E}$는 닫힌 몰입
$$
\mathbf{P}(\mathcal{E}) \longrightarrow
\mathbf{P}(\mathcal{O}_S^{\oplus n}) = \mathbf{P}^{n - 1}_S
$$
을 유도하므로 원하는 결론이 따른다. 일반적으로 $\mathcal{L}$을 $S$
위의 충분한 가역층이라 하자. 성질, 명제
\ref{properties-proposition-characterize-ample}에 의해 정수 $n$이
존재하여
$\mathcal{E} \otimes_{\mathcal{O}_S} \mathcal{L}^{\otimes n}$이 유한 개의
절단으로 대역 생성된다. 구성, 보조정리
\ref{constructions-lemma-twisting-and-proj}에 의해
$\mathbf{P}(\mathcal{E}) =
\mathbf{P}(\mathcal{E} \otimes_{\mathcal{O}_S} \mathcal{L}^{\otimes n})$이므로
증명이 끝난다.
\end{proof}

\begin{lemma}
\label{morphisms-lemma-proper-ample-is-proj}
$f : X \to S$를 보편적으로 닫힌 사상이라 하자. $\mathcal{L}$을
$f$-충분한 가역 $\mathcal{O}_X$-가군이라 하자. 그러면 보조정리
\ref{morphisms-lemma-characterize-relatively-ample}의 표준 사상
$$
r : X
\longrightarrow
\underline{\text{Proj}}_S
\left(
\bigoplus\nolimits_{d \geq 0} f_*\mathcal{L}^{\otimes d}
\right)
$$
은 동형이다.
\end{lemma}

\begin{proof}
$f$-충분한 가역 가군의 존재는 $f$가 준콤팩트임을 강제하므로 $f$는
준콤팩트임에 유의하라. 인용한 보조정리에 의해 사상 $r$은 열린
몰입이다. 한편 보조정리 \ref{morphisms-lemma-image-proper-scheme-closed}에 의해
$r$의 상은 닫혀 있다($r$의 표적은 구성, 보조정리
\ref{constructions-lemma-relative-proj-separated}에 의해 $S$ 위에서
분리이다). 마지막으로 성질, 보조정리
\ref{properties-lemma-ample-immersion-into-proj}에 의해 $r$의 상은
조밀하다(여기서 보조정리
\ref{morphisms-lemma-characterize-relatively-ample}의 증명에서 사상 $r$이 $S$의
아핀 열린집합 위에서 성질, 보조정리
\ref{properties-lemma-map-into-proj}의 표준 사상으로 주어진다는 것도
사용한다). 따라서 $r$은 전사인 열린 몰입, 즉 동형이다.
\end{proof}

\begin{lemma}
\label{morphisms-lemma-proper-ample-delete-affine}
$f : X \to S$를 보편적으로 닫힌 사상이라 하자. $\mathcal{L}$을
$f$-충분한 가역 $\mathcal{O}_X$-가군이라 하자.
$s \in \Gamma(X, \mathcal{L})$라 하자. 그러면 $X_s \to S$는 아핀
사상이다.
\end{lemma}

\begin{proof}
문제는 $S$에서 국소적이므로(보조정리 \ref{morphisms-lemma-characterize-affine})
$S$가 아핀이라고 가정해도 된다. 보조정리
\ref{morphisms-lemma-proper-ample-is-proj}에 의해 $X = \text{Proj}(A)$라고 쓸 수
있으며, 여기서 $A$는 등급환이고 $s$는 $f \in A_1$에 대응하며
$X_s = D_+(f)$이다(성질, 보조정리
\ref{properties-lemma-map-into-proj}). 이는 $\text{Proj}(A)$의 구성으로
보조정리를 증명한다. 구성, 절 \ref{constructions-section-proj}를 보라.
\end{proof}

\section{정수적 사상과 유한 사상}
\label{morphisms-section-integral}

\noindent
환 사상 $R \to A$에 대하여 $A$의 모든 원소가 $R$의 계수를 갖는
일계수 방정식을 만족하면 이 사상을 {\it 정수적}이라고 함을 상기하자.
환 사상 $R \to A$에 대하여 $A$가 유한 $R$-가군이면 이 사상을
{\it 유한}이라고 함을 상기하자. 대수학, 정의
\ref{algebra-definition-integral-ring-map}를 보라.

\begin{definition}
\label{morphisms-definition-integral}
$f : X \to S$를 스킴의 사상이라 하자.
\begin{enumerate}
\item $f$를 {\it 정수적}이라고 하는 것은 $f$가 아핀이고, 모든 아핀
열린집합 $\Spec(R) = V \subset S$에 대하여 그 역상
$\Spec(A) = f^{-1}(V) \subset X$가 아핀이며 연관된 환 사상
$R \to A$가 정수적이라는 뜻이다.
\item $f$를 {\it 유한}이라고 하는 것은 $f$가 아핀이고, 모든 아핀
열린집합 $\Spec(R) = V \subset S$에 대하여 그 역상
$\Spec(A) = f^{-1}(V) \subset X$가 아핀이며 연관된 환 사상
$R \to A$가 유한이라는 뜻이다.
\end{enumerate}
\end{definition}

\noindent
정수적 사상과 유한 사상이 분리이고 준콤팩트임은 명백하다. 또한 유한
사상이 유한형 사상임도 명백하다. 이 절의 보조정리 대부분은 완전히
표준적이다. 다만 절 끝의 흥미로운 보조정리
\ref{morphisms-lemma-integral-universally-closed}에 주목하라.

\begin{lemma}
\label{morphisms-lemma-integral-local}
$f : X \to S$를 스킴의 사상이라 하자. 다음 조건들은 동치이다.
\begin{enumerate}
\item 사상 $f$는 정수적이다.
\item 아핀 열린 덮개 $S = \bigcup U_i$가 존재하여 각 $f^{-1}(U_i)$가
아핀이고
$\mathcal{O}_S(U_i) \to \mathcal{O}_X(f^{-1}(U_i))$가 정수적이다.
\item 열린 덮개 $S = \bigcup U_i$가 존재하여 각
$f^{-1}(U_i) \to U_i$가 정수적이다.
\end{enumerate}
더욱이 $f$가 정수적이면 모든 열린 부분스킴 $U \subset S$에 대하여
사상 $f : f^{-1}(U) \to U$는 정수적이다.
\end{lemma}

\begin{proof}
대수학, 보조정리 \ref{algebra-lemma-integral-local}를 보라.
몇 가지 세부사항은 생략한다.
\end{proof}

\begin{lemma}
\label{morphisms-lemma-finite-local}
$f : X \to S$를 스킴의 사상이라 하자. 다음 조건들은 동치이다.
\begin{enumerate}
\item 사상 $f$는 유한이다.
\item 아핀 열린 덮개 $S = \bigcup U_i$가 존재하여 각 $f^{-1}(U_i)$가
아핀이고
$\mathcal{O}_S(U_i) \to \mathcal{O}_X(f^{-1}(U_i))$가 유한이다.
\item 열린 덮개 $S = \bigcup U_i$가 존재하여 각
$f^{-1}(U_i) \to U_i$가 유한이다.
\end{enumerate}
더욱이 $f$가 유한이면 모든 열린 부분스킴 $U \subset S$에 대하여
사상 $f : f^{-1}(U) \to U$는 유한이다.
\end{lemma}

\begin{proof}
대수학, 보조정리 \ref{algebra-lemma-integral-local}를 보라.
몇 가지 세부사항은 생략한다.
\end{proof}

\begin{lemma}
\label{morphisms-lemma-finite-integral}
유한 사상은 정수적이다. 국소적으로 유한형인 정수적 사상은 유한이다.
\end{lemma}

\begin{proof}
대수학, 보조정리 \ref{algebra-lemma-finite-is-integral} 및 보조정리
\ref{algebra-lemma-characterize-finite-in-terms-of-integral}를 보라.
\end{proof}

\begin{lemma}
\label{morphisms-lemma-composition-finite}
유한 사상들의 합성은 유한이다. 정수적 사상들에 대해서도 마찬가지이다.
\end{lemma}

\begin{proof}
대수학, 보조정리 \ref{algebra-lemma-finite-transitive} 및
\ref{algebra-lemma-integral-transitive}를 보라.
\end{proof}

\begin{lemma}
\label{morphisms-lemma-base-change-finite}
유한 사상의 기저변환은 유한이다. 정수적 사상에 대해서도 마찬가지이다.
\end{lemma}

\begin{proof}
대수학, 보조정리 \ref{algebra-lemma-base-change-integral}을 보라.
\end{proof}

\begin{lemma}
\label{morphisms-lemma-integral-universally-closed}
\begin{slogan}
정수적 $=$ 아핀 $+$ 보편적으로 닫힌
\end{slogan}
$f : X \to S$를 스킴의 사상이라 하자. 다음 조건들은 동치이다.
\begin{enumerate}
\item $f$는 정수적이다.
\item $f$는 아핀이고 보편적으로 닫혀 있다.
\end{enumerate}
\end{lemma}

\begin{proof}
(1)을 가정하자. 정수적 사상은 정의에 의해 아핀이다. 정수적 사상의
기저변환은 정수적이므로 (2)를 증명하려면 정수적 사상이 닫혀 있음을
보이면 충분하다. 이는 대수학, 보조정리
\ref{algebra-lemma-integral-going-up} 및
\ref{algebra-lemma-going-up-closed}에서 따른다.

\medskip\noindent
(2)를 가정하자. $f$가 사상 $f : \Spec(A) \to \Spec(R)$이고 이것이
환 사상 $R \to A$에서 온다고 가정해도 된다. $a$를 $A$의 원소라 하자.
$a$가 $R$ 위에서 정수적임, 즉 핵 $I$ 안에 일계수 다항식이 있음을
보여야 한다. 여기서 사상 $R[x] \to A$는 $x$를 $a$로 보낸다.
환 $B = A[x]/(ax -1)$을 생각하고 $J$를 합성
$R[x]\to A[x] \to B$의 핵이라 하자. $f\in J$이면
$q\in A[x]$가 존재하여 $f = (ax-1)q$가 $A[x]$에서 성립한다. 따라서
$f = \sum_i f_ix^i$이고 $q = \sum_iq_ix^i$이면 모든 $i \geq 0$에
대하여 $f_i = aq_{i-1} - q_i$이다. $n \geq \deg q + 1$에 대하여 다항식
$$
\sum\nolimits_{i \geq 0} f_i x^{n - i} =
\sum\nolimits_{i \geq 0} (a q_{i - 1} - q_i) x^{n - i} =
(a - x) \sum\nolimits_{i \geq 0} q_i x^{n - i - 1}
$$
은 분명히 $I$에 속한다. $f_0 = 1$이면 이 다항식은 일계수이기도 하므로,
$J$가 상수항 $1$인 다항식을 포함함을 증명하는 것으로 귀착된다.
$\Spec(R[x]/(J + (x))$가 공집합임을 증명하여 이를 보이겠다.


\medskip\noindent
$f$가 보편적으로 닫혀 있으므로 기저변환
$\Spec(A[x]) \to \Spec(R[x])$는 닫혀 있다. 따라서 닫힌 부분집합
$\Spec(B) \subset \Spec(A[x])$의 상은
$\Spec(R[x]/J) \subset \Spec(R[x])$라는 닫힌 부분집합이다. 예
\ref{morphisms-example-scheme-theoretic-image} 및 보조정리
\ref{morphisms-lemma-quasi-compact-scheme-theoretic-image}를 보라. 특히
$\Spec(B) \to \Spec(R[x]/J)$는 전사이다. 모든 정사각형이 당김인 다음
도표를 생각하자.
$$
\xymatrix{
\Spec(B) \ar@{->>}[r]^g &
\Spec(R[x]/J)  \ar[r] &
\Spec(R[x])\\
\emptyset \ar[u] \ar[r] &
\Spec(R[x]/(J + (x)))\ar[u] \ar[r] &
\Spec(R) \ar[u]^0
}
$$
왼쪽 아래 모서리는 공집합이다. 실제로 이는 $R\otimes_{R[x]} B$의
스펙트럼인데, 사상 $R[x]\to B$는 $x$를 가역원으로 보내고 사상
$R[x]\to R$는 $x$를 $0$으로 보낸다. $g$가 전사이므로 이는
$\Spec(R[x]/(J + (x)))$가 공집합임을 함의하며, 이것이 보이고자 한
것이다.
\end{proof}

\begin{lemma}
\label{morphisms-lemma-integral-fibres}
$f : X \to S$를 정수적 사상이라 하자. 그러면 $X$의 모든 점은 그
올에서 닫힌점이다.
\end{lemma}

\begin{proof}
대수학, 보조정리 \ref{algebra-lemma-integral-no-inclusion}을 보라.
\end{proof}

\begin{lemma}
\label{morphisms-lemma-integral-dimension}
$f : X \to Y$를 정수적 사상이라 하자. 그러면
$\dim(X) \leq \dim(Y)$이다. $f$가 전사이면 $\dim(X) = \dim(Y)$이다.
\end{lemma}

\begin{proof}
$X$와 $Y$의 차원은 아핀 열린 덮개의 원소들의 차원들의 상한이므로,
$Y$와 $X$가 아핀이라고 가정해도 된다. 부등식은 대수학, 보조정리
\ref{algebra-lemma-integral-dim-up}에서 따른다. 그러면 등식은 대수학,
보조정리 \ref{algebra-lemma-dimension-going-up} 및
\ref{algebra-lemma-integral-going-up}에서 따른다.
\end{proof}

\begin{lemma}
\label{morphisms-lemma-finite-quasi-finite}
유한 사상은 준유한이다.
\end{lemma}

\begin{proof}
이는 대수학, 보조정리 \ref{algebra-lemma-quasi-finite}와 보조정리
\ref{morphisms-lemma-quasi-finite-locally-quasi-compact}에서 함의된다. 또는
보조정리 \ref{morphisms-lemma-integral-fibres}에 의해 올의 모든 점은 닫힌점이고
(유한 사상이 정수적이라는 사실도 사용한다), 보조정리
\ref{morphisms-lemma-quasi-finite-at-point-characterize} (3)을 사용하면 $f$가
$x$에서 준유한임을 알 수 있고, 이는 모든 $x \in X$에 대해 성립한다.
\end{proof}

\begin{lemma}
\label{morphisms-lemma-finite-proper}
$f : X \to S$를 스킴의 사상이라 하자. 다음 조건들은 동치이다.
\begin{enumerate}
\item $f$는 유한이다.
\item $f$는 아핀이고 고유이다.
\end{enumerate}
\end{lemma}

\begin{proof}
이는 보조정리 \ref{morphisms-lemma-integral-universally-closed}, 유한 사상이
정수적이고 분리라는 사실, 고유 사상은 유한형이고 분리이고 보편적으로
닫힌 사상과 같다는 사실, 그리고 유한형인 정수적 사상은 유한이라는
사실(보조정리 \ref{morphisms-lemma-finite-integral})에서 형식적으로 따른다.
\end{proof}

\begin{lemma}
\label{morphisms-lemma-closed-immersion-finite}
닫힌 몰입은 유한이다(따라서 더욱이 정수적이다).
\end{lemma}

\begin{proof}
닫힌 몰입은 아핀이고(보조정리 \ref{morphisms-lemma-closed-immersion-affine}),
전사 환 사상은 유한이고 정수적이므로 참이다.
\end{proof}

\begin{lemma}
\label{morphisms-lemma-finite-union-finite}
$X_i \to Y$, $i = 1, \ldots, n$을 스킴의 유한 사상들이라 하자.
그러면 $X_1 \amalg \ldots \amalg X_n \to Y$도 유한이다.
\end{lemma}

\begin{proof}
$R \to A_i$, $i = 1, \ldots, n$이 유한 환 사상들이면
$R \to A_1 \times \ldots \times A_n$도 유한이라는 대수적 사실에서
따른다.
\end{proof}

\begin{lemma}
\label{morphisms-lemma-finite-permanence}
$f : X \to Y$와 $g : Y \to Z$를 사상들이라 하자.
\begin{enumerate}
\item $g \circ f$가 유한이고 $g$가 분리이면 $f$는 유한이다.
\item $g \circ f$가 정수적이고 $g$가 분리이면 $f$는 정수적이다.
\end{enumerate}
\end{lemma}

\begin{proof}
$g \circ f$가 유한이고(각각 정수적이고) $g$가 분리라고 가정하자.
기저변환 $X \times_Z Y \to Y$는 보조정리
\ref{morphisms-lemma-base-change-finite}에 의해 유한이다(각각 정수적이다).
사상 $X \to X \times_Z Y$는 $Y \to Z$가 분리이므로 닫힌 몰입이다.
스킴, 보조정리 \ref{schemes-lemma-section-immersion}을 보라. 닫힌 몰입은
유한이다(각각 정수적이다). 보조정리
\ref{morphisms-lemma-closed-immersion-finite}를 보라. 유한 사상들의(각각 정수적
사상들의) 합성은 유한이다(각각 정수적이다). 보조정리
\ref{morphisms-lemma-composition-finite}를 보라. 따라서 결론을 얻는다.
\end{proof}

\begin{lemma}
\label{morphisms-lemma-finite-monomorphism-closed}
$f : X \to Y$를 스킴의 사상이라 하자. $f$가 유한이고 모노사상이면
$f$는 닫힌 몰입이다.
\end{lemma}

\begin{proof}
이는 대수학, 보조정리
\ref{algebra-lemma-finite-epimorphism-surjective}로 귀착된다.
\end{proof}

\begin{lemma}
\label{morphisms-lemma-finite-projective}
유한 사상은 사영이다.
\end{lemma}

\begin{proof}
$f : X \to S$를 유한 사상이라 하자. 그러면 $f_*\mathcal{O}_X$는
준연접 $\mathcal{O}_S$-가군이고(보조정리
\ref{morphisms-lemma-affine-equivalence-algebras}), 유한형이다(유한 사상의
정의와 성질, 보조정리
\ref{properties-lemma-finite-type-module}에 의한다). 닫힌 몰입
$$
\sigma :
X
\longrightarrow
\mathbf{P}(f_*\mathcal{O}_X) =
\underline{\text{Proj}}_S(\text{Sym}^*_{\mathcal{O}_S}(f_*\mathcal{O}_X))
$$
이 $S$ 위에 존재한다고 주장하며, 이로써 증명이 끝난다. 실제로
$\sigma$를 (구성, 보조정리
\ref{constructions-lemma-apply-relative}를 통해) 전사
$$
f^*f_*\mathcal{O}_X \longrightarrow \mathcal{O}_X
$$
에 대응하는 사상으로 잡는다. 이 전사는 수반 사상
$f^*f_* \to \text{id}$에서 온다. 그러면 $\sigma$는 닫힌 몰입이다.
실제로 $S$에서 아핀 국소적으로 $X = \Spec(A)$ 및 $S = \Spec(R)$로
쓸 수 있다. $X$가 $S$ 위에서 유한이므로,
$a_1, \ldots, a_n \in A$를 택하여 $A$를 $R$-가군으로서 생성하게 할
수 있다. $R$-대수 사상
$R[T_0, T_1, \ldots, T_n] \to \text{Sym}_R^*(A)$ 가운데 $T_0$를
$1$로 보내고 $T_i$를 $a_i$로 보내는 사상은 전사인데, 여기서
$1 \leq i \leq n$이다. 따라서 구성, 보조정리
\ref{constructions-lemma-surjective-graded-rings-map-proj}에 의해
$\mathbf{P}(f_*\mathcal{O}_X)$는 $\mathbf{P}^n_S$의 닫힌 부분스킴이다.
그러므로 유도된 사상 $X \to \mathbf{P}^n_S$가 닫힌 몰입임을 보이면
충분하다(예를 들어 보조정리
\ref{morphisms-lemma-closed-immersion-ideals}를 보라). 독자는
$X \to \mathbf{P}^n_S$의 상이 열린집합
$D_+(T_0) \cong \mathbf{A}^n_S$에 포함되고,
$X \to \mathbf{A}^n_S$가 전사 $R$-대수 사상
$R[x_1, \ldots, x_n] \to \text{Sym}_R^*(A)$에 대응하며 이 사상이
$x_i$를 $a_i$로 보냄을 확인할 수 있다. 따라서
$X \to \mathbf{P}^n_S$는 몰입이다(닫힌 몰입과 열린
몰입의 합성으로서). $X$가 $S$ 위에서 유한이므로 이 몰입의 상은
닫혀 있다(여기서는 보조정리 \ref{morphisms-lemma-finite-proper}, 보조정리
\ref{morphisms-lemma-image-proper-scheme-closed}, 구성, 보조정리
\ref{constructions-lemma-projective-space-separated}을 사용한다).
스킴, 보조정리 \ref{schemes-lemma-immersion-when-closed}에 의해
결론을 얻는다.
\end{proof}

\section{보편 위상동형사상}
\label{morphisms-section-universal-homeomorphisms}

\noindent
보편 위상동형사상은 실제로 정수적이고 보편 단사이며 전사인 사상에
지나지 않으므로 다음 정의는 사실 불필요하다. 보조정리
\ref{morphisms-lemma-universal-homeomorphism}을 보라.

\begin{definition}
\label{morphisms-definition-universal-homeomorphism}
스킴의 사상 $f : X \to Y$에 대하여 기저변환
$f' : Y' \times_Y X \to Y'$가 모든 사상 $Y' \to Y$에 대해
위상동형이면 이 사상을 {\it 보편 위상동형사상}이라고 한다.
\end{definition}

\noindent
먼저 필수적인 보조정리들을 서술한다.

\begin{lemma}
\label{morphisms-lemma-base-change-universal-homeomorphism}
스킴의 보편 위상동형사상을 스킴의 임의의 사상에 따라 기저변환한
사상은 보편 위상동형사상이다.
\end{lemma}

\begin{proof}
정의에서 즉시 따른다.
\end{proof}

\begin{lemma}
\label{morphisms-lemma-composition-universal-homeomorphism}
스킴의 보편 위상동형사상 두 개를 합성하면 보편 위상동형사상이다.
\end{lemma}

\begin{proof}
생략한다.
\end{proof}

\noindent
다음 간단한 보조정리가 보편 위상동형사상을 특징짓는 핵심이다.

\begin{lemma}
\label{morphisms-lemma-homeomorphism-affine}
$f : X \to Y$를 스킴의 사상이라 하자. $f$가 $Y$의 어떤 닫힌
부분집합 위로의 위상동형이면 $f$는 아핀이다.
\end{lemma}

\begin{proof}
$y \in Y$를 점이라 하자. $y \not \in f(X)$이면, $y$의 아핀 근방
가운데 $f(X)$와 서로소인 것이 존재한다. $y \in f(X)$이면 $x \in X$를
$X$의 점 가운데 $y$로 사상되는 유일한 점이라 하자. $y \in V$를 아핀
열린 근방이라 하자. $U \subset X$를 $x$의 아핀 열린 근방 가운데
$V$ 안으로 사상되는 것으로 잡자. $f(U) \subset V \cap f(X)$는
$f$에 대한 가정에 의해
유도된 위상에서 열려 있으므로, $h \in \Gamma(V, \mathcal{O}_Y)$를
택하여 $y \in D(h)$이고 $D(h) \cap f(X) \subset f(U)$가 되게 할 수
있다. $h' \in \Gamma(U, \mathcal{O}_X)$를 $f^\sharp(h)$의 $U$로의
제한이라 하자. 그러면 $D(h') \subset U$는 $f^{-1}(D(h))$와 같음을
알 수 있다. 다시 말해 $Y$의 모든 점은 역상이 아핀인 열린 근방을
갖는다. 따라서 $f$는 아핀이다. 보조정리
\ref{morphisms-lemma-characterize-affine}을 보라.
\end{proof}

\begin{lemma}
\label{morphisms-lemma-universal-homeomorphism}
$f : X \to Y$를 스킴의 사상이라 하자. 다음 조건들은 동치이다.
\begin{enumerate}
\item $f$는 보편 위상동형사상이다.
\item $f$는 정수적이고 보편 단사이며 전사이다.
\end{enumerate}
\end{lemma}

\begin{proof}
$f$가 보편 위상동형사상이라고 가정하자. 보조정리
\ref{morphisms-lemma-homeomorphism-affine}에 의해 $f$는 아핀이다. $f$가 명백히
보편적으로 닫혀 있으므로 보조정리
\ref{morphisms-lemma-integral-universally-closed}에 의해 $f$는 정수적이다.
또한 $f$가 보편 단사이고 전사임도 명백하다.

\medskip\noindent
$f$가 정수적이고 보편 단사이며 전사라고 가정하자. 보조정리
\ref{morphisms-lemma-integral-universally-closed}에 의해 $f$는 보편적으로
닫혀 있다. 또한 이는 보편 전단사이므로(보조정리
\ref{morphisms-lemma-base-change-surjective}를 보라), 보편 위상동형사상임을
알 수 있다.
\end{proof}

\begin{lemma}
\label{morphisms-lemma-reduction-universal-homeomorphism}
$X$를 스킴이라 하자. 표준 닫힌 몰입 $X_{red} \to X$는 보편
위상동형사상이다(스킴, 정의
\ref{schemes-definition-reduced-induced-scheme}를 보라).
\end{lemma}

\begin{proof}
생략한다.
\end{proof}

\begin{lemma}
\label{morphisms-lemma-check-closed-infinitesimally}
$f : X \to S$와 $S' \to S$를 스킴의 사상들이라 하자. 다음을
가정하자.
\begin{enumerate}
\item $S' \to S$는 닫힌 몰입이다.
\item $S' \to S$는 점들 위에서 전단사이다.
\item $X \times_S S' \to S'$는 닫힌 몰입이다.
\item $X \to S$는 유한형이거나 $S' \to S$는 유한표현이다.
\end{enumerate}
그러면 $f : X \to S$는 닫힌 몰입이다.
\end{lemma}

\begin{proof}
가정 (1)과 (2)는 $S' \to S$가 보편 위상동형사상임을 함의한다
(예를 들어 $S_{red} = S'_{red}$이고 보조정리
\ref{morphisms-lemma-reduction-universal-homeomorphism}를 사용하기 때문이다).
따라서 (3)은 $X \to S$가 $S$의 닫힌 부분집합 위로의 위상동형임을
함의한다. 그러면 보조정리 \ref{morphisms-lemma-homeomorphism-affine}에 의해
$X \to S$는 아핀이다. $U \subset S$를 아핀 열린집합이라 하고
$U = \Spec(A)$로 쓰자. 그러면 (1)에 의해
$S' = \Spec(A/I)$이며, (2)에 의해 $I$는 국소 멱영 아이디얼이다.
$f$가 아핀이므로 $f^{-1}(U) = \Spec(B)$임을 알 수 있다. 가정 (4)는
$B$가 유한형 $A$-대수이거나(보조정리
\ref{morphisms-lemma-locally-finite-type-characterize}) $I$가 유한생성임을
말해 준다(보조정리
\ref{morphisms-lemma-closed-immersion-finite-presentation}). 가정 (3)은
$A/I \to B/IB$가 전사라는 것이다. $A \to B$가 유한형이면 대수학,
보조정리 \ref{algebra-lemma-surjective-mod-locally-nilpotent}에서,
$I$가 유한생성이고 따라서 멱영이면 대수학, 보조정리
\ref{algebra-lemma-NAK}에서 $A \to B$가 전사임을 얻는다. 이는 $f$가
닫힌 몰입이라는 뜻이다. 보조정리 \ref{morphisms-lemma-closed-immersion}을 보라.
\end{proof}

\begin{lemma}
\label{morphisms-lemma-permanence-universal-homeomorphism}
$f : X \to Z$를 두 사상 $g : X \to Y$와 $h : Y \to Z$의 합성이라
하자. 사상들 $\{f, g, h\}$ 가운데 둘이 보편 위상동형사상이면 나머지
하나도 그러하다.
\end{lemma}

\begin{proof}
$g$와 $h$가 모두 보편 위상동형사상이면, 보조정리
\ref{morphisms-lemma-composition-universal-homeomorphism}에 의해 $f$도 그러하다.

\medskip\noindent
$f$와 $g$가 모두 보편 위상동형사상이라고 가정하자. $h$도
그러함을 보이고자 한다. 이제 스킴의 임의의 사상
$\alpha : Z' \to Z$에 따라 도표를 기저변환하면 모든 정사각형이
카르테시안인 다음 도표를 얻는다.
$$
\xymatrix{
X' \ar[r]^{g'} \ar[d] & Y' \ar[r]^{h'} \ar[d] & Z' \ar[d] \\
X \ar[r]^{g} & Y \ar[r]^{h} & Z.
}
$$
가정에 의해 합성 $f'= h' \circ g' : X' \to Z'$와
$g' : X' \to Y'$는 위상동형이므로 $h'$도 위상동형이다. 이로써
$h$가 보편 위상동형사상임이 증명된다.

\medskip\noindent
마지막으로 $f$와 $h$가 보편 위상동형사상이라고 가정하자. $g$가 보편
위상동형사상임을 보이고자 한다. $\beta : Y' \to Y$를 스킴의 임의의
사상이라 하자. 모든 정사각형이 카르테시안인 다음 도표를 얻는다.
$$
\xymatrix{
X' \ar[r]^{g'} \ar[d] & Y' \ar[d]^{\gamma} & \\
X'' \ar[r]^{g''} \ar[d] & Y'' \ar[r]^{h''} \ar[d] & Y' \ar[d]^{h \circ \beta} \\
X \ar[r]^{g} & Y \ar[r]^{h} & Z.
}
$$
여기서 사상 $\gamma : Y' \to Y''$는 올곱의 보편 성질과 두 사상
$id_{Y'} : Y' \to Y'$ 및 $\beta : Y' \to Y$로 정의된다. $g'$가
위상동형임을 증명하겠다. 위상동형이라는 성질은 3개 중 2개 성질을
가지므로 $g''$가 위상동형임을 알 수 있다. 위쪽 정사각형을 보면
$\gamma$가 보편 위상동형사상임을 보이는 것으로 충분하다. $h''$가
위상동형이므로 보조정리 \ref{morphisms-lemma-homeomorphism-affine}에 의해 아핀
사상이고, 따라서 더욱이 분리이다(보조정리
\ref{morphisms-lemma-affine-separated}). $h'' \circ \gamma$가 항등사상이므로,
스킴, 보조정리 \ref{schemes-lemma-section-immersion}에 의해
$\gamma$는 닫힌 몰입임을 알 수 있다. $h''$가 전단사이므로
$\gamma$는 전단사인 닫힌 몰입이며, 따라서 원하는 대로 보편
위상동형사상이다(예를 들어 보조정리
\ref{morphisms-lemma-universal-homeomorphism}의 특징지음을 사용한다).
\end{proof}

\section{아핀 스킴의 보편 위상동형사상}
\label{morphisms-section-universal-homeomorphisms-affine}

\noindent
이 절에서는 아핀 스킴의 보편 위상동형사상을 특징짓는다.

\begin{lemma}
\label{morphisms-lemma-subalgebra-inherits}
$A \to B$를 환 사상이라 하고 유도된 스킴 사상
$f : \Spec(B) \to \Spec(A)$가 각각 보편 위상동형사상, 잉여체 위의
동형을 유도하는 보편 위상동형사상, 보편적으로 닫힌 사상, 또는
보편적으로 닫히고 보편 단사인 사상이라고 하자. 그러면 모든
$A$-부분대수 $B' \subset B$에 대하여
$f' : \Spec(B') \to \Spec(A)$에도 같은 명제가 성립한다.
\end{lemma}

\begin{proof}
$f$가 보편적으로 닫혀 있으면 보조정리
\ref{morphisms-lemma-integral-universally-closed}에 의해 $B$는 $A$ 위에서
정수적이다. 따라서 $B'$은 $A$ 위에서 정수적이고 $f'$은
보편적으로 닫혀 있다(같은 보조정리에 의한다). 이로써 $f$가
보편적으로 닫힌 경우가 증명된다.

\medskip\noindent
계속해서 $B$가 $B'$ 위에서 정수적임을 안다(대수학, 보조정리
\ref{algebra-lemma-integral-permanence}). 이는
$\Spec(B) \to \Spec(B')$가 전사임을 함의한다(대수학, 보조정리
\ref{algebra-lemma-integral-overring-surjective}). 따라서 $A \to B$가
잉여체들의 순수 비분리 확대를 유도하면 $A \to B'$도 그러하다. 이로써
$f$가 보편적으로 닫히고 보편 단사인 경우가 증명된다. 보조정리
\ref{morphisms-lemma-universally-injective}를 보라.

\medskip\noindent
$f$가 보편 위상동형사상인 경우는 위의 고찰, 보조정리
\ref{morphisms-lemma-universal-homeomorphism}, 그리고 $f$가 전사이면 $f'$도
전사라는 명백한 관찰에서 따른다.

\medskip\noindent
$A \to B$가 잉여체 위의 동형들을 유도하면 $A \to B'$도 그러하다
(둘째 문단의 논증을 보라). 이로써 나머지 경우에도 보조정리가
성립함을 안다.
\end{proof}

\begin{lemma}
\label{morphisms-lemma-colimit-inherits}
$A$를 환이라 하자. $B = \colim B_\lambda$를 $A$-대수들의 여과
여극한이라 하자. 각 $f_\lambda : \Spec(B_\lambda) \to \Spec(A)$가
각각 보편 위상동형사상, 잉여체 위의 동형을 유도하는 보편
위상동형사상, 보편적으로 닫힌 사상, 또는 보편적으로 닫히고 보편
단사인 사상이라 하자. 그러면 $f : \Spec(B) \to \Spec(A)$에도 같은
명제가 성립한다.
\end{lemma}

\begin{proof}
$f_\lambda$가 보편적으로 닫혀 있으면 보조정리
\ref{morphisms-lemma-integral-universally-closed}에 의해 $B_\lambda$는 $A$
위에서 정수적이다. 따라서 $B$는 $A$ 위에서 정수적이고 $f$는
보편적으로 닫혀 있다(같은 보조정리에 의한다). 이로써 각
$f_\lambda$가 보편적으로 닫힌 경우가 증명된다.

\medskip\noindent
$\mathfrak q \subset B$를 $\mathfrak p \subset A$ 위에 놓인
소아이디얼이라 하고, 그 역상을
$\mathfrak q_\lambda \subset B_\lambda$로 나타내자. 그러면
$\kappa(\mathfrak q) = \colim \kappa(\mathfrak q_\lambda)$이다. 따라서
$A \to B_\lambda$가 잉여체들의 순수 비분리 확대를 유도하면
$A \to B$도 그러하다. 이로써 $f_\lambda$가 보편적으로 닫히고 보편
단사인 경우가 증명된다. 보조정리 \ref{morphisms-lemma-universally-injective}를
보라.

\medskip\noindent
$f$가 보편 위상동형사상인 경우는 위의 고찰과 보조정리
\ref{morphisms-lemma-universal-homeomorphism}, 그리고 $B$의 소아이디얼이 모든
$B_\lambda$의 양립하는 소아이디얼 열과 같은 것이라는 사실을 함께
사용하면 따른다.

\medskip\noindent
$A \to B_\lambda$가 잉여체 위의 동형들을 유도하면 $A \to B$도
그러하다(둘째 문단의 논증을 보라). 이로써 나머지 경우에도
보조정리가 성립함을 안다.
\end{proof}

\begin{lemma}
\label{morphisms-lemma-special-elements-and-localization}
$A \subset B$를 환 확대라 하자. $S \subset A$를 곱셈적 부분집합이라
하자. $n \geq 1$이고 $b_i \in B$이며 $1 \leq i \leq n$이라 하자.
다음을 만족하는 임의의 $x \in S^{-1}B$는
$$
x \not \in S^{-1}A\text{ and } b_i x^i \in S^{-1}A\text{ for }i = 1, \ldots, n
$$
$s^{-1}y$와 같으며, 여기서 $s \in S$이고 $y \in B$는 다음을
만족한다.
$$
y \not \in A\text{ and } b_i y^i \in A\text{ for }i = 1, \ldots, n
$$
\end{lemma}

\begin{proof}
생략한다. 도움말: 분모를 없애라.
\end{proof}

\begin{lemma}
\label{morphisms-lemma-nth-and-nplusone-implies-square-and-cube}
$A \subset B$를 환 확대라 하자. $b \in B$, $b \not \in A$와 정수
$n \geq 2$가 존재하여 $b^n \in A$ 및 $b^{n + 1} \in A$이면,
$b' \in B$, $b' \not \in A$가 존재하여 $(b')^2 \in A$ 및
$(b')^3 \in A$이다.
\end{lemma}

\begin{proof}
$b$와 $n$을 보조정리에서와 같이 잡자. 그러면 $b$의 충분히 큰 모든
거듭제곱은 $A$에 속한다. 실제로
$(b^n)^k(b^{n + 1})^i = b^{(k + i)n + i}$이며, 이는 모든 거듭제곱
$b^m$이 $m \geq n^2$일 때 $A$에 속함을 함의한다.
따라서 $i \geq 1$이 $b^i \not \in A$를 만족하는 가장 큰 정수이면
$(b^i)^2 \in A$ 및 $(b^i)^3 \in A$이다.
\end{proof}

\begin{lemma}
\label{morphisms-lemma-square-and-cube}
$A \subset B$를 환 확대라 하고
$\Spec(B) \to \Spec(A)$가 잉여체 위의 동형들을 유도하는 보편
위상동형사상이라고 하자. $A \not = B$이면 $b \in B$,
$b \not \in A$가 존재하여 $b^2 \in A$ 및 $b^3 \in A$이다.
\end{lemma}

\begin{proof}
$A \subset B$가 정수적임을 상기하자(보조정리
\ref{morphisms-lemma-integral-universally-closed}). 보조정리
\ref{morphisms-lemma-subalgebra-inherits}에 의해 $B$가 $A$ 위에서 하나의
원소로 생성된다고 가정해도 된다. 따라서 $B$는 $A$ 위에서 유한이다
(대수학, 보조정리
\ref{algebra-lemma-characterize-finite-in-terms-of-integral}). 따라서
$B/A$의 지지집합은 $A$-가군으로서 닫혀 있고 공집합이 아니다(대수학,
보조정리 \ref{algebra-lemma-support-closed} 및
\ref{algebra-lemma-support-zero}). $\mathfrak p \subset A$를 이
지지집합의 극소 소아이디얼이라 하자. $A \subset B$를
$A_\mathfrak p \subset B_\mathfrak p$로 대체한 뒤에는(보조정리
\ref{morphisms-lemma-special-elements-and-localization}에 의해 허용된다)
$(A, \mathfrak m)$이 국소환이고, $B$가 $A$ 위에서 유한이며,
$B/A$의 지지집합이 $\{\mathfrak m\}$이고 이를 $A$-가군으로 본다고
가정해도 된다. $B/A$가 유한 가군이므로
$I = \text{Ann}_A(B/A)$는 $\mathfrak m = \sqrt{I}$를 만족한다
(대수학, 보조정리 \ref{algebra-lemma-support-closed}).
$\mathfrak m' \subset B$를 $\mathfrak m$ 위에 놓인 유일한
소아이디얼이라 하자. $\Spec(B) \to \Spec(A)$가 위상동형이므로
$\mathfrak m' = \sqrt{IB}$를 얻는다. $f \in \mathfrak m'$에 대해
$n \geq 1$을 택하여 $f^n \in IB$가 되게 하자. 그러면
$f^{n + 1} \in IB$이기도 하다. $IB \subset A$이며 이는 $I$의 선택에
의한 것이므로
$f^n, f^{n + 1} \in A$라고 결론짓는다. 보조정리
\ref{morphisms-lemma-nth-and-nplusone-implies-square-and-cube}를 사용하면
$\mathfrak m' \not \subset A$일 때 보조정리의 결론이 참임을 얻는다.
그러나 $\mathfrak m' \subset A$이면 $\mathfrak m' = \mathfrak m$이고,
가정에 의해 잉여체들도 동형이므로 $A = B$라고 결론짓는다. 이 모순으로
증명이 끝난다.
\end{proof}

\begin{lemma}
\label{morphisms-lemma-pth-power-and-multiple}
$A \subset B$를 환 확대라 하고 $\Spec(B) \to \Spec(A)$가 보편
위상동형사상이라고 하자. $A \not = B$이면 $b \in B$,
$b \not \in A$가 존재하여 $b^2 \in A$ 및 $b^3 \in A$이거나,
소수 $p$와 $b \in B$, $b \not \in A$가 존재하여 $pb \in A$ 및
$b^p \in A$이다.
\end{lemma}

\begin{proof}
논증은 보조정리 \ref{morphisms-lemma-square-and-cube}의 증명과 거의 완전히
같지만, 제대로 작동함을 확실히 하기 위해 모든 내용을 적는다.

\medskip\noindent
$A \subset B$가 정수적임을 상기하자(보조정리
\ref{morphisms-lemma-integral-universally-closed}). 보조정리
\ref{morphisms-lemma-subalgebra-inherits}에 의해 $B$가 $A$ 위에서 하나의
원소로 생성된다고 가정해도 된다. 따라서 $B$는 $A$ 위에서 유한이다
(대수학, 보조정리
\ref{algebra-lemma-characterize-finite-in-terms-of-integral}). 따라서
$B/A$의 지지집합은 $A$-가군으로서 닫혀 있고 공집합이 아니다(대수학,
보조정리 \ref{algebra-lemma-support-closed} 및
\ref{algebra-lemma-support-zero}). $\mathfrak p \subset A$를 이
지지집합의 극소 소아이디얼이라 하자. $A \subset B$를
$A_\mathfrak p \subset B_\mathfrak p$로 대체한 뒤에는(보조정리
\ref{morphisms-lemma-special-elements-and-localization}에 의해 허용된다)
$(A, \mathfrak m)$이 국소환이고, $B$가 $A$ 위에서 유한이며,
$B/A$의 지지집합이 $\{\mathfrak m\}$이고 이를 $A$-가군으로 본다고
가정해도 된다. $B/A$가 유한 가군이므로
$I = \text{Ann}_A(B/A)$는 $\mathfrak m = \sqrt{I}$를 만족한다
(대수학, 보조정리 \ref{algebra-lemma-support-closed}).
$\mathfrak m' \subset B$를 $\mathfrak m$ 위에 놓인 유일한
소아이디얼이라 하자. $\Spec(B) \to \Spec(A)$가 위상동형이므로
$\mathfrak m' = \sqrt{IB}$를 얻는다. $f \in \mathfrak m'$에 대해
$n \geq 1$을 택하여 $f^n \in IB$가 되게 하자. 그러면
$f^{n + 1} \in IB$이기도 하다. $IB \subset A$이며 이는 $I$의 선택에
의한 것이므로
$f^n, f^{n + 1} \in A$라고 결론짓는다. 보조정리
\ref{morphisms-lemma-nth-and-nplusone-implies-square-and-cube}를 사용하면
$\mathfrak m' \not \subset A$일 때 보조정리의 결론이 참임을 얻는다.
$\mathfrak m' \subset A$이면 $\mathfrak m' = \mathfrak m$이다.
$A \not = B$이므로 잉여체의 사상
$\kappa = A/\mathfrak m \to B/\mathfrak m' = \kappa'$는 동형일 수 없다.
보조정리 \ref{morphisms-lemma-universally-injective}에 의해 $\kappa$의 표수가
소수 $p$이고 확대 $\kappa'/\kappa$가 순수 비분리임을 얻는다.
$b \in B$를 그 $\kappa'$에서의 상이 $\kappa$에는 속하지 않지만
$p$제곱은 $\kappa$에 속하는 것으로 택하자. 그러면
$b \not \in A$, $b^p \in A$, $pb \in A$이다
($pb \in \mathfrak m' = \mathfrak m \subset A$이기 때문이다).
이는 원하는 결론이다.
\end{proof}

\begin{proposition}
\label{morphisms-proposition-universal-homeomorphism-equal-residue-fields}
$A \subset B$를 환 확대라 하자. 다음 조건들은 동치이다.
\begin{enumerate}
\item $\Spec(B) \to \Spec(A)$는 잉여체 위의 동형들을 유도하는 보편
위상동형사상이다.
\item 모든 유한 부분집합 $E \subset B$는 다음과 같은 확대에 포함된다.
$$
A[b_1, \ldots, b_n] \subset B
$$
즉 $b_i^2, b_i^3 \in A[b_1, \ldots, b_{i - 1}]$이며, 이는
$i = 1, \ldots, n$에 대해 성립한다.
\end{enumerate}
\end{proposition}

\begin{proof}
(1)을 가정하자. 초한 재귀를 사용하여 각 순서수 $\alpha$에 대해
$A$-부분대수 $B_\alpha \subset B$를 다음과 같이 구성한다.
$B_0 = A$로 둔다. $\alpha$가 극한 순서수이면
$B_\alpha = \colim_{\beta < \alpha} B_\beta$로 둔다.
$\alpha = \beta + 1$이면, $B_\beta = B$인 경우에는
$B_\alpha = B_\beta$로 두고, $B_\beta \not = B$인 경우에는 보조정리
\ref{morphisms-lemma-square-and-cube}를 적용하여
$b_\alpha \in B$, $b_\alpha \not \in B_\beta$이고
$b_\alpha^2, b_\alpha^3 \in B_\beta$인 것을 택한 뒤
$B_\alpha = B_\beta[b_\alpha] \subset B$로 둔다. 명백히
$B = \colim B_\alpha$이다(사실 기수를 살펴보면 $B = B_\alpha$이고
그러한 순서수 $\alpha$가 존재함을 알 수 있다).
$A \to B_\alpha$가 모든 순서수 $\alpha$에 대해 (2)를 만족함을 초한 귀납으로
증명하겠다. $\alpha = 0$일 때는 명백하다. 모든
$\beta < \alpha$에 대해 명제가 성립한다고 가정하고
$E \subset B_\alpha$를 유한 부분집합이라 하자. $\alpha$가 극한
순서수이면 $B_\alpha = \bigcup_{\beta < \alpha} B_\beta$이고,
$E \subset B_\beta$이며 그러한 $\beta < \alpha$가 존재함을 알 수 있으므로 이
경우의 결과가 증명된다. $\alpha = \beta + 1$이면
$B_\alpha = B_\beta[b_\alpha]$이다. 따라서 모든 $e \in E$는 다항식
$e = \sum d_{e, i}b_\alpha^i$로 쓸 수 있으며, 여기서
$d_{e, i} \in B_\beta$이다.
$D \subset B_\beta$를 집합
$D = \{d_{e, i}\} \cup \{b_\alpha^2, b_\alpha^3\}$이라 하자.
귀납가정에 의해 보조정리의 명제에서와 같은
$A$-부분대수 $A[b_1, \ldots, b_n] \subset B_\beta$가 존재하여
$D$를 포함한다. 그러면
$A[b_1, \ldots, b_n, b_\alpha] \subset B_\alpha$는 보조정리의
명제에서와 같은 $A$-부분대수로서 $B_\alpha$에 들어가며 $E$를 포함한다.

\medskip\noindent
(2)를 가정하자. $B = \colim B_\lambda$를 그 유한
$A$-부분대수들의 여극한으로 쓰자. 보조정리
\ref{morphisms-lemma-colimit-inherits}에 의해
$\Spec(B_\lambda) \to \Spec(A)$가 잉여체 위의 동형들을 유도하는
보편 위상동형사상임을 보이면 충분하다. 보편적으로 닫힌 사상들의
합성은 보편적으로 닫혀 있고, 잉여체 위의 동형을 유도하는
사상들에 대해서도 같은 명제가 성립한다. 따라서 $A \subset B$이고
$B$가 하나의 원소 $b$로 생성되며 $b^2, b^3 \in A$이면 (1)이
성립함을 보이면 충분하다. 이러한 확대는 정수적이므로
$\Spec(B) \to \Spec(A)$는 보편적으로 닫혀 있고 전사이다(보조정리
\ref{morphisms-lemma-integral-universally-closed} 및 대수학, 보조정리
\ref{algebra-lemma-integral-overring-surjective}).
$(b^2)^3 = (b^3)^2$라는 등식이 $A$에서 성립함에 유의하라. 임의의
환 사상 $\varphi : A \to K$가 주어지면 체 $K$ 안에
$\lambda \in K$가 존재하여 $\varphi(b^2) = \lambda^2$이고
$\varphi(b^3) = \lambda^3$임을 안다. 실제로 $\lambda = 0$으로 두는
경우는 $\varphi(b^2) = 0$일 때이고, 그렇지 않으면
$\lambda = \varphi(b^3)/\varphi(b^2)$로 둔다. 따라서
$B \otimes_A K$는
$K[x]/(x^2 - \lambda^2, x^3 - \lambda^3)$의 몫이다. 이 환은 잉여체가
$K$인 소아이디얼을 정확히 하나 갖는다. 이는
$\Spec(B) \to \Spec(A)$가 전단사이고 잉여체 위의 동형들을 유도함을
함의한다. 보편 닫힘과 함께 사용하면 (1)이 성립함을 얻는다. 보조정리
\ref{morphisms-lemma-universal-homeomorphism} 및
\ref{morphisms-lemma-universally-injective}를 보라.
\end{proof}

\begin{proposition}
\label{morphisms-proposition-universal-homeomorphism}
$A \subset B$를 환 확대라 하자. 다음 조건들은 동치이다.
\begin{enumerate}
\item $\Spec(B) \to \Spec(A)$는 보편 위상동형사상이다.
\item 모든 유한 부분집합 $E \subset B$는 다음과 같은 확대에 포함된다.
$$
A[b_1, \ldots, b_n] \subset B
$$
그리고 $i = 1, \ldots, n$에 대해 다음 중 하나가 성립한다.
\begin{enumerate}
\item $b_i^2, b_i^3 \in A[b_1, \ldots, b_{i - 1}]$이다.
\item 소수 $p$가 존재하여
$pb_i, b_i^p \in A[b_1, \ldots, b_{i - 1}]$이다.
\end{enumerate}
\end{enumerate}
\end{proposition}

\begin{proof}
증명은 명제
\ref{morphisms-proposition-universal-homeomorphism-equal-residue-fields}의
증명과 정확히 같고, 다음과 같이 바뀐다.
\begin{enumerate}
\item 보조정리 \ref{morphisms-lemma-pth-power-and-multiple}을 사용하되 보조정리
\ref{morphisms-lemma-square-and-cube} 대신 사용한다. 따라서 각 후속 순서수
$\alpha = \beta + 1$에 대해
$b_\alpha^2, b_\alpha^3 \in B_\beta$이거나 소수 $p$가 존재하여
$pb_\alpha, b_\alpha^p \in B_\beta$이다.
\item $\alpha$가 후속 순서수이면 경우에 따라
$D = \{d_{e, i}\} \cup \{b_\alpha^2, b_\alpha^3\}$로 잡거나
$D = \{d_{e, i}\} \cup \{pb_\alpha, b_\alpha^p\}$로 잡는다. 이는
$\alpha$가 어느 경우에 속하는지에 달려 있다.
\item (2) $\Rightarrow$ (1)의 증명에서는 $B$가 $A$ 위에서 하나의
원소 $b$로 생성되고 $pb, b^p \in B$이며, 여기서 $p$는 어떤 소수인
경우도 고려해야 한다. 여기서 $A \subset B$는 예를 들어 대수학, 보조정리
\ref{algebra-lemma-p-ring-map}에 의해 스펙트럼 위의 보편
위상동형사상을 유도한다.
\end{enumerate}
이로써 증명이 끝난다.
\end{proof}

\begin{lemma}
\label{morphisms-lemma-universal-homeo-iso-if-invert-p}
$p$를 소수라 하자. $A \to B$를 동형 $A[1/p] \to B[1/p]$를 유도하는
환 사상이라 하자(예를 들어 $p$가 $A$에서 멱영인 경우이다). 다음
조건들은 동치이다.
\begin{enumerate}
\item $\Spec(B) \to \Spec(A)$는 보편 위상동형사상이다.
\item $A \to B$의 핵은 국소 멱영 아이디얼이고, 모든 $b \in B$에
대해 어떤 $p$-거듭제곱 $q$가 존재하여 $qb$와 $b^q$가
$A \to B$의 상에 속한다.
\end{enumerate}
\end{lemma}

\begin{proof}
(2)가 성립하면 대수학, 보조정리 \ref{algebra-lemma-p-ring-map}에 의해
(1)이 성립한다. (1)을 가정하자. 그러면 대수학, 보조정리
\ref{algebra-lemma-image-dense-generic-points}에 의해 $A \to B$의
핵은 멱영원들로 이루어진다. 따라서 $A$를 $A \to B$의 상으로
대체하여 $A \subset B$라고 가정해도 된다. 대수학, 보조정리
\ref{algebra-lemma-help-with-powers}에 의해 집합
$$
B' = \{b \in B \mid p^nb, b^{p^n} \in A\text{ for some }n \geq 0\}
$$
은 $A$-부분대수이고 $B$에 포함된다(곱에 대해 닫혀 있음은 자명하다).
$B' = B$임을 보여야 한다. 그렇지 않으면 보조정리
\ref{morphisms-lemma-pth-power-and-multiple}에 따라 $b \in B$,
$b \not \in B'$가 존재하여 $b^2, b^3 \in B'$이거나, 소수 $\ell$이
존재하여 $\ell b, b^\ell \in B'$이다. 두 경우가 모두 모순으로
이어짐을 보여 보조정리를 증명하겠다.

\medskip\noindent
$A[1/p] = B[1/p]$이므로 $p$-거듭제곱 $q$를 택하여 $qb \in A$가
되게 할 수 있다.

\medskip\noindent
$b^2, b^3 \in B'$이면 $b^q \in B'$이기도 하다. $B'$의 정의에 의해
$(b^q)^{q'} \in A$가 되게 하는 어떤 $p$-거듭제곱 $q'$가 존재함을 얻는다. 그러면
$qq'b, b^{qq'} \in A$이므로 $b \in B'$이며, 이는 모순이다.

\medskip\noindent
이제 소수 $\ell$이 존재하여 $\ell b, b^\ell \in B'$이라고 가정하자.
$\ell \not = p$이면 $\ell b \in B'$와 $qb \in A \subset B'$에서
$b \in B'$를 얻어 모순이다. 따라서 $\ell = p$이고 $b^p \in B'$이며,
앞과 정확히 같은 방식으로 모순을 얻는다.
\end{proof}

\begin{lemma}
\label{morphisms-lemma-make-universal-homeo}
$A$를 환이라 하자. $x, y \in A$라 하자.
\begin{enumerate}
\item $x^3 = y^2$가 $A$에서 성립하면
$A \to B = A[t]/(t^2 - x, t^3 - y)$는 잉여체 위의 전단사와
스펙트럼 위의 보편 위상동형사상을 유도한다.
\item 소수 $p$가 존재하여 $p^px = y^p$가 $A$에서 성립하면
$A \to B = A[t]/(t^p - x, pt - y)$는 스펙트럼 위의 보편
위상동형사상을 유도한다.
\end{enumerate}
\end{lemma}

\begin{proof}
보조정리 \ref{morphisms-lemma-universal-homeomorphism}의 판정법으로 이를
확인하겠다. 두 경우 모두 환 사상은 정수적이다. 따라서 체 $k$와
환 사상 $\varphi : A \to k$가 주어졌을 때 $k$-대수
$B \otimes_A k$가 유일한 소아이디얼을 가지며, 그 잉여체가 (1)의
경우에는 $k$와 같고 (2)의 경우에는 $k$ 위에서 순수 비분리임을
보이면 충분하다. 보조정리 \ref{morphisms-lemma-universally-injective}를 보라.

\medskip\noindent
(1)의 경우 $\lambda = 0$으로 두는 것은 $\varphi(x) = 0$일 때이고, 그렇지
않으면 $\lambda = \varphi(y)/\varphi(x)$로 둔다. 그러면
$B = k[t]/(t^2 - \lambda^2, t^3 - \lambda^2)$이다. 따라서 결과는
명백하다.

\medskip\noindent
(2)의 경우 $k$의 표수가 $p$이면 $\varphi(y) = 0$을 얻고
$B = k[t]/(t^p - \varphi(x))$이다. 이는 잉여체가 $k$이거나 $k$의
순수 비분리 $p$차 확대인 국소 아르틴 $k$-대수이다. $k$의 표수가
$p$가 아니면 $\lambda = \varphi(y)/p$로 둠으로써
$B = k[t]/(t - \lambda) = k$임을 알고, 역시 결론을 얻는다.
\end{proof}

\begin{lemma}
\label{morphisms-lemma-universal-homeo-limit}
$A \to B$를 환 사상이라 하자.
\begin{enumerate}
\item $A \to B$가 스펙트럼 위의 보편 위상동형사상을 유도하면
$B = \colim B_i$는 유한표현 $A$-대수 $B_i$들의 여과 여극한으로
쓸 수 있으며, 각 $A \to B_i$는 스펙트럼 위의 보편
위상동형사상을 유도한다.
\item $A \to B$가 잉여체 위의 동형들과 스펙트럼 위의 보편
위상동형사상을 유도하면 $B = \colim B_i$는 유한표현 $A$-대수
$B_i$들의 여과 여극한으로 쓸 수 있으며, 각 $A \to B_i$는 잉여체
위의 동형들과 스펙트럼 위의 보편 위상동형사상을 유도한다.
\end{enumerate}
\end{lemma}

\begin{proof}
(1)의 증명. 대수학, 보조정리 \ref{algebra-lemma-when-colimit}의
판정법을 사용하겠다. $A \to C$가 유한표현이고
$\varphi : C \to B$가 $A$-대수 사상이라고 하자. 그 상을
$B' = \varphi(C) \subset B$라 하자. 그러면 보조정리
\ref{morphisms-lemma-subalgebra-inherits}에 의해 $A \to B'$는 스펙트럼 위의
보편 위상동형사상을 유도한다. 대수학, 보조정리
\ref{algebra-lemma-ring-colimit-fp}에 의해
$B' = \colim_{i \in I} B_i$로 쓸 수 있는데, $A \to B_i$는
유한표현이고 전이 사상들은 전사이다. 대수학, 보조정리
\ref{algebra-lemma-characterize-finite-presentation}에 의해 지표
$0 \in I$와 분해 $C \to B_0 \to B'$를 사상 $C \to B'$에 대해 택할 수
있다. 사상 $\Spec(B_i) \to \Spec(A)$가 충분히 큰 $i$에 대해
보편 위상동형사상이라고 주장한다.
이 주장이 (1)의 증명을 끝낸다.

\medskip\noindent
주장의 증명. 보조정리 \ref{morphisms-lemma-reduction-universal-homeomorphism}에
의해 환 사상 $A_{red} \to B'_{red}$는 스펙트럼 위의 보편
위상동형사상을 유도한다. 따라서 대수학, 보조정리
\ref{algebra-lemma-image-dense-generic-points}에 의해
$A_{red} \subset B'_{red}$이다. $A' = \Im(A \to B')$로 두면 전사
$A \to A' \to A_{red}$를 가지며, 이들은 전단사
$\Spec(A_{red}) = \Spec(A') = \Spec(A)$를 유도한다. 따라서
$A' \subset B'$는 스펙트럼 위의 보편 위상동형사상을 유도한다.
명제 \ref{morphisms-proposition-universal-homeomorphism}와 $B'$가 $A'$ 위에서
유한형이라는 사실에 의해, $n$과
$b'_1, \ldots, b'_n \in B'$를 찾아
$B' = A'[b'_1, \ldots, b'_n]$가 되게 할 수 있고, 각
$j = 1, \ldots, n$에 대해 다음 중 하나가 성립한다.
\begin{enumerate}
\item $(b'_j)^2, (b'_j)^3 \in A'[b'_1, \ldots, b'_{j - 1}]$이다.
\item 소수 $p$가 존재하여
$pb'_j, (b'_j)^p \in A'[b'_1, \ldots, b'_{j - 1}]$이다.
\end{enumerate}
$b_1, \ldots, b_n \in B_0$를 $b'_1, \ldots, b'_n$의 올림으로 택하자.
$i \geq 0$에 대해 $b_{j, i}$로 $b_j$의 $B_i$ 안에서의 상을 나타내자.
충분히 큰 $i$에 대하여 각 $j = 1, \ldots, n$마다 다음 중 하나가
성립한다.
\begin{enumerate}
\item $b_{j, i}^2, b_{j, i}^3 \in A_i[b_{1, i}, \ldots, b_{j - 1, i}]$이다.
\item 소수 $p$가 존재하여
$pb_{j, i}, b_{j, i}^p \in A_i[b_{1, i}, \ldots, b_{j - 1, i}]$이다.
\end{enumerate}
여기서 $A_i \subset B_i$는 $A \to B_i$의 상이다. $A \to A_i$는
핵이 국소 멱영 아이디얼인 전사 환 사상임에 유의하라. 필요하다면
$i$를 더 키워서 $B_i$가 $b_1, \ldots, b_n$으로 $A_i$ 위에서
생성된다고, 다시 말해 $B_i = A_i[b_1, \ldots, b_n]$이라고 가정해도
된다. 대수학, 보조정리 \ref{algebra-lemma-p-ring-map} 및
\ref{algebra-lemma-2-3-ring-map}에 의해
$A \to A_i \to A_i[b_1] \to \ldots \to A_i[b_1, \ldots, b_n] = B_i$는
스펙트럼 위의 보편 위상동형사상들을 유도한다고 결론짓는다. 이로써
주장의 증명이 끝난다.

\medskip\noindent
(2)의 증명도 정확히 같다.
\end{proof}

\section{절대 약정규화와 반정규화}
\label{morphisms-section-seminormalization}

\noindent
앞 절에서 증명한 결과들에 동기를 얻어 다음 정의를 제시한다.

\begin{definition}
\label{morphisms-definition-seminormal-ring}
$A$를 환이라 하자.
\begin{enumerate}
\item $A$가 {\it 반정규}라는 것은 모든 $x, y \in A$에 대하여
$x^3 = y^2$이면 유일한
$a \in A$가 존재하여 $x = a^2$이고 $y = a^3$일 때,
성립한다고 정의한다.
\item $A$가 {\it 절대 약정규}라는 것은 (a) $A$가 반정규이고,
(b) 임의의 소수 $p$와
$x, y \in A$에 대하여 $p^px = y^p$이면 유일한 $a \in A$가
존재하여 $x = a^p$이고 $y = pa$일 때 성립한다고 정의한다.
\end{enumerate}
\end{definition}

\noindent
흥미로운 관찰은, \cite{Costa}를 보라, 반정규환의 정의에서
존재성만으로 충분하다는 것이다.\footnote{$A$를 다음 성질을 가지는
환이라 하자. 모든 $x, y \in A$에 대하여 $x^3 = y^2$이면
$a \in A$가 존재하여 $x = a^2$이고 $y = a^3$이다. 그러면
$A$는 축약환이다. 실제로 $x^2 = 0$이면 $x^2 = x^3$이므로
$a$가 존재하여 $x = a^3$이고 $x = a^2$이다. 따라서
$x = a^3 = ax = a^4 = x^2 = 0$이다. 마지막으로
$a_1^2 = a_2^2$이고 $a_1^3 = a_2^3$인 축약환 안의
$a_1, a_2$에 대하여
$(a_1 - a_2)^3 = a_1^3 - 3a_1^2a_2 +3a_1a_2^2 - a_2^3 =
(1 - 3 + 3 - 1)a_1^3 = 0$이므로 $a_1 = a_2$이다.}
$a$의 존재만 가정하면 된다. 절대 약정규 스킴은
\cite[Appendix B]{rydh_descent}에서 정의되었다.

\begin{lemma}
\label{morphisms-lemma-seminormal-local-property}
반정규성 또는 절대 약정규성은 환의 국소 성질이다.
성질, 정의 \ref{properties-definition-property-local}를 보라.
\end{lemma}

\begin{proof}
$A$가 반정규이고 $f \in A$라고 하자. $x', y' \in A_f$가
$(x')^3 = (y')^2$를 만족한다고 하자. $x' = x/f^{2n}$ 및
$y' = y/f^{3n}$으로 쓰되, 여기서 어떤 $n \geq 0$과
$x, y \in A$를 택한다. $x, y$를 각각
$f^{2m}x, f^{3m}y$로, $n$을 $n + m$으로 바꾸면
$x^3 = y^2$가 $A$에서 성립함을 알 수 있다. 그러면 유일한
$a \in A$를 얻어 $x = a^2$이고 $y = a^3$이 된다.
$a' = a/f^n$으로 두면
원하는 대로 $x' = (a')^2$이고 $y' = (a')^3$이다. $a'$의 유일성은
$a$의 유일성에서 따른다. 정확히 같은 방식으로, $A$가 절대
약정규이면 $A_f$도 절대 약정규임을 독자가 보일 수 있다.

\medskip\noindent
$A$를 환이라 하고 $f_1, \ldots, f_n \in A$가 단위 아이디얼을
생성한다고 하자. 각 $A_{f_i}$가 반정규라고 가정하되 이는 모든
$i$에 대해 성립한다고 하자. $x, y \in A$가 $x^3 = y^2$를
만족한다고 하자. 각 $i$에 대해 유일한 $a_i \in A_{f_i}$를 얻어
$x = a_i^2$이고 $y = a_i^3$이 $A_{f_i}$에서 성립한다. 유일성과 첫 문단의 결과
(즉 $A_{f_if_j}$가 반정규라는 결과)에 의해 $a_i$와 $a_j$는
$A_{f_if_j}$의 같은 원소로 사상된다. 대수학, 보조정리
\ref{algebra-lemma-standard-covering}에 의해 유일한 $a \in A$를
얻으며, 이는 $a_i$로 $A_{f_i}$에서 사상되고 모든 $i$에 대해
그렇게 된다.
같은 이유로 $x = a^2$이고 $y = a^3$이다. 이 $a$가 유일함은
명백하다. 따라서 $A$는 반정규이다. $A_{f_i}$가 절대 약정규라고
가정하면, 정확히 같은 논증으로 $A$가 절대 약정규임을 알 수 있다.
\end{proof}

\noindent
이제 반정규 스킴과 절대 약정규 스킴을 정의한다.

\begin{definition}
\label{morphisms-definition-seminormal}
$X$를 스킴이라 하자.
\begin{enumerate}
\item $X$의 모든 $x \in X$가 아핀 열린 근방
$\Spec(R) = U \subset X$를 가지며 환 $R$이 반정규일 때,
이를 {\it 반정규}라고 한다.
\item $X$의 모든 $x \in X$가 아핀 열린 근방
$\Spec(R) = U \subset X$를 가지며 환 $R$이 절대 약정규일 때,
이를
{\it 절대 약정규}라고 한다.
\end{enumerate}
\end{definition}

\noindent
다음은 당연히 필요한 보조정리이다.

\begin{lemma}
\label{morphisms-lemma-locally-seminormal}
$X$를 스킴이라 하자. 다음 조건들은 동치이다.
\begin{enumerate}
\item 스킴 $X$는 반정규이다.
\item 모든 아핀 열린집합 $U \subset X$에 대해 환 $\mathcal{O}_X(U)$는
반정규이다.
\item 아핀 열린 덮개 $X = \bigcup U_i$가 존재하여 각
$\mathcal{O}_X(U_i)$가 반정규이다.
\item 열린 덮개 $X = \bigcup X_j$가 존재하여 각 열린 부분스킴
$X_j$가 반정규이다.
\end{enumerate}
더욱이 $X$가 반정규이면 모든 열린 부분스킴이 반정규이다.
“반정규”를 “절대 약정규”로 바꾸어도 같은 명제들이 성립한다.
\end{lemma}

\begin{proof}
성질, 보조정리 \ref{properties-lemma-locally-P}와
보조정리 \ref{morphisms-lemma-seminormal-local-property}를 결합하라.
\end{proof}

\begin{lemma}
\label{morphisms-lemma-seminormal-reduced}
반정규 스킴 또는 환은 축약이다. 따라서 절대 약정규 스킴 또는
환에 대해서도 같은 명제가 성립한다.
\end{lemma}

\begin{proof}
$A$를 환이라 하자. $a \in A$가 영이 아니지만 $a^2 = 0$이면
$a^2 = 0^2$이고 $a^3 = 0^3$이므로 $A$는 반정규가 아니다.
\end{proof}

\begin{lemma}
\label{morphisms-lemma-seminormalization-ring}
$A$를 환이라 하자.
\begin{enumerate}
\item 스펙트럼 위의 보편 위상동형사상을 유도하는 환 사상
$A \to B$들의 범주는 종대상 $A \to A^{awn}$을 가진다.
\item (1)의 범주에 속하는 $A \to B$가 주어졌을 때, 이에 따른
사상 $B \to A^{awn}$이 동형인 것은 $B$가 절대 약정규인 것과
동치이다.
\item 잉여체 위의 동형들과 스펙트럼 위의 보편 위상동형사상을
유도하는 환 사상 $A \to B$들의 범주는 종대상 $A \to A^{sn}$을
가진다.
\item (3)의 범주에 속하는 $A \to B$가 주어졌을 때, 이에 따른
사상 $B \to A^{sn}$이 동형인 것은 $B$가 반정규인 것과 동치이다.
\end{enumerate}
임의의 환 사상 $\varphi : A \to A'$에 대해
$\varphi^{awn} : A^{awn} \to (A')^{awn}$과
$\varphi^{sn} : A^{sn} \to (A')^{sn}$ 중 $\varphi$와 호환되는
사상들이 유일하게 존재한다.
\end{lemma}

\begin{proof}
(1)과 (2)를 증명하고, (3), (4), 마지막 명제의 증명은 생략한다.
다음 꼴의 $A$-대수들의 범주를 생각하자.
$$
B = A[x_1, \ldots, x_n]/J
$$
여기서 $J$는 유한생성 아이디얼이고 $A \to B$는 스펙트럼 위의 보편
위상동형사상을 정의한다. 이 범주가 유향임을 주장한다
(범주론, 정의 \ref{categories-definition-directed}). 실제로
$$
B = A[x_1, \ldots, x_n]/J
\quad\text{and}\quad
B' = A[x_1, \ldots, x_{n'}]/J'
$$
가 주어지면 다음을 생각할 수 있다.
$$
B'' = A[x_1, \ldots, x_{n + n'}]/J''
$$
여기서 $J''$는 $J$의 원소들과
$f(x_{n + 1}, \ldots, x_{n + n'})$ 꼴의 원소들로 생성되며,
후자에서 $f \in J'$이다.
그러면 $A$-대수 준동형 $B \to B''$ 및 $B' \to B''$을 얻으며,
이들은 동형 $B \otimes_A B' \to B''$을 유도한다. 보조정리
\ref{morphisms-lemma-base-change-universal-homeomorphism}와
\ref{morphisms-lemma-composition-universal-homeomorphism}에 의해
$\Spec(B'') \to \Spec(A)$는 보편 위상동형사상이고, 따라서
$A \to B''$은 우리 범주에 속한다.
마지막으로 우리 범주 안의 $\varphi, \varphi' : B \to B'$가
주어지고 $B$가 위에 표시한 꼴이라고 하자. 몫환 $B''$을 생각하되,
이는 $B'$를 원소 $\varphi(x_i) - \varphi'(x_i)$,
$i = 1, \ldots, n$으로 생성되는 아이디얼로 나눈 것이다.
$\Spec(B') = \Spec(B)$이므로 $\Spec(B'') \to \Spec(B')$는
전단사 닫힌 몰입이고, 따라서 보편 위상동형사상이다. 그러므로
$B''$은 우리 범주에 속하고, $\varphi, \varphi'$는
$B' \to B''$에 의해 같아진다. 이로써 주장의 증명이 끝난다.
다음과 같이 둔다.
$$
A^{awn} = \colim B
$$
여기서 여극한은 방금 설명한 범주 전체에 걸친다.
$A \to A^{awn}$은 보조정리 \ref{morphisms-lemma-colimit-inherits}에 의해
스펙트럼 위의 보편 위상동형사상을 유도함에 유의하라
(여기서 범주가 유향이라는 사실을 사용한다).

\medskip\noindent
스펙트럼 위의 보편 위상동형사상을 유도하는 유한표현 환 사상
$A \to B$가 주어지면, 표준 사상 $B \to A^{awn}$을
$A^{awn}$의 바로 그 구성으로부터 얻는다. (1)에 나오는 모든 $A \to B$는
유한표현이고 (1)에 나오는 $A \to B$들의 여과 여극한이므로
(보조정리 \ref{morphisms-lemma-universal-homeo-limit}),
$A \to A^{awn}$이 범주 (1)의 종대상임을 알 수 있다.

\medskip\noindent
$x, y \in A^{awn}$을 $x^3 = y^2$인 원소들이라 하자. 보조정리
\ref{morphisms-lemma-make-universal-homeo}에 의해
$A^{awn} \to A^{awn}[t]/(t^2 - x, t^3 - y)$는 스펙트럼 위의 보편
위상동형사상을 유도한다. 따라서
$A \to A^{awn}[t]/(t^2 - x, t^3 - y)$는 범주 (1)에 속하고,
유일한 $A$-대수 사상
$A^{awn}[t]/(t^2 - x, t^3 - y) \to A^{awn}$을 얻는다.
$a \in A^{awn}$를 $t$의 상이라 하자. 이는 따라서
$a^2 = x$이고 $a^3 = y$인 $A^{awn}$의 유일한 원소이다.
정확히 같은 방식으로,
소수 $p$와 $x, y \in A^{awn}$가 $p^px = y^p$를 만족하도록
주어지면 $a \in A^{awn}$ 중 $a^p = x$이고 $pq = y$인 원소를
유일하게 얻는다. 따라서 정의에 의해 $A^{awn}$은 절대 약정규이다.

\medskip\noindent
마지막으로 $A \to B$가 범주 (1)에 속하고 $B$가 절대 약정규라고
하자. $A^{awn} \to B^{awn}$은 스펙트럼 위의 보편 위상동형사상을
유도하고 $A^{awn}$은 축약이므로(보조정리
\ref{morphisms-lemma-seminormal-reduced}), 다음을 얻는다.
$A^{awn} \subset B^{awn}$
(대수학, 보조정리 \ref{algebra-lemma-image-dense-generic-points}를
보라). 이 포함이 등식이 아니면 보조정리
\ref{morphisms-lemma-pth-power-and-multiple}에 의해
$b \in B^{awn}$, $b \not \in A^{awn}$인 원소가 존재하여
$b^2, b^3 \in A^{awn}$이거나 $pb, b^p \in A^{awn}$이며,
후자의 경우 어떤 소수 $p$가 존재한다. 그러나 정의
\ref{morphisms-definition-seminormal-ring}의 존재성과 유일성에 의해
$b \in A^{awn}$이어야 하므로, 증명을 끝내는 모순을 얻는다.
\end{proof}

\begin{lemma}
\label{morphisms-lemma-seminormalization}
$X$를 스킴이라 하자.
\begin{enumerate}
\item 보편 위상동형사상 $Y \to X$들의 범주는 시초대상
$X^{awn} \to X$를 가진다.
\item (1)의 범주에 속하는 $Y \to X$가 주어졌을 때, 이에 따른
사상 $X^{awn} \to Y$가 동형인 것은 $Y$가 절대 약정규인 것과
동치이다.
\item 잉여체 위의 동형들을 유도하는 보편 위상동형사상
$Y \to X$들의 범주는 시초대상 $X^{sn} \to X$를 가진다.
\item (3)의 범주에 속하는 $Y \to X$가 주어졌을 때, 이에 따른
사상 $X^{sn} \to Y$가 동형인 것은 $Y$가 반정규인 것과 동치이다.
\end{enumerate}
스킴의 임의의 사상 $h : X' \to X$에 대해
$h^{awn} : (X')^{awn} \to X^{awn}$과
$h^{sn} : (X')^{sn} \to X^{sn}$ 중 $h$와 호환되는 사상들이
유일하게 존재한다.
\end{lemma}

\begin{proof}
(1)과 (2)를 증명하고 (3)과 (4)의 증명은 생략하겠다.
$h : X' \to X$를 스킴의 사상이라 하자. (1)이 $X$와 $X'$에 대해
성립하면 $X' \times_X X^{awn} \to X'$는 보편 위상동형사상이고,
따라서 유일한 사상 $(X')^{awn} \to X' \times_X X^{awn}$을
$X'$ 위에서
$(X')^{awn} \to X'$의 보편 성질에 의해 얻는다. 이를 사영
$X' \times_X X^{awn} \to X^{awn}$과 합성하여 $h^{awn}$을 얻는다.
또한 (2)가 $X$와 $X'$에 대해 성립하고 $h$가 열린 몰입이면,
$X' \times_X X^{awn}$은 절대 약정규이고
(보조정리 \ref{morphisms-lemma-locally-seminormal}), 따라서
$(X')^{awn} \to X' \times_X X^{awn}$이 동형임을 얻는다.

\medskip\noindent
모든 보편 위상동형사상이 아핀임을 상기하라. 보조정리
\ref{morphisms-lemma-homeomorphism-affine}를 보라. 따라서 $X$가 아핀이면
(1)과 (2)는 보조정리 \ref{morphisms-lemma-seminormalization-ring}에서 바로
따른다. $X$를 스킴이라 하고 $\mathcal{B}$를 $X$의 아핀
열린집합들의 집합이라 하자. 각 $U \in \mathcal{B}$에 대해
$U^{awn} \to U$를 얻고, $V \subset U$, $V, U \in \mathcal{B}$에
대해 앞 문단의 논의로부터 표준 동형
$\rho_{V, U} : V^{awn} \to V \times_U U^{awn}$을 얻는다.
따라서 상대적 붙이기(구성, 보조정리
\ref{constructions-lemma-relative-glueing})에 의해 사상
$X^{awn} \to X$를 얻는데, 이는 $U^{awn}$으로 $U$ 위에서
$\rho_{V, U}$와 호환되게 제한된다. 다음으로 $Y \to X$를 보편
위상동형사상이라 하자. 그러면 $U \times_X Y \to U$는
$U \in \mathcal{B}$에 대해 보편 위상동형사상이고, 유일한
사상 $g_U : U^{awn} \to U \times_X Y$를 $U$ 위에서 얻는다. 이 $g_U$들은
사상 $\rho_{V, U}$들과 호환된다. 자세한 내용은 생략한다.
그러므로 유일한 사상 $g : X^{awn} \to Y$가 $X$ 위에서 존재하여
$g_U$와 $U$ 위에서 일치한다. 구성, 주석
\ref{constructions-remark-relative-glueing-functorial}을 보라.
이는 $X$에 대해 (1)을 증명한다. (2)는 아핀 국소적으로 성립하므로
따른다.
\end{proof}

\begin{definition}
\label{morphisms-definition-seminormalization}
$X$를 스킴이라 하자.
\begin{enumerate}
\item 보조정리 \ref{morphisms-lemma-seminormalization}에서 구성한 사상
$X^{sn} \to X$를 $X$의 {\it 반정규화}라고 한다.
\item 보조정리 \ref{morphisms-lemma-seminormalization}에서 구성한 사상
$X^{awn} \to X$를 $X$의 {\it 절대 약정규화}라고 한다.
\end{enumerate}
\end{definition}

\noindent
실제로 반정규화 $X^{sn}$은 $X$의 반정규 스킴이고, 절대 약정규화
$X^{awn}$은 절대 약정규 스킴이다. 더욱이 스킴의 임의의 사상
$h : Y \to X$에 대해 다음과 같은 표준 가환 그림을 얻는다.
$$
\xymatrix{
Y^{awn} \ar[d]^{h^{awn}} \ar[r] &
Y^{sn} \ar[d]^{h^{sn}} \ar[r] &
Y \ar[d]^h \\
X^{awn} \ar[r] &
X^{sn} \ar[r] &
X
}
$$
사상 $h^{sn}$과 $h^{awn}$은 $h$와 호환되는 유일한 사상들이다.

\begin{lemma}
\label{morphisms-lemma-characterize-seminormal}
$X$를 스킴이라 하자. 다음 조건들은 동치이다.
\begin{enumerate}
\item $X$는 반정규이다.
\item $X$는 자신의 반정규화와 같다. 즉 사상
$X^{sn} \to X$는 동형이다.
\item $\pi : Y \to X$가 잉여체 위의 동형들을 유도하는 보편
위상동형사상이고 $Y$가 축약이면, $\pi$는 동형이다.
\end{enumerate}
\end{lemma}

\begin{proof}
(1)과 (2)의 동치는 보조정리 \ref{morphisms-lemma-seminormalization}에서
명백하다. (3)이 성립하면 $X^{sn} \to X$는 동형이므로 (2)가
성립함을 알 수 있다.

\medskip\noindent
(2)가 성립한다고 가정하고, $\pi : Y \to X$를 잉여체 위의 동형들을
유도하는 보편 위상동형사상이라 하되 $Y$가 축약이라고 하자.
보조정리 \ref{morphisms-lemma-seminormalization}에 의해 분해
$X \to Y \to X$가 $\text{id}_X$에 대해 존재한다. 그러면 $X \to Y$는 닫힌
몰입이다(스킴, 보조정리 \ref{schemes-lemma-section-immersion}와,
예를 들어 보조정리 \ref{morphisms-lemma-universally-injective-separated}에
의해 $\pi$가 분리라는 사실을 사용한다). 또한 $X \to Y$는 점들
위의 전단사이므로, $Y$의 축약성에 의해 동형이어야 한다.
이로써 증명이 끝난다.
\end{proof}

\begin{lemma}
\label{morphisms-lemma-characterize-absolutely-weakly-normal}
$X$를 스킴이라 하자. 다음 조건들은 동치이다.
\begin{enumerate}
\item $X$는 절대 약정규이다.
\item $X$는 자신의 절대 약정규화와 같다. 즉 사상
$X^{awn} \to X$는 동형이다.
\item $\pi : Y \to X$가 보편 위상동형사상이고 $Y$가 축약이면,
$\pi$는 동형이다.
\end{enumerate}
\end{lemma}

\begin{proof}
보조정리 \ref{morphisms-lemma-characterize-seminormal}와 정확히 같은 방식으로
증명한다.
\end{proof}

\section{유한 국소 자유 사상}
\label{morphisms-section-finite-locally-free}

\noindent
많은 논문에서 저자들은 실제로 유한 국소 자유 사상을 뜻하면서
유한 평탄 사상이라는 말을 사용한다. 밑이 국소 뇌터이면 두 개념이
같기 때문이다. 그러나 일반적으로는 같지 않다. 연습문제, 연습문제
\ref{exercises-exercise-flat-not-projective}를 보라.

\begin{definition}
\label{morphisms-definition-finite-locally-free}
$f : X \to S$를 스킴의 사상이라 하자. $f$를
{\it 유한 국소 자유}라고 하는 것은 $f$가 아핀이고
$f_*\mathcal{O}_X$가 유한 국소 자유 $\mathcal{O}_S$-가군이면
성립한다고 정의한다. 이 경우 $f$가 {\it 계수} 또는
{\it 차수} $d$를 가진다는 것은 층 $f_*\mathcal{O}_X$가
차수 $d$의 유한 국소 자유일 때 성립한다고 정의한다.
\end{definition}

\noindent
$f : X \to S$가 유한 국소 자유이면 $S$는 열린닫힌 부분스킴
$S_d$들의 서로소 합이고, $f^{-1}(S_d) \to S_d$는 차수 $d$의
유한 국소 자유이다.

\begin{lemma}
\label{morphisms-lemma-finite-flat}
$f : X \to S$를 스킴의 사상이라 하자. 다음 조건들은 동치이다.
\begin{enumerate}
\item $f$는 유한 국소 자유이다.
\item $f$는 유한이고 평탄이며 국소 유한표현이다.
\end{enumerate}
$S$가 국소 뇌터이면 이 조건들은 다음 조건과도 동치이다.
\begin{enumerate}
\item[(3)] $f$는 유한이고 평탄이다.
\end{enumerate}
\end{lemma}

\begin{proof}
$V \subset S$를 아핀 열린집합이라 하자. 세 경우 모두 사상은
아핀이므로 $f^{-1}(V)$도 아핀이다. 따라서
$V = \Spec(R)$이고 $f^{-1}(V) = \Spec(A)$라고 쓸 수 있으며,
여기서 $R$-대수 $A$를 택한다.
(1)을 가정하자. 이는 $S$를 아핀 열린집합 $V = \Spec(R)$들로
덮어서 $A$가 유한 자유 $R$-가군이 되게 할 수 있다는 뜻이다.
그러면 대수학, 보조정리 \ref{algebra-lemma-finite-finite-type}에
의해 $R \to A$는 유한표현이다. 따라서 (2)가 성립한다.
역으로 (2)를 가정하자. 모든 아핀 열린집합 $V = \Spec(R)$에 대해,
그 열린집합이 속한 $S$를 생각하면 환 사상 $R \to A$는 유한이고 유한표현이며,
$A$는 평탄 $R$-가군이다. 대수학, 보조정리
\ref{algebra-lemma-finite-finitely-presented-extension}에 의해
$A$는 유한표현 $R$-가군이다. 따라서 대수학, 보조정리
\ref{algebra-lemma-finite-projective}에 의해 $A$는 유한 국소
자유이다. 그러므로 (1)이 성립한다. 뇌터인 경우는 뇌터환 위의
유한 가군이 유한표현 가군이라는 사실에서 따른다. 대수학, 보조정리
\ref{algebra-lemma-Noetherian-finite-type-is-finite-presentation}를 보라.
\end{proof}

\begin{lemma}
\label{morphisms-lemma-composition-finite-locally-free}
유한 국소 자유 사상들의 합성은 유한 국소 자유이다.
\end{lemma}

\begin{proof}
생략한다.
\end{proof}

\begin{lemma}
\label{morphisms-lemma-base-change-finite-locally-free}
유한 국소 자유 사상의 밑변환은 유한 국소 자유이다.
\end{lemma}

\begin{proof}
생략한다.
\end{proof}

\begin{lemma}
\label{morphisms-lemma-finite-locally-free}
$f : X \to S$를 스킴의 유한 국소 자유 사상이라 하자.
열린닫힌 부분스킴들에 의한 서로소 합 분해
$S = \coprod_{d \geq 0} S_d$가 존재하여, $X_d = f^{-1}(S_d)$로
둘 때 제한 사상 $f|_{X_d}$는 유한 국소 자유 사상
$X_d \to S_d$이고 그 차수는 $d$이다.
\end{lemma}

\begin{proof}
유한 국소 자유 층은 국소적으로 잘 정의된 계수를 가지므로
성립한다. 자세한 내용은 생략한다.
\end{proof}

\begin{lemma}
\label{morphisms-lemma-massage-finite}
$f : Y \to X$를 유한 사상이라 하고 $X$가 아핀이라고 하자.
다음 그림이 존재한다.
$$
\xymatrix{
Z' \ar[rd] &
Y' \ar[l]^i \ar[d] \ar[r] &
Y \ar[d] \\
 & X' \ar[r] & X
}
$$
여기서 다음이 성립한다.
\begin{enumerate}
\item $Y' \to Y$와 $X' \to X$는 전사인 유한 국소 자유 사상이다.
\item $Y' = X' \times_X Y$이다.
\item $i : Y' \to Z'$는 닫힌 몰입이다.
\item $Z' \to X'$는 유한 국소 자유이다.
\item $Z' = \bigcup_{j = 1, \ldots, m} Z'_j$는 닫힌 부분스킴들의
(집합론적) 유한 합이며, 각각은 $X'$로 동형사상된다.
\end{enumerate}
\end{lemma}

\begin{proof}
$X = \Spec(A)$이고 $Y = \Spec(B)$라고 쓰자. 또한 더 많은 대수학,
절 \ref{more-algebra-section-descent-flatness-integral}을 보라.
$x_1, \ldots, x_n \in B$를 $B$의 $A$ 위 생성원들이라 하자.
각 $i$에 대해 $P_i(T) \in A[T]$인 일계수 다항식을 택하여
$P(x_i) = 0$이 $B$에서 성립하게 할 수 있다. 대수학, 보조정리
\ref{algebra-lemma-adjoin-roots}를 $n$번 적용하면, 유한 국소 자유
환 확대 $A \subset A'$가 존재하여 각 $P_i$가 완전히 분해된다.
$$
P_i(T) = \prod\nolimits_{k = 1, \ldots, d_i} (T - \alpha_{ik})
$$
여기서 어떤 $\alpha_{ik} \in A'$들이 존재한다. 다음과 같이 두자.
$$
C = A'[T_1, \ldots, T_n]/(P_1(T_1), \ldots, P_n(T_n))
$$
그리고 $B' = A' \otimes_A B$라 하자. 사상 $C \to B'$,
$T_i \mapsto 1 \otimes x_i$는 $A'$-대수 전사이다.
$X' = \Spec(A')$, $Y' = \Spec(B')$, $Z' = \Spec(C)$로 두면
(1)--(4)가 성립함을 알 수 있다. (5)가 성립하는 이유는 집합론적으로
$\Spec(C)$가 다음 아이디얼들로 정의되는 닫힌 부분스킴들의
합이기 때문이다.
$$
(T_1 - \alpha_{1k_1}, T_2 - \alpha_{2k_2}, \ldots, T_n - \alpha_{nk_n})
$$
여기서 임의의 $1 \leq k_i \leq d_i$를 택한다.
\end{proof}

\noindent
다음 보조정리는 아래의 보조정리
\ref{morphisms-lemma-image-nowhere-dense-quasi-finite}에서 올바른 일반성으로
진술된다.

\begin{lemma}
\label{morphisms-lemma-image-nowhere-dense-finite}
$f : Y \to X$를 스킴의 유한 사상이라 하자.
$T \subset Y$를 $Y$의 닫힌 조밀하지 않은 부분집합이라 하자.
그러면 $f(T) \subset X$는 $X$의 닫힌 조밀하지 않은 부분집합이다.
\end{lemma}

\begin{proof}
보조정리 \ref{morphisms-lemma-finite-proper}에 의해 $f(T) \subset X$가
닫혀 있음을 안다. $X = \bigcup X_i$를 아핀 덮개라 하자.
$T$가 $Y$에서 조밀하지 않으므로
$T \cap f^{-1}(X_i)$도 $f^{-1}(X_i)$에서 조밀하지 않다.
따라서 아핀인 경우에 정리를 증명할 수 있으면
$f(T) \cap X_i$가 조밀하지 않음을 알 수 있다. 그러면 위상수학,
보조정리 \ref{topology-lemma-nowhere-dense-local}에 의해
$T$가 $X$에서 조밀하지 않음을 얻는다.

\medskip\noindent
$X$가 아핀이라고 가정하자. 보조정리
\ref{morphisms-lemma-massage-finite}에 나오는 다음 그림을 택하자.
$$
\xymatrix{
Z' \ar[rd] &
Y' \ar[l]^i \ar[d]^{f'} \ar[r]_a &
Y \ar[d]^f \\
 & X' \ar[r]^b & X
}
$$
사상 $a$, $b$는 유한 국소 자유이므로 열린 사상이다
(보조정리 \ref{morphisms-lemma-finite-flat} 및 \ref{morphisms-lemma-fppf-open}).
따라서 $T' = a^{-1}(T)$는 조밀하지 않다. 위상수학, 보조정리
\ref{topology-lemma-open-inverse-image-closed-nowhere-dense}를 보라.
사상 $b$는 전사이고 열린 사상이다. 따라서
$f'(T') = b^{-1}(f(T))$가 조밀하지 않음을 증명할 수 있다면
$f(T)$도 조밀하지 않다. 위상수학, 보조정리
\ref{topology-lemma-open-inverse-image-closed-nowhere-dense}를 보라.
$i$가 닫힌 몰입이므로 위상수학, 보조정리
\ref{topology-lemma-closed-image-nowhere-dense}에 의해
$i(T') \subset Z'$가 닫혀 있고 조밀하지 않음을 알 수 있다.
따라서 문제를 다음 문단에서 논의하는 경우로 환원하였다.

\medskip\noindent
$Y = \bigcup_{i = 1, \ldots, n} Y_i$가 닫힌 부분집합들의 유한
합이고 각각이 $X$로 동형사상된다고 가정하자.
$T_i = Y_i \cap T$를 생각하자. 각 $T_i$가 $Y_i$에서 조밀하지
않으면, 각 $f(T_i)$는 $X$에서 조밀하지 않다. 이는
$Y_i \to X$가 동형이기 때문이다.
따라서 $f(T) = f(T_i)$는 $X$의 조밀하지 않은
닫힌 부분집합들의 유한 합이고 결론을 얻는다. 위상수학, 보조정리
\ref{topology-lemma-nowhere-dense}를 보라.
그렇지 않다고 하자. 이를테면 $T_1$이 공집합이 아닌 열린집합
$V \subset Y_1$을 포함한다고 하자. 이것이 모순임을 보이겠다.
$Y_2 \cap V \subset V$를 생각하자. 이는 진부분 닫힌집합이거나
$V$와 같다. 첫 번째 경우 $V$를 $V \setminus V \cap Y_2$로
바꾸어서 $V \subset T_1$이 $Y_1$에서 열려 있고 $Y_2$와 만나지
않게 한다. 두 번째 경우 $V \subset Y_1 \cap Y_2$이고 이는
$Y_1$과 $Y_2$ 모두에서 열린다.
$i = 3, \ldots, n$에 대해 차례로 반복하자. 그 결과 서로소 합 분해
$$
\{1, \ldots, n\} = I_1 \amalg I_2, \quad 1 \in I_1
$$
와, 열린집합 $V$를 $Y_1$에서 얻어서 $T_1$에 포함되게 하고,
$V \subset Y_i$가 $i \in I_1$에 대해 성립하며
$V \cap Y_i = \emptyset$이 $i \in I_2$에 대해 성립하게 할 수 있다.
$U = f(V)$로 두자. 이는 $X$의 열린집합인데,
$f|_{Y_1} : Y_1 \to X$가 동형이기 때문이다. 그러면
$$
f^{-1}(U) = V\ \amalg\ \bigcup\nolimits_{i \in I_2} (Y_i \cap f^{-1}(U))
$$
$\bigcup_{i \in I_2} Y_i$가 닫혀 있으므로
$V \subset f^{-1}(U)$가 열려 있고, 따라서 $V \subset Y$가
열려 있음을 얻는다. 이는 $T$가 $Y$에서 조밀하지 않다는
가정에 모순이며, 원하는 결론이다.
\end{proof}

\section{유리 사상}
\label{morphisms-section-rational-maps}

\noindent
$X$를 스킴이라 하자. $U$, $V$가 $X$에서 조밀한 열린집합들이면
$U \cap V$도 그러함에 유의하라.

\begin{definition}
\label{morphisms-definition-rational-map}
$X$, $Y$를 스킴이라 하자.
\begin{enumerate}
\item $f : U \to Y$, $g : V \to Y$를 조밀한 열린 부분집합
$U$, $V$ 위에서 정의된 스킴의 사상들이라 하되 이 부분집합들은
$X$에 속한다고 하자. $f$가 {\it 동치}인 사상을 $g$라고 하는 것은
$f|_W = g|_W$인 어떤 $W \subset U \cap V$가 $X$에서 조밀하게
열릴 때를 말한다.
\item {\it $X$에서 $Y$로 가는 유리 사상}은 (1)에서 정의한
동치관계의 동치류이다.
\item $X$, $Y$가 밑스킴 $S$ 위의 스킴이라고 하자.
$X$에서 $Y$로 가는 유리 사상이 {\it $S$-유리 사상, 즉
$X$에서 $Y$로 가는 밑에 상대적인 유리 사상}이라는 것은 그 동치류가
대표원 $f : U \to Y$를 가지고 이 대표원이 $S$-사상일 때를 말한다.
\end{enumerate}
\end{definition}

\noindent
정의의 (1)에 나오는 두 사상 $f$, $g$가 동치라고 말하는 대신
같은 유리 사상을 정의한다고 말하겠다. 어떤 경우에는 유리 사상이
일반점에서의 국소환 사이의 사상들로 결정된다.

\begin{lemma}
\label{morphisms-lemma-rational-map-finite-presentation}
$S$를 스킴이라 하자. $X$와 $Y$를 $S$ 위의 스킴들이라 하자.
$X$가 일반점 $x_1, \ldots, x_n$을 가지는 유한 개의 기약 성분을
가진다고 가정하자. $s_i \in S$를 $x_i$의 상이라 하자. 다음 사상을
생각하자.
$$
\left\{
\begin{matrix}
S\text{-rational maps} \\
\text{from }X\text{ to }Y
\end{matrix}
\right\}
\longrightarrow
\left\{
\begin{matrix}
(y_1, \varphi_1, \ldots, y_n, \varphi_n)\text{ where }
y_i \in Y\text{ lies over }s_i\text{ and}\\
\varphi_i : \mathcal{O}_{Y, y_i} \to \mathcal{O}_{X, x_i}
\text{ is a local }\mathcal{O}_{S, s_i}\text{-algebra map}
\end{matrix}
\right\}
$$
이 사상은 $f : U \to Y$를 $2n$-쌍으로 보내며, 이 쌍은
$y_i = f(x_i)$이고 $\varphi_i = f^\sharp_{x_i}$이다. 그러면
\begin{enumerate}
\item $Y \to S$가 국소 유한형이면 이 사상은 단사이다.
\item $Y \to S$가 국소 유한표현이면 이 사상은 전단사이다.
\item $Y \to S$가 국소 유한형이고 $X$가 축약이면 이 사상은
전단사이다.
\end{enumerate}
\end{lemma}

\begin{proof}
$X$의 임의의 조밀한 열린집합은 점들 $x_i$를 포함하므로 이 구성은
의미가 있다. (1) 또는 (2)를 증명할 때 $X$를 임의의 조밀한
열린집합으로 바꾸어도 된다. 따라서 $Z_1, \ldots, Z_n$가 $X$의
기약 성분들이면 $X$를
$X \setminus \bigcup_{i \not = j} Z_i \cap Z_j$로 바꾸어도 된다.
그렇게 바꾸면 $X$는 자신의 기약 성분들의 서로소 합이다
(이들을 열린닫힌 부분스킴으로 본다). 그러면 화살표의 우변과 좌변이
모두 기약 성분들에 걸친 곱이므로 $X$가 기약인 경우로 환원된다.

\medskip\noindent
$X$가 기약이고 그 일반점 $x$가 $s \in S$ 위에 놓인다고 하자.
(1)은 보조정리 \ref{morphisms-lemma-morphism-defined-local-ring}의 (1)에서
따르고, (2)와 (3)은 같은 보조정리의 (2)에서 따른다.
\end{proof}

\begin{definition}
\label{morphisms-definition-rational-function}
$X$를 스킴이라 하자. {\it $X$ 위의 유리함수}란 $X$에서
$\mathbf{A}^1_{\mathbf{Z}}$로 가는 유리 사상을 말한다.
\end{definition}

\noindent
아핀 직선 $\mathbf{A}^1$의 정의는 구성, 정의
\ref{constructions-definition-affine-n-space}를 보라.
$X$를 $S$ 위의 스킴이라 하자. 임의의 열린집합 $U \subset X$에
대해 사상 $U \to \mathbf{A}^1_{\mathbf{Z}}$은 사상
$U \to \mathbf{A}^1_S$와 같으며, 후자는 $S$ 위의 사상이다.
따라서 유리함수는 $S$-유리 사상 중 $X$에서
$\mathbf{A}^1_S$로 가는 것과도 같다.

\medskip\noindent
표준적인 식별
$\Mor(T, \mathbf{A}^1_{\mathbf{Z}}) = \Gamma(T, \mathcal{O}_T)$이
임의의 스킴 $T$에 대해 있음을 상기하라. 스킴, 예
\ref{schemes-example-global-sections}를 보라. 따라서
$\mathbf{A}^1_{\mathbf{Z}}$은 스킴의 범주에서 환 대상이다.
더 정확히 말해 사상들
\begin{eqnarray*}
+ : \mathbf{A}^1_{\mathbf{Z}} \times \mathbf{A}^1_{\mathbf{Z}}
& \longrightarrow &
\mathbf{A}^1_{\mathbf{Z}} \\
(f, g) & \longmapsto & f + g \\
* : \mathbf{A}^1_{\mathbf{Z}} \times \mathbf{A}^1_{\mathbf{Z}}
& \longrightarrow &
\mathbf{A}^1_{\mathbf{Z}} \\
(f, g) & \longmapsto & fg
\end{eqnarray*}
은 환의 덧셈과 곱셈의 모든 공리(언제나처럼 $1$을 가지는 가환
환의 공리)를 만족한다. 따라서 $\mathbf{A}^1_{\mathbf{Z}}$로 가는
유리 사상들의 집합도 자연스러운 환 구조를 가진다.

\begin{definition}
\label{morphisms-definition-ring-of-rational-functions}
$X$를 스킴이라 하자. {\it $X$ 위의 유리함수환}은, 원소가
유리함수이고 덧셈과 곱셈은 방금 설명한 것인 환 $R(X)$이다.
\end{definition}

\noindent
기약 성분이 유한 개뿐인 스킴에 대해서는 이를 계산할 수 있다.

\begin{lemma}
\label{morphisms-lemma-integral-scheme-rational-functions}
$X$를 기약 성분 $X_1, \ldots, X_n$이 유한 개인 스킴이라 하자.
$\eta_i \in X_i$가 일반점이면
$$
R(X) = \mathcal{O}_{X, \eta_1} \times \ldots \times \mathcal{O}_{X, \eta_n}
$$
$X$가 축약이면 이는 $\prod \kappa(\eta_i)$와 같다.
$X$가 정수적이면 $R(X) = \mathcal{O}_{X, \eta} = \kappa(\eta)$는
체이다.
\end{lemma}

\begin{proof}
$U \subset X$를 조밀한 열린 부분집합이라 하자. 그러면
$U_i = (U \cap X_i) \setminus (\bigcup_{j \not = i} X_j)$는
$\eta_i$를 포함하므로 공집합이 아닌 열린집합이고, $X_i$에
포함되며, $\bigcup U_i \subset U \subset X$는 조밀하다.
따라서 보조정리의 식별은 다음 등식들의 연쇄에서 나온다.
\begin{align*}
R(X)
& =
\colim_{U \subset X\text{ open dense}} \Mor(U, \mathbf{A}^1_\mathbf{Z}) \\
& =
\colim_{U \subset X\text{ open dense}} \mathcal{O}_X(U) \\
& =
\colim_{\eta_i \in U_i \subset X\text{ open}} \prod \mathcal{O}_X(U_i) \\
& =
\prod \colim_{\eta_i \in U_i \subset X\text{ open}} \mathcal{O}_X(U_i) \\
& =
\prod \mathcal{O}_{X, \eta_i}
\end{align*}
두 번째 등식은 스킴, 예
\ref{schemes-example-global-sections}이다. 마지막 명제는 대수학,
보조정리 \ref{algebra-lemma-minimal-prime-reduced-ring}에서 따른다.
\end{proof}

\begin{definition}
\label{morphisms-definition-function-field}
$X$를 정수적 스킴이라 하자. $X$의 {\it 함수체} 또는
{\it 유리함수체}는 체 $R(X)$이다.
\end{definition}

\noindent
때로는 이 체를 $k(X)$로 나타내어 $R(X)$ 표기를 대신하기도 한다.
함수체 개념을 사용하여 정수적 스킴의 분리 조건을 해명할 수 있다.
보조정리 \ref{morphisms-lemma-integral-scheme-rational-functions}에 의해
정수적 스킴의 모든 국소환 $\mathcal{O}_{X, x}$을 $R(X)$의
국소 부분환으로 볼 수 있음에 유의하라.

\begin{lemma}
\label{morphisms-lemma-distinct-local-rings}
$X$를 정수적 분리 스킴이라 하자. $Z_1$, $Z_2$를 $X$의 서로 다른
기약 닫힌 부분집합들이라 하고, $\eta_i$를 $Z_i$의 일반점이라 하자.
$Z_1 \not\subset Z_2$이면
$\mathcal{O}_{X, \eta_1} \not \subset \mathcal{O}_{X, \eta_2}$이며,
이는 $R(X)$의 부분환으로서의 포함관계이다.
특히 $Z_1 = \{x\}$가 하나의 닫힌점 $x$로 이루어지면,
$x$의 근방에서 정칙이지만 $\mathcal{O}_{X, \eta_{2}}$에는 속하지
않는 함수가 존재한다.
\end{lemma}

\begin{proof}
먼저 $X$가 분리라는 가정 아래, 스킴 사상
$\Spec(\mathcal{O}_{X, \eta_2}) \to X$ 중 $X$ 위에서 그 합성이
$$
\Spec(R(X)) \longrightarrow
\Spec(\mathcal{O}_{X, \eta_2}) \longrightarrow X
$$
이고 표준 사상 $\Spec(R(X)) \to X$와 일치하는 것이 유일함을
관찰하라. 실제로 표준 사상
$can : \Spec(\mathcal{O}_{X, \eta_2}) \to X$가 존재한다. 스킴, 식
(\ref{schemes-equation-canonical-morphism})을 보라.
$a$라는 두 번째 사상이 $X$로 주어지면, 가정에 의해 $a$는
$can$과 $\Spec(\mathcal{O}_{X, \eta_2})$의 일반점에서 일치한다.
$X$가 분리이므로 $\Spec(\mathcal{O}_{X, \eta_2})$에서 이 두 사상이
일치하는 부분집합은 닫혀 있다. 스킴, 보조정리
\ref{schemes-lemma-where-are-they-equal}를 보라. 따라서
$a = can$이고 이는 $\Spec(\mathcal{O}_{X, \eta_2})$ 전체에서
성립한다.

\medskip\noindent
$Z_1 \not \subset Z_2$라고 가정하고, 반대로
$\mathcal{O}_{X, \eta_{1}} \subset \mathcal{O}_{X, \eta_{2}}$가
$R(X)$의 부분환으로서 성립한다고 하자. 그러면 두 번째 사상
$$
\Spec(\mathcal{O}_{X, \eta_{2}}) \longrightarrow
\Spec(\mathcal{O}_{X, \eta_{1}}) \longrightarrow
X.
$$
을 얻는다. 위의 결과에 의해 이 합성은 $can$과 같아야 한다.
이는 $\eta_2$가 $\eta_1$로 특수화됨을 뜻한다(스킴, 보조정리
\ref{schemes-lemma-specialize-points}를 보라). 그러나 이는
$Z_1 \not \subset Z_2$라는 가정에 모순이다.
\end{proof}

\begin{definition}
\label{morphisms-definition-domain-of-definition}
$\varphi$를 두 스킴 $X$와 $Y$ 사이의 유리 사상이라 하자.
$\varphi$가 {\it 점 $x \in X$에서 정의된다}는 것은
대표원 $(U, f)$가 $\varphi$에 대해 존재하고 $x \in U$일 때를 말한다.
$\varphi$의 {\it 정의역}은 $\varphi$가 정의되는 모든 점들의
집합이다.
\end{definition}

\noindent
이 정의에 따르면 일반적으로 $\varphi$가 자신의 정의역 전체에서
정의된 대표원을 가진다는 명제는 참이 아니다.

\begin{lemma}
\label{morphisms-lemma-rational-map-from-reduced-to-separated}
$X$와 $Y$를 스킴이라 하자. $X$가 축약이고 $Y$가 분리라고 가정하자.
$\varphi$를 $X$에서 $Y$로 가는 유리 사상이라 하고, 그 정의역을
$U \subset X$라 하자. 그러면 유일한 사상 $f : U \to Y$가
존재하여 $\varphi$를 나타낸다. $X$와 $Y$가 분리 스킴 $S$ 위의
스킴들이고 $\varphi$가 $S$-유리 사상이면, $f$는 $S$ 위의
사상이다.
\end{lemma}

\begin{proof}
$(V, g)$와 $(V', g')$를 $\varphi$의 대표원들이라 하자.
그러면 $g, g'$는 어떤 조밀한 열린 부분스킴
$W \subset V \cap V'$에서 일치한다. 한편
$E$를 $g|_{V \cap V'}$와 $g'|_{V \cap V'}$의 동등자라 하자.
이는 $V \cap V'$의 닫힌 부분스킴이다(스킴, 보조정리
\ref{schemes-lemma-where-are-they-equal}). 이제 $W \subset E$이므로
집합론적으로 $E = V \cap V'$이다. $V \cap V'$가 축약이므로
스킴론적으로도 $E = V \cap V'$라고 결론짓는다. 즉
$g|_{V \cap V'} = g'|_{V \cap V'}$이다. 따라서 대표원들
$g : V \to Y$를 $\varphi$에 대해 붙여 사상 $f : U \to Y$를 얻을 수 있다.
스킴, 보조정리 \ref{schemes-lemma-glue}를 보라.
마지막 명제의 증명은 생략한다.
\end{proof}

\noindent
일반적으로 유리 사상들을 합성하는 것은 의미가 없다. 첫 번째
유리 사상의 대표원의 상이 두 번째 유리 사상의 정의역과 만나지
않을 수 있기 때문이다. 그러나 스킴들이 기약이라고 가정하고
우세 유리 사상만을 생각하면 유리 사상들을 합성할 수 있다.

\begin{definition}
\label{morphisms-definition-dominant-rational}
$X$와 $Y$를 기약 스킴들이라 하자. $X$에서 $Y$로 가는 유리 사상의
임의의 대표원 $f : U \to Y$가 스킴의 우세 사상이면, 그 유리
사상을 {\it 우세}라고 한다.
\end{definition}

\noindent
보조정리 \ref{morphisms-lemma-dominant-finite-number-irreducible-components}에
의해 이는 다음 조건과 동치이다. 즉 일반점 $\eta \in X$가 일반점
$\xi$로 사상되는데, 후자는 $Y$의 일반점이다. 다시 말해
$f(\eta) = \xi$이고 이는 임의의 대표원 $f : U \to Y$에 대해
성립한다. 우세 유리 사상 $\varphi$가 기약 스킴 $X$와 $Y$
사이에 있고, 임의의 유리 사상 $\psi$가 $Y$에서 $Z$로 가면
둘을 합성할 수 있다.
실제로 대표원 $f : U \to Y$를 택하되 $U \subset X$가 조밀하게
열리게 하고, 대표원 $g : V \to Z$를 택하되 $V \subset Y$가
조밀하게 열리게 하자. 그러면
$W = f^{-1}(V) \subset X$는 공집합이 아닌 열린집합이다
($X$의 일반점을 포함하기 때문이다). 이제 $\psi \circ \varphi$를
$g \circ f|_W : W \to Z$의 동치류로 정의한다. 이것이 잘 정의됨을
확인하는 일은 생략한다.

\medskip\noindent
이 방식으로 대상들이 기약 스킴이고 사상들이 우세 유리 사상인
범주를 얻는다. 밑스킴 $S$가 주어지면, 대상들이 $S$ 위의 기약
스킴이고 사상들이 우세 $S$-유리 사상인 범주도 같은 방식으로
정의할 수 있다.

\begin{definition}
\label{morphisms-definition-birational-schemes}
$X$와 $Y$를 기약 스킴들이라 하자.
\begin{enumerate}
\item $X$와 $Y$가 기약 스킴과 우세 유리 사상의 범주에서
동형이면 $X$와 $Y$를 {\it 쌍유리}라고 한다.
\item $X$와 $Y$가 밑스킴 $S$ 위의 스킴들이라고 가정하자.
$X$와 $Y$를 {\it $S$-쌍유리}라고 하는 것은 $X$와 $Y$가
$S$ 위의 기약 스킴과 우세 $S$-유리 사상의 범주에서 동형일 때를
말한다.
\end{enumerate}
\end{definition}

\noindent
$X$와 $Y$가 쌍유리 기약 스킴들이면, $X$에서 $Z$로 가는
유리 사상들의 집합과 $Y$에서 $Z$로 가는 유리 사상들의 집합
사이에 모든 스킴 $Z$에 대해 전단사가 있다(그리고 $Z$에 대해
함자적이다). “일반적인” 기약 스킴에 대해서는 이것이 가능한
정의 가운데 하나일 뿐이다. 다른 정의는 $X$와 $Y$가 동형인
유리함수환을 가진다고 요구하는 것이다. 다양체에 대해서는 이
조건들이 동치이다. 보조정리
\ref{morphisms-lemma-criterion-birational-finite-presentation}을 보라.

\begin{lemma}
\label{morphisms-lemma-birational-integral}
$X$와 $Y$를 기약 스킴들이라 하자.
\begin{enumerate}
\item 스킴 $X$와 $Y$가 쌍유리인 것은 이들이 동형인 공집합이
아닌 열린집합들을 가지는 것과 동치이다.
\item $X$와 $Y$가 밑스킴 $S$ 위의 스킴들이라고 가정하자.
$X$와 $Y$가 $S$-쌍유리인 것은 공집합이 아닌 열린집합
$U \subset X$와 $V \subset Y$가 존재하여 $S$-동형인 것과 동치이다.
\end{enumerate}
\end{lemma}

\begin{proof}
$X$와 $Y$가 쌍유리라고 가정하자. $f : U \to Y$와
$g : V \to X$가 각각 $X$에서 $Y$로, $Y$에서 $X$로 가는 서로
역인 우세 유리 사상들을 정의한다고 하자. $V$가 아핀이라고
가정해도 된다. $U$를 $f^{-1}(V)$의 아핀 열린집합으로 바꾸어도
된다. $g \circ f$가 우세 유리 사상으로서 항등사상이므로 합성
$U \to V \to X$는 $U$의 조밀한 열린집합에서 항등사상이다.
따라서 $U$를 더 작은 아핀 열린집합으로 바꾸어
$U \to V \to X$가 $U$를 $X$에 넣는 포함사상이 되게 해도 된다.
따라서 $U \to V$는 몰입이다
(스킴, 보조정리 \ref{schemes-lemma-section-immersion}을
$U \to g^{-1}(U) \to U$에 적용하라).
그러나 $U$와 $V$의 역할을 바꾸어 위 논증을 다시 적용하면,
공집합이 아닌 아핀 열린집합 $V' \subset V$가 존재하여 포함사상이
$V' \to U \to V$로 분해됨을 알 수 있다. 그러면
$V' \to U$는 반드시 열린 몰입이다. 실제로
$V' \to f^{-1}(V') \to V'$는 모노사상들이며
(스킴, 보조정리 \ref{schemes-lemma-immersions-monomorphisms}),
그 합성이 항등사상이므로 모두 동형이다. 따라서 $V'$은
$X$와 $Y$ 각각의 열린집합과 동형이다. $S$-유리 사상의 경우에도
정확히 같은 논증이 적용된다.
\end{proof}

\begin{remark}
\label{morphisms-remark-generalize-category}
기약 스킴과 우세 유리 사상의 범주를 다음과 같이 일반화할 수 있다.
스킴 $X$에 대해 $X^0$으로 점 $x \in X$ 중
$\dim(\mathcal{O}_{X, x}) = 0$인 것들의 집합을 나타내자. 다시 말해 $X^0$은
$X$의 기약 성분들의 일반점 집합이다. 이제 다음 범주를 생각한다.
\begin{enumerate}
\item 대상은 모든 준콤팩트 열린집합의 기약 성분이 유한 개인
스킴 $X$이다.
\item $X$에서 $Y$로 가는 사상은 유리 사상 $f : U \to Y$ 중
$X$에서 $Y$로 가고 $f(U^0) = Y^0$를 만족하는 것이다.
\end{enumerate}
$U \subset X$가 스킴의 조밀한 열린집합이면
$U^0 \subset X^0$이 등식일 필요는 없지만, $X$가 우리 범주의
대상이면 등식이다. 따라서 우리 범주의 두 사상이 주어졌을 때
그 합성은 잘 정의되고 다시 우리 범주의 사상이다.
\end{remark}

\begin{remark}
\label{morphisms-remark-pseudo-morphisms}
정의 \ref{morphisms-definition-rational-map}에는 다음 변형이 있다.
사상 $U \to Y$ 중 스킴론적으로 조밀한 열린 부분스킴
$U \subset X$ 위에서 정의된 것만을 생각한다. 보조정리
\ref{morphisms-lemma-intersection-scheme-theoretically-dense}에 의해 이것이
동치관계를 줌을 알 수 있다. 이들의 동치류를
{\it $X$에서 $Y$로 가는 유사 사상}이라고 한다. $X$가 축약이면
두 개념은 일치한다.
\end{remark}

\section{쌍유리 사상}
\label{morphisms-section-birational}

\noindent
다양체의 쌍유리 사상은 목표의 열린 부분집합 위에서 동형이라는
성질을 가진다는 개념에 익숙할 수도 있다. 그러나 일반적으로
쌍유리 사상은 공집합이 아닌 어떤 열린집합 위에서도 동형이 아닐
수 있다. 예 \ref{morphisms-example-birational-not-iso-over-open}을 보라.
다음이 형식적인 정의이다.

\begin{definition}
\label{morphisms-definition-birational}
\begin{reference}
\cite[(2.2.9)]{EGA1}
\end{reference}
$X$, $Y$를 스킴들이라 하자. $X$와 $Y$가 유한 개의 기약 성분을
가진다고 가정하자. 사상 $f : X \to Y$가 {\it 쌍유리}라는 것은
다음 조건들이 성립할 때를 말한다.
\begin{enumerate}
\item $f$는 $X$의 기약 성분들의 일반점 집합과 $Y$의 기약
성분들의 일반점 집합 사이의 전단사를 유도한다.
\item 모든 일반점 $\eta \in X$ 가운데 $X$의 기약 성분의 일반점인 것에 대해 국소환
사상 $\mathcal{O}_{Y, f(\eta)} \to \mathcal{O}_{X, \eta}$는
동형이다.
\end{enumerate}
\end{definition}

\noindent
아래에서 쌍유리 사상의 일반점 위 올들이 한원소 집합임을 보겠다.
더욱이 실제로 마주치는 대부분의 경우, 기약 스킴 $X$와 $Y$
사이에 쌍유리 사상이 존재하면 $X$와 $Y$가 쌍유리 스킴임을
보겠다.

\begin{lemma}
\label{morphisms-lemma-birational-dominant}
$f : X \to Y$를 기약 성분이 유한 개인 스킴들의 사상이라 하자.
$f$가 쌍유리이면 $f$는 우세이다.
\end{lemma}

\begin{proof}
보조정리 \ref{morphisms-lemma-generic-points-in-image-dominant}와 정의에서
따른다.
\end{proof}

\begin{lemma}
\label{morphisms-lemma-birational-generic-fibres}
$f : X \to Y$를 기약 성분이 유한 개인 스킴들의 쌍유리 사상이라
하자. $y \in Y$가 어떤 기약 성분의 일반점이면 밑변환
$X \times_Y \Spec(\mathcal{O}_{Y, y}) \to \Spec(\mathcal{O}_{Y, y})$는
동형이다.
\end{lemma}

\begin{proof}
$Y = \Spec(B)$가 아핀이고 기약이라고 가정해도 된다. 그러면
$X$도 기약이다. 공집합이 아닌 임의의 아핀 열린집합
$U \subset X$에 대해 결과를 증명하면 $X$에 대해서도 성립한다
(짧은 논증은 생략한다). 따라서 $X$도 아핀이라고 가정해도 되며,
$X = \Spec(A)$라고 쓰자. $y \in Y$가 극소 소아이디얼
$\mathfrak q \subset B$에 대응한다고 하자. 가정에 의해 $A$는
유일한 극소 소아이디얼 $\mathfrak p$를 가지고, 이는 $\mathfrak q$
위에 놓이며, $B_\mathfrak q \to A_\mathfrak p$는 동형이다.
따라서 $A_\mathfrak q \to \kappa(\mathfrak p)$는 전사이고,
그러므로 $\mathfrak p A_\mathfrak q$는 극대 아이디얼이다. 한편
$\mathfrak p A_\mathfrak q$는 $A_\mathfrak q$의 유일한 극소
소아이디얼이다. 따라서 $\mathfrak p A_\mathfrak q$는
$A_\mathfrak q$의 유일한 소아이디얼이고
$A_\mathfrak q = A_\mathfrak p$이다.
$A_\mathfrak q = A \otimes_B B_\mathfrak q$이므로 보조정리가
따른다.
\end{proof}

\begin{example}
\label{morphisms-example-birational-not-iso-over-open}
다음은 목표의 어떤 열린집합 위에서도 동형이 아닌 쌍유리 사상의
두 예이다.

\medskip\noindent
첫 번째 예. $k$를 무한체라 하자. $A = k[x]$라 하자. 다음과 같이
두자.
$B = k[x, \{y_{\alpha}\}_{\alpha \in k}]/
((x-\alpha)y_\alpha, y_\alpha y_\beta)$.
포함 $A \subset B$가 존재하고, 수축 $B \to A$가 존재하며 이는
모든 $y_\alpha$를 영으로 보낸다. 사상 $\Spec(A) \to \Spec(B)$와 사상
$\Spec(B) \to \Spec(A)$는 모두 쌍유리이지만 어떤 열린집합
위에서도 동형이 아니다.

\medskip\noindent
두 번째 예. $A$를 정역이라 하자. $S \subset A$를 $0$을 포함하지
않는 곱셈적 부분집합이라 하자. $B = S^{-1}A$로 두면 사상
$f : \Spec(B) \to \Spec(A)$는 쌍유리이다. 열린집합 $U$가
$\Spec(A)$에 존재하여 $f^{-1}(U) \to U$가 동형이면, 어떤 $a \in A$가
존재하여 $S$의 모든 원소가 주국소화 $A_a$에서 가역원이 된다.
$A = \mathbf{Z}$로 잡고 $S$를 홀수들의 집합으로 잡으면 반례를
얻는다.
\end{example}

\begin{lemma}
\label{morphisms-lemma-birational-birational}
$f : X \to Y$를 밑스킴 $S$ 위에서 기약 성분이 유한 개인 스킴들의
쌍유리 사상이라 하자. 다음 조건들 가운데 하나가 성립한다고 하자.
\begin{enumerate}
\item $f$는 국소 유한형이고 $Y$는 축약이다.
\item $f$는 국소 유한표현이다.
\end{enumerate}
조밀한 열린집합 $U \subset X$와 $V \subset Y$가 존재하여
$f(U) \subset V$이고 $f|_U : U \to V$는 동형이다.
특히 $X$와 $Y$가 기약이면 $X$와 $Y$는 $S$-쌍유리이다.
\end{lemma}

\begin{proof}
$X$와 $Y$가 기약인 경우로 즉시 환원할 수 있으며 그 과정은
생략한다. 더 축소한 뒤 $X$와 $Y$가 아핀이라고 가정해도 된다.
$X = \Spec(A)$와 $Y = \Spec(B)$라고 쓰자. 가정에 의해 각각
$A$, $B$는 각각 유일한 극소 소아이디얼 $\mathfrak p$,
$\mathfrak q$를 가지고, 소아이디얼 $\mathfrak p$는
$\mathfrak q$ 위에 놓이며, $B_\mathfrak q = A_\mathfrak p$이다.
보조정리 \ref{morphisms-lemma-birational-generic-fibres}에 의해
$B_\mathfrak q = A_\mathfrak q = A_\mathfrak p$이다.

\medskip\noindent
$B \to A$가 유한형이라고 하자. 이를테면
$A = B[x_1, \ldots, x_n]$라고 하자. 각 지표에 대해
$b_i \in B$와 $g_i \in B \setminus \mathfrak q$가 존재하여
$b_i/g_i$가 $x_i$의 $A_\mathfrak q$에서의 상으로 사상된다.
따라서 $b_i  - g_ix_i$는 $A_{g_i'}$에서 영으로 사상되며, 여기서
$g_i' \in B \setminus \mathfrak q$이다.
$g = \prod g_i g'_i$로 두면 $B_g \to A_g$가 전사임을 알 수 있다.
더욱이 $Y$가 축약이면 사상 $B_g \to B_\mathfrak q$가 단사이고,
따라서 $B_g \to A_g$도 단사이다. 이로써 (1)을 증명하였다.

\medskip\noindent
(2)의 증명. 앞 문단의 논증에 의해 $B \to A$가 전사라고 가정해도
된다. $f$가 국소 유한표현이므로 핵 $J \subset B$는 유한생성
아이디얼이다. $J = (b_1, \ldots, b_r)$라고 하자.
$B_\mathfrak q = A_\mathfrak q$이므로
$g_i \in B \setminus \mathfrak q$가 존재하여
$g_i b_i = 0$이 되게 할 수 있다. $g = \prod g_i$로 두면
$B_g \to A_g$가 동형임을 알 수 있다.
\end{proof}

\begin{lemma}
\label{morphisms-lemma-criterion-birational-finite-presentation}
$S$를 스킴이라 하자. $X$와 $Y$를 $S$ 위에서 국소 유한표현인
기약 스킴들이라 하자. $x \in X$와 $y \in Y$를 일반점들이라 하자.
다음 조건들은 동치이다.
\begin{enumerate}
\item $X$와 $Y$는 $S$-쌍유리이다.
\item $X$와 $Y$의 공집합이 아닌 열린집합들 가운데
$S$-동형인 것들이 존재한다.
\item $x$와 $y$는 같은 점 $s$로 사상되고, 이 점은 $S$에 속하며,
$\mathcal{O}_{X, x}$와 $\mathcal{O}_{Y, y}$는
$\mathcal{O}_{S, s}$-대수로서 동형이다.
\end{enumerate}
\end{lemma}

\begin{proof}
(1)과 (2)의 동치는 보조정리
\ref{morphisms-lemma-birational-integral}에서 보았다. (2)가 (3)을 함의함은
즉시 알 수 있다. 이제 (3)을 가정하고 (1)을 증명하자.
보조정리 \ref{morphisms-lemma-rational-map-finite-presentation}에 의해 유리 사상
$f : U \to Y$가 존재하여 $x \in U$를 $y$로 보내고 주어진 동형
$\mathcal{O}_{Y, y} \cong \mathcal{O}_{X, x}$를 유도한다. 따라서
$f$는 쌍유리 사상이고, 보조정리
\ref{morphisms-lemma-birational-birational}에 의해 공집합이 아닌 열린집합들
위의 동형을 유도한다. 이로써 증명이 끝난다.
\end{proof}

\begin{lemma}
\label{morphisms-lemma-common-open}
$S$를 스킴이라 하자. $X$와 $Y$를 $S$ 위에서 국소 유한형인 정수적
스킴들이라 하자. $x \in X$와 $y \in Y$를 일반점들이라 하자.
다음 조건들은 동치이다.
\begin{enumerate}
\item $X$와 $Y$는 $S$-쌍유리이다.
\item $X$와 $Y$의 공집합이 아닌 열린집합들 가운데 $S$-동형인
것들이 존재한다.
\item $x$와 $y$는 같은 점 $s \in S$로 사상되고
$\kappa(x) \cong \kappa(y)$가 $\kappa(s)$-확대로서 성립한다.
\end{enumerate}
\end{lemma}

\begin{proof}
(1)과 (2)의 동치는 보조정리
\ref{morphisms-lemma-birational-integral}에서 보았다. (2)가 (3)을 함의함은
즉시 알 수 있다. 이제 (3)을 가정하고 (1)을 증명하자.
$\mathcal{O}_{X, x} = \kappa(x)$이고
$\mathcal{O}_{Y, y} = \kappa(y)$임에 유의하라. 이는 대수학,
보조정리 \ref{algebra-lemma-minimal-prime-reduced-ring}에서 따른다.
보조정리 \ref{morphisms-lemma-rational-map-finite-presentation}에 의해 유리 사상
$f : U \to Y$가 존재하여 $x \in U$를 $y$로 보내고 주어진 동형
$\mathcal{O}_{Y, y} \cong \mathcal{O}_{X, x}$를 유도한다. 따라서
$f$는 쌍유리 사상이고, 보조정리
\ref{morphisms-lemma-birational-birational}에 의해 공집합이 아닌 열린집합들
위의 동형을 유도한다. 이로써 증명이 끝난다.
\end{proof}

\section{일반적으로 유한인 사상}
\label{morphisms-section-generically-finite}

\noindent
이 절에서는 국소적으로 유한형이고 어떤 의미에서 ``일반적으로 유한''인
스킴 사이의 사상들을 특징짓는다.

\begin{lemma}
\label{morphisms-lemma-generically-finite}
$X$, $Y$를 스킴들이라 하자.
$f : X \to Y$가 국소적으로 유한형이라고 하자.
$\eta \in Y$를 $Y$의 한 기약 성분의 일반점이라 하자.
다음 조건들은 동치이다.
\begin{enumerate}
\item 집합 $f^{-1}(\{\eta\})$는 유한이다.
\item 아핀 열린집합 $U_i \subset X$, $i = 1, \ldots, n$과
$V \subset Y$가 존재하여 $f(U_i) \subset V$,
$\eta \in V$, $f^{-1}(\{\eta\}) \subset \bigcup U_i$이고,
각 $f|_{U_i} : U_i \to V$는 유한이다.
\end{enumerate}
$f$가 준분리이면 이 조건들은 다음 조건과도 동치이다.
\begin{enumerate}
\item[(3)] 아핀 열린집합 $V \subset Y$와 $U \subset X$가 존재하여
$f(U) \subset V$, $\eta \in V$, $f^{-1}(\{\eta\}) \subset U$이고,
$f|_U : U \to V$는 유한이다.
\end{enumerate}
$f$가 준콤팩트이고 준분리이면 이 조건들은 다음 조건과도 동치이다.
\begin{enumerate}
\item[(4)] 아핀 열린집합 $V \subset Y$, $\eta \in V$가 존재하여
$f^{-1}(V) \to V$가 유한이다.
\end{enumerate}
\end{lemma}

\begin{proof}
문제는 밑에서 국소적이다. 따라서 $Y$를 $\eta$의 아핀 근방으로
대체할 수 있으며, 아래의 증명 전체에서 $Y$가 아핀이라고, 이를테면
$Y = \Spec(R)$이라고 가정한다.

\medskip\noindent
(2)가 (1)을 함의함은 명백하다.
$f^{-1}(\{\eta\}) = \{\xi_1, \ldots, \xi_n\}$가 유한이라고 가정하자.
아핀 열린집합 $U_i \subset X$를 택하되 $\xi_i \in U_i$가 되게 하자.
대수학, 보조정리 \ref{algebra-lemma-generically-finite}에 의해 $Y$를
$Y$의 표준 열린집합으로 대체하면 각 사상 $U_i \to Y$가 유한임을
알 수 있다. 달리 말하면 (2)가 성립한다.

\medskip\noindent
(3)이 (1)을 함의함은 명백하다. $f$가 준분리이고 (1)이 성립한다고
가정하자. $f^{-1}(\{\eta\}) = \{\xi_1, \ldots, \xi_n\}$라고 쓰자.
보조정리 \ref{morphisms-lemma-finite-fibre}에 의해 $\xi_i$들 사이에는 특수화가
없다. 각 $\xi_i$는 일반점 $\eta$로 사상되고 이 점은 $Y$의 한 기약
성분의 일반점이므로, 자명하지 않은 특수화 $\xi \leadsto \xi_i$가
$X$ 안에 있을 수 없다
($\xi$도 $\eta$로 사상될 것이기 때문이다). 따라서 각 $\xi_i$는
$X$의 한 기약 성분의 일반점이다. $Y$가 아핀이고 $f$가 준분리이므로
$X$가 준분리임을 알 수 있다. 성질, 보조정리
\ref{properties-lemma-maximal-points-affine}에 의해 아핀 열린집합
$U \subset X$를 찾을 수 있고, 이는 모든 $\xi_i$를 포함한다. 대수학,
보조정리 \ref{algebra-lemma-generically-finite}에 의해 $Y$를 $Y$의
표준 열린집합으로 대체하면 사상 $U \to Y$가 유한임을 알 수 있다.
달리 말하면 (3)이 성립한다.

\medskip\noindent
$f$에 아무 추가 가정도 두지 않아도 (4)는 (1)--(3)을 모두
함의함이 명백하다. $f$가 준콤팩트이고 준분리라고 가정하자. 동치인
조건 (1)--(3)이 (4)를 함의함을 보이면 된다. $U$, $V$를 (3)에서와
같이 잡고 $Y$를 $V$로 대체하자. $f$가 준콤팩트이고 $Y$가
준콤팩트이므로(아핀이므로) $X$도 준콤팩트이다. 따라서
$Z = X \setminus U$는 준콤팩트이고, 따라서 사상
$f|_Z : Z \to Y$도 준콤팩트이다. $Z$의 구성에 의해
$\eta \not \in f(Z)$이다. 그러므로 보조정리
\ref{morphisms-lemma-quasi-compact-generic-point-not-in-image}에 의해 아핀 열린 근방
$V'$가 존재하고, 이는 $\eta$의 $Y$ 안에서의 근방이며,
$f^{-1}(V') \cap Z = \emptyset$가 된다. 이때
$f^{-1}(V') \subset U$이고, 이는 $f^{-1}(V') \to V'$가 유한이라는
뜻이다.
\end{proof}

\begin{example}
\label{morphisms-example-counter-generically-finite}
$A = \prod_{n \in \mathbf{N}} \mathbf{F}_2$라 하자.
$A$의 모든 원소는 멱등원이다. 따라서 모든 소아이디얼은 잉여체가
$\mathbf{F}_2$인 극대 아이디얼이다. 그러므로
$X = \Spec(A)$의 위상은 완전 비연결이고 준콤팩트이다. 사영 사상
$A \to \mathbf{F}_2$는 $\Spec(A)$의 열린점들을 정의한다. $X$가
준콤팩트이므로 $X$의 모든 점이 열릴 수는 없다. 열린점이 아닌 닫힌점
$x \in X$를 택하자. 그러면 스킴 $Y$를 만들 수 있는데, 이는 $X$의
두 복사본을 $X \setminus \{x\}$를 따라 붙인 것이다. 달리 말하면
이는 $X$에서 $x$를 두 겹으로 만든 것이다. 스킴, 예
\ref{schemes-example-affine-space-zero-doubled}와 비교하라. 사상
$f : Y \to X$는 준콤팩트이고 유한형이며 유한 올을 가지지만
준분리는 아니다. 점 $x \in X$는 $X$의 한 기약 성분의 일반점이다
($X$가 완전 비연결이기 때문이다). 그러나 보조정리
\ref{morphisms-lemma-generically-finite}의 성질 (3)과 (4)는 성립하지 않는다.
그 이유는 임의의 열린 근방 $x \in U \subset X$에 대해 역상
$f^{-1}(U)$가 아핀이 아니기 때문이다. 실제로 $f^{-1}(U)$ 위의
함수들은 $x$ 위에 놓이는 두 점을 분리할 수 없다(증명은 생략한다.
이는 좋은 연습문제이다). 따라서 보조정리의 (3)과 (4)에서 $f$가
준분리라는 조건이 필요하다.
\end{example}

\begin{remark}
\label{morphisms-remark-quasi-finite-finite-over-dense-open}
보조정리 \ref{morphisms-lemma-generically-finite}의 대안으로, 준유한 사상은
목표의 한 조밀한 열린집합 위에서 유한이라는 명제를 사용할 수 있다.
이는 사상에 관한 추가 내용, 보조정리
\ref{more-morphisms-lemma-quasi-finite-finite-over-dense-open}에서
증명한다.
\end{remark}

\begin{lemma}
\label{morphisms-lemma-quasi-finiteness-over-generic-point}
$X$, $Y$를 스킴들이라 하고 $f : X \to Y$가 국소적으로 유한형이라고
하자. $X^0$, 각각 $Y^0$가 $X$, 각각 $Y$의 기약 성분들의 일반점
집합을 나타낸다고 하자. $\eta \in Y^0$라 하자. 다음 조건들은
동치이다.
\begin{enumerate}
\item $f^{-1}(\{\eta\}) \subset X^0$이다.
\item $f$는 $\eta$ 위에 놓이는 모든 점에서 준유한이다.
\item $f$는 모든 $\xi \in X^0$ 가운데 $\eta$ 위에 놓이는 점에서
준유한이다.
\end{enumerate}
\end{lemma}

\begin{proof}
조건 (1)이 성립하면 올 $X_\eta$의 점들 사이에는 특수화가 없다.
따라서 보조정리 \ref{morphisms-lemma-quasi-finite-at-point-characterize}에 의해
(2)가 성립한다. (2) $\Rightarrow$ (3)은 즉시 알 수 있다. $\eta$가
$Y$의 일반점이므로 $X_\eta$의 일반점들은 $X$의 일반점들이다.
그러므로 (3)과 보조정리
\ref{morphisms-lemma-quasi-finite-at-point-characterize}에 의해 $X_\eta$의
일반점들은 닫힌점이기도 하다. 따라서 $X_\eta$의 모든 점이
일반점이고, (1)이 성립함을 알 수 있다.
\end{proof}

\begin{lemma}
\label{morphisms-lemma-finite-over-dense-open}
$X$, $Y$를 스킴들이라 하고 $f : X \to Y$가 국소적으로 유한형이라고
하자. $X^0$, 각각 $Y^0$가 $X$, 각각 $Y$의 기약 성분들의 일반점
집합을 나타낸다고 하자. 다음을 가정하자.
\begin{enumerate}
\item $X^0$와 $Y^0$는 유한이고 $f^{-1}(Y^0) = X^0$이다.
\item $f$가 준콤팩트이거나 $f$가 분리이다.
\end{enumerate}
그러면 조밀한 열린집합 $V \subset Y$가 존재하여
$f^{-1}(V) \to V$가 유한이다.
\end{lemma}

\begin{proof}
$Y$의 기약 성분이 유한 개이므로, 그 기약 성분들의 분리합인 조밀한
열린집합을 찾을 수 있다. 따라서 $Y$가 일반점 $\eta$를 갖는 기약
아핀 스킴이라고 가정해도 된다. $\eta$ 위의 올은 $X^0$가 유한이므로
유한이다.

\medskip\noindent
$f$가 분리이고 $Y$가 기약 아핀이라고 가정하자. 보조정리
\ref{morphisms-lemma-generically-finite}의 (3)에서와 같이 $V \subset Y$와
$U \subset X$를 택하자. $f|_U : U \to V$가 유한이므로
$U \subset f^{-1}(V)$는 열리고 닫혀 있기도 하다(보조정리
\ref{morphisms-lemma-image-proper-scheme-closed}와 \ref{morphisms-lemma-finite-proper}).
따라서 $f^{-1}(V) = U \amalg W$이며, 여기서 $W$는 $X$의 열린
부분스킴이다. 그러나 $U$가 $X$의 모든 일반점을
포함하므로 원하는 대로 $W = \emptyset$이다.

\medskip\noindent
$f$가 준콤팩트이고 $Y$가 기약 아핀이라고 가정하자. 그러면 $X$는
준콤팩트이므로, 분리인 조밀한 열린 부분스킴 $U \subset X$가
존재한다(성질, 보조정리
\ref{properties-lemma-quasi-compact-dense-open-separated}). 일반점 집합
$X^0$가 유한이므로 $X^0 \subset U$이다. 따라서
$\eta \not \in f(X \setminus U)$이다. $X \setminus U \to Y$가
준콤팩트이므로, 공집합이 아닌 열린집합 $V \subset Y$가 존재하여
$f^{-1}(V) \subset U$가 됨을 알 수 있다. 보조정리
\ref{morphisms-lemma-quasi-compact-dominant}를 보라. $X$를 $f^{-1}(V)$로,
$Y$를 $V$로 대체하면 앞 문단에서 다룬 분리인 경우로 환원된다.
\end{proof}

\begin{lemma}
\label{morphisms-lemma-birational-isomorphism-over-dense-open}
$X$, $Y$를 스킴들이라 하자. $f : X \to Y$를 기약 성분이 유한 개인
스킴 사이의 쌍유리 사상이라 하자. 다음을 가정하자.
\begin{enumerate}
\item $f$가 준콤팩트이거나 $f$가 분리이다.
\item $f$가 국소적으로 유한형이고 $Y$가 축약이거나, $f$가
국소적으로 유한표현이다.
\end{enumerate}
그러면 조밀한 열린집합 $V \subset Y$가 존재하여
$f^{-1}(V) \to V$가 동형이다.
\end{lemma}

\begin{proof}
보조정리 \ref{morphisms-lemma-finite-over-dense-open}에 의해 $f$가 유한이라고
가정해도 된다. $Y$의 기약 성분이 유한 개이므로, 그 기약 성분들의
분리합인 조밀한 열린집합을 찾을 수 있다. 따라서 $Y$가 기약이라고
가정해도 된다. 보조정리 \ref{morphisms-lemma-birational-birational}에 의해
공집합이 아닌 열린집합 $U \subset X$가 존재하여
$f|_U : U \to Y$가 열린몰입이 되게 할 수 있다. $f$가 유한이므로
닫힌 부분집합 $f(X \setminus U)$를 $Y$에서 제거하면 $f$가
동형임을 알 수 있다.
\end{proof}

\begin{lemma}
\label{morphisms-lemma-finite-degree}
$X$, $Y$를 정수적 스킴들이라 하자.
$f : X \to Y$가 국소적으로 유한형이라고 하자.
$f$가 우세라고 가정하자. 다음 조건들은 동치이다.
\begin{enumerate}
\item 확장 $R(Y) \subset R(X)$의 초월차수는 $0$이다.
\item 확장 $R(Y) \subset R(X)$는 유한이다.
\item 공집합이 아닌 아핀 열린집합 $U \subset X$와 $V \subset Y$가
존재하여 $f(U) \subset V$이고 $f|_U : U \to V$는 유한이다.
\item $X$의 일반점은 $X$의 점들 가운데 $Y$의 일반점으로 사상되는
유일한 점이다.
\end{enumerate}
$f$가 분리이거나 $f$가 준콤팩트이면 이 조건들은 다음 조건과도
동치이다.
\begin{enumerate}
\item[(5)] 공집합이 아닌 아핀 열린집합 $V \subset Y$가 존재하여
$f^{-1}(V) \to V$가 유한이다.
\end{enumerate}
\end{lemma}

\begin{proof}
임의의 아핀 열린집합 $\Spec(A) = U \subset X$와
$\Spec(R) = V \subset Y$를 택하되 $f(U) \subset V$가 되게 하자.
정의에 의해 $R$과 $A$는 정역이다. 환 사상 $R \to A$는 유한형이다
(보조정리 \ref{morphisms-lemma-locally-finite-type-characterize}). 보조정리
\ref{morphisms-lemma-dominant-finite-number-irreducible-components}에 의해 $X$의
일반점은 $Y$의 일반점으로 사상되므로 $R \to A$는 단사이다.
$K = R(Y)$를 $R$의 분수체라 하고 $L = R(X)$를 $A$의 분수체라 하자.
그러면 $L/K$는 유한생성 체확장이다. 따라서 (1)과 (2)는 동치이다.

\medskip\noindent
(2)가 성립한다고 하자. $x_1, \ldots, x_n \in A$를 $A$의 $R$ 위
생성원으로 택하자. 가정에 의해 영이 아닌 다항식
$P_i(X) \in R[X]$가 존재하여 $P_i(x_i) = 0$이다.
$f_i \in R$를 $P_i$의 최고차항 계수라 하자. 그러면
$R_{f_1 \ldots f_n} \to A_{f_1 \ldots f_n}$가 유한임을 알 수 있다.
즉 (3)이 성립한다. (3)이 (2)를 함의함에 유의하라. 따라서 이제
(1), (2), (3)이 모두 동치임을 알 수 있다.

\medskip\noindent
$\eta$를 $X$의 일반점이라 하고 $\eta' \in Y$를 $Y$의 일반점이라
하자. (4)를 가정하자. 그러면 $\dim_\eta(X_{\eta'}) = 0$이고,
보조정리 \ref{morphisms-lemma-dimension-fibre-at-a-point}에 의해
$R(X) = \kappa(\eta)$의 초월차수가 $0$임을 알 수 있다. 여기서
이는 $R(Y) = \kappa(\eta')$ 위의 초월차수이다. 달리 말하면 (1)이
성립한다. 이제 동치인 조건
(1), (2), (3)이 성립한다고 하자. $x \in X$가 $\eta'$로 사상되는
점이라고 하자. $x$가 $\eta$의 특수화이므로 포함관계
$R(Y) \subset \mathcal{O}_{X, x} \subset R(X)$를 얻는다. 따라서
$\mathcal{O}_{X, x}$는 체이다. 대수학, 보조정리
\ref{algebra-lemma-integral-over-field}를 보라. 그러므로
$x = \eta$이다. 따라서 (1)--(4)가 모두 동치임을 알 수 있다.

\medskip\noindent
$f$에 아무 추가 가정이 없어도 (5)가 (3)을 함의함은 명백하다.
남은 것은 $f$가 분리이거나 준콤팩트이면 동치인 조건 (1)--(4)가
(5)를 함의함을 증명하는 일이다. 이는 보조정리
\ref{morphisms-lemma-finite-over-dense-open}에서 따른다.
\end{proof}

\begin{definition}
\label{morphisms-definition-degree}
$X$와 $Y$를 정수적 스킴들이라 하자.
$f : X \to Y$가 국소적으로 유한형이고 우세라고 하자.
$[R(X) : R(Y)] < \infty$이거나, 보조정리
\ref{morphisms-lemma-finite-degree}의 동치인 조건 (1)--(4) 가운데 어느 하나가
성립한다고 가정하자. 그러면 양의 정수
$$
\deg(X/Y) = [R(X) : R(Y)]
$$
를 {\it $X$의 $Y$ 위 차수}라고 한다.
\end{definition}

\noindent
이 개념은 다음을 만족하는 사상 $f : X \to Y$로 확장할 수 있다.
(a) $Y$가 일반점 $\eta$를 갖는 정수적 스킴이고, (b) $f$가
국소적으로 유한형이며, (c) $f^{-1}(\{\eta\})$가 유한이다. 이 경우
다음과 같이 정의할 수 있다.
$$
\deg(X/Y)
=
\sum\nolimits_{\xi \in X, \ f(\xi) = \eta}
\dim_{R(Y)} (\mathcal{O}_{X, \xi}).
$$
실제로 $R(Y) = \kappa(\eta) = \mathcal{O}_{Y, \eta}$이고
(보조정리 \ref{morphisms-lemma-integral-scheme-rational-functions}), 위의 차원들은
앞의 보조정리 \ref{morphisms-lemma-generically-finite}에 의해 유한이다. 그러나
대부분의 응용에는 위에서 제시한 정의가 적절하다.

\begin{lemma}
\label{morphisms-lemma-degree-composition}
$X$, $Y$, $Z$를 정수적 스킴들이라 하자.
$f : X \to Y$와 $g : Y \to Z$를 국소적으로 유한형인 우세 사상들이라
하자. $[R(X) : R(Y)] < \infty$이고 $[R(Y) : R(Z)] < \infty$라고
가정하자. 그러면
$$
\deg(X/Z) = \deg(X/Y) \deg(Y/Z).
$$
\end{lemma}

\begin{proof}
이는 체의 유한 확장들의 탑에서 차수가 곱셈적이라는 사실에서 따른다.
체론, 보조정리 \ref{fields-lemma-multiplicativity-degrees}를 보라.
\end{proof}

\begin{remark}
\label{morphisms-remark-definition-generically-finite}
$f : X \to Y$를 국소적으로 유한형인 스킴들의 사상이라 하자.
{\it 일반적으로 유한}인 사상을 정의하는 데 사용할 수 있는 성질은
(적어도) 두 가지이다. 어느 성질을 택하는지는 그 성질이 원천에서
국소적이기를 원하는지, 목표에서 국소적이기를 원하는지에 대응한다.
\begin{enumerate}
\item (목표에서 국소적이며, Ravi Vakil이 제안하였다.) $Y$의 모든
준콤팩트 열린집합이 유한 개의 기약 성분을 가진다고 가정하자(예를 들어
$Y$가 국소 뇌터이면 그렇다). 각 일반점의 역상이 유한일 것을 요구한다.
보조정리 \ref{morphisms-lemma-generically-finite}를 보라.
\item (원천에서 국소적이다.) 조밀한 열린집합 $U \subset X$가
존재하여 $U \to Y$가 국소적으로 준유한일 것을 요구한다.
\end{enumerate}
(1)의 경우에는 $Y$에 관한 보조 조건 없이도 요구사항을 공식화할 수
있지만, 이는 일반 스킴에 적절한 개념을 주지 않을 가능성이 크다.
위와 같이 공식화한 성질 (2)는 일반점 위의 올이 유한임을 함의하지
않는다. 그러나 $f$가 준콤팩트이고 $Y$가 (1)에서와 같으면 함의한다.
\end{remark}

\begin{definition}
\label{morphisms-definition-modification}
$X$를 정수적 스킴이라 하자. {\it $X$의 수정}은 쌍유리 고유 사상
$f : X' \to X$ 가운데 $X'$가 정수적인 것이다.
\end{definition}

\noindent
$f : X' \to X$를 정의에서와 같은 수정이라 하자. 보조정리
\ref{morphisms-lemma-finite-degree}에 의해 공집합이 아닌 $U \subset X$가
존재하여 $f^{-1}(U) \to U$가 유한이다. 일반 평탄성
(명제 \ref{morphisms-proposition-generic-flatness})에 의해
$f^{-1}(U) \to U$가 평탄하고 유한표현이라고 가정해도 된다. 따라서
$f^{-1}(U) \to U$는 유한 국소 자유이다(보조정리
\ref{morphisms-lemma-finite-flat}). $f$가 쌍유리이므로 $X'$의 $X$ 위 차수는
$1$이다. 따라서 $f^{-1}(U) \to U$는 차수 $1$인 유한 국소 자유
사상이며, 달리 말하면 동형이다. 그러므로 수정을 다음과 같이
{\it 다시 정의}할 수 있다. 정수적 스킴들의 고유 사상
$f : X' \to X$ 가운데 $f^{-1}(U) \to U$가 어떤 공집합이 아닌
열린집합 $U \subset X$에 대해 동형인 것이다.

\begin{definition}
\label{morphisms-definition-alteration}
\begin{reference}
\cite[Definition 2.20]{alterations}
\end{reference}
$X$를 정수적 스킴이라 하자. {\it $X$의 변경}은 고유 우세 사상
$f : Y \to X$ 가운데 $Y$가 정수적이고 $f^{-1}(U) \to U$가 어떤
공집합이 아닌 열린집합 $U \subset X$에 대해 유한인 것이다.
\end{definition}

\noindent
이는 \cite{alterations}에 제시된 정의와 같지만, 여기서는 $X$와 $Y$가
뇌터일 것을 요구하지 않는다. 위와 같이 논증하면 변경은 함수체의
유한 확장을 유도하는 정수적 스킴들의 고유 우세 사상
$f : Y \to X$, 즉 보조정리 \ref{morphisms-lemma-finite-degree}의 동치인 조건들을
만족하는 사상임을 알 수 있다.

\section{차원 공식}
\label{morphisms-section-dimension-formula}

\noindent
뇌터 스킴들 사이의 사상에 대해서는 국소환의 차원에 관해 조금 더
말할 수 있다. 다음은 중요하면서도 그다지 증명하기 어렵지 않은
결과이다. $R(X)$가 정수적 스킴 $X$의 함수체를 나타냄을 상기하라.

\begin{lemma}
\label{morphisms-lemma-dimension-formula}
$S$를 스킴이라 하자.
$f : X \to S$를 스킴들의 사상이라 하자.
$x \in X$라 하고 $s = f(x)$로 두자. 다음을 가정하자.
\begin{enumerate}
\item $S$는 국소적으로 뇌터이다.
\item $f$는 국소적으로 유한형이다.
\item $X$와 $S$는 정수적이다.
\item $f$는 우세이다.
\end{enumerate}
다음 부등식이 성립한다.
\begin{equation}
\label{morphisms-equation-dimension-formula}
\dim(\mathcal{O}_{X, x})
\leq
\dim(\mathcal{O}_{S, s}) + \text{trdeg}_{R(S)}R(X)
- \text{trdeg}_{\kappa(s)} \kappa(x).
\end{equation}
더욱이 $S$가 보편적으로 카테나리이면 등호가 성립한다.
\end{lemma}

\begin{proof}
이에 대응하는 대수학의 명제는 대수학, 보조정리
\ref{algebra-lemma-dimension-formula}이다.
\end{proof}

\begin{lemma}
\label{morphisms-lemma-dimension-formula-general}
$S$를 스킴이라 하고 $f : X \to S$를 스킴들의 사상이라 하자.
$x \in X$라 하고 $s = f(x)$로 두자. $S$가 국소적으로 뇌터이고
$f$가 국소적으로 유한형이라고 가정하자. 다음 부등식이 성립한다.
\begin{equation}
\label{morphisms-equation-dimension-formula-general}
\dim(\mathcal{O}_{X, x})
\leq
\dim(\mathcal{O}_{S, s}) + E - \text{trdeg}_{\kappa(s)} \kappa(x).
\end{equation}
여기서 $E$는 $\text{trdeg}_{\kappa(f(\xi))}(\kappa(\xi))$가 갖는
값의 최댓값이고, 이때 $\xi$는 $X$의 기약 성분들의 일반점 가운데
$x$를 포함하는 성분의 일반점을 달린다.
\end{lemma}

\begin{proof}
$X_1, \ldots, X_n$을 $X$의 $x$를 포함하는 기약 성분들이라 하고,
각각에 축약 유도 스킴 구조를 주자. 이들은 극소 소아이디얼
$\mathfrak q_i$들에 대응하며, 이 아이디얼들은 $\mathcal{O}_{X, x}$의
극소 소아이디얼이므로
유한 개이다(스킴, 보조정리 \ref{schemes-lemma-specialize-points}와
대수학, 보조정리
\ref{algebra-lemma-Noetherian-irreducible-components}). 그러면
$\dim(\mathcal{O}_{X, x}) = \max \dim(\mathcal{O}_{X, x}/\mathfrak q_i)
= \max \dim(\mathcal{O}_{X_i, x})$이다. $\xi$들 가운데 $E$의 정의에
나타나는 것은 정확히 일반점 $\xi_i \in X_i$들이다.
$Z_i = \overline{\{f(\xi_i)\}} \subset S$라 하고 축약 유도 스킴
구조를 주자. 합성 $X_i \to X \to S$는 $Z_i$를 통해 분해된다
(스킴, 보조정리 \ref{schemes-lemma-map-into-reduction}). 따라서 차원
공식(보조정리 \ref{morphisms-lemma-dimension-formula})을 적용하여
$\dim(\mathcal{O}_{X_i, x}) \leq \dim(\mathcal{O}_{Z_i, x}) +
\text{trdeg}_{\kappa(f(\xi))}(\kappa(\xi)) -
\text{trdeg}_{\kappa(s)} \kappa(x)$임을 알 수 있다. 이들을 모두
합치면 보조정리를 얻는다.
\end{proof}

\noindent
한 가지 응용은 차원함수가 주어진 보편적으로 카테나리인 스킴 위의
임의의 유한형 스킴에 차원함수를 구성하는 것이다. 차원함수의 정의는
위상수학, 정의 \ref{topology-definition-dimension-function}를 보라.

\begin{lemma}
\label{morphisms-lemma-dimension-function-propagates}
$S$를 국소적으로 뇌터이고 보편적으로 카테나리인 스킴이라 하자.
$\delta : S \to \mathbf{Z}$를 차원함수라 하자.
$f : X \to S$를 스킴들의 사상이라 하자.
$f$가 국소적으로 유한형이라고 가정하자. 그러면 사상
\begin{align*}
\delta = \delta_{X/S} : X & \longrightarrow \mathbf{Z} \\
x & \longmapsto \delta(f(x)) + \text{trdeg}_{\kappa(f(x))} \kappa(x)
\end{align*}
은 $X$ 위의 차원함수이다.
\end{lemma}

\begin{proof}
$f : X \to S$가 국소적으로 유한형이라고 하자.
$x \leadsto y$, $x \not = y$를 $X$ 안의 특수화라 하자.
$\delta_{X/S}(x) > \delta_{X/S}(y)$임을 보이고, 또한
$\delta_{X/S}(x) = \delta_{X/S}(y) + 1$임을 보여야 한다. 뒤의 등식은
$y$가 $x$의 직접 특수화인 경우에 관한 것이다.

\medskip\noindent
아핀 열린집합 $V \subset S$를 택하여 $y$의 상을 포함하게 하고,
아핀 열린집합 $U \subset X$를 택하여 $V$로 사상되며 $y$를
포함하게 하자.
$X$를 $U$로, $S$를 $V$로 대체해도 명백히 무방하다. 따라서
$X = \Spec(A)$이고 $S = \Spec(R)$이며 $f$가 환 사상
$R \to A$로 주어진다고 가정해도 된다. 환 $R$은 보편적으로
카테나리이고(보조정리 \ref{morphisms-lemma-universally-catenary-local}), 사상
$R \to A$는 유한형이다(보조정리
\ref{morphisms-lemma-locally-finite-type-characterize}).

\medskip\noindent
$\mathfrak q \subset A$를 점 $x$에 대응하는 소아이디얼이라 하고,
$\mathfrak p \subset R$을 $f(x)$에 대응하는 소아이디얼이라 하자.
$\delta'$를 $\delta$의 $S' = \Spec(R/\mathfrak p) \subset S$에 대한
제한이라 하자. 이는 차원함수이다. 환 $R/\mathfrak p$는 보편적으로
카테나리이다. $\delta_{X/S}$를 $X' = \Spec(A/\mathfrak q)$에
제한한 함수는 $\delta_{X'/S'}$와 명백히 같으며, 뒤의 함수는
차원함수 $\delta'$를 사용하여 구성한 것이다. 따라서 위의 가정들에 더하여
$R \subset A$가 정역들이라고, 달리 말하면 $X$와 $S$가 정수적
스킴들이라고 가정해도 된다. 또한 $x$는 $X$의 일반점이고
$f(x)$는 $S$의 일반점이라고 가정해도 된다.

\medskip\noindent
$\mathcal{O}_{X, x} = R(X)$임에 유의하라. 또한
$x \leadsto y$, $x \not = y$이므로 $\mathcal{O}_{X, y}$의 스펙트럼은
적어도 두 점을 가진다(스킴, 보조정리
\ref{schemes-lemma-specialize-points}). 따라서
$\dim(\mathcal{O}_{X, y}) > 0$이다. $y$가 $x$의 직접 특수화이면
$\Spec(\mathcal{O}_{X, y}) = \{x, y\}$이고
$\dim(\mathcal{O}_{X, y}) = 1$이다.

\medskip\noindent
$s = f(x)$, $t = f(y)$라고 쓰자. 다음을 계산한다.
\begin{align*}
\delta_{X/S}(x) - \delta_{X/S}(y)
& =
\delta(s) + \text{trdeg}_{\kappa(s)} \kappa(x)
- \delta(t) - \text{trdeg}_{\kappa(t)} \kappa(y) \\
& =
\delta(s) - \delta(t) +
\text{trdeg}_{R(S)} R(X) - \text{trdeg}_{\kappa(t)} \kappa(y) \\
& =
\delta(s) - \delta(t) + \dim(\mathcal{O}_{X, y})
- \dim(\mathcal{O}_{S, t})
\end{align*}
마지막 단계에서 (\ref{morphisms-equation-dimension-formula})의 등호를 사용하였다.
$\delta$가 스킴 $S$ 위의 차원함수이고 $s \in S$가 일반점이므로,
차이 $\delta(s) - \delta(t)$는 위상수학, 보조정리
\ref{topology-lemma-dimension-function-catenary}에 의해
$\text{codim}(\overline{\{t\}}, S)$와 같다. 이는 성질, 보조정리
\ref{properties-lemma-codimension-local-ring}에 의해
$\dim(\mathcal{O}_{S, t})$와 같다. 따라서
$$
\delta_{X/S}(x) - \delta_{X/S}(y) = \dim(\mathcal{O}_{X, y})
$$
이고, 위에서 $\dim(\mathcal{O}_{X, y})$에 관해 말한 것으로부터
보조정리가 따른다.
\end{proof}

\noindent
차원 공식의 또 다른 응용은 차원이 ``변경'' 아래에서 바뀌지 않는다는
것이다(변경은 뒤에서 정의한다).

\begin{lemma}
\label{morphisms-lemma-alteration-dimension}
$f : X \to Y$를 스킴들의 사상이라 하자. 다음을 가정하자.
\begin{enumerate}
\item $Y$는 국소적으로 뇌터이다.
\item $X$와 $Y$는 정수적 스킴이다.
\item $f$는 우세이다.
\item $f$는 국소적으로 유한형이다.
\end{enumerate}
그러면
$$
\dim(X) \leq \dim(Y) + \text{trdeg}_{R(Y)} R(X).
$$
더 나아가 $f$가 닫힌 사상이면\footnote{예를 들어 $f$가 고유이면
정의 \ref{morphisms-definition-proper}를 보라.} 등호가 성립한다.
\end{lemma}

\begin{proof}
$f : X \to Y$를 보조정리에서와 같이 잡자.
$\xi_0 \leadsto \xi_1 \leadsto \ldots \leadsto \xi_e$를 $X$ 안의
특수화들의 열이라 하자. $x = \xi_e$, $y = f(x)$로 두자. 주어진
$e \leq \dim(\mathcal{O}_{X, x})$임에 유의하라. 이는 주어진
특수화들이 $\mathcal{O}_{X, x}$의 스펙트럼 안에서 일어나기 때문이다.
스킴, 보조정리
\ref{schemes-lemma-specialize-points}를 보라. 차원 공식인 보조정리
\ref{morphisms-lemma-dimension-formula}에 의해 다음을 얻는다.
\begin{align*}
e & \leq \dim(\mathcal{O}_{X, x}) \\
& \leq \dim(\mathcal{O}_{Y, y}) +
\text{trdeg}_{R(Y)} R(X) - \text{trdeg}_{\kappa(y)} \kappa(x) \\
& \leq \dim(\mathcal{O}_{Y, y}) + \text{trdeg}_{R(Y)} R(X)
\end{align*}
따라서 원하는 대로 $e \leq \dim(Y) + \text{trdeg}_{R(Y)} R(X)$임을
얻는다.

\medskip\noindent
이제 $f$가 닫힌 사상이기도 하다고 가정하자.
$\overline{\xi}_0 \leadsto \overline{\xi}_1 \leadsto \ldots
\leadsto \overline{\xi}_d$를 $Y$ 안의 특수화들의 열이라 하자.
$\dim(X) \geq d + r$임을 보이고자 한다.
$\overline{\xi}_0 = \eta$가 $Y$의 일반점이라고 가정해도 된다.
일반올 $X_\eta$는 $\kappa(\eta) = R(Y)$ 위에서 국소적으로 유한형인
스킴이다. $f$가 우세이므로 이는 공집합이 아니다. 따라서 보조정리
\ref{morphisms-lemma-ubiquity-Jacobson-schemes}에 의해 야콥슨 스킴이다.
그러므로 보조정리 \ref{morphisms-lemma-jacobson-finite-type-points}에 의해
닫힌점 $\xi_0 \in X_\eta$를 찾을 수 있고 확장
$\kappa(\eta) \subset \kappa(\xi_0)$은 유한이다.
$\mathcal{O}_{X, \xi_0} = \mathcal{O}_{X_\eta, \xi_0}$임에 유의하라.
이는 $\eta$가 $Y$의 일반점이기 때문이다. 따라서
$\dim(\mathcal{O}_{X, \xi_0}) = r$임을 알 수 있다. 실제로 보조정리
\ref{morphisms-lemma-dimension-formula}을 스킴 $X_\eta$, 보편적으로 카테나리인
밑스킴 $\Spec(\kappa(\eta))$, 그리고 점 $\xi_0$에 적용하면 된다
(보조정리 \ref{morphisms-lemma-ubiquity-uc}도 보라). 이는
$\xi_{-r} \leadsto \ldots \leadsto \xi_{-1} \leadsto \xi_0$가 $X$ 안에
존재한다는 뜻이다. 한편 $f$가 닫힌 사상이므로 특수화는 $f$를 따라
올려진다. 위상수학, 보조정리
\ref{topology-lemma-closed-open-map-specialization}를 보라. 따라서
$\xi_0$가 $\eta = \overline{\xi}_0$ 위에 놓이므로, 특수화
$\xi_0 \leadsto \xi_1 \leadsto \ldots \leadsto \xi_d$를 찾을 수 있고,
이는 $\overline{\xi}_0 \leadsto \overline{\xi}_1 \leadsto \ldots
\leadsto \overline{\xi}_d$ 위에 놓인다.
달리 말하면
$$
\xi_{-r} \leadsto \ldots \leadsto \xi_{-1} \leadsto \xi_0
\leadsto \xi_1 \leadsto \ldots \leadsto \xi_d
$$
이고, 이는 원하는 대로 $\dim(X) \geq d + r$이라는 뜻이다.
\end{proof}

\begin{lemma}
\label{morphisms-lemma-alteration-dimension-general}
$f : X \to Y$를 스킴들의 사상이라 하자. $Y$가 국소적으로 뇌터이고
$f$가 국소적으로 유한형이라고 가정하자. 그러면
$$
\dim(X) \leq \dim(Y) + E
$$
이다. 여기서 $E$는 $\text{trdeg}_{\kappa(f(\xi))}(\kappa(\xi))$의
상한이고, 이때 $\xi$는 $X$의 기약 성분들의 일반점을 달린다.
\end{lemma}

\begin{proof}
보조정리 \ref{morphisms-lemma-dimension-formula-general}과 성질, 보조정리
\ref{properties-lemma-dimension}에서 즉시 따른다.
\end{proof}

\section{상대 정규화}
\label{morphisms-section-normalization-X-in-Y}

\noindent
이 절에서는 {\it 한 스킴의 다른 스킴 안에서의 정규화}를 구성한다.

\begin{lemma}
\label{morphisms-lemma-integral-closure}
$X$를 스킴이라 하자. $\mathcal{A}$를 $\mathcal{O}_X$-대수들의
준연접 층이라 하자. 다음 규칙으로 정의한 부분층
$\mathcal{A}' \subset \mathcal{A}$를 생각하자.
$$
U \longmapsto \{f \in \mathcal{A}(U) \mid
f_x \in \mathcal{A}_x \text{ integral over } \mathcal{O}_{X, x}
\text{ for all }x \in U\}
$$
그러면 이는 준연접 $\mathcal{O}_X$-대수이고, 줄기 $\mathcal{A}'_x$는
$\mathcal{O}_{X, x}$의 $\mathcal{A}_x$ 안에서의 정수적 폐포이다.
또한 임의의 아핀 열린집합 $U \subset X$에 대해 환
$\mathcal{A}'(U) \subset \mathcal{A}(U)$는 $\mathcal{O}_X(U)$의
$\mathcal{A}(U)$ 안에서의 정수적 폐포이다.
\end{lemma}

\begin{proof}
조건들이 국소적이므로 이는 부분층이다. 대수학, 보조정리
\ref{algebra-lemma-integral-closure-is-ring}에 의해
$\mathcal{O}_X$-대수이다. $U \subset X$를 아핀 열린집합이라 하자.
$U = \Spec(R)$이라고 쓰고, $\mathcal{A}$가 $R$-대수 $A$에 연관된
준연접 층이라고 하자. 그러면 대수학, 보조정리
\ref{algebra-lemma-integral-closure-stalks}에 의해 $\mathcal{A}'$의
$U$ 위 값은 정수적 폐포 $A'$로 주어지며, 이는 $R$의 $A$ 안에서의
정수적 폐포이다. 이는 보조정리의 마지막 주장을 증명한다.
$\mathcal{A}'$가 준연접임을 증명하려면
$\mathcal{A}'(D(f)) = A'_f$임을 보이면 충분하다. 이는 정수적 폐포와
국소화가 교환한다는 사실에서 따른다. 대수학, 보조정리
\ref{algebra-lemma-integral-closure-localize}를 보라. 같은 사실로부터
줄기들에 관한 주장도 따른다.
\end{proof}

\begin{definition}
\label{morphisms-definition-integral-closure}
$X$를 스킴이라 하자. $\mathcal{A}$를 $\mathcal{O}_X$-대수들의
준연접 층이라 하자. {\it $\mathcal{O}_X$의 $\mathcal{A}$ 안에서의
정수적 폐포}는 위의 보조정리
\ref{morphisms-lemma-integral-closure}에서 구성한 준연접
$\mathcal{O}_X$-부분대수 $\mathcal{A}' \subset \mathcal{A}$이다.
\end{definition}

\noindent
위 정의의 상황에서는 상대 스펙트럼의 사상
$$
\xymatrix{
Y = \underline{\Spec}_X(\mathcal{A}) \ar[rr] \ar[rd] & &
X' = \underline{\Spec}_X(\mathcal{A}') \ar[ld] \\
& X &
}
$$
을 생각할 수 있다. 보조정리 \ref{morphisms-lemma-affine-equivalence-algebras}를
보라. 스킴 $X' \to X$는 $X$의 스킴 $Y$ 안에서의 정규화가 될 것이다.
이제 조금 더 일반적인 상황을 생각하자. 스킴들의 준콤팩트 준분리 사상
$f : Y \to X$가 주어졌다고 하자. 이 경우 $\mathcal{O}_X$-대수들의 층
$f_*\mathcal{O}_Y$는 준연접이다. 스킴, 보조정리
\ref{schemes-lemma-push-forward-quasi-coherent}를 보라. 정수적 폐포
$\mathcal{O}' \subset f_*\mathcal{O}_Y$를 취하면 $\mathcal{O}_X$-대수들의
준연접 층을 얻으며, 그 상대 스펙트럼이 $X$의 $Y$ 안에서의
정규화이다. 다음이 형식적인 정의이다.

\begin{definition}
\label{morphisms-definition-normalization-X-in-Y}
$f : Y \to X$를 스킴들의 준콤팩트 준분리 사상이라 하자.
$\mathcal{O}'$를 $\mathcal{O}_X$의 $f_*\mathcal{O}_Y$ 안에서의 정수적
폐포라 하자. {\it $X$의 $Y$ 안에서의 정규화}는 스킴
\footnote{스킴 $X'$는 정규일 필요가 없다. 예를 들어 $Y = X$이고
$f = \text{id}_X$이면 $X' = X$이다.}
$$
\nu : X' = \underline{\Spec}_X(\mathcal{O}') \to X
$$
이며 $X$ 위에 놓인다. 여기에는 원래 사상 $f$의 자연스러운 분해
$$
Y \xrightarrow{f'} X' \xrightarrow{\nu} X
$$
가 함께 주어진다.
\end{definition}

\noindent
이 분해는 표준 사상 $Y \to \underline{\Spec}_X(f_*\mathcal{O}_Y)$
(구성, 보조정리 \ref{constructions-lemma-canonical-morphism}를 보라)과
포함 사상 $\mathcal{O}' \to f_*\mathcal{O}_Y$에서 오는 상대
스펙트럼의 사상의 합성이다. 정규화는 다음과 같이 특징지을 수 있다.

\begin{lemma}
\label{morphisms-lemma-characterize-normalization}
$f : Y \to X$를 스킴들의 준콤팩트 준분리 사상이라 하자.
분해 $f = \nu \circ f'$를 생각하자. 여기서 $\nu : X' \to X$는
$X$의 $Y$ 안에서의 정규화이다. 이 분해는 다음 두 성질로
특징지어진다.
\begin{enumerate}
\item 사상 $\nu$는 정수적이다.
\item 임의의 분해 $f = \pi \circ g$에서 $\pi : Z \to X$가
정수적이면, 어떤 유일한 사상 $h : X' \to Z$에 대해 가환 그림
$$
\xymatrix{
Y \ar[d]_{f'} \ar[r]_g & Z \ar[d]^\pi \\
X' \ar[ru]^h \ar[r]^\nu & X
}
$$
이 존재한다.
\end{enumerate}
더욱이 사상 $f' : Y \to X'$는 우세이고, (2)의 사상
$h : X' \to Z$는 $Z$의 $Y$ 안에서의 정규화이다.
\end{lemma}

\begin{proof}
$\mathcal{O}' \subset f_*\mathcal{O}_Y$를 정의
\ref{morphisms-definition-normalization-X-in-Y}에서와 같이 $\mathcal{O}_X$의
정수적 폐포라 하자. 사상 $\nu$는 구성에 의해 정수적이므로 (1)을
증명한다. (2)에서와 같이 분해 $f = \pi \circ g$가 주어지고
$\pi : Z \to X$가 정수적이라고 하자. 정의
\ref{morphisms-definition-integral}에 의해 $\pi$는 아핀이며, 따라서 $Z$는
$\mathcal{O}_X$-대수들의 어떤 준연접 층 $\mathcal{B}$의 상대
스펙트럼이다. 사상 $g : Y \to Z$는 $\mathcal{O}_X$-대수들의 사상
$\chi : \mathcal{B} \to f_*\mathcal{O}_Y$에 대응한다.
$\mathcal{B}(U)$가 $\mathcal{O}_X(U)$ 위에서 정수적임은 임의의 아핀
열린집합 $U \subset X$에 대해 성립하므로(정의
\ref{morphisms-definition-integral}), 보조정리 \ref{morphisms-lemma-integral-closure}에서
$\chi(\mathcal{B}) \subset \mathcal{O}'$임을 알 수 있다. 상대
스펙트럼의 함자성인 보조정리
\ref{morphisms-lemma-affine-equivalence-algebras}에 의해 유일한 사상
$h : X' \to Z$를 얻는다. 그림이 가환이라는 확인은 생략한다.

\medskip\noindent
(1)과 (2)는 이 분해를 어떤 범주의 시작 대상이라는 조건으로
특징짓기 때문에, 분해 $f = \nu \circ f'$를 특징지음이 명백하다.

\medskip\noindent
(2)의 보편 성질에서 $f'$가 $X'$의 진정한 닫힌 부분스킴을 통해
분해되지 않음을 알 수 있다. 따라서 $f'$의 스킴론적 상은 $X'$이다.
$f'$는 준콤팩트이다(스킴, 보조정리
\ref{schemes-lemma-quasi-compact-permanence}와 $\nu$가 아핀 사상이어서
분리라는 사실에 의한다). 따라서 $f'(Y)$는 $X'$에서 조밀하다.
그러므로 $f'$는 우세이다.

\medskip\noindent
$g$가 준콤팩트이고 준분리임에 유의하라. 이는 스킴, 보조정리
\ref{schemes-lemma-compose-after-separated}와
\ref{schemes-lemma-quasi-compact-permanence}에서 따른다. 따라서
보조정리의 마지막 명제는 의미가 있다. (2)의 사상 $h$는 보조정리
\ref{morphisms-lemma-finite-permanence}에 의해 정수적이다. 분해
$g = \pi' \circ g'$가 주어지고 $\pi' : Z' \to Z$가 정수적이면, 분해
$f = (\pi \circ \pi') \circ g'$를 얻고 사상 $h' : X' \to Z'$도
얻는다. 유일성에 의해 $\pi' \circ h' = h$이다. 따라서 특징지음
(1), (2)를 사상 $h : X' \to Z$에 적용하면 보조정리의 마지막 주장을
얻는다.
\end{proof}

\begin{lemma}
\label{morphisms-lemma-functoriality-normalization}
다음이 스킴 사상들의 가환 그림이라고 하자.
$$
\xymatrix{
Y_2 \ar[d]_{f_2} \ar[r] & Y_1 \ar[d]^{f_1} \\
X_2 \ar[r] & X_1
}
$$
$f_1$, $f_2$가 준콤팩트이고 준분리라고 가정하자.
$f_i = \nu_i \circ f_i'$, $i = 1, 2$를 표준 분해라 하자. 여기서
$\nu_i : X_i' \to X_i$는 $X_i$의 $Y_i$ 안에서의 정규화이다. 그러면
다음 가환 그림에 들어가는 유일한 화살표 $X'_2 \to X'_1$가 존재한다.
$$
\xymatrix{
Y_2 \ar[d]_{f_2'} \ar[r] & Y_1 \ar[d]^{f_1'} \\
X_2' \ar[d]_{\nu_2} \ar[r] & X_1' \ar[d]^{\nu_1} \\
X_2 \ar[r] & X_1
}
$$
\end{lemma}

\begin{proof}
보조정리 \ref{morphisms-lemma-characterize-normalization}의 (1)과 보조정리
\ref{morphisms-lemma-base-change-finite}에 의해 밑변환
$X_2 \times_{X_1} X'_1 \to X_2$는 정수적이다. $f_2$가 이 사상을
통해 분해됨에 유의하라. 따라서 보조정리
\ref{morphisms-lemma-characterize-normalization}의 (2)에 의해 유일한 사상
$X'_2 \to X_2 \times_{X_1} X'_1$를 얻는다. 이는 가환 그림에 들어가는
화살표 $X'_2 \to X'_1$를 주며, 유일성도 함께 따른다.
\end{proof}

\begin{lemma}
\label{morphisms-lemma-normalization-localization}
$f : Y \to X$를 스킴들의 준콤팩트 준분리 사상이라 하자.
$U \subset X$를 열린 부분스킴이라 하고 $V = f^{-1}(U)$로 두자.
그러면 $U$의 $V$ 안에서의 정규화는 $U$의 역상이며, 이 역상은
$X$의 $Y$ 안에서의 정규화에서 취한다.
\end{lemma}

\begin{proof}
구성에서 명백하다.
\end{proof}

\begin{lemma}
\label{morphisms-lemma-normalization-is-normalization}
$f : Y \to X$를 스킴들의 준콤팩트 준분리 사상이라 하자.
$X'$를 $X$의 $Y$ 안에서의 정규화라 하자. 그러면 $X'$의 $Y$ 안에서의
정규화는 $X'$이다.
\end{lemma}

\begin{proof}
$Y \to X'' \to X'$가 $X'$의 $Y$ 안에서의 정규화이면, 합성
$X'' \to X$에 보조정리 \ref{morphisms-lemma-characterize-normalization}을
적용하여 표준 사상 $h : X' \to X''$를 얻으며 이는 $X$ 위에 있다. 사상
$h$와 $X'' \to X'$가 서로 역이라는 확인은 생략한다(보조정리의
분해의 유일성을 사용한다).
\end{proof}

\begin{lemma}
\label{morphisms-lemma-normalization-in-reduced}
$f : Y \to X$를 스킴들의 준콤팩트 준분리 사상이라 하자.
$X' \to X$를 $X$의 $Y$ 안에서의 정규화라 하자. $Y$가 축약이면
$X'$도 축약이다.
\end{lemma}

\begin{proof}
축약환의 부분환이 축약이라는 사실에서 따른다. 몇몇 세부사항은
생략한다.
\end{proof}

\begin{lemma}
\label{morphisms-lemma-normalization-generic}
$f : Y \to X$를 스킴들의 준콤팩트 준분리 사상이라 하자.
$X' \to X$를 $X$의 $Y$ 안에서의 정규화라 하자. $X'$의 기약 성분의
모든 일반점은 $Y$의 기약 성분의 어떤 일반점의 상이다.
\end{lemma}

\begin{proof}
보조정리 \ref{morphisms-lemma-normalization-localization}에 의해
$X = \Spec(A)$가 아핀이라고 가정해도 된다. 유한 아핀 열린덮개
$Y = \bigcup \Spec(B_i)$를 택하자. 그러면 $X' = \Spec(A')$이고,
사상들 $\Spec(B_i) \to Y \to X'$는 함께 단사인 $A$-대수 사상
$A' \to \prod B_i$를 정의한다. 따라서 대수학, 보조정리
\ref{algebra-lemma-injective-minimal-primes-in-image}에서 보조정리가
따른다.
\end{proof}

\begin{lemma}
\label{morphisms-lemma-normalization-in-disjoint-union}
$f : Y \to X$를 스킴들의 준콤팩트 준분리 사상이라 하자.
$Y = Y_1 \amalg Y_2$가 두 스킴의 분리합이라고 가정하자.
$f_i = f|_{Y_i}$라고 쓰자. $X_i'$를 $X$의 $Y_i$ 안에서의 정규화라
하자. 그러면 $X_1' \amalg X_2'$는 $X$의 $Y$ 안에서의 정규화이다.
\end{lemma}

\begin{proof}
정수적 폐포의 말로 쓰면 이는 다음 사실에 대응한다. $A \to B$를
환 사상이라 하자. $B = B_1 \times B_2$라고 가정하자. $A_i'$를
$A$의 $B_i$ 안에서의 정수적 폐포라 하자. 그러면
$A_1' \times A_2'$는 $A$의 $B$ 안에서의 정수적 폐포이다. 이것이
성립하는 이유는 원소 $(1, 0)$과 $(0, 1)$이 $B$에 속하는 멱등원이어서
$A$ 위에서 정수적이기 때문이다. 따라서 정수적 폐포 $A'$, 즉
$A$의 $B$ 안에서의 정수적 폐포는 곱이고, 그 인자들이 위에서 기술한 정수적 폐포
$A'_i$임을 어렵지 않게 알 수 있다(몇몇 세부사항은 생략한다).
\end{proof}

\begin{lemma}
\label{morphisms-lemma-normalization-in-universally-closed}
$f : X \to S$를 스킴들의 준콤팩트 준분리 보편 닫힌 사상이라 하자.
그러면 $f_*\mathcal{O}_X$는 $\mathcal{O}_S$ 위에서 정수적이다.
달리 말하면 $S$의 $X$ 안에서의 정규화는 구성, 보조정리
\ref{constructions-lemma-canonical-morphism}의 분해
$$
X \longrightarrow \underline{\Spec}_S(f_*\mathcal{O}_X)
\longrightarrow S
$$
와 같다.
\end{lemma}

\begin{proof}
문제는 $S$에서 국소적이므로 $S = \Spec(R)$가 아핀이라고 가정해도
된다. $h \in \Gamma(X, \mathcal{O}_X)$라 하자. $h$가 $R$ 위의
일계수 방정식을 만족함을 보여야 한다. $h$를 다음 가환 그림의
사상으로 생각하자.
$$
\xymatrix{
X \ar[rr]_h \ar[rd]_f & & \mathbf{A}^1_S \ar[ld] \\
& S &
}
$$
$Z \subset \mathbf{A}^1_S$를 $h$의 스킴론적 상이라 하자. 정의
\ref{morphisms-definition-scheme-theoretic-image}를 보라. $h$는 준콤팩트이다.
실제로 $f$가 준콤팩트이고 $\mathbf{A}^1_S \to S$가 분리이다. 스킴,
보조정리 \ref{schemes-lemma-quasi-compact-permanence}를 보라.
보조정리 \ref{morphisms-lemma-quasi-compact-scheme-theoretic-image}에 의해 사상
$X \to Z$는 우세이다. 보조정리
\ref{morphisms-lemma-image-proper-scheme-closed}에 의해 사상 $X \to Z$는
닫힌 사상이다. 따라서 집합론적으로 $h(X) = Z$이다. 그러므로
보조정리 \ref{morphisms-lemma-image-proper-is-proper}를 사용하여 $Z \to S$가
보편 닫힌 사상이고 심지어 고유임을 얻는다. $Z \subset \mathbf{A}^1_S$
이므로 $Z \to S$는 아핀이고 고유이며, 따라서 보조정리
\ref{morphisms-lemma-integral-universally-closed}에 의해 정수적이다.
$\mathbf{A}^1_S = \Spec(R[T])$라고 쓰면 아이디얼 $I \subset R[T]$,
즉 $Z$의 아이디얼이 일계수 다항식 $P(T) \in R[T]$를 포함함을 얻는다.
따라서 $P(h) = 0$이고 증명이 끝난다.
\end{proof}

\begin{lemma}
\label{morphisms-lemma-normalization-in-integral}
$f : Y \to X$를 정수적 사상이라 하자. 그러면 $X$의 $Y$ 안에서의
정규화는 $Y$와 같다.
\end{lemma}

\begin{proof}
보조정리 \ref{morphisms-lemma-integral-universally-closed}에 의해 이는 보조정리
\ref{morphisms-lemma-normalization-in-universally-closed}의 특수한 경우이다.
\end{proof}

\begin{lemma}
\label{morphisms-lemma-normal-normalization}
$f : Y \to X$를 스킴들의 준콤팩트 준분리 사상이라 하자. $X'$를
$X$의 $Y$ 안에서의 정규화라 하자. 다음을 가정하자.
\begin{enumerate}
\item $Y$는 정규 스킴이다.
\item $Y$의 준콤팩트 열린집합은 유한 개의 기약 성분을 가진다.
\end{enumerate}
그러면 $X'$는 정수적 정규 스킴들의 분리합이다. 더욱이 사상
$Y \to X'$는 우세이고 기약 성분들의 전단사를 유도한다.
\end{lemma}

\begin{proof}
$U \subset X$를 아핀 열린집합이라 하자. 역상 $U'$를 생각하자. 이는
$U$의 $X'$에서의 역상이다. $V = f^{-1}(U)$로 두자. 보조정리
\ref{morphisms-lemma-normalization-localization}에 의해 $V \to U' \to U$는
$U$의 $V$ 안에서의 정규화이다. $U = \Spec(A)$라고 쓰자. 그러면
$V$는 준콤팩트이므로 가정에 의해 유한 개의 기약 성분을 가진다.
따라서 성질, 보조정리
\ref{properties-lemma-normal-locally-finite-nr-irreducibles}에 의해
$V = \coprod_{i = 1, \ldots n} V_i$는 정규 정수적 스킴들의 유한
분리합이다. 보조정리 \ref{morphisms-lemma-normalization-in-disjoint-union}에 의해
$U' = \coprod_{i = 1, \ldots, n} U_i'$임을 알 수 있다. 여기서
$U'_i$는 $U$의 $V_i$ 안에서의 정규화이다. 성질, 보조정리
\ref{properties-lemma-normal-integral-sections}에 의해
$B_i = \Gamma(V_i, \mathcal{O}_{V_i})$는 정규 정역이다.
$U_i' = \Spec(A_i')$이고 $A_i' \subset B_i$는 $A$의 $B_i$ 안에서의
정수적 폐포임에 유의하라. 보조정리
\ref{morphisms-lemma-integral-closure}를 보라. 대수학, 보조정리
\ref{algebra-lemma-integral-closure-in-normal}에 의해
$A_i' \subset B_i$는 정규 정역이다. 따라서
$U' = \coprod U_i'$는 정규 정수적 스킴들의 유한 합집합이고, 그러므로
정규이다.

\medskip\noindent
$X'$는 스킴 $U'$들에 의한 열린덮개를 가지므로, 성질, 보조정리
\ref{properties-lemma-locally-normal}에서 $X'$가 정규임을 얻는다.
한편 각 $U'$는 기약 스킴들의 유한 분리합이므로, $X'$의 모든
준콤팩트 열린집합은 유한 개의 기약 성분을 가진다(생략하는 위상적
논증에 의한다). 따라서 성질, 보조정리
\ref{properties-lemma-normal-locally-finite-nr-irreducibles}에 의해
$X'$는 정규 정수적 스킴들의 분리합이다. 위의 $X'$의 기술에서
$Y \to X'$가 우세이고, 기약 성분들의 전단사 $V \to U'$를 모든
아핀 열린집합 $U \subset X$에 대해 유도함이 명백하다. 사상
$Y \to X'$에 대한 기약 성분들의 전단사는 이로부터 위상적 논증으로
따른다(생략한다).
\end{proof}

\begin{lemma}
\label{morphisms-lemma-nagata-normalization-finite-general}
$f : X \to S$를 사상이라 하자. 다음을 가정하자.
\begin{enumerate}
\item $S$는 나가타 스킴이다.
\item $f$는 준콤팩트이고 준분리이다.
\item $X$의 준콤팩트 열린집합은 유한 개의 기약 성분을 가진다.
\item $x \in X$가 한 기약 성분의 일반점이면 체확장
$\kappa(x)/\kappa(f(x))$는 유한생성이다.
\item $X$는 축약이다.
\end{enumerate}
그러면 정규화 $\nu : S' \to S$, 즉 $S$의 $X$ 안에서의 정규화는
유한이다.
\end{lemma}

\begin{proof}
가정 (1)에 의해 $S = \Spec(R)$이고 $R$이 나가타 환인 경우로 즉시
환원된다. 정수적 폐포 $A$, 즉 $R$의 $\Gamma(X, \mathcal{O}_X)$
안에서의 정수적 폐포가 $R$ 위에서 유한임을 보여야 한다. 가정 (2)에
의해 $f$가 준콤팩트이므로 $X = \bigcup_{i = 1, \ldots, n} U_i$로
쓸 수 있고, 각 $U_i$는 아핀이다.
$U_i = \Spec(B_i)$라고 하자. 각 $B_i$는 가정 (5)에 의해 축약이고,
가정 (3)과 대수학, 보조정리 \ref{algebra-lemma-irreducible}에 의해
유한 개의 극소 소아이디얼
$\mathfrak q_{i1}, \ldots, \mathfrak q_{im_i}$를 가진다. 다음 포함들이
성립한다.
$$
\Gamma(X, \mathcal{O}_X) \subset B_1 \times \ldots \times B_n
\subset
\prod\nolimits_{i = 1, \ldots, n}
\prod\nolimits_{j = 1, \ldots, m_i} (B_i)_{\mathfrak q_{ij}}
$$
두 번째 포함은 대수학, 보조정리
\ref{algebra-lemma-reduced-ring-sub-product-fields}에서 따른다.
대수학, 보조정리
\ref{algebra-lemma-minimal-prime-reduced-ring}에 의해
$\kappa(\mathfrak q_{ij}) = (B_i)_{\mathfrak q_{ij}}$이다. 따라서
정수적 폐포 $A$, 즉 $R$의 $\Gamma(X, \mathcal{O}_X)$ 안에서의
정수적 폐포는 정수적 폐포 $A_{ij}$들의 곱에 포함된다. 여기서 이들은
$R$의 $\kappa(\mathfrak q_{ij})$ 안에서의 정수적 폐포이다.
$R$이 뇌터이므로 $A_{ij}$가 유한 $R$-가군임을 각 $i, j$에 대해
보이면 충분하다. $\mathfrak p_{ij} \subset R$을
$\mathfrak q_{ij}$의 상이라 하자. 가정 (4)에 의해
$\kappa(\mathfrak q_{ij})/\kappa(\mathfrak p_{ij})$는 유한생성
체확장이므로 $R \to \kappa(\mathfrak q_{ij})$는 본질적으로
유한형이다. 따라서 대수학, 보조정리
\ref{algebra-lemma-nagata-in-reduced-finite-type-finite-integral-closure}에
의해 $R \to A_{ij}$는 유한이다.
\end{proof}

\begin{lemma}
\label{morphisms-lemma-nagata-normalization-finite}
$f : X \to S$를 사상이라 하자. 다음을 가정하자.
\begin{enumerate}
\item $S$는 나가타 스킴이다.
\item $f$는 유한형이다.
\item $X$는 축약이다.
\end{enumerate}
그러면 정규화 $\nu : S' \to S$, 즉 $S$의 $X$ 안에서의 정규화는
유한이다.
\end{lemma}

\begin{proof}
이는 보조정리 \ref{morphisms-lemma-nagata-normalization-finite-general}의 특수한
경우이다. 실제로 유한형 사상 $f$는 정의에 의해 준콤팩트이고,
보조정리 \ref{morphisms-lemma-finite-type-Noetherian-quasi-separated}에 의해
준분리이므로 (2)가 성립한다. $X$는 보조정리
\ref{morphisms-lemma-finite-type-noetherian}에 의해 국소적으로 뇌터이므로
조건 (3)이 성립한다. 마지막으로 유한형 사상은 유한생성 잉여체
확장을 유도하므로 조건 (4)가 성립한다.
\end{proof}

\begin{lemma}
\label{morphisms-lemma-relative-normalization-normal-codim-1}
$f : Y \to X$를 스킴들의 유한형 사상이라 하되 $Y$는 축약이고 $X$는
나가타라고 하자. $X'$를 $X$의 $Y$ 안에서의 정규화라 하자.
$x' \in X'$가 다음을 만족하는 점이라고 하자.
\begin{enumerate}
\item $\dim(\mathcal{O}_{X', x'}) = 1$이다.
\item $Y \to X'$의 $x'$ 위 올은 공집합이다.
\end{enumerate}
그러면 $\mathcal{O}_{X', x'}$는 이산 값매김환이다.
\end{lemma}

\begin{proof}
$X$를 $x'$의 상의 아핀 근방으로 대체할 수 있다. 따라서
$X = \Spec(A)$라고 가정해도 되며 $A$는 나가타이다. 보조정리
\ref{morphisms-lemma-nagata-normalization-finite}에 의해 사상 $X' \to X$는
유한이다. 따라서 $X' = \Spec(A')$로 쓸 수 있고, 여기서 $A$-대수
$A'$는 유한이다. 보조정리 \ref{morphisms-lemma-normalization-is-normalization}에 의해
$X$를 $X'$로 대체하면 다음 문단에서 기술하는 경우로 환원된다.

\medskip\noindent
$X = X' = \Spec(A)$이고 $A$가 뇌터인 경우를 생각하자.
$\mathfrak p \subset A$를 점 $x'$에 대응하는 소아이디얼이라 하자.
$g \in \mathfrak p$를 택하되 $A$의 어느 극소 소아이디얼에도
포함되지 않게 하자(소아이디얼 회피와 $A$가 유한 개의 극소
소아이디얼을 갖는다는 사실을 사용한다. 대수학, 보조정리
\ref{algebra-lemma-silly}와
\ref{algebra-lemma-Noetherian-irreducible-components}를 보라).
$Z = f^{-1}V(g) \subset Y$로 두자. 이는 $Y$의 닫힌 부분스킴이다.
그러면 $f(Z)$는 $g$의 선택에 의해 어떤 일반점도 포함하지 않고,
$x'$도 포함하지 않는다. 이는 $x'$가 $f$의 상에 없기 때문이다. $f(Z)$의 폐포는
보조정리 \ref{morphisms-lemma-reach-points-scheme-theoretic-image}에 의해
$f(Z)$의 점들의 특수화들로 이루어진 집합이다. 따라서 $f(Z)$의 폐포는
$x'$를 포함하지 않는다. 실제로 조건
$\dim(\mathcal{O}_{X', x'}) = 1$에 의해 $X = X'$의 일반점들만
$x'$로 특수화한다. 달리 말하면 $X$를 $x'$의 아핀 열린 근방으로
대체한 뒤 $f^{-1}V(g) = \emptyset$라고 가정해도 된다. 따라서 $g$는
$Y$ 위의 가역인 대역 함수로 사상되고, 다음 분해를 얻는다.
$$
A \to A_g \to \Gamma(Y, \mathcal{O}_Y)
$$
$X = X'$이므로 이는 $A$가 $A$의 $A_g$ 안에서의 정수적 폐포와 같음을
뜻한다. 대수학, 보조정리
\ref{algebra-lemma-integral-closure-localize}에 의해 $A_\mathfrak p$는
$A_\mathfrak p$의 $A_\mathfrak p[1/g]$ 안에서의 정수적 폐포이다.
$g$의 선택에 의해, $\dim(A_\mathfrak p) = 1$이고 $A$가 축약이므로
$A_\mathfrak p[1/g]$는 체들의 유한곱이다($\mathfrak p$에 포함된 극소
소아이디얼들의 잉여체들의 곱이다). 따라서 $A_\mathfrak p$는
정규이다(대수학, 보조정리
\ref{algebra-lemma-characterize-reduced-ring-normal}). 이로써 증명이
끝난다. 몇몇 세부사항은 생략한다.
\end{proof}

\section{정규화}
\label{morphisms-section-normalization}

\noindent
이제 스킴 $X$의 정규화를 살펴본다.
$X$가 국소적으로 유한 개의 기약 성분을 가질 때에만 이를 정의하고
구성한다. 각 준콤팩트 열린집합이 유한 개의 기약 성분을 갖는 스킴
$X$를 택하자. $X^{(0)} \subset X$를 $X$의 기약 성분들의 일반점
집합이라 하자. 다음을 생각하자.
\begin{equation}
\label{morphisms-equation-generic-points}
f :
Y = \coprod\nolimits_{\eta \in X^{(0)}} \Spec(\kappa(\eta))
\longrightarrow
X
\end{equation}
이는 스킴, 절 \ref{schemes-section-points}의 표준 사상들을 사용하여
일반점들을 $X$에 포함시키는 사상이다. 이 사상은 가정에 의해
준콤팩트이고, $Y$가 분리되어 있으므로 준분리임에 유의하라(스킴, 절
\ref{schemes-section-separation-axioms}를 보라).

\begin{definition}
\label{morphisms-definition-normalization}
각 준콤팩트 열린집합이 유한 개의 기약 성분을 갖는 스킴 $X$를
생각하자. $X$의 {\it 정규화}를 다음 사상으로 정의한다.
$$
\nu : X^\nu \longrightarrow X
$$
이는 $X$의, 위에서 구성한 사상 $f : Y \to X$에서의 정규화이다
((\ref{morphisms-equation-generic-points}) 참조).
\end{definition}

\noindent
모든 국소 뇌터 스킴은 국소적으로 유한한 기약 성분 집합을 가지므로
이 정의가 적용된다. 보통 정규화는 축약 스킴에 대해서만 정의한다.
위의 정의에서 $X$의 정규화는 $X_{red}$, 즉 $X$의 축약의 정규화와 같다.

\begin{lemma}
\label{morphisms-lemma-normalization-reduced}
각 준콤팩트 열린집합이 유한 개의 기약 성분을 갖는 스킴 $X$를
생각하자. 정규화 사상 $\nu$는 축약 $X_{red}$를 통해 분해되고,
$X^\nu \to X_{red}$는 $X_{red}$의 정규화이다.
\end{lemma}

\begin{proof}
$f : Y \to X$를 (\ref{morphisms-equation-generic-points})의 사상이라 하자.
스킴, 보조정리 \ref{schemes-lemma-map-into-reduction}에서 분해
$Y \to X_{red} \to X$, 즉 $f$의 분해를 얻는다. 보조정리
\ref{morphisms-lemma-characterize-normalization}에 의해 표준 사상
$X^\nu \to X_{red}$를 얻고, $X^\nu$는 $X_{red}$의 $Y$ 안에서의
정규화이다. $Y \to X_{red}$가 $X_{red}$에 대해 구성한
(\ref{morphisms-equation-generic-points})의 사상과 동일하므로 보조정리가 따른다.
\end{proof}

\noindent
$X$가 축약이면 $X$의 정규화는 $\mathcal{O}_X$의 유리형 함수층
$\mathcal{K}_X$ 안에서의 정수적 폐포의 상대 스펙트럼과 같다(인자, 절
\ref{divisors-section-meromorphic-functions}를 보라). 실제로 이 경우
$\mathcal{K}_X = f_*\mathcal{O}_Y$이다. 인자, 보조정리
\ref{divisors-lemma-reduced-finite-irreducible}와 그 증명을 보라.
이를 여기서 명시적으로 기술한다.

\begin{lemma}
\label{morphisms-lemma-description-normalization}
각 준콤팩트 열린집합이 유한 개의 기약 성분을 갖는 축약 스킴 $X$를
생각하자. $\Spec(A) = U \subset X$를 아핀 열린집합이라 하자. 그러면
\begin{enumerate}
\item $A$는 유한 개의 극소 소아이디얼
$\mathfrak q_1, \ldots, \mathfrak q_t$를 가진다.
\item 전분수환 $Q(A)$, 즉 $A$의 전분수환은
$Q(A/\mathfrak q_1) \times \ldots \times Q(A/\mathfrak q_t)$이다.
\item 정수적 폐포 $A'$, 즉 $A$의 $Q(A)$ 안에서의 정수적 폐포는 정역
$A/\mathfrak q_i$의 체 $Q(A/\mathfrak q_i)$ 안에서의 정수적 폐포들의
곱이다.
\item $\nu^{-1}(U)$는 $A'$의 스펙트럼과 동일시되며, 여기서
$\nu : X^\nu \to X$는 정규화 사상이다.
\end{enumerate}
\end{lemma}

\begin{proof}
극소 소아이디얼들은 기약 성분들과 대응하므로(대수학, 보조정리
\ref{algebra-lemma-irreducible}), 가정에 의해 (1)을 얻는다.
$(0) = \mathfrak q_1 \cap \ldots \cap \mathfrak q_t$이다. 이는 $A$가
축약이기 때문이다(대수학,
보조정리 \ref{algebra-lemma-Zariski-topology}). 그러면 대수학, 보조정리
\ref{algebra-lemma-total-ring-fractions-no-embedded-points}와
\ref{algebra-lemma-minimal-prime-reduced-ring}에 의해
$Q(A) = \prod A_{\mathfrak q_i} = \prod \kappa(\mathfrak q_i)$를 얻는다.
이는 (2)를 증명한다. (3)은 대수학, 보조정리
\ref{algebra-lemma-characterize-reduced-ring-normal} 또는 보조정리
\ref{morphisms-lemma-normalization-in-disjoint-union}에서 따른다. (4)는
$f^{-1}(U) \to U$가 다음 사상임이 명백하므로 성립한다. 여기서
$f : Y \to X$는 (\ref{morphisms-equation-generic-points})의 사상이다.
$$
\Spec\left(\prod \kappa(\mathfrak q_i)\right)
\longrightarrow
\Spec(A)
$$
\end{proof}

\begin{lemma}
\label{morphisms-lemma-stalk-normalization}
각 준콤팩트 열린집합이 유한 개의 기약 성분을 갖는 스킴 $X$를
생각하자. $\nu : X^\nu \to X$를 $X$의 정규화라 하고 $x \in X$라
하자. 다음은 $\mathcal{O}_{X, x}$-대수로서 표준적으로 동형이다.
\begin{enumerate}
\item 줄기 $(\nu_*\mathcal{O}_{X^\nu})_x$,
\item $\mathcal{O}_{X, x}$의, $(\mathcal{O}_{X, x})_{red}$의
전분수환 안에서의 정수적 폐포,
\item $\mathcal{O}_{X, x}$의 극소 소아이디얼들의 잉여체의 곱 안에서
$\mathcal{O}_{X, x}$의 정수적 폐포(이러한 극소 소아이디얼은 유한 개이다).
\end{enumerate}
\end{lemma}

\begin{proof}
$X$를 $x$의 아핀 열린 근방으로 대체한 뒤, $X$가 유한 개의 기약
성분을 갖고 $x$가 그 모두에 포함된다고 가정해도 된다. 그러면 줄기
$(\nu_*\mathcal{O}_{X^\nu})_x$는 $A = \mathcal{O}_{X, x}$의, 곱
$L$, 즉 $A$의 극소 소아이디얼들의 잉여체의 곱 안에서의 정수적
폐포이다. 이는 정규화의
구성과 보조정리 \ref{morphisms-lemma-integral-closure}에서 따른다. 또는 보조정리
\ref{morphisms-lemma-description-normalization} 및 정규화가 국소화와 가환한다는
사실(대수학, 보조정리
\ref{algebra-lemma-integral-closure-localize})을 사용할 수 있다.
$A_{red}$는 유한 개의 극소 소아이디얼을 가지므로(이들은 정확히
$X$에서 $x$를 지나는 기약 성분들의 일반점들과 대응한다), $L$은
$A_{red}$의 전분수환이다(대수학, 보조정리
\ref{algebra-lemma-total-ring-fractions-no-embedded-points}). 따라서
우리의 환은 $A$의, $A_{red}$의 전분수환 안에서의 정수적 폐포이기도 하다.
\end{proof}

\begin{lemma}
\label{morphisms-lemma-normalization-normal}
각 준콤팩트 열린집합이 유한 개의 기약 성분을 갖는 스킴 $X$를
생각하자.
\begin{enumerate}
\item 정규화 $X^\nu$는 정수적 정규 스킴들의 분리합이다.
\item 사상 $\nu : X^\nu \to X$는 정수적이고 전사이며 기약 성분들의
전단사를 유도한다.
\item 정수적 사상 $\alpha : X' \to X$에 대하여, $U \subset X$가
준콤팩트 열린집합일 때마다 역상 $\alpha^{-1}(U)$가 유한 개의 기약
성분을 가지며 $\alpha|_{\alpha^{-1}(U)} : \alpha^{-1}(U) \to U$가
쌍유리라고 가정하자.\footnote{원천과 대상이 유한
개의 기약 성분을 가질 때에만 사상이 쌍유리라는 의미를 정의했으므로
이 어색한 표현이 필요하다. $X'_{red} \to X_{red}$가 이 조건을
만족하면 충분하다.} 그러면 분해 $X^\nu \to X' \to X$가 존재하고,
$X^\nu \to X'$는 $X'$의 정규화이다.
\item 모든 사상 $Z \to X$에 대하여, $Z$가 정규 스킴이고 $Z$의 각
기약 성분이 $X$의 한 기약 성분을 우세하게 사상하면 유일한 분해
$Z \to X^\nu \to X$가 존재한다.
\end{enumerate}
\end{lemma}

\begin{proof}
$f : Y \to X$를 (\ref{morphisms-equation-generic-points})에서와 같이 두자. 스킴
$X^\nu$는 정수적 정규 스킴들의 분리합이다. 이는 $Y$가 정규이고
$Y$의 모든 아핀 열린집합이 유한 개의 기약 성분을 가진다는 사실과
보조정리 \ref{morphisms-lemma-normal-normalization}에서 따른다. 따라서 (1)이
증명된다. 또는 보조정리
\ref{morphisms-lemma-normalization-reduced}와
\ref{morphisms-lemma-description-normalization}에서 (1)을 도출할 수도 있다.

\medskip\noindent
사상 $\nu$는 보조정리 \ref{morphisms-lemma-characterize-normalization}에 의해
정수적이다. 보조정리 \ref{morphisms-lemma-normal-normalization}에 의해 사상
$Y \to X^\nu$는 기약 성분들의 전단사를 유도하며, $Y$의 구성에 의해
이는 $X^\nu \to X$가 기약 성분들의 전단사를 유도함을 뜻한다. 구성상
$f : Y \to X$는 우세하므로 $\nu$도 우세하다. 정수적 사상은 닫힌
사상이므로(보조정리 \ref{morphisms-lemma-integral-universally-closed}), 따라서
$\nu$는 전사이다. 이는 (2)를 증명한다.

\medskip\noindent
$\alpha : X' \to X$가 (3)에서와 같다고 하자. $X'$가 정규화가
정의되는 가정을 만족함은 명백하다. $f' : Y' \to X'$를 $X'$에서
시작하여 구성한 (\ref{morphisms-equation-generic-points})의 사상이라 하자.
$\alpha$가 국소적으로 쌍유리이므로 $Y' = Y$이고
$f = \alpha \circ f'$임이 명백하다. 따라서 분해
$X^\nu \to X' \to X$가 존재하고, 보조정리
\ref{morphisms-lemma-characterize-normalization}에 의해 $X^\nu \to X'$는
$X'$의 정규화이다. 이는 (3)을 증명한다.

\medskip\noindent
$g : Z \to X$를 원천이 정규 스킴이고 그 모든 기약 성분이 $X$의 한
기약 성분을 우세하게 사상하는 사상이라 하자. 보조정리
\ref{morphisms-lemma-normalization-reduced}에 의해 $X^\nu = X_{red}^\nu$이고,
스킴, 보조정리 \ref{schemes-lemma-map-into-reduction}에 의해
$Z \to X$는 $X_{red}$를 통해 분해된다. 따라서 $X$를 $X_{red}$로
대체하여 $X$가 축약이라고 가정해도 된다. 더욱이 분해는 유일하므로
$Z$ 위에서 국소적으로 이를 구성하면 충분하다. $W \subset Z$와
$U \subset X$를 $g(W) \subset U$인 아핀 열린집합들이라 하자.
$U = \Spec(A)$와 $W = \Spec(B)$로 쓰고, $g|_W$는
$\varphi : A \to B$로 주어진다고 하자. 보조정리
\ref{morphisms-lemma-description-normalization}의 결과들을 자유롭게 사용하겠다.
$\mathfrak p_1, \ldots, \mathfrak p_t$를 $A$의 극소 소아이디얼들이라
하자. $Z$가 정규이므로 $B$는 정규환이고, 특히 축약이다. 더욱이 가정에
의해 모든 극소 소아이디얼 $\mathfrak q \subset B$에 대해
$\varphi^{-1}(\mathfrak q)$는 $A$의 극소 소아이디얼이다. 따라서
$x \in A$가 영인자가 아닌 원소, 즉
$x \not \in \bigcup \mathfrak p_i$이면 $\varphi(x)$는 $B$에서
영인자가 아닌 원소이다. 따라서 표준 환 사상 $Q(A) \to Q(B)$를 얻는다.
$B$가 정규이므로 $Q(B)$ 안에서 자신의 정수적 폐포와 같다(대수학,
보조정리 \ref{algebra-lemma-normal-ring-integrally-closed}를 보라).
그러므로 정수적 폐포 $A' \subset Q(A)$, 즉 $A$의 정수적 폐포는
$B$로 사상되며, 이는 표준 사상 $Q(A) \to Q(B)$에 의한다.
$\nu^{-1}(U) = \Spec(A')$이므로 이는 표준 분해
$W \to \nu^{-1}(U) \to U$, 즉 $\nu|_W$의 분해를 준다.
그 유일성의 확인은 생략한다.
\end{proof}

\begin{lemma}
\label{morphisms-lemma-normalization-in-terms-of-components}
각 준콤팩트 열린집합이 유한 개의 기약 성분을 갖는 스킴 $X$를
생각하자. $Z_i \subset X$, $i \in I$를 축약 유도 구조를 부여한
$X$의 기약 성분들이라 하자. $Z_i^\nu \to Z_i$를 정규화라 하자.
그러면 $\coprod_{i \in I} Z_i^\nu \to X$는 $X$의 정규화이다.
\end{lemma}

\begin{proof}
보조정리 \ref{morphisms-lemma-normalization-reduced}에 의해 $X$가 축약이라고
가정해도 된다. 그러면 보조정리는 보조정리
\ref{morphisms-lemma-description-normalization}의 국소적 기술에서 따르거나,
$\coprod Z_i \to X$가 정수적이고 국소적으로 쌍유리이므로($X$가
축약이고 국소적으로 유한 개의 기약 성분을 가지기 때문이다) 보조정리
\ref{morphisms-lemma-normalization-normal}의 (3)에서 따른다.
\end{proof}

\begin{lemma}
\label{morphisms-lemma-normalization-birational}
유한 개의 기약 성분을 갖는 축약 스킴 $X$를 생각하자. 그러면 정규화
사상 $X^\nu \to X$는 쌍유리이다.
\end{lemma}

\begin{proof}
보조정리 \ref{morphisms-lemma-normalization-normal}에 의해 정규화는 기약 성분들의
전단사를 유도한다. $\eta \in X$를 $X$의 한 기약 성분의 일반점이라
하고, $\eta^\nu \in X^\nu$를 $X^\nu$의 대응하는 기약 성분의
일반점이라 하자. 그러면 $\eta^\nu \mapsto \eta$이고, 증명을 끝내려면
정의 \ref{morphisms-definition-birational}에 따라
$\mathcal{O}_{X, \eta} \to \mathcal{O}_{X^\nu, \eta^\nu}$가 동형임을
보여야 한다. $X$와 $X^\nu$가 축약이므로 두 국소환 모두 자신의
잉여체와 같다(대수학, 보조정리
\ref{algebra-lemma-minimal-prime-reduced-ring}). 한편 정규화를 $X$의
$Y = \coprod \Spec(\kappa(\eta))$ 안에서의 정규화로 구성했으므로
$\kappa(\eta) \subset \kappa(\eta^\nu) \subset \kappa(\eta)$를 얻고,
증명이 끝난다.
\end{proof}

\begin{lemma}
\label{morphisms-lemma-finite-birational-over-normal}
정수적 스킴들의 유한(또는 단지 정수적인) 쌍유리 사상 $f : X \to Y$에서
$Y$가 정규이면 이 사상은 동형사상이다.
\end{lemma}

\begin{proof}
$V \subset Y$를 아핀 열린집합이라 하고, 그 역상 $U \subset X$도
아핀 열린집합이라고 하자. $f$는 정수적 스킴들의 쌍유리 사상이므로
준동형 $\mathcal{O}_Y(V) \to \mathcal{O}_X(U)$는 분수체들의 동형을
유도하는 정역들의 단사 사상이다. $Y$가 정규이므로 환
$\mathcal{O}_Y(V)$는 그 분수체 안에서 정수적으로 닫혀 있다. $f$가
유한(또는 정수적)이므로 $\mathcal{O}_X(U)$의 모든 원소는
$\mathcal{O}_Y(V)$ 위에서 정수적이다. 따라서
$\mathcal{O}_Y(V) = \mathcal{O}_X(U)$를 얻는다. 이는 원하는 대로
$f$가 동형사상임을 증명한다.
\end{proof}

\begin{lemma}
\label{morphisms-lemma-normalization-trivial}
국소적으로 유한 개의 기약 성분을 갖는 스킴 $X$를 생각하자. 정규화
사상 $\nu : X^\nu \to X$가 동형사상인 것은 $X$가 정규인 것과
동치이다.
\end{lemma}

\begin{proof}
$\nu$가 동형사상이면 보조정리 \ref{morphisms-lemma-normalization-normal}에 의해
$X$는 정규이다. 반대로 $X$가 정규라고 가정하자. 보조정리
\ref{morphisms-lemma-normalization-localization}과 성질, 보조정리
\ref{properties-lemma-normal-locally-finite-nr-irreducibles}에 의해
$X$가 정수적이라고 가정해도 된다. 보조정리
\ref{morphisms-lemma-normalization-normal}에 의해 사상 $\nu$는 정수적이고
$X^\nu$는 유일한 기약 성분을 가진다(따라서 정수적 스킴이다).
보조정리 \ref{morphisms-lemma-normalization-birational}와
\ref{morphisms-lemma-finite-birational-over-normal}에서 결론을 얻는다.
\end{proof}

\begin{lemma}
\label{morphisms-lemma-Japanese-normalization}
$X$를 정수적이고 일본적인 스킴이라 하자. 정규화
$\nu : X^\nu \to X$는 유한 사상이다.
\end{lemma}

\begin{proof}
이는 정의(성질, 정의 \ref{properties-definition-nagata})와 보조정리
\ref{morphisms-lemma-description-normalization}에서 따른다. 즉 이 경우 보조정리는
$\nu^{-1}(\Spec(A))$가 $A$의 분수체 안에서 취한 그 환의 정수적
폐포의 스펙트럼이라는 뜻이다.
\end{proof}

\begin{lemma}
\label{morphisms-lemma-nagata-normalization}
$X$를 나가타 스킴이라 하자. 정규화 $\nu : X^\nu \to X$는 유한
사상이다.
\end{lemma}

\begin{proof}
나가타 스킴은 국소 뇌터이므로 정의 \ref{morphisms-definition-normalization}이
적용됨에 유의하라. 이제 보조정리는 보조정리
\ref{morphisms-lemma-nagata-normalization-finite-general}의 특수한 경우이지만,
다음과 같이 직접 증명할 수도 있다. $X^\nu \to X$를 합성
$X^\nu \to X_{red} \to X$로 쓰자. $X_{red} \to X$는 닫힌 몰입이므로
유한이다. 따라서 보조정리 \ref{morphisms-lemma-composition-finite}에 의해 축약
나가타 스킴의 경우에 보조정리를 증명하면 충분하다. 아핀 열린집합
$\Spec(A) = U \subset X$를 택하자. 보조정리
\ref{morphisms-lemma-description-normalization}에 의해
$\nu^{-1}(U) = \Spec(\prod A_i')$이다. 여기서 $A_i'$는
$A/\mathfrak q_i$의 분수체 안에서의 정수적 폐포이다. $A$는 나가타환이므로
(성질, 보조정리 \ref{properties-lemma-locally-nagata}를 보라), 각 환확장
$A/\mathfrak q_i \subset A'_i$는 유한이다. 따라서
$A \to \prod A'_i$는 유한 환 사상이고 결론을 얻는다.
\end{proof}

\begin{lemma}
\label{morphisms-lemma-normalization-geom-unibranch-univ-homeo}
$X$를 기약이고 기하적으로 단분지인 스킴이라 하자. 정규화 사상
$\nu : X^\nu \to X$는 보편 위상동형사상이다.
\end{lemma}

\begin{proof}
보조정리 \ref{morphisms-lemma-universal-homeomorphism}에 따라 $\nu$가 정수적이고,
보편 단사이며, 전사임을 보여야 한다. 보조정리
\ref{morphisms-lemma-normalization-normal}에 의해 사상 $\nu$는 정수적이다.
$x \in X$라 하고 $A = \mathcal{O}_{X, x}$로 두자. $X$가 기약이므로
$A$는 유일한 극소 소아이디얼 $\mathfrak p$를 가지며
$A_{red} = A/\mathfrak p$이다. 보조정리
\ref{morphisms-lemma-stalk-normalization}에 의해 줄기
$A' = (\nu_*\mathcal{O}_{X^\nu})_x$는 $A$의, $A_{red}$의 분수체
안에서의 정수적 폐포이다. 대수 추가, 정의
\ref{more-algebra-definition-unibranch}에 의해 $A'$는 유일한 소아이디얼
$\mathfrak m'$를 가지며, 이는 $\mathfrak m_x \subset A$ 위에 놓이고
$\kappa(\mathfrak m')/\kappa(x)$는 순수
비분리이다. 따라서 $\nu$는 전단사(그러므로 전사)이고, 보조정리
\ref{morphisms-lemma-universally-injective}에 의해 보편 단사이다.
\end{proof}







\section{약정규화}
\label{morphisms-section-weak-normalization}

\noindent
정규화의 경우와 마찬가지로, 스킴이 국소적으로 유한 개의 기약 성분을
가질 때에만 그 약정규화를 정의하겠다.

\begin{lemma}
\label{morphisms-lemma-relative-weak-normalization-algebra}
환 사상 $A \to B$가 스펙트럼의 우세 사상
$\Spec(B) \to \Spec(A)$를 유도한다고 하자. 다음을 만족하는
$A$-부분대수 $B' \subset B$가 존재한다.
\begin{enumerate}
\item $\Spec(B') \to \Spec(A)$는 보편 위상동형사상이다.
\item 분해 $A \to C \to B$가 주어지고
$\Spec(C) \to \Spec(A)$가 보편 위상동형사상이면, $C \to B$의 상은
$B'$에 포함된다.
\end{enumerate}
\end{lemma}

\begin{proof}
보조정리 \ref{morphisms-lemma-reduction-universal-homeomorphism}을 더 언급하지 않고
사용하겠다. 다음 가환 그림을 생각하자.
$$
\xymatrix{
B \ar[r] & B_{red} \\
A \ar[u] \ar[r] & A_{red} \ar[u]
}
$$
(2)에서와 같은 임의의 분해 $A \to C \to B$, 즉 $A \to B$의 분해에 대하여
$A_{red} \to C_{red} \to B_{red}$는 (2)에서와 같은
$A_{red} \to B_{red}$의 분해임을 알 수 있다. 따라서 보조정리가
$A_{red} \to B_{red}$에 대해 성립하여 $A_{red}$-부분대수
$B'_{red} \subset B_{red}$를 산출한다면, $B' \subset B$를
$B'_{red}$의 역상과 같게 두는 것으로 $A \to B$에 대한 보조정리를
해결할 수 있다. 이로써 다음 문단에서 다루는 경우로 환원된다.

\medskip\noindent
$A$와 $B$가 축약이라고 가정하자. 이 경우 대수학, 보조정리
\ref{algebra-lemma-image-dense-generic-points}에 의해 $A \subset B$이다.
$A \to C \to B$를 (2)에서와 같은 분해라 하자. 그러면
$A \subset C$에 명제 \ref{morphisms-proposition-universal-homeomorphism}을 적용하여
$C$의 모든 원소가 다음과 같은 확대 $A[c_1, \ldots, c_n] \subset C$에
포함됨을 알 수 있다. $i = 1, \ldots, n$에 대해
\begin{enumerate}
\item $c_i^2, c_i^3 \in A[c_1, \ldots, c_{i - 1}]$이거나,
\item 소수 $p$가 존재하여
$pc_i, c_i^p \in A[c_1, \ldots, c_{i - 1}]$이다.
\end{enumerate}
따라서 $B' \subset B$를 다음 원소들의 부분집합으로 정의하면 성질
(2)가 성립한다. 즉 $b \in B$가 다음 확대에 포함되어야 한다.
$
A[b_1, \ldots, b_n] \subset B
$
여기서 (*)는 다음 조건을 뜻한다. $i = 1, \ldots, n$에 대해
\begin{enumerate}
\item $b_i^2, b_i^3 \in A[b_1, \ldots, b_{i - 1}]$이거나,
\item 소수 $p$가 존재하여
$pb_i, b_i^p \in A[b_1, \ldots, b_{i - 1}]$이다.
\end{enumerate}
확인할 것은 두 가지뿐이다. (a) $B'$가 $A$-부분대수라는 것과
(b) $\Spec(B') \to \Spec(A)$가 보편 위상동형사상이라는 것이다.
(a)는 다음 사실에서 따른다. (*)를 만족하는 $n \geq 0$과
$b_1, \ldots, b_n \in B$, 그리고 (*)를 만족하는 $m \geq 0$과
$b'_1, \ldots, b'_m \in B$가 주어지면, 정수 $n + m$ 및
$b_1, \ldots, b_n, b'_1, \ldots, b'_m \in B$도 (*)를 만족한다.
마지막으로 (b)는 명제 \ref{morphisms-proposition-universal-homeomorphism}과
$B'$의 구성에서 따른다.
\end{proof}

\begin{lemma}
\label{morphisms-lemma-relative-weak-normalization-localization}
환 사상 $A \to B$가 스펙트럼의 우세 사상
$\Spec(B) \to \Spec(A)$를 유도한다고 하자. 보조정리
\ref{morphisms-lemma-relative-weak-normalization-algebra}의 $A$-부분대수
$B' \subset B$의 구성은 국소화와 가환한다(설명은 증명을 보라).
\end{lemma}

\begin{proof}
$S \subset A$를 곱셈적 부분집합이라 하자. 그러면
$S^{-1}A \to S^{-1}B$ 역시 우세 사상
$\Spec(S^{-1}B) \to \Spec(S^{-1}A)$를 유도하는 환 사상이다(보조정리
\ref{morphisms-lemma-base-change-dominant}와
\ref{morphisms-lemma-generalizations-lift-flat}을 보라). 따라서 보조정리
\ref{morphisms-lemma-relative-weak-normalization-algebra}은 $S^{-1}A$-부분대수
$(S^{-1}B)' \subset S^{-1}B$를 산출한다. 명제의 의미는
$S^{-1}B' = (S^{-1}B)'$가 $S^{-1}A$-부분대수로서 $S^{-1}B$ 안에서
성립한다는 것이다.

\medskip\noindent
이를 확인하기 위해 보조정리
\ref{morphisms-lemma-relative-weak-normalization-algebra}의 증명에서 사용한
$B'$와 $(S^{-1}B)'$의 구성을 사용하겠다. 첫 단계에서 $B'$는
$A_{red}$-부분대수 $B'_{red} \subset B_{red}$의 역상인데, 이는 환 사상
$A_{red} \to B_{red}$에 대해 구성한 것이며, $(S^{-1}B)'$도 마찬가지이다.
$S^{-1}B_{red} = (S^{-1}B)_{red}$임에 유의하면 다음 문단의 경우로
환원된다.

\medskip\noindent
$A$와 $B$가 축약이면, 우리는 $B'$를 부분대수
$A[b_1, \ldots, b_n]$들의 합집합으로 구성했다. 여기서
$i = 1, \ldots, n$에 대해
\begin{enumerate}
\item $b_i^2, b_i^3 \in A[b_1, \ldots, b_{i - 1}]$이거나,
\item 소수 $p$가 존재하여
$pb_i, b_i^p \in A[b_1, \ldots, b_{i - 1}]$이다.
\end{enumerate}
$(S^{-1}B)' \subset S^{-1}B$에 대해서도 마찬가지이다. 따라서
$B' \to B \to S^{-1}B$의 상이 $(S^{-1}B)'$에 포함됨은 명백하다.
대응하는 사상 $S^{-1}B' \to (S^{-1}B)'$가 전사임을 보이려면 보조정리
\ref{morphisms-lemma-special-elements-and-localization}을 사용하여 분모를 차례로
없애면 된다. 세부사항은 생략한다.
\end{proof}

\begin{lemma}
\label{morphisms-lemma-relative-seminormalization-algebra}
환 사상 $A \to B$가 스펙트럼의 우세 사상
$\Spec(B) \to \Spec(A)$를 유도한다고 하자. 다음을 만족하는
$A$-부분대수 $B' \subset B$가 존재한다.
\begin{enumerate}
\item $\Spec(B') \to \Spec(A)$는 잉여체들 위의 동형을 유도하는
보편 위상동형사상이다.
\item 분해 $A \to C \to B$가 주어지고 $\Spec(C) \to \Spec(A)$가
잉여체들 위의 동형을 유도하는 보편 위상동형사상이면,
$C \to B$의 상은 $B'$에 포함된다.
\end{enumerate}
\end{lemma}

\begin{proof}
이 증명은 보조정리 \ref{morphisms-lemma-relative-weak-normalization-algebra}의
증명과 정확히 같지만, 명제
\ref{morphisms-proposition-universal-homeomorphism-equal-residue-fields}를 명제
\ref{morphisms-proposition-universal-homeomorphism} 대신 사용한다.
\end{proof}

\begin{lemma}
\label{morphisms-lemma-relative-seminormalization-localization}
환 사상 $A \to B$가 스펙트럼의 우세 사상
$\Spec(B) \to \Spec(A)$를 유도한다고 하자. 보조정리
\ref{morphisms-lemma-relative-seminormalization-algebra}의 $A$-부분대수
$B' \subset B$의 구성은 국소화와 가환한다(설명은 증명을 보라).
\end{lemma}

\begin{proof}
증명은 보조정리 \ref{morphisms-lemma-relative-weak-normalization-localization}의
증명과 같다.
\end{proof}

\begin{lemma}
\label{morphisms-lemma-relative-seminormalization}
$f : Y \to X$를 스킴들의 준콤팩트하고 준분리이며 우세한 사상이라 하자.
\begin{enumerate}
\item 분해 $Y \to X' \to X$에서 $X' \to X$가 보편 위상동형사상인
그러한 분해들의 범주는 시작 대상 $Y \to X^{Y/wn} \to X$를 가진다.
\item 분해 $Y \to X' \to X$에서 $X' \to X$가 잉여체들 위의 동형을
유도하는 보편 위상동형사상인 그러한 분해들의 범주는 시작 대상
$Y \to X^{Y/sn} \to X$를 가진다.
\end{enumerate}
더욱이 분해 $Y \to X^{Y/wn} \to X$와
$Y \to X^{Y/sn} \to X$의 구성은 $X$의 열린 부분스킴으로의
기저변환과 가환한다.
\end{lemma}

\begin{proof}
(1)을 증명하고 (2)의 증명은 생략하겠다. 마지막 주장도 분해의 구성에서
따를 것이다. 보조정리 \ref{morphisms-lemma-universal-homeomorphism}을 더 언급하지
않고 사용하겠다. 먼저 $Y \to X^{Y/n} \to X$를 $X$의 $Y$ 안에서의
정규화라 하자. 정의 \ref{morphisms-definition-normalization-X-in-Y}를 보라.
(1)에서와 같은 $Y \to X' \to X$에 대해 주어진 사상들과 가환하는
유일한 사상 $X^{Y/n} \to X'$를 얻는다. 보조정리
\ref{morphisms-lemma-characterize-normalization}을 보라. 따라서 $f$를
$X^{Y/n} \to X$로 대체한 경우의 보조정리를 증명하면 충분하다.
이로써 다음 문단에서 연구하는 경우로 환원된다.

\medskip\noindent
$f$가 정수적이라고 가정하자(증명의 나머지는 $f$가 아핀인 경우에도
성립한다). $U = \Spec(A)$를 $X$의 아핀 열린집합이라 하고,
$V = f^{-1}(U) = \Spec(B)$를 $Y$에서의 역상이라 하자. 그러면
$A \to B$는 스펙트럼 위의 우세 사상을 유도하는 환 사상이다. 보조정리
\ref{morphisms-lemma-relative-weak-normalization-algebra}에 의해 $A$-부분대수
$B' \subset B$를 얻는다. $U^{V/wn} = \Spec(B')$로 두면 분해
$V \to U^{V/wn} \to U$는 분해 $V \to U' \to U$ 중에서
$U' \to U$가 보편 위상동형사상인 것들의 범주의 시작 대상이다.

\medskip\noindent
$U_1 \subset U_2 \subset X$를 아핀 열린집합들이라 하자.
$V_i = f^{-1}(U_i)$로 두면 표준 사상
$$
\rho_{U_1}^{U_2} : U_1^{V_1/wn} \to U_1 \times_{U_2} U_2^{V_2/wn}
$$
을 $U_1$ 위에서 $U_1^{V_1/wn}$의 보편 성질에 의해 얻는다. 이 사상들은
자연스러운 함자성을 만족하며, 그 정식화와 증명은 독자에게 맡긴다.
더욱이 사상 $\rho_{U_1}^{U_2}$는 동형사상이다. 이는
$U_1 \subset U_2$가 표준 열린집합일 때에는 보조정리
\ref{morphisms-lemma-relative-weak-normalization-localization}에서 따르고,
일반적인 경우에는 이 사상들의 함자성과 스킴, 보조정리
\ref{schemes-lemma-standard-open-two-affines}를 사용하여 그 경우로
환원된다(세부사항은 생략한다). 따라서 상대 붙이기(구성, 보조정리
\ref{constructions-lemma-relative-glueing})에 의해 사상
$X^{Y/wn} \to X$를 얻는데, 이는 $U^{V/wn} \to U$로 제한되며 그 제한은
$U$ 위에서 $\rho_{U_1}^{U_2}$들과 가환한다. 물론 사상
$V \to U^{V/wn}$들은 사상 $Y \to X^{Y/wn}$으로 붙는다(구성, 주석
\ref{constructions-remark-relative-glueing-functorial}을 보라). 이로써
두 번째 사상이 보편 위상동형사상인 분해
$Y \to X^{Y/wn} \to X$를 얻는다.

\medskip\noindent
마지막으로 $Y \to X' \to X$를 (1)에서와 같은 분해라 하자. 위와 같이
$V \to U^{V/wn} \to U \subset X$가 주어지면 두 번째 화살표가 보편
위상동형사상인 분해 $V \to U \times_X X' \to U$를 얻고, 유일한 사상
$g_U : U^{V/wn} \to U \times_X X'$를 얻는다. 이는 $U$ 위의 사상이며
$U^{V/wn}$의 보편 성질에서 나온다. 이 $g_U$들은 사상
$\rho_{U_1}^{U_2}$들과 가환한다. 세부사항은 생략한다. 따라서
유일한 사상 $g : X^{Y/wn} \to X'$가 존재한다. 이는 $X$ 위에서
$g_U$와 $U$ 위에서 일치한다. 구성, 주석
\ref{constructions-remark-relative-glueing-functorial}을 보라. 이는
$Y \to X^{Y/wn} \to X$가 우리 범주의 시작 대상임을 증명하며,
증명이 끝난다.
\end{proof}

\begin{definition}
\label{morphisms-definition-relative-seminormalization}
$f : Y \to X$를 스킴들의 준콤팩트하고 준분리이며 우세한 사상이라 하자.
\begin{enumerate}
\item 보조정리 \ref{morphisms-lemma-relative-seminormalization}의 (2)에서 구성한
분해 $Y \to X^{Y/sn} \to X$를 {\it $X$의 $Y$ 안에서의 반정규화}라 한다.
\item 보조정리 \ref{morphisms-lemma-relative-seminormalization}의 (1)에서 구성한
분해 $Y \to X^{Y/wn} \to X$를 {\it $X$의 $Y$ 안에서의 약정규화}라 한다.
\end{enumerate}
\end{definition}

\noindent
국소적으로 유한 개의 기약 성분을 갖는 스킴의 반정규화를 다시
해석하는 한 방법을 제시한다.

\begin{lemma}
\label{morphisms-lemma-seminormal-and-normalization}
각 준콤팩트 열린집합이 유한 개의 기약 성분을 갖는 스킴 $X$를
생각하자. $\nu : X^\nu \to X$를 $X$의 정규화라 하자. 그러면
$X$의 $X^\nu$ 안에서의 반정규화는 $X$의 반정규화이다. 식으로 쓰면
$X^{sn} = X^{X^\nu/sn}$이다.
\end{lemma}

\begin{proof}
$f : Y \to X$를 (\ref{morphisms-equation-generic-points})에서와 같이 두어
$X^\nu$가 $X$의 $Y$ 안에서의 정규화가 되게 하자. 반정규화
$X^{sn} \to X$, 즉 $X$의 반정규화는 잉여체들 위의 동형을 유도하는 보편 위상동형사상
$X' \to X$들의 범주의 시작 대상이다. $Y$는 $X$의 기약 성분들의
일반점에서의 잉여체들의 스펙트럼들의 분리합이므로, 이 범주의 모든
$X' \to X$에 대해 표준 올림 $f' : Y \to X'$, 즉 $f$의 올림을 얻는다. 그러면
보조정리 \ref{morphisms-lemma-characterize-normalization}에 의해 표준 사상
$X^\nu \to X'$를 얻는다. 이어서 표준 사상
$X^{X^\nu/sn} \to X'$를 얻는데, 이는 $X^{X^\nu/sn}$의 보편 성질에서
나온다. 이로써
$X^{X^\nu/sn}$가 $X^{sn}$과 같은 보편 성질을 만족함을 알 수 있고,
결론이 따른다.
\end{proof}

\noindent
보조정리 \ref{morphisms-lemma-seminormal-and-normalization}는 다음 정의의 동기를
준다. $X$가 국소적으로 유한 개의 기약 성분을 갖는 경우에만 정규화를
구성했으므로, 약정규화에 대해서도 이 경우로 제한하겠다.

\begin{definition}
\label{morphisms-definition-weak-normalization}
각 준콤팩트 열린집합이 유한 개의 기약 성분을 갖는 스킴 $X$를
생각하자. $X$의 {\it 약정규화}를 다음과 같은 약정규화로 정의한다.
$$
X^\nu \longrightarrow X^{wn} \longrightarrow X
$$
이는 $X$의 정규화 $X^\nu$ 안에서의 $X$의 약정규화이다
(정의 \ref{morphisms-definition-normalization}). 식으로 쓰면
$X^{wn} = X^{X^\nu/wn}$이다.
\end{definition}

\noindent
보조정리 \ref{morphisms-lemma-seminormal-and-normalization}와 결합하면, 국소적으로
유한 개의 기약 성분을 갖는 스킴 $X$에 대해 표준 사상들이 존재함을
알 수 있다.
$$
X^\nu \to X^{wn} \to X^{sn} \to X
$$
이 정의를 했으므로 이제 스킴이 약정규라는 의미를 말할 수 있다(그
스킴이 국소적으로 유한 개의 기약 성분을 갖는다는 전제 아래에서이다).

\begin{definition}
\label{morphisms-definition-weakly-normal}
각 준콤팩트 열린집합이 유한 개의 기약 성분을 갖는 스킴 $X$를
생각하자. $X$가 {\it 약정규}라는 것은 약정규화
$X^{wn} \to X$가 동형사상이라는 뜻이다
(정의 \ref{morphisms-definition-weak-normalization}).
\end{definition}

\noindent
정의들로부터, 각 준콤팩트 열린집합이 유한 개의 기약 성분을 갖는
스킴 $X$에 대해 다음이 즉시 따른다.
$$
X\text{ normal} \Rightarrow
X\text{ weakly normal} \Rightarrow
X\text{ seminormal}
$$
아핀인 경우 약정규성의 의미는 다음과 같이 구체화할 수 있다.

\begin{lemma}
\label{morphisms-lemma-affine-weakly-normal}
유한 개의 기약 성분을 갖는 아핀 스킴 $X = \Spec(A)$를 생각하자.
$X$가 약정규인 것은 다음과 동치이다.
\begin{enumerate}
\item $A$가 반정규이다(정의 \ref{morphisms-definition-seminormal-ring}).
\item 소수 $p$와 $z, w \in A$가 주어져서 (a) $z$가 영인자가 아니고,
(b) $w^p$가 $z^p$로 나누어지며, (c) $pw$가 $z$로 나누어지면,
$w$는 $z$로 나누어진다.
\end{enumerate}
\end{lemma}

\begin{proof}
$X$가 약정규라고 가정하자. 약정규 스킴은 반정규이므로 (약정규 스킴의
정의에 의해) (1)이 성립한다. 특히 $A$는 축약이다. $p, z, w$를
(2)에서와 같이 택하자. $x, y \in A$를 $z^p x = w^p$이고
$zy = pw$가 되도록 택하자. 그러면 $p^p x = y^p$이다. 환 사상
$A \to C = A[t]/(t^p - x, pt - y)$는 스펙트럼 위의 보편 위상동형사상을
유도한다. 정규화 $X^\nu$, 즉 $X$의 정규화는 정수적 폐포 $A'$, 즉
$A$의 전분수환 안에서 취한 $A$의 정수적 폐포의 스펙트럼이다. 보조정리
\ref{morphisms-lemma-description-normalization}을 보라. $a = w/z \in A'$임에
유의하라. 이는 $a^p = x$이기 때문이다. 따라서 $A$-대수 준동형
$A \to C \to A'$를 얻으며, 이는 $t$를 $a$로 보낸다. 이때 약정규라는
정의 성질 $X = X^{wn} = X^{X^\nu/wn}$에 의해 $C \to A'$는 $A$ 안으로
사상된다. 따라서 원하는 대로 $a \in A$를 얻는다.

\medskip\noindent
반대로 (1)과 (2)를 가정하자. $A'$를 앞 문단에서와 같이 두자.
$X^{X^\nu/wn} = X$임을 보여야 한다. 보조정리
\ref{morphisms-lemma-relative-weak-normalization-algebra}의 증명에서의 구성에 의해
스킴 $X^{X^\nu/wn}$은 $A'$의 다음 부분환의 스펙트럼이다. 즉
$A[a_1, \ldots, a_n] \subset A'$ 중 $i = 1, \ldots, n$에 대해
아래 조건을 만족하는 것들의 합집합이다.
\begin{enumerate}
\item[(a)] $a_i^2, a_i^3 \in A[a_1, \ldots, a_{i - 1}]$이거나,
\item[(b)] 소수 $p$가 존재하여
$pa_i, a_i^p \in A[a_1, \ldots, a_{i - 1}]$이다.
\end{enumerate}
그러면 (1)과 (2)를 사용하여 귀납적으로
$a_1, \ldots, a_n \in A$임을 알 수 있다. 세부사항은 생략한다.
따라서 $X = X^{X^\nu/wn}$이고, 그러므로 $X$는 약정규이다.
\end{proof}

\noindent
다음은 빠질 수 없는 보조정리이다.

\begin{lemma}
\label{morphisms-lemma-locally-weakly-normal}
각 준콤팩트 열린집합이 유한 개의 기약 성분을 갖는 스킴 $X$를
생각하자. 다음은 동치이다.
\begin{enumerate}
\item 스킴 $X$는 약정규이다.
\item 모든 아핀 열린집합 $U \subset X$에 대해 환 $\mathcal{O}_X(U)$는
보조정리 \ref{morphisms-lemma-affine-weakly-normal}의 조건 (1)과 (2)를 만족한다.
\item 아핀 열린덮개 $X = \bigcup U_i$가 존재하여 각 환
$\mathcal{O}_X(U_i)$가 보조정리 \ref{morphisms-lemma-affine-weakly-normal}의
조건 (1)과 (2)를 만족한다.
\item 열린덮개 $X = \bigcup X_j$가 존재하여 각 열린 부분스킴
$X_j$가 약정규이다.
\end{enumerate}
더욱이 $X$가 약정규이면 모든 열린 부분스킴도 약정규이다.
\end{lemma}

\begin{proof}
$X$가 약정규일 조건은 사상 $X^{wn} = X^{X^\nu/wn} \to X$가
동형사상인 것이다. $X^\nu \to X$의 구성은 열린 부분스킴으로의
기저변환과 가환하고, $X^{X^\nu/wn}$의 구성은 $X$의 열린 부분스킴으로의
기저변환과 가환하므로(보조정리
\ref{morphisms-lemma-relative-seminormalization}), 보조정리는 명백하다.
\end{proof}









\section{자리스키 주정리(대수적 판)}
\label{morphisms-section-Zariski}

\noindent
이는 순수하게 대수적인 방법으로 증명할 수 있는 판이다. 더 강력한 판들
(아핀이 아닌 사상들에 대한 판들)을 증명하려면 도구를 더 개발해야 한다.
스킴의 코호몰로지, 절
\ref{coherent-section-applications-formal-functions}와 사상의 추가 내용, 절
\ref{more-morphisms-section-application-etale-neighbourhoods}을 보라.

\begin{theorem}[자리스키 주정리의 대수적 판]
\label{morphisms-theorem-main-theorem}
$f : Y \to X$를 스킴들의 아핀 사상이라 하자. $f$가 유한형이라고
가정하자. $X'$를 $X$의 $Y$ 안에서의 정규화라 하자. 그림은 다음과 같다.
$$
\xymatrix{
Y \ar[rd]_f \ar[rr]_{f'} & & X' \ar[ld]^\nu \\
& X &
}
$$
그러면 다음을 만족하는 열린 부분스킴 $U' \subset X'$가 존재한다.
\begin{enumerate}
\item $(f')^{-1}(U') \to U'$는 동형사상이다.
\item $(f')^{-1}(U') \subset Y$는 $f$가 준유한인 점들의 집합이다.
\end{enumerate}
\end{theorem}

\begin{proof}
$X$, 따라서 $Y$가 아핀인 경우로 즉시 환원된다. $X = \Spec(R)$이고
$Y = \Spec(A)$라고 하자. 그러면 $X' = \Spec(A')$이며, 여기서 $A'$는
$R$의 $A$ 안에서의 정수적 폐포이다. 정의
\ref{morphisms-definition-integral-closure}와
\ref{morphisms-definition-normalization-X-in-Y}를 보라. 대수학, 정리
\ref{algebra-theorem-main-theorem}에 의해, 모든 $y \in Y$ 중 $f$가
준유한인 점에 대해 열린집합 $U'_y \subset X'$가 존재하여
$(f')^{-1}(U'_y) \to U'_y$가 동형사상이 된다.
$U' = \bigcup U'_y$로 두되, 여기서 $y \in Y$는 $f$가 준유한인
모든 점을 달린다. 이제 $f$가
$(f')^{-1}(U')$의 모든 점에서 준유한임을 보이면 된다.
$y \in (f')^{-1}(U')$이고 그 상이 $x \in X$이면,
$Y_x \to X'_x$는 $y$의 한 근방에서 동형사상임을 알 수 있다. 따라서
$Y_x$에는 $y$로 특수화되는 점이 없다. 실제로 보조정리
\ref{morphisms-lemma-integral-fibres}에 의해 이는 $f'(y)$에 대해 $X'_x$에서
성립한다. 보조정리 \ref{morphisms-lemma-quasi-finite-at-point-characterize}의
(3)에 의해 이는 $f$가 $y$에서 준유한임을 뜻한다.
\end{proof}

\noindent
자리스키 주정리의 대수적 판을 사용하여 사상이 준유한인 점들의 집합이
열린집합임을 보일 수 있다.

\begin{lemma}
\label{morphisms-lemma-quasi-finite-points-open}
\begin{slogan}
사상의 국소 준유한 자취는 열린집합이다.
\end{slogan}
$f : X \to S$를 스킴들의 사상이라 하자. $X$의 점들 중 $f$가
준유한인 것들의 집합은 열린집합 $U \subset X$이다. 유도된 사상 $U \to S$는
국소 준유한이다.
\end{lemma}

\begin{proof}
$f$가 $x$에서 준유한이라고 가정하자. 정의
\ref{morphisms-definition-quasi-finite}에서와 같은 아핀 열린집합
$x \in U = \Spec(A) \subset X$, $V = \Spec(R) \subset S$를 택하자.
위의 정리 \ref{morphisms-theorem-main-theorem} 또는 대수학, 보조정리
\ref{algebra-lemma-quasi-finite-open}에 의해 소아이디얼들 $\mathfrak q$ 중
$R \to A$가 준유한인 것들의 집합은 $\Spec(A)$에서 열린집합이다.
이들은 모두 $X$의 점들 중 $f$가 준유한인 것들에 대응하므로 첫 번째 주장을
얻는다. 두 번째 주장은 명백하다.
\end{proof}

\noindent
뒤에서 다음 보조정리를 일반적인 준유한 분리 사상들로 개선한다.
사상의 추가 내용, 보조정리
\ref{more-morphisms-lemma-quasi-finite-separated-pass-through-finite}를 보라.

\begin{lemma}
\label{morphisms-lemma-quasi-finite-affine}
$f : Y \to X$를 스킴들의 사상이라 하자. 다음을 가정하자.
\begin{enumerate}
\item $X$와 $Y$는 아핀이다.
\item $f$는 준유한이다.
\end{enumerate}
그러면 다음 그림이 존재한다.
$$
\xymatrix{
Y \ar[rd]_f \ar[rr]_j & & Z \ar[ld]^\pi \\
& X &
}
$$
여기서 $Z$는 아핀이고, $\pi$는 유한이며, $j$는 열린 몰입이다.
\end{lemma}

\begin{proof}
이는 대수학, 보조정리
\ref{algebra-lemma-quasi-finite-open-integral-closure}를 스킴의 언어로
다시 쓴 것이다.
\end{proof}

\begin{lemma}
\label{morphisms-lemma-image-nowhere-dense-quasi-finite}
$f : Y \to X$를 스킴들의 준유한 사상이라 하자. $T \subset Y$를
$Y$의 닫힌 조밀한 곳 없는 부분집합이라 하자. 그러면
$f(T) \subset X$는 $X$의 조밀한 곳 없는 부분집합이다.
\end{lemma}

\begin{proof}
보조정리 \ref{morphisms-lemma-image-nowhere-dense-finite}의 증명에서와 마찬가지로,
밑 $X$가 아핀인 경우로 즉시 환원된다. 이 경우
$Y = \bigcup_{i = 1, \ldots, n} Y_i$는 아핀 열린집합들의 유한 합집합이다
($f$가 준콤팩트이기 때문이다). 각 $T \cap Y_i$가 조밀한 곳 없고,
조밀한 곳 없는 집합들의 유한 합집합도 조밀한 곳 없으므로(위상수학,
보조정리 \ref{topology-lemma-nowhere-dense}를 보라), 상
$f(T \cap Y_i)$가 $X$에서 조밀한 곳 없음을 보이면 충분하다. 이로써
$X$와 $Y$가 모두 아핀인 경우로 환원된다. 이제 위의 보조정리
\ref{morphisms-lemma-quasi-finite-affine}를 적용하여 다음 그림을 얻는다.
$$
\xymatrix{
Y \ar[rd]_f \ar[rr]_j & & Z \ar[ld]^\pi \\
& X &
}
$$
여기서 $Z$는 아핀이고, $\pi$는 유한이며, $j$는 열린 몰입이다.
$\overline{T} = \overline{j(T)} \subset Z$로 두자. 위상수학, 보조정리
\ref{topology-lemma-image-nowhere-dense-open}에 의해 $\overline{T}$는
$Z$에서 조밀한 곳 없다. $f(T) \subset \pi(\overline{T})$이므로 유한인
경우의 대응하는 결과에서 보조정리가 따른다. 보조정리
\ref{morphisms-lemma-image-nowhere-dense-finite}를 보라.
\end{proof}




\section{보편적으로 유계인 올}
\label{morphisms-section-universally-bounded}

\noindent
스킴 $X$를 생각하되 이는 체 $k$ 위의 스킴이라고 하자. $X$가 $k$ 위에서 유한이면
$X = \Spec(A)$이고, 여기서 $A$는 유한 $k$-대수이다. 이를 달리 말하면
$X$가 $\Spec(k)$ 위에서 유한 국소 자유라는 뜻이다. 정의
\ref{morphisms-definition-finite-locally-free}를 보라. 따라서
$X \to \Spec(k)$는 정수 $d \geq 0$인 {\it 차수}를 가지며, 이는
$d = \dim_k(A)$이다. 때때로 이를 (유한) 스킴 $X$의 $k$ 위에서의
{\it 차수}라고 부른다.

\begin{definition}
\label{morphisms-definition-universally-bounded}
$f : X \to Y$를 스킴들의 사상이라 하자.
\begin{enumerate}
\item 정수 $n$이 {\it $f$의 올들의 차수에 대한 상계}라는 것은 모든
$y \in Y$에 대해 올 $X_y$가 $\kappa(y)$ 위의 유한 스킴이고
$\kappa(y)$ 위의 차수가 $\leq n$이라는 뜻이다.
\item $f$의 올들의 차수에 대한 상계가 되는 정수 $n$이 존재하면
{\it $f$의 올들이 보편적으로 유계이다}라고 한다.\footnote{이는 아마
표준적이지 않은 표기법이다.}
\end{enumerate}
\end{definition}

\noindent
특히 한 올 안의 점들의 수도 $n$으로 제약됨에 유의하라. (모든 올이
유한 축약 스킴이어도 그 역은 성립하지 않는다.)

\begin{lemma}
\label{morphisms-lemma-characterize-universally-bounded}
$f : X \to Y$를 스킴들의 사상이라 하고 $n \geq 0$이라 하자. 다음은
동치이다.
\begin{enumerate}
\item 정수 $n$은 $f$의 올들의 차수에 대한 상계이다.
\item 모든 사상 $\Spec(k) \to Y$에 대해, 여기서 $k$는 체이고,
올곱 $X_k = \Spec(k) \times_Y X$는 $k$ 위에서 유한하고 차수가
$\leq n$이다.
\end{enumerate}
이 경우 $f$의 올들은 보편적으로 유계이고 스킴 $X_k$들은 많아야
$n$개의 점을 가진다. 더 정확히, $X_k = \{x_1, \ldots, x_t\}$이면
다음을 얻는다.
$$
n \geq \sum\nolimits_{i = 1, \ldots, t} [\kappa(x_i) : k]
$$
\end{lemma}

\begin{proof}
(2) $\Rightarrow$ (1)은 자명하다. 다른 함의는 다음 사실에서 따른다.
$\Spec(k) \to Y$의 상이 $y$이면
$X_k = \Spec(k) \times_{\Spec(\kappa(y))} X_y$이다. 정의에 의해
$f$의 올들이 보편적으로 유계라는 것은 어떤 $n$이 존재한다는 뜻이다.
마지막으로 $X_k = \Spec(A)$라고 가정하자. 그러면
$\dim_k A = n$이다. 따라서 $A$는 아르틴이고, 모든 소아이디얼은 극대
아이디얼 $\mathfrak m_i$이며, $A$는 이 극대 아이디얼들에서의
국소화들의 곱이다. 대수학, 보조정리
\ref{algebra-lemma-finite-dimensional-algebra}와
\ref{algebra-lemma-artinian-finite-length}을 보라. 그러면
$\mathfrak m_i$는 $x_i$에 대응하고,
$A_{\mathfrak m_i} = \mathcal{O}_{X_k, x_i}$이며, 따라서 전사
$A \to  \bigoplus \kappa(\mathfrak m_i) = \bigoplus \kappa(x_i)$가
존재한다. 선형대수에 의해 이는 보조정리의 명제에 있는 부등식을
함의한다.
\end{proof}

\begin{lemma}
\label{morphisms-lemma-finite-locally-free-universally-bounded}
$f$가 차수 $d$의 유한 국소 자유 사상이면, $d$는 $f$의 올들의 차수에
대한 상계이다.
\end{lemma}

\begin{proof}
$f$의 모든 기저변환은 차수 $d$의 유한 국소 자유 사상이므로(보조정리
\ref{morphisms-lemma-base-change-finite-locally-free}), $f$의 올들은 모두 차수
$d$를 가진다. 따라서 주장이 성립한다.
\end{proof}

\begin{lemma}
\label{morphisms-lemma-composition-universally-bounded}
올들이 보편적으로 유계인 사상들의 합성은 올들이 보편적으로 유계인
사상이다. 더 정확히, $n$이 $f : X \to Y$의 올들의 차수에 대한
상계이고, $m$이 $g : Y \to Z$의 차수에 대한 상계라고 가정하자.
그러면 $nm$은 $g \circ f : X \to Z$의 올들의 차수에 대한 상계이다.
\end{lemma}

\begin{proof}
$f : X \to Y$와 $g : Y \to Z$의 올들이 보편적으로 유계라고 하자.
$\deg(X_y/\kappa(y)) \leq n$이 모든 $y \in Y$에 대해 성립하고,
$\deg(Y_z/\kappa(z)) \leq m$이 모든 $z \in Z$에 대해 성립한다고 하자.
점 $z \in Z$를 택하자. 가정에 의해 스킴 $Y_z$는
$\Spec(\kappa(z))$ 위에서 유한이다. 특히 $Y_z$의 바탕 위상공간은
유한 이산집합이다. 사상 $f_z : X_z \to Y_z$의 올들은 $f$의 올들이며,
이들은 $Y$의 대응하는 점들 위에 놓이고 위 논증에 의해 유한 이산집합이다.
따라서 $X_z$의 바탕 위상공간도 유한 이산집합이라고 결론 내린다.
그러므로 $X_z$는 아핀 스킴이다(이는 좋은 연습문제이며, 예를 들어
$X_z$의 모든 점들의 집합에 성질, 보조정리
\ref{properties-lemma-maximal-points-affine}를 적용해도 따른다).
$X_z = \Spec(A)$, $Y_z = \Spec(B)$, $k = \kappa(z)$라고 쓰자. 그러면
$k \to B \to A$이고, (a) $\dim_k(B) \leq m$이며, (b) 모든 극대
아이디얼 $\mathfrak m \subset B$에 대해
$\dim_{\kappa(\mathfrak m)}(A/\mathfrak mA) \leq n$임을 안다.
이로부터 $\dim_k(A) \leq nm$이 따른다고 주장한다. $B$는 자신의
국소화들 $B_{\mathfrak m}$의 곱임에 유의하라. 예를 들어 이는
$Y_z$가 $1$-점 스킴들의 분리합이기 때문이거나 대수학, 보조정리
\ref{algebra-lemma-finite-dimensional-algebra}와
\ref{algebra-lemma-artinian-finite-length}에서 따른다. 따라서
$\dim_k(B) = \sum_{\mathfrak m} \dim_k(B_{\mathfrak m})$이고
$\dim_k(A) = \sum_{\mathfrak m} \dim_k(A_{\mathfrak m})$이며, 두 경우
모두 $\mathfrak m$은 $B$의 극대 아이디얼들을 달린다($A$의 극대
아이디얼들이 아니다). 위의 사실과 나카야마 보조정리(대수학, 보조정리
\ref{algebra-lemma-NAK})에 의해 각 $A_{\mathfrak m}$은
$B_{\mathfrak m}^{\oplus n}$의 $B_{\mathfrak m}$-가군으로서의
몫이다. 따라서
$\dim_k(A_{\mathfrak m}) \leq n \dim_k(B_{\mathfrak m})$이다. 모두
합치면 다음을 얻는다.
$$
\dim_k(A) = \sum\nolimits_{\mathfrak m} \dim_ta	(A_{\mathfrak m})
\leq \sum\nolimits_{\mathfrak m} n \dim_k(B_{\mathfrak m})
= n \dim_k(B) \leq nm
$$
원하는 결론이다.
\end{proof}

\begin{lemma}
\label{morphisms-lemma-base-change-universally-bounded}
올들이 보편적으로 유계인 사상의 기저변환은 올들이 보편적으로 유계인
사상이다. 더 정확히, $n$이 $f : X \to Y$의 올들의 차수에 대한
상계이고 $Y' \to Y$가 임의의 사상이면, 기저변환
$f' : Y' \times_Y X \to Y'$의 올들의 차수도 $n$으로 제약된다.
\end{lemma}

\begin{proof}
이는 보조정리 \ref{morphisms-lemma-characterize-universally-bounded}의 결과에서
명백하다.
\end{proof}

\begin{lemma}
\label{morphisms-lemma-descent-universally-bounded}
$f : X \to Y$를 스킴들의 사상이라 하자. $Y' \to Y$를 스킴들의
사상이라 하고, $f' : X' = X_{Y'} \to Y'$를 $f$의 기저변환이라 하자.
$Y' \to Y$가 전사이고 $f'$의 올들이 보편적으로 유계이면, $f$의
올들도 보편적으로 유계이다. 더 정확히, $n$이 $f'$의 올들의 차수에
대한 상계이면 $n$은 $f$의 올들의 차수에 대한 상계이기도 하다.
\end{lemma}

\begin{proof}
$n \geq 0$을 $f'$의 올들의 차수에 대한 상계인 정수라 하자. 이 $n$은
$f$에 대해서도 쓸 수 있다고 주장한다. 실제로 점 $y \in Y$가
주어지면 점 $y' \in Y'$를 택하되 $y$ 위에 놓이게 하고 다음을 관찰하자.
$$
X'_{y'} = \Spec(\kappa(y')) \times_{\Spec(\kappa(y))} X_y.
$$
$X'_{y'}$가 가정에 의해 차수 $\leq n$이고 $\kappa(y')$ 위에서 유한인
스킴이므로, $X_y$도 차수 $\leq n$이고 $\kappa(y)$ 위에서 유한인
스킴임이 따른다. (몇몇 세부사항은 생략한다.)
\end{proof}

\begin{lemma}
\label{morphisms-lemma-immersion-universally-bounded}
몰입의 올들은 보편적으로 유계이다.
\end{lemma}

\begin{proof}
정수 $n = 1$을 정의에서 사용할 수 있다.
\end{proof}

\begin{lemma}
\label{morphisms-lemma-etale-universally-bounded}
$f : X \to Y$를 스킴들의 에탈 사상이라 하자. $n \geq 0$이라 하자.
다음은 동치이다.
\begin{enumerate}
\item 정수 $n$은 올들의 차수에 대한 상계이다.
\item 모든 체 $k$와 사상 $\Spec(k) \to Y$에 대해 기저변환
$X_k = \Spec(k) \times_Y X$는 많아야 $n$개의 점을 가진다.
\item 모든 $y \in Y$와 모든 분리가능 대수적 폐포
$\kappa(y) \subset \kappa(y)^{sep}$에 대해 스킴
$X_{\kappa(y)^{sep}}$은 많아야 $n$개의 점을 가진다.
\end{enumerate}
\end{lemma}

\begin{proof}
이는 보조정리 \ref{morphisms-lemma-characterize-universally-bounded}와 올들 $X_y$가
$\kappa(y)$의 유한 분리가능 체확장들의 스펙트럼들의 분리합이라는
사실에서 따른다. 보조정리 \ref{morphisms-lemma-etale-over-field}를 보라.
\end{proof}

\noindent
올들이 보편적으로 유계라는 것은 절대적 개념이지 상대적 개념이 아니다.
이 때문에 다음 보조정리의 조건은 $X$가 준콤팩트라는 것이며,
$f$가 준콤팩트라는 것이 아니다.

\begin{lemma}
\label{morphisms-lemma-locally-quasi-finite-qc-source-universally-bounded}
$f : X \to Y$를 스킴들의 사상이라 하자. 다음을 가정하자.
\begin{enumerate}
\item $f$는 국소 준유한이다.
\item $X$는 준콤팩트이다.
\end{enumerate}
그러면 $f$의 올들은 보편적으로 유계이다.
\end{lemma}

\begin{proof}
$X$가 준콤팩트이므로 유한 아핀 열린덮개
$X = \bigcup_{i = 1, \ldots, n} U_i$와 아핀 열린집합들
$V_i \subset Y$, $i = 1, \ldots, n$을 택하여
$f(U_i) \subset V_i$가 되게 할 수 있다. ``국소 준유한성''의 국소적
성격(보조정리 \ref{morphisms-lemma-quasi-finite-at-point-characterize}의 (4)를
보라)에 의해 사상들 $f|_{U_i} : U_i \to V_i$는 아핀 사이의 국소
준유한 사상이고, 따라서 준유한이다. 보조정리
\ref{morphisms-lemma-quasi-finite-locally-quasi-compact}를 보라. 모든
$y \in Y$에 대해 $X_y = \bigcup_{y \in V_i} (U_i)_y$가 열린덮개임은
명백하다. 따라서 아핀 사이의 준유한 사상에 대해 보조정리를 증명하면
충분하다(즉 $n_i$를 사상 $f|_{U_i} : U_i \to V_i$에 사용할 수 있으면
$\sum n_i$를 $f$에 사용할 수 있다).

\medskip\noindent
$f : X \to Y$를 아핀 사이의 준유한 사상이라 가정하자. 보조정리
\ref{morphisms-lemma-quasi-finite-affine}에 의해 다음 그림을 얻을 수 있다.
$$
\xymatrix{
X \ar[rd]_f \ar[rr]_j & & Z \ar[ld]^\pi \\
& Y &
}
$$
여기서 $Z$는 아핀이고, $\pi$는 유한이며, $j$는 열린 몰입이다. $j$의
올들은 보편적으로 유계이므로(보조정리
\ref{morphisms-lemma-immersion-universally-bounded}), $\pi$의 올들이 보편적으로
유계임을 보이는 것으로 환원된다(보조정리
\ref{morphisms-lemma-composition-universally-bounded}).

\medskip\noindent
이로써 형태 $\Spec(B) \to \Spec(A)$의 사상으로 환원되며, 여기서
$A \to B$는 유한이다. $B$가 $x_1, \ldots, x_n$에 의해 $A$-가군으로서 생성된다고
하자. 그러면
$$
A^{\oplus n} \longrightarrow B, \quad
(a_1, \ldots, a_n) \longmapsto \sum a_i x_i
$$
는 전사 $A$-가군 사상이다. 모든 소아이디얼 $\mathfrak p \subset A$에
대해 이는 $\kappa(\mathfrak p)$-벡터공간들의 전사 사상
$$
\kappa(\mathfrak p)^{\oplus n} \longrightarrow B \otimes_A \kappa(\mathfrak p)
$$
을 유도한다. 달리 말하면 정수 $n$을 올들이 보편적으로 유계인 사상의
정의에서 사용할 수 있다.
\end{proof}

\begin{lemma}
\label{morphisms-lemma-universally-bounded-permanence}
다음 스킴 사상들의 가환 그림을 생각하자.
$$
\xymatrix{
X \ar[rd]_g \ar[rr]_f & & Y \ar[ld]^h \\
& Z &
}
$$
$g$의 올들이 보편적으로 유계이고 $f$가 전사이며 평탄이면, $h$의
올들도 보편적으로 유계이다. 더 정확히, $n$이 $g$의 올들의 차수에
대한 상계이면 $n$은 $h$의 올들의 차수에 대한 상계이기도 하다.
\end{lemma}

\begin{proof}
$g$의 올들이 보편적으로 유계이고 $f$가 전사이며 평탄이라고 가정하자.
$g$의 올들의 차수가 $n \in \mathbf{N}$으로 제약된다고 하자. $n$을
$h$에도 사용할 수 있다고 주장한다. $z \in Z$를 택하자. 스킴 사상
$X_z \to Y_z$를 생각하자. 이는 평탄하고 전사이다. 가정에 의해
$X_z$는 $\kappa(z)$ 위의 유한 스킴이고, 특히 아르틴환의 스펙트럼이다
(대수학, 보조정리 \ref{algebra-lemma-finite-dimensional-algebra}에 의한다).
보조정리 \ref{morphisms-lemma-Artinian-affine}에 의해 사상 $X_z \to Y_z$는
아핀이고, 특히 준콤팩트이다. 보조정리
\ref{morphisms-lemma-fpqc-quotient-topology}에 의해 $Y_z$는 유한 이산이다. 이는
$X_z$에 대해 성립하기 때문이다. 따라서 $Y_z$는 아핀 스킴이다(이는 좋은
연습문제이며, 예를 들어 $Y_z$의 모든 점들의 집합에 성질, 보조정리
\ref{properties-lemma-maximal-points-affine}를 적용해도 따른다).
$Y_z = \Spec(B)$와 $X_z = \Spec(A)$로 쓰자. 그러면 $A$는 $B$ 위에서
충실 평탄이므로 $B \subset A$이다. 따라서 원하는 대로
$\dim_k(B) \leq \dim_k(A) \leq n$이다.
\end{proof}





\section{기타 결과}
\label{morphisms-section-misc}

\noindent
다른 곳에 들어맞지 않는 결과들을 모은다.

\begin{lemma}
\label{morphisms-lemma-factor-through-point}
$f : Y \to X$를 스킴들의 사상이라 하고 $x \in X$를 한 점이라 하자.
$Y$가 축약이고 $f(Y)$가 집합론적으로 $\{x\}$에 포함된다고 가정하자.
그러면 $f$는 표준 사상 $x = \Spec(\kappa(x)) \to X$를 통해 분해된다.
\end{lemma}

\begin{proof}
생략한다. 힌트: 아핀에서 국소적으로 작업하면 가환대수 보조정리로
환원된다. 환 사상 $A \to B$가 주어지고 $B$가 축약이며, 유일한
소아이디얼 $\mathfrak p \subset A$가 존재하여
$\Spec(B) \to \Spec(A)$의 상에 속한다면, $A \to B$는
$\kappa(\mathfrak p)$를 통해 분해된다. 이는 좋은 연습문제이다.
\end{proof}

\begin{lemma}
\label{morphisms-lemma-factor-through-set-of-points}
$f : Y \to X$를 스킴들의 사상이라 하고 $E \subset X$라 하자.
$X$가 국소 뇌터이고, $E$의 원소들 사이에 자명하지 않은 특수화가
없으며, $Y$가 축약이고, $f(Y) \subset E$라고 가정하자. 그러면
$f$는 $\coprod_{x \in E} x \to X$를 통해 분해된다.
\end{lemma}

\begin{proof}
$E$가 한원소집합이면 보조정리 \ref{morphisms-lemma-factor-through-point}에서
따른다. $E$가 유한이면 특수화에 관한 가정에 의해 $E$는 ($X$에서
유도된 위상을 갖춘) 유한 이산공간이다. 따라서 이 경우는 한원소집합의
경우로 환원된다. 일반적인 경우에는 $X$와 $Y$가 아핀 스킴인 경우로
환원된다. $f : Y \to X$가 환 사상 $\varphi : A \to B$에 대응한다고
하자. $A' \subset B$를 $\varphi$의 상이라 하자.
$E' \subset \Spec(A') \subset \Spec(A)$를 $A'$의 극소 소아이디얼들의
집합이라 하자. 대수학, 보조정리
\ref{algebra-lemma-injective-minimal-primes-in-image}에 의해 집합 $E'$는
$\Spec(B) \to \Spec(A') \subset \Spec(A)$의 상에 포함된다. 따라서
$E' \subset E$라고 결론 내린다. $A'$가 뇌터이므로 대수학, 보조정리
\ref{algebra-lemma-Noetherian-irreducible-components}에 의해 $E'$는
유한이다. $\Spec(B) \to \Spec(A)$의 상에 있는 다른 모든 점은
$E'$의 한 원소의 특수화이고 $E$에도 속하므로, 그 상이 $E'$에 포함된다고
결론 내린다($E$의 점들 사이의 특수화에 관한 가정에 의한다). 따라서
위에서 다룬 $E$가 유한인 경우로 환원된다.
\end{proof}
